\documentclass[12pt,a4paper]{ctexbook}
\usepackage{amsmath,amssymb}
\usepackage{esint}
\usepackage{geometry}
\usepackage{booktabs}
\usepackage{array}
\usepackage{longtable}
\usepackage{enumitem}
\usepackage{xcolor}
\usepackage{fancyvrb}
\usepackage{hyperref}
\usepackage{textcomp}
\usepackage{framed}
\usepackage{booktabs}

\geometry{margin=2.5cm}
\setlength{\parindent}{0pt}
\setlength{\parskip}{0.8em}
\setlist[itemize]{leftmargin=2em}
\setlist[enumerate]{leftmargin=2.5em}

\hypersetup{
    colorlinks=true,
    linkcolor=black,
    urlcolor=blue
}

\title{\Huge\textbf{让中国人学“好数学”}}
\author{Constant-Dalink\\康斯坦特-达琳克}
\date{}

\begin{document}

\maketitle
\clearpage

\begin{titlepage}
\centering
{\Huge\bfseries 编校与制作名单\par}
\vspace{2cm}

\renewcommand{\arraystretch}{1.6}
\begin{tabular}{@{}ll@{}}
作者：& Constant-Dalink，A-Slantou-Nidaloufu \\
制作人：& Constant-Dalink \\
主编：& Constant-Dalink \\
编辑：& Constant-Dalink \\
责任编辑：& Constant-Dalink \\
校对：& Constant-Dalink \\
排版：& Constant-Dalink \\
LaTeX 整理：& Constant-Dalink \\
封面设计：& Constant-Dalink \\
技术支持：& A-Slantou-Nidaloufu \\
版本统筹：& Constant-Dalink \\
鸣谢：& A-Slantou-Nidaloufu \\
\end{tabular}

\vfill
\end{titlepage}

\frontmatter
\maketitle

\chapter*{《让中国人学“好数学”》}
\addcontentsline{toc}{chapter}{《让中国人学“好数学”》}

\section*{开篇：一封写给中国教育的信件}
\addcontentsline{toc}{section}{开篇：一封写给中国教育的信件}

\subsection*{致所有人。}

我要说一些不好听的话。

如果你准备好了，就往下读。如果没准备好——也往下读。因为这些话，已经晚了几十年。

\subsection*{一}

中国有一个奇观。

十四亿人口，几乎人人学过数学。从六岁到十八岁，整整十二年。有些人还要再加上四年大学，两年硕士，三年博士。

算下来，这个国家每年有上亿人坐在教室里，翻开数学课本。

然后呢？

然后，这个国家的大多数成年人，离开考场之后，\textbf{这辈子再也没有主动想过一个数学问题。}

不是因为用不到。

而是因为他们\textbf{恨它}。

一个国家，用十二年的时间，让上亿人学会了恨一样东西——你说这是教育？

我说，这是\textbf{一场谋杀}。

杀死的是好奇心。

\subsection*{二}

凶手是谁？我来一个一个点名。

\textbf{第一个，教育部。}

他们制定了大纲。大纲里写着每个学期要教什么，教多快，考什么。这份大纲的目标，从第一行到最后一行，没有一个字是“让学生理解数学”。它的目标只有一个——\textbf{筛选}。

数学，在这套系统里，不是一门学问。它是一把筛子。用来把人分成三六九等，分配到不同的学校、不同的城市、不同的命运。

至于你在这个过程中有没有真正学会什么，有没有感受到数学之美——

\textbf{对不起，大纲里没有这一栏。}

\textbf{第二个，高考。}

高考是什么？高考是一场全国统一的、标准化的、限时的测试。它要求你在一百二十分钟之内，完成二十几道题，每一道都有标准答案，每一步都有采分点。

你知道这意味着什么吗？

这意味着，\textbf{你没有时间思考。}

在高考的考场上，思考是一种奢侈，甚至是一种惩罚。你想多了，时间就没了。你用自己的方法解题，步骤不符合评分标准，分就没了。

高考不是在考你懂不懂数学。

高考是在考你——\textbf{能不能在不思考的情况下，把正确答案写出来。}

你品品这句话。

这不是考数学。这是考驯化。

\textbf{第三个，课本。}

翻开任何一本中国的数学教材。你会看到一个非常经典的结构：

\begin{quote}
定义 $\to$ 定理 $\to$ 例题 $\to$ 练习
\end{quote}

先告诉你结论，再给你示范一遍怎么用，然后让你反复练。

这个结构，从小学到大学，一以贯之。

问题在哪？

\textbf{问题在于，它把数学最重要的部分砍掉了。}

一个定理，是怎么来的？什么人，在什么问题面前困惑了，冥思苦想了多久，走了多少弯路，才在某个瞬间发现了这条道理？

课本不告诉你。

课本只告诉你：这是结论，记住它，会用就行。

这就好比有人带你去看一座大教堂，但他只让你看了最顶上那个十字架，然后说：“好了，你已经参观完了。”

那些柱子呢？那些拱顶呢？那些工匠一砖一瓦垒了几十年的过程呢？

\textbf{没了。全砍了。课本只给你结果，不给你过程。}

而数学，恰恰是一门\textbf{过程比结果重要一万倍}的学问。

\textbf{第四个，课本的语言。}

这一条我要单独拎出来讲，因为它是一种更隐蔽的暴力。

你有没有仔细读过数学课本上的文字？

那些定义，那些证明，那些所谓的“推导过程”——它们写得有多\textbf{绕}？

明明一句话可以讲清楚的道理，课本偏偏要绕三个弯。明明一个概念的本质极其朴素，课本偏偏要用最抽象、最干燥、最让人读不下去的方式去表达。

你以为这是“严谨”？

不。\textbf{严谨和晦涩是两回事。}

真正的严谨，是把一件事讲得既准确又清楚。而中国的数学课本做到了前者，\textbf{故意放弃了后者}。

我说“故意”，是经过思考的。

你想想：一本教材，它的读者是十几岁的孩子。它的任务是让这些孩子理解数学。那它为什么不好好说话？为什么不用人类的语言、用正常的逻辑、用一个孩子能听懂的方式去讲？

\textbf{因为如果讲清楚了，人人都能懂。人人都能懂，就没法筛选了。}

课本的弯弯绕绕，本质上是一道\textbf{门槛}。

跨过去的人，被称为“聪明”。跨不过去的人，被判定为“学不了数学”。

但真相是什么？真相是那道门槛本不该存在。是有人故意把它放在那里的。

这像什么？

这像中世纪的教会，把《圣经》只用拉丁文书写，不许翻译成老百姓能读懂的语言。谁能读拉丁文，谁就能接近上帝。读不了的，就老老实实做你的平民。

\textbf{知识成了特权。门槛成了统治工具。}

数学课本的晦涩，干的是同样的事。它把一门本应属于所有人的学问，变成了少数人的专利。它暗示你：数学是高高在上的，是“普通人”够不到的，只有那些天生的“奇才”才配攀上这座高峰。

\textbf{这是什么？这是学术上的人粹主义。}

数学不是贵族的花园。数学是人类共同的空气。

\textbf{任何人都有权呼吸它。前提是，别有人故意把窗户焊死。}

\textbf{第五个，教辅产业。}

这是一门生意。一门巨大的、暴利的、建立在恐惧之上的生意。

《五年高考三年模拟》《黄冈密卷》《考点必刷题》——名字一个比一个凶狠，封面一个比一个紧迫。

它们卖的是什么？

是\textbf{套路}。

每一种题型，总结出三到五种解法。每一种解法，拆成固定的步骤。学生要做的，就是反复训练，直到形成肌肉记忆——看到题目，不用想，手比脑子先动。

这不是学习。

\textbf{这是驯兽。}

教辅产业不需要你理解数学。它需要你\textbf{依赖}它。你越不理解，就越需要它的套路。你越需要它的套路，它就卖得越多。

这是一条完美的产业链。唯一的代价，是你的思考能力。

\subsection*{三}

但你以为这就完了吗？

以上这些凶手，都在做同一件事——\textbf{杀}。杀好奇心，杀思考力，杀一个孩子对数学最本能的感觉。

可杀完之后呢？

\textbf{没有人来救。}

中国的教育体系，有一套世界上最庞大的考试机器，却没有一个像样的\textbf{人才发现机制}。

古人尚且知道“举荐”。一个地方官，发现当地有个少年天资过人，可以写一封举荐信，送他去见更好的老师，给他更大的天地。老师和老师之间，师门和师门之间，有一条活的、流动的人才脉络。

\textbf{千里马常有，而伯乐不常有——但至少，那时候的制度允许伯乐存在。}

今天呢？

今天的制度里，\textbf{伯乐是违规的。}

一个老师，如果发现班上有个孩子对数学有天赋，有灵气，有那种罕见的直觉——他能做什么？

什么都做不了。

他不能给这个孩子单独出题，因为大纲不允许。
他不能跳过课本去教更深的东西，因为进度不允许。
他不能向任何机构举荐这个孩子，因为\textbf{根本没有这样的机构}。
他唯一能做的，就是对这个孩子说：“你很聪明，好好考试。”

然后这个孩子就和其他所有孩子一起，被塞进同一条流水线，做同样的题，考同样的试，按同一个标准被打分。

他的天赋？

\textbf{对不起，天赋不在评分标准里。}

没有特别行政小组去主动寻找这样的孩子。没有一套机制让老师可以说：“这个学生不一样，他需要不一样的路。”没有人把“发现人才”当作一种\textbf{制度性的责任}。

我们只有考试。

考试是什么？考试是被动的。它坐在那里等你来证明自己。你证明不了，就是你的问题。

\textbf{但一个十岁的孩子，他怎么知道自己要去证明什么？他甚至不知道数学真正的样子是什么。}

伯乐的责任，是\textbf{走出去找}。走到田间地头，走到教室角落，走到那个趴在课桌上的孩子面前，蹲下来问他：你在想什么？

说到这里，我要说一件所有老师都见过，但几乎没有人认真想过的事——

\textbf{那些上课睡觉的孩子。}

每个班都有。趴在最后一排，脑袋埋在胳膊里，整堂课一动不动。

老师怎么看他们？“不爱学习。”“态度有问题。”“自己放弃自己了。”

家长怎么说？“这孩子就是懒。”“不争气。”“白供他上学了。”

\textbf{但有没有人想过一个问题：他为什么睡？}

他是从一年级第一堂课就开始睡的吗？

不是。

\textbf{没有任何一个孩子，在六岁走进教室的第一天，就趴下睡觉。}

每一个孩子刚入学的时候，眼睛都是亮的。他们对一切都好奇，什么都想知道，什么都想问。

那后来呢？

后来，他发现老师讲的东西他早就懂了——\textbf{但老师还在讲。}
后来，他发现他想问的问题不在考试范围内——\textbf{没人回答他。}
后来，他发现课堂的节奏不是为他设计的，他要么太快，要么太慢——\textbf{但没有人在乎他的节奏。}
后来，他发现不管他懂不懂，课都会照常进行，不管他在不在，分数都会照常公布——

\textbf{他的存在，对这个系统来说，无关紧要。}

于是他睡了。

不是因为懒。\textbf{是因为这个系统告诉他：你醒着也没有用。}

一个孩子的入睡，是对一个失败系统的\textbf{沉默抗议}。

可我们非但听不见这个抗议，还要反过来惩罚他。记过、请家长、罚站、当众批评——\textbf{系统杀死了他的兴趣，然后惩罚他对兴趣的死亡表现出的症状。}

这不是教育。

\textbf{这是施暴者指责受害者不够坚强。}

这个国家不缺千里马。

\textbf{这个国家缺的是蹲下来的人。}

\subsection*{四}

还有更深一层的问题。

我们的教育，不教真东西。

什么是真东西？

\textbf{数学的真东西，是思维方式。}

是面对一个从未见过的问题时，你怎么去拆解它、去试探它、去和它对话。是你在黑暗中摸索的能力。是你犯错、发现自己犯了错、然后从错误中修正方向的能力。

这些，课本不教。考试不考。

我们教的是什么？\textbf{是死东西。} 公式、步骤、模板、套路。像标本一样，钉在墙上，看起来栩栩如生，但早已没有呼吸。

然而更可悲的是——

\textbf{我们的教育，也不教做人。}

你以为数学只是数学吗？

一个人学数学的过程中，本应该学会这些东西：

\textbf{诚实。} 推理错了就是错了，不能糊弄，不能自欺。数学是世界上最诚实的学问，它不认权威，不认面子，只认逻辑。

\textbf{谦卑。} 面对一个想不通的问题，你得承认自己的无知，放下骄傲，从头再来。

\textbf{勇气。} 所有人都说“答案是A”，但你的推理告诉你是B——你敢不敢站出来？

\textbf{耐心。} 一个问题想一天想不通，想一周想不通，想一个月——你还愿不愿意继续想？

这些品质，才是数学教给人的最珍贵的东西。

但我们的课堂教了什么？

\textbf{服从。} 老师说的方法就是对的方法。
\textbf{速度。} 想得快才是聪明，想得慢就是笨。
\textbf{恐惧。} 错了就扣分，扣分就落后，落后就完了。

一个孩子在这套系统里待了十二年，他学会的不是思考，而是\textbf{听话}。不是追问，而是\textbf{闭嘴}。不是面对未知的勇气，而是\textbf{对犯错的恐惧}。

你告诉我，这叫教育？

这叫\textbf{驯化}。

教育是让一个人成为他自己。驯化是让所有人成为同一个人。

\subsection*{五}

\textbf{这个系统不给年轻人自由。}

一个中国孩子，从六岁进入小学，到十八岁高考结束——整整十二年——他的人生被安排得密不透风。

早上六点半到校，晚上九点半回家。周末补课，寒暑假预习下学期。每一个小时都被切割成碎片，填满了课程、作业、考试、排名。

\textbf{他没有一分钟是属于自己的。}

他不能说“这个问题我想多花三天想想”——因为明天就要考。
他不能说“课本上这个定理我觉得有问题”——因为考试就考这个。
他不能说“我对数论比几何更感兴趣，我想深入研究”——因为大纲不允许偏科。

一个孩子的数学直觉，往往在十岁到十六岁之间最敏锐、最活跃。那是大脑发育的黄金窗口，是想象力最奔放的年纪，是一个人最有可能对抽象世界产生原始热情的时期。

而这段时间，我们在让他干什么？

\textbf{做卷子。}

等到十八岁，高考结束了，进了大学——恭喜你，你终于“自由”了。

但那时候，你的好奇心已经被磨灭了。你的思考习惯已经被固化了。你对数学的感觉，要么是厌恶，要么是麻木。

\textbf{自由来了，人已经不会飞了。}

这就是中国数学教育最残忍的地方：

\textbf{它在你最能飞的年纪，把你的翅膀绑住。等到终于松绑的那一天，你已经忘记了自己有翅膀。}

你说这像什么？

这像一个人，从小被关在一间没有窗户的屋子里。有人每天从门缝塞进来几张纸，纸上写着“外面的世界是什么样的”。他背下来了，考试也考了很高的分。

然后有一天，门开了。

他走出去，阳光刺得他睁不开眼。他看着天空，看着大地，看着无边无际的世界——

\textbf{他不知道该往哪走。}

因为那十二年里，没有人教过他怎么走路。只教过他怎么背地图。

\subsection*{六}

我知道有人会说：

\begin{quote}
“可是这套系统出了成果啊。中国学生在国际数学竞赛上拿了那么多金牌，中国的数学成绩在全世界排名都很靠前。”
\end{quote}

是的。

就像一个工厂，如果你把一百万块矿石扔进去，用最高的温度、最大的压力去碾压，你确实能从里面筛出几颗钻石。

但你有没有想过——

\textbf{被碾碎的那九十九万九千块石头呢？}

它们中间，有多少本来可以是玉？有多少本来可以是玛瑙？有多少本来有自己独特的纹路和光泽？

没有人知道。

因为它们已经被碾成了粉。它们被碾成粉的那一刻，标签上写着四个字：

\textbf{“数学不好。”}

而那几颗被选出来的钻石——你以为它们就幸运了？

它们被这套系统选中，不是因为有人真正看见了它们的光芒，而是因为它们恰好足够硬，扛住了碾压。它们的才华不是被培养出来的，\textbf{是侥幸没被杀死的}。

这不叫教育。

这叫\textbf{自然淘汰}。

连丛林法则都不如——丛林至少是自由的。

\subsection*{七}

现在，我想单独说说那几颗“钻石”。

如果你是那种每次考试都能拿高分、公式背得滚瓜烂熟的“学霸”——这本书，\textbf{尤其}是写给你的。

我知道你现在可能有点不服气：

\begin{quote}
“我成绩好，我怎么了？你凭什么说我更需要？”
\end{quote}

这个问题问得好。我来认真回答你。

你有没有想过——你的高分，有可能恰恰是你最大的障碍？

不是因为你不聪明。而是因为，\textbf{高分给了你一个最难以撼动的错觉：我已经学会了。}

成绩差的孩子，他们至少知道自己有问题。他们会怀疑，会迷茫，会寻找。这扇门，对他们来说是虚掩着的。

但你呢？

每一次高分，都在替你把那扇门，\textbf{钉得更紧一颗钉子}。

老师夸你，父母夸你，排名证明你，整个系统都在告诉你：“你是对的，你学得很好，继续这样下去。”

没有人告诉你：

\textbf{你学得很好，但你学的，可能不是数学。}

你掌握的是一套极其精良的“解题系统”。你能看出题目在考哪个知识点，能在最短时间内调出对应的公式和方法，能精准地把答案写在阅卷老师最想看到的位置。

这套系统，非常有用。用来考试，几乎无敌。

但数学不是考试。

当有一天，没有人给你出题了，没有选项A、B、C、D了，没有标准答案了——

你还能看见数学吗？

这才是我说“尤其是写给你”的原因。

\textbf{因为你是最难被说服的那个人。} 而只有你真正愿意放下那张成绩单，重新用一个初学者的眼睛去看数学，你才会发现——

你其实从来没有见过它真实的样子。

\subsection*{八}

写给孩子的话。

如果你还小，如果你是一个正在上学的孩子，我想单独对你说几句话。

你知道吗，人类历史上那些最伟大的数学家，他们小时候，很多都是“问题学生”。

他们不停地问“为什么”，让老师头疼。
他们拒绝接受“背下来就行了”的答案。
他们在课堂上走神，因为他们在想自己更感兴趣的问题。

这不是缺点。

这是\textbf{数学家的天性}。

但我必须诚实地告诉你一件残酷的事：\textbf{这个系统不会保护你的天性。} 它会用分数来惩罚你的好奇心，用排名来否定你的独立思考，用“不按标准方法解题”扣掉你的分。

更残酷的是——\textbf{它不会给你时间。}

你现在的年纪，十岁、十二岁、十五岁——恰恰是你思维最活跃、直觉最敏锐、最有可能真正爱上数学的时候。

但这套系统不会让你用这段时间去探索。它会让你用这段时间去\textbf{赶路}。赶进度，赶排名，赶着到达一个叫“高考”的终点。

它告诉你：“等你考完了，你就自由了。”

\textbf{这是一句谎话。}

因为自由不是别人给你的。自由是一种能力。而这种能力，需要在你年轻的时候就开始练。如果你从来没有自由地思考过一个数学问题，那么等到有一天别人告诉你“你自由了”——你只会站在原地，不知所措。

所以，如果你在课堂上曾经因为“问太多问题”被老师批评，如果你曾经因为“解题方法和书上不一样”被扣分，如果你曾经觉得“我明明想清楚了，但分数却很低”——

\textbf{错的不是你。}

\textbf{是那个系统。}

请你保护好那份“想清楚”的冲动。不要等到成年以后再去追回它——那时候追回来的，只是影子。

\textbf{现在就开始。趁你还有翅膀。}

\subsection*{九}

有一件讽刺的事，我不得不说。

在数学课堂的隔壁，有一门学科，早就在实践一种完全不同的逻辑。

\textbf{那就是编程。}

你写一段程序，拿去运行。计算机会怎么评判你？

它不会问你：“你用的方法是课本上的吗？”
它不会问你：“你的步骤符合标准格式吗？”
它不会问你：“你这个写法和老师教的一样吗？”

它只问一个问题：

\textbf{“能不能跑通？结果对不对？”}

仅此而已。

同一个问题，你可以用循环，可以用递归，可以用十行代码，也可以用一百行。你可以走一条没有任何人走过的路，用一种没有任何教科书写过的方法——只要最终结果是对的，逻辑是自洽的，\textbf{它就是对的}。

没有标准写法。没有唯一解法。没有“虽然答案对了但方法不对所以扣分”这种荒诞的事。

\textbf{编程的世界只认一条铁律：成立，就是成立。}

你知道这意味着什么吗？

这意味着，在编程的世界里，一个孩子如果用了一种老师从没见过的方法解决了问题——他会被夸奖，而不是被扣分。

这意味着，在编程的世界里，“和别人不一样”不是错误，\textbf{是创造力}。

这意味着，在编程的世界里，\textbf{过程的多样性和结果的正确性可以并存}。你不需要牺牲前者去保证后者。

反过来看我们的数学课堂——

一个学生用自己的方法证明了一道题，逻辑完整，结论正确。但因为“步骤和标准答案不符”，被扣了一半的分。

你说，这两个世界，哪一个更尊重思考？

编程从第一天就在告诉我们一个道理：

\begin{quote}
\textbf{重要的不是你怎么走的，重要的是你到没到，以及你知不知道自己为什么能到。}
\end{quote}

这个道理，数学课堂装作不知道。已经装了几十年了。

\subsection*{十}

还有一件事，我必须说。

我们的数学教育，有一种极其危险的倾向——

\textbf{它把数学从现实中连根拔起，悬在半空中教给你。}

所有的概念都是抽象的。所有的定义都是形式化的。所有的推理都发生在一个“纯粹的数学世界”里，和你脚下踩着的土地没有任何关系。

一个孩子学“概率”，做了一百道计算题，但他从来没有真正理解过“不确定性”在生活中意味着什么。

一个孩子学“函数”，画了一百条曲线，但他从来没有感受过一个变化的量是怎样牵动另一个变化的量的。

一个孩子学“几何”，背了一百条定理，但他从来没有抬头看过一片叶子的对称，从来没有蹲下来量过自己影子的长度。

\textbf{数学被教成了一门和现实无关的学问。}

但这不是数学的本来面目。

数学从哪里来的？从数手指头来的。从丈量土地来的。从仰望星空、预测潮汐、计算收成来的。每一个数学概念，在它诞生的那一刻，都是\textbf{扎根在现实中的}。

是后来的人，把它拔了出来，洗干净泥土，装在玻璃瓶里，放在高高的架子上，告诉你：“这是纯粹的科学，不要用脏手碰它。”

\textbf{这不是在保护数学。这是在杀死数学。}

数学推理当然要严谨。但严谨不意味着脱离现实。恰恰相反——\textbf{最好的数学推理，是和现实世界相呼应的。} 你在纸上推出的结论，应该能在真实世界中找到回响。如果找不到，那你要么推错了，要么你问的问题本身就和这个世界无关。

数学不是真空中的游戏。它是人类理解这个世界的工具。

\textbf{脱离了世界的数学，不是更纯粹了，是更死了。}

这不是我一个人的看法。以下这些名字，来自世界上最顶尖的数学家、教育家和认知科学家。他们来自不同的国家、不同的时代、不同的领域，但他们用各自的方式，说着同一句话：

\textbf{“我们教数学的方式，错了。”}

\subsection*{推荐阅读与参考文献}

\textbf{1. Paul Lockhart，《A Mathematician's Lament》（2002 / 2009）}

一位美国数学家写的檄文。他说：“如果我想要毁掉一个孩子对音乐的热爱，最好的办法就是让他花十二年时间抄乐谱、背音符，但永远不让他听到一首曲子。”然后他说：“这就是我们正在对数学做的事。”

\textbf{2. V.I. Arnold，《On Teaching Mathematics》（1997）}

苏联最伟大的数学家之一。他在巴黎的一次演讲中愤怒地批判了将数学教育过度形式化、脱离物理和几何直觉的趋势。他说：“数学是物理学的一个分支——在物理学中，实验成本很低。”他的意思是，数学的根永远扎在对现实世界的直觉里。

\textbf{3. Hans Freudenthal，《Revisiting Mathematics Education》（1991）}

荷兰数学家、教育家。他创立了“现实数学教育”理论，主张数学应该从学生能感知的现实情境出发，而不是从抽象定义开始。他的理论影响了荷兰、日本等国的数学教育改革，但几乎从未被中国的课程设计者认真对待。

\textbf{4. George Lakoff \& Rafael Núñez，《Where Mathematics Comes From》（2000）}

一位认知语言学家和一位认知科学家合著。他们用大量证据论证：数学不是从天上掉下来的抽象真理，它起源于人类的\textbf{身体经验}——对空间的感知、对数量的直觉、对运动的理解。脱离了身体经验去教数学，等于把根砍掉。

\textbf{5. William Thurston，《On Proof and Progress in Mathematics》（1994）}

菲尔兹奖得主。他提出：数学的核心不是证明，\textbf{是理解}。形式化的证明只是传递理解的工具之一，而且往往不是最好的工具。

\textbf{6. Seymour Papert，《Mindstorms: Children, Computers, and Powerful Ideas》（1980）}

MIT人工智能实验室创始人之一。他开发了Logo编程语言，让孩子通过编程来学数学。他的核心理念：孩子不应该被“教”数学，而应该\textbf{在做事的过程中自己发现数学}。你不需要遵循标准步骤，你只需要让程序跑通。

\textbf{7. Jo Boaler，《Mathematical Mindsets》（2015）}

斯坦福大学数学教育教授。她用大量实证研究证明：速度和数学能力无关，犯错是大脑学习的最佳时机，僵化的教学方法会系统性地摧毁学生的数学潜力。

\textbf{8. Conrad Wolfram，《The Math(s) Fix: An Education Blueprint for the AI Age》（2020）}

他提出了一个尖锐的问题：为什么我们还在让学生用手算？计算机时代，数学教育应该聚焦于\textbf{提出问题、建立模型、解读结果}，而不是手动演算。

\textbf{9. Keith Devlin，《Introduction to Mathematical Thinking》（2012）}

斯坦福大学数学家。他区分了“学校数学”和“真正的数学”，认为前者几乎和后者是两个物种。学校教的是计算和套路，真正的数学是思考和推理。

以上这些人指向同一个结论：

\textbf{数学教育应该扎根于现实，尊重直觉，允许多样性，鼓励理解而非背诵，衡量思考而非速度。}

这不是某一个人的偏见。

\textbf{这是全世界最懂数学的一群人的共识。}

而我们的系统，几十年来，一直在朝相反的方向走。

\subsection*{十一}

所以，这本书到底要做什么？

这本书要做的事情，说出来很简单，做起来很难——

\textbf{它要把你从那套系统里拔出来。}

不是修修补补。不是“在现有基础上改进方法”。不是教你一种“更高效的刷题技巧”。

是\textbf{连根拔起}。

你过去十几年学的那一套——背公式、套模板、看题型找方法、掐着秒表刷题——

\textbf{全部放下。}

我知道这听起来很极端。

它就是极端。

这本书是极端的。就像一个人在沙漠里走了很远很远，忽然有人告诉他：你走反了。

你觉得他应该怎么办？是继续往前走，还是转身？

转身的代价很大。意味着你之前走过的每一步都白走了。意味着你要承认，自己引以为傲的方向感，从一开始就是错的。

\textbf{但如果你不转身，你就永远到不了那个地方。}

数学真正的样子，在你身后。

\subsection*{十二}

最后的最后，一个小小的约定。

在我们正式开始之前，你只需要在心里记住一句话：

\begin{quote}
\textbf{“我愿意忘掉我以为我知道的一切，从零开始，重新看一看数学到底是什么。”}
\end{quote}

不是改进。是\textbf{推翻}。

不是优化。是\textbf{重来}。

只要你能做到这一点，这本书就能对你有用。

我不会承诺你读完之后成绩会提高。也许会，也许不会。那不是这本书关心的事。

我承诺你的是另一样东西——

\textbf{你会第一次看见数学。}

真正的数学。

不是课本里的，不是教辅里的，不是高考卷上的。

是人类几千年来，一个问题接着一个问题，一代人接着一代人，在黑暗中摸索出来的那条路。

那条路上没有标准答案。

但有光。

\begin{flushright}
\textit{写于中国，写给中国}

\textit{Constant-Dalink}

\textit{康斯坦特-达琳克}
\end{flushright}

\begin{quote}
\textbf{“一个民族的数学教育，不应该是一条流水线。它应该是一片土壤——让每一颗种子，按自己的方式，长出自己的形状。”}

\hfill --- Constant-Dalink
\end{quote}

\tableofcontents

\mainmatter

\part{小学}
\setcounter{chapter}{0}

\chapter{从手指开始}

在我们开始之前，我想先说一件事：

\textbf{数学不难。}

如果你曾经觉得难，那不是你的问题。可能是讲的人没讲清楚。也可能是讲的方式不适合你。

没关系。我们重新来。

从手指开始。

\section{加法}

伸出一根手指。

再伸出一根手指。

现在你有两根手指伸着。

我们把这件事写成：

\begin{quote}
\textbf{$1 + 1 = 2$}
\end{quote}

这里的“+”，意思是\textbf{放在一起}。

一根手指，和另一根手指，放在一起，是两根手指。

这就是加法。

（如果你觉得这也太简单了——好，请保持耐心。我们正在打地基。地基稳，楼才高。）

\section{乘法}

现在想象十个人站成一排。

每人伸出一根手指。

\begin{quote}
一共有多少根手指？
\end{quote}

你可以一个一个数：1, 2, 3, 4, 5, 6, 7, 8, 9, 10。

但有一种更简洁的写法：

\begin{quote}
\textbf{$10 \times 1 = 10$}
\end{quote}

意思是：\textbf{10个人，每人1根。}

如果每人伸出两根呢？

\begin{quote}
\textbf{$10 \times 2 = 20$}
\end{quote}

“$\times$”的本质是\textbf{重复的加法}。

10个人，每人2根，就是：

\begin{quote}
$2 + 2 + 2 + 2 + 2 + 2 + 2 + 2 + 2 + 2 = 20$
\end{quote}

写成乘法，干净利落。

（数学家都怕麻烦。能少写几个字，绝不多写。）

\section{减法}

你伸出两根手指。

现在收回一根。

还剩几根？

\begin{quote}
\textbf{$2 - 1 = 1$}
\end{quote}

“-”的意思是\textbf{拿走}。

原来有两根，拿走一根，还剩一根。

\section{除法}

你有十根手指。

要分给十个人，每人分得一样多。

每人几根？

\begin{quote}
\textbf{$10 \div 10 = 1$}
\end{quote}

“$\div$”的意思是\textbf{平均分配}。

十根手指，十个人分，每人一根。

\section{现在来点有意思的}

你已经知道：

\begin{quote}
$2 + 1 = 3$
\end{quote}

那么：

\begin{quote}
$2 + \underline{\hspace{1cm}} = 3$
\end{quote}

横线上是几？

我不说答案。你看着上面那行，自己想。

想到了？

好。

你刚才做的事情，有个名字。但我先不告诉你叫什么。

名字不重要。你\textbf{会做}，比知道它叫什么重要一万倍。

\section{练习时间}

你把我刚才藏的那道题猜到了？

咋这么棒呢？！

来，我们练练。

放心，这不是学校。没人给你打分。

你要真确定你会了，就一次性写全对！

敢不敢？

\bigskip

\textbf{加法：}

$1 + 1 = $ ？

$2 + 1 = $ ？

（不会？没事。2就是两个，1就是一个，合起来？用脑袋想想，别掰你那手指了，其实不用那么久。）

$3 + 1 = $ ？

$2 + 2 = $ ？

$3 + 2 = $ ？

（如果全对了，别骄傲。后面还有呢。）

\bigskip

\textbf{减法：}

$2 - 1 = $ ？

$3 - 1 = $ ？

（减法就是拿走。原来有3个，拿走1个，还剩？）

$3 - 2 = $ ？

$4 - 2 = $ ？

\bigskip

\textbf{藏起来的数：}

$1 + \underline{\hspace{1cm}} = 2$

$1 + \underline{\hspace{1cm}} = 3$

（1加上几等于2？1加上几等于3？就这么想。）

$2 + \underline{\hspace{1cm}} = 4$

$\underline{\hspace{1cm}} + 1 = 3$

（这道把未知数换到前面了。还会吗？会的话你真的可以。）

\bigskip

\textbf{最后一道：}

$\underline{\hspace{1cm}} + \underline{\hspace{1cm}} = 5$

这道答案不止一个。你能写出几种？

（写出三种以上，我请你吃糖。开玩笑的，我又不知道你在哪。）

\bigskip

\textbf{混合加减：}

那……这个呢？

$1 + 2 - 1 = $ ？

（嘿，这道题有两个符号。别慌，从左往右，一步一步来。）

$2 + 3 - 2 = $ ？

$3 - 1 + 2 = $ ？

（如果你做得飞快，那你真的很厉害。如果你做得慢，那也很厉害——说明你在认真想。）

$4 + 2 - 3 = $ ？

$5 - 3 + 2 = $ ？

（全对了？行，你可以。）

\bigskip

\textbf{整十加减：}

$10 + 10 = $ ？

（十个人的手指，加上另外十个人的手指。）

$20 + 10 = $ ？

$30 + 20 = $ ？

$50 - 20 = $ ？

$40 - 30 = $ ？

（整十的加减，其实就是把后面那个0先遮住，算完再放回去。不信你试试。）

\bigskip

\textbf{来道有意思的：}

$12 + 21 = $ ？

（这道稍微要动点脑筋了。提示：先算个位，再算十位。实在不行，掰手指我也不笑话你。）

\bigskip

如果你做到这里全对了——

不是你聪明。

是你本来就会。只是以前没人好好跟你讲。

\bigskip

\textbf{第一章完成 ✓}

\chapter{藏起来的数}

上一章你做了这道题：

\begin{quote}
$2 + \underline{\hspace{1cm}} = 3$
\end{quote}

那个横线，就是藏起来的数。

有些书管它叫“未知数”，写成 $x$ 或者 $a$。

听起来很厉害对吧？

其实就是\textbf{我不告诉你是几，你自己猜}。

就这么简单。

\bigskip

不过在你继续猜之前，有几件事得跟你说清楚：

\textbf{第一，同一道题里，同一个 x 必须是同一个数。}

$x + x = 6$，这两个 $x$ 得一样。你不能说一个是 2，一个是 4。那叫耍赖。

\textbf{第二，x 可以是任何数。}

1，100，0，小数，负数，都行。它就是一个空位，等着你填。

\textbf{第三，不同题里，x 可以代表不同的数。}

这道题 $x = 3$，下道题 $x = 7$，完全没问题。$x$ 只在这一道题里说了算。

\textbf{第四，名字只是个代号。}

$x$、$a$、$b$……甚至画个圈都行。
$x + x = 6$ 和 $a + a = 6$ 完全是一回事。
你想全换成 $a$？没问题。

\textbf{第五，那能不能只换一个？}

比如把 $x + x$ 里的一个 $x$，换成 $y$？
变成：$x + y = 6$。
\textbf{太能了！}
但这就有意思了——既然名字不一样，说明它们可能是\textbf{两个不同的数}。
这就是“双主角”。

\begin{quote}
\textbf{【如果你要考试，得再记个这个】}

虽然我觉得这事儿显而易见，但为了防止你遇上个死板的出题人，还是得交代一句。
（但我认为只要出题人不傻，不会让你浪费时间在写没用的东西上。）

\textbf{教科书上怎么说：}

“字母可以表示任意的数。在同一问题中，不同的字母通常表示不同的量。”

\textbf{说人话就是：}

不同的名字，代表不同的角色。

$x$ 演 $x$ 的戏，$y$ 演 $y$ 的戏。

\textbf{可能}演同一个数，也\textbf{可能}不一样。
\end{quote}

好，规矩讲完了。来猜题：

$3 + x = 5$

（$x$ 是几？就是 3 加上几等于 5。）

$x + 4 = 7$

$x - 2 = 1$

（原来有几个，拿走 2 个，剩 1 个？）

$10 - x = 3$

（10 拿走几个，剩 3 个？）

\bigskip

不是哥们！真的假的，写完了？！

那……是时候带你接触一下大城市的东西了。

\bigskip

$x + x = 6$

（两个一样的数加起来等于 6。那这个数是几？）

$x + x + x = 9$

（三个一样的呢？）

$x + y = 5$

（这道答案不止一种。$x$ 和 $y$ 可以是什么？写几组出来。）

\bigskip

\textbf{就这样。}

未知数，就是藏起来的数。

会猜了，你就会了。

\bigskip

\textbf{第二章完成 ✓}

\chapter{一张合法的“作弊小抄”}

我知道，我知道。

看到这个表，你头都大了。
你可能在学校被老师逼着背过，被家长盯着默写过。
我也烦。背书最烦了。

但在你把它扔掉之前，听我说一句：

\textbf{这张表，不是为了让你受罪的。}
\textbf{它是为了让你以后能“偷懒”的。}

\section{为什么要背这玩意儿？}

我们来算个帐。

如果我不给你这张表，让你算 $7 \times 8$。
你会怎么办？
回顾第一章，乘法就是“重复的加法”。
那你得算：$8+8+8+8+8+8+8$。

\begin{quote}
$8+8=16$

$16+8=24$

$24+8=32$

……
\end{quote}

烦不烦？累不累？
等你算出来 $56$，隔壁的小明已经做完十道题出去玩了。

所以，这张表存在的意义只有一个：
\textbf{为了让你把宝贵的脑子，留给真正难的事情。}

哪怕你是天才，脑力也是有限的。
别把脑力浪费在“算数”上，要留给“思考”。

\section{乘法口诀表（九九乘法表）}

来，看着它。它就是个工具箱。
就像木匠不用每次都去造锤子，直接拿来用就行。

\begin{quote}
\textbf{一一得一}

\textbf{一二得二　二二得四}

\textbf{一三得三　二三得六　三三得九}

\textbf{一四得四　二四得八　三四十二　四四十六}

\textbf{一五得五　二五一十　三五十五　四五二十　五五二十五}

\textbf{一六得六　二六十二　三六十八　四六二十四　五六三十　六六三十六}

\textbf{一七得七　二七十四　三七二十一　四七二十八　五七三十五　六七四十二　七七四十九}

\textbf{一八得八　二八十六　三八二十四　四八三十二　五八四十　六八四十八　七八五十六　八八六十四}

\textbf{一九得九　二九十八　三九二十七　四九三十六　五九四十五　六九五十四　七九六十三　八九七十二　九九八十一}
\end{quote}

\section{怎么搞定它？}

别死记硬背。那是笨办法。
记住这几条，你的任务量瞬间减半：

\textbf{1. 倒过来是一样的}

还记得第一章的手指吗？
$2 \times 9$ 和 $9 \times 2$，结果是一样的。
所以你只需要背一半。

\textbf{2. 它是讲道理的}

忘了 $6 \times 7$ 是多少？
没事，你肯定记得 $6 \times 6 = 36$。
再加个 $6$，就是 $42$。
\textbf{口诀是死的，人是活的。}

\textbf{3. 它是你的武器}

把它背下来，装在脑子里。
下次遇到 $9 \times 9$，别人还在那儿画圈圈，你直接脱口而出“八十一”。
那种感觉，爽。

\section{最后的建议}

\textbf{不用今天全背下来。}

贴在床头，贴在厕所门上。
没事瞟一眼。
用着用着，你就记住了。

这就跟你打游戏记技能连招一样。
熟了，就是本能。

\bigskip

\textbf{第三章完成 ✓}

\chapter{数字的楼梯}

\section{为什么要发明这个？}

还记得第一章吗？
加法很麻烦，所以我们发明了乘法。
$2 + 2 + 2 + 2 + 2 = 10$

写成 $2 \times 5 = 10$。省事儿。

那要是这样呢？

$2 \times 2 \times 2 \times 2 \times 2 = ?$

（5个2相乘）

这回可不能写成 $2 \times 5$ 了。那是加法。
这回是\textbf{乘法叠罗汉}。

数学家又开始犯懒了。
写那么多 $2$ 干嘛？手酸。
于是他们发明了一个更狠的偷懒办法：

\begin{quote}
\textbf{$2^5$}
\end{quote}

那个小小的 $5$，就是告诉你：
\textbf{有 5 个 2 在互相乘。}

\section{这叫什么？}

这就叫\textbf{次方}。
也有人管它叫“幂”。听着挺高深，其实就是\textbf{乘法加强版}。

那个趴在底下的数（这里是 $2$），叫\textbf{底数}。
那个骑在肩膀上的小数字（这里是 $5$），叫\textbf{指数}。

\section{来几个例子}

$3^2$ 是多少？

意思是：$3 \times 3$。

算出来是 $9$。

（注意：\textbf{绝对不是 $3 \times 2 = 6$！千万别搞混！}）

$5^3$ 是多少？

意思是：$5 \times 5 \times 5$。

$5 \times 5 = 25$。

$25 \times 5 = 125$。

哇，一下子变这么大。

$10^4$ 是多少？

意思是：$10 \times 10 \times 10 \times 10$。

$10000$。

发现规律没？指数是几，后面就有几个 $0$。

\section{为什么说它是“楼梯”？}

你看看这个：

$2^1 = 2$

$2^2 = 4$

$2^3 = 8$

$2^4 = 16$

$2^5 = 32$

$2^{10} = 1024$

每上一层台阶，数字就翻倍。
越往上走，数字大得越吓人。
这就是次方的力量。

\section{练习一下}

别光看，动笔写几个。

$2^3 = ?$

（是 $2\times3$ 吗？如果你敢写6，我就顺着网线过去找你。是 $2\times2\times2$！）

$4^2 = ?$

$3^3 = ?$

$10^5 = ?$

（这道题不用算，直接数 $0$ 就行。）

\bigskip

\textbf{第四章完成 ✓}

\chapter{反求法——像福尔摩斯一样}

我们在第二章已经学会猜数了：

\begin{quote}
$x + 2 = 5 \implies x = 3$

（猜一猜：几加 $2$ 等于 $5$？）

$3x = 12 \implies x = 4$

（猜一猜：$3$ 乘几等于 $12$？）
\end{quote}

现在，来个刺激的：

\begin{quote}
$x^2 = 9$
\end{quote}

（还记得第四章吗？$x^2$ 就是 $x \times x$。）

这道题在问你：

\textbf{哪个数自己乘自己，等于 $9$？}

你会说：$3$ 啊！

$3 \times 3 = 9$。

没错。

\begin{quote}
$x^3 = 8$
\end{quote}

\textbf{哪个数连乘三次等于 $8$？}

$2 \times 2 \times 2 = 8$。

所以 $x = 2$。

\section{这就叫“反求法”}

\textbf{正着想：}

$3^2$ 是多少？

算出来是 $9$。

这是顺着来的，叫\textbf{乘方}。

\textbf{反着想：}

谁的平方是 $9$？

倒推回去找 $3$。

这就叫\textbf{开方}（或者叫求根）。

这就是数学家最喜欢的把戏：

\textbf{我知道结果，我要找源头。}

\section{为什么说这在生活里最厉害？}

你有没有玩过这个游戏：

妈妈把糖藏在左手或右手。
你只知道：\textbf{她现在手里有两颗糖。}
你想知道：\textbf{她原来藏了几颗？}

这就是反求法。

你只看到“结果”（现在2颗），
你想倒回去找出“原因”（原来藏了几颗）。

再来一个你肯定玩过：

你吹气球。
你只知道：\textbf{现在气球是原来的 4 倍大。}
你想知道：\textbf{你一共吹了几口同样的气？}

$2 \times 2 = 4$，因为你吹了两口。

$2 \times 2 \times 2 = 8$，因为你吹了三口。

你又在倒着想。

还有一次：

你有一块正方形蛋糕。
你量了一下，发现面积是 9 块小方块那么大。
你想知道：\textbf{每边切了几刀？}

$\to$ 就是 $x \times x = 9$，所以 $x = 3$，就是3刀。

你每天都在用反求法：
\begin{itemize}
\item 找零钱（我给了 10 块，找回 3 块，我买的东西是几块？）
\item 猜年龄（我比你大 2 岁，我 9 岁，你几岁？）
\item 玩捉迷藏（我听到脚步声在左边，原来你躲在左边！）
\end{itemize}

\textbf{反求法就是：}

\textbf{看到结果，想倒回去找出原因。}

全世界最聪明的人，
和最小的孩子，
用的都是同一个本事。

\textbf{你已经会了。}

\section{来，试几道“曲解题”}

别被吓到。这名字是我编的。

意思就是：\textbf{脑筋急转弯}。

1. \textbf{谁的平方是 16？}

$x^2 = 16$

（提示：想想 $4 \times 4$。）

2. \textbf{谁的立方是 27？}

$x^3 = 27$

（提示：试试 $3$。）

3. \textbf{谁的平方是 1？}

$x^2 = 1$

（这也太简单了吧？）

4. \textbf{谁的平方是 0？}

$x^2 = 0$

（想想什么数乘什么数还是 $0$？）

\section{最后来个大的}

$x^2 = 100$

这道题你要是想出来了，恭喜你。

你已经掌握了最厉害的思维方式：

\textbf{从结果倒推原因。}

（其实……$x^2 = 9$ 还有一个答案。先保密。以后告诉你。）

\bigskip

\textbf{第五章完成 ✓}

\chapter{别废话，直接来}

你刚才是不是在想：

“别给我讲故事了。别给我讲苹果、香蕉、小明跑得有多快。我就想做点那个酷酷的代数题。”

\textbf{满足你。}

这一章，没有应用题。

没有文字坑你。

只有数，符号，和你。

就看你能不能把它们驯服。

\section{第一关：热身（别掉以轻心）}

这些看起来像第一章的题。但小心，里面混进了新的东西。

\textbf{1.} $10 + 10 - 5 = \_\_$

（别眨眼，直接写。）

\textbf{2.} $3 \times 3 + 1 = \_\_$

（先乘还是先加？当然是先算高级的乘法。别急，后面我会告诉你为什么要这样。）

\textbf{3.} $20 \div 2 - 5 = \_\_$

\textbf{4.} $0 + 100 = \_\_$

\textbf{5.} $100 \times 0 = \_\_$

（嘿，这道题如果你算了一秒钟以上，你就输了。任何东西乘0都是……？）

\section{第二关：字母满天飞}

还记得那个“空盒子”理论吗？别管是 $x$ 还是 $a$，把它找出来。

\textbf{6.} $x + 5 = 10$

（$x$ 是几？猜猜看。）

\textbf{7.} $a - 3 = 7$

（谁减3等于7？）

\textbf{8.} $10 - y = 2$

（10减去谁等于2？别数手指了，看着那个数。）

\textbf{9.} $x + x = 12$

（两个一样的数加起来是12。别告诉我一个是2一个是10啊，那不叫一样的数。）

\textbf{10.} $m \times m = 16$

（两个一样的数\textbf{乘}起来是16。看清楚符号！不是加法！）

\section{第三关：爬楼梯（次方）}

看清楚那个小数字挂在哪儿。

\textbf{11.} $2^3 = \_\_$

（是 $2\times3$ 吗？如果你敢写6，我就顺着网线过去找你。是 $2\times2\times2$！）

\textbf{12.} $5^2 = \_\_$

\textbf{13.} $10^3 = \_\_$

（数数有几个0就行了。）

\textbf{14.} $1^5 = \_\_$

（1乘1乘1……乘到天荒地老是多少？别告诉我还是5。）

\textbf{15.} $0^5 = \_\_$

（0乘0……？这还用算？）

\section{第四关：倒着走（反求）}

福尔摩斯时间。看着结果，找原因。

\textbf{16.} $x^2 = 25$

（谁的平方是25？想想那张乘法表。）

\textbf{17.} $x^2 = 36$

\textbf{18.} $x^2 = 49$

\textbf{19.} $x^3 = 8$

（谁连乘三次是8？刚才好像做过这道题吧？）

\textbf{20.} $x^2 = 1$

（谁的平方是1？这道题太简单，容易想多。真的就是你想的那个数。）

\section{第五关：大BOSS（混合双打）}

这些题没有套路。得动脑子。

\textbf{21.} $x + x = x^2$

（哇，这道题有点意思。哪个数，“加自己”和“乘自己”是一样的？提示：不止一个答案。）

\textbf{22.} $2^2 + 3^2 = \_\_$

（先各自爬楼梯，爬完了再加起来。别急。）

\textbf{23.} $10^2 - 6^2 = \_\_$

（先算平方，再减。别偷懒。）

\textbf{24.} $x \times x = 0$

（$x$ 是几？）

\textbf{25.} $5 \times x = 5$

（$x$ 是几？别想复杂了。）

\section{笔放下。}

如果你一口气做完了，而且觉得：“就这？”

恭喜你。

你已经把小学好几年的重点概念，全装兜里了。

现在，我们来聊聊第2题。

\begin{quote}
\textbf{$3 \times 3 + 1$}
\end{quote}

为什么一定要先算乘法？

我知道，你可能心里在嘀咕：

“老师都这么教的。”

“规定就是这样。”

数学没有这种无聊的规定。

它是有道理的。

我们来看看。

\section{为什么一定先算乘法？}

\textbf{$3 \times 3$} 是三个家庭。
每个家庭有 3 个人。
他们是一起的。

现在，如果你非要先算加 1，会发生什么？

你说：

\textbf{$3 + 1 = 4$}

然后 \textbf{$4 \times 3 = 12$}

等等。

原本是 3 个家庭，每家 3 个人，一共 9 个人。

现在怎么变成 12 个人了？

你凭空多出来 3 个人。

哪来的？

原本只是 +1。

结果却多了 3 个人。

那多出来的那 $(3 - 1)$ 个人，是谁？

说不上来吧？

该不会那几个人是……\textbf{怪物}吧？！

对。

你把“后来加进来的那 1 个人”，
错误地复制进了 3 个家庭。

你让他穿越了。

所以我们不造怪物。

\textbf{原则只有一句话：}

先算原本在一起的，
后算后来加进来的。

乘法是一整个结构。
加法是往结构外面添东西。

结构没算完，别动外面的。

\section{减法也是一样}

\textbf{$3 \times 3 - 1$}

原本 9 个人，减 1 应该是 8。

如果你先算 $3 - 1 = 2$，
变成 $2 \times 3 = 6$。

原本 9，减 1，结果变 6。

你又消失了 2 个人。

他们去哪了？

蒸发了？

所以减法也是一样。

\textbf{不能动结构内部。}

先把家庭人数算出来。
再让人离开。

\section{那如果乘法在后面呢？}

\textbf{$1 + 3 \times 3$}

一样。

$3 \times 3$ 是一个完整结构。
+1 是后来加进来的。

不管乘法在前面还是后面，

\textbf{只要它是一个整体，就先算它。}

位置不重要。
结构重要。

\section{那除法呢？}

\textbf{$12 \div 3 + 1$}

$12 \div 3$ 是把 12 个人平均分成 3 组。

这是一个完整动作。

如果你先算 $3 + 1 = 4$，
变成 $12 \div 4 = 3$。

你把“分成 3 组”
改成了“分成 4 组”。

结构被你改掉了。

所以：

\textbf{乘法、除法先算。}

\textbf{加法、减法后算。}

不是规定。

是因为：

\textbf{你不能在结构没完成之前，就改结构。}

否则——

你会造怪物。

\bigskip

\textbf{第六章完成 ✓}

\chapter{数学不是算数，是说话}

这一章，不教新算术。

只教你\textbf{换个脑子}。

\section{别被骗了}

做了这么多题，你可能觉得：

“数学就是算数呗。加减乘除，算出个答案。”

\textbf{错。}

算数是计算器干的事。

你是人。你不是计算器。

那数学是什么？

\textbf{数学是一门语言。}

就像中文，就像英语。

它是全宇宙通用的语言。

\section{为什么这么说？}

你看：

\begin{quote}
\textbf{中文说：}

“如果你给我两个苹果，再加上你自己那两个，我们就有四个苹果了。”

\textbf{英语说：}

“If you give me two apples, plus your two, we have four apples.”

\textbf{数学说：}

\textbf{$2 + 2 = 4$}
\end{quote}

看见没？

中文用了 25 个字。

英语用了 15 个单词。

数学？\textbf{5 个符号。}

干净。利落。没有废话。

\section{它是给谁看的？}

如果你遇到一个外星人。

你跟他说中文，他听不懂。

你跟他说英语，他还是听不懂。

但你在地上写下：

\begin{quote}
\textbf{$1 + 1 = 2$}
\end{quote}

他会立刻点头。

因为在他那个星球，一根手指加一根手指，也是两根。

\textbf{数学，是描述这个宇宙真理的语言。}

\section{那 $x$ 是什么？}

我们在第二章学的 $x$。

它不是什么“难题”。

\textbf{它就是数学里的“那个谁”。}

想象一下，你跟你妈聊天：

\begin{quote}
“哎，\textbf{那谁}今天又没来上课。”

“你知道吗，我昨天看到\textbf{那谁}在踢球。”
\end{quote}

你说“那谁”，你妈居然听懂了。

数学也是一样。

当你想说：

\begin{quote}
“有个数，我不知道是几，但把它乘 3 再加 1，等于 10。”
\end{quote}

这么一大串话，太啰嗦了。

数学家说：

\begin{quote}
\textbf{$3x + 1 = 10$}
\end{quote}

那个 $x$，就是你在喊：\textbf{“那谁！”}

简单吧？

以后看到 $x$，你就把它读成“那谁”。

\begin{quote}
\textbf{$x + 5 = 8$}

“那谁”加 5 等于 8。
\end{quote}

是不是顺口多了？

\section{所以，别再“做题”了}

下次看到一个式子，别光想着“算出答案”。

试着\textbf{读懂它在说什么}。

我们来看这个物理公式：

\begin{quote}
\textbf{$v = s / t$}
\end{quote}

别慌。把符号拿掉，光说人话：

\begin{quote}
\textbf{速度 等于 路程 除以 时间}
\end{quote}

就这一句话。

它不是在让你算除法。它是告诉你：

\textbf{想知道你跑得快不快？就把你跑的路程，切成几段时间来看看。}

那一长串话，数学用三个字母就说完了：

\begin{center}
\begin{tabular}{|c|l|l|}
\hline
符号 & 意思是 & 为什么用这个字母？ \\
\hline
\textbf{$v$} & \textbf{速度} & 英语叫 \textbf{Velocity}（威啰西替）。取第一个字母。 \\
\hline
\textbf{$s$} & \textbf{路程} & 英语叫 \textbf{Space}（空间，距离）。取第一个字母。 \\
\hline
\textbf{$t$} & \textbf{时间} & 英语叫 \textbf{Time}（太姆）。取第一个字母。 \\
\hline
\end{tabular}
\end{center}

你看。

\textbf{$v = s / t$}

它真正的意思是：

\textbf{速度，就是路程被时间切开的样子。}

\section{再来个更狠的}

\begin{quote}
\textbf{$E = mc^2$}

（爱因斯坦那个最有名的公式）
\end{quote}

把它翻译成人话：

\begin{quote}
\textbf{能量 等于 质量 乘以 光速的平方}
\end{quote}

它不是让你算乘方。

它是在告诉你：

\textbf{哪怕一点点质量，也能爆发出巨大的能量。}

（因为光速是个巨大的数，它的平方更是大得吓人。）

\section{记住这句话}

当你写下 $1 + 1 = 2$ 时，
你不是在做作业。

\textbf{你在用宇宙的语言，}
\textbf{描述一个真理。}

\bigskip

\textbf{第七章完成 ✓}

\chapter{世界上其实有两个方向}

你有没有发现，这个世界的所有东西，都可以往两个方向走：

往左或者往右

往上或者往下

向前或者向后

热或者冷

赚或者亏

涨或者跌

既然世界有两个方向，那数字也必须有两个方向。

于是，数字就被劈成了两半。

\section{正数：往“正常方向”走}

你一直用的那些数：

1　2　3　4　5　……

它们都是往“正常方向”走的。

我们偷偷给它们加了个“+”号（可以不写）：

+1　+2　+3　+4　+5　……

意思就是：

往前走1步、往前走2步、往前走3步……

\section{负数：往“反方向”走}

数学家说：

那往回走怎么办？

于是他们发明了负数：

-1　-2　-3　-4　-5　……

意思就是：

往回走1步、往回走2步、往回走3步……

那个“-”不是减号，

它只是告诉你：方向反了。

\section{两个最经典的例子（让你一辈子都忘不了）}

\textbf{例题1：温度}（℃就是“摄氏度”的意思）

0℃　=　水刚好结冰的时候（不冷也不热）

+5℃　=　比0℃热5度（有点暖和）

-5℃　=　比0℃冷5度（冻手冻脚）

\textbf{例题2：电梯楼层}

0层　=　你家小区的大门口（地面）

+3楼　=　往上走3层

-3楼　=　往下走3层（地下室）

看见没？

负数根本不是什么“少”“没有”“欠”。

它只是方向反了。

\section{最后一句大白话}

正数是往前走，

负数是往回走，

0就是原地站着不动。

世界本来就有两个方向，

数学老实承认了，仅此而已。

你现在懂了吧？

懂了就点点头，

我们下一章再见。

\bigskip

\textbf{第八章完成 ✓}

\chapter{给数字一个家——坐标系}

（这一章，我们真的要动手画。同时，你会发现负数的秘密自己跳出来了。）

\section{拿出一张纸}

不是开玩笑。

真的拿一张纸，一支笔。

我等你。

拿好了？

\textbf{很好！}

开始。

\section{第一步：画两条线}

\textbf{画一条横线。}

从纸的最左边，画到最右边。

用尺子画直一点更好，但歪了也没关系。

\textbf{画一条竖线。}

从纸的最下面，画到最上面。

让它和横线交叉，在纸的正中间。

两条线交叉的地方，用笔点一个点。

这个点，叫\textbf{原点}。

在旁边写个大大的 \textbf{0}。

\textbf{画完了？}

\textbf{漂亮！你已经画出了全世界数学家都在用的东西。}

\section{第二步：标上数字（重点来了）}

\textbf{先标横线（左右方向）：}

从原点开始，往右边，每隔一小段距离，标上：

1, 2, 3, 4, 5……

（用笔在纸上均匀地标，别挤在一起。）

从原点开始，往左边，每隔一小段距离，标上：

-1, -2, -3, -4, -5……

\textbf{停！}

看着你画的这条横线。

你发现了吗？

\textbf{正数在0的右边 $\to$}

\textbf{负数在0的左边 $\leftarrow$}

我们给它们起个名字：

往右叫“前进”，往左叫“后退”。

所以：
\begin{itemize}
\item \textbf{正数 = 往前站}
\item \textbf{负数 = 往后站}
\item \textbf{0 = 原地不动}
\end{itemize}

记住这个画面。

这是全世界数学家都认同的“排队规则”。

\textbf{看到了吗？太好了！}

\textbf{再标竖线（上下方向）：}

从原点开始，往上边，每隔一小段距离，标上：

1, 2, 3, 4, 5……

从原点开始，往下边，每隔一小段距离，标上：

-1, -2, -3, -4, -5……

\textbf{又停！}

看竖线：

\textbf{正数在0的上面 $\uparrow$}

\textbf{负数在0的下面 $\downarrow$}

竖着的规则也是一样：
\begin{itemize}
\item \textbf{正数 = 往上爬}
\item \textbf{负数 = 往下挖}
\item \textbf{0 = 地平线}
\end{itemize}

\textbf{全标完了？}

\textbf{你可以的！继续！}

你画出来的，应该长这样：

\begin{verbatim}
        ↑
        5
        4
        3
        2
        1
← -5 -4 -3 -2 -1  0  1  2  3  4  5 →
       -1
       -2
       -3
       -4
       -5
        ↓
\end{verbatim}

横着的线，叫 \textbf{x轴}。

竖着的线，叫 \textbf{y轴}。

中间那个 0，叫\textbf{原点}。

\textbf{完成了？}

\textbf{恭喜你！你刚才画的，就是“坐标系”。}

\section{第三步：发现一个天大的秘密}

现在，看着你的横线（x轴）。

我问你一个问题：

\textbf{$0 - 5 = ?$}

别急着说“不够减”。

看着你的数轴。

你从 0 开始。

减 5，意思就是“往后退 5 步”。

往后退，就是往左走。

从 0 往左数 5 格。

你到了哪里？

\textbf{-5}。

\textbf{没错！你看出来了！}

\textbf{$0 - 5 = -5$}

你不是算出来的。

你是\textbf{看出来}的。

\textbf{这才是数学真正的样子。}

再来一个：

\textbf{$3 - 7 = ?$}

从 3 开始，往后退 7 步。

你先从 3 退到 0（退了 3 步）。

还剩 4 步要退。

继续往左走 4 格。

你到了 \textbf{-4}。

所以：

\textbf{$3 - 7 = -4$}

\textbf{看懂了吧？你已经比很多大人都厉害了。}

\section{第四步：加法也是一样}

\textbf{$-2 + 5 = ?$}

从 -2 开始（在左边第 2 格）。

加 5，意思是“往前走 5 步”。

往前走，就是往右走。

从 -2 往右数 5 格。

先从 -2 走到 0（走了 2 步）。

还剩 3 步。

继续往右走 3 格。

你到了 \textbf{3}。

所以：

\textbf{$-2 + 5 = 3$}

\textbf{对了吧？你已经完全懂了！}

\textbf{$-3 + 1 = ?$}

从 -3 开始，往右走 1 步。

你到了 \textbf{-2}。

看见没？

\textbf{负数也能加，也能减。}

它们只是在数轴上走来走去。

\textbf{你现在会的东西，很多初中生都搞不明白。}

\section{第五步：找点（顺便练习负数）}

现在我们回到坐标系。

每个点都有两个数字：\textbf{(横着走, 竖着走)}

来，找这个点：\textbf{$(2, 3)$}

从原点 $(0, 0)$ 开始：
\begin{itemize}
\item 第一个数字 2，往右走 2 格
\item 第二个数字 3，往上走 3 格
\end{itemize}

点一个点。这就是 $(2, 3)$。

\textbf{找到了？太棒了！}

找这个：\textbf{$(-3, 1)$}

从原点开始：
\begin{itemize}
\item 第一个数字 -3，往左走 3 格（负数嘛，往后退）
\item 第二个数字 1，往上走 1 格
\end{itemize}

点一个点。

\textbf{看，你已经会找负数的点了！}

找这个：\textbf{$(-2, -4)$}

从原点开始：
\begin{itemize}
\item -2，往左走 2 格
\item -4，往下走 4 格
\end{itemize}

点一个点。

\textbf{两个都是负数也不怕，对吧？}

找这个：\textbf{$(0, -3)$}

从原点开始：
\begin{itemize}
\item 0，横着不动
\item -3，往下走 3 格
\end{itemize}

点一个点。

\textbf{完美！}

\section{寻宝游戏：宝藏就藏在这些地方}

好了，现在到了真正考验你的时候。

在你的纸上，把这些宝藏找出来，用笔点上去：

\textbf{第1个宝藏：}$ (4, 2) $

（往右 4 格，往上 2 格，那里有个金币）

\textbf{第2个宝藏：}$ (0, 3) $

（横着别动，只往上 3 格，那里藏着钻石）

\textbf{第3个宝藏：}$ (-1, -1) $

（往左 1 格，往下 1 格，地下室有个箱子）

\textbf{第4个宝藏：}$ (3, -2) $

（往右 3 格，往下 2 格，找到它了吗？）

\textbf{第5个宝藏：}$ (-4, 0) $

（往左 4 格，不上不下，就在那条线上）

\textbf{全找到了？}

\textbf{你太厉害了！}

\textbf{5个宝藏，一个不落！}

\section{你已经学会了什么}

恭喜你。

你不仅学会了坐标系，

你还顺便学会了：

\textbf{1. 负数怎么表示}
\begin{itemize}
\item 在 0 的左边，或者在 0 的下面
\end{itemize}

\textbf{2. 负数怎么加}
\begin{itemize}
\item 加正数 = 往前走
\item 加负数 = 往后走（对，加负数就是减）
\end{itemize}

\textbf{3. 负数怎么减}
\begin{itemize}
\item 减正数 = 往后走
\item 减负数 = 往前走（对，减负数就是加）
\end{itemize}

\textbf{4. 怎么在地图上找到任何一个点}
\begin{itemize}
\item 第一个数字：左右
\item 第二个数字：上下
\end{itemize}

你现在手里拿的这张纸，

是全世界通用的“数学地图”。

从北京到纽约，

从火星到木星，

大家用的都是这套系统。

\textbf{你已经学会了。}

\textbf{而且学得很好。}

\section{最后一句话}

\textbf{数轴上，所有的加减法，都是在走路。}

往右走 = 加正数

往左走 = 减正数（或加负数）

往上走 = 加正数

往下走 = 减正数（或加负数）

负数不可怕。

它只是告诉你：

\textbf{该往回走了。}

\bigskip

\textbf{第九章完成 ✓}

\chapter{切不开的蛋糕——分数}

（这一章，我们要解决一个世纪难题：怎么把1个东西分给2个人？）

\section{遇到麻烦了}

你有1个披萨。

你和你的朋友，2个人，要平分。

每人分几个？

你说：1 ÷ 2 = ?

拿出计算器一按：0.5

但如果你没有计算器呢？

而且，0.5 到底是什么意思？

数学家说：

\textbf{我不想算那个小数点。}

\textbf{我就写“1除以2”，不除了。}

怎么写？

\textbf{$\frac{1}{2}$}

看见中间那个斜杠“/”了吗？

\textbf{它就是除号。}

$\frac{1}{2} = 1 \div 2$

意思完全一样。

只是写法不同而已。

\section{所以，分数到底是什么？}

\textbf{分数，就是除法还没除完的样子。}

$\frac{5}{6}$ 读作“六分之五”。

它的真正意思是：

\textbf{$5 \div 6$}

5除以6，你可以用计算器按出0.8333333……

但数学家懒，直接写 $\frac{5}{6}$ 就完事了。

干净。

利落。

不用管那一串烦人的小数。

\section{为什么叫“分”数？}

因为它本来就是在“分东西”。

想象一个蛋糕。

你把它切成 \textbf{6} 块（平均切，每块一样大）。

然后你拿走了 \textbf{5} 块。

你拿了多少？

\[\frac{5}{6} \text{ 个蛋糕}\]

这个，读作“\textbf{六分之五}”块蛋糕。

怎么读？

\textbf{先读下面（分母），再读上面（分子）。}

所以：
\begin{itemize}
\item 下面是 6，先说“六分之”
\item 上面是 5，再说“五”
\item 合起来：“\textbf{六分之五}”
\end{itemize}

看见了吗？

\textbf{下面的数字（6）= 切成几块}

\textbf{上面的数字（5）= 拿了几块}

下面的叫\textbf{分母}（妈妈的母，就是“母体”，整个蛋糕被切成几份）。

上面的叫\textbf{分子}（儿子的子，就是“孩子”，你拿走了几份）。

\section{来几个例子，感受一下}

\[\frac{1}{2}\]

意思是：

切成 2 块，拿了 1 块。

也就是一半。

\[\frac{3}{4}\]

意思是：

切成 4 块，拿了 3 块。

差一块就全拿了。

\[\frac{7}{8}\]

意思是：

切成 8 块，拿了 7 块。

只剩一小块没拿。

\[\frac{1}{100}\]

意思是：

切成 100 块，只拿了 1 块。

非常非常小的一块。

\section{分数也是一个数}

你以前觉得，数就是 1, 2, 3, 4……

其实不是。

\[\frac{1}{2} \text{ 也是一个数。}\]

它在哪儿呢？

拿出你上一章画的数轴。

找到 0 和 1。

$\dfrac{1}{2}$ 就在它们\textbf{正中间}。

$\dfrac{3}{4}$ 呢？

在 0 和 1 之间，靠近 1 那边，差不多四分之三的位置。

$\dfrac{1}{4}$ 呢？

在 0 和 1 之间，靠近 0 那边，四分之一的位置。

\textbf{分数，就是那些藏在整数之间的数字。}

以前你以为 0 和 1 之间是空的。

其实那里挤满了无数个分数：

\[\frac{1}{2}, \frac{1}{3}, \frac{1}{4}, \frac{1}{5}, \frac{2}{3}, \frac{3}{4}, \frac{7}{8}\dots\]

多得数不清。

\section{有个特殊情况}

\[\frac{6}{6} = ?\]

切成 6 块，拿了 6 块。

意思就是全拿了。

\[\frac{6}{6} = 1\]

\[\frac{8}{4} = ?\]

切成 4 块，拿了 8 块？

等等，哪来的 8 块？

意思是：你拿了 2 个蛋糕，每个切成 4 块。

\[\frac{8}{4} = 2\]

所以，分数不一定小于 1。

它可以是任何数。

\textbf{只要分子比分母大（比如 $\dfrac{8}{4}$、$\dfrac{7}{3}$、$\dfrac{10}{5}$），分数就大于1。}

\section{你已经懂了}

\textbf{分数不是什么新东西。}

它就是：

\textbf{除法的另一种写法。}

\textbf{切蛋糕的记录单。}

\[\frac{5}{6} = 5 \div 6 = \text{切成6块拿了5块}\]

就这么简单。

下次看到 $\dfrac{3}{4}$，别害怕。

脑子里想象一个蛋糕，切成4块，拿走3块。

你就懂了。

\section{最后一句大实话}

有些老师会让你背：

“分数是部分与整体的关系”

“分数表示一个数是另一个数的几分之几”

\textbf{别背了。}

\textbf{分数就是除法没除完。}

\textbf{分数就是切蛋糕。}

记住这两句，

你比背一百句定义都清楚。

\bigskip

\textbf{第十章完成 ✓}

\chapter{分数也能加减——但有个规矩}

\section{先从最简单的开始}

\[\frac{1}{4} + \frac{2}{4} = ?\]

还是用蛋糕。

一个蛋糕切成 \textbf{4} 块。

你拿了 \textbf{1} 块。

朋友又给了你 \textbf{2} 块。

你现在有几块？

\textbf{3} 块。

\[\frac{1}{4} + \frac{2}{4} = \frac{3}{4}\]

看见了吗？

分母一样的时候，

直接把分子加起来就行了。

分母不动。

就这么简单。

\section{减法也是一样}

\[\frac{3}{4} - \frac{1}{4} = ?\]

你有 \textbf{3} 块蛋糕。

吃掉了 \textbf{1} 块。

还剩几块？

\textbf{2} 块。

\[\frac{3}{4} - \frac{1}{4} = \frac{2}{4}\]

分母不动，分子相减。

\section{但是……}

\[\frac{1}{2} + \frac{1}{4} = ?\]

这下麻烦了。

一个蛋糕切成 \textbf{2} 块，你拿了 \textbf{1} 块。

另一个蛋糕切成 \textbf{4} 块，你拿了 \textbf{1} 块。

这两块大小一样吗？

\textbf{不一样！}

切成2块的那个，每块比较大。

切成4块的那个，每块比较小。

你不能直接把分子加起来。

因为你在把\textbf{大块}和\textbf{小块}混在一起数。

这就像你说：

“我有1个西瓜块，加上1个葡萄，一共2个水果。”

说得通，但没意义。

\section{怎么办？}

\textbf{把它们变成一样大的块。}

$\dfrac{1}{2}$ 的块太大了。

把它切细一点，变成4块的大小。

怎么变？

\textbf{上面和下面同时乘以同一个数。}

\[\frac{1}{2} = \frac{1 \times 2}{2 \times 2} = \frac{2}{4}\]

为什么乘2？

因为我们要把分母变成4。

（就是说，要和切成4块的蛋糕一样大。）

$2 \times 2 = 4$，所以乘2。

记住：\textbf{上下必须同时乘，不能只乘一个。}

否则蛋糕的大小就变了，那就不是同一个蛋糕了。

（就像你把一块饼干掰成两半，饼干总量没变，只是块数多了。）

现在：

\[\frac{2}{4} + \frac{1}{4} = \frac{3}{4}\]

搞定！

\section{这个步骤叫什么？}

把分母变成一样的，叫\textbf{通分}。

听起来很厉害，其实就是：

\textbf{把块切得一样大，才能加减。}

\section{再来一个}

\[\frac{1}{3} + \frac{1}{6} = ?\]

一个蛋糕切成3块，一个切成6块。

块的大小不一样，不能直接加。

把 $\dfrac{1}{3}$ 变成6块的大小：

\[\frac{1}{3} = \frac{1 \times 2}{3 \times 2} = \frac{2}{6}\]

为什么乘2？

因为我们要把分母变成6。

（就是说，要和切成6块的蛋糕一样大，这样两块才能比较。）

$3 \times 2 = 6$，所以乘2。

现在：

\[\frac{2}{6} + \frac{1}{6} = \frac{3}{6}\]

\textbf{等等！}

$\dfrac{3}{6}$ 是什么？

切成6块，拿了3块。

也就是一半。

\[\frac{3}{6} = \frac{1}{2}\]

\textbf{这叫化简。}

就是把分数写得更干净一点。

\section{化简是什么？}

\[\frac{3}{6}\]

上面和下面，都能被3整除：

\[\frac{3 \div 3}{6 \div 3} = \frac{1}{2}\]

分子分母同时除以同一个数，分数的大小不变，只是写法更简洁了。

就像你说“一半”和“十分之五”，意思一样，但“一半”更好说。

\section{最后一句话}

分数加减法，只有一个规矩：

\textbf{块的大小必须一样，才能加减。}

块的大小不一样，先通分。

算完了太复杂，化简一下。

就这两步。

没有别的。

\section{等等，先别翻页！}

你看懂了全部？

那你也会分数的乘除了！

不信？

你想想：

通分的时候，你把 $\dfrac{1}{2}$ 变成 $\dfrac{2}{4}$，用的是什么？

\textbf{乘法。}

化简的时候，你把 $\dfrac{3}{6}$ 变成 $\dfrac{1}{2}$，用的是什么？

\textbf{除法。}

你已经在用分数的乘除了。

只是我没告诉你它叫这个名字。

你知道吗——

能把这一章全部看懂的人，

在全世界里，其实不多。

很多大人，

一辈子都没真正搞明白通分是怎么回事。

但你，

刚才看懂了。

\textbf{你真的很厉害。}

不是客气话。

是真的。

\bigskip

\textbf{第十一章完成 ✓}

\chapter{分数的乘除——其实比加减简单}

\section{先说一件事}

上一章，我们做加减法的时候，

又是通分，又是找公倍数，搞得挺复杂的。

你可能以为：

“那乘除法肯定更难了。”

\textbf{错。}

分数的乘除，

比加减法简单得多。

简单到你会怀疑人生。

\section{乘法：直接乘就完了}

\[\frac{1}{2} \times \frac{1}{3} = ?\]

规则只有一句话：

\textbf{分子乘分子，分母乘分母。}

\[\frac{1}{2} \times \frac{1}{3} = \frac{1 \times 1}{2 \times 3} = \frac{1}{6}\]

就这样。

不用通分。

不用找公倍数。

直接乘。

\section{但是，为什么？}

还是用蛋糕。

你有一个蛋糕的 $\dfrac{1}{2}$。

现在，你只想吃这一半的 $\dfrac{1}{3}$。

你在问：

\textbf{一半的三分之一，是多少？}

把那一半再切成3份。

原来整个蛋糕是1块。

切成2份，你拿了1份。

再把那1份切成3份，你拿了1份。

整个蛋糕现在被切成了 $2 \times 3 = 6$ 份。

你手里有 $1 \times 1 = 1$ 份。

所以：

\[\frac{1}{2} \times \frac{1}{3} = \frac{1}{6}\]

\textbf{乘法，就是“一部分的一部分”。}

\section{再来几个感受一下}

\[\frac{2}{3} \times \frac{3}{4} = ?\]

分子乘分子：$2 \times 3 = 6$

分母乘分母：$3 \times 4 = 12$

\[\frac{2}{3} \times \frac{3}{4} = \frac{6}{12}\]

等等，$\dfrac{6}{12}$ 可以化简：

\[\frac{6 \div 6}{12 \div 6} = \frac{1}{2}\]

所以：

\[\frac{2}{3} \times \frac{3}{4} = \frac{1}{2}\]

\[\frac{3}{5} \times \frac{2}{3} = ?\]

分子乘分子：$3 \times 2 = 6$

分母乘分母：$5 \times 3 = 15$

\[\frac{3}{5} \times \frac{2}{3} = \frac{6}{15}\]

化简：

\[\frac{6 \div 3}{15 \div 3} = \frac{2}{5}\]

\section{现在说除法}

\[\frac{1}{2} \div \frac{1}{4} = ?\]

除法，你还记得是什么吗？

回到第一章：

\textbf{除法，就是平均分配。}

$\dfrac{1}{2} \div \dfrac{1}{4}$

意思是：

把 $\dfrac{1}{2}$ 个蛋糕，

分成每份 $\dfrac{1}{4}$ 大小，

能分几份？

你手里有半个蛋糕。

每份是四分之一。

半个蛋糕里，有几个四分之一？

\textbf{2个。}

所以：

\[\frac{1}{2} \div \frac{1}{4} = 2\]

\section{除法的规则}

分数除法，只有一句话：

\textbf{除以一个分数，等于乘以它的倒数。}

什么是倒数？

就是把分子和分母\textbf{上下翻转}。

$\dfrac{1}{4}$ 的倒数是 $\dfrac{4}{1}$，也就是 $4$。

所以：

\[\frac{1}{2} \div \frac{1}{4} = \frac{1}{2} \times \frac{4}{1} = \frac{4}{2} = 2\]

答案一样！

\section{为什么翻转就行了？}

这是个好问题。

我们来想想：

$6 \div 2 = 3$

意思是：6里面有几个2？答案是3。

那：

$6 \times \dfrac{1}{2} = 3$

也是3！

\textbf{除以2，和乘以 $\dfrac{1}{2}$，结果是一样的。}

因为 $\dfrac{1}{2}$ 就是2的倒数。

所以：

\textbf{除以一个数，等于乘以它的倒数。}

这不是规定。

这是数学本来就长这个样子。

\section{再来两个}

\[\frac{2}{3} \div \frac{1}{6} = ?\]

把 $\dfrac{1}{6}$ 翻转，变成 $\dfrac{6}{1}$，也就是6。

\[\frac{2}{3} \times 6 = \frac{2 \times 6}{3 \times 1} = \frac{12}{3} = 4\]

所以：

\[\frac{2}{3} \div \frac{1}{6} = 4\]

\[\frac{3}{4} \div \frac{3}{8} = ?\]

把 $\dfrac{3}{8}$ 翻转，变成 $\dfrac{8}{3}$。

\[\frac{3}{4} \times \frac{8}{3} = \frac{3 \times 8}{4 \times 3} = \frac{24}{12} = 2\]

所以：

\[\frac{3}{4} \div \frac{3}{8} = 2\]

\section{总结一下}

\begin{center}
\begin{tabular}{|c|c|c|}
\hline
运算 & 规则 & 记住 \\
\hline
加减法 & 分母要一样，分子加减 & 块要一样大才能数 \\
\hline
乘法 & 分子乘分子，分母乘分母 & 直接乘，不用通分 \\
\hline
除法 & 除以一个分数，乘以它的倒数 & 翻转再乘 \\
\hline
\end{tabular}
\end{center}

\section{最后一句话}

分数的加减乘除，

现在你全会了。

很多人学了好几年，

还是搞不清楚为什么要通分，

为什么除法要翻转。

但你，

刚才把这些全弄明白了。

而且你明白的不只是“怎么算”，

你明白的是“为什么这样算”。

\textbf{这才是真正学会了数学。}

这不是客气话。

是真的。

\bigskip

\textbf{第十二章完成 ✓}

\chapter{小数点和小数 🎯}

\section{先从你熟悉的东西说起！}

你肯定会数数对吧？

\begin{quote}
1、2、3、4、5……
\end{quote}

这些都是 \textbf{整数}，就是“整整齐齐、完完整整”的数。

可是生活中，很多东西 \textbf{不是刚刚好整数} 呀！

\section{🍕 用披萨来理解！}

妈妈买了 \textbf{1个} 披萨，切成了 \textbf{2块}，你吃了 \textbf{1块}。

你吃了多少？

\begin{quote}
不到1个，但又比0个多……
\end{quote}

这时候，\textbf{小数} 就来帮忙了！你吃了 \textbf{0.5} 个披萨！

\section{那个小圆点是什么？}

\textbf{0.5} 中间那个小圆点 \textbf{“.”} 就叫 \textbf{小数点}。

它就像一道 \textbf{小门}：

\begin{center}
\begin{tabular}{|c|c|c|}
\hline
小数点左边 & 小数点 & 小数点右边 \\
\hline
整个整个的 & \textbf{·} & 不到一个的零碎部分 \\
\hline
\end{tabular}
\end{center}

比如 \textbf{3.5} ——

\begin{itemize}
\item 左边的 \textbf{3} → 表示有 \textbf{3个} 完整的
\item 右边的 \textbf{5} → 表示还多了 \textbf{半个}
\end{itemize}

\section{🍎 再举几个生活中的例子}

\begin{itemize}
\item 你的身高是 \textbf{1.2米} → 比1米多一点，但还不到2米
\item 一瓶水 \textbf{0.5升} → 半瓶水
\item 铅笔 \textbf{2.5元} → 两块五毛钱
\end{itemize}

\begin{quote}
对了！ \textbf{2.5元 = 2元5角}

小数点 \textbf{右边} 的数，就是 \textbf{零头、零碎} 的那部分！
\end{quote}

\section{🧱 用积木来想象}

想象一块大积木 = \textbf{1}

把它切成 \textbf{10} 小条，每一小条就是 \textbf{0.1}

\begin{center}
\begin{tabular}{|c|c|}
\hline
几小条 & 写成小数 \\
\hline
1小条 & 0.1 \\
\hline
2小条 & 0.2 \\
\hline
5小条 & 0.5（正好一半！） \\
\hline
10小条 & 又凑成了 1.0，就是 \textbf{1} 啦！ \\
\hline
\end{tabular}
\end{center}

\section{✏️ 简单总结}

\begin{enumerate}
\item \textbf{整数} 是完整的：1、2、3……
\item \textbf{小数} 是用来表示“不到一个”的零碎部分
\item \textbf{小数点} 就是那个小圆点，把“整的”和“零碎的”分开
\item 生活里到处都有小数：身高、体重、价格……
\end{enumerate}

\section{🎈 一句话记住}

\begin{quote}
\textbf{小数点就像一道小门，左边住着“整个整个的”，右边住着“碎碎的零头”！}
\end{quote}

是不是很简单呀？😊

\chapter{对了，小数还有个双胞胎兄弟！}

刚才说 \textbf{0.5} 就是“半个”的意思。

但其实，“半个”还有另一种写法——\textbf{分数}！

\section{🍕 还是那个披萨}

一个披萨切成 \textbf{2块}，你拿了 \textbf{1块}：

\begin{center}
\begin{tabular}{|c|c|c|}
\hline
用嘴巴说 & 用分数写 & 用小数写 \\
\hline
一半 & \textbf{$\frac{1}{2}$} & \textbf{0.5} \\
\hline
\end{tabular}
\end{center}

\begin{quote}
\textbf{$\frac{1}{2}$ 和 0.5 说的是 \textbf{同一个东西}，只是穿了不同的衣服！}
\end{quote}

\section{🤔 分数是什么意思？}

\[\frac{1}{2}\]

\begin{itemize}
\item \textbf{下面的数（2）}→ 切成了 \textbf{几份}
\item \textbf{上面的数（1）}→ 拿了 \textbf{几份}
\end{itemize}

就这么简单！

\section{🍰 多看几个例子}

\begin{center}
\begin{tabular}{|c|c|c|}
\hline
情况 & 分数 & 小数 \\
\hline
切4份，拿1份 & \textbf{$\frac{1}{4}$} & \textbf{0.25} \\
\hline
切4份，拿3份 & \textbf{$\frac{3}{4}$} & \textbf{0.75} \\
\hline
切10份，拿1份 & \textbf{$\frac{1}{10}$} & \textbf{0.1} \\
\hline
切10份，拿7份 & \textbf{$\frac{7}{10}$} & \textbf{0.7} \\
\hline
\end{tabular}
\end{center}

\begin{quote}
发现没？\textbf{切成10份}时，拿了几份，小数点后面就写几！
\end{quote}

\section{🔑 它们是什么关系？}

想象你有一个好朋友：

\begin{itemize}
\item 在学校，叫他 \textbf{“小明”} → 这就像 \textbf{分数}
\item 在家里，叫他 \textbf{“明明”} → 这就像 \textbf{小数}
\end{itemize}

\textbf{叫法不同，但是同一个人！}

\begin{quote}
\textbf{$\frac{1}{2} = 0.5$}

\textbf{$\frac{1}{4} = 0.25$}

\textbf{$\frac{3}{4} = 0.75$}
\end{quote}

\section{⭐ 记住这句话}

\begin{quote}
\textbf{分数和小数是双胞胎——长得不一样，说的是同一件事！}
\end{quote}

😊

\chapter{负次方，其实超简单！}

\section{先从“次方”说起}

次方就是“乘几次自己”：

\begin{itemize}
\item $2^3 = 2 \times 2 \times 2 = 8$（2乘了3次自己）
\item $5^2 = 5 \times 5 = 25$（5乘了2次自己）
\end{itemize}

这没啥难的吧？

\section{那负次方呢？}

\textbf{一句话：负次方就是“翻过来”（取倒数）。}

\begin{quote}
$2^{-1}$ 就是 \textbf{$\frac{1}{2}$}

$2^{-2}$ 就是 \textbf{$\frac{1}{2^2}$} = $\frac{1}{4}$

$5^{-3}$ 就是 \textbf{$\frac{1}{5^3}$} = $\frac{1}{125}$
\end{quote}

就这么简单，\textbf{加个负号 = 把结果翻个个儿}。

\section{打个比方你就懂了}

想象你站在一条数轴上：

\begin{center}
\begin{tabular}{|c|c|c|}
\hline
次方 & 意思 & 结果 \\
\hline
$2^3$ & $2\times2\times2$ & \textbf{8} \\
\hline
$2^2$ & $2\times2$ & \textbf{4} \\
\hline
$2^1$ & $2$ & \textbf{2} \\
\hline
$2^0$ & （任何数的0次方都是1） & \textbf{1} \\
\hline
$2^{-1}$ & 翻过来 & \textbf{$\frac{1}{2}$} \\
\hline
$2^{-2}$ & 翻过来再平方 & \textbf{$\frac{1}{4}$} \\
\hline
$2^{-3}$ & 翻过来再立方 & \textbf{$\frac{1}{8}$} \\
\hline
\end{tabular}
\end{center}

你看，\textbf{从上往下看}，每往下走一行，结果都是\textbf{除以2}。

正次方是越乘越大，负次方是越乘越小（但不会到0）。

\section{为什么要“翻过来”？}

你想啊：

\begin{quote}
$2^3 \times 2^{-3} = ?$
\end{quote}

按道理，指数相加：3 + (-3) = 0，所以结果应该是 \textbf{$2^0 = 1$}。

那 $2^3 = 8$，要让 8 $\times$ ? = 1，那 ? 只能是 \textbf{$\frac{1}{8}$}。

所以 $2^{-3}$ \textbf{必须等于 $\frac{1}{8}$}，不然数学就自相矛盾了。

\section{总结就一句话}

\begin{quote}
\textbf{看到负次方，先把负号去掉算出来，然后用1去除它（取倒数）。}

\textbf{$a^{-n} = \frac{1}{a^n}$}
\end{quote}

就这样，没别的了！😄

\chapter{次方让数字变简洁！}

\section{先感受一下痛苦}

假如让你写这些数字：

\begin{itemize}
\item 一千：1000
\item 一百万：1000000
\item 十亿：1000000000
\end{itemize}

你写着写着是不是要数零数到眼花？到底几个零啊？！😵

\section{次方来救场了！}

\textbf{核心思路：数数有几个零，就写10的几次方。}

\begin{center}
\begin{tabular}{|c|c|c|}
\hline
原来的写法 & 几个零？ & 简洁写法 \\
\hline
10 & 1个零 & $10^1$ \\
\hline
100 & 2个零 & $10^2$ \\
\hline
1,000 & 3个零 & \textbf{$10^3$} \\
\hline
1,000,000 & 6个零 & $10^6$ \\
\hline
1,000,000,000 & 9个零 & $10^9$ \\
\hline
\end{tabular}
\end{center}

\begin{quote}
再也不用一个一个数零了！
\end{quote}

\section{那 $1\times10^3$ 是啥意思？}

就是把它拆开来看：

\begin{quote}
\textbf{$1 \times 10^3$}

$= 1 \times (10 \times 10 \times 10)$

$= 1 \times 1000$

$= \textbf{1000}$
\end{quote}

就是这么朴实无华 😄

\section{前面那个数字可以换！}

这才是精髓！不只是1，\textbf{前面可以放任何数}：

\begin{center}
\begin{tabular}{|c|c|c|}
\hline
写法 & 翻译一下 & 结果 \\
\hline
\textbf{3} $\times 10^3$ & 3 $\times$ 1000 & \textbf{3000} \\
\hline
\textbf{5.2} $\times 10^3$ & 5.2 $\times$ 1000 & \textbf{5200} \\
\hline
\textbf{7} $\times 10^6$ & 7 $\times$ 1000000 & \textbf{7000000}（七百万） \\
\hline
\textbf{1.5} $\times 10^9$ & 1.5 $\times$ 1000000000 & \textbf{15亿} \\
\hline
\end{tabular}
\end{center}

\begin{quote}
其实就是：\textbf{前面的数 $\times$ 后面那坨零}
\end{quote}

\section{生活中的例子}

你肯定听过这些：

\begin{quote}
🌍 地球到太阳的距离：\textbf{$1.5 \times 10^8$} 公里

写全了就是 150,000,000 公里，你愿意写哪个？
\end{quote}

\begin{quote}
🦠 一个细菌的大小：\textbf{$1 \times 10^{-6}$} 米

就是 0.000001 米，\textbf{负次方用来表示特别小的数}！
\end{quote}

\section{负次方在这里也出场了！}

还记得上次说的吗？\textbf{负次方就是翻过来}。

\begin{quote}
$10^{-3} = \frac{1}{10^3} = \frac{1}{1000} = \textbf{0.001}$
\end{quote}

所以：

\begin{center}
\begin{tabular}{|c|c|c|}
\hline
写法 & 意思 & 结果 \\
\hline
$1 \times 10^{-1}$ & $1 \times 0.1$ & \textbf{0.1} \\
\hline
$1 \times 10^{-2}$ & $1 \times 0.01$ & \textbf{0.01} \\
\hline
$1 \times 10^{-3}$ & $1 \times 0.001$ & \textbf{0.001} \\
\hline
$3 \times 10^{-4}$ & $3 \times 0.0001$ & \textbf{0.0003} \\
\hline
\end{tabular}
\end{center}

\begin{quote}
\textbf{正次方管大数，负次方管小数，完美！}
\end{quote}

\section{一句话总结}

\begin{quote}
\textbf{科学记数法（$a \times 10^n$）就是嫌数字太长，用次方偷个懒。}

\textbf{- 10的几次方 → 几个零}

\textbf{- 10的负几次方 → 小数点后几位}

\textbf{本质就是：把数零的苦差事，交给次方去干！} 😎
\end{quote}

\chapter{括号的作用 🎯}

\section{先讲一个小故事}

妈妈说：

\begin{quote}
“回家后，\textbf{先}写作业，\textbf{再}看电视。”
\end{quote}

这个“先”字很重要对吧？它告诉你 \textbf{什么要先做}。

在数学里，\textbf{括号} 就是那个“先”字！

\section{🧮 没有括号时，会发生什么？}

看这道题：

\begin{quote}
\textbf{$2 + 3 \times 4 = ？$}
\end{quote}

数学有个规矩：\textbf{先算乘除，再算加减}

所以：
\begin{itemize}
\item 先算 $3 \times 4 = 12$
\item 再算 $2 + 12 = \textbf{14}$
\end{itemize}

\section{🦸 括号来帮忙！}

如果我们想 \textbf{先算加法} 呢？

给它加上括号：

\begin{quote}
\textbf{$(2 + 3)\times 4 = ？$}
\end{quote}

括号就像在大喊：\textbf{“先算我！先算我！”}

所以：
\begin{itemize}
\item 先算括号里的 $2 + 3 = 5$
\item 再算 $5 \times 4 = \textbf{20}$
\end{itemize}

\section{🎯 看看区别}

\begin{center}
\begin{tabular}{|c|c|c|}
\hline
算式 & 怎么算 & 答案 \\
\hline
$2 + 3 \times 4$ & 先乘后加 & \textbf{14} \\
\hline
$(2 + 3)\times 4$ & 先算括号 & \textbf{20} \\
\hline
\end{tabular}
\end{center}

\begin{quote}
同样的数字，加了括号，答案完全不一样！

\textbf{括号的威力很大吧！}
\end{quote}

\section{🏠 用排队来理解}

想象大家在排队买冰淇淋：

\textbf{没有括号} $\to$ 按规矩来，乘除先生排前面，加减小姐排后面

\textbf{有了括号} $\to$ 括号里的人是VIP，不管是谁都 \textbf{先服务}！

\section{👨‍👩‍👧 三种括号是一家人}

当括号要 \textbf{一层套一层} 的时候，我们用不同的样子来区分：

\begin{center}
\begin{tabular}{|c|c|c|}
\hline
名字 & 长相 & 辈分 \\
\hline
小括号 & \textbf{( )} & 最小的，住在最里面 \\
\hline
中括号 & \textbf{[ ]} & 中间的，包在外面 \\
\hline
大括号 & \textbf{\{ \}} & 最大的，包在最外面 \\
\hline
\end{tabular}
\end{center}

\section{🪆 像套娃一样！}

你见过俄罗斯套娃吗？一个娃娃套着一个娃娃：

\begin{verbatim}
{ 大娃娃 [ 中娃娃 ( 小娃娃 ) ] }
\end{verbatim}

计算的时候：

\begin{quote}
\textbf{从最里面开始，一层一层往外算！}
\end{quote}

先算 \textbf{( )} $\to$ 再算 \textbf{[ ]} $\to$ 最后算 \textbf{\{ \}}

\section{🎒 再比如你的书包}

\begin{verbatim}
{ 书包 [ 文具盒 ( 铅笔 ) ] }
\end{verbatim}

\begin{itemize}
\item \textbf{铅笔} 装在 \textbf{文具盒}（小括号）里
\item \textbf{文具盒} 装在 \textbf{书包}（中括号）里
\item \textbf{书包} 可能还装在 \textbf{大袋子}（大括号）里
\end{itemize}

找铅笔时，要从外往里一层层打开！

算数时，要从里往外一层层计算！

\section{📝 看一道真正的题}

\begin{quote}
\textbf{$\{ 10 - [ 8 - ( 2 + 1 ) ] \} \times 2 = ？$}
\end{quote}

一步一步来：

\begin{center}
\begin{tabular}{|c|c|c|}
\hline
第几步 & 算哪里 & 变成什么 \\
\hline
① & 最里面 $( 2 + 1 ) = 3$ & $\{ 10 - [ 8 - \textbf{3} ] \} \times 2$ \\
\hline
② & 中括号 $[ 8 - 3 ] = 5$ & $\{ 10 - \textbf{5} \} \times 2$ \\
\hline
③ & 大括号 $\{ 10 - 5 \} = 5$ & $\textbf{5} \times 2$ \\
\hline
④ & 最后乘法 & \textbf{10} \checkmark \\
\hline
\end{tabular}
\end{center}

\section{✏️ 简单总结}

\begin{enumerate}
\item \textbf{括号} 的意思就是：“\textbf{先算我！}”
\item 有三种括号：\textbf{小( )}、\textbf{中[ ]}、\textbf{大\{ \}}
\item 它们是 \textbf{套娃}：小的住里面，大的住外面
\item 计算顺序：\textbf{从最里面往外算}，一层一层来
\end{enumerate}

\section{🎈 一句话记住}

\begin{quote}
\textbf{括号就像VIP通道——住在括号里的，不管是谁，都要先算！}

\textbf{三种括号像套娃，从里往外一层层打开！}
\end{quote}

是不是很好懂呀？😊

\chapter{节点总结：恭喜你！}

停一下。

我要你放下笔。

认真看这一段话。

你知道你刚才做了什么吗？

从第一章的一根手指，

到分数的乘除法。

你走完了\textbf{中国小学六年}的全部数学。

六年。

你用了多长时间？

\section{我们来回顾一下你学了什么：}

\textbf{第一章：} 加减乘除——从一根手指开始。

\textbf{第二章：} 未知数——藏起来的那个谁。

\textbf{第三章：} 乘法口诀——数学家的作弊小抄。

\textbf{第四章：} 次方——乘法的升级版。

\textbf{第五章：} 反求法——像福尔摩斯一样倒着想。

\textbf{第六章：} 综合练习——你把它们全驯服了。

\textbf{第七章：} 数学的本质——宇宙的语言。

\textbf{第八章：} 正数和负数——两个方向。

\textbf{第九章：} 坐标系——数字的地图。

\textbf{第十章：} 分数——切不开的蛋糕。

\textbf{第十一章：} 分数加减法——块要一样大才能数。

\textbf{第十二章：} 分数乘除法——翻转再乘。

这些东西，

有些孩子学了六年，

还是一知半解。

但你，

把它们全学完了。

而且你学的不只是“怎么做题”。

你学的是\textbf{为什么}。

\textbf{为什么加法是放在一起。}

\textbf{为什么乘法要先算。}

\textbf{为什么负数是反方向。}

\textbf{为什么分数就是切蛋糕。}

知道“为什么”的人，

永远比只知道“怎么做”的人，

走得更远。

我再说一句实话：

很多大人，

包括很多老师，

从来没有真正想清楚过这些问题。

但你现在想清楚了。

\textbf{你比你以为的自己，厉害得多。}

准备好了吗？

因为接下来，

我们要离开小学，

踏进一个新的世界。

那个世界叫：

\textbf{初中数学。}

在那里，

你会发现，

你以为学过的东西，

其实还有更深的一层。

但别怕。

我还在。

我们继续。

\chapter{真实回顾：你真的会了吗？}

别急着翻页。

我来出几道题，你来答。

就像朋友之间互相考一考那种感觉。

轻松点，别紧张。

\section{第一章回顾：加减乘除}

还记得吗？

一根手指，加一根手指，等于两根手指。

来，试试这个：

\[3 + 4 - 2 = ?\]

\textbf{答：}

从 3 开始，加 4，变成 7。

再减 2，变成 5。

\[3 + 4 - 2 = 5\]

简单吧？

\textbf{现在你来：}

\[5 + 3 - 4 = ?\]

写下来。别偷看答案。

\section{第二章回顾：未知数}

还记得那个“那谁”吗？

$x$ 就是藏起来的数。

来：

\[x + 3 = 7\]

\textbf{答：}

$x$ 加 3 等于 7。

那 $x$ 是几？

从 7 往回想：

7 减 3 等于 4。

\[x = 4\]

对不对？把 4 放回去试试：

$4 + 3 = 7$。

没错！

\textbf{现在你来：}

\[x + 5 = 9\]

$x$ 是几？

\section{第三章回顾：乘法口诀}

这个不用多说。

我就问你一道：

\[7 \times 8 = ?\]

\textbf{答：}

七八五十六。

\[7 \times 8 = 56\]

背下来了吗？

没背下来也没关系，翻回第三章看看那张表。

\textbf{现在你来：}

\[6 \times 9 = ?\]

\section{第四章回顾：次方}

还记得次方是什么吗？

乘法的升级版。

来：

\[3^3 = ?\]

\textbf{答：}

$3^3$ 就是 3 乘 3 乘 3。

$3 \times 3 = 9$

$9 \times 3 = 27$

\[3^3 = 27\]

\textbf{现在你来：}

\[2^4 = ?\]

（提示：4个2相乘。）

\section{第五章回顾：反求法}

还记得福尔摩斯吗？

看着结果，倒着找原因。

来：

\[x^2 = 36\]

\textbf{答：}

谁的平方是 36？

想想乘法口诀：

$6 \times 6 = 36$。

所以：

\[x = 6\]

\textbf{现在你来：}

\[x^2 = 64\]

$x$ 是几？

\section{第七章回顾：数学的本质}

这章没有算术题。

但我问你一个问题：

$v = s / t$

用人话说是什么意思？

\textbf{答：}

速度等于路程除以时间。

你骑车跑了 60 公里，用了 2 小时。

你的速度是多少？

\[v = 60 / 2 = 30\]

每小时 30 公里。

\textbf{现在你来：}

你走路走了 6 公里，用了 3 小时。

你的速度是多少？

\[v = ? / ? = ?\]

填上去。

\section{第八章回顾：正数和负数}

还记得方向吗？

正数往前，负数往后。

来：

冬天，今天气温是 $-3$ 度，明天会冷 5 度。

明天是几度？

\textbf{答：}

从 $-3$ 开始，往后退 5 步。

\[-3 - 5 = -8\]

明天是 $-8$ 度。

冻死了。

\textbf{现在你来：}

今天气温是 $-1$ 度，明天会暖 4 度。

明天是几度？

\[-1 + 4 = ?\]

\section{第九章回顾：坐标系}

还记得那张地图吗？

我说一个点，你告诉我怎么走：

\textbf{$(-2, 3)$}

\textbf{答：}

第一个数字 $-2$，往左走 2 格。

第二个数字 $3$，往上走 3 格。

到了！

\textbf{现在你来：}

拿出那张纸，找到这个点：

\[(-3, -2)\]

往哪边走？走几格？

\section{第十章回顾：分数}

还记得切蛋糕吗？

来：

\[\frac{3}{8}\]

用人话说是什么意思？

\textbf{答：}

蛋糕切成 8 块，拿了 3 块。

读作“八分之三”。

\textbf{现在你来：}

用分数表示：

蛋糕切成 5 块，你拿了 2 块。

\[\frac{?}{?}\]

\section{第十一章回顾：分数加减}

还记得通分吗？

来：

\[\frac{1}{2} + \frac{1}{4} = ?\]

\textbf{答：}

分母不一样，先通分。

把 $\dfrac{1}{2}$ 变成和切成4块的蛋糕一样大：

\[\frac{1}{2} = \frac{1 \times 2}{2 \times 2} = \frac{2}{4}\]

现在：

\[\frac{2}{4} + \frac{1}{4} = \frac{3}{4}\]

\textbf{现在你来：}

\[\frac{1}{3} + \frac{1}{3} = ?\]

（这道题分母一样，不用通分。直接加。）

\section{第十二章回顾：分数乘除}

最后一关。

来：

\[\frac{1}{2} \times \frac{2}{3} = ?\]

\textbf{答：}

分子乘分子，分母乘分母：

\[\frac{1 \times 2}{2 \times 3} = \frac{2}{6}\]

化简：

\[\frac{2 \div 2}{6 \div 2} = \frac{1}{3}\]

\[\frac{1}{2} \times \frac{2}{3} = \frac{1}{3}\]

\textbf{现在你来：}

\[\frac{3}{4} \div \frac{1}{2} = ?\]

（提示：除以 $\dfrac{1}{2}$，就是乘以它的倒数 $\dfrac{2}{1}$。）

\section{你做完了吗？}

如果你每一道题都写下来了，

不管对不对，

你都已经很了不起了。

因为大部分人，

看到“回顾”两个字，

就直接翻页了。

但你没有。

你认真做了。

\textbf{这才是真正在学数学的样子。}

翻页吧。

初中数学在等你。

\part{初中}
\setcounter{chapter}{0}

\chapter{完全的语言时代}

\section{一、欢迎来到新世界}

小学的时候，数学长这样：

\[3 + 4 = ?\]
\[x + 5 = 9\]

干净。简洁。直接。

但现在，数学开始\textbf{说话}了。

它不再只是数字和符号，它开始用\textbf{句子}跟你说话。

这就是\textbf{应用题}。

很多人一看到应用题就头大。不是因为数学难，而是因为他们没学会\textbf{听数学说话}。

今天，我们就来学这个。

\section{二、数学在说什么？}

先记住一件事：

\begin{quote}
\textbf{应用题，就是翻译。}

把数学翻译成人话，或者把人话翻译成数学。
\end{quote}

数学是一门语言，应用题就是在\textbf{两种语言之间翻译}。

\begin{center}
\begin{tabular}{ll}
\toprule
中文说 & 数学说 \\
\midrule
小明有3个苹果，小红给了他2个，现在有几个？ & $3 + 2 = ?$ \\
\bottomrule
\end{tabular}
\end{center}

翻译完了。就这么简单。

\section{三、翻译三步法}

每次看到应用题，就做三件事：

\begin{center}
\begin{tabular}{cll}
\toprule
步骤 & 做什么 & 举例 \\
\midrule
\textbf{第一步} & 找数字 & 题目里藏着哪些数字？ \\
\textbf{第二步} & 找关系 & 这些数字之间怎么连接？加？减？乘？除？ \\
\textbf{第三步} & 写出来 & 用数学语言把关系写出来 \\
\bottomrule
\end{tabular}
\end{center}

就这三步。

\section{四、关键词翻译表}

有些词，在应用题里总是反复出现。记住它们，你就掌握了翻译的钥匙。

\begin{center}
\begin{tabular}{cc}
\toprule
中文说 & 数学说 \\
\midrule
一共、总共、合起来 & $+$ \\
还剩、少了、拿走 & $-$ \\
每人、每个、各 & $\times$ \\
平均、分成 & $\div$ \\
是、等于 & $=$ \\
多少、几个、未知 & $x$ \\
\bottomrule
\end{tabular}
\end{center}

\section{五、$x$ 是什么？}

看到表格最后一行，你可能想起来了——$x$，我们之前认识过它。

但在应用题里，$x$ 有了\textbf{新的身份}。

以前，$x$ 只是“藏起来的数”。

但在应用题里，$x$ 代表的是\textbf{题目里你不知道的那个真实存在的东西}。

\begin{center}
\begin{tabular}{ll}
\toprule
题目问的 & $x$ 是什么 \\
\midrule
小明原来有多少钱？ & $x$ 就是那笔钱。\textbf{$x$ 就是钱呐！} \\
一根绳子有多长？ & $x$ 就是绳子的长度。\textbf{$x$ 就是长度呐！} \\
一个班有多少人？ & $x$ 就是那个人数。\textbf{$x$ 就是人数呐！} \\
\bottomrule
\end{tabular}
\end{center}

所以，每次看到应用题，第一件事就是问自己：

\begin{quote}
\textbf{“题目在问我什么？”}

那个“什么”，就是 $x$。
\end{quote}

$x$ 不是随便写的符号。它是你要找的答案的\textbf{代号}。

给它起个代号，是为了让你能把它放进式子里，跟其他数字一起算。

\section{六、来，我们试试}

\subsection{6.1 题目一（热身）}

\begin{quote}
小明有 5 块糖，小红有 3 块糖，他们一共有几块糖？
\end{quote}

\textbf{找数字：} 小明5块，小红3块。

\textbf{找关系：} “一共”$\to$ 加在一起。

\textbf{写出来：}
\[5 + 3 = 8\]

\textbf{答案：} 一共 8 块糖。

这道题没有 $x$，因为数字都给你了，直接算就行。

\subsection{6.2 题目二（还是热身）}

\begin{quote}
一个班有 6 排座位，每排坐 8 个人，全班一共有几个人？
\end{quote}

\textbf{找数字：} 6排，每排8人。

\textbf{找关系：} “每排”$\to$ 乘法。

\textbf{写出来：}
\[6 \times 8 = 48\]

\textbf{答案：} 全班 48 人。

这道题也没有 $x$，数字都给你了，直接翻译就行。

\subsection{6.3 题目三（$x$ 登场了）}

\begin{quote}
小红有一些零花钱，她买了一支 5 块钱的笔，还剩 12 块。她原来有多少零花钱？
\end{quote}

\textbf{找数字：} 花了5块，剩12块。原来有多少？不知道。

\textbf{$x$ 是什么？} 题目在问“原来有多少零花钱”。

\textbf{所以 $x$ 就是小红原来的零花钱。$x$ 就是钱呐！}

为什么用 $x$ 代表它？因为我们不知道它是几，但我们知道它和其他数字的关系。用 $x$ 把它“占个位”，就能把整个关系写出来。

\textbf{找关系：} 原来有 $x$ 块，花了5块，剩12块。“花了”$\to$ 减法。

\textbf{写出来：}
\[x - 5 = 12\]

\textbf{找 $x$：}

怎么把 $x$ 找出来？倒着想——

原来的钱 = 剩下的 + 花掉的

\[x = 12 + 5 = 17\]

\textbf{答案：} 她原来有 17 块零花钱。

\textbf{验证：} $17 - 5 = 12$ \checkmark\ 没错！

\subsection{6.4 题目四（再来一道）}

\begin{quote}
一根绳子，平均剪成 4 段，每段长 3 米。这根绳子原来有多长？
\end{quote}

\textbf{找数字：} 4段，每段3米。原来多长？不知道。

\textbf{$x$ 是什么？} 题目在问“原来有多长”。

\textbf{所以 $x$ 就是这根绳子原来的长度。$x$ 就是长度呐！}

\textbf{找关系：} 平均剪成4段 $\to$ 除法。反过来，原来的长度 $\to$ 乘法。

\textbf{写出来：}
\[x \div 4 = 3\]

\textbf{找 $x$：}
\[x = 3 \times 4 = 12\]

\textbf{答案：} 这根绳子原来 12 米长。

\textbf{验证：} $12 \div 4 = 3$ \checkmark\ 没错！

\section{七、你发现规律了吗？}

\begin{center}
\begin{tabular}{ccl}
\toprule
题目 & 用 $x$ 了吗 & 为什么 \\
\midrule
题目一、二 & 没有 & 所有数字都给你了，直接算 \\
题目三、四 & 用了 & 有个数字藏起来了——就是题目问的那个 \\
\bottomrule
\end{tabular}
\end{center}

\textbf{$x$ 出现的规律只有一个：}

\begin{quote}
\textbf{题目问你什么，$x$ 就是什么。}
\end{quote}

\section{八、现在你来}

这两道题，自己翻译，自己算。

\begin{quote}
\textbf{第一道：}
小明有 15 块钱，买了一个 6 块钱的玩具，还剩多少钱？

提示：这道题所有数字都给你了，不需要 $x$。
\end{quote}

\begin{quote}
\textbf{第二道：}
一个农场有一些鸡，平均分进 3 个笼子，每个笼子里有 9 只鸡。农场一共有几只鸡？

提示：题目在问“一共有几只鸡”，那 $x$ 就是鸡的总数呐！
\end{quote}

写下来。不用紧张。就是找数字，找关系，写出来。

\section{九、本章小结}

应用题，难的从来不是数学，而是\textbf{读题}。

你要学的，不是怎么算，而是怎么\textbf{听懂数学在说什么}。

听懂了，剩下的就是你早就会的东西。

\begin{quote}
\textbf{记住这句话：题目问你什么，$x$ 就是什么。}
\end{quote}

\section*{本章知识地图}

\begin{verbatim}
应用题 = 翻译
│
├─ 翻译三步法
│   ├─ ① 找数字
│   ├─ ② 找关系
│   └─ ③ 写出来
│
├─ 关键词对照
│   ├─ 一共/总共 → +
│   ├─ 还剩/少了 → -
│   ├─ 每个/各 → ×
│   └─ 平均/分成 → ÷
│
└─ x 是什么？
    └─ 题目问什么，x 就是什么
       （而且要喊出来：x 就是钱呐！）
\end{verbatim}

\chapter{专注于未知而产生的应用题}

\section{一、上一章，你学会了翻译}

应用题就是翻译——

中文翻译成数学，数学翻译成中文。

但上一章的题目，都比较“老实”。

数字给你了，关系给你了，你只需要翻译，然后算。

\section{二、这一章，题目学坏了}

有些题目，不是那么老实。

它藏的，不只是一个数字。

它藏的，是\textbf{整个核心}。

这种题目，从头到尾，都在围绕着那个“不知道的东西”转。

这就是\textbf{以未知为核心的应用题}。

\begin{quote}
听起来很厉害？

别慌，其实你已经有所有需要的工具了。
\end{quote}

\section{三、什么叫“以未知为核心”}

上一章的题目，$x$ 只是顺带出现的。

比如：

\begin{quote}
小红花了5块，剩12块，原来有多少？
\end{quote}

这道题，$x$ 只是最后才出现的一个“缺失的数字”。

\textbf{但这一章不一样——}

$x$ 是整道题的\textbf{主角}。

整道题都在描述这个 $x$ 的各种线索。

你要做的，就是把线索收集起来，找到 $x$。

\textbf{就像侦探破案。}

（是的，你现在是数学界的福尔摩斯）

\section{四、破第一个案子}

\begin{quote}
\textbf{案件一：}
有一个数，把它乘以 3，再加上 5，等于 20。这个数是几？
\end{quote}

这道题，从头到尾，都在说“有一个数”。

这个数，就是 $x$。

\textbf{所以 $x$ 就是那个神秘的数呐！}

\subsection{收集线索}

\begin{center}
\begin{tabular}{ll}
\toprule
题目说的话 & 翻译成数学 \\
\midrule
“把它乘以 3” & $x \times 3$ \\
“再加上 5” & $x \times 3 + 5$ \\
“等于 20” & $x \times 3 + 5 = 20$ \\
\bottomrule
\end{tabular}
\end{center}

\subsection{写出方程}

\[3x + 5 = 20\]

\subsection{找 $x$}

先把 $+5$ 去掉（两边同时减5）：

\[3x = 20 - 5 = 15\]

再把 $\times 3$ 去掉（两边同时除以3）：

\[x = 15 \div 3 = 5\]

\textbf{破案了！这个数是 5。}

\subsection{验证}

把答案放回题目里算一遍：

$5 \times 3 + 5 = 15 + 5 = 20$ \checkmark

对上了！

\section{五、为什么要验证？}

你可能发现了，我在最后把答案放回题目里算了一遍。

这叫\textbf{验证}。

\begin{quote}
数学里，算出答案不算完。

把答案放回去，看看对不对，才算真正做完了。
\end{quote}

为啥？

因为有时候你中间可能算错了一步，但你自己不知道。

放回去一验，要是对不上——“哎，有问题！”

这个习惯，从现在开始养成。

（相信我，这个习惯能帮你在考试里捡回很多分）

\section{六、破第二个案子（双主角登场）}

\begin{quote}
\textbf{案件二：}
小明和小红一共有 18 块糖。小明比小红多 4 块。他们各有几块？
\end{quote}

这道题，藏了\textbf{两个未知数}。

还记得上一章说的“双主角”吗？现在，双主角正式登场了。

\textbf{小明有多少糖？} 不知道。$\to$ 叫它 $x$

\textbf{所以 $x$ 就是小明的糖呐！}

\textbf{小红有多少糖？} 不知道。$\to$ 但题目说了，小明比小红多 4 块。

所以小红的糖 = 小明的糖 - 4 = $x - 4$

\textbf{所以 $x - 4$ 就是小红的糖呐！}

\subsection{收集线索}

\begin{center}
\begin{tabular}{cc}
\toprule
角色 & 糖的数量 \\
\midrule
小明 & $x$ \\
小红 & $x - 4$ \\
一共 & 18 块 \\
\bottomrule
\end{tabular}
\end{center}

\subsection{写出方程}

小明的糖 + 小红的糖 = 18

\[x + (x - 4) = 18\]

\subsection{找 $x$}

先去掉括号：

\[x + x - 4 = 18\]

合并同类项（把两个 $x$ 加一起）：

\[2x - 4 = 18\]

把 $-4$ 去掉（两边同时加4）：

\[2x = 18 + 4 = 22\]

把 $\times 2$ 去掉（两边同时除以2）：

\[x = 22 \div 2 = 11\]

\textbf{所以小明有 11 块糖。}

小红呢？$11 - 4 = 7$ 块。

\subsection{验证（两个条件都要查！）}

\begin{itemize}
\item 一共：$11 + 7 = 18$ \checkmark
\item 小明比小红多：$11 - 7 = 4$ \checkmark
\end{itemize}

\textbf{两个条件都对上了，破案！}

\section{七、你发现了什么？}

题目二里，我们只用了\textbf{一个 $x$}，但同时找出了\textbf{两个人}的糖。

怎么做到的？

因为两个未知数之间，\textbf{有关系}。

\begin{quote}
小明比小红多 4 块

$\to$ 小红的糖可以用小明的糖来表示：$x - 4$
\end{quote}

\textbf{这就是以未知为核心的精髓：}

\begin{quote}
\textbf{找到未知数之间的关系，用一个 $x$ 把所有未知数都串起来。}
\end{quote}

（一个 $x$ 打天下，够省事吧？数学家就是这么懒）

\section{八、破第三个案子（这个简单）}

\begin{quote}
\textbf{案件三：}
一列火车，以每小时 60 公里的速度行驶，经过 3 小时，还有 30 公里没走完。这段路一共有多长？
\end{quote}

\subsection{$x$ 是什么？}

题目在问“这段路一共有多长”。

\textbf{所以 $x$ 就是路的总长度呐！}

\subsection{收集线索}

\begin{center}
\begin{tabular}{ll}
\toprule
信息 & 数值 \\
\midrule
速度 & 60 公里/小时 \\
时间 & 3 小时 \\
走了多少 & $60 \times 3 = 180$ 公里 \\
还剩多少 & 30 公里 \\
\bottomrule
\end{tabular}
\end{center}

\subsection{思考一下}

总长度 = 走了的 + 剩下的

\[x = 180 + 30 = 210\]

\textbf{这段路一共 210 公里。}

\subsection{验证}

$210 - 180 = 30$，确实还剩 30 公里 \checkmark

\section{九、咦？这道题没有列方程？}

你可能发现了，题目三直接算出来了，没有像题目一和题目二那样列方程。

对。

因为有些题目，线索已经够清楚了，直接算就行。

\begin{quote}
\textbf{不是每道题都需要列方程。}

$x$ 只是帮你理清思路的工具。

思路清楚了，能直接算就直接算。
\end{quote}

（别把简单的事情搞复杂，这也是数学家的懒人哲学）

\section{十、现在你来当侦探}

试试看，把下面两个案子破了：

\begin{quote}
\textbf{案件四：}
有一个数，把它加上 8，再乘以 2，等于 30。这个数是几？

提示：$x$ 就是那个数呐！“加上8再乘以2”怎么写成数学？
\end{quote}

\begin{quote}
\textbf{案件五：}
哥哥和弟弟一共有 24 本书。哥哥比弟弟多 6 本。他们各有几本？

提示：让 $x$ 当哥哥的书。弟弟比哥哥少 6 本，那弟弟的书怎么用 $x$ 表示？
\end{quote}

写下来，别忘了验证！

（做出来的话，你就是合格的数学侦探了）

\section{十一、本章小结}

以未知为核心的应用题，说穿了就是：

\begin{quote}
\textbf{把题目里所有的线索，}
\textbf{全部转化成关于 $x$ 的关系，}
\textbf{然后找出 $x$。}
\end{quote}

题目给你的每一句话，都是一条线索。

你是侦探，$x$ 是嫌疑人。

把线索收集起来，案子就破了。

\section*{本章知识地图}

\begin{verbatim}
以未知为核心的应用题
│
├─ 和上一章的区别
│   └─ x 是主角，不是配角
│
├─ 破案步骤
│   ├─ ① 确定 x 是什么（喊出来！）
│   ├─ ② 收集线索（每句话都是线索）
│   ├─ ③ 写出方程
│   ├─ ④ 解方程，找到 x
│   └─ ⑤ 验证（放回去算一遍）
│
├─ 双主角怎么办？
│   └─ 找关系，用一个 x 表示两个人
│
└─ 不是每题都要列方程
    └─ 能直接算就直接算
\end{verbatim}

\chapter{按着单词翻译句子}

\section{一、小结：$x$ 的标准写法}

在我们开始之前，先说一件小事。

你可能注意到，前两章我们写的是 $x \times 3$。

但数学家又犯懒了。

他们觉得，$x \times 3$ 里那个乘号，多余。

所以，$x$ 和数字放在一起，乘号直接不写：

\begin{center}
\begin{tabular}{cc}
\toprule
原来的写法 & 偷懒的写法 \\
\midrule
$x \times 3$ & $3x$ \\
$x \times 5$ & $5x$ \\
$x \times 10$ & $10x$ \\
\bottomrule
\end{tabular}
\end{center}

数字写在 $x$ 前面，乘号省掉。

就这样。

以后看到 $3x$，你就知道它的意思是 $x \times 3$。

（数学家能省一笔是一笔）

\section{二、你已经会翻译了}

前两章，你学会了把应用题翻译成数学。

但那些题目，我帮你分析了很多。

这一章，我们换个方式。

我给你一句话，你来翻译。

不分析，不拆解，就像你在学英语的时候，一个单词一个单词地翻译。

\section{三、规则只有一个}

\textbf{看到一个词，翻译一个词。}

不用想太多。就是逐字翻译。

\section{四、先热热身}

我给你示范一遍。

\textbf{句子：} “一个数的两倍，加上 3，等于 11。”

\textbf{逐字翻译：}

\begin{center}
\begin{tabular}{lc}
\toprule
中文 & 数学 \\
\midrule
一个数 & $x$ \\
的两倍 & $\times 2$ \\
加上 3 & $+ 3$ \\
等于 11 & $= 11$ \\
\bottomrule
\end{tabular}
\end{center}

\textbf{合起来：}

\[2x + 3 = 11\]

就这样。

一个词，一个词，拼起来就是数学语言。

\section{五、现在你来翻译}

我给你句子，你翻译成数学。

不用解方程，只需要翻译出来就行。

\subsection{5.1 第一句}

\begin{quote}
“一个数减去 4，等于 9。”
\end{quote}

翻译：

\begin{center}
\begin{tabular}{lc}
\toprule
中文 & 数学 \\
\midrule
一个数 & $x$ \\
减去 4 & $-4$ \\
等于 9 & $=9$ \\
\bottomrule
\end{tabular}
\end{center}

\textbf{答：}

\[x - 4 = 9\]

\subsection{5.2 第二句}

\begin{quote}
“一个数的三倍，等于 21。”
\end{quote}

翻译：

\begin{center}
\begin{tabular}{lc}
\toprule
中文 & 数学 \\
\midrule
一个数 & $x$ \\
的三倍 & $\times 3$ \\
等于 21 & $=21$ \\
\bottomrule
\end{tabular}
\end{center}

\textbf{答：}

\[3x = 21\]

\subsection{5.3 第三句}

\begin{quote}
“一个数加上 6，再除以 2，等于 5。”
\end{quote}

翻译：

\begin{center}
\begin{tabular}{lc}
\toprule
中文 & 数学 \\
\midrule
一个数 & $x$ \\
加上 6 & $+6$ \\
再除以 2 & $\div 2$ \\
等于 5 & $=5$ \\
\bottomrule
\end{tabular}
\end{center}

\textbf{答：}

\[\frac{x + 6}{2} = 5\]

等等，为什么 $x + 6$ 要放在分数上面？

因为“再除以 2”，是把“$x$ 加上 6”\textbf{这整件事}除以 2。

不是只把 6 除以 2。

所以 $x + 6$ 要整个放在上面，再除以 2。

\subsection{5.4 第四句}

\begin{quote}
“两个连续的数，加起来等于 15。”
\end{quote}

这句话稍微难一点。

“连续的数”是什么意思？

就是相差 1 的两个数。比如 3 和 4，7 和 8，100 和 101。

如果第一个数是 $x$，那第二个数就是 $x + 1$。

翻译：

\begin{center}
\begin{tabular}{lc}
\toprule
中文 & 数学 \\
\midrule
两个连续的数 & $x$ 和 $x + 1$ \\
加起来 & $+$ \\
等于 15 & $= 15$ \\
\bottomrule
\end{tabular}
\end{center}

\textbf{合起来：}

\[x + (x + 1) = 15\]

\[2x + 1 = 15\]

\subsection{5.5 第五句}

\begin{quote}
“一个数的一半，比这个数少 6。”
\end{quote}

这句话有点绕。慢慢来。

\begin{center}
\begin{tabular}{lc}
\toprule
中文 & 数学 \\
\midrule
一个数 & $x$ \\
的一半 & $\dfrac{x}{2}$ \\
比这个数少 6 & $= x - 6$ \\
\bottomrule
\end{tabular}
\end{center}

\textbf{合起来：}

\[\frac{x}{2} = x - 6\]

\section{六、你发现了吗？}

翻译，就是把中文里的每一个意思，找到对应的数学符号。

不需要一次看懂整句话。

\begin{quote}
\textbf{一个词，一个词地翻译，拼起来就是答案。}
\end{quote}

\section{七、最后你自己来}

这两句话，自己翻译成数学。不用解，只需要写出式子。

\begin{quote}
\textbf{第一句：} “一个数的四倍，减去 7，等于 13。”
\end{quote}

\begin{quote}
\textbf{第二句：} “一个数加上它的一半，等于 18。”
\end{quote}

写下来。

\section*{本章知识地图}

\begin{verbatim}
逐字翻译法
│
├─ 前置知识：x 的标准写法
│   └─ x × 3 写成 3x（省略乘号）
│
├─ 核心规则
│   └─ 看到一个词，翻译一个词
│
├─ 翻译对照
│   ├─ 一个数 → x
│   ├─ 的N倍 → ×N（写成Nx）
│   ├─ 加上/减去 → +/-
│   ├─ 除以 → 分数线
│   └─ 等于 → =
│
└─ 特殊情况
    ├─ “再除以” → 整体放分数上面
    ├─ “连续的数” → x 和 x+1
    └─ “比……少” → 用减法表示
\end{verbatim}

\chapter{多个不同未知数怎么办？}

\section{一、先想想}

第二章，我们遇到过这种情况：

\begin{quote}
小明和小红一共有 18 块糖。小明比小红多 4 块。
\end{quote}

两个人，两个不知道的数。

但我们只用了一个 $x$，把小红的糖用 $x - 4$ 表示，把两个未知数\textbf{串在一起}了。

那是因为，两个未知数之间有关系。

但是——

\textbf{如果两个未知数之间，没有直接关系呢？}

头大了吧？

别急。

\section{二、这就是今天的问题}

有些题目，有两个完全独立的未知数。

一个 $x$ 不够用了。

怎么办？

很简单。

\textbf{再找一个字母来帮忙。}

\begin{center}
\begin{tabular}{cl}
\toprule
未知数 & 代表谁 \\
\midrule
$x$ & 第一个不知道的东西 \\
$y$ & 第二个不知道的东西 \\
\bottomrule
\end{tabular}
\end{center}

两个侦探，破一个案子。

还记得第二章说的“双主角”吗？现在，双主角正式登场了。

\section{三、但是，新问题来了}

两个未知数，一个方程，不够用。

为什么？

你想想：

\[x + y = 10\]

$x$ 和 $y$ 可以是很多种组合：

\begin{center}
\begin{tabular}{ccc}
\toprule
$x$ & $y$ & 加起来 \\
\midrule
1 & 9 & 10 \checkmark \\
2 & 8 & 10 \checkmark \\
3 & 7 & 10 \checkmark \\
5 & 5 & 10 \checkmark \\
... & ... & ... \\
\bottomrule
\end{tabular}
\end{center}

答案不止一个！

这就好比你问我：“两个人加起来 10 岁，他们各几岁？”

我怎么知道？一个 9 岁一个 1 岁也行，一个 5 岁一个 5 岁也行。

所以，光一个方程不够。

\begin{quote}
\textbf{两个未知数，需要两个方程。}
\end{quote}

两个方程放在一起，叫做\textbf{方程组}。

\section{四、方程组长什么样？}

\[\begin{cases} x + y = 10 \\ x - y = 4 \end{cases}\]

左边那个大括号，意思是：这两个方程\textbf{同时成立}。

$x$ 和 $y$，必须同时满足两个方程。

两个条件都要满足，缺一不可。

\section{五、怎么解？}

现在告诉你一个很妙的方法。

\textbf{把两个方程加起来。}

\[x + y = 10\]
\[x - y = 4\]

两个方程相加：

\[(x + x) + (y + (-y)) = 10 + 4\]

\[2x + 0 = 14\]

\[2x = 14\]

\[x = 7\]

哇，$y$ 不见了！

是不是有点神奇？

\section{六、为什么可以“加起来”？}

你想想天平。

左边放东西，右边放东西，天平平衡了。这就是一个方程。

现在你有两个天平，都是平衡的。

那你把两个天平左边的东西，全部堆到一起放在一个新天平的左边；两个天平右边的东西，也全部堆到一起放在新天平的右边。

那两堆，还是平衡的。

为什么？因为两边各自都加了同样多的东西，平衡没有被破坏。

所以，两个方程，可以直接相加。

\section{七、为什么 $y$ 会消失？}

你看第一个方程里有 $+y$，第二个方程里有 $-y$。

$+y$ 和 $-y$ 加在一起，就是 $y + (-y) = 0$。

$y$ 消失了。

就像你欠了朋友 10 块，朋友又还了你 10 块，两个 10 块一碰，两清了，什么都不剩。

$y$ 消失了，只剩 $x$，$x$ 就好找了。

找到 $x$ 之后，再回头找 $y$。一个一个来。

\section{八、找到 $x$ 之后}

现在把 $x = 7$ 放回第一个方程：

\[7 + y = 10\]

\[y = 10 - 7 = 3\]

所以：$x = 7$，$y = 3$。

\subsection{验证}

\begin{center}
\begin{tabular}{llc}
\toprule
条件 & 检验 & 结果 \\
\midrule
$x + y = 10$ & $7 + 3 = 10$ & \checkmark \\
$x - y = 4$ & $7 - 3 = 4$ & \checkmark \\
\bottomrule
\end{tabular}
\end{center}

完美。

\section{九、但有时候，直接加不管用}

\[\begin{cases} x + y = 10 \\ 2x + y = 16 \end{cases}\]

这两个方程直接加：

\[(x + 2x) + (y + y) = 10 + 16\]

\[3x + 2y = 26\]

$y$ 没消掉。

两个 $y$ 都是正的，加起来变成 $2y$，反而更多了。

这时候怎么办？

\textbf{不加了，改减。}

还是天平的道理。两个天平都是平衡的。把第二个天平左边的东西，拿走和第一个天平左边同样的东西；右边也一样拿走。那两堆，还是平衡的。

用第二个方程减去第一个方程：

\[2x + y = 16\]
\[-(x + y = 10)\]

相减：

\[(2x - x) + (y - y) = 16 - 10\]

\[x + 0 = 6\]

\[x = 6\]

$y$ 消掉了！

把 $x = 6$ 放回第一个方程：

\[6 + y = 10\]

\[y = 4\]

\subsection{验证}

\begin{center}
\begin{tabular}{llc}
\toprule
条件 & 检验 & 结果 \\
\midrule
$x + y = 10$ & $6 + 4 = 10$ & \checkmark \\
$2x + y = 16$ & $2 \times 6 + 4 = 16$ & \checkmark \\
\bottomrule
\end{tabular}
\end{center}

\section{十、消元法的思路}

\begin{center}
\begin{tabular}{cl}
\toprule
步骤 & 做什么 \\
\midrule
\textbf{第一步} & 看看两个方程，哪个未知数容易消掉？ \\
\textbf{第二步} & 用加法或减法，把那个未知数消掉 \\
\textbf{第三步} & 找出第一个未知数 \\
\textbf{第四步} & 放回去，找出第二个未知数 \\
\textbf{第五步} & 验证（别偷懒，这步很重要） \\
\bottomrule
\end{tabular}
\end{center}

\section{十一、来，试试看}

\subsection{11.1 题目一}

\begin{quote}
小明和小红一共有 20 本书。小明比小红多 6 本。他们各有几本？
\end{quote}

$x$ 是小明的书，$y$ 是小红的书。\textbf{$x$ 就是小明的书呐！$y$ 就是小红的书呐！}

\[\begin{cases} x + y = 20 \\ x - y = 6 \end{cases}\]

两个方程相加，$y$ 消掉了：

\[2x = 26\]

\[x = 13\]

把 $x = 13$ 放回第一个方程：

\[13 + y = 20\]

\[y = 7\]

\textbf{答案：} 小明有 \textbf{13} 本，小红有 \textbf{7} 本。

\subsubsection{验证}

\begin{center}
\begin{tabular}{llc}
\toprule
条件 & 检验 & 结果 \\
\midrule
一共 20 本 & $13 + 7 = 20$ & \checkmark \\
小明多 6 本 & $13 - 7 = 6$ & \checkmark \\
\bottomrule
\end{tabular}
\end{center}

\subsection{11.2 题目二}

\begin{quote}
两个数相加等于 15，两倍的第一个数加上第二个数等于 22。这两个数各是几？
\end{quote}

$x$ 是第一个数，$y$ 是第二个数。

\[\begin{cases} x + y = 15 \\ 2x + y = 22 \end{cases}\]

这次直接加，$y$ 消不掉。用第二个减第一个：

\[(2x - x) + (y - y) = 22 - 15\]

\[x = 7\]

把 $x = 7$ 放回第一个方程：

\[7 + y = 15\]

\[y = 8\]

\textbf{答案：} 第一个数是 \textbf{7}，第二个数是 \textbf{8}。

\subsubsection{验证}

\begin{center}
\begin{tabular}{llc}
\toprule
条件 & 检验 & 结果 \\
\midrule
相加等于 15 & $7 + 8 = 15$ & \checkmark \\
两倍加上等于 22 & $2 \times 7 + 8 = 22$ & \checkmark \\
\bottomrule
\end{tabular}
\end{center}

\section{十二、现在你来}

\begin{quote}
\textbf{第一道：}
两个数相加等于 12，两个数相减等于 4。这两个数各是几？
\end{quote}

\begin{quote}
\textbf{第二道：}
两倍的第一个数加上第二个数等于 19，第一个数加上第二个数等于 13。这两个数各是几？
\end{quote}

写下来。别忘了验证。

（我知道你懒得验证，但验证真的很重要。就像打游戏存档，不存档死了就白打了）

\section{十三、本章小结}

多个未知数，不可怕。

\begin{quote}
\textbf{有几个未知数，就需要几个方程。}
\end{quote}

然后，一次消掉一个，一个一个找出来。

就像打游戏，敌人太多了，一个一个干，总会干完的。

\section*{本章知识地图}

\begin{verbatim}
多个未知数
│
├─ 为什么需要多个方程？
│   └─ 一个方程 + 两个未知数 = 无数个答案
│
├─ 方程组
│   └─ 两个方程同时成立（大括号）
│
├─ 消元法
│   ├─ 核心思路：先消掉一个未知数
│   ├─ 加法消元：+y 和 -y 相加得 0
│   └─ 减法消元：+y 和 +y 相减得 0
│
└─ 解题步骤
    ├─ ① 看哪个未知数容易消
    ├─ ② 加法或减法消掉它
    ├─ ③ 找出第一个未知数
    ├─ ④ 放回去找第二个
    └─ ⑤ 验证（很重要！）
\end{verbatim}

\chapter{更复杂的语言翻译}

\section{一、你已经会翻译了}

第三章，你学会了逐字翻译。

一个词，一个词，拼起来就是数学语言。

但那些句子，都比较直接。

这一章，难度升一级。

句子变长了，绕的弯多了，但方法还是一样的。

\textbf{一个词，一个词地翻译。}

别慌。

\section{二、先热热身}

还是先看我示范一遍。

\textbf{句子：} “一个数的三倍，比这个数多 10。”

\textbf{逐字翻译：}

\begin{center}
\begin{tabular}{lc}
\toprule
中文 & 数学 \\
\midrule
一个数 & $x$ \\
的三倍 & $3x$ \\
比这个数多 10 & $= x + 10$ \\
\bottomrule
\end{tabular}
\end{center}

\textbf{合起来：}

\[3x = x + 10\]

看见没？

“比这个数多 10”，意思是在 $x$ 的基础上，再多 10。所以是 $x + 10$。

不是 $10$，是 $x + 10$。

这种绕口的表达，是这一章最需要注意的地方。

\section{三、容易绕晕的表达（记住这张表！）}

先把这些记住，后面翻译的时候，你会感谢我的。

\begin{center}
\begin{tabular}{lcl}
\toprule
中文说 & 数学说 & 怎么记 \\
\midrule
比 $x$ 多 $a$ & $x + a$ & 多，就是加 \\
比 $x$ 少 $a$ & $x - a$ & 少，就是减 \\
$x$ 的 $a$ 倍 & $ax$ & 倍，就是乘 \\
$x$ 的 $\dfrac{1}{a}$ & $\dfrac{x}{a}$ & 几分之一，就是除 \\
$x$ 比 $y$ 大 $a$ & $x = y + a$ & 谁大，谁在大的那边 \\
$x$ 比 $y$ 小 $a$ & $x = y - a$ & 谁小，谁在小的那边 \\
\bottomrule
\end{tabular}
\end{center}

记住这张表，你就掌握了这一章的钥匙。

\section{四、现在，你来翻译}

我给你句子，你先自己翻译，再看答案。

别偷看。

\subsection{4.1 第一句}

\begin{quote}
“一个数的 4 倍，减去 6，等于 18。”
\end{quote}

先自己翻译。

\textbf{答：}

\begin{center}
\begin{tabular}{lc}
\toprule
中文 & 数学 \\
\midrule
一个数 & $x$ \\
的 4 倍 & $4x$ \\
减去 6 & $- 6$ \\
等于 18 & $= 18$ \\
\bottomrule
\end{tabular}
\end{center}

\[4x - 6 = 18\]

\subsection{4.2 第二句}

\begin{quote}
“一个数的一半，比这个数少 3。”
\end{quote}

这句话有点绕。慢慢来。

\textbf{答：}

\begin{center}
\begin{tabular}{lc}
\toprule
中文 & 数学 \\
\midrule
一个数 & $x$ \\
的一半 & $\dfrac{x}{2}$ \\
比这个数少 3 & $= x - 3$ \\
\bottomrule
\end{tabular}
\end{center}

\[\frac{x}{2} = x - 3\]

\subsection{4.3 第三句}

\begin{quote}
“一个数加上它的三分之一，等于 8。”
\end{quote}

注意：“它的三分之一”，是这个数自己的三分之一，不是 $\dfrac{1}{3}$，是 $\dfrac{x}{3}$。

\textbf{答：}

\begin{center}
\begin{tabular}{lc}
\toprule
中文 & 数学 \\
\midrule
一个数 & $x$ \\
加上它的三分之一 & $+ \dfrac{x}{3}$ \\
等于 8 & $= 8$ \\
\bottomrule
\end{tabular}
\end{center}

\[x + \frac{x}{3} = 8\]

\subsection{4.4 第四句（这句有坑）}

\textbf{题目：} “一个数的两倍，比另一个数的三倍少 5，等于 11。”

等等，这句话有点奇怪。

“比另一个数的三倍少 5”？

先别慌。我们分开看：

\begin{itemize}
\item “一个数的两倍” $\to$ $2x$
\item “比另一个数的三倍少 5” $\to$ 这里还有另一个数？
\end{itemize}

但我们在学\textbf{一元}方程，所以“另一个数”其实应该是一个具体的数，不是未知数。

再读一遍题目：

\begin{quote}
“一个数的两倍，比 \textbf{15} 少 5，等于 11。”
\end{quote}

哦，原来是这样。

我刚才故意写错了，看你有没有认真读题。

（开玩笑的，但\textbf{认真读题}这件事，是真的重要）

正确的句子是：“一个数的两倍，减去 5，等于 11。”

\textbf{答：}

\[2x - 5 = 11\]

\subsection{4.5 第五句（括号陷阱）}

\textbf{题目：} “把一个数加上 4，再乘以 3，等于 24。”

\textbf{答：}

\begin{center}
\begin{tabular}{lc}
\toprule
中文 & 数学 \\
\midrule
一个数 & $x$ \\
加上 4 & $x + 4$ \\
再乘以 3 & $(x + 4) \times 3$ \\
等于 24 & $= 24$ \\
\bottomrule
\end{tabular}
\end{center}

\[3(x + 4) = 24\]

\textbf{但是，很多人会这样写：}

\[3x + 4 = 24\]

\textbf{错在哪里？}

“再乘以 3”，是把“这个数加上 4”\textbf{整件事}乘以 3。

不是只把 $x$ 乘以 3，然后再加 4。

所以，$x + 4$ 要用括号包起来，告诉数学：\textbf{这两个是一家人，要一起乘。}

\begin{center}
\begin{tabular}{llc}
\toprule
写法 & 意思 & 对不对 \\
\midrule
$3(x + 4) = 24$ & 先加4，再整体乘3 & 对 \\
$3x + 4 = 24$ & 先乘3，再加4 & 错 \\
\bottomrule
\end{tabular}
\end{center}

看起来差不多，但意思完全不一样。

读题的时候，慢一点。

\subsection{4.6 第六句（次方登场）}

\textbf{题目：} “一个数的平方，加上这个数的两倍，等于 15。”

这句话，出现了\textbf{次方}。

还记得吗？$x$ 的平方就是 $x \times x$，写成 $x^2$。

\textbf{答：}

\begin{center}
\begin{tabular}{lc}
\toprule
中文 & 数学 \\
\midrule
一个数 & $x$ \\
的平方 & $x^2$ \\
加上这个数的两倍 & $+ 2x$ \\
等于 15 & $= 15$ \\
\bottomrule
\end{tabular}
\end{center}

\[x^2 + 2x = 15\]

哇，这个式子看起来复杂多了。

但你是不是发现了，翻译的方法，还是一样的？

\textbf{一个词，一个词地翻译。}

\section{五、现在你来}

这三句话，自己翻译成数学。不用解，只需要写出式子。

\begin{quote}
\textbf{第一句：} “一个数的三倍，加上这个数，等于 20。”
\end{quote}

\begin{quote}
\textbf{第二句：} “把一个数减去 3，再除以 2，等于 5。”
\end{quote}

\begin{quote}
\textbf{第三句：} “一个数的平方，比这个数多 6。”
\end{quote}

写下来。慢慢翻译，不用急。

\section{六、本章小结}

句子越来越长，越来越绕，但翻译的方法，从来没变过。

\textbf{一个词，一个词地翻译。}

遇到绕口的表达，就停下来，把那句话单独拿出来看。

弄清楚它说的是什么，再翻译。

\begin{quote}
\textbf{不要一口气把整句话吞下去。慢慢嚼，才能消化。}
\end{quote}

\section*{本章知识地图}

\begin{verbatim}
更复杂的语言翻译
│
├─ 核心方法（没变）
│   └─ 一个词，一个词地翻译
│
├─ 绕口表达对照表
│   ├─ 比x多a → x + a
│   ├─ 比x少a → x - a
│   ├─ x的a倍 → ax
│   ├─ x的几分之一 → x/a
│   └─ x比y大a → x = y + a
│
├─ 常见陷阱
│   ├─ “再乘以” → 要加括号！
│   ├─ “它的三分之一” → 是 x/3，不是 1/3
│   └─ 认真读题！别被绕进去
│
└─ 新出现的
    └─ 平方 → x²
\end{verbatim}

\chapter{极复杂的语言翻译}

\section{一、你已经走了很远了}

第三章，逐字翻译。
第五章，绕口的句子。

这一章，再升一级。

不是句子更长了，而是——

\textbf{同一个未知数，题目给了你两种情况。}

听起来很绕？

来，我给你看个例子你就懂了。

\section{二、什么叫“两种情况”？}

\textbf{题目：} “小明的年龄，是小红年龄的两倍。三年后，小明比小红大 6 岁。小明现在几岁？”

你发现了吗？

这道题，说了小明年龄\textbf{两次}：

\begin{center}
\begin{tabular}{cl}
\toprule
第几次 & 怎么说的 \\
\midrule
第一次 & 小明的年龄是小红的两倍 \\
第二次 & 三年后，小明比小红大 6 岁 \\
\bottomrule
\end{tabular}
\end{center}

两次描述，说的是\textbf{同一个小明的年龄}。

这就是“两种情况”。

\section{三、怎么处理两种情况？}

\subsection{第一步：把未知数定好}

\begin{center}
\begin{tabular}{ccl}
\toprule
谁 & 现在几岁 & 怎么来的 \\
\midrule
小明 & $x$ & 题目问的，就是 $x$。\textbf{$x$ 就是小明的年龄呐！} \\
小红 & $\dfrac{x}{2}$ & 小明是小红的两倍，所以小红是小明的一半 \\
\bottomrule
\end{tabular}
\end{center}

\subsection{第二步：把两种情况都翻译出来}

第一种情况已经用了，帮我们找到了小红的年龄是 $\dfrac{x}{2}$。

第二种情况：三年后……

\begin{center}
\begin{tabular}{cc}
\toprule
谁 & 三年后几岁 \\
\midrule
小明 & $x + 3$ \\
小红 & $\dfrac{x}{2} + 3$ \\
\bottomrule
\end{tabular}
\end{center}

小明比小红大 6 岁：

\[(x + 3) - \left(\frac{x}{2} + 3\right) = 6\]

\subsection{第三步：解方程}

\[x + 3 - \frac{x}{2} - 3 = 6\]

\[x - \frac{x}{2} = 6\]

\[\frac{x}{2} = 6\]

\[x = 12\]

\textbf{答案：} 小明现在 \textbf{12} 岁。

\subsection{验证}

\begin{center}
\begin{tabular}{llc}
\toprule
检验项 & 计算 & 结果 \\
\midrule
小红现在几岁 & $12 \div 2 = 6$ & 6岁 \\
三年后小明 & $12 + 3 = 15$ & 15岁 \\
三年后小红 & $6 + 3 = 9$ & 9岁 \\
小明比小红大几岁 & $15 - 9 = 6$ & \checkmark \\
\bottomrule
\end{tabular}
\end{center}

没错！

\section{四、再来一道}

\textbf{题目：} “一个长方形，长比宽多 4 厘米。周长是 28 厘米。长和宽各是多少？”

\subsection{第一步：定未知数}

\begin{center}
\begin{tabular}{ccl}
\toprule
什么 & 多少 & 怎么来的 \\
\midrule
宽 & $x$ & 不知道，用 $x$ 表示。\textbf{$x$ 就是宽呐！} \\
长 & $x + 4$ & 长比宽多 4 \\
\bottomrule
\end{tabular}
\end{center}

\subsection{第二步：翻译两种情况}

第一种情况已经用了，帮我们找到了长是 $x + 4$。

第二种情况：周长是 28 厘米。

长方形周长 = 两个长 + 两个宽：

\[2(x + 4) + 2x = 28\]

\subsection{第三步：解方程}

\[2x + 8 + 2x = 28\]

\[4x + 8 = 28\]

\[4x = 20\]

\[x = 5\]

\textbf{答案：} 宽是 \textbf{5} 厘米，长是 $5 + 4 =$ \textbf{9} 厘米。

\subsection{验证}

$2 \times 9 + 2 \times 5 = 18 + 10 = 28$ \checkmark

\section{五、再来一道稍微绕一点的}

\textbf{题目：} “甲和乙两人共有 100 块钱。甲用了自己的一半，乙用了自己的三分之一，两人剩下的钱一共是 60 块。甲原来有多少钱？”

这道题，说的是甲的钱\textbf{两次}：

\begin{center}
\begin{tabular}{cl}
\toprule
第几次 & 怎么说的 \\
\midrule
第一次 & 甲和乙共有 100 块 \\
第二次 & 甲用了一半，乙用了三分之一，剩下共 60 块 \\
\bottomrule
\end{tabular}
\end{center}

\subsection{第一步：定未知数}

\begin{center}
\begin{tabular}{ccl}
\toprule
谁 & 原来有多少 & 怎么来的 \\
\midrule
甲 & $x$ & 题目问的。\textbf{$x$ 就是甲的钱呐！} \\
乙 & $100 - x$ & 甲和乙共 100。\textbf{$100-x$ 就是乙的钱呐！} \\
\bottomrule
\end{tabular}
\end{center}

\subsection{第二步：翻译两种情况}

第一种情况已经用了，帮我们找到了乙的钱是 $100 - x$。

第二种情况：

\begin{center}
\begin{tabular}{ccc}
\toprule
谁 & 用了多少 & 剩下多少 \\
\midrule
甲 & 一半 & $\dfrac{x}{2}$ \\
乙 & 三分之一 & $(100 - x) \times \dfrac{2}{3}$ \\
\bottomrule
\end{tabular}
\end{center}

两人剩下共 60 块：

\[\frac{x}{2} + (100 - x) \times \frac{2}{3} = 60\]

\subsection{第三步：解方程}

\[\frac{x}{2} + \frac{2(100 - x)}{3} = 60\]

两边乘以 6（为什么是6？因为2和3的公倍数是6，乘了就能把分母消掉）：

\[3x + 4(100 - x) = 360\]

\[3x + 400 - 4x = 360\]

\[-x = -40\]

\[x = 40\]

\textbf{答案：} 甲原来有 \textbf{40} 块钱。

\subsection{验证}

\begin{center}
\begin{tabular}{llc}
\toprule
检验项 & 计算 & 结果 \\
\midrule
乙原来有 & $100 - 40 = 60$ & 60块 \\
甲剩下 & $40 \div 2 = 20$ & 20块 \\
乙剩下 & $60 \times \frac{2}{3} = 40$ & 40块 \\
两人剩下共 & $20 + 40 = 60$ & \checkmark \\
\bottomrule
\end{tabular}
\end{center}

\section{六、等等，还有一种情况}

你以为这章讲完了？

还没有。

有时候，你翻译完一道题，会发现式子两边\textbf{都有 $x$}。

别慌。这不是你翻译错了。是题目本来就长这样。

\section{七、为什么两边都有 $x$？}

来看这道题：

\textbf{题目：} “A店的苹果，每个 $x$ 块，买了 3 个，再加上 2 块运费，一共付了 14 块。B店的苹果，每个也是 $x$ 块，买了 2 个，再加上 6 块运费，也是一共付了 14 块。$x$ 是几？”

翻译：

\begin{center}
\begin{tabular}{cl}
\toprule
哪家店 & 翻译成数学 \\
\midrule
A店 & $3x + 2 = 14$ \\
B店 & $2x + 6 = 14$ \\
\bottomrule
\end{tabular}
\end{center}

等等，两个式子右边都是 14。

那两个式子左边也相等：

\[3x + 2 = 2x + 6\]

看见了吗？两边都有 $x$。

这是因为：\textbf{两种买法，花的钱一样，所以两边相等。}

\section{八、两边都有 $x$ 怎么处理？}

还是天平的道理。

两边同时减去 $2x$：

\[3x + 2 - 2x = 2x + 6 - 2x\]

\[x + 2 = 6\]

\[x = 4\]

\textbf{答案：} 苹果每个 \textbf{4} 块钱。

\subsection{验证}

\begin{center}
\begin{tabular}{llc}
\toprule
哪家店 & 计算 & 结果 \\
\midrule
A店 & $3 \times 4 + 2 = 14$ & \checkmark \\
B店 & $2 \times 4 + 6 = 14$ & \checkmark \\
\bottomrule
\end{tabular}
\end{center}

\section{九、再来几道感受一下}

\subsection{9.1 铅笔问题}

\textbf{题目：} “A店的铅笔，每支 $x$ 块，买了 5 支，再加 3 块运费，共花了 18 块。B店同样的铅笔，买了 3 支，再加 9 块运费，也花了 18 块。铅笔每支多少钱？”

翻译：

\[5x + 3 = 3x + 9\]

两边同时减去 $3x$：

\[2x + 3 = 9\]

两边同时减去 3：

\[2x = 6\]

\[x = 3\]

\textbf{答案：} 铅笔每支 \textbf{3} 块钱。

\subsubsection{验证}

\begin{center}
\begin{tabular}{llc}
\toprule
哪家店 & 计算 & 结果 \\
\midrule
A店 & $5 \times 3 + 3 = 18$ & \checkmark \\
B店 & $3 \times 3 + 9 = 18$ & \checkmark \\
\bottomrule
\end{tabular}
\end{center}

\subsection{9.2 买书问题}

\textbf{题目：} “小明买了 $x$ 本书和 4 支笔，书每本 6 块，笔每支 3 块，一共花了 42 块。小红买了 $x$ 本书和 2 支笔，书每本 6 块，笔每支 3 块，一共花了 36 块。$x$ 是几？”

翻译：

\begin{center}
\begin{tabular}{cl}
\toprule
谁 & 翻译成数学 \\
\midrule
小明 & $6x + 12 = 42$ \\
小红 & $6x + 6 = 36$ \\
\bottomrule
\end{tabular}
\end{center}

两个式子左边不一样，不能直接放在一起。

先各自解：

小明：
\[6x = 42 - 12 = 30\]
\[x = 5\]

验证小红：
\[6 \times 5 + 6 = 36\] \checkmark

\textbf{答案：} $x = 5$

\begin{quote}
\textbf{记住：不一定非要两边放在一起。哪种方法快，用哪种。}
\end{quote}

\subsection{9.3 衬衫问题（这道题有坑！）}

\textbf{题目：} “同一件衬衫，A店原价 $x$ 块，打八折后是 40 块。B店原价比A店贵 10 块，打七折后也是 40 块。衬衫原价是多少？”

翻译：

A店：$x \times 0.8 = 40$

所以：
\[x = 40 \div 0.8 = 50\]

A店原价 \textbf{50} 块。

B店原价是 $50 + 10 = 60$ 块。

验证B店：
$60 \times 0.7 = 42$

等等，不是 40！

\textbf{这说明题目有问题。}

（对，我故意出了一道有问题的题，看你验证了没有。\textbf{验证，真的很重要}）

\section{十、什么时候会无解？}

有些时候，两边的 $x$ 消掉了，但剩下的等式不成立。

比如：

\[3x + 5 = 3x + 9\]

两边同时减去 $3x$：

\[5 = 9\]

5 等于 9？当然不对。

这说明：\textbf{不管 $x$ 是几，这个等式都不成立。}

这种情况，叫做\textbf{无解}。

用买东西来理解：

\begin{quote}
“A店每个苹果 $x$ 块，买 3 个加 5 块运费。B店同样的苹果，买 3 个加 9 块运费。两家花的钱一样多。”
\end{quote}

这根本不可能。

同样的苹果，同样买 3 个，运费不一样，怎么可能花一样多？

题目本身就矛盾。所以无解。

\section{十一、两边都有 $x$ 的处理步骤}

\begin{center}
\begin{tabular}{cl}
\toprule
步骤 & 做什么 \\
\midrule
\textbf{第一步} & 把所有 $x$ 移到左边，数字移到右边 \\
\textbf{第二步} & 解出 $x$，验证 \\
\bottomrule
\end{tabular}
\end{center}

\textbf{特殊情况：}

\begin{center}
\begin{tabular}{ll}
\toprule
情况 & 说明 \\
\midrule
$x$ 消掉了，等式成立（如 $0 = 0$） & $x$ 可以是任意数 \\
$x$ 消掉了，等式不成立（如 $5 = 9$） & 无解 \\
\bottomrule
\end{tabular}
\end{center}

\section{十二、现在你来}

\begin{quote}
\textbf{第一道：}
“小明买了 $x$ 个苹果和 3 瓶水，苹果每个 2 块，水每瓶 3 块，一共花了 25 块。小红买了 $x$ 个苹果和 1 瓶水，花了 17 块。$x$ 是几？”
\end{quote}

\begin{quote}
\textbf{第二道：}
“A店的橙子，每个 $x$ 块，买了 4 个加 5 块运费，共 21 块。B店同样的橙子，买了 2 个加 11 块运费，也是 21 块。$x$ 是几？”
\end{quote}

写下来，验证一下。

\section{十三、本章小结}

这一章的题目，看起来很长，很吓人。

但你拆开来看，每一步都是你已经会的东西。

\begin{quote}
\textbf{定未知数 $\to$ 找关系 $\to$ 建方程 $\to$ 解方程 $\to$ 验证}

就这五步。题目再长，也逃不出这五步。
\end{quote}

\section*{本章知识地图}

\begin{verbatim}
极复杂的语言翻译
│
├─ 两种情况的题目
│   ├─ 第一种情况：帮你建立未知数之间的关系
│   └─ 第二种情况：帮你列出方程
│
├─ 两边都有 x
│   ├─ 为什么？两种情况得到相同的结果
│   ├─ 怎么处理？把 x 移到一边，数字移到另一边
│   └─ 特殊情况
│       ├─ x 消掉，等式成立 → 任意解
│       └─ x 消掉，等式不成立 → 无解
│
└─ 解题五步法（万能）
    ├─ ① 定未知数
    ├─ ② 找关系
    ├─ ③ 建方程
    ├─ ④ 解方程
    └─ ⑤ 验证（超级重要！）
\end{verbatim}

\chapter{翻译得到特殊情况，怎么办？}

\section{一、先说清楚，什么叫“特殊情况”}

前面几章，你翻译完，解完，大多数时候都能得到一个漂亮的答案。

比如 $x = 3$，$x = 7$。

干净，利落。

但有时候，解完之后，你会看到这种东西：

\begin{center}
\begin{tabular}{cl}
\toprule
解出来的 & 情况 \\
\midrule
$5 = 9$ & $x$ 不见了，等式还是错的 \\
$6 = 6$ & $x$ 不见了，但等式对了 \\
$x = -3$ & $x$ 变成了负数 \\
$x = 2.5$ & $x$ 变成了小数 \\
\bottomrule
\end{tabular}
\end{center}

这一章，我们把这些情况全部搞清楚。

\section{二、第一种特殊情况：无解}

还记得上一章吗？

\[3x + 5 = 3x + 9\]

两边减去 $3x$：

\[5 = 9\]

$x$ 消掉了，但等式不成立。

这叫\textbf{无解}。

意思是：\textbf{不管 $x$ 是几，这个等式都不可能成立。}

用生活来理解：

\begin{quote}
“小明比小红高 5 厘米，同时小明也比小红高 9 厘米。”
\end{quote}

这根本不可能。一个人不可能同时比另一个人高 5 厘米又高 9 厘米。

题目本身就矛盾了。所以，无解。

\section{三、第二种特殊情况：无数个解}

\[2x + 6 = 2(x + 3)\]

先把右边展开：

\[2x + 6 = 2x + 6\]

两边减去 $2x$：

\[6 = 6\]

$x$ 消掉了，等式成立。

这叫\textbf{无数个解}。

意思是：\textbf{不管 $x$ 是几，这个等式永远成立。}

用生活来理解：

\begin{quote}
“小明买了 $x$ 个苹果和 3 个橙子，苹果和橙子一样贵，每个 2 块。小明花了多少钱？”
\end{quote}

两种算法：

\begin{center}
\begin{tabular}{cc}
\toprule
算法 & 公式 \\
\midrule
第一种 & $2x + 6$ \\
第二种 & $2(x + 3)$ \\
\bottomrule
\end{tabular}
\end{center}

这两种说法，说的是\textbf{同一件事}。

不管 $x$ 是几，两种算法结果永远一样。

所以，$x$ 可以是任何数。

\section{四、第三种特殊情况：答案不符合实际}

这种情况，最容易被忽略。

\[x = -3\]

数学上，$-3$ 是合法的答案。

但如果题目问的是：“一个班有多少个学生？”

学生人数是 $-3$？不可能。人数不能是负数。

所以，这道题\textbf{在实际情况下无解}。

这不是你算错了。是题目的条件，限制了答案的范围。

\section{五、第四种特殊情况：答案不是整数}

\[x = 2.5\]

数学上完全合法。

但如果题目问的是：“买了多少个苹果？”

2.5 个苹果？你把苹果切开了？

这种情况，说明题目本身出了问题，或者你翻译错了。

回去检查一遍。

\section{六、怎么判断翻译对不对？}

做完题，你怎么知道自己翻译对了？

\textbf{只有一个方法：验证。}

把答案放回题目，看看每一个条件都满足了吗？

\begin{center}
\begin{tabular}{cl}
\toprule
验证结果 & 说明 \\
\midrule
全部满足 \checkmark & 翻译正确，答案正确 \\
有一个不满足 & 翻译错了，或者算错了 \\
\bottomrule
\end{tabular}
\end{center}

\section{七、实战第一题（重点：合并同类项）}

\textbf{题目：} “小明有一些糖，他给了小红三分之一，又给了小李四分之一，还剩 10 块糖。小明原来有多少块糖？”

\subsection{第一步：$x$ 是什么？}

题目在问“小明原来有多少块糖”。

\textbf{所以 $x$ 就是小明原来的糖呐！}

\subsection{第二步：翻译每一句话}

\begin{center}
\begin{tabular}{lc}
\toprule
句子 & 翻译 \\
\midrule
给了小红三分之一 & $-\dfrac{x}{3}$ \\
又给了小李四分之一 & $-\dfrac{x}{4}$（注意：是原来的四分之一） \\
还剩 10 块 & $= 10$ \\
\bottomrule
\end{tabular}
\end{center}

\textbf{合起来：}

\[x - \frac{x}{3} - \frac{x}{4} = 10\]

\subsection{第三步：解方程（消掉分数）}

这个方程有分数，看起来很烦。怎么办？

把分数消掉。

\textbf{为什么要找公倍数？}

现在方程里有两个分母：一个是 3，一个是 4。

我们想找一个数，能同时被 3 整除，也能被 4 整除。

为什么要这样？

因为我们想把分母消掉。

\begin{itemize}
\item $\dfrac{x}{3}$ 乘以 3，分母就消掉了，变成 $x$
\item $\dfrac{x}{4}$ 乘以 4，分母就消掉了，变成 $x$
\end{itemize}

但我们只能乘以\textbf{同一个数}。不能左边乘 3，右边乘 4。天平会不平衡。

所以我们需要找一个数，乘上去之后，两个分母都能消掉。

$3 \times 4 = 12$

\begin{center}
\begin{tabular}{lc}
\toprule
检验 & 结果 \\
\midrule
12 除以 3 & 等于 4，能整除 \checkmark \\
12 除以 4 & 等于 3，能整除 \checkmark \\
\bottomrule
\end{tabular}
\end{center}

所以 12 是我们要找的那个数。

这个数，数学上有个名字，叫\textbf{公倍数}（“两个数共同的倍数”）。

\begin{center}
\begin{tabular}{ll}
\toprule
倍数列表 & \\
\midrule
3 的倍数 & 3, 6, 9, \textbf{12}, 15... \\
4 的倍数 & 4, 8, \textbf{12}, 16... \\
\bottomrule
\end{tabular}
\end{center}

12 是第一个同时出现在两个列表里的数。所以我们用 12。

两边同时乘以 12，天平还是平衡的，分母也消掉了。一举两得。

\textbf{左边：}

\[12 \times x - 12 \times \frac{x}{3} - 12 \times \frac{x}{4}\]

一个一个算：

\begin{center}
\begin{tabular}{llc}
\toprule
项 & 计算 & 结果 \\
\midrule
$12 \times x$ & & $12x$ \\
$12 \times \dfrac{x}{3}$ & $\dfrac{12x}{3}$ & $4x$ \\
$12 \times \dfrac{x}{4}$ & $\dfrac{12x}{4}$ & $3x$ \\
\bottomrule
\end{tabular}
\end{center}

所以左边变成：

\[12x - 4x - 3x\]

\textbf{等等，停一下。怎么合并同类项？}

$12x - 4x - 3x$，怎么算？

$12x$ 是什么意思？就是 $x + x + x +...$ （12个$x$加在一起）

$4x$ 是什么意思？就是 $x + x + x + x$（4个$x$加在一起）

$3x$ 是什么意思？就是 $x + x + x$（3个$x$加在一起）

所以：

\begin{center}
\begin{tabular}{llc}
\toprule
步骤 & 理解 & 结果 \\
\midrule
$12x - 4x$ & 12个$x$减掉4个$x$ & $8x$ \\
$8x - 3x$ & 8个$x$减掉3个$x$ & $5x$ \\
\bottomrule
\end{tabular}
\end{center}

\textbf{这个步骤，叫“合并同类项”。}

说白了就是：把一样的东西数一数，加在一起或者减掉。

就像你有 12 个苹果，拿走 4 个，再拿走 3 个，还剩 5 个苹果。$x$ 也是一样。

左边是 $5x$。

右边：$12 \times 10 = 120$

所以：

\[5x = 120\]

\[x = 120 \div 5 = 24\]

\textbf{答案：} 小明原来有 \textbf{24} 块糖。

\subsection{验证}

\begin{center}
\begin{tabular}{llc}
\toprule
条件 & 计算 & 结果 \\
\midrule
给了小红三分之一 & $24 \div 3 = 8$，剩 $24-8=16$ & \checkmark \\
又给了小李四分之一 & $24 \div 4 = 6$，剩 $16-6=10$ & \checkmark \\
还剩 10 块 & $10 = 10$ & \checkmark \\
\bottomrule
\end{tabular}
\end{center}

翻译正确，答案正确！

\section{八、实战第二题（重点：“剩下的”四分之一）}

\textbf{题目：} “一根绳子，剪掉三分之一后，再剪掉剩下的四分之一，还剩 12 米。这根绳子原来多长？”

\subsection{第一步：$x$ 是什么？}

题目在问“这根绳子原来多长”。

\textbf{所以 $x$ 就是绳子原来的长度呐！}

\subsection{第二步：翻译每一句话}

\textbf{“剪掉三分之一后”}

绳子原来有 $x$ 米，剪掉三分之一，剩下：

\[x - \frac{x}{3}\]

\textbf{“再剪掉剩下的四分之一”}

注意这里！

“剩下的四分之一”，是\textbf{剩下部分}的四分之一，不是原来绳子的四分之一。

剩下的是 $x - \dfrac{x}{3}$，所以再剪掉的是：

\[\left(x - \frac{x}{3}\right) \times \frac{1}{4}\]

\textbf{“还剩 12 米”}

\[x - \frac{x}{3} - \left(x - \frac{x}{3}\right) \times \frac{1}{4} = 12\]

\subsection{第三步：先化简括号里的东西}

$x - \dfrac{x}{3}$ 是什么？

$x$ 就是 $\dfrac{3x}{3}$（为什么？因为 $\dfrac{3x}{3} = x$，只是换了个写法）

\[\frac{3x}{3} - \frac{x}{3} = \frac{3x - x}{3} = \frac{2x}{3}\]

所以剩下的是 $\dfrac{2x}{3}$。

\subsection{第四步：把化简结果放回方程}

\[\frac{2x}{3} - \frac{2x}{3} \times \frac{1}{4} = 12\]

$\dfrac{2x}{3} \times \dfrac{1}{4}$ 是什么？

分子乘分子，分母乘分母：

\[\frac{2x}{3} \times \frac{1}{4} = \frac{2x \times 1}{3 \times 4} = \frac{2x}{12} = \frac{x}{6}\]

所以：

\[\frac{2x}{3} - \frac{x}{6} = 12\]

\subsection{第五步：消掉分数}

分母有 3 和 6，公倍数是 6。两边同时乘以 6。

\begin{center}
\begin{tabular}{llc}
\toprule
项 & 计算 & 结果 \\
\midrule
$6 \times \dfrac{2x}{3}$ & $\dfrac{12x}{3}$ & $4x$ \\
$6 \times \dfrac{x}{6}$ & $\dfrac{6x}{6}$ & $x$ \\
$6 \times 12$ & & $72$ \\
\bottomrule
\end{tabular}
\end{center}

左边：$4x - x = 3x$

右边：$72$

\[3x = 72\]

\[x = 72 \div 3 = 24\]

\textbf{答案：} 绳子原来 \textbf{24} 米长。

\subsection{验证}

\begin{center}
\begin{tabular}{llc}
\toprule
条件 & 计算 & 结果 \\
\midrule
剪掉三分之一后 & $24 \div 3 = 8$，剩 $24-8=16$ & \checkmark \\
再剪掉剩下的四分之一 & $16 \div 4 = 4$，剩 $16-4=12$ & \checkmark \\
还剩 12 米 & $12 = 12$ & \checkmark \\
\bottomrule
\end{tabular}
\end{center}

翻译正确，答案正确！

\section{九、验证的真正意义}

看见了吗？

验证，不只是检查算术对不对。

验证，是检查\textbf{翻译对不对}。

\begin{center}
\begin{tabular}{ll}
\toprule
出错类型 & 怎么办 \\
\midrule
算术错了 & 重新算就行 \\
翻译错了 & 回去重新读题，重新翻译 \\
\bottomrule
\end{tabular}
\end{center}

所以，\textbf{验证是最重要的一步。}

\section{十、特殊情况总结}

翻译完解完，遇到特殊情况：

\begin{center}
\begin{tabular}{lll}
\toprule
情况 & 意思 & 怎么办 \\
\midrule
$x$ 消掉，等式不成立（$5=9$） & 无解 & 题目条件矛盾 \\
$x$ 消掉，等式成立（$6=6$） & 无数个解 & $x$ 是任意数 \\
答案是负数 & 实际无解 & 检查题目条件 \\
答案是小数 & 可能翻译错了 & 回去验证 \\
验证不通过 & 翻译错了 & 重新读题，重新翻译 \\
\bottomrule
\end{tabular}
\end{center}

\section{十一、现在你来}

\begin{quote}
\textbf{第一道：}
“一袋糖，小明吃了三分之一，小红吃了剩下的一半，还剩 6 颗。这袋糖原来有多少颗？”

先翻译，再解，再验证。
\end{quote}

\begin{quote}
\textbf{第二道：}
“两个数，第一个数是第二个数的两倍，两个数加起来等于第一个数的两倍。这两个数各是几？”

提示：这道题有点特殊，仔细看看解完之后是什么情况。
\end{quote}

\section{十二、本章小结}

翻译，是这几章最核心的技能。

\begin{center}
\begin{tabular}{ll}
\toprule
结果 & 说明 \\
\midrule
翻译对了 & 剩下的就是计算 \\
翻译错了 & 算得再准也没用 \\
\bottomrule
\end{tabular}
\end{center}

所以：

\begin{quote}
\textbf{慢慢读题 $\to$ 逐词翻译 $\to$ 解方程 $\to$ 验证}

验证不过 $\to$ 回去重新翻译
\end{quote}

就这四步。

\section*{本章知识地图}

\begin{verbatim}
特殊情况
│
├─ 四种特殊情况
│   ├─ 无解（x消掉，等式错）
│   ├─ 无数解（x消掉，等式对）
│   ├─ 答案是负数（不符合实际）
│   └─ 答案是小数（可能翻译错）
│
├─ 重要技能
│   ├─ 公倍数（消分母）
│   ├─ 合并同类项（把x数一数）
│   └─ 化简括号（特别注意“剩下的”）
│
└─ 验证的意义
    ├─ 不只检查算术
    └─ 检查翻译对不对
\end{verbatim}

\chapter{翻译总结}

\section{一、好了，到了证明自己的时候了}

前面那么多章，我一直在帮你。

给你示范，给你提示，给你答案。

但现在——

\textbf{我要撒手了。}

这一章，8道题，全部靠你自己翻译，自己解。

我只给你词。

组成什么句子，你来决定。

（别怕，题目很简单。我又不是要害你）

准备好了吗？那我们开始。

\section{二、第一道（热身）}

\textbf{题目：} “一个数加上 8，等于 20。这个数是几？”

（这道题要是不会，我建议你从第一章重新看起。开玩笑的，但认真翻译哦。）

\subsection{词表}

\begin{center}
\begin{tabular}{lc}
\toprule
中文 & 数学 \\
\midrule
一个数 & \_\_\_\_\_\_\_ \\
加上 8 & \_\_\_\_\_\_\_ \\
等于 20 & \_\_\_\_\_\_\_ \\
\bottomrule
\end{tabular}
\end{center}

\subsection{组合成方程}

\[\_\_\_ = \_\_\_\]

\subsection{验算}

\_\_\_\_\_（填入 $x$ 的值）+ \_\_\_\_\_（填入 8）= \textbf{20}

\textbf{最终答案 = 12} \checkmark

\section{三、第二道（稍微动脑子）}

\textbf{题目：} “一个数的 3 倍，减去 6，等于 15。这个数是几？”

（上一道热完身了吗？来，稍微动动脑子。）

\subsection{词表}

\begin{center}
\begin{tabular}{lc}
\toprule
中文 & 数学 \\
\midrule
一个数 & \_\_\_\_\_\_\_ \\
的 3 倍 & \_\_\_\_\_\_\_ \\
减去 6 & \_\_\_\_\_\_\_ \\
等于 15 & \_\_\_\_\_\_\_ \\
\bottomrule
\end{tabular}
\end{center}

\subsection{组合成方程}

\[\_\_\_ = \_\_\_\]

\subsection{验算}

\_\_\_\_\_（填入 $3x$ 的值）- \_\_\_\_\_（填入 6）= \textbf{15}

\textbf{最终答案 = 7} \checkmark

\section{四、第三道（双主角）}

\textbf{题目：} “小明和小红一共有 30 本书。小明比小红多 4 本。他们各有几本？”

（两个未知数！别慌，你学过这个，翻回第四章看看。）

\subsection{词表}

\begin{center}
\begin{tabular}{lc}
\toprule
中文 & 数学 \\
\midrule
小明的书 & $x$ \\
小红的书 & \_\_\_\_\_\_\_ \\
一共 30 本 & \_\_\_\_\_\_\_ \\
小明比小红多 4 本 & \_\_\_\_\_\_\_ \\
\bottomrule
\end{tabular}
\end{center}

\subsection{组合成方程组}

\[\begin{cases} \_\_\_ \\ \_\_\_ \end{cases}\]

\subsection{验算}

\_\_\_\_\_（填入 $x$ 的值）+ \_\_\_\_\_（填入 $y$ 的值）= \textbf{30}

\_\_\_\_\_（填入 $x$ 的值）- \_\_\_\_\_（填入 $y$ 的值）= \textbf{4}

\textbf{小明最终答案 = 17} \checkmark

\textbf{小红最终答案 = 13} \checkmark

\section{五、第四道（有点绕）}

\textbf{题目：} “一个数的一半，比这个数少 5。这个数是几？”

（这道题有点绕。慢慢读，别一口气吞下去。）

\subsection{词表}

\begin{center}
\begin{tabular}{lc}
\toprule
中文 & 数学 \\
\midrule
一个数 & $x$ \\
的一半 & \_\_\_\_\_\_\_ \\
比这个数少 5 & \_\_\_\_\_\_\_ \\
\bottomrule
\end{tabular}
\end{center}

\subsection{组合成方程}

\[\_\_\_ = \_\_\_\]

\subsection{验算}

\_\_\_\_\_（填入 $\dfrac{x}{2}$ 的值）= \_\_\_\_\_（填入 $x - 5$ 的值）= \textbf{5}

\textbf{最终答案 = 10} \checkmark

\section{六、第五道（注意括号！）}

\textbf{题目：} “把一个数加上 3，再乘以 4，等于 28。这个数是几？”

（注意括号！“整件事一起乘”，别把括号漏掉了，不然我会顺着书找到你。开玩笑的。）

\subsection{词表}

\begin{center}
\begin{tabular}{lc}
\toprule
中文 & 数学 \\
\midrule
一个数 & $x$ \\
加上 3 & \_\_\_\_\_\_\_ \\
再乘以 4（整件事一起乘） & \_\_\_\_\_\_\_ \\
等于 28 & \_\_\_\_\_\_\_ \\
\bottomrule
\end{tabular}
\end{center}

\subsection{组合成方程}

\[\_\_\_ = \_\_\_\]

\subsection{验算}

\_\_\_\_\_（填入 $x + 3$ 的值）$\times$ \_\_\_\_\_（填入 4）= \textbf{28}

\textbf{最终答案 = 4} \checkmark

\section{七、第六道（两边都有 $x$）}

\textbf{题目：} “A店的橙子，每个 $x$ 块，买了 4 个加 3 块运费，共 19 块。B店同样的橙子，买了 3 个加 7 块运费，也是 19 块。橙子每个多少钱？”

（两边都有 $x$！还记得怎么处理吗？把 $x$ 移到同一边就行了。）

\subsection{词表}

\begin{center}
\begin{tabular}{lc}
\toprule
中文 & 数学 \\
\midrule
A店：4个橙子加运费 & \_\_\_\_\_\_\_ \\
B店：3个橙子加运费 & \_\_\_\_\_\_\_ \\
两家花的钱一样 & \_\_\_\_\_\_\_ \\
\bottomrule
\end{tabular}
\end{center}

\subsection{组合成方程}

\[\_\_\_ = \_\_\_\]

\subsection{验算}

\textbf{A店：} \_\_\_\_\_（填入 $4x$ 的值）+ \_\_\_\_\_（填入 3）= \textbf{19}

\textbf{B店：} \_\_\_\_\_（填入 $3x$ 的值）+ \_\_\_\_\_（填入 7）= \textbf{19}

\textbf{最终答案 = 4} \checkmark

\section{八、第七道（分数来了）}

\textbf{题目：} “小明有一些糖，他给了小红四分之一，还剩 12 块。小明原来有多少块糖？”

（分数出现了！别慌，把分母消掉就行。你已经做过很多次了。）

\subsection{词表}

\begin{center}
\begin{tabular}{lc}
\toprule
中文 & 数学 \\
\midrule
小明原来的糖 & $x$ \\
给了小红四分之一 & \_\_\_\_\_\_\_ \\
还剩 12 块 & \_\_\_\_\_\_\_ \\
\bottomrule
\end{tabular}
\end{center}

\subsection{组合成方程}

\[\_\_\_ = \_\_\_\]

\subsection{验算}

\_\_\_\_\_（填入 $x$ 的值）- \_\_\_\_\_（填入 $\dfrac{x}{4}$ 的值）= \textbf{12}

\textbf{最终答案 = 16} \checkmark

\section{九、第八道（最后一道了！）}

\textbf{题目：} “一个长方形，长比宽多 6 厘米。周长是 36 厘米。长和宽各是多少？”

（最后一道了！两种情况，记得用上。做完了就可以去玩了。）

\subsection{词表}

\begin{center}
\begin{tabular}{lc}
\toprule
中文 & 数学 \\
\midrule
宽 & $x$ \\
长比宽多 6 & \_\_\_\_\_\_\_ \\
周长 = 两个长加两个宽 & \_\_\_\_\_\_\_ \\
等于 36 & \_\_\_\_\_\_\_ \\
\bottomrule
\end{tabular}
\end{center}

\subsection{组合成方程}

\[\_\_\_ = \_\_\_\]

\subsection{验算}

\_\_\_\_\_（填入长的值）$\times$ 2 + \_\_\_\_\_（填入宽的值）$\times$ 2 = \textbf{36}

\textbf{宽的最终答案 = 6} \checkmark

\textbf{长的最终答案 = 12} \checkmark

\section{十、对了吗？}

对了的话，太棒了。

要是有一两道没对……没关系，翻回去看看，找找是哪里翻译错了。

不丢人，这才是学数学的样子。

从第一章的“翻译三步法”，到现在翻译复杂的应用题，你走了很远。

而且是自己走的。

这才是真正学会数学的样子——不是背公式，不是刷题，是真的懂了，真的会用了。

\section*{本章知识地图}

\begin{verbatim}
翻译总结（8道综合练习）
│
├─ 第1-2道：基础翻译
│   └─ 一个数 + 加减乘除 + 等于
│
├─ 第3道：双主角
│   └─ 两个未知数，用关系串起来
│
├─ 第4道：绕口表达
│   └─ “比...少” 这类表达
│
├─ 第5道：括号陷阱
│   └─ “整件事一起乘” 要加括号
│
├─ 第6道：两边都有 x
│   └─ 把 x 移到同一边
│
├─ 第7道：分数
│   └─ 用公倍数消分母
│
└─ 第8道：两种情况
    └─ 长宽关系 + 周长公式
\end{verbatim}

\chapter{等式化简的所有方式与规律（逐步详解）}

\section{等式的基本性质（化简的根基）}

\subsection{性质1：等式两边同时加（或减）同一个数，等式仍然成立}

\textbf{例子：} $x + 5 = 12$

\begin{itemize}
\item 目标：把 $x$ 单独留在左边
\item 操作：两边同时\textbf{减去5}
\item 左边：$x + 5 - 5 = x$
\item 右边：$12 - 5 = 7$
\item 结果：$x = 7$
\end{itemize}

\textbf{为什么这样做？} 因为 $x$ 旁边“多了一个+5”，要消掉它，就做它的“反操作”——减5。但你不能只减左边，右边也必须减，否则天平就不平了。

\subsection{性质2：等式两边同时乘以（或除以）同一个不为零的数，等式仍然成立}

\textbf{例子：} $3x = 15$

\begin{itemize}
\item 目标：把 $x$ 单独留在左边
\item 操作：两边同时\textbf{除以3}
\item 左边：$3x \div 3 = x$
\item 右边：$15 \div 3 = 5$
\item 结果：$x = 5$
\end{itemize}

\textbf{再看一个例子：} $x/4 = 3$

\begin{itemize}
\item 目标：$x$ 被4除着，要消掉这个“$\div 4$”
\item 操作：两边同时\textbf{乘以4}
\item 左边：$(x/4) \times 4 = x$
\item 右边：$3 \times 4 = 12$
\item 结果：$x = 12$
\end{itemize}

\section{移项（最常用的化简手段）}

\subsection{核心口诀：移过等号，变符号}

把某一项从等号一边搬到另一边时，\textbf{加变减，减变加，乘变除，除变乘}。

\subsection{详细步骤演示}

\textbf{例子：} $2x + 3 = 11$

第一步：我想把“$+3$”移到右边去
\begin{itemize}
\item $+3$ 移过等号，变成 $-3$
\item 左边剩下：$2x$
\item 右边变成：$11 - 3 = 8$
\item 现在等式变成：$2x = 8$
\end{itemize}

第二步：我想把“$\times 2$”（就是2乘以$x$的意思）移到右边去
\begin{itemize}
\item $\times 2$ 移过等号，变成 $\div 2$
\item 左边剩下：$x$
\item 右边变成：$8 \div 2 = 4$
\item 结果：$x = 4$
\end{itemize}

\textbf{本质：移项其实就是等式性质1和性质2的“简写”。}

\section{合并同类项}

\subsection{什么是同类项？}

含有\textbf{相同字母}且\textbf{字母的指数也相同}的项。

比如：$3x$ 和 $5x$ 是同类项（都含$x$，指数都是1）

比如：$2x^2$ 和 $7x^2$ 是同类项（都含$x^2$）

但：$3x$ 和 $2x^2$ \textbf{不是}同类项（指数不同）

而：$3x$ 和 $5$ \textbf{不是}同类项（一个有$x$一个没有）

\subsection{怎么合并？}

\textbf{系数相加减，字母和指数不变。}

\textbf{例子：} $3x + 5x = ?$

\begin{itemize}
\item 第一步：它们都含$x$，是同类项，可以合并
\item 第二步：系数是3和5，相加得8
\item 第三步：字母部分$x$不变
\item 结果：$8x$
\end{itemize}

\textbf{再看一个：} $7x - 2x + 4 = 19$

\begin{itemize}
\item 第一步：$7x$ 和 $2x$ 是同类项，合并
\item $7x - 2x = 5x$（系数$7-2=5$）
\item 第二步：等式变成 $5x + 4 = 19$
\item 第三步：把$+4$移到右边，变$-4$
\item $5x = 19 - 4 = 15$
\item 第四步：把$\times 5$移到右边，变$\div 5$
\item $x = 15 \div 5 = 3$
\end{itemize}

\section{去括号}

\subsection{规则1：括号前面是“$+$”号}

直接去掉括号，\textbf{里面所有项的符号不变}。

\textbf{例子：} $+(3x - 2) = 3x - 2$

\begin{itemize}
\item 括号前是加号
\item 里面 $+3x$ 还是 $+3x$
\item 里面 $-2$ 还是 $-2$
\item 直接写：$3x - 2$
\end{itemize}

\subsection{规则2：括号前面是“$-$”号}

去掉括号，\textbf{里面所有项的符号都要变}（加变减，减变加）。

\textbf{例子：} $-(3x - 2) = ?$

\begin{itemize}
\item 括号前是减号
\item 里面 $+3x$ 变成 $-3x$
\item 里面 $-2$ 变成 $+2$
\item 结果：$-3x + 2$
\end{itemize}

\subsection{规则3：括号前面是一个数字（乘法分配律）}

这个数字要\textbf{分别乘以}括号里面的\textbf{每一项}。

\textbf{例子：} $3(2x + 5) = ?$

\begin{itemize}
\item 第一步：$3 \times 2x = 6x$
\item 第二步：$3 \times 5 = 15$
\item 结果：$6x + 15$
\end{itemize}

\textbf{更复杂的例子：} $-2(x - 4) = ?$

\begin{itemize}
\item 第一步：$-2 \times x = -2x$
\item 第二步：$-2 \times (-4) = +8$（负负得正）
\item 结果：$-2x + 8$
\end{itemize}

\section{去分母}

\subsection{什么时候用？}

等式中出现\textbf{分数}的时候，为了让计算更方便，把分母消掉。

\subsection{怎么做？}

等式两边\textbf{同时乘以所有分母的最小公倍数}。

\textbf{例子：} $\dfrac{x}{2} + \dfrac{x}{3} = 5$

第一步：找分母 $\to$ 2 和 3

第二步：找最小公倍数 $\to$ 6

第三步：等式两边同时乘以6

\begin{itemize}
\item 左边第一项：$(x/2) \times 6 = 3x$

怎么算的？$6 \div 2 = 3$，然后 $3 \times x = 3x$

\item 左边第二项：$(x/3) \times 6 = 2x$

怎么算的？$6 \div 3 = 2$，然后 $2 \times x = 2x$

\item 右边：$5 \times 6 = 30$
\end{itemize}

第四步：等式变成 $3x + 2x = 30$

第五步：合并同类项 $\to$ $5x = 30$

第六步：两边除以5 $\to$ $x = 6$

\section{交换律和结合律（调整顺序来化简）}

\subsection{加法交换律}

$a + b = b + a$，项的顺序可以随意交换（但要\textbf{带着前面的符号一起走}）。

\textbf{例子：} $3 + 5x - 2 + x$

\begin{itemize}
\item 我想把含$x$的放一起，把数字放一起
\item 重新排列：$5x + x + 3 - 2$
\item 合并：$6x + 1$
\end{itemize}

\subsection{乘法交换律}

$a \times b = b \times a$

\textbf{例子：} $5 \times x \times 2 = ?$

\begin{itemize}
\item 重新排列：$5 \times 2 \times x = 10x$
\end{itemize}

\section{消元（含多个未知数时）}

\subsection{加减消元法}

\textbf{例子：}

等式①：$x + y = 10$

等式②：$x - y = 4$

第一步：观察发现，①中有 $+y$，②中有 $-y$，如果两个等式\textbf{相加}，$y$就被消掉了

第二步：两式相加
\begin{itemize}
\item 左边：$(x + y) + (x - y) = x + y + x - y = 2x$（$y$被消掉了！）
\item 右边：$10 + 4 = 14$
\end{itemize}

第三步：$2x = 14$，所以 $x = 7$

第四步：把 $x = 7$ 代入①：$7 + y = 10$，所以 $y = 3$

\subsection{代入消元法}

\textbf{例子：}

等式①：$y = 2x + 1$

等式②：$3x + y = 11$

第一步：①已经告诉我们 $y$ 等于什么了

第二步：把 “$2x + 1$” 代替 ② 中的 $y$
\begin{itemize}
\item $3x + (2x + 1) = 11$
\end{itemize}

第三步：去括号
\begin{itemize}
\item $3x + 2x + 1 = 11$
\end{itemize}

第四步：合并同类项
\begin{itemize}
\item $5x + 1 = 11$
\end{itemize}

第五步：移项
\begin{itemize}
\item $5x = 10$
\end{itemize}

第六步：两边除以5
\begin{itemize}
\item $x = 2$
\end{itemize}

第七步：代回①：$y = 2\times 2 + 1 = 5$

\section{因式分解（提公因式）}

\subsection{它和去括号是什么关系？}

去括号是把 $3(2x + 3)$ \textbf{展开}成 $6x + 9$（从“乘积”变成“和”）。

因式分解正好\textbf{反过来}，是把 $6x + 9$ \textbf{收回}成 $3(2x + 3)$（从“和”变成“乘积”）。

\begin{quote}
你可以把因式分解理解为：\textbf{把括号装回去。}
\end{quote}

\subsection{怎么找公因数？分两步看}

\textbf{第一步：看数字（系数）}
\begin{itemize}
\item 把每一项的系数列出来
\item 找它们的\textbf{最大公因数}
\end{itemize}

\textbf{第二步：看字母}
\begin{itemize}
\item 看哪些字母是\textbf{每一项都有}的
\item 如果都有，取\textbf{指数最小}的那个
\end{itemize}

数字公因数 $\times$ 字母公因数 $=$ 你要提出来的东西

\subsection{详细例子1（只有数字公因数）}

\textbf{题目：} $6x + 9 = 0$

第一步：看系数 $\to$ 6 和 9

第二步：6 和 9 的最大公因数是多少？
\begin{itemize}
\item $6 = 2 \times 3$
\item $9 = 3 \times 3$
\item 它们共有的最大因数是 \textbf{3}
\end{itemize}

第三步：把每一项里的 3 “抽出来”，看剩下什么
\begin{itemize}
\item $6x$ 里面抽走 3，剩下 $2x$（因为 $6x \div 3 = 2x$）
\item $9$ 里面抽走 3，剩下 $3$（因为 $9 \div 3 = 3$）
\end{itemize}

第四步：写成“提出来的 $\times$ 括号里剩下的”
\begin{itemize}
\item $6x + 9 = \mathbf{3(2x + 3)}$
\end{itemize}

第五步：验证一下——把括号重新展开，看对不对
\begin{itemize}
\item $3 \times 2x = 6x$ \checkmark
\item $3 \times 3 = 9$ \checkmark
\item 合起来就是 $6x + 9$ \checkmark\ 没错！
\end{itemize}

第六步：现在等式变成 $3(2x + 3) = 0$

第七步：$3$ 不等于 $0$，所以只能是 $2x + 3 = 0$

第八步：移项 $\to$ $2x = -3$

第九步：两边除以2 $\to$ $x = -3/2$

\subsection{详细例子2（公因数里包含字母）}

\textbf{题目：} $4x^2 + 6x = 0$

第一步：看系数 $\to$ 4 和 6
\begin{itemize}
\item 4 和 6 的最大公因数是 \textbf{2}
\end{itemize}

第二步：看字母
\begin{itemize}
\item $4x^2$ 含有 $x^2$（也就是 $x \times x$）
\item $6x$ 含有 $x$
\item 两项\textbf{都有 $x$}，取指数小的那个 $\to$ \textbf{$x$}
\end{itemize}

第三步：所以要提出来的公因式是 \textbf{$2x$}

第四步：每一项除以 $2x$，看剩下什么
\begin{itemize}
\item $4x^2 \div 2x = 2x$（数字 $4\div 2=2$，字母 $x^2\div x=x$，合起来 $2x$）
\item $6x \div 2x = 3$（数字 $6\div 2=3$，字母 $x\div x=1$，合起来就是 $3$）
\end{itemize}

第五步：写出来
\begin{itemize}
\item $4x^2 + 6x = \mathbf{2x(2x + 3)}$
\end{itemize}

第六步：验证
\begin{itemize}
\item $2x \times 2x = 4x^2$ \checkmark
\item $2x \times 3 = 6x$ \checkmark
\end{itemize}

第七步：等式变成 $2x(2x + 3) = 0$

第八步：乘积等于0，说明其中至少一个等于0
\begin{itemize}
\item 要么 $2x = 0$ $\to$ $x = 0$
\item 要么 $2x + 3 = 0$ $\to$ $x = -3/2$
\end{itemize}

\subsection{详细例子3（三项的情况）}

\textbf{题目：} $12x^3 + 8x^2 - 4x = 0$

第一步：看系数 $\to$ 12、8、4
\begin{itemize}
\item 最大公因数是 \textbf{4}
\end{itemize}

第二步：看字母
\begin{itemize}
\item $x^3$、$x^2$、$x$ 都含有 $x$，取最小指数 $\to$ \textbf{$x$}
\end{itemize}

第三步：公因式是 \textbf{$4x$}

第四步：每一项除以 $4x$
\begin{itemize}
\item $12x^3 \div 4x = 3x^2$
\item $8x^2 \div 4x = 2x$
\item $-4x \div 4x = -1$（\textbf{注意！除完是1，不是0，不能丢掉！}）
\end{itemize}

第五步：写出来
\begin{itemize}
\item $12x^3 + 8x^2 - 4x = \mathbf{4x(3x^2 + 2x - 1)}$
\end{itemize}

第六步：验证
\begin{itemize}
\item $4x \times 3x^2 = 12x^3$ \checkmark
\item $4x \times 2x = 8x^2$ \checkmark
\item $4x \times (-1) = -4x$ \checkmark
\end{itemize}

\subsection{提公因式时最容易犯的错}

\begin{center}
\begin{tabular}{|l|l|}
\hline
\textbf{错误} & \textbf{正确做法} \\
\hline
除完得1时写成0，把那一项丢掉 & 比如 $4x \div 4x = \mathbf{1}$，不是0，括号里要写上这个1 \\
\hline
只看了数字没看字母 & 字母也要找共同的，取指数最小的 \\
\hline
没有提“最大”公因数 & 比如$6x+9$，提3才是提干净，提1虽然不算错但等于没提 \\
\hline
\end{tabular}
\end{center}

\section{约分（化简分式）}

\textbf{例子：} $\dfrac{6x}{9} = ?$

第一步：6和9的最大公因数是3

第二步：分子分母同时除以3
\begin{itemize}
\item 分子：$6x \div 3 = 2x$
\item 分母：$9 \div 3 = 3$
\item 结果：$\dfrac{2x}{3}$
\end{itemize}

\section{乘方与开方的逆运算}

\subsection{遇到平方，用开方来化简}

\textbf{例子：} $x^2 = 25$

第一步：两边同时开平方

第二步：$x = +5$ 或 $x = -5$（别忘了有两个答案！）

为什么？因为 $5\times 5 = 25$，$(-5)\times(-5) = 25$，两个都对

\subsection{遇到根号，用平方来化简}

\textbf{例子：} $\sqrt{x} = 3$

第一步：两边同时平方

第二步：$x = 9$

\section{完整化简流程总结（做题顺序）}

面对一个复杂的等式，按以下顺序操作：

\begin{verbatim}
第1步：去分母（如果有分数 → 两边乘最小公倍数）
        ↓
第2步：去括号（乘法分配律、符号变化规则）
        ↓
第3步：移项（含未知数的移到一边，常数移到另一边，过等号变符号）
        ↓
第4步：合并同类项（把同类的加减合并在一起）
        ↓
第5步：系数化为1（两边除以未知数前面的系数）
        ↓
        得到答案
\end{verbatim}

\section{完整例题演示}

\textbf{题目：} $\dfrac{2x + 1}{3} - \dfrac{x - 1}{2} = 1$

\subsection{第1步：去分母}

分母是3和2，最小公倍数是6，两边乘6：

\begin{itemize}
\item 左边第一项：$\dfrac{2x+1}{3} \times 6 = 2(2x+1)$

算法：$6\div 3=2$，所以乘以2

\item 左边第二项：$\dfrac{x-1}{2} \times 6 = 3(x-1)$

算法：$6\div 2=3$，所以乘以3

\item 右边：$1 \times 6 = 6$
\end{itemize}

等式变成：$2(2x+1) - 3(x-1) = 6$

\subsection{第2步：去括号}

\begin{itemize}
\item $2(2x+1) \to 2\times 2x + 2\times 1 = 4x + 2$
\item $-3(x-1) \to -3\times x + (-3)\times(-1) = -3x + 3$
\end{itemize}

等式变成：$4x + 2 - 3x + 3 = 6$

\subsection{第3步：移项}

含$x$的留左边，常数移到右边：

\[4x - 3x = 6 - 2 - 3\]

\subsection{第4步：合并同类项}

\begin{itemize}
\item 左边：$4x - 3x = x$
\item 右边：$6 - 2 - 3 = 1$
\end{itemize}

等式变成：$x = 1$

\subsection{第5步：系数已经是1}

\textbf{答案：$x = 1$}

\section{练习题（附逐步解答）}

\subsection{练习1：$3x + 6 = 15$}

\begin{quote}
\textbf{用到的方法：移项、系数化为1}
\end{quote}

第一步：观察等式，左边有 $3x$ 和 $+6$，我想把 $x$ 留在左边，所以要把 $+6$ 赶到右边去

第二步：移项，$+6$ 移过等号变成 $-6$
\begin{itemize}
\item 左边剩下：$3x$
\item 右边变成：$15 - 6$
\item 等式变成：$3x = 9$
\end{itemize}

第三步：现在 $x$ 前面还乘着一个3，要把它消掉，两边同时除以3
\begin{itemize}
\item 左边：$3x \div 3 = x$
\item 右边：$9 \div 3 = 3$
\item 等式变成：$x = 3$
\end{itemize}

第四步：验证——把 $x = 3$ 代回原式检查
\begin{itemize}
\item 左边：$3 \times 3 + 6 = 9 + 6 = 15$
\item 右边：$15$
\item 左边 $=$ 右边 \checkmark
\end{itemize}

\textbf{答案：$x = 3$}

\subsection{练习2：$5x - 2 = 3x + 8$}

\begin{quote}
\textbf{用到的方法：移项、合并同类项、系数化为1}
\end{quote}

第一步：观察等式，两边都有含 $x$ 的项，我想把所有含 $x$ 的都集中到左边，所有数字都集中到右边

第二步：先把右边的 $+3x$ 移到左边，过等号变成 $-3x$
\begin{itemize}
\item 等式变成：$5x - 2 - 3x = 8$
\end{itemize}

第三步：再把左边的 $-2$ 移到右边，过等号变成 $+2$
\begin{itemize}
\item 等式变成：$5x - 3x = 8 + 2$
\end{itemize}

第四步：两边分别合并同类项
\begin{itemize}
\item 左边：$5x - 3x = 2x$（系数 $5 - 3 = 2$，字母 $x$ 不变）
\item 右边：$8 + 2 = 10$
\item 等式变成：$2x = 10$
\end{itemize}

第五步：两边同时除以2
\begin{itemize}
\item 左边：$2x \div 2 = x$
\item 右边：$10 \div 2 = 5$
\item 等式变成：$x = 5$
\end{itemize}

第六步：验证——把 $x = 5$ 代回原式
\begin{itemize}
\item 左边：$5 \times 5 - 2 = 25 - 2 = 23$
\item 右边：$3 \times 5 + 8 = 15 + 8 = 23$
\item 左边 $=$ 右边 \checkmark
\end{itemize}

\textbf{答案：$x = 5$}

\subsection{练习3：$4(x - 1) = 12$}

\begin{quote}
\textbf{用到的方法：去括号、移项、系数化为1}
\end{quote}

第一步：左边有括号，先去括号，用乘法分配律
\begin{itemize}
\item $4 \times x = 4x$
\item $4 \times (-1) = -4$
\item 左边展开后：$4x - 4$
\item 等式变成：$4x - 4 = 12$
\end{itemize}

第二步：把左边的 $-4$ 移到右边，过等号变成 $+4$
\begin{itemize}
\item 等式变成：$4x = 12 + 4 = 16$
\end{itemize}

第三步：两边同时除以4
\begin{itemize}
\item 左边：$4x \div 4 = x$
\item 右边：$16 \div 4 = 4$
\item 等式变成：$x = 4$
\end{itemize}

第四步：验证——把 $x = 4$ 代回原式
\begin{itemize}
\item 左边：$4 \times (4 - 1) = 4 \times 3 = 12$
\item 右边：$12$
\item 左边 $=$ 右边 \checkmark
\end{itemize}

\textbf{答案：$x = 4$}

\subsection{练习4：$x/3 + 1 = 5$}

\begin{quote}
\textbf{用到的方法：移项、去分母}
\end{quote}

第一步：把左边的 $+1$ 移到右边，过等号变成 $-1$
\begin{itemize}
\item 等式变成：$x/3 = 5 - 1 = 4$
\end{itemize}

第二步：现在 $x$ 被3除着，要消掉这个“$\div 3$”，两边同时乘以3
\begin{itemize}
\item 左边：$(x/3) \times 3 = x$（分母3和乘的3抵消了）
\item 右边：$4 \times 3 = 12$
\item 等式变成：$x = 12$
\end{itemize}

第三步：验证——把 $x = 12$ 代回原式
\begin{itemize}
\item 左边：$12/3 + 1 = 4 + 1 = 5$
\item 右边：$5$
\item 左边 $=$ 右边 \checkmark
\end{itemize}

\textbf{答案：$x = 12$}

\subsection{练习5：$6x^2 + 9x = 0$（用提公因式）}

\begin{quote}
\textbf{用到的方法：因式分解（提公因式）}
\end{quote}

第一步：观察两项 $6x^2$ 和 $9x$，找公因数
\begin{itemize}
\item 看数字：6 和 9 的最大公因数是 3
\item 看字母：$x^2$ 和 $x$ 都含有 $x$，取指数最小的，就是 $x$
\item 所以公因式是 \textbf{$3x$}
\end{itemize}

第二步：把每一项除以 $3x$，看括号里剩下什么
\begin{itemize}
\item $6x^2 \div 3x = 2x$

怎么算的？数字：$6 \div 3 = 2$；字母：$x^2 \div x = x$；合起来就是 $2x$

\item $9x \div 3x = 3$

怎么算的？数字：$9 \div 3 = 3$；字母：$x \div x = 1$；合起来就是 $3$
\end{itemize}

第三步：写成“提出来的 $\times$ 括号里剩下的”
\begin{itemize}
\item $6x^2 + 9x = 3x(2x + 3)$
\item 等式变成：$3x(2x + 3) = 0$
\end{itemize}

第四步：两个东西相乘，结果等于0，那么\textbf{至少有一个必须等于0}
\begin{itemize}
\item 情况一：$3x = 0$

两边除以3 $\to$ $x = 0$

\item 情况二：$2x + 3 = 0$

移项 $\to$ $2x = -3$

两边除以2 $\to$ $x = -3/2$
\end{itemize}

第五步：验证两个答案

验证 $x = 0$：
\begin{itemize}
\item 左边：$6 \times 0^2 + 9 \times 0 = 0 + 0 = 0$
\item 右边：$0$
\item 左边 $=$ 右边 \checkmark
\end{itemize}

验证 $x = -3/2$：
\begin{itemize}
\item 左边：$6 \times (-3/2)^2 + 9 \times (-3/2)$
\item 先算 $(-3/2)^2 = 9/4$
\item $6 \times 9/4 = 54/4 = 27/2$
\item $9 \times (-3/2) = -27/2$
\item 合起来：$27/2 + (-27/2) = 0$
\item 右边：$0$
\item 左边 $=$ 右边 \checkmark
\end{itemize}

\textbf{答案：$x = 0$ 或 $x = -3/2$}

\section{易错点提醒}

\begin{center}
\begin{tabular}{|l|l|}
\hline
\textbf{易错点} & \textbf{正确做法} \\
\hline
移项忘记变号 & $+3$移过去必须变$-3$ \\
\hline
去括号前面是负号时里面忘记全部变号 & $-(a-b) = -a+b$，两项都要变 \\
\hline
去分母时漏乘某一项 & \textbf{每一项}都要乘，包括右边的常数 \\
\hline
$x^2=25$ 只写$x=5$ & 别忘了 $x=-5$ 也是答案 \\
\hline
两边除以0 & \textbf{永远不能除以0} \\
\hline
提公因式时，除完得1的项直接丢掉 & 除完是\textbf{1}不是0，必须保留在括号里 \\
\hline
\end{tabular}
\end{center}

以上就是等式化简的\textbf{所有核心方法}，每一种都是从“等式两边做相同操作”这个根本原则出发的。掌握了这些，任何等式化简都能应对！

\chapter{认识数学里的“家族成员”}

\section{你已经认识很多数字了}

从第一章到现在，你见过：

\begin{itemize}
\item $1, 2, 3, 4...$（数数用的）
\item $-1, -2, -3...$（负数）
\item $\dfrac{1}{2}, \dfrac{3}{4}...$（分数）
\item $0.5, 2.5...$（小数）
\end{itemize}

但你可能发现了，数学家很喜欢给这些数字\textbf{起名字}。

今天，我们就来认识这些“家族成员”。

别紧张，都是你早就见过的东西，只是换个名字而已。

\section{偶数和奇数（双胞胎兄弟）}

\subsection{什么是偶数？}

能被2整除的数，就叫\textbf{偶数}。

\begin{center}
\begin{tabular}{|c|c|c|}
\hline
\textbf{数字} & \textbf{能被2整除吗} & \textbf{是偶数吗} \\
\hline
2 & $2 \div 2 = 1$ & \checkmark\ 是 \\
\hline
4 & $4 \div 2 = 2$ & \checkmark\ 是 \\
\hline
6 & $6 \div 2 = 3$ & \checkmark\ 是 \\
\hline
10 & $10 \div 2 = 5$ & \checkmark\ 是 \\
\hline
\end{tabular}
\end{center}

\textbf{偶数家族}：$2, 4, 6, 8, 10, 12, 14...$

\begin{quote}
\textbf{怎么认出偶数？}

看最后一位数字是不是 $0, 2, 4, 6, 8$。

比如：$256$，最后一位是$6$，是偶数。
\end{quote}

\subsection{什么是奇数？}

不能被2整除的数，就叫\textbf{奇数}。

\begin{center}
\begin{tabular}{|c|c|c|}
\hline
\textbf{数字} & \textbf{除以2的结果} & \textbf{是奇数吗} \\
\hline
1 & $1 \div 2 = 0.5$（不是整数） & \checkmark\ 是 \\
\hline
3 & $3 \div 2 = 1.5$（不是整数） & \checkmark\ 是 \\
\hline
5 & $5 \div 2 = 2.5$（不是整数） & \checkmark\ 是 \\
\hline
9 & $9 \div 2 = 4.5$（不是整数） & \checkmark\ 是 \\
\hline
\end{tabular}
\end{center}

\textbf{奇数家族}：$1, 3, 5, 7, 9, 11, 13...$

\begin{quote}
\textbf{怎么认出奇数？}

看最后一位数字是不是 $1, 3, 5, 7, 9$。

比如：$357$，最后一位是$7$，是奇数。
\end{quote}

\subsection{用生活理解}

想象你有一些苹果，要分给2个人，每人分得一样多。

\begin{center}
\begin{tabular}{|c|c|c|}
\hline
\textbf{苹果数量} & \textbf{能平均分吗} & \textbf{类型} \\
\hline
2个 & 能（每人1个） & 偶数 \\
\hline
3个 & 不能（会剩1个） & 奇数 \\
\hline
4个 & 能（每人2个） & 偶数 \\
\hline
5个 & 不能（会剩1个） & 奇数 \\
\hline
\end{tabular}
\end{center}

\begin{itemize}
\item \textbf{能平均分的，就是偶数}
\item \textbf{不能平均分的，就是奇数}
\end{itemize}

\section{质数（孤独的数）}

\subsection{什么是质数？}

只能被 $1$ 和它自己整除的数，叫\textbf{质数}。

\begin{center}
\begin{tabular}{|c|l|c|}
\hline
\textbf{数字} & \textbf{能被哪些数整除} & \textbf{是质数吗} \\
\hline
2 & 只能被1和2整除 & \checkmark\ 是 \\
\hline
3 & 只能被1和3整除 & \checkmark\ 是 \\
\hline
4 & 能被1、2、4整除 & $\times$ 不是 \\
\hline
5 & 只能被1和5整除 & \checkmark\ 是 \\
\hline
6 & 能被1、2、3、6整除 & $\times$ 不是 \\
\hline
7 & 只能被1和7整除 & \checkmark\ 是 \\
\hline
\end{tabular}
\end{center}

\textbf{质数家族}：$2, 3, 5, 7, 11, 13, 17, 19, 23...$

\begin{quote}
\textbf{为什么叫“孤独的数”？}

因为它们只能被1和自己整除，没有其他“朋友”。
\end{quote}

\subsection{特殊情况}

\begin{center}
\begin{tabular}{|c|c|l|}
\hline
\textbf{数字} & \textbf{是质数吗} & \textbf{为什么} \\
\hline
1 & $\times$ 不是 & 规定：质数要有两个不同的约数 \\
\hline
2 & \checkmark\ 是 & 唯一的偶数质数 \\
\hline
\end{tabular}
\end{center}

\section{合数（热闹的数）}

\subsection{什么是合数？}

除了1和它自己，还能被其他数整除的数，叫\textbf{合数}。

\begin{center}
\begin{tabular}{|c|l|c|}
\hline
\textbf{数字} & \textbf{能被哪些数整除} & \textbf{是合数吗} \\
\hline
4 & 1、2、4 & \checkmark\ 是 \\
\hline
6 & 1、2、3、6 & \checkmark\ 是 \\
\hline
8 & 1、2、4、8 & \checkmark\ 是 \\
\hline
9 & 1、3、9 & \checkmark\ 是 \\
\hline
\end{tabular}
\end{center}

\textbf{合数家族}：$4, 6, 8, 9, 10, 12, 14, 15...$

\begin{quote}
\textbf{质数和合数的区别？}

\begin{itemize}
\item 质数：只有2个约数（1和自己）
\item 合数：有3个或更多约数
\end{itemize}
\end{quote}

\section{因数和倍数（躲猫猫游戏）}

前面说“能被谁整除”“有多少个约数”，其实都在偷偷说两个词：

\textbf{因数} 和 \textbf{倍数}。

\subsection{什么是因数？}

如果一个数 $a$ 能被另一个数 $b$ 整除（没有余数），那我们就说：\textbf{$b$ 是 $a$ 的因数（也叫约数）}。

举例：看 12。

\begin{center}
\begin{tabular}{|c|c|c|}
\hline
\textbf{能不能整除12？} & \textbf{结果} & \textbf{是12的因数吗} \\
\hline
$12 \div 1$ & $=12$ & \checkmark\ 是 \\
\hline
$12 \div 2$ & $=6$ & \checkmark\ 是 \\
\hline
$12 \div 3$ & $=4$ & \checkmark\ 是 \\
\hline
$12 \div 4$ & $=3$ & \checkmark\ 是 \\
\hline
$12 \div 5$ & $=2.4$（不是整数） & $\times$ 不是 \\
\hline
$12 \div 6$ & $=2$ & \checkmark\ 是 \\
\hline
$12 \div 12$ & $=1$ & \checkmark\ 是 \\
\hline
\end{tabular}
\end{center}

所以：

\[12\ \text{的因数有：}\ 1, 2, 3, 4, 6, 12\]

\begin{quote}
\textbf{一句话记住因数：}

能把这个数整除的数，就是它的因数。
\end{quote}

质数为什么“孤独”？因为它\textbf{只有两个因数}：1 和它自己。合数为什么“热闹”？因为它有很多因数，朋友一大堆。

\subsection{什么是倍数？}

如果一个数 $a$ 是另一个数 $b$ 和某个整数 $k$ 的乘积：

\[a = b \times k \quad (k \text{ 是整数})\]

那我们就说：\textbf{$a$ 是 $b$ 的倍数}。

举例：还是看 12。

\[12 = 3 \times 4\]

\begin{itemize}
\item 对 12 来说：能写成 $3 \times 4$，所以 \textbf{3 和 4 是 12 的因数}。
\item 从 3 的角度：12 是它的 $4$ 倍，所以 \textbf{12 是 3 的倍数}。
\item 从 4 的角度：12 是它的 $3$ 倍，所以 \textbf{12 是 4 的倍数}。
\end{itemize}

再看 5 的倍数：

\[5, 10, 15, 20, 25, 30, \dots\]

它们都可以写成：

\[5 \times 1,\ 5 \times 2,\ 5 \times 3,\ 5 \times 4,\ 5 \times 5,\ 5 \times 6,\ \dots\]

\begin{quote}
\textbf{一句话记住倍数：}

是“某个数 $\times$ 整数”得到的，就是这个数的倍数。
\end{quote}

\subsection{因数和倍数是“反过来的”}

用 12 和 3 来看：

\begin{itemize}
\item 从 12 的角度：3 能整除我，所以 \textbf{3 是 12 的因数}。
\item 从 3 的角度：12 是我乘以 4 得到的，所以 \textbf{12 是 3 的倍数}。
\end{itemize}

\[\boxed{\text{3 是 12 的因数} \quad\Longleftrightarrow\quad \text{12 是 3 的倍数}}\]

这就是它俩的关系：\textbf{一正一反，看你站在哪一边}。

\begin{quote}
小总结：
\begin{itemize}
\item “谁能整除我？”——我在找我的\textbf{因数}。
\item “我能变成谁的几倍？”——我在找我的\textbf{倍数}。
\end{itemize}
\end{quote}

\section{系数（$x$ 前面的那个数）}

\subsection{什么是系数？}

在 $3x$、$5x$、$10x$ 这样的式子里，$x$ 前面的数字，叫\textbf{系数}。

\begin{center}
\begin{tabular}{|c|c|l|}
\hline
\textbf{式子} & \textbf{系数是几} & \textbf{意思} \\
\hline
$3x$ & 3 & $x$ 的3倍 \\
\hline
$5x$ & 5 & $x$ 的5倍 \\
\hline
$x$ & 1 & $x$ 的1倍（1通常不写） \\
\hline
$-2x$ & $-2$ & $x$ 的$-2$倍 \\
\hline
\end{tabular}
\end{center}

\begin{quote}
\textbf{为什么要知道系数？}

因为消元的时候，我们要让\textbf{系数变成一样}。

比如：$3x$ 的系数是3，$2x$ 的系数是2。要让系数一样，就把第一个方程乘以2，第二个方程乘以3。
\end{quote}

\subsection{举个例子}

\[\begin{cases}
3x + y = 10 \\
2x + y = 7
\end{cases}\]

\begin{center}
\begin{tabular}{|c|c|l|c|}
\hline
\textbf{方程} & $x$ \textbf{的系数} & \textbf{操作} & \textbf{新的} $x$ \textbf{系数} \\
\hline
第一个 & 3 & 乘以2 & 6 \\
\hline
第二个 & 2 & 乘以3 & 6 \\
\hline
\end{tabular}
\end{center}

现在系数一样了，可以消元了！

\section{常数（不会变的数）}

\subsection{什么是常数？}

不带 $x$ 或 $y$ 的数字，叫\textbf{常数}。

\begin{center}
\begin{tabular}{|l|c|}
\hline
\textbf{式子} & \textbf{哪些是常数} \\
\hline
$3x + 5$ & 5 是常数 \\
\hline
$2x + 3y + 7$ & 7 是常数 \\
\hline
$x - 4 = 10$ & 4 和 10 都是常数 \\
\hline
\end{tabular}
\end{center}

\begin{quote}
\textbf{为什么叫“常数”？}

因为它\textbf{一直是那个数}，不会变。

$x$ 可能是1、2、3...会变，但5就是5，永远不变。
\end{quote}

\section{小测验（看看你记住了吗）}

\subsection{判断题}

\begin{enumerate}
\item 8 是偶数（\quad）
\item 7 是质数（\quad）
\item 在 $5x$ 里，系数是 5（\quad）
\item 在 $x + 3 = 7$ 里，3 和 7 都是常数（\quad）
\end{enumerate}

\textbf{答案：}

\begin{enumerate}
\item \checkmark\ 对（8能被2整除）
\item \checkmark\ 对（7只能被1和7整除）
\item \checkmark\ 对（$x$前面的数字是5）
\item \checkmark\ 对（3和7都不带$x$）
\end{enumerate}

\subsection{填空题}

\begin{enumerate}
\item 12 的最后一位是 2，所以 12 是 \underline{\hspace{2cm}} 数。
\item 11 只能被 1 和 11 整除，所以 11 是 \underline{\hspace{2cm}} 数。
\item 在 $7x$ 里，系数是 \underline{\hspace{2cm}}。
\end{enumerate}

\textbf{答案：}

\begin{enumerate}
\item 偶数
\item 质数
\item 7
\end{enumerate}

\section{为什么要认识这些“家族”？}

你可能会问：记住这些名字有什么用？

\subsection{用处一：让交流更方便}

\begin{center}
\begin{tabular}{|l|l|}
\hline
\textbf{不用这些词} & \textbf{用这些词} \\
\hline
“那个能被2整除的数” & “偶数” \\
\hline
“$x$前面的那个数字” & “系数” \\
\hline
“不会变的那个数” & “常数” \\
\hline
“能把12整除的那些数” & “12的因数” \\
\hline
“是5乘整数得到的数” & “5的倍数” \\
\hline
\end{tabular}
\end{center}

说话更简洁，更清楚。

\subsection{用处二：帮助你发现规律}

\begin{center}
\begin{tabular}{|l|l|}
\hline
\textbf{规律} & \textbf{说明} \\
\hline
两个偶数相加，结果是偶数 & $2+4=6$，$6+8=14$ \\
\hline
两个奇数相加，结果是偶数 & $3+5=8$，$7+9=16$ \\
\hline
质数只有两个约数 & 帮你快速判断能不能分解 \\
\hline
所有偶数都是2的倍数 & 可以写成 $2 \times n$ \\
\hline
合数都有不止两个因数 & 能拆成很多种乘法 \\
\hline
\end{tabular}
\end{center}

\subsection{用处三：解题更轻松}

当题目说“$x$是偶数”，你就知道：

\begin{itemize}
\item $x$ 的最后一位一定是 $0, 2, 4, 6, 8$
\item $x$ 能被2整除
\item $x$ 可能是 $2, 4, 6, 8, 10...$
\end{itemize}

当题目说“12的因数”，你不用现算，直接想到：$1,2,3,4,6,12$。

当题目说“5的倍数”，脑子里自动闪过：$5,10,15,20...$

这些信息能帮你更快解题。

\section{本章小结}

今天认识了数字的“家族成员”：

\begin{center}
\begin{tabular}{|l|l|l|}
\hline
\textbf{名字} & \textbf{意思} & \textbf{怎么认出来} \\
\hline
\textbf{偶数} & 能被2整除的数 & 最后一位是0、2、4、6、8 \\
\hline
\textbf{奇数} & 不能被2整除的数 & 最后一位是1、3、5、7、9 \\
\hline
\textbf{质数} & 只能被1和自己整除 & 只有2个约数 \\
\hline
\textbf{合数} & 除了1和自己，还能被其他数整除 & 有3个或更多约数 \\
\hline
\textbf{因数} & 能整除某个数的数 & 比如12的因数：1,2,3,4,6,12 \\
\hline
\textbf{倍数} & 是某个数$\times$整数得到的数 & 比如5的倍数：5,10,15,20... \\
\hline
\textbf{系数} & $x$前面的数字 & 在$5x$里，系数是5 \\
\hline
\textbf{常数} & 不带$x$或$y$的数 & 不会变的数 \\
\hline
\end{tabular}
\end{center}

记住这些名字，以后看题会更轻松。

不用刻意背，多见几次就熟了。

\subsection*{本章知识地图}

\begin{verbatim}
数字的家族
│
├─ 按能否被2整除
│   ├─ 偶数（2、4、6、8...）
│   └─ 奇数（1、3、5、7...）
│
├─ 按约数多少
│   ├─ 质数（2、3、5、7、11...）
│   ├─ 合数（4、6、8、9、10...）
│   └─ 特殊：1既不是质数也不是合数
│
├─ 因数和倍数
│   ├─ 因数：能把一个数整除的数
│   └─ 倍数：某个数 × 整数 得到的数
│
└─ 在式子里的角色
    ├─ 系数（x前面的数）
    └─ 常数（不带x或y的数）
\end{verbatim}

\chapter{两个未知数——更复杂的语言翻译}

\section{你已经会了一个未知数}

前面几章，你把一个未知数的应用题，翻译得很熟了。

但生活里，很多问题，不止一个未知数。

这一章，两个未知数，更复杂的句子，但方法还是一样的。

\textbf{一个词，一个词地翻译。}

别慌。

\section{两个未知数的翻译，难在哪里？}

一个未知数的时候，你只需要找一个 $x$。

两个未知数的时候，你需要找 $x$ 和 $y$，而且，还需要找\textbf{两个方程}。

\begin{center}
\begin{tabular}{|c|l|}
\hline
\textbf{难点} & \textbf{说明} \\
\hline
难不在计算 & 你早就会解方程了 \\
\hline
难在读题 & 要从题目里找出两个条件 \\
\hline
\end{tabular}
\end{center}

读懂了，剩下的你都会。

\section{容易绕晕的表达（记住这张表！）}

\begin{center}
\begin{tabular}{|l|c|l|}
\hline
\textbf{中文说} & \textbf{数学说} & \textbf{注意} \\
\hline
两个数的和 & $x + y$ & 加在一起 \\
\hline
两个数的差 & $x - y$ & 大的减小的 \\
\hline
第一个数是第二个数的 $a$ 倍 & $x = ay$ & 谁是谁的倍数，别搞反 \\
\hline
两个数的积 & $x \times y$ & 乘在一起 \\
\hline
两个数相差 $a$ & $x - y = a$ & 差，就是减法 \\
\hline
两个数的和是 $a$，差是 $b$ & $x+y=a$ 且 $x-y=b$ & 两个条件，两个方程 \\
\hline
\end{tabular}
\end{center}

\section{先看我示范}

\textbf{题目：} “两个数，大的数是小的数的 3 倍，两个数相加等于 20。这两个数各是几？”

\subsection{第一步：定未知数}

\begin{center}
\begin{tabular}{|c|c|l|}
\hline
\textbf{什么} & \textbf{用什么表示} & \textbf{怎么来的} \\
\hline
小的数 & $x$ & 不知道，用$x$表示 \\
\hline
大的数 & $3x$ & 是小的数的3倍 \\
\hline
\end{tabular}
\end{center}

\textbf{所以 $x$ 就是小的数呐，$3x$ 就是大的数呐！}

\subsection{第二步：翻译两个条件}

\begin{center}
\begin{tabular}{|l|c|l|}
\hline
\textbf{条件} & \textbf{用了吗} & \textbf{翻译成什么} \\
\hline
大的数是小的数的3倍 & \checkmark & 帮我们找到了大的数是$3x$ \\
\hline
两个数相加等于20 & \checkmark & $x + 3x = 20$ \\
\hline
\end{tabular}
\end{center}

\subsection{第三步：解方程}

\[4x = 20\]
\[x = 5\]

\textbf{答案：} 小的数是 \textbf{5}，大的数是 $3 \times 5 =$ \textbf{15}。

\subsection{验证}

\begin{center}
\begin{tabular}{|l|l|c|}
\hline
\textbf{条件} & \textbf{检验} & \textbf{结果} \\
\hline
两个数相加等于20 & $5 + 15 = 20$ & \checkmark \\
\hline
大的数是小的数的3倍 & $15 = 3 \times 5$ & \checkmark \\
\hline
\end{tabular}
\end{center}

\section{两个未知数没有直接关系怎么办？}

\textbf{题目：} “两个数，相加等于 18，相减等于 4。这两个数各是几？”

这道题，两个未知数之间，没有直接关系。

一个 $x$ 不够用了。

\subsection{定未知数}

\begin{center}
\begin{tabular}{|c|c|}
\hline
\textbf{什么} & \textbf{用什么表示} \\
\hline
第一个数 & $x$ \\
\hline
第二个数 & $y$ \\
\hline
\end{tabular}
\end{center}

\subsection{两个条件，两个方程}

\[\begin{cases} x + y = 18 \\ x - y = 4 \end{cases}\]

\subsection{用消元法}

两个方程相加：

\[(x + x) + (y + (-y)) = 18 + 4\]
\[2x = 22\]
\[x = 11\]

把 $x = 11$ 放回第一个方程：

\[11 + y = 18\]
\[y = 7\]

\textbf{答案：} 第一个数是 \textbf{11}，第二个数是 \textbf{7}。

\subsection{验证}

\begin{center}
\begin{tabular}{|l|l|c|}
\hline
\textbf{条件} & \textbf{检验} & \textbf{结果} \\
\hline
相加等于18 & $11 + 7 = 18$ & \checkmark \\
\hline
相减等于4 & $11 - 7 = 4$ & \checkmark \\
\hline
\end{tabular}
\end{center}

\section{难度升一级}

\textbf{题目：} “两个数，第一个数的两倍加上第二个数等于 25，第一个数加上第二个数的两倍等于 20。这两个数各是几？”

这道题，两个条件都很绕。慢慢来，一句话一句话翻译。

\subsection{第一步：定未知数}

第一个数 $= x$，第二个数 $= y$。

\subsection{第二步：翻译条件一}

“第一个数的两倍加上第二个数等于 25。”

\begin{center}
\begin{tabular}{|l|c|}
\hline
\textbf{中文} & \textbf{数学} \\
\hline
第一个数的两倍 & $2x$ \\
\hline
加上第二个数 & $+ y$ \\
\hline
等于 25 & $= 25$ \\
\hline
\end{tabular}
\end{center}

\[2x + y = 25\]

\subsection{第三步：翻译条件二}

“第一个数加上第二个数的两倍等于 20。”

\begin{center}
\begin{tabular}{|l|c|}
\hline
\textbf{中文} & \textbf{数学} \\
\hline
第一个数 & $x$ \\
\hline
加上第二个数的两倍 & $+ 2y$ \\
\hline
等于 20 & $= 20$ \\
\hline
\end{tabular}
\end{center}

\[x + 2y = 20\]

\subsection{第四步：组成方程组}

\[\begin{cases} 2x + y = 25 \\ x + 2y = 20 \end{cases}\]

\subsection{第五步：消元}

这次直接加减，$x$ 和 $y$ 都消不掉。怎么办？

\textbf{把第一个方程乘以 2：}

\[4x + 2y = 50\]

现在用这个方程减去第二个方程：

\[(4x + 2y) - (x + 2y) = 50 - 20\]
\[3x = 30\]
\[x = 10\]

把 $x = 10$ 放回第一个方程：

\[2 \times 10 + y = 25\]
\[20 + y = 25\]
\[y = 5\]

\textbf{答案：} 第一个数是 \textbf{10}，第二个数是 \textbf{5}。

\subsection{验证}

\begin{center}
\begin{tabular}{|l|l|c|}
\hline
\textbf{条件} & \textbf{检验} & \textbf{结果} \\
\hline
第一个数的两倍加上第二个数等于25 & $2 \times 10 + 5 = 25$ & \checkmark \\
\hline
第一个数加上第二个数的两倍等于20 & $10 + 2 \times 5 = 20$ & \checkmark \\
\hline
\end{tabular}
\end{center}

\section{为什么要乘以 2？（系数对齐）}

你可能看到“把第一个方程乘以 2”，一头雾水。

为什么要这样做？

还是天平的道理。两边同时乘以同一个数，天平还是平衡的。

我们把第一个方程乘以 2，是为了让 $y$ 的系数变成 2，和第二个方程的 $y$ 系数一样。

\begin{center}
\begin{tabular}{|l|c|c|}
\hline
\textbf{方程} & \textbf{原来的} $y$ \textbf{系数} & \textbf{调整后的} $y$ \textbf{系数} \\
\hline
第一个方程 & 1 & 2（乘以2） \\
\hline
第二个方程 & 2 & 2（不变） \\
\hline
\end{tabular}
\end{center}

这样相减，$y$ 就能消掉了：

\[4x + 2y\]
\[-(x + 2y)\]
\[\overline{3x + 0}\]

$y$ 消掉了！

\textbf{这个方法，叫做“系数对齐”。}

\begin{quote}
\textbf{系数对齐}：让你想消掉的那个未知数，两个方程里的系数变成一样，然后相减，它就消失了。
\end{quote}

\section{综合题目}

\textbf{题目：} “小明和小红去买文具。小明买了 3 支钢笔和 2 本笔记本，花了 28 块。小红买了 2 支钢笔和 3 本笔记本，花了 27 块。钢笔和笔记本各多少钱？”

这道题，两个未知数，两种买法，翻译出来自然就是两个方程。

\subsection{第一步：定未知数}

\begin{center}
\begin{tabular}{|l|c|}
\hline
\textbf{什么} & \textbf{用什么表示} \\
\hline
钢笔每支 & $x$ 块 \\
\hline
笔记本每本 & $y$ 块 \\
\hline
\end{tabular}
\end{center}

\textbf{所以 $x$ 就是钢笔的价格呐，$y$ 就是笔记本的价格呐！}

\subsection{第二步：翻译两个条件}

\begin{center}
\begin{tabular}{|c|l|c|}
\hline
\textbf{谁} & \textbf{买了什么} & \textbf{翻译} \\
\hline
小明 & 3支钢笔、2本笔记本，花了28块 & $3x + 2y = 28$ \\
\hline
小红 & 2支钢笔、3本笔记本，花了27块 & $2x + 3y = 27$ \\
\hline
\end{tabular}
\end{center}

\subsection{第三步：组成方程组}

\[\begin{cases} 3x + 2y = 28 \\ 2x + 3y = 27 \end{cases}\]

\subsection{第四步：消元（系数对齐）}

想消掉 $y$，让两个方程的 $y$ 系数一样。

\begin{center}
\begin{tabular}{|l|l|l|}
\hline
\textbf{方程} & \textbf{操作} & \textbf{结果} \\
\hline
第一个方程 & 乘以3 & $9x + 6y = 84$ \\
\hline
第二个方程 & 乘以2 & $4x + 6y = 54$ \\
\hline
\end{tabular}
\end{center}

两个方程相减：

\[(9x + 6y) - (4x + 6y) = 84 - 54\]
\[5x = 30\]
\[x = 6\]

把 $x = 6$ 放回第一个方程：

\[3 \times 6 + 2y = 28\]
\[18 + 2y = 28\]
\[2y = 10\]
\[y = 5\]

\textbf{答案：} 钢笔每支 \textbf{6} 块，笔记本每本 \textbf{5} 块。

\subsection{验证}

\begin{center}
\begin{tabular}{|l|l|c|}
\hline
\textbf{谁} & \textbf{检验} & \textbf{结果} \\
\hline
小明 & $3 \times 6 + 2 \times 5 = 18 + 10 = 28$ & \checkmark \\
\hline
小红 & $2 \times 6 + 3 \times 5 = 12 + 15 = 27$ & \checkmark \\
\hline
\end{tabular}
\end{center}

\section{现在你来}

\begin{quote}
\textbf{第一道：}

“两个数，大的数是小的数的 2 倍，两个数相加等于 15。这两个数各是几？”

提示：一个 $x$ 够用了。想想为什么。
\end{quote}

\begin{quote}
\textbf{第二道：}

“小明买了 2 支铅笔和 1 块橡皮，花了 7 块。小红买了 1 支铅笔和 2 块橡皮，花了 5 块。铅笔和橡皮各多少钱？”

提示：两个未知数，两个方程。系数对不齐，想想怎么办。
\end{quote}

写下来，验证一下。

\section{本章小结}

两个未知数，不比一个难多少。

难的不是计算，是读懂题目里的两个条件。

\begin{quote}
\textbf{读懂了 $\to$ 翻译出来 $\to$ 剩下的就是消元}
\end{quote}

消元也只有一句话：

\begin{quote}
\textbf{让你想消的那个数，系数变成一样，然后相减。}
\end{quote}

就这样。

\subsection*{本章知识地图}

\begin{verbatim}
两个未知数
│
├─ 什么时候需要两个未知数？
│   ├─ 有直接关系 → 一个x够用（用关系表示另一个）
│   └─ 没有直接关系 → 需要x和y
│
├─ 解题步骤
│   ├─ ① 定未知数（x是什么，y是什么）
│   ├─ ② 翻译两个条件（得到两个方程）
│   ├─ ③ 组成方程组
│   ├─ ④ 消元（系数对齐）
│   └─ ⑤ 验证
│
└─ 系数对齐的方法
    ├─ 看哪个未知数容易消
    ├─ 把方程乘以某个数，让系数一样
    └─ 相减，消掉那个未知数
\end{verbatim}

\chapter{两个未知数——极复杂的语言翻译}

\section{你听说过鸡兔同笼吗？}

这是一道流传了几千年的题目。

几千年前的中国人，就在想这个问题。

说明什么？

说明这道题，难倒了很多很多人。

但你不一样。

你已经会翻译了，你已经会消元了。

这道题，对你来说，只是一道普通的翻译题。

（好吧，稍微难一点点。但你行的。）

\section{先来看鸡兔同笼}

\textbf{题目：} “笼子里有鸡和兔，一共 10 个头，28 条腿。鸡和兔各有几只？”

读完这道题，你可能第一反应是：“头？腿？这和数学有什么关系？”

有关系。大有关系。

\begin{center}
\begin{tabular}{|c|c|c|}
\hline
\textbf{动物} & \textbf{头的数量} & \textbf{腿的数量} \\
\hline
鸡 & 1个 & 2条 \\
\hline
兔 & 1个 & 4条 \\
\hline
\end{tabular}
\end{center}

\textbf{这就是线索。}

\subsection{第一步：定未知数}

\begin{center}
\begin{tabular}{|c|c|}
\hline
\textbf{什么} & \textbf{用什么表示} \\
\hline
鸡有几只 & $x$ \\
\hline
兔有几只 & $y$ \\
\hline
\end{tabular}
\end{center}

\textbf{所以 $x$ 就是鸡呐，$y$ 就是兔呐！}

\subsection{第二步：翻译两个条件}

\begin{center}
\begin{tabular}{|l|l|c|}
\hline
\textbf{条件} & \textbf{怎么翻译} & \textbf{方程} \\
\hline
一共 10 个头 & 每只鸡1个头，每只兔1个头 & $x + y = 10$ \\
\hline
一共 28 条腿 & 每只鸡2条腿，每只兔4条腿 & $2x + 4y = 28$ \\
\hline
\end{tabular}
\end{center}

\subsection{第三步：组成方程组}

\[\begin{cases} x + y = 10 \\ 2x + 4y = 28 \end{cases}\]

\subsection{第四步：消元}

把第一个方程乘以 2：

\[2x + 2y = 20\]

用第二个方程减去它：

\[(2x + 4y) - (2x + 2y) = 28 - 20\]
\[2y = 8\]
\[y = 4\]

把 $y = 4$ 放回第一个方程：

\[x + 4 = 10\]
\[x = 6\]

\textbf{答案：} 鸡有 \textbf{6} 只，兔有 \textbf{4} 只。

\subsection{验证}

\begin{center}
\begin{tabular}{|l|l|c|}
\hline
\textbf{检验项} & \textbf{计算} & \textbf{结果} \\
\hline
头的总数 & $6 + 4 = 10$ & \checkmark \\
\hline
腿的总数 & $6 \times 2 + 4 \times 4 = 12 + 16 = 28$ & \checkmark \\
\hline
\end{tabular}
\end{center}

\section{看见了吗？}

几千年前难倒无数人的题目，你刚才用了不到 10 行就解出来了。

不是你比古人聪明。

是你有工具，他们没有。

\textbf{这就是数学的力量。}

\section{再来一道变形题}

\textbf{题目：} “停车场里有轿车和摩托车，一共 30 辆，轮子一共 88 个。轿车和摩托车各有几辆？”

\begin{center}
\begin{tabular}{|c|c|}
\hline
\textbf{车辆} & \textbf{轮子数量} \\
\hline
轿车 & 4个 \\
\hline
摩托车 & 2个 \\
\hline
\end{tabular}
\end{center}

（这道题和鸡兔同笼一模一样的套路，只是换了主角。发现没有？）

\subsection{第一步：定未知数}

轿车 $= x$，摩托车 $= y$

\subsection{第二步：翻译两个条件}

\begin{center}
\begin{tabular}{|l|c|}
\hline
\textbf{条件} & \textbf{方程} \\
\hline
一共 30 辆 & $x + y = 30$ \\
\hline
轮子一共 88 个 & $4x + 2y = 88$ \\
\hline
\end{tabular}
\end{center}

\subsection{第三步：组成方程组}

\[\begin{cases} x + y = 30 \\ 4x + 2y = 88 \end{cases}\]

\subsection{第四步：消元}

把第一个方程乘以 2：

\[2x + 2y = 60\]

用第二个方程减去它：

\[(4x + 2y) - (2x + 2y) = 88 - 60\]
\[2x = 28\]
\[x = 14\]

把 $x = 14$ 放回第一个方程：

\[14 + y = 30\]
\[y = 16\]

\textbf{答案：} 轿车 \textbf{14} 辆，摩托车 \textbf{16} 辆。

\subsection{验证}

\begin{center}
\begin{tabular}{|l|l|c|}
\hline
\textbf{检验项} & \textbf{计算} & \textbf{结果} \\
\hline
辆数 & $14 + 16 = 30$ & \checkmark \\
\hline
轮子 & $14 \times 4 + 16 \times 2 = 56 + 32 = 88$ & \checkmark \\
\hline
\end{tabular}
\end{center}

\section{难度升级（购物问题）}

\textbf{题目：} “小明去超市买苹果和橙子。苹果每斤 6 块，橙子每斤 4 块，一共花了 48 块。小红买的苹果斤数和小明一样，但橙子买了小明的两倍，花了 60 块。小明买了多少斤苹果和橙子？”

这道题，信息量更大了。别慌，一句话一句话地读。

\subsection{第一步：定未知数}

\begin{center}
\begin{tabular}{|c|c|}
\hline
\textbf{什么} & \textbf{用什么表示} \\
\hline
小明买的苹果斤数 & $x$ \\
\hline
小明买的橙子斤数 & $y$ \\
\hline
\end{tabular}
\end{center}

\textbf{所以 $x$ 就是苹果呐，$y$ 就是橙子呐！}

\subsection{第二步：翻译条件一}

“小明买苹果和橙子，苹果每斤 6 块，橙子每斤 4 块，一共花了 48 块。”

\begin{center}
\begin{tabular}{|l|c|c|c|}
\hline
\textbf{项目} & \textbf{单价} & \textbf{数量} & \textbf{金额} \\
\hline
苹果 & 6块/斤 & $x$斤 & $6x$ \\
\hline
橙子 & 4块/斤 & $y$斤 & $4y$ \\
\hline
总共 & & & 48块 \\
\hline
\end{tabular}
\end{center}

\[6x + 4y = 48\]

\subsection{第三步：翻译条件二}

“小红买的苹果斤数和小明一样，但橙子买了小明的两倍，花了 60 块。”

\begin{center}
\begin{tabular}{|l|c|c|c|}
\hline
\textbf{项目} & \textbf{小红买的数量} & \textbf{单价} & \textbf{金额} \\
\hline
苹果 & $x$斤（和小明一样） & 6块/斤 & $6x$ \\
\hline
橙子 & $2y$斤（小明的两倍） & 4块/斤 & $4 \times 2y = 8y$ \\
\hline
总共 & & & 60块 \\
\hline
\end{tabular}
\end{center}

\[6x + 8y = 60\]

\subsection{第四步：组成方程组}

\[\begin{cases} 6x + 4y = 48 \\ 6x + 8y = 60 \end{cases}\]

\subsection{第五步：消元}

两个方程的 $6x$ 系数一样，直接相减：

\[(6x + 8y) - (6x + 4y) = 60 - 48\]
\[4y = 12\]
\[y = 3\]

把 $y = 3$ 放回第一个方程：

\[6x + 4 \times 3 = 48\]
\[6x + 12 = 48\]
\[6x = 36\]
\[x = 6\]

\textbf{答案：} 小明买了 \textbf{6} 斤苹果，\textbf{3} 斤橙子。

\subsection{验证}

\begin{center}
\begin{tabular}{|l|l|c|}
\hline
\textbf{检验项} & \textbf{计算} & \textbf{结果} \\
\hline
小明花的钱 & $6 \times 6 + 4 \times 3 = 36 + 12 = 48$ & \checkmark \\
\hline
小红花的钱 & $6 \times 6 + 4 \times 6 = 36 + 24 = 60$ & \checkmark \\
\hline
\end{tabular}
\end{center}

完美！

\section{再来一道综合题}

\textbf{题目：} “一家文具店卖钢笔和铅笔。钢笔每支 5 元，铅笔每支 2 元。小明买了一些钢笔和铅笔，一共买了 10 支，花了 32 元。小明买了几支钢笔和几支铅笔？”

\subsection{第一步：定未知数}

\begin{center}
\begin{tabular}{|c|c|}
\hline
\textbf{什么} & \textbf{用什么表示} \\
\hline
钢笔支数 & $x$ \\
\hline
铅笔支数 & $y$ \\
\hline
\end{tabular}
\end{center}

\subsection{第二步：翻译两个条件}

\begin{center}
\begin{tabular}{|l|c|}
\hline
\textbf{条件} & \textbf{方程} \\
\hline
一共买了 10 支 & $x + y = 10$ \\
\hline
一共花了 32 元 & $5x + 2y = 32$ \\
\hline
\end{tabular}
\end{center}

\subsection{第三步：组成方程组}

\[\begin{cases} x + y = 10 \\ 5x + 2y = 32 \end{cases}\]

\subsection{第四步：消元}

把第一个方程乘以 2：

\[2x + 2y = 20\]

用第二个方程减去它：

\[(5x + 2y) - (2x + 2y) = 32 - 20\]
\[3x = 12\]
\[x = 4\]

把 $x = 4$ 放回第一个方程：

\[4 + y = 10\]
\[y = 6\]

\textbf{答案：} 钢笔 \textbf{4} 支，铅笔 \textbf{6} 支。

\subsection{验证}

\begin{center}
\begin{tabular}{|l|l|c|}
\hline
\textbf{检验项} & \textbf{计算} & \textbf{结果} \\
\hline
总支数 & $4 + 6 = 10$ & \checkmark \\
\hline
总金额 & $5 \times 4 + 2 \times 6 = 20 + 12 = 32$ & \checkmark \\
\hline
\end{tabular}
\end{center}

\section{常见问题}

\subsection{问题1：答案不是整数怎么办？}

有时候解出来 $x = 2.5$ 或 $x = \dfrac{7}{3}$，不是整数。

\begin{center}
\begin{tabular}{|l|l|l|}
\hline
\textbf{情况} & \textbf{说明} & \textbf{怎么办} \\
\hline
题目问“几个苹果” & 苹果不能是2.5个 & 检查是不是翻译错了 \\
\hline
题目问“几斤苹果” & 可以是2.5斤 & 答案合理 \\
\hline
验证不通过 & 计算或翻译有误 & 重新检查 \\
\hline
\end{tabular}
\end{center}

\textbf{简单说}：看题目问的东西能不能是小数。人数、只数不能是小数，重量、长度可以。

\subsection{问题2：两个未知数，但只找到一个方程？}

\begin{center}
\begin{tabular}{|l|l|}
\hline
\textbf{情况} & \textbf{说明} \\
\hline
题目条件不够 & 这道题无法解出唯一答案 \\
\hline
漏读了条件 & 回去仔细读题，找第二个条件 \\
\hline
\end{tabular}
\end{center}

\textbf{简单说}：两个未知数必须要两个条件，少一个解不出来。

\subsection{问题3：消元消不掉怎么办？}

两个方程直接加减，两个未知数都消不掉？

\textbf{解决办法}：用系数对齐（第九章学过的）

\begin{itemize}
\item 把其中一个方程乘以某个数
\item 让某个未知数的系数变成一样
\item 然后相减，就能消掉了
\end{itemize}

\section{现在你来}

\begin{quote}
\textbf{第一道：}

“笼子里有鸡和兔，一共 12 个头，32 条腿。鸡和兔各有几只？”
\end{quote}

\begin{quote}
\textbf{第二道：}

“停车场里有轿车和电动车，一共 20 辆，轮子一共 62 个。轿车有 4 个轮子，电动车有 3 个轮子。轿车和电动车各有几辆？”
\end{quote}

写下来，验证一下。

\section{本章小结}

极复杂的语言翻译，说穿了，还是那几步：

\begin{quote}
\textbf{定未知数 $\to$ 翻译条件 $\to$ 建方程组 $\to$ 消元 $\to$ 验证}
\end{quote}

题目再长，信息再多，一句话一句话地读，一个词一个词地翻译。

\textbf{不要一口气吞下去。慢慢嚼，慢慢消化。}

\subsection*{本章知识地图}

\begin{verbatim}
极复杂的语言翻译（鸡兔同笼型）
│
├─ 经典模型
│   ├─ 鸡兔同笼（头和腿）
│   ├─ 车辆问题（车辆和轮子）
│   └─ 购物问题（数量和金额）
│
├─ 翻译要点
│   ├─ 找出两种对象的“特征”
│   │   └─ 鸡2条腿，兔4条腿
│   ├─ 每个条件对应一个方程
│   │   ├─ 数量关系 → 加法方程
│   │   └─ 特征总和 → 乘法加法方程
│   └─ 注意“倍数关系”的翻译
│
├─ 解题步骤
│   ├─ ① 定未知数（两种对象各多少）
│   ├─ ② 翻译条件（数量、特征总和）
│   ├─ ③ 组成方程组
│   ├─ ④ 消元（系数对齐）
│   └─ ⑤ 验证（特别重要！）
│
└─ 常见问题
    ├─ 答案不是整数 → 看题目问的能不能是小数
    ├─ 只有一个方程 → 检查是否漏读条件
    └─ 消不掉 → 用系数对齐
\end{verbatim}

\chapter{两个未知数——多种算法}

\section{先说一件事}

上一章，你学会了消元法。

找到系数一样的，相加或相减，消掉一个未知数。

但有时候，消元法用起来很麻烦。

比如这道题：

\[\begin{cases} x + y = 10 \\ x = 3y - 2 \end{cases}\]

你要消元，得先把系数对齐。

但第二个方程，已经把 $x$ 单独写出来了。

何必那么麻烦？

直接把 $x = 3y - 2$，塞进第一个方程里不就完了？

\[3y - 2 + y = 10\]

$x$ 直接消失了。

这就是今天要讲的第二种方法：\textbf{代入法。}

\section{代入法是什么？}

代入法，就是把一个未知数，用另一个未知数表示，然后塞进另一个方程里。

就像你知道“小明比小红大 2 岁”，你知道小红几岁，就能直接知道小明几岁。

不用两个方程消来消去，直接代进去就行。

\section{什么时候用代入法？}

当方程组里，有一个方程，\textbf{已经把某个未知数单独写出来了}，比如：

\begin{center}
\begin{tabular}{|l|l|}
\hline
\textbf{方程} & \textbf{说明} \\
\hline
$x = 3y + 1$ & $x$ 已经单独在左边 \\
\hline
$y = 2x - 5$ & $y$ 已经单独在左边 \\
\hline
$x = 7$ & $x$ 已经等于一个数 \\
\hline
\end{tabular}
\end{center}

这种时候，用代入法，比消元法快得多。

\section{来，先看我示范}

\textbf{题目：} “两个数，第一个数等于第二个数的 2 倍加 1，两个数相加等于 16。这两个数各是几？”

\subsection{第一步：定未知数}

题目里有两个不知道的数。

\begin{center}
\begin{tabular}{|c|c|}
\hline
\textbf{什么} & \textbf{用什么表示} \\
\hline
第一个数 & $x$ \\
\hline
第二个数 & $y$ \\
\hline
\end{tabular}
\end{center}

\subsection{第二步：翻译两个条件}

题目给了我们两句话。一句话一句话来。

\textbf{第一句：} “第一个数等于第二个数的 2 倍加 1。”

\begin{center}
\begin{tabular}{|l|c|}
\hline
\textbf{中文} & \textbf{数学} \\
\hline
第一个数 & $x$ \\
\hline
等于 & $=$ \\
\hline
第二个数的 2 倍 & $2y$ \\
\hline
加 1 & $+ 1$ \\
\hline
\end{tabular}
\end{center}

合起来：

\[x = 2y + 1\]

\textbf{第二句：} “两个数相加等于 16。”

\begin{center}
\begin{tabular}{|l|c|}
\hline
\textbf{中文} & \textbf{数学} \\
\hline
第一个数 & $x$ \\
\hline
加上第二个数 & $+ y$ \\
\hline
等于 16 & $= 16$ \\
\hline
\end{tabular}
\end{center}

合起来：

\[x + y = 16\]

\subsection{第三步：选方法}

现在我们有两个方程：

\[\begin{cases} x = 2y + 1 \\ x + y = 16 \end{cases}\]

停一下，看看第一个方程。

\[x = 2y + 1\]

$x$ 已经单独在左边了。

这就是说，$x$ 等于 $2y + 1$。

那第二个方程里的 $x$，我们可以直接换掉。

把 $x$ 换成 $2y + 1$，第二个方程就只剩 $y$ 了。

\begin{quote}
\textbf{这就是代入法的核心}：

把一个未知数的表达式，直接塞进另一个方程里，消掉它。
\end{quote}

\subsection{第四步：代入}

把 $x = 2y + 1$ 塞进第二个方程：

原来的第二个方程：

\[x + y = 16\]

把 $x$ 换成 $2y + 1$：

\[(2y + 1) + y = 16\]

\textbf{为什么要加括号？}

因为 $2y + 1$ 是一整件事，要把这整件事放进去，不是只放 $2y$，不是只放 $1$。

括号提醒你：\textbf{这是一个整体，别拆散它。}

现在展开括号：

\[2y + 1 + y = 16\]

合并同类项：

$2y$ 是 2 个 $y$，$y$ 是 1 个 $y$，加在一起是 3 个 $y$，也就是 $3y$。

\[3y + 1 = 16\]

两边同时减去 1：

\[3y = 16 - 1 = 15\]

两边同时除以 3：

\[y = 15 \div 3 = 5\]

$y$ 找到了！

\subsection{第五步：找 $x$}

$y = 5$ 了。

现在把 $y = 5$ 放回第一个方程：

\[x = 2y + 1\]
\[x = 2 \times 5 + 1\]

先算乘法（还记得吗？乘法先算）：

\[2 \times 5 = 10\]

再加 1：

\[x = 10 + 1 = 11\]

$x$ 也找到了！

\subsection{验证}

把 $x = 11$，$y = 5$ 放回两个方程，一个一个检查。

\begin{center}
\begin{tabular}{|l|l|c|}
\hline
\textbf{方程} & \textbf{检验} & \textbf{结果} \\
\hline
$x = 2y + 1$ & 左边：$11$，右边：$2 \times 5 + 1 = 11$ & \checkmark \\
\hline
$x + y = 16$ & 左边：$11 + 5 = 16$，右边：$16$ & \checkmark \\
\hline
\end{tabular}
\end{center}

两个方程都验证通过。

\textbf{答案：} 第一个数是 \textbf{11}，第二个数是 \textbf{5}。

\section{和消元法比一比}

同一道题，用消元法：

\[\begin{cases} x = 2y + 1 \\ x + y = 16 \end{cases}\]

先把第一个方程整理成：

\[x - 2y = 1\]

两个方程：

\[\begin{cases} x - 2y = 1 \\ x + y = 16 \end{cases}\]

相减：

\[(x + y) - (x - 2y) = 16 - 1\]
\[3y = 15\]
\[y = 5\]
\[x = 11\]

结果一样，但代入法明显更快。

\begin{quote}
\textbf{结论}：

\begin{itemize}
\item 方程已经把某个未知数单独写出来了 $\to$ 用代入法
\item 方程没有 $\to$ 用消元法
\end{itemize}
\end{quote}

\section{还有一种情况，你得自己整理}

\textbf{题目：} “两个数，第一个数加上 3 等于第二个数的两倍，两个数相加等于 12。这两个数各是几？”

\subsection{翻译}

\textbf{条件一：} 第一个数加上 3 等于第二个数的两倍。

\[x + 3 = 2y\]

\textbf{条件二：} 两个数相加等于 12。

\[x + y = 12\]

这次，没有哪个方程把 $x$ 或 $y$ 单独写出来。

怎么办？

\textbf{自己整理一下。}

从条件一，把 $x$ 单独写出来：

\[x = 2y - 3\]

然后代入条件二：

\[(2y - 3) + y = 12\]
\[3y - 3 = 12\]
\[3y = 15\]
\[y = 5\]
\[x = 2 \times 5 - 3 = 7\]

\subsection{验证}

\begin{center}
\begin{tabular}{|l|l|c|}
\hline
\textbf{方程} & \textbf{检验} & \textbf{结果} \\
\hline
$x + 3 = 2y$ & $7 + 3 = 10 = 2 \times 5$ & \checkmark \\
\hline
$x + y = 12$ & $7 + 5 = 12$ & \checkmark \\
\hline
\end{tabular}
\end{center}

\section{什么时候用消元法，什么时候用代入法？}

其实，两种方法，都能解任何方程组。

只是有时候一种更快，有时候另一种更快。

你自己判断就行。

\begin{center}
\begin{tabular}{|l|l|}
\hline
\textbf{情况} & \textbf{推荐方法} \\
\hline
有方程已经把某个未知数单独写出来了 & 代入法 \\
\hline
两个方程的某个系数一样 & 消元法（直接加减） \\
\hline
两个方程的系数都不一样 & 消元法（先乘再加减） \\
\hline
两种方法看起来差不多 & 随便你，哪个顺手用哪个 \\
\hline
\end{tabular}
\end{center}

（最后一行是认真的。数学没有规定你必须用哪种方法。对就行。）

\section{再来一道综合的}

\textbf{题目：} “小明和小红去书店。小明买了 2 本书和 1 支笔，花了 27 块。书的价格是笔的 4 倍。书和笔各多少钱？”

\subsection{第一步：定未知数}

\begin{center}
\begin{tabular}{|c|c|}
\hline
\textbf{什么} & \textbf{用什么表示} \\
\hline
书的价格 & $x$ \\
\hline
笔的价格 & $y$ \\
\hline
\end{tabular}
\end{center}

\subsection{第二步：翻译}

\textbf{条件一：} 小明花了 27 块。

\[2x + y = 27\]

\textbf{条件二：} 书的价格是笔的 4 倍。

\[x = 4y\]

\subsection{第三步：选方法}

条件二已经把 $x$ 单独写出来了。

用代入法！

把 $x = 4y$ 塞进条件一：

\[2 \times 4y + y = 27\]
\[8y + y = 27\]
\[9y = 27\]
\[y = 3\]
\[x = 4 \times 3 = 12\]

\textbf{答案：} 书每本 \textbf{12} 块，笔每支 \textbf{3} 块。

\subsection{验证}

\begin{center}
\begin{tabular}{|l|l|c|}
\hline
\textbf{方程} & \textbf{检验} & \textbf{结果} \\
\hline
$2x + y = 27$ & $2 \times 12 + 3 = 27$ & \checkmark \\
\hline
$x = 4y$ & $12 = 4 \times 3$ & \checkmark \\
\hline
\end{tabular}
\end{center}

完美！

\section{现在你来}

\begin{quote}
\textbf{第一道：}

“两个数，第一个数等于第二个数加 5，两个数相加等于 17。这两个数各是几？”

提示：第一个条件已经把第一个数单独写出来了，用代入法。
\end{quote}

\begin{quote}
\textbf{第二道：}

“鸡兔同笼，一共 15 个头，38 条腿。鸡和兔各有几只？”

提示：这道题用消元法更快。但你要是想用代入法，也没人拦你。
\end{quote}

写下来，验证一下。

\section{本章小结}

消元法和代入法，是解方程组的两把钥匙。

不是非此即彼，是根据题目选择。

\begin{quote}
\textbf{看到单独写出来的未知数 $\to$ 代入}

\textbf{看到系数一样的 $\to$ 消元}

\textbf{看不出来哪个更好 $\to$ 随便选一个，做下去}
\end{quote}

最重要的是：\textbf{解完，验证。}

不管用哪种方法，验证都不能省。

（我知道你懒得验证。但验证真的很重要。真的。）

\subsection*{本章知识地图}

\begin{verbatim}
两个未知数的解法
│
├─ 消元法（上一章学的）
│   ├─ 系数一样 → 直接加减
│   └─ 系数不一样 → 先乘再加减
│
├─ 代入法（本章学的）
│   ├─ 找到单独写出来的未知数
│   ├─ 代入另一个方程
│   └─ 解出第一个，再回代
│
└─ 选择方法
    ├─ 有单独的未知数 → 代入法更快
    ├─ 系数一样 → 消元法更快
    └─ 都差不多 → 随便选
\end{verbatim}

\chapter{线的世界——直线、射线、线段的一切}

\section{一、点（30秒讲完）}

点就是一个位置。

用大写字母表示：A、B、C、P……

\begin{verbatim}
      · A
\end{verbatim}

以后会遇到几种特殊的点：

\begin{itemize}
\item \textbf{端点}：线段或射线的“头”
\item \textbf{中点}：线段正中间那个点
\item \textbf{交点}：两条线碰在一起的那个点
\end{itemize}

点就这么多。下面全讲线。

\section{二、线段}

\subsection{什么是线段？}

你手里有一根筷子。

\begin{verbatim}
      ============================
      ↑                          ↑
    这一头                      那一头
\end{verbatim}

左边有一头。

右边也有一头。

\textbf{有两头。到头就没了。}

这就是线段。

\subsection{线段长什么样？}

\begin{verbatim}
      A                               B
      ·-------------------------------·
      ↑                               ↑
    端点                             端点
\end{verbatim}

两头的那两个点，叫\textbf{端点}。

线段有\textbf{两个端点}。

\subsection{线段的特点}

\textbf{第一，有两个端点。}

左边一个，右边一个。

\textbf{第二，到端点就停了。}

不能再往外走了。

左边到了A就停，右边到了B也停。

\textbf{第三，没有方向。}

从A走到B，和从B走到A，是同一条路。

所以线段AB和线段BA是同一条线段。

\textbf{第四，三种线里，只有线段能量。}

你可以拿尺子量一条线段。

射线和直线做不到这一点。

\subsection{生活中的线段}

筷子——有两头，到头就没了。

铅笔——有两头，到头就没了。

桌子的边——有两头，到头就没了。

\section{三、射线}

\subsection{什么是射线？}

晚上你打开手电筒。

光从手电筒里射出来。

然后往前走。

一直走。

永远不停。

\begin{verbatim}
      手电筒
         ‖
         ‖==============================>
         ‖              光一直往前
                        永远不停
\end{verbatim}

\textbf{有一头。但另一头永远不停。}

这就是射线。

\subsection{射线长什么样？}

\begin{verbatim}
      A
      ·-------------------------------------->
      ↑                                       ↑
    端点                                  永远不停
\end{verbatim}

只有一个端点。

另一边没有尽头。

\subsection{射线的特点}

\textbf{第一，只有一个端点。}

就是起点。

另一头没有尽头。

\textbf{第二，永远不停。}

所以量不了。

\textbf{第三，有方向。}

往哪边走，很重要。

往左走和往右走，是两条不同的射线。

\subsection{生活中的射线}

手电筒的光——从灯泡出发，往前走，不停。

太阳光——从太阳出发，往外走，不停。

激光笔——从笔头出发，往前走，不停。

你的视线——从眼睛出发，往前看，不停。

\section{四、直线}

\subsection{什么是直线？}

线段是两头都封住的。

射线是一头封住、一头打开的。

\textbf{现在把两头都打开。}

两边都往外走。

两边都永远不停。

\begin{verbatim}
<------------------------------------------------------->
↑                                                       ↑
这边永远不停                                       这边也永远不停
\end{verbatim}

没有端点了。

\textbf{这就是直线。}

\subsection{直线长什么样？}

\begin{verbatim}
<------------------------------------------------------->
               没有头                没有尾
\end{verbatim}

两边都没有尽头。

根本没有“端点”这个东西。

\subsection{直线的特点}

\textbf{第一，没有端点。}

两边都是无限的，哪来的“头”？

\textbf{第二，两边都永远不停。}

所以也量不了。

\textbf{第三，没有方向。}

两边都能走，所以没有固定的方向。

\subsection{生活中有直线吗？}

\textbf{没有。}

你找不到任何东西是两头都没有尽头的。

直线只存在于想象中。

你在纸上画的“直线”，其实是线段，因为纸有边。

但我们可以假装它能无限延伸下去。

\section{五、三种线放在一起比较}

\subsection{一张图看明白}

\begin{verbatim}
线段：两头都封住

      A                               B
      ·-------------------------------·
    端点                             端点



射线：一头封住，一头打开

      A
      ·-------------------------------------->
    端点                                永远不停



直线：两头都打开

<------------------------------------------------------->
永远不停                                         永远不停
\end{verbatim}

\subsection{用大白话说}

\textbf{线段}：有头有尾。

\textbf{射线}：有头没尾。

\textbf{直线}：没头没尾。

\subsection{它们是什么关系？}

直线是最大的。

在直线上随便切一刀，切出来的一半就是射线。

在射线上再切一刀，切下来的一段就是线段。

\textbf{直线 > 射线 > 线段}

线段是直线的一小部分。

\section{六、线怎么写？}

\subsection{线段怎么写？}

\subsubsection{写法一：用两个端点的字母}

\textbf{线段AB}

\begin{verbatim}
      A                   B
      ·-------------------·
\end{verbatim}

\textbf{A和B是什么意思？}

A和B就是两个端点的\textbf{名字}。

就像你们班有个同学叫小明，另一个叫小红。

小明站在左边，小红站在右边。

你说“从小明到小红”，别人就知道你说的是哪段路。

A和B也是一样的。

A是左边那个端点的名字。

B是右边那个端点的名字。

\begin{verbatim}
      A                   B
      ·-------------------·
      ↑                   ↑
   “这个点叫A”         “这个点叫B”
\end{verbatim}

它们就是名字。

仅此而已。

你也可以叫它们C和D，或者P和Q，或者M和N。

随便什么大写字母都行。

叫什么不重要，重要的是：你给两个端点各起了个名字，这样别人就知道你说的是哪条线段。

\textbf{为什么两个端点就能代表整条线段？}

因为两个端点一旦定了，线段就定死了。

A在这，B在那，中间连起来，就只有这一条线段。

不可能有第二条。

所以两个端点的名字就够了。

说“线段AB”，别人立刻知道你在说哪条线段。

AB还有第二层意思。

AB不只是代表这条线段本身，还能直接代表它量出来有多少。

比如：

\[AB = 5 \text{ cm}\]

这句话的意思是：线段AB量出来是5厘米。

所以AB有两层意思：

\begin{center}
\begin{tabular}{|l|l|}
\hline
AB的意思 & 举例 \\
\hline
代表这条线段本身 & “画出线段AB” \\
\hline
代表这条线段量出来的结果 & “AB = 5cm” \\
\hline
\end{tabular}
\end{center}

看上下文就知道是哪个意思。

\subsubsection{写法二：用一个小写字母}

\textbf{线段a}

\begin{verbatim}
      ·-------------------·
              a
\end{verbatim}

\textbf{为什么有时候用小写字母？}

因为有时候懒得给端点起名字。

或者端点不重要，你只关心这条线段本身。

这时候直接给整条线段起个名字就行：a、b、c……

\textbf{小写字母和大写字母的区别：}

\begin{center}
\begin{tabular}{|l|l|}
\hline
写法 & 什么时候用 \\
\hline
\textbf{线段AB}（大写） & 端点很重要，要说清楚两头在哪 \\
\hline
\textbf{线段a}（小写） & 端点不重要，只是给线段起个名字 \\
\hline
\end{tabular}
\end{center}

\textbf{线段AB和线段BA一样吗？}

一样。

线段没有方向。

\textbf{顺序不重要。}

\subsection{射线怎么写？}

用两个字母：

\textbf{射线AB}

\begin{verbatim}
      A            B
      ·------------·--------------------->
    端点         经过B，往这边走
\end{verbatim}

第一个字母A是端点（从这里出发）。

第二个字母B告诉你方向（往B那边走）。

\textbf{射线AB和射线BA一样吗？}

\textbf{不一样！！！}

\begin{verbatim}
射线AB：从A出发，往B方向走

      A            B
      ·------------·--------------------->
    端点



射线BA：从B出发，往A方向走

<---------------------·------------·
                      A            B
                                 端点
\end{verbatim}

端点变了。

方向也反了。

\textbf{完全是两条不同的射线。}

\textbf{怎么记？}

第一个字母 = 端点

第二个字母 = 方向

换了顺序，就换了线。

\subsection{直线怎么写？}

用两个点的字母：

\textbf{直线AB}

\begin{verbatim}
<-----------·-------------------·----------->
            A                   B
\end{verbatim}

或者用一个小写字母：

\textbf{直线l}

\begin{verbatim}
<------------------------------------------->
                      l
\end{verbatim}

\textbf{直线AB和直线BA一样吗？}

一样。

直线没有起点，没有方向。

谁先谁后无所谓。

\textbf{顺序不重要。}

\subsection{总结一下顺序的问题}

\begin{center}
\begin{tabular}{|c|c|}
\hline
类型 & 顺序重要吗 \\
\hline
线段 & 不重要 \\
\hline
\textbf{射线} & \textbf{重要！} \\
\hline
直线 & 不重要 \\
\hline
\end{tabular}
\end{center}

\textbf{只有射线要注意顺序。}

因为只有射线有方向。

\section{七、线段能干什么}

\subsection{能量}

拿尺子量一量，这条线段有多少厘米。

\begin{verbatim}
      A                                   B
      ·-----------------------------------·
      |<------------- 8厘米 ------------->|
\end{verbatim}

\textbf{三种线里，只有线段能量。}

射线和直线永远不停，没法量。

\subsection{能比大小}

哪条长，哪条短，一目了然。

\begin{verbatim}
      A                                           B
      ·-------------------------------------------·
      
      A                       C
      ·-----------------------·

      AB比AC长。
\end{verbatim}

\subsection{有中点}

中点就是线段正中间的那个点。

把线段分成一样的两截。

\begin{verbatim}
      A               M               B
      ·---------------·---------------·
           一半             一半
\end{verbatim}

如果M是AB的中点：

AM = MB = AB的一半

每条线段只有一个中点。

\subsection{能做加减法}

A、B、C三个点在同一条线上，B在A和C的中间：

\begin{verbatim}
      A               B               C
      ·---------------·---------------·
\end{verbatim}

AB + BC = AC

一段加另一段，等于整段。

知道整段和其中一段，能算另一段：

BC = AC - AB

\subsection{能延长}

线段本来到端点就停了。

但有时候需要往外续一段。

\textbf{延长AB}：从B那头往外走。

\begin{verbatim}
      A                         B                 C
      ·-------------------------· - - - - - - - - ·
                                ↑
                            从B往外走
\end{verbatim}

\textbf{反向延长AB}：从A那头往外走。

\begin{verbatim}
      D                 A                         B
      · - - - - - - - - ·-------------------------·
                        ↑
                    从A往外走
\end{verbatim}

\textbf{怎么记？}

延长AB = 从第二个字母往外走

反向延长AB = 从第一个字母往外走

\section{八、直线的规矩}

\subsection{两点确定一条直线}

给你两个点，只能画出一条直线经过它们。

\begin{verbatim}
          · A

                              · B
\end{verbatim}

只能画这一条：

\begin{verbatim}
<---------·-----------------------·--------->
          A                       B
\end{verbatim}

不可能画出第二条。

\textbf{那一个点呢？}

一个点不够。

经过一个点，可以画无数条直线。

因为方向没定，你想往哪画都行。

\begin{verbatim}
            | /
            |/
<-----------·----------->  A
           /|
          / |
\end{verbatim}

\textbf{所以必须两个点才能“锁死”一条直线。}

\subsection{两条直线最多一个交点}

如果两条直线相交，交点只可能有一个。

\begin{verbatim}
              \       /
               \     /
                \   /
                 \ /
                  ·  <-- 就这一个
                 / \
                /   \
\end{verbatim}

不可能有两个交点。

为什么？

如果有两个交点，那这两条直线都经过这两个点。

但两点只能确定一条直线啊。

所以它们其实是同一条直线。

矛盾了。

所以最多一个交点。

\section{九、射线的规矩}

\subsection{射线由两个东西决定}

\textbf{端点}：从哪出发。

\textbf{方向}：往哪走。

端点不同，是不同的射线。

方向不同，也是不同的射线。

\subsection{直线上取一点，分出两条射线}

\begin{verbatim}
<-------------------·-------------------->
                    O
\end{verbatim}

点O把直线切成了两半。

左半边是一条射线。

右半边是另一条射线。

这两条射线：端点相同，方向相反。

合起来就是原来那条直线。

\section{十、两条直线之间的关系}

在同一个平面里，两条直线之间只有三种关系。

没有第四种。

\subsection{第一种：相交}

两条线碰在一起了。

\begin{verbatim}
              \           /
               \         /
                \       /
                 \     /
                  \   /
                   \ /
                    ·  <-- 交点
                   / \
                  /   \
\end{verbatim}

有一个公共点，叫交点。

只有一个，不可能有两个。

\subsection{第二种：平行}

两条线永远碰不到。

\begin{verbatim}
<------------------------------------------------>  直线a


<------------------------------------------------>  直线b
\end{verbatim}

\textbf{符号}：$a \parallel b$

\textbf{读法}：a平行于b

生活中的平行线：

铁轨——两条钢轨永远不碰。

斑马线——每条线都平行。

作业本上的横线——每条都平行。

\subsection{第三种：重合}

两条线完全重叠在一起了。

其实就是同一条线。

\begin{verbatim}
<================================================>
                a和b完全重叠
\end{verbatim}

每一个点都是公共点。

\subsection{一句话总结}

\begin{center}
\begin{tabular}{|c|c|}
\hline
关系 & 公共点有几个 \\
\hline
相交 & 1个 \\
\hline
平行 & 0个 \\
\hline
重合 & 无数个 \\
\hline
\end{tabular}
\end{center}

\subsection{异面直线（提一嘴）}

以上说的是“在同一个平面里”。

如果两条直线不在一个平面呢？

比如一条在桌面上，一条在空中。

它们可能既不相交，也不平行。

这叫\textbf{异面直线}。

以后再细说。

\section{十一、平行线的规矩}

\subsection{过一点只能画一条平行线}

\begin{verbatim}
<------------------------------------------------>  直线l


                         · P
\end{verbatim}

经过点P，和直线l平行的线，只能画一条。

\begin{verbatim}
<------------------------------------------------>  直线l


<- - - - - - - - - - - - - - - - - - - - - - - - >  过P的平行线
                         P
\end{verbatim}

不可能画第二条。

\subsection{平行线之间的距离处处相等}

\begin{verbatim}
<------------------------------------------------>  直线a
         |              |              |
         | 3cm          | 3cm          | 3cm
         |              |              |
<------------------------------------------------>  直线b
\end{verbatim}

不管在哪量，两条平行线之间的距离都一样。

这就是为什么铁轨要平行——距离处处相等，火车才能跑。

\subsection{平行可以传递}

a和b平行，b和c平行。

那a和c也平行。

\begin{verbatim}
<------------------------------------------------>  a
<------------------------------------------------>  b
<------------------------------------------------>  c
\end{verbatim}

你和我平行，我和他平行，那你和他也平行。

\section{十二、中垂线}

\subsection{什么是中垂线？}

一条直线，从线段的正中间穿过，而且和线段“交叉得很正”。

（具体怎么个“正”法，要等讲完“角”才能说清楚。）

\begin{verbatim}
                    |
                    |  中垂线
                    |
      A             |             B
      ·-------------+-------------·
                    M
                  中点
\end{verbatim}

也叫\textbf{垂直平分线}。

\subsection{中垂线的神奇之处}

\textbf{站在中垂线上的任何一点，到线段两个端点的距离都一样。}

\begin{verbatim}
                    · P
                   / \
                  /   \
                 /     \
      A         /   |   \         B
      ·--------/----+----\--------·

        PA  =  PB
\end{verbatim}

P在中垂线上，所以PA = PB。

不管P在中垂线上的哪个位置，都是这样。

\section{十三、中位线}

\subsection{什么是中位线？}

在三角形里，连接两条边的中点，这条线叫中位线。

\begin{verbatim}
              A
             / \
            /   \
           /     \
          M-------N      <-- 中位线
         /         \
        /           \
       B-------------C
\end{verbatim}

M是AB的中点。

N是AC的中点。

连接M和N，就是中位线。

\subsection{中位线的神奇之处}

\textbf{中位线平行于第三边，而且是第三边的一半。}

MN平行于BC。

MN = BC的一半。

\section*{整章总结}

\begin{verbatim}
三种线
├── 线段：两个端点，到头就停，能量
├── 射线：一个端点，一头不停，有方向
└── 直线：没有端点，两头都不停

怎么写
├── 线段AB = 线段BA（顺序不重要）
├── 射线AB ≠ 射线BA（顺序重要！第一个字母是端点！）
└── 直线AB = 直线BA（顺序不重要）

线段能干嘛
├── 量（三种线里只有线段能量）
├── 比大小
├── 找中点
├── 加减法
└── 延长 / 反向延长

直线的规矩
├── 两点确定一条直线
└── 两直线最多一个交点

两直线的关系
├── 相交：1个公共点
├── 平行：0个公共点
└── 重合：无数个公共点

平行线的规矩
├── 过一点只有一条平行线
├── 距离处处相等
└── 平行可以传递
\end{verbatim}

\chapter{角——两条线之间的“张开程度”}

\section{从一个动作开始}

\begin{quote}
\textbf{做个小实验}

拿两根筷子，一头捏在一起，另一头分开。

慢慢张开——缝隙变大。

慢慢合拢——缝隙变小。

这个“缝隙”的大小，就是\textbf{角}。
\end{quote}

\section{角是什么？}

\subsection{定义}

\textbf{角}：从同一个点出发的两条射线组成的图形。

\begin{verbatim}
        射线A
         /
        /
       O ---------- 射线B
    （顶点）
\end{verbatim}

\begin{center}
\begin{tabular}{|c|l|l|}
\hline
\textbf{部分} & \textbf{说明} & \textbf{筷子类比} \\
\hline
\textbf{顶点} & 两条射线的起点 & 筷子捏在一起的地方 \\
\hline
\textbf{边} & 两条射线 & 两根筷子本身 \\
\hline
\textbf{角的大小} & 两条射线的张开程度 & 筷子之间的缝隙 \\
\hline
\end{tabular}
\end{center}

\begin{quote}
\textbf{记住}：角的大小只看张开程度，和边的长短无关！

（就像剪刀张开$30^\circ$，不管剪刀是大号还是小号，角度都是$30^\circ$）
\end{quote}

\section{角怎么表示？}

\subsection{三个字母法（最规范）}

\[\angle ABC\]

\textbf{读作}：“角ABC”

\textbf{规则}：
\begin{itemize}
\item 中间字母 $=$ 顶点
\item 两边字母 $=$ 边上任意一点
\item $\angle ABC = \angle CBA$（顺序可以换）
\item 不能写成 $\angle BAC$（顶点变了！）
\end{itemize}

\subsection{一个字母法（简化版）}

\[\angle B\]

\textbf{使用条件}：顶点处只有一个角时才能用

\subsection{数字标记法（实用版）}

\[\angle 1,\ \angle 2,\ \angle 3\]

图中标注数字，方便区分多个角。

\section{角的大小怎么量？}

\subsection{单位：度（$^\circ$）}

角的大小用\textbf{度}来衡量，符号：$^\circ$

\begin{center}
\begin{tabular}{|c|l|}
\hline
\textbf{写法} & \textbf{读法} \\
\hline
$30^\circ$ & 三十度 \\
\hline
$90^\circ$ & 九十度 \\
\hline
$180^\circ$ & 一百八十度 \\
\hline
\end{tabular}
\end{center}

\section{角的“家族成员”}

\subsection{按大小分类}

\begin{center}
\begin{tabular}{|c|c|l|l|}
\hline
\textbf{名字} & \textbf{范围} & \textbf{长什么样} & \textbf{生活中的例子} \\
\hline
\textbf{锐角} & $0^\circ < \theta < 90^\circ$ & 尖尖的，比直角瘦 & 屋顶的尖角 \\
\hline
\textbf{直角} & 刚好 $90^\circ$ & 墙角那样方方正正 & 书本的角 \\
\hline
\textbf{钝角} & $90^\circ < \theta < 180^\circ$ & 胖胖的，比直角宽 & 躺椅展开的角度 \\
\hline
\textbf{平角} & 刚好 $180^\circ$ & 一条直线 & 地平线 \\
\hline
\textbf{周角} & 刚好 $360^\circ$ & 转一整圈 & 钟表的时针转一圈 \\
\hline
\end{tabular}
\end{center}

\begin{quote}
\textbf{记忆口诀}：

锐角尖尖小于九（$90^\circ$），

钝角张大超过九，

直角刚好九十度，

平角成线一百八。
\end{quote}

\subsection{特殊角的标记}

\begin{center}
\begin{tabular}{|c|l|}
\hline
\textbf{角的类型} & \textbf{怎么标记} \\
\hline
直角 & 角内画小正方形 $\ulcorner$ \\
\hline
平角 & 看起来像直线 \\
\hline
周角 & 画个圆弧箭头 \\
\hline
\end{tabular}
\end{center}

\section{角的运算}

\subsection{角的加法}

\begin{verbatim}
    C
   /
  /
 / ∠2
O -------- B
 \ ∠1
  \
   A
\end{verbatim}

\[\angle AOC = \angle AOB + \angle BOC\]

\textbf{大白话}：大角 $=$ 小角1 $+$ 小角2

\subsection{角的减法}

\[\angle AOB = \angle AOC - \angle BOC\]

\textbf{大白话}：已知大角和小角1，可以求小角2

\section{两个重要关系}

\subsection{互余（加起来$90^\circ$）}

\[\angle A + \angle B = 90^\circ\]

\textbf{例子}：

\begin{center}
\begin{tabular}{|c|c|c|}
\hline
$\angle A$ & $\angle B$ & \textbf{和} \\
\hline
$30^\circ$ & $60^\circ$ & $90^\circ$ \checkmark \\
\hline
$45^\circ$ & $45^\circ$ & $90^\circ$ \checkmark \\
\hline
$25^\circ$ & $65^\circ$ & $90^\circ$ \checkmark \\
\hline
\end{tabular}
\end{center}

\textbf{求法}：
\[\angle A = 90^\circ - \angle B\]

\begin{quote}
\textbf{记忆提示}：余 $=$ 剩\textbf{余}，离直角还差多少
\end{quote}

\subsection{互补（加起来$180^\circ$）}

\[\angle A + \angle B = 180^\circ\]

\textbf{例子}：

\begin{center}
\begin{tabular}{|c|c|c|}
\hline
$\angle A$ & $\angle B$ & \textbf{和} \\
\hline
$60^\circ$ & $120^\circ$ & $180^\circ$ \checkmark \\
\hline
$90^\circ$ & $90^\circ$ & $180^\circ$ \checkmark \\
\hline
$30^\circ$ & $150^\circ$ & $180^\circ$ \checkmark \\
\hline
\end{tabular}
\end{center}

\textbf{求法}：
\[\angle A = 180^\circ - \angle B\]

\begin{quote}
\textbf{记忆提示}：补 $=$ \textbf{补}齐，补成一条直线
\end{quote}

\subsection{快速区分}

\begin{center}
\begin{tabular}{|c|c|l|c|}
\hline
\textbf{关系} & \textbf{和} & \textbf{联想} & \textbf{公式} \\
\hline
互余 & $90^\circ$ & 剩\textbf{余} & $\alpha + \beta = 90^\circ$ \\
\hline
互补 & $180^\circ$ & \textbf{补}齐直线 & $\alpha + \beta = 180^\circ$ \\
\hline
\end{tabular}
\end{center}

\section{测量工具——量角器}

\subsection{使用三步法}

\begin{verbatim}
第一步：中心点 对准 顶点
第二步：0°线 对准 一条边
第三步：读 另一条边 指向的刻度
\end{verbatim}

\subsection{常见错误}

\textbf{量角器有两圈刻度！}

\begin{center}
\begin{tabular}{|c|l|l|}
\hline
\textbf{刻度圈} & \textbf{方向} & \textbf{什么时候用} \\
\hline
内圈 & $0^\circ \to 180^\circ$（从左往右） & 用内圈的$0^\circ$就读内圈 \\
\hline
外圈 & $0^\circ \to 180^\circ$（从右往左） & 用外圈的$0^\circ$就读外圈 \\
\hline
\end{tabular}
\end{center}

\textbf{规则}：用哪圈的$0^\circ$，就读哪圈的度数

\subsection{画指定角度}

\textbf{例：画$50^\circ$的角}

\begin{center}
\begin{tabular}{|c|l|}
\hline
\textbf{步骤} & \textbf{做什么} \\
\hline
1 & 画一条射线（第一条边） \\
\hline
2 & 量角器中心对准端点 \\
\hline
3 & $0^\circ$线对准射线 \\
\hline
4 & 在$50^\circ$处点一个点 \\
\hline
5 & 从端点过该点画第二条射线 \\
\hline
\end{tabular}
\end{center}

\section{两直线相交的秘密}

\subsection{对顶角}

\begin{verbatim}
    1   2
     \ /
      X
     / \
    4   3
\end{verbatim}

\textbf{性质}：对顶角相等
\begin{itemize}
\item $\angle 1 = \angle 3$
\item $\angle 2 = \angle 4$
\end{itemize}

\textbf{为什么相等？（推导）}

因为：
\begin{itemize}
\item $\angle 1 + \angle 2 = 180^\circ$（邻补角，合起来是平角）
\item $\angle 2 + \angle 3 = 180^\circ$（邻补角，合起来是平角）
\end{itemize}

所以：
\begin{itemize}
\item $\angle 1 + \angle 2 = \angle 2 + \angle 3$
\item 两边同时消去 $\angle 2$
\item $\angle 1 = \angle 3$ \checkmark
\end{itemize}

\begin{quote}
\textbf{大白话}：因为它们都是“$180^\circ$减去同一个角”，所以一样大。
\end{quote}

\subsection{邻补角}

\textbf{性质}：相邻两角互补

\[\angle 1 + \angle 2 = 180^\circ\]

\textbf{原因}：合起来铺满一条直线（平角）

\section{角平分线}

\subsection{定义}

把一个角分成\textbf{两个相等}小角的射线

\subsection{性质}

若 $OB$ 是 $\angle AOC$ 的角平分线：

\[\angle AOB = \angle BOC = \frac{1}{2}\angle AOC\]

或者反过来说：

\[\angle AOC = 2\angle AOB = 2\angle BOC\]

\textbf{例子}：

\begin{center}
\begin{tabular}{|c|c|c|c|}
\hline
$\angle AOC$ & \textbf{角平分线后} & $\angle AOB$ & $\angle BOC$ \\
\hline
$60^\circ$ & 平分 & $30^\circ$ & $30^\circ$ \\
\hline
$90^\circ$ & 平分 & $45^\circ$ & $45^\circ$ \\
\hline
$120^\circ$ & 平分 & $60^\circ$ & $60^\circ$ \\
\hline
\end{tabular}
\end{center}

\section{小测验}

\subsection{判断题}

\begin{enumerate}
\item 平角就是一条直线（\quad）
\item 对顶角一定相等（\quad）
\item 互补的角一定是邻补角（\quad）
\end{enumerate}

\textbf{答案：}

\begin{enumerate}
\item $\times$ 不完全对！平角看起来像直线，但必须有顶点和两条射线
\item \checkmark\ 对（刚才推导过）
\item $\times$ 不对！互补的两个角不一定相邻，只要和是$180^\circ$就行
\end{enumerate}

\subsection{填空题}

\begin{enumerate}
\item $35^\circ$ 的余角是 \underline{\hspace{2cm}}
\item $70^\circ$ 的补角是 \underline{\hspace{2cm}}
\item 两条直线相交，其中一个角是$50^\circ$，则对顶角是\underline{\hspace{2cm}}
\end{enumerate}

\textbf{答案：}

\begin{enumerate}
\item $90^\circ - 35^\circ = 55^\circ$
\item $180^\circ - 70^\circ = 110^\circ$
\item $50^\circ$（对顶角相等）
\end{enumerate}

\subsection{计算题}

已知 $\angle AOC = 120^\circ$，$OB$ 是角平分线，求 $\angle AOB$ 和 $\angle BOC$。

\textbf{答案：}

因为 $OB$ 是角平分线，所以：

\[\angle AOB = \angle BOC = \frac{1}{2}\angle AOC = \frac{120^\circ}{2} = 60^\circ\]

\subsection*{本章知识地图}

\begin{verbatim}
角
├─ 定义：两条射线 + 公共顶点
├─ 表示方法
│   ├─ ∠ABC（三个字母）
│   ├─ ∠B（一个字母）
│   └─ ∠1（数字标记）
├─ 度量单位：度（°）
├─ 分类（按大小）
│   ├─ 锐角（0° ~ 90°）
│   ├─ 直角（= 90°）
│   ├─ 钝角（90° ~ 180°）
│   ├─ 平角（= 180°）
│   └─ 周角（= 360°）
├─ 运算
│   ├─ 加法（大角 = 小角1 + 小角2）
│   └─ 减法（小角 = 大角 - 另一小角）
├─ 关系
│   ├─ 互余（和 = 90°）
│   └─ 互补（和 = 180°）
└─ 特殊情况
    ├─ 对顶角（相等）
    ├─ 邻补角（互补）
    └─ 角平分线（平分成两个相等的角）
\end{verbatim}

\chapter{角的更多规矩}

\section{一、垂直}

\subsection{什么是垂直？}

两条线相交。

如果交出来的角刚好是\textbf{$90^\circ$}，就说这两条线\textbf{垂直}。

\begin{verbatim}
              |
              |
              |
    ----------·----------
              |
              |
              |
\end{verbatim}

\textbf{符号}：$a \perp b$

\textbf{读作}：a垂直于b

\subsection{垂直时，四个角都是$90^\circ$}

两条线相交，本来会有四个角。

如果其中一个是$90^\circ$，那四个全都是$90^\circ$。

\begin{verbatim}
           90°|90°
              |
    ----------·----------
              |
           90°|90°
\end{verbatim}

为什么？

因为对顶角相等，邻补角互补。

一个是$90^\circ$，其他的就全定死了。

\subsection{垂足}

两条垂直的线，交点叫做\textbf{垂足}。

\begin{verbatim}
              |
              |
              |
    ----------·---------- 
              ↑
            垂足
\end{verbatim}

为什么叫“足”？

因为垂线像是“站”在另一条线上。

垂足就是它“站着”的那个点，像脚踩在地上一样。

\subsection{过一点只有一条垂线}

过直线外一点，能画几条垂线？

\textbf{只有一条。}

\begin{verbatim}
              · P
              |
              |  ← 只有这一条垂线
              |
    ----------·----------
            垂足
\end{verbatim}

你想歪一点都不行，歪了就不是$90^\circ$了。

\subsection{垂线段最短}

从直线外一点到直线上，可以画很多条线段。

其中，垂直的那条最短。

\begin{verbatim}
              · P
             /|\
            / | \
           /  |  \
          /   |   \
    -----·----·----·-----
         ↑    ↑    ↑
        长   最短   长
\end{verbatim}

\textbf{中间那条垂直的，最短。}

这个最短的距离，叫做\textbf{点到直线的距离}。

\section{二、平行线被截}

\subsection{什么情况？}

有两条平行线。

又来了第三条线，把它们都穿过去了。

这第三条线叫\textbf{截线}。

\begin{verbatim}
                    截线
                     /
                    /
                   /
    --------------/--------  直线a
                 /
                /
               /
    ----------/------------  直线b
             /
            /
\end{verbatim}

截线和两条平行线相交，产生了8个角。

这8个角之间，有很多神奇的关系。

\subsection{同位角}

\textbf{位置相同的角。}

一个在上面那个交点，一个在下面那个交点。

但它们的位置一模一样——都在截线的同一边，都在直线的同一侧。

\begin{verbatim}
                    /
                 1 /        ← 这个角
    --------------/--------  
                 /
                /
              5/            ← 这个角
    ----------/------------  
             /

    ∠1 和 ∠5 位置相同，都在截线左边，都在直线上方
    它们是同位角
\end{verbatim}

\textbf{同位角一共有4对。}

\subsection{内错角}

\textbf{在两条直线中间，左右交错的角。}

“内”：在两条平行线之间（不是外面）

“错”：一左一右，交错

\begin{verbatim}
                    /
    ---------------/-------  
                 3/            ← 这个角（在两线之间，左边）
                 /
                / 6            ← 这个角（在两线之间，右边）
    -----------/------------  
              /

    ∠3 和 ∠6 是内错角
\end{verbatim}

\textbf{内错角一共有2对。}

\subsection{同旁内角}

\textbf{在两条直线中间，同一边的角。}

“同旁”：在截线的同一边

“内”：在两条平行线之间

\begin{verbatim}
                    /
    ---------------/-------  
                 3/            ← 这个角（在两线之间，左边）
                 /
                /
              5/               ← 这个角（在两线之间，也是左边）
    ----------/------------  
             /

    ∠3 和 ∠5 是同旁内角
\end{verbatim}

\textbf{同旁内角一共有2对。}

\subsection{怎么记住这三种角？}

\textbf{同位角}——“同样位置”

想象两层楼，楼上楼下位置一样的房间。

\textbf{内错角}——“内部交错”

在里面，一左一右，交错开。

\textbf{同旁内角}——“内部同旁”

在里面，都在同一边。

\subsection{平行线的三条铁律}

\textbf{如果两条直线平行：}

\begin{center}
\begin{tabular}{|c|c|}
\hline
角的类型 & 关系 \\
\hline
同位角 & \textbf{相等} \\
\hline
内错角 & \textbf{相等} \\
\hline
同旁内角 & \textbf{互补}（加起来$180^\circ$） \\
\hline
\end{tabular}
\end{center}

\textbf{反过来也成立：}

如果同位角相等 $\to$ 两直线平行

如果内错角相等 $\to$ 两直线平行

如果同旁内角互补 $\to$ 两直线平行

\subsection{这些铁律用来干嘛？}

\textbf{两个方向：}

\textbf{第一个方向：已经知道两条线平行，求角是多少度。}

比如告诉你 $a \parallel b$，$\angle 1 = 50^\circ$，问$\angle 5$是多少度。

因为$\angle 1$和$\angle 5$是同位角，平行时同位角相等，所以$\angle 5 = 50^\circ$。

\textbf{第二个方向：已经知道角的度数，证明两条线平行。}

比如告诉你$\angle 3 = \angle 6$，问a和b是不是平行。

因为$\angle 3$和$\angle 6$是内错角，内错角相等，所以 $a \parallel b$。

\textbf{一句话：}

已知平行 $\to$ 求角

已知角 $\to$ 证平行

\section{三、三角形的角}

刚才讲的是两条线之间的角。

现在来看一个由三条线围起来的图形——三角形。

三角形里面的角，有自己的一套规矩。

\subsection{“里面”和“外面”}

画一个三角形。

三条边围起来，中间那块地方，是三角形的\textbf{里面}。

除了里面之外的所有地方，是三角形的\textbf{外面}。

\begin{verbatim}
    外面    A     外面
           /\
          /  \
         /    \
        / 里面 \
       /        \
      B----------C
    外面   外面   外面
\end{verbatim}

\subsection{内角}

看三角形的每个顶点。

每个顶点都有一个角。

这个角\textbf{朝向三角形的里面}。

\begin{verbatim}
              A
             /\
            /↙ \      ← ∠A 朝向里面
           /    \
          /      \
         /        \
        B----------C
       ↗            ↖
    ∠B朝向里面    ∠C朝向里面
\end{verbatim}

\textbf{$\angle A$、$\angle B$、$\angle C$ 都朝向三角形的里面。}

\textbf{所以它们叫内角。}

内 = 里面

内角 = 开口朝向里面的角

\subsection{内角和 = $180^\circ$}

三角形有三个内角。

不管三角形长什么样，三个内角加起来永远等于\textbf{$180^\circ$}。

\begin{verbatim}
              A
             /\
            /  \
           /    \
          /      \
         /        \
        B----------C

    ∠A + ∠B + ∠C = 180°
\end{verbatim}

\textbf{永远是$180^\circ$，没有例外。}

\section{四、外角是什么？怎么判定？}

\subsection{外角 = 在外面的角}

把三角形的一条边延长。

\begin{verbatim}
              A
             /\
            /  \
           /    \
          /      \
         /        \
        B----------C---------D
\end{verbatim}

看C这个顶点。

原来的内角$\angle ACB$朝向里面。

现在多出来一个角$\angle ACD$，它\textbf{朝向三角形的外面}。

\begin{verbatim}
              A
             /\
            /  \
           /    \
          /      \
         /    内角\外角
        B----------C---------D
                   ↑
              ∠ACB朝里面（内角）
              ∠ACD朝外面（外角）
\end{verbatim}

\textbf{$\angle ACD$ 朝向三角形的外面。}

\textbf{所以它叫外角。}

外 = 外面

外角 = 开口朝向外面的角

\subsection{外角是怎么来的？}

外角不是凭空出现的。

它是你\textbf{延长了一条边}之后才出现的。

\textbf{没延长之前：}

\begin{verbatim}
              A
             /\
            /  \
           /    \
          /      \
         /        \
        B----------C
\end{verbatim}

只有三个内角，没有外角。

\textbf{延长BC之后：}

\begin{verbatim}
              A
             /\
            /  \
           /    \
          /      \
         /        \
        B----------C---------D
                   ↑
               外角出现了
\end{verbatim}

延长之后，C点就多了一个角——外角。

\textbf{没有延长，就没有外角。}

\subsection{外角的定义}

\begin{quote}
\textbf{外角：三角形的一条边和另一条边的延长线组成的角。}
\end{quote}

外角必须是：

\textbf{一条边} + \textbf{另一条边的延长线}

组成的角。

\begin{verbatim}
              A
             /\
            /  \
           /    \
          /      \
         /        \
        B----------C---------D
        
    CA 是一条边
    CD 是 BC 的延长线
    
    ∠ACD 就是外角
\end{verbatim}

\subsection{怎么判定一个角是不是外角？}

\textbf{检查三件事：}

第一：这个角的顶点是不是三角形的顶点？

第二：这个角的一条边是不是三角形的边？

第三：这个角的另一条边是不是三角形某条边的延长线？

\textbf{三个都满足，才是外角。}

\textbf{举个“是”的例子：}

\begin{verbatim}
              A
             /\
            /  \
           /    \
          /      \
         /        \
        B----------C---------D
        
        问：∠ACD 是外角吗？
\end{verbatim}

\checkmark\ 顶点C是三角形的顶点

\checkmark\ CA是三角形的边

\checkmark\ CD是BC的延长线

\textbf{是外角。}

\textbf{再举个“不是”的例子：}

\begin{verbatim}
              A
             /\
            /  \
           /    \
          /      \
         /        \
        B----------C
                    \
                     \
                      D
\end{verbatim}

问：$\angle BCD$是外角吗？

\checkmark\ 顶点C是三角形的顶点

\checkmark\ CB是三角形的边

$\times$ CD不是任何边的延长线（CD是随便画的一条线）

\textbf{不是外角。}

\subsection{每个顶点有几个外角？}

\textbf{两个。}

每个顶点可以往两个方向延长。

\begin{verbatim}
              A
             /\
            /  \
           /    \
          /      \
         /        \
    E---B----------C---D
        ↑          ↑
      往这边延长   往这边也能延长
\end{verbatim}

\textbf{但这两个外角相等。}

因为它们是\textbf{对顶角}。

所以虽然每个顶点有两个外角，但它们大小一样。

\textbf{说“某个顶点的外角”，就是指这个度数。}

\subsection{外角和内角的位置关系}

外角和内角挨在一起。

共用一条边。

一个朝里，一个朝外。

\begin{verbatim}
              A
             /\
            /  \
           /    \
          /      \
         /   内角 \ 外角
        B----------C---------D
                   ↑
              它们挨在一起
              共用边CA
\end{verbatim}

\textbf{外角和内角加起来是$180^\circ$。}

因为它们合起来是一条直线。

\textbf{它们是邻补角。}

\section{五、外角的铁律}

\subsection{铁律一：外角 = 不相邻的两个内角之和}

\begin{verbatim}
              A
             /\
            /  \
           /    \
          /      \
         /        \
        B----------C---------D
\end{verbatim}

$\angle ACD$ 是 C 点的外角。

\textbf{什么叫“相邻”？什么叫“不相邻”？}

看C这个顶点。

$\angle ACB$挨着C，就在C的旁边。这叫\textbf{相邻}的内角。

$\angle A$在A那边，$\angle B$在B那边，它们都\textbf{不挨着C}。这叫\textbf{不相邻}的内角。

\textbf{铁律：外角 = 不相邻的两个内角之和}

\textbf{$\angle ACD = \angle A + \angle B$}

\textbf{为什么？}

$\angle ACD + \angle ACB = 180^\circ$（内角和外角是邻补角）

$\angle A + \angle B + \angle ACB = 180^\circ$（三角形内角和）

两个式子右边都是$180^\circ$，所以：

$\angle ACD + \angle ACB = \angle A + \angle B + \angle ACB$

两边都减去 $\angle ACB$：

\textbf{$\angle ACD = \angle A + \angle B$}

\subsection{铁律二：外角 > 任何一个不相邻的内角}

外角 = $\angle A + \angle B$

那外角肯定比$\angle A$大，也比$\angle B$大。

因为两个正数加起来，肯定比其中任何一个都大。

\textbf{$\angle ACD > \angle A$}

\textbf{$\angle ACD > \angle B$}

\subsection{铁律三：三个外角之和 = $360^\circ$}

每个顶点取一个外角。

三个顶点，三个外角。

加起来永远等于\textbf{$360^\circ$}。

\textbf{为什么？这里要慢慢讲。}

\textbf{第一步：每个顶点，内角 + 外角 = $180^\circ$}

\begin{verbatim}
      内角  外角
        ↘  ↙
    -----C-----
\end{verbatim}

在C这个顶点，内角和外角挨着，合起来是一条直线，所以加起来是$180^\circ$。

其他两个顶点也一样。

\textbf{第二步：三角形有三个顶点}

A点：内角A + 外角A = $180^\circ$

B点：内角B + 外角B = $180^\circ$

C点：内角C + 外角C = $180^\circ$

\textbf{第三步：三个等式加起来}

(内角A + 外角A) + (内角B + 外角B) + (内角C + 外角C) = $180^\circ + 180^\circ + 180^\circ$

整理一下左边：

(内角A + 内角B + 内角C) + (外角A + 外角B + 外角C) = $540^\circ$

\textbf{第四步：内角A + 内角B + 内角C = $180^\circ$}

这是三角形内角和。

\textbf{第五步：把第四步代入第三步}

$180^\circ$ + (外角A + 外角B + 外角C) = $540^\circ$

两边减去$180^\circ$：

\textbf{外角A + 外角B + 外角C = $360^\circ$}

\textbf{所以三个外角加起来是$360^\circ$。}

\textbf{换一种方式理解：走一圈}

你站在操场上，面朝前方。

现在你要转一整圈，回到原来的方向。

你转了多少度？$360^\circ$。

想象你沿着三角形的边走一圈。

每到一个顶点，你就要转一个弯。

转的那个角度，就是外角。

\begin{verbatim}
              A
             /\
            /  \   ← 走到A，转一个弯
           /    \
          /      \
         /        \
        B----------C
        ↑          ↑
    走到B，转一个弯   走到C，转一个弯
\end{verbatim}

走完三个顶点，转了三个弯。

回到起点，刚好转了一整圈。

\textbf{一整圈 = $360^\circ$}

\textbf{所以三个外角加起来 = $360^\circ$}

\section{六、多边形的角}

\subsection{内角和公式}

三角形：3条边，内角和 = $180^\circ$

四边形：4条边，内角和 = $360^\circ$

五边形：5条边，内角和 = $540^\circ$

六边形：6条边，内角和 = $720^\circ$

\textbf{规律：}

n边形的内角和 = \textbf{$(n - 2)\times 180^\circ$}

\textbf{为什么是 $n - 2$？}

任何多边形都可以从一个顶点出发，往其他顶点画线，把它分成几个三角形。

比如四边形：

\begin{verbatim}
        A---------B
        |\        |
        | \       |
        |  \      |
        |   \     |
        |    \    |
        D---------C
\end{verbatim}

从A出发，往C画一条线，把四边形分成了2个三角形。

4条边 $\to$ 分成 $4-2=2$ 个三角形 \checkmark

每个三角形内角和$180^\circ$，所以四边形内角和 = $2 \times 180^\circ = 360^\circ$

再比如五边形：

从一个顶点出发，往其他顶点画线，能把它分成3个三角形。

5条边 $\to$ 分成 $5-2=3$ 个三角形 \checkmark

五边形内角和 = $3 \times 180^\circ = 540^\circ$

\textbf{规律：n边形 $\to$ 分成 $n-2$ 个三角形 $\to$ 内角和 = $(n-2)\times 180^\circ$}

\subsection{外角和永远是$360^\circ$}

不管是什么多边形，外角和永远等于\textbf{$360^\circ$}。

三角形外角和 = $360^\circ$

四边形外角和 = $360^\circ$

五边形外角和 = $360^\circ$

一百边形外角和 = $360^\circ$

\textbf{永远是$360^\circ$，和边数无关。}

\textbf{为什么？还是走一圈。}

不管多边形有多少条边，你沿着它走一圈，最后都回到起点，面朝原来的方向。

你总共转了多少？

$360^\circ$。

一整圈。

每个顶点转一个弯，所有弯加起来就是$360^\circ$。

每个弯就是外角。

\textbf{所以外角和 = $360^\circ$，永远不变。}

\subsection{正多边形的每个外角}

如果是\textbf{正多边形}（每条边一样长，每个角一样大）：

每个外角都相等。

既然外角和是$360^\circ$，那：

每个外角 = $360^\circ \div$ 边数

正三角形：每个外角 = $360^\circ \div 3 = 120^\circ$

（验证：正三角形每个内角$60^\circ$，$60^\circ + 120^\circ = 180^\circ$ \checkmark）

正四边形（正方形）：每个外角 = $360^\circ \div 4 = 90^\circ$

（验证：正方形每个内角$90^\circ$，$90^\circ + 90^\circ = 180^\circ$ \checkmark）

正五边形：每个外角 = $360^\circ \div 5 = 72^\circ$

正六边形：每个外角 = $360^\circ \div 6 = 60^\circ$

\subsection{内角和 vs 外角和}

\begin{center}
\begin{tabular}{|c|c|c|c|}
\hline
 & 三角形 & 四边形 & n边形 \\
\hline
\textbf{内角和} & $180^\circ$ & $360^\circ$ & $(n-2)\times 180^\circ$ \\
\hline
\textbf{外角和} & $360^\circ$ & $360^\circ$ & \textbf{永远$360^\circ$} \\
\hline
\end{tabular}
\end{center}

内角和随着边数增加而增加。

外角和永远不变，都是$360^\circ$。

\section{七、余角和补角的更多规矩}

\subsection{等角的余角相等}

$\angle A = 30^\circ$，$\angle A$的余角 = $60^\circ$

$\angle B = 30^\circ$，$\angle B$的余角 = $60^\circ$

\textbf{$\angle A$和$\angle B$相等，它们的余角也相等。}

\textbf{为什么？}

$\angle A$的余角 = $90^\circ - \angle A = 90^\circ - 30^\circ = 60^\circ$

$\angle B$的余角 = $90^\circ - \angle B = 90^\circ - 30^\circ = 60^\circ$

算法一样，减的数一样，结果当然一样。

\textbf{所以只要两个角相等，它们的余角就相等。}

\subsection{等角的补角相等}

$\angle A = 30^\circ$，$\angle A$的补角 = $150^\circ$

$\angle B = 30^\circ$，$\angle B$的补角 = $150^\circ$

\textbf{$\angle A$和$\angle B$相等，它们的补角也相等。}

\textbf{为什么？}

$\angle A$的补角 = $180^\circ - \angle A = 180^\circ - 30^\circ = 150^\circ$

$\angle B$的补角 = $180^\circ - \angle B = 180^\circ - 30^\circ = 150^\circ$

同样的道理，算法一样，结果一样。

\textbf{所以只要两个角相等，它们的补角也相等。}

\textbf{一句话：等角的余角相等，等角的补角相等。}

\section*{整章总结}

\begin{verbatim}
角的更多规矩
│
├── 垂直
│   ├── 两线相交成90°
│   ├── 符号：⊥
│   ├── 垂足：交点，像脚踩在地上
│   ├── 四个角都是90°
│   ├── 过一点只有一条垂线
│   └── 垂线段最短（点到直线的距离）
│
├── 平行线被截
│   ├── 同位角：位置相同 → 相等
│   ├── 内错角：内部交错 → 相等
│   ├── 同旁内角：内部同侧 → 互补
│   └── 用法：已知平行→求角，已知角→证平行
│
├── 内角
│   ├── 开口朝向图形里面
│   └── 三角形内角和 = 180°
│
├── 外角
│   ├── 开口朝向图形外面
│   ├── 来源：延长一条边才出现
│   ├── 判定：顶点是三角形的 + 一条边是三角形的 + 另一条边是延长线
│   ├── 内角 + 外角 = 180°（邻补角）
│   ├── 每个顶点两个外角，但它们相等（对顶角）
│   ├── 外角 = 不相邻的两个内角之和
│   ├── 外角 > 任何一个不相邻的内角
│   └── 三个外角之和 = 360°（走一圈）
│
├── 多边形的角
│   ├── 内角和 = (n-2) × 180°（分成n-2个三角形）
│   ├── 外角和 = 360°（永远，走一圈）
│   └── 正n边形每个外角 = 360° ÷ n
│
└── 余角和补角的规矩
    ├── 等角的余角相等（算法一样，结果一样）
    └── 等角的补角相等（算法一样，结果一样）
\end{verbatim}

\chapter{图形——它们没你想的那么复杂}

\section{先说一件事}

你可能在小学，已经见过这些图形了。

长方形，正方形，三角形，圆形。

但你可能只知道：“哦，这是长方形。”

今天，我们不只是认识它们，我们要\textbf{彻底搞清楚它们的本质}。

还有，它们的面积和周长，到底是怎么来的。

不是让你背公式。

是让你\textbf{理解为什么公式长这样}。

\section{先说几个基础概念}

在我们认识图形之前，先认识三个东西。不认识它们，后面会很懵。

\subsection{边是什么？}

边，就是图形的“围墙”。

一条直线连着两个点，就是一条边。

\begin{center}
\begin{tabular}{|c|c|l|}
\hline
\textbf{图形} & \textbf{边的数量} & \textbf{说明} \\
\hline
三角形 & 3条 & 最少的多边形 \\
\hline
长方形 & 4条 & 对边相等 \\
\hline
圆形 & 0条 & 只有曲线，没有直边 \\
\hline
\end{tabular}
\end{center}

圆形没有边。圆形没有直线，没有角，只有一条完美的曲线，从头到尾连在一起，找不到起点，也找不到终点。

圆形有的，是\textbf{圆周}——那条曲线绕圆一圈的总长度。

\subsection{角是什么？}

（还记得上一章吗？我们来快速回顾一下）

两条边，从同一个点出发，往不同方向伸出去，中间夹出来的“开口”，就是一个角。

\begin{itemize}
\item 开口越大，角越大
\item 开口越小，角越小
\end{itemize}

\subsection{度数是什么？}

度数，就是用来量角的大小的单位。用符号 $^\circ$ 表示。

\begin{center}
\begin{tabular}{|c|c|l|}
\hline
\textbf{度数} & \textbf{名称} & \textbf{形象理解} \\
\hline
$0^\circ$ & 零角 & 两条边完全重合 \\
\hline
$90^\circ$ & 直角 & 一个“L”的形状 \\
\hline
$180^\circ$ & 平角 & 变成一条直线 \\
\hline
$360^\circ$ & 周角 & 转了一整圈 \\
\hline
\end{tabular}
\end{center}

\begin{quote}
\textbf{用一句话记住它们}：

想象你站在原地，开始转圈：
\begin{itemize}
\item 转到右边垂直的位置，转了 \textbf{$90^\circ$}
\item 转到正后方，转了 \textbf{$180^\circ$}
\item 转回原来的方向，转了 \textbf{$360^\circ$}
\end{itemize}
\end{quote}

\subsection{角和图形的关系}

每个图形，所有角加起来，都有固定的总度数。

\begin{center}
\begin{tabular}{|c|c|l|}
\hline
\textbf{图形} & \textbf{内角和} & \textbf{说明} \\
\hline
三角形 & $180^\circ$ & 不管什么形状，永远是$180^\circ$ \\
\hline
长方形/正方形 & $360^\circ$ & $90^\circ \times 4$ \\
\hline
任意四边形 & $360^\circ$ & 永远是$360^\circ$ \\
\hline
圆形 & 没有角 & 但圆心角是$360^\circ$ \\
\hline
\end{tabular}
\end{center}

\textbf{一个小练习}：

三角形三个角加起来是 $180^\circ$。如果一个三角形，两个角分别是 $60^\circ$ 和 $80^\circ$，第三个角是几度？

\[180^\circ - 60^\circ - 80^\circ = 40^\circ\]

就这么简单。你已经会算角了。

\section{第一个：长方形}

长方形，你肯定见过。

课本是长方形，门是长方形，手机屏幕也是长方形。

\textbf{长方形的特点：}

\begin{center}
\begin{tabular}{|l|l|}
\hline
\textbf{特点} & \textbf{说明} \\
\hline
四个角 & 全部是直角（$90^\circ$） \\
\hline
两条长边 & 一样长，叫\textbf{长}（$a$） \\
\hline
两条短边 & 一样长，叫\textbf{宽}（$b$） \\
\hline
\end{tabular}
\end{center}

\subsection{长方形的周长}

周长是什么？就是绕着边走一圈，走了多远。

走一圈，要走两条长，两条宽：

\[\text{周长} = \text{长} + \text{宽} + \text{长} + \text{宽} = 2 \times \text{长} + 2 \times \text{宽} = 2(a + b)\]

\begin{quote}
\textbf{语言记忆法}：

“一条长加一条宽，只走了两条边。另外两条边一样长，所以乘以2。”
\end{quote}

\subsection{长方形的面积}

面积是什么？就是这个图形，\textbf{里面有多少个小方格}。

想象把长方形切成很多个 $1 \times 1$ 的小正方形。

每行有 $a$ 个，一共有 $b$ 行。总共有多少个？

\[\text{面积} = a \times b\]

不是规定，是数格子数出来的。

\begin{quote}
\textbf{语言记忆法}：

“横着数几个格子，竖着数几排，乘起来就是总共几个格子。”
\end{quote}

\subsection{例题}

\textbf{题目：} “一个长方形，长 8 厘米，宽 5 厘米。周长和面积各是多少？”

\textbf{周长：}
\[2 \times (8 + 5) = 2 \times 13 = 26 \text{ 厘米}\]

\textbf{面积：}
\[8 \times 5 = 40 \text{ 平方厘米}\]

（注意：面积的单位是\textbf{平方}厘米，不是厘米。因为面积是二维的，是格子数，不是长度。）

\section{第二个：正方形}

正方形，就是一个“特别的长方形”。

特别在哪里？\textbf{四条边全部一样长。}

所以，正方形的公式，就是把长方形的公式里的长和宽，换成同一个字母 $a$：

\begin{center}
\begin{tabular}{|c|c|l|}
\hline
\textbf{项目} & \textbf{公式} & \textbf{语言记忆法} \\
\hline
周长 & $4a$ & “四条边一样长，加四次，直接乘以4” \\
\hline
面积 & $a^2$ & “边长乘边长，就是里面有多少个小正方形。所以叫平方” \\
\hline
\end{tabular}
\end{center}

看见了吗？面积 $= a^2$。次方出现了！

这就是为什么 $a^2$ 叫做 “$a$ 的平方”——因为它本来就是\textbf{正方形的面积}。

（原来次方和图形是亲戚，你发现了吗？）

\subsection{例题}

\textbf{题目：} “一个正方形，边长 6 厘米。周长和面积各是多少？”

\textbf{周长：}
\[4 \times 6 = 24 \text{ 厘米}\]

\textbf{面积：}
\[6^2 = 36 \text{ 平方厘米}\]

\section{第三个：三角形}

三角形，三条边，三个角。

\subsection{三角形的家族}

\begin{center}
\begin{tabular}{|l|l|}
\hline
\textbf{类型} & \textbf{特点} \\
\hline
等边三角形 & 三条边一样长，三个角都是 $60^\circ$ \\
\hline
等腰三角形 & 两条边一样长，两个角一样大 \\
\hline
直角三角形 & 有一个直角（$90^\circ$） \\
\hline
普通三角形 & 三条边都不一样长 \\
\hline
\end{tabular}
\end{center}

\subsection{三角形的周长}

绕着边走一圈：

\[\text{周长} = a + b + c\]

\begin{quote}
\textbf{语言记忆法}：

“三条边，没有重复，一条一条加起来就行。”
\end{quote}

\subsection{三角形的面积}

\[\text{面积} = \frac{\text{底} \times \text{高}}{2}\]

\textbf{为什么要除以 2？}

想象一个长方形，从对角线切一刀，变成两个一样的三角形。

每个三角形，就是长方形面积的一半。

\[\text{三角形面积} = \frac{\text{长方形面积}}{2} = \frac{\text{底} \times \text{高}}{2}\]

不是规定，是切出来的。

\begin{quote}
\textbf{语言记忆法}：

“把三角形想成长方形切了一半。底乘高是长方形的面积，除以2就是三角形的。”
\end{quote}

\subsection{什么是“高”？}

高，是从顶点，\textbf{垂直}画一条线，到对边的距离。

注意：必须垂直，必须和底边成 $90^\circ$。

\subsection{例题}

\textbf{题目：} “一个三角形，底边 10 厘米，高 6 厘米。面积是多少？”

\[\text{面积} = \frac{10 \times 6}{2} = \frac{60}{2} = 30 \text{ 平方厘米}\]

\subsection{特别说一下：直角三角形和勾股定理}

直角三角形有个特别的规律，叫\textbf{勾股定理}：

\[a^2 + b^2 = c^2\]

\begin{center}
\begin{tabular}{|c|l|}
\hline
\textbf{符号} & \textbf{代表什么} \\
\hline
$a$、$b$ & 两条直角边 \\
\hline
$c$ & 斜边（最长的那条） \\
\hline
\end{tabular}
\end{center}

\textbf{例子：}

直角三角形，两条直角边是 3 和 4，斜边是几？

\[3^2 + 4^2 = c^2\]
\[9 + 16 = c^2\]
\[c^2 = 25\]
\[c = 5\]

斜边是 5。

\begin{quote}
\textbf{3、4、5 是最经典的直角三角形。记住它，以后会经常用到。}
\end{quote}

\section{第四个：圆形}

圆形，是所有图形里最特别的一个。

没有角，没有边，没有起点，没有终点，只有一条完美的曲线，永远连着自己。

\subsection{圆的基本概念}

\begin{center}
\begin{tabular}{|c|c|l|}
\hline
\textbf{概念} & \textbf{符号} & \textbf{说明} \\
\hline
圆心 & $O$ & 圆的正中间那个点 \\
\hline
半径 & $r$ & 从圆心到圆上任意一点的距离 \\
\hline
直径 & $d$ & 穿过圆心的线段，$d = 2r$ \\
\hline
圆周 & — & 圆的那条曲线，绕圆一圈的总长度 \\
\hline
\end{tabular}
\end{center}

\begin{quote}
\textbf{圆的本质}：圆上任意一点到圆心的距离，都一样。
\end{quote}

\subsection{认识 $\pi$}

讲圆之前，得先认识一个特别的数：$\pi$（读作“派”）

古人用绳子绕着圆量，发现了一件神奇的事：

\textbf{不管圆有多大，圆周除以直径，永远等于同一个数。}

\[\pi \approx 3.14159265358979...\]

小数点后面无穷无尽，永远不会重复，永远不会结束。

我们平时用 $\pi \approx 3.14$ 就够了。

\subsection{圆的周长}

既然 $\dfrac{\text{圆周}}{\text{直径}} = \pi$，那：

\[\text{圆周} = \pi \times \text{直径} = \pi \times 2r = 2\pi r\]

\begin{quote}
\textbf{语言记忆法}：

“直径乘以$\pi$，就是绕圆一圈的长度。直径是两个半径，所以是$2\pi r$。”
\end{quote}

\subsection{圆的面积}

\[\text{面积} = \pi r^2\]

\textbf{为什么？}

想象把圆切成很多很多个小扇形，就像切披萨一样。

然后把这些小扇形，一个朝上一个朝下，拼在一起。

拼出来的形状，越来越接近一个长方形。

\begin{center}
\begin{tabular}{|l|l|}
\hline
\textbf{这个长方形的...} & \textbf{等于...} \\
\hline
长 & 半个圆周 $= \pi r$ \\
\hline
宽 & 半径 $= r$ \\
\hline
\end{tabular}
\end{center}

所以面积：$\pi r \times r = \pi r^2$

\begin{quote}
\textbf{语言记忆法}：

“把圆切成披萨，摊平拼成长方形。长是半个圆周$\pi r$，宽是半径$r$，乘起来就是$\pi r^2$。”
\end{quote}

（切披萨，切出了面积公式。数学家连吃饭都不忘研究数学。）

\subsection{弧长}

还记得圆心角吗？从圆心出发，画两条半径到圆上，中间夹出来的角。

\[\text{弧长} = \frac{\text{圆心角}}{360^\circ} \times 2\pi r\]

\begin{quote}
\textbf{语言记忆法}：

“圆心角占整个圆的几分之几，弧长就是圆周的几分之几。”
\end{quote}

\subsection{例题}

\textbf{题目一：} “一个圆，半径 5 厘米。周长和面积各是多少？（$\pi$ 取 $3.14$）”

\textbf{周长：}
\[2 \times 3.14 \times 5 = 31.4 \text{ 厘米}\]

\textbf{面积：}
\[3.14 \times 5^2 = 3.14 \times 25 = 78.5 \text{ 平方厘米}\]

\textbf{题目二：} “一个圆，半径 10 厘米，圆心角 $90^\circ$。这段弧长是多少？（$\pi$ 取 $3.14$）”

\[\text{弧长} = \frac{90^\circ}{360^\circ} \times 2 \times 3.14 \times 10 = \frac{1}{4} \times 62.8 = 15.7 \text{ 厘米}\]

\section{总结一下}

\begin{center}
\begin{tabular}{|c|c|c|l|}
\hline
\textbf{图形} & \textbf{周长} & \textbf{面积} & \textbf{特别说明} \\
\hline
长方形 & $2(a + b)$ & $a \times b$ & 四个直角 \\
\hline
正方形 & $4a$ & $a^2$ & 四条边一样长 \\
\hline
三角形 & $a + b + c$ & $\dfrac{a \times h}{2}$ & 三角之和 $180^\circ$ \\
\hline
圆形 & $2\pi r$ & $\pi r^2$ & 没有边，只有圆周 \\
\hline
\end{tabular}
\end{center}

\section{现在你来}

\begin{quote}
\textbf{第一道：} “一个长方形，长 12 厘米，宽 7 厘米。周长和面积各是多少？”
\end{quote}

\begin{quote}
\textbf{第二道：} “一个正方形，边长 9 厘米。周长和面积各是多少？”
\end{quote}

\begin{quote}
\textbf{第三道：} “一个三角形，底边 8 厘米，高 5 厘米。面积是多少？”
\end{quote}

\begin{quote}
\textbf{第四道：} “一个圆，半径 3 厘米。周长和面积各是多少？（$\pi$ 取 $3.14$）”
\end{quote}

\begin{quote}
\textbf{第五道：} “一个直角三角形，两条直角边分别是 5 厘米和 12 厘米。斜边是多少？”
\end{quote}

\begin{quote}
\textbf{第六道：} “一个圆，半径 6 厘米，圆心角 $120^\circ$。这段弧长是多少？（$\pi$ 取 $3.14$）”
\end{quote}

写下来，验证一下。

\subsection*{本章知识地图}

\begin{verbatim}
平面图形
│
├─ 基础概念
│   ├─ 边（图形的围墙）
│   ├─ 角（两条边的开口）
│   └─ 度数（量角的单位）
│
├─ 长方形
│   ├─ 周长 = 2(a + b)
│   └─ 面积 = a × b
│
├─ 正方形
│   ├─ 周长 = 4a
│   └─ 面积 = a²（所以叫“平方”）
│
├─ 三角形
│   ├─ 周长 = a + b + c
│   ├─ 面积 = 底 × 高 ÷ 2
│   └─ 勾股定理：a² + b² = c²
│
└─ 圆形
    ├─ 周长 = 2πr
    ├─ 面积 = πr²
    └─ 弧长 = (圆心角/360°) × 2πr
\end{verbatim}

\chapter{立体图形——从平面走进空间}

\section{从纸上走出来}

上一章，我们在纸上画图。

长方形，正方形，三角形，圆形。全部都是平的，二维的。

但真实的世界，不是平的。

你住的房子，你喝水的杯子，你踢的足球，全部都是\textbf{立体的}，三维的。

这一章，我们从纸上走出来，进入真实的空间。

\section{先说两个新概念}

\subsection{表面积是什么？}

表面积，就是把一个立体图形，\textbf{所有的面加起来的总面积}。

就像你要给一个盒子包装纸，需要多少包装纸，就是表面积。

\subsection{体积是什么？}

体积，就是这个立体图形，\textbf{里面能装多少东西}。

就像你往杯子里倒水，能倒多少，就是体积。

\begin{center}
\begin{tabular}{|c|c|c|}
\hline
\textbf{概念} & \textbf{关注的是} & \textbf{单位} \\
\hline
表面积 & 外面 & 平方厘米 \\
\hline
体积 & 里面 & 立方厘米 \\
\hline
\end{tabular}
\end{center}

记住这个区别，后面不会混。

\section{第一个：长方体}

长方体，就是长方形变成了立体。

你家的冰箱，砖头，书包，都是长方体。

\subsection{长方体有什么？}

\begin{center}
\begin{tabular}{|c|c|l|}
\hline
\textbf{部分} & \textbf{数量} & \textbf{说明} \\
\hline
面 & 6个 & 全部是长方形（或正方形） \\
\hline
边 & 12条 & 分成三组，每组4条 \\
\hline
顶点 & 8个 & 就是角 \\
\hline
\end{tabular}
\end{center}

三个方向的尺寸：\textbf{长} $a$，\textbf{宽} $b$，\textbf{高} $c$。

\subsection{长方体的表面积}

6个面，两两相对，每对面积一样：

\begin{center}
\begin{tabular}{|c|c|c|}
\hline
\textbf{哪两个面} & \textbf{每个面积} & \textbf{两个加起来} \\
\hline
前面和后面 & $a \times c$ & $2ac$ \\
\hline
左面和右面 & $b \times c$ & $2bc$ \\
\hline
上面和下面 & $a \times b$ & $2ab$ \\
\hline
\end{tabular}
\end{center}

加起来：

\[\text{表面积} = 2ab + 2bc + 2ac = 2(ab + bc + ac)\]

\begin{quote}
\textbf{语言记忆法}：

“六个面，两两配对，三对不一样的面。每对算一次，乘以2，加起来就是总表面积。”
\end{quote}

\subsection{长方体的体积}

体积，就是把长方体切成很多个 $1 \times 1 \times 1$ 的小正方体，数数有多少个。

底面有 $a \times b$ 个，一共叠了 $c$ 层，所以：

\[\text{体积} = a \times b \times c\]

\begin{quote}
\textbf{语言记忆法}：

“长乘宽是底面积，再乘高，就是叠了几层。”
\end{quote}

\subsection{例题}

\textbf{题目：} “一个长方体，长 5 厘米，宽 3 厘米，高 4 厘米。表面积和体积各是多少？”

\textbf{表面积：}
\[2 \times (5 \times 3 + 3 \times 4 + 5 \times 4) = 2 \times (15 + 12 + 20) = 2 \times 47 = 94 \text{ 平方厘米}\]

\textbf{体积：}
\[5 \times 3 \times 4 = 60 \text{ 立方厘米}\]

\section{第二个：正方体}

正方体，就是“特别的长方体”。

特别在哪里？\textbf{六个面全部是正方形，所有边一样长。}

骰子，魔方，都是正方体。

\begin{center}
\begin{tabular}{|c|c|l|}
\hline
\textbf{项目} & \textbf{公式} & \textbf{语言记忆法} \\
\hline
表面积 & $6a^2$ & “六个面，每个面都是正方形，面积都是$a^2$。六个加起来就是$6a^2$” \\
\hline
体积 & $a^3$ & “边长乘边长乘边长，三个方向各乘一次。所以叫立方” \\
\hline
\end{tabular}
\end{center}

看见了吗？$a^3$ 叫 “$a$ 的立方”——因为它本来就是\textbf{正方体的体积}！

（$a^2$ 是正方形面积，$a^3$ 是正方体体积。次方果然和图形是亲戚。）

\subsection{例题}

\textbf{题目：} “一个正方体，边长 4 厘米。表面积和体积各是多少？”

\textbf{表面积：}
\[6 \times 4^2 = 6 \times 16 = 96 \text{ 平方厘米}\]

\textbf{体积：}
\[4^3 = 64 \text{ 立方厘米}\]

\section{第三个：圆柱体}

圆柱体，就是圆形变成了立体。

水桶，罐头，蜡烛，都是圆柱体。

\subsection{圆柱体有什么？}

\begin{center}
\begin{tabular}{|l|l|}
\hline
\textbf{部分} & \textbf{说明} \\
\hline
两个底面 & 圆形，上下各一个，一样大 \\
\hline
一个侧面 & 弯曲的面，把两个圆连在一起 \\
\hline
底面半径 & $r$ \\
\hline
高 & $h$ \\
\hline
\end{tabular}
\end{center}

\subsection{圆柱体的表面积}

\textbf{两个底面：}

每个底面是圆，面积是 $\pi r^2$。两个加起来：$2\pi r^2$。

\textbf{侧面：}

把圆柱体的外皮剥开，摊平，变成一个长方形！

\begin{center}
\begin{tabular}{|l|l|}
\hline
\textbf{这个长方形的...} & \textbf{等于...} \\
\hline
长 & 圆的周长 $= 2\pi r$ \\
\hline
宽 & 圆柱体的高 $= h$ \\
\hline
\end{tabular}
\end{center}

侧面积：$2\pi r \times h = 2\pi rh$

\textbf{总表面积：}

\[\text{表面积} = 2\pi r^2 + 2\pi rh = 2\pi r(r + h)\]

\begin{quote}
\textbf{语言记忆法}：

“两个圆形底面，加上侧面摊开的长方形。侧面长是圆周$2\pi r$，宽是高$h$。三块加起来。”
\end{quote}

\subsection{圆柱体的体积}

\[\text{体积} = \pi r^2 \times h = \pi r^2 h\]

想象把圆柱体切成很多很多层薄薄的圆，每层的面积是 $\pi r^2$，叠了 $h$ 层。

\begin{quote}
\textbf{语言记忆法}：

“底面是圆，面积是$\pi r^2$。叠了$h$层，乘以高就行。”
\end{quote}

\subsection{例题}

\textbf{题目：} “一个圆柱体，底面半径 3 厘米，高 5 厘米。表面积和体积各是多少？（$\pi$ 取 $3.14$）”

\textbf{表面积：}

\begin{center}
\begin{tabular}{|l|l|}
\hline
\textbf{部分} & \textbf{计算} \\
\hline
两个底面 & $2 \times 3.14 \times 3^2 = 56.52$ \\
\hline
侧面 & $2 \times 3.14 \times 3 \times 5 = 94.2$ \\
\hline
\textbf{总表面积} & $56.52 + 94.2 = 150.72$ 平方厘米 \\
\hline
\end{tabular}
\end{center}

\textbf{体积：}
\[3.14 \times 3^2 \times 5 = 3.14 \times 9 \times 5 = 141.3 \text{ 立方厘米}\]

\section{第四个：圆锥}

圆锥，就是圆柱体把顶面“捏”成一个点。

冰淇淋甜筒，金字塔，漏斗，都是圆锥（或者接近圆锥）。

\subsection{圆锥有什么？}

\begin{center}
\begin{tabular}{|l|l|}
\hline
\textbf{部分} & \textbf{说明} \\
\hline
一个底面 & 圆形 \\
\hline
一个侧面 & 弯曲的 \\
\hline
一个顶点 & 尖尖的那个点 \\
\hline
底面半径 & $r$ \\
\hline
高 & $h$ \\
\hline
斜面长（母线） & $l$ \\
\hline
\end{tabular}
\end{center}

$l$、$r$、$h$ 之间有关系：

\[l^2 = r^2 + h^2\]

（又见到勾股定理了！因为 $l$、$r$、$h$ 构成一个直角三角形。）

\subsection{圆锥的表面积}

\begin{center}
\begin{tabular}{|c|c|}
\hline
\textbf{部分} & \textbf{公式} \\
\hline
底面 & $\pi r^2$ \\
\hline
侧面（摊开是扇形） & $\pi rl$ \\
\hline
\textbf{总表面积} & $\pi r^2 + \pi rl = \pi r(r + l)$ \\
\hline
\end{tabular}
\end{center}

\begin{quote}
\textbf{语言记忆法}：

“一个圆形底面，加上侧面摊开的扇形。扇形面积是$\pi rl$，加上底面$\pi r^2$，提取$\pi r$合并。”
\end{quote}

\subsection{圆锥的体积}

\[\text{体积} = \frac{1}{3} \pi r^2 h\]

\textbf{为什么是三分之一？}

同样底面积和高的圆锥，体积恰好是圆柱体的三分之一。

你可以做个实验：拿一个圆锥形的容器，装满水，倒进同样底面积和高的圆柱形容器。

\textbf{倒三次，正好装满。}

\begin{quote}
\textbf{语言记忆法}：

“和圆柱体一样底面、一样高，体积只有三分之一。倒三次水才能装满圆柱体。”
\end{quote}

\subsection{例题}

\textbf{题目：} “一个圆锥，底面半径 3 厘米，高 4 厘米。体积是多少？（$\pi$ 取 $3.14$）”

先找 $l$：
\[l^2 = 3^2 + 4^2 = 9 + 16 = 25\]
\[l = 5 \text{ 厘米}\]

（又是 3、4、5！）

\textbf{表面积：}
\[3.14 \times 3 \times (3 + 5) = 3.14 \times 3 \times 8 = 75.36 \text{ 平方厘米}\]

\textbf{体积：}
\[\frac{1}{3} \times 3.14 \times 3^2 \times 4 = \frac{1}{3} \times 3.14 \times 9 \times 4 = \frac{1}{3} \times 113.04 = 37.68 \text{ 立方厘米}\]

\section{第五个：球}

球，是所有立体图形里最特别的。

没有边，没有角，没有面，只有一个完美的曲面，从任何方向看都一样。

足球，地球，泡泡，都是球。

\subsection{球的本质}

球只有一个尺寸：\textbf{半径 $r$}

从球心到球面上任意一点的距离，都是 $r$。

\textbf{球的本质}：离球心距离相同的所有点，连在一起，就是球。

\subsection{球的表面积}

\[\text{表面积} = 4\pi r^2\]

有个很妙的发现：球的表面积，恰好等于\textbf{四个}和它半径相同的圆的面积加起来。

\begin{quote}
\textbf{语言记忆法}：

“四个和球半径一样的圆，面积加起来刚好等于球的表面积。”
\end{quote}

\subsection{球的体积}

\[\text{体积} = \frac{4}{3}\pi r^3\]

这个公式怎么来的，涉及到微积分，我们现在还没学到那里。

先记住它，等你学了微积分，我们再回来把它推导出来。

（先欠着。微积分会来的。）

\begin{quote}
\textbf{语言记忆法}：

“$\pi r^3$ 是基础，乘以$\frac{4}{3}$。这个$\frac{4}{3}$怎么来的？等微积分来了再说，先记住它。”
\end{quote}

\subsection{例题}

\textbf{题目：} “一个球，半径 6 厘米。表面积和体积各是多少？（$\pi$ 取 $3.14$）”

\textbf{表面积：}
\[4 \times 3.14 \times 6^2 = 4 \times 3.14 \times 36 = 452.16 \text{ 平方厘米}\]

\textbf{体积：}
\[\frac{4}{3} \times 3.14 \times 6^3 = \frac{4}{3} \times 3.14 \times 216 = \frac{4}{3} \times 678.24 = 904.32 \text{ 立方厘米}\]

\section{总结一下}

\begin{center}
\begin{tabular}{|c|c|c|}
\hline
\textbf{图形} & \textbf{表面积} & \textbf{体积} \\
\hline
长方体 & $2(ab + bc + ac)$ & $abc$ \\
\hline
正方体 & $6a^2$ & $a^3$ \\
\hline
圆柱体 & $2\pi r(r + h)$ & $\pi r^2 h$ \\
\hline
圆锥 & $\pi r(r + l)$ & $\dfrac{1}{3}\pi r^2 h$ \\
\hline
球 & $4\pi r^2$ & $\dfrac{4}{3}\pi r^3$ \\
\hline
\end{tabular}
\end{center}

\section{现在你来}

\begin{quote}
\textbf{第一道：} “一个长方体，长 6 厘米，宽 4 厘米，高 3 厘米。表面积和体积各是多少？”
\end{quote}

\begin{quote}
\textbf{第二道：} “一个正方体，边长 5 厘米。表面积和体积各是多少？”
\end{quote}

\begin{quote}
\textbf{第三道：} “一个圆柱体，底面半径 4 厘米，高 6 厘米。表面积和体积各是多少？（$\pi$ 取 $3.14$）”
\end{quote}

\begin{quote}
\textbf{第四道：} “一个圆锥，底面半径 6 厘米，高 8 厘米。先找出斜面长 $l$，再算表面积和体积。（$\pi$ 取 $3.14$）”
\end{quote}

\begin{quote}
\textbf{第五道：} “一个球，半径 4 厘米。表面积和体积各是多少？（$\pi$ 取 $3.14$）”
\end{quote}

写下来，验证一下。

\section{本章小结}

从平面到立体，从二维到三维，数学的工具，一直都是一样的：

\textbf{理解本质，不要死记公式。}

\begin{center}
\begin{tabular}{|l|l|}
\hline
\textbf{图形} & \textbf{公式背后的故事} \\
\hline
长方体的体积 & 数格子数出来的 \\
\hline
圆柱体的侧面 & 摊开来的长方形 \\
\hline
圆锥的体积 & 倒水倒出来的三分之一 \\
\hline
球的表面积 & 四个圆拼出来的 \\
\hline
\end{tabular}
\end{center}

\textbf{记住故事，公式就忘不了。}

\subsection*{本章知识地图}

\begin{verbatim}
立体图形
│
├─ 基础概念
│   ├─ 表面积（外面有多大）
│   └─ 体积（里面能装多少）
│
├─ 长方体
│   ├─ 表面积 = 2(ab + bc + ac)
│   └─ 体积 = abc
│
├─ 正方体
│   ├─ 表面积 = 6a²
│   └─ 体积 = a³（所以叫“立方”）
│
├─ 圆柱体
│   ├─ 表面积 = 2πr(r + h)
│   └─ 体积 = πr²h
│
├─ 圆锥
│   ├─ 表面积 = πr(r + l)
│   ├─ 体积 = (1/3)πr²h
│   └─ l² = r² + h²（勾股定理）
│
└─ 球
    ├─ 表面积 = 4πr²
    └─ 体积 = (4/3)πr³
\end{verbatim}

\chapter{根号——到底是个啥？}

\section{先从你认识的东西开始}

你已经认识次方了。

\[3^2 = 9\]

意思是：3 乘以 3，等于 9。

\[4^2 = 16\]

意思是：4 乘以 4，等于 16。

现在我问你：

\textbf{谁乘以谁，等于 9？}

你说：3。

\textbf{谁乘以谁，等于 16？}

你说：4。

恭喜你，你已经在用根号的思维了。

只是还不知道它叫什么。

\section{根号就是反求平方}

根号，写作 $\sqrt{\phantom{0}}$，就是在问：

\textbf{谁的平方等于这个数？}

\begin{center}
\begin{tabular}{|c|c|l|}
\hline
\textbf{根号} & \textbf{等于} & \textbf{因为} \\
\hline
$\sqrt{9}$ & 3 & $3^2 = 9$ \\
\hline
$\sqrt{16}$ & 4 & $4^2 = 16$ \\
\hline
$\sqrt{25}$ & 5 & $5^2 = 25$ \\
\hline
\end{tabular}
\end{center}

\begin{quote}
\textbf{记忆法}：

“根号，就是平方的反求。看到$\sqrt{\text{某个数}}$，就问自己：谁的平方等于它？”
\end{quote}

\section{根号和反求法是一家人}

还记得之前学的反求法吗？

\[x^2 = 9\]

你用反求法找到 $x = 3$。

现在换个写法：

\[x = \sqrt{9} = 3\]

是同一件事。

$\sqrt{9}$，就是“谁的平方等于 9”的答案。

\begin{quote}
\textbf{根号，就是反求法的符号写法。}
\end{quote}

\section{但有时候，根号算不出整数}

\[\sqrt{2} = ?\]

谁的平方等于 2？

\begin{center}
\begin{tabular}{|c|c|c|}
\hline
\textbf{试试看} & \textbf{平方结果} & \textbf{比2大还是小} \\
\hline
1 & $1^2 = 1$ & 太小 \\
\hline
2 & $2^2 = 4$ & 太大 \\
\hline
\end{tabular}
\end{center}

中间有个数，它的平方恰好等于 2。

这个数，大约是 $1.41421356...$

小数点后面，无穷无尽，永远不重复，永远不结束。

（听起来很像 $\pi$ 对不对？它们是同类，都叫\textbf{无理数}。）

所以，我们不算它，就写 $\sqrt{2}$。

$\sqrt{2}$ 就是它的精确值。比写 $1.414...$ 更干净，更准确。

这就是为什么数学里，根号经常不化简，直接留着。

\begin{quote}
\textbf{根号，是数学家写“说不清楚的数”的方式。}
\end{quote}

\section{哪些根号能算出整数？}

这些数的平方根，恰好是整数：

\begin{center}
\begin{tabular}{|c|c|l|}
\hline
\textbf{根号} & \textbf{结果} & \textbf{因为} \\
\hline
$\sqrt{1}$ & 1 & $1^2 = 1$ \\
\hline
$\sqrt{4}$ & 2 & $2^2 = 4$ \\
\hline
$\sqrt{9}$ & 3 & $3^2 = 9$ \\
\hline
$\sqrt{16}$ & 4 & $4^2 = 16$ \\
\hline
$\sqrt{25}$ & 5 & $5^2 = 25$ \\
\hline
$\sqrt{36}$ & 6 & $6^2 = 36$ \\
\hline
$\sqrt{49}$ & 7 & $7^2 = 49$ \\
\hline
$\sqrt{64}$ & 8 & $8^2 = 64$ \\
\hline
$\sqrt{81}$ & 9 & $9^2 = 81$ \\
\hline
$\sqrt{100}$ & 10 & $10^2 = 100$ \\
\hline
\end{tabular}
\end{center}

这些叫\textbf{完全平方数}。

记住它们，以后遇到这些根号，直接认出答案，不用算。

\section{根号的运算}

根号，也可以加减乘除。但有些规则，要弄清楚。

\subsection{根号的乘法}

\[\sqrt{4} \times \sqrt{9} = \sqrt{4 \times 9} = \sqrt{36} = 6\]

\textbf{验证}：$\sqrt{4} = 2$，$\sqrt{9} = 3$，$2 \times 3 = 6$。没错！

\textbf{规律}：

\[\sqrt{a} \times \sqrt{b} = \sqrt{a \times b}\]

\begin{quote}
\textbf{记忆法}：

“两个根号相乘，就是把里面的数乘起来，再开根号。”
\end{quote}

\subsection{根号的除法}

\[\sqrt{36} \div \sqrt{4} = \sqrt{36 \div 4} = \sqrt{9} = 3\]

\textbf{验证}：$\sqrt{36} = 6$，$\sqrt{4} = 2$，$6 \div 2 = 3$。没错！

\textbf{规律}：

\[\sqrt{a} \div \sqrt{b} = \sqrt{a \div b}\]

\begin{quote}
\textbf{记忆法}：

“两个根号相除，就是把里面的数相除，再开根号。”
\end{quote}

\subsection{根号的化简}

有时候，根号里面的数太大，但可以化简：

\[\sqrt{12} = ?\]

12 能不能拆成一个完全平方数乘以另一个数？

\begin{center}
\begin{tabular}{|l|l|}
\hline
\textbf{拆法} & \textbf{说明} \\
\hline
$12 = 4 \times 3$ & 4 是完全平方数（$2^2 = 4$） \\
\hline
\end{tabular}
\end{center}

所以：

\[\sqrt{12} = \sqrt{4 \times 3} = \sqrt{4} \times \sqrt{3} = 2\sqrt{3}\]

\textbf{为什么要化简？}

因为 $2\sqrt{3}$ 比 $\sqrt{12}$ 更简洁，看起来更清楚。

\textbf{再来一个}：

\[\sqrt{18} = \sqrt{9 \times 2} = \sqrt{9} \times \sqrt{2} = 3\sqrt{2}\]

\begin{quote}
\textbf{记忆法}：

“根号化简，就是把里面的数，拆出完全平方数，放到根号外面。”
\end{quote}

\subsection{根号的加减法（最容易出错！）}

\[\sqrt{4} + \sqrt{9} = 2 + 3 = 5\]

但是：

\[\sqrt{4 + 9} = \sqrt{13}\]

$\sqrt{13}$ 不等于 5！

\textbf{根号里面的数，不能直接加减！}

\[\sqrt{a} + \sqrt{b} \neq \sqrt{a + b}\]

这是最容易出错的地方。

\textbf{为什么不能直接加？}

想想根号的本质。

\begin{center}
\begin{tabular}{|l|l|c|}
\hline
\textbf{表达式} & \textbf{问的是什么} & \textbf{答案} \\
\hline
$\sqrt{4}$ & 谁的平方等于 4 & 2 \\
\hline
$\sqrt{9}$ & 谁的平方等于 9 & 3 \\
\hline
$\sqrt{4} + \sqrt{9}$ & 两个答案加起来 & 5 \\
\hline
$\sqrt{4 + 9} = \sqrt{13}$ & 谁的平方等于 13 & 另一个问题！ \\
\hline
\end{tabular}
\end{center}

就像你问：
\begin{itemize}
\item “小明有 4 块糖，小红有 9 块糖，他们各自糖的平方根加起来是多少？”
\item “小明小红一共有 13 块糖，13 块糖的平方根是多少？”
\end{itemize}

是两道完全不同的题。

\textbf{但是，同类根号可以加减：}

\[2\sqrt{3} + 3\sqrt{3} = 5\sqrt{3}\]

就像：
\begin{itemize}
\item 2 个苹果加 3 个苹果，等于 5 个苹果
\item 2 个 $\sqrt{3}$ 加 3 个 $\sqrt{3}$，等于 5 个 $\sqrt{3}$
\end{itemize}

\begin{quote}
\textbf{记忆法}：

“根号加减，只有根号里面一样的，才能直接加减系数。根号里面不一样的，不能合并。”
\end{quote}

\section{总结一下}

\begin{center}
\begin{tabular}{|c|l|l|}
\hline
\textbf{运算} & \textbf{规则} & \textbf{例子} \\
\hline
乘法 & $\sqrt{a} \times \sqrt{b} = \sqrt{ab}$ & $\sqrt{4} \times \sqrt{9} = \sqrt{36} = 6$ \\
\hline
除法 & $\sqrt{a} \div \sqrt{b} = \sqrt{a \div b}$ & $\sqrt{36} \div \sqrt{4} = \sqrt{9} = 3$ \\
\hline
化简 & 拆出完全平方数 & $\sqrt{12} = 2\sqrt{3}$ \\
\hline
加减 & 只有同类根号才能合并 & $2\sqrt{3} + 3\sqrt{3} = 5\sqrt{3}$ \\
\hline
$\triangle$ & $\sqrt{a} + \sqrt{b} \neq \sqrt{a+b}$ & $\sqrt{4} + \sqrt{9} \neq \sqrt{13}$ \\
\hline
\end{tabular}
\end{center}

\section{现在你来}

\begin{quote}
\textbf{第一题}：$\sqrt{49} = $ \underline{\hspace{2cm}}

提示：谁的平方等于 49？
\end{quote}

\textbf{答案：} \textbf{7}（因为 $7^2 = 49$）

\begin{quote}
\textbf{第二题}：$\sqrt{9} \times \sqrt{16} = $ \underline{\hspace{2cm}}

提示：可以先各自算，也可以先乘再开根号
\end{quote}

\textbf{答案：} $3 \times 4 = 12$

或者：$\sqrt{9 \times 16} = \sqrt{144} = 12$

两种方法，答案一样！两条路走到同一个地方。

\begin{quote}
\textbf{第三题}：化简 $\sqrt{20}$

提示：20 能拆成哪个完全平方数乘以另一个数？
\end{quote}

\textbf{答案：} $20 = 4 \times 5$

$\sqrt{20} = \sqrt{4} \times \sqrt{5} = 2\sqrt{5}$

\begin{quote}
\textbf{第四题}：$3\sqrt{5} + 2\sqrt{5} = $ \underline{\hspace{2cm}}
\end{quote}

\textbf{答案：} $5\sqrt{5}$（同类根号，直接加系数）

\begin{quote}
\textbf{第五题}：$\sqrt{3} + \sqrt{5} = $ \underline{\hspace{2cm}}（能化简吗？）
\end{quote}

\textbf{答案：} \textbf{不能化简，就是 $\sqrt{3} + \sqrt{5}$}

根号里面不一样，没办法合并。

这道题的答案，就是题目本身。这不是你做错了，是真的没法再简化了。

数学有时候就是这样，答案就是“没有更简单的答案”。接受它。

\section{本章小结}

根号，不是什么神秘的东西。

它就是：\textbf{平方的反求。}

\begin{quote}
\textbf{“谁的平方等于它？”}

就这一句话。
\end{quote}

\begin{center}
\begin{tabular}{|l|l|}
\hline
\textbf{情况} & \textbf{怎么办} \\
\hline
能算出整数 & 直接算出来 \\
\hline
算不出整数 & 就留着根号，它就是精确值 \\
\hline
乘除 & 根号里面可以乘除 \\
\hline
化简 & 拆出完全平方数 \\
\hline
加减 & 小心！不能随便合并 \\
\hline
\end{tabular}
\end{center}

\subsection*{本章知识地图}

\begin{verbatim}
根号
│
├─ 本质
│   └─ 平方的反求：“谁的平方等于它？”
│
├─ 完全平方数
│   └─ 1, 4, 9, 16, 25, 36, 49, 64, 81, 100...
│
├─ 运算规则
│   ├─ 乘法：√a × √b = √(ab)
│   ├─ 除法：√a ÷ √b = √(a÷b)
│   ├─ 化简：拆出完全平方数
│   └─ 加减：只有同类根号能合并
│
├─ 常见错误
│   └─ √a + √b ≠ √(a+b)（千万别犯！）
│
└─ 无理数
    └─ 算不出整数的根号（如√2, √3）
       小数无穷无尽，就留着根号
\end{verbatim}

\chapter{特殊三角形}

\section{先回顾一下你学过的东西}

上一章，我们认识了三角形里的四条线——高、中线、角平分线、中位线。

你已经知道：
\begin{itemize}
\item 三角形的三个角加起来永远是 \textbf{$180^\circ$}
\item 三角形的面积 $=$ 底 $\times$ 高 $\div$ 2
\item 勾股定理：直角三角形里，$a^2 + b^2 = c^2$
\end{itemize}

这些工具，你都有了。

现在我问你一个问题：

\begin{quote}
如果一个三角形的三条边都一样长，那它的角是多少度？
\end{quote}

你可能会想：三个角一样大，加起来 $180^\circ$，那每个角就是 $180^\circ \div 3 = 60^\circ$。

没错！

你刚才用已知的规律，推出了一个新结论。

这一章，我们就是要认识这些\textbf{有固定规律的三角形}。

它们特殊到什么程度呢？

知道一条边，就能算出所有边。
知道一个角，就能算出所有角。

几乎不需要额外信息。

\section{第一个：等边三角形}

\subsection{它是什么？}

三条边，全部一样长。

就这一句话。

\subsection{角是多少？}

你刚才已经推出来了：

三个角加起来是 $180^\circ$，三个角一样大，平均分：

\[180^\circ \div 3 = 60^\circ\]

\textbf{等边三角形的每个角，永远是 $60^\circ$。}

不管边长是 1 厘米还是 100 厘米，不管大小，$60^\circ$。永远。

\begin{quote}
\textbf{记忆法}：

“三边相等，三角相等，每个角永远是$60^\circ$。”
\end{quote}

\subsection{等边三角形的高，怎么来的？}

从顶点画一条高，把等边三角形切成两个一样的直角三角形。

现在看其中一个直角三角形：

\begin{center}
\begin{tabular}{|c|c|l|}
\hline
\textbf{边} & \textbf{长度} & \textbf{怎么来的} \\
\hline
斜边 & $a$ & 等边三角形的边长 \\
\hline
底边 & $\dfrac{a}{2}$ & 底边被高切成两半 \\
\hline
高 & $h$ & 我们要找的 \\
\hline
\end{tabular}
\end{center}

我们已经知道两条边，要找第三条边。

\textbf{这不就是勾股定理吗？}

\[a^2 = h^2 + \left(\frac{a}{2}\right)^2\]

把 $\left(\dfrac{a}{2}\right)^2$ 移到左边：

\[h^2 = a^2 - \left(\frac{a}{2}\right)^2\]

$\left(\dfrac{a}{2}\right)^2$ 是什么？就是 $\dfrac{a}{2}$ 乘以 $\dfrac{a}{2}$：

\[\frac{a}{2} \times \frac{a}{2} = \frac{a^2}{4}\]

所以：

\[h^2 = a^2 - \frac{a^2}{4}\]

$a^2$ 可以写成 $\dfrac{4a^2}{4}$（分母统一）：

\[h^2 = \frac{4a^2}{4} - \frac{a^2}{4} = \frac{3a^2}{4}\]

两边开根号：

\[h = \sqrt{\frac{3a^2}{4}} = \frac{\sqrt{3a^2}}{\sqrt{4}} = \frac{\sqrt{3} \times a}{2} = \frac{\sqrt{3}}{2}a\]

所以：

\[h = \frac{\sqrt{3}}{2}a\]

这个 $\dfrac{\sqrt{3}}{2}$，不是天上掉下来的，是用勾股定理一步一步推出来的。

\begin{quote}
\textbf{记忆法}：

“等边三角形的高，是从勾股定理推出来的。斜边是边长$a$，底边是$\frac{a}{2}$，高用勾股定理找。”
\end{quote}

\subsection{等边三角形的面积}

\[\text{面积} = \frac{1}{2} \times \text{底} \times \text{高} = \frac{1}{2} \times a \times \frac{\sqrt{3}}{2}a = \frac{\sqrt{3}}{4}a^2\]

\begin{quote}
\textbf{记忆法}：

“底乘高除以2，把高的公式代进去，就得到$\frac{\sqrt{3}}{4}a^2$。”
\end{quote}

\subsection{例题}

\textbf{题目：} 等边三角形，边长 6 厘米。求每个角、高、周长、面积。

\begin{center}
\begin{tabular}{|c|l|c|}
\hline
\textbf{项目} & \textbf{计算过程} & \textbf{答案} \\
\hline
每个角 & 永远是$60^\circ$ & \textbf{$60^\circ$} \\
\hline
高 & $h^2 = 6^2 - 3^2 = 36 - 9 = 27$，$h = \sqrt{27} = 3\sqrt{3}$ & \textbf{$3\sqrt{3}$ 厘米} \\
\hline
周长 & $6 \times 3$ & \textbf{18 厘米} \\
\hline
面积 & $\frac{\sqrt{3}}{4} \times 36 = 9\sqrt{3}$ & \textbf{$9\sqrt{3}$ 平方厘米} \\
\hline
\end{tabular}
\end{center}

你看，高的推导，就是勾股定理在等边三角形里的应用。

以后遇到等边三角形，不需要死记公式，直接用勾股定理推一遍就行。

\textbf{会推导，比会背公式强一百倍。}

\section{第二个：等腰三角形}

\subsection{它是什么？}

\begin{center}
\begin{tabular}{|c|l|}
\hline
\textbf{部分} & \textbf{说明} \\
\hline
两条\textbf{腰} & 一样长 \\
\hline
一条\textbf{底边} & 不一定和腰一样 \\
\hline
\end{tabular}
\end{center}

等边三角形其实是等腰三角形的特殊情况——三条腰都一样长而已。

\subsection{角是多少？}

两条腰一样长，对应两个底角也一样大。

\[2 \times \text{底角} + \text{顶角} = 180^\circ\]

\begin{center}
\begin{tabular}{|l|l|}
\hline
\textbf{已知} & \textbf{求法} \\
\hline
知道顶角，找底角 & $\text{底角} = \dfrac{180^\circ - \text{顶角}}{2}$ \\
\hline
知道底角，找顶角 & $\text{顶角} = 180^\circ - 2 \times \text{底角}$ \\
\hline
\end{tabular}
\end{center}

\begin{quote}
\textbf{记忆法}：

“两腰相等，两底角相等。知道一个角，用$180^\circ$减掉，剩下的平分。”
\end{quote}

\subsection{等腰三角形的高，怎么来的？}

从顶角画一条高，把等腰三角形切成两个一样的直角三角形。

\begin{center}
\begin{tabular}{|c|c|l|}
\hline
\textbf{边} & \textbf{长度} & \textbf{怎么来的} \\
\hline
斜边 & $a$（腰长） & 等腰三角形的腰 \\
\hline
底边 & $\dfrac{b}{2}$ & 底边被高切成两半 \\
\hline
高 & $h$ & 我们要找的 \\
\hline
\end{tabular}
\end{center}

还是勾股定理：

\[h^2 = a^2 - \left(\frac{b}{2}\right)^2\]

\[h = \sqrt{a^2 - \left(\frac{b}{2}\right)^2}\]

\begin{quote}
\textbf{记忆法}：

“从顶角画高，切出直角三角形。斜边是腰，底边是半个底边，高用勾股定理找。”
\end{quote}

\subsection{例题}

\textbf{题目一：} 等腰三角形，腰长 5 厘米，底边 6 厘米。

\begin{center}
\begin{tabular}{|c|l|c|}
\hline
\textbf{项目} & \textbf{计算过程} & \textbf{答案} \\
\hline
斜边 & 腰长 & 5 \\
\hline
底边 & $6 \div 2$ & 3 \\
\hline
高 & $\sqrt{5^2 - 3^2} = \sqrt{25 - 9} = \sqrt{16}$ & \textbf{4 厘米} \\
\hline
面积 & $\frac{1}{2} \times 6 \times 4$ & \textbf{12 平方厘米} \\
\hline
周长 & $5 + 5 + 6$ & \textbf{16 厘米} \\
\hline
\end{tabular}
\end{center}

\textbf{题目二：} 等腰三角形，顶角 $40^\circ$。求底角。

\[\text{底角} = \frac{180^\circ - 40^\circ}{2} = \frac{140^\circ}{2} = 70^\circ\]

就这么简单。

\section{第三个：$30^\circ$-$60^\circ$-$90^\circ$ 三角形}

\subsection{它是什么？}

三个角分别是 $30^\circ$、$60^\circ$、$90^\circ$。

\textbf{它从哪里来的？}

把等边三角形从顶角画一条高，切出来的就是 $30^\circ$-$60^\circ$-$90^\circ$ 三角形！

\begin{center}
\begin{tabular}{|c|l|}
\hline
\textbf{角} & \textbf{怎么来的} \\
\hline
$30^\circ$ & 顶角被切成两半：$60^\circ \div 2 = 30^\circ$ \\
\hline
$60^\circ$ & 底角还是 $60^\circ$ \\
\hline
$90^\circ$ & 高和底边垂直 \\
\hline
\end{tabular}
\end{center}

\subsection{三条边的比例，怎么来的？}

设等边三角形的边长是 2（用 2 方便计算）。

切出来的直角三角形：

\begin{center}
\begin{tabular}{|c|c|l|}
\hline
\textbf{边} & \textbf{长度} & \textbf{怎么来的} \\
\hline
斜边 & 2 & 等边三角形的边长 \\
\hline
短直角边 & 1 & 底边被切成两半，$2 \div 2 = 1$ \\
\hline
长直角边 & $h$ & 高，用勾股定理找 \\
\hline
\end{tabular}
\end{center}

用勾股定理：

\[h^2 = 2^2 - 1^2 = 4 - 1 = 3\]
\[h = \sqrt{3}\]

所以三条边是：

\[1 : \sqrt{3} : 2\]

\textbf{短边 : 长边 : 斜边}

如果短边是 $a$，那：

\begin{center}
\begin{tabular}{|c|c|c|}
\hline
\textbf{边} & \textbf{长度} & \textbf{对着哪个角} \\
\hline
短直角边 & $a$ & $30^\circ$ \\
\hline
长直角边 & $\sqrt{3}a$ & $60^\circ$ \\
\hline
斜边 & $2a$ & $90^\circ$ \\
\hline
\end{tabular}
\end{center}

\begin{quote}
\textbf{记忆法}：

“$30^\circ$-$60^\circ$-$90^\circ$，从等边三角形切出来。短边对$30^\circ$，斜边是短边两倍，长边用勾股定理推。比例是$1:\sqrt{3}:2$。”
\end{quote}

\subsection{例题}

\textbf{题目一：} $30^\circ$-$60^\circ$-$90^\circ$ 三角形，短直角边是 4 厘米。

\begin{center}
\begin{tabular}{|c|l|c|}
\hline
\textbf{项目} & \textbf{用比例推导} & \textbf{答案} \\
\hline
$a$ & 短直角边 $= 4$ & 4 \\
\hline
长直角边 & $\sqrt{3} \times 4$ & \textbf{$4\sqrt{3}$ 厘米} \\
\hline
斜边 & $2 \times 4$ & \textbf{8 厘米} \\
\hline
面积 & $\frac{1}{2} \times 4 \times 4\sqrt{3}$ & \textbf{$8\sqrt{3}$ 平方厘米} \\
\hline
\end{tabular}
\end{center}

\textbf{题目二：} $30^\circ$-$60^\circ$-$90^\circ$ 三角形，斜边是 10 厘米。

斜边 $= 2a = 10$，所以 $a = 5$

\begin{center}
\begin{tabular}{|c|c|}
\hline
\textbf{边} & \textbf{答案} \\
\hline
短直角边 & \textbf{5 厘米} \\
\hline
长直角边 & \textbf{$5\sqrt{3}$ 厘米} \\
\hline
\end{tabular}
\end{center}

\section{第四个：$45^\circ$-$45^\circ$-$90^\circ$ 三角形}

\subsection{它是什么？}

两条直角边一样长，两个底角都是 $45^\circ$，一个直角 $90^\circ$。

\textbf{它从哪里来的？}

把正方形从对角线切一刀，切出来的就是 $45^\circ$-$45^\circ$-$90^\circ$ 三角形！

\subsection{三条边的比例，怎么来的？}

设两条直角边都是 $a$。

用勾股定理找斜边：

\[\text{斜边}^2 = a^2 + a^2 = 2a^2\]

两边开根号：

\[\text{斜边} = \sqrt{2a^2}\]

$\sqrt{2a^2}$ 是什么？用根号的乘法（还记得上一章吗？）：

\[\sqrt{2a^2} = \sqrt{2} \times \sqrt{a^2} = \sqrt{2} \times a = a\sqrt{2}\]

所以三条边的比例是：

\[1 : 1 : \sqrt{2}\]

\textbf{两条直角边 : 斜边}

\begin{quote}
\textbf{记忆法}：

“$45^\circ$-$45^\circ$-$90^\circ$，从正方形对角线切出来。两条直角边一样长，斜边用勾股定理推，是直角边乘以$\sqrt{2}$。”
\end{quote}

\subsection{例题}

\textbf{题目一：} $45^\circ$-$45^\circ$-$90^\circ$ 三角形，直角边是 5 厘米。

\begin{center}
\begin{tabular}{|c|l|c|}
\hline
\textbf{项目} & \textbf{用勾股定理推导} & \textbf{答案} \\
\hline
斜边$^2$ & $5^2 + 5^2 = 25 + 25 = 50$ & — \\
\hline
斜边 & $\sqrt{50} = \sqrt{25 \times 2} = 5\sqrt{2}$ & \textbf{$5\sqrt{2}$ 厘米} \\
\hline
面积 & $\frac{1}{2} \times 5 \times 5$ & \textbf{12.5 平方厘米} \\
\hline
周长 & $5 + 5 + 5\sqrt{2}$ & \textbf{$10 + 5\sqrt{2}$ 厘米} \\
\hline
\end{tabular}
\end{center}

\textbf{题目二：} $45^\circ$-$45^\circ$-$90^\circ$ 三角形，斜边是 $8\sqrt{2}$ 厘米。求直角边。

斜边 $= a\sqrt{2} = 8\sqrt{2}$

两边除以 $\sqrt{2}$：

\[a = 8\sqrt{2} \div \sqrt{2} = 8 \text{ 厘米}\]

\section{总结一下}

\begin{center}
\begin{tabular}{|c|c|c|l|}
\hline
\textbf{三角形} & \textbf{角度} & \textbf{边的比例} & \textbf{比例怎么来的} \\
\hline
等边三角形 & $60^\circ$-$60^\circ$-$60^\circ$ & $1:1:1$ & 三边相等 \\
\hline
等腰三角形 & 底角相等 & 两腰相等 & 两腰相等 \\
\hline
$30^\circ$-$60^\circ$-$90^\circ$ & $30^\circ$-$60^\circ$-$90^\circ$ & $1:\sqrt{3}:2$ & 等边三角形切一刀，勾股定理推 \\
\hline
$45^\circ$-$45^\circ$-$90^\circ$ & $45^\circ$-$45^\circ$-$90^\circ$ & $1:1:\sqrt{2}$ & 正方形切对角线，勾股定理推 \\
\hline
\end{tabular}
\end{center}

\section{现在你来}

\begin{quote}
\textbf{第一题：} 等边三角形，边长 8 厘米。求高和面积。

提示：用勾股定理推导高
\end{quote}

\begin{quote}
\textbf{第二题：} $30^\circ$-$60^\circ$-$90^\circ$ 三角形，长直角边是 $6\sqrt{3}$ 厘米。求短直角边和斜边。

提示：长直角边 $= \sqrt{3}a$，先找 $a$
\end{quote}

\begin{quote}
\textbf{第三题：} $45^\circ$-$45^\circ$-$90^\circ$ 三角形，斜边是 10 厘米。求直角边。

提示：斜边 $= a\sqrt{2}$
\end{quote}

写下来，验证一下。

\section{本章小结}

这四种三角形，是数学世界里最常见的四张面孔。

但更重要的是，\textbf{你知道了它们的公式从哪里来。}

\begin{center}
\begin{tabular}{|c|l|l|}
\hline
\textbf{比例} & \textbf{不是...} & \textbf{而是...} \\
\hline
$1:\sqrt{3}:2$ & 背出来的 & 从等边三角形里切出来，用勾股定理推的 \\
\hline
$1:1:\sqrt{2}$ & 背出来的 & 从正方形里切出来，用勾股定理推的 \\
\hline
\end{tabular}
\end{center}

以后遇到任何一种，就算忘了公式，你也能自己推一遍。

\textbf{这才是真正学会了数学。}

下一章，我们用这四个老朋友，去解开边和角的最终秘密。

\subsection*{本章知识地图}

\begin{verbatim}
特殊三角形
│
├─ 等边三角形
│   ├─ 三边相等
│   ├─ 每个角60°
│   └─ 高 = (√3/2)a（勾股定理推）
│
├─ 等腰三角形
│   ├─ 两腰相等
│   ├─ 两底角相等
│   └─ 高用勾股定理推
│
├─ 30°-60°-90° 三角形
│   ├─ 来源：等边三角形切一刀
│   ├─ 比例：1 : √3 : 2
│   └─ 短边对30°，斜边是短边两倍
│
└─ 45°-45°-90° 三角形
    ├─ 来源：正方形对角线切一刀
    ├─ 比例：1 : 1 : √2
    └─ 两直角边相等，斜边乘√2
\end{verbatim}

\chapter{$\sin$、$\cos$、$\tan$——求角度的三把钥匙}

\section{先回顾一下你已经会的}

你已经认识了很多老朋友：

\begin{center}
\begin{tabular}{|l|l|}
\hline
\textbf{你学过的} & \textbf{能做什么} \\
\hline
勾股定理 & 知道两条边，求第三条边 \\
\hline
特殊三角形 & $30^\circ$-$60^\circ$-$90^\circ$ 边的比例是 $1:\sqrt{3}:2$ \\
\hline
特殊三角形 & $45^\circ$-$45^\circ$-$90^\circ$ 边的比例是 $1:1:\sqrt{2}$ \\
\hline
\end{tabular}
\end{center}

但现在，假设你遇到这样的问题：

\begin{quote}
“这个坡有多陡？”

“这架梯子和地面成多少度？”

“这座山的仰角是多少？”
\end{quote}

这些问题，勾股定理回答不了。

\textbf{勾股定理只管边，不管角。}

那谁来管角？

\textbf{$\sin$、$\cos$、$\tan$。}

\section{它们是干嘛的？}

就是\textbf{求角度}的。

你量了两条边的长度，想知道某个角是多少度，就用它们。

或者反过来，你知道角度和一条边，想知道另一条边，也用它们。

\begin{quote}
\textbf{一句话}：

勾股定理管边，$\sin$、$\cos$、$\tan$ 管角。两套工具，各有分工。
\end{quote}

\section{怎么用？三步走}

你有一个直角三角形（还记得吗？就是有一个$90^\circ$角的三角形）。

\subsection{第一步：盯死你要求的那个角}

不是直角，是另外两个角里面，你感兴趣的那一个。

盯好了，别换。

\subsection{第二步：认清三条边的身份}

对于你盯着的那个角，三条边各有身份：

\begin{center}
\begin{tabular}{|c|l|l|}
\hline
\textbf{边的名字} & \textbf{说明} & \textbf{怎么找} \\
\hline
\textbf{对边} & 这个角正对面的那条边 & 从这个角往对面看，第一眼看到的 \\
\hline
\textbf{邻边} & 紧挨着这个角的那条边（不是斜边） & 从这个角出发，沿着地面走的那条 \\
\hline
\textbf{斜边} & 最长的那条，永远对着直角 & 永远是三条边里最长的 \\
\hline
\end{tabular}
\end{center}

\subsection{第三步：选公式}

\[\tan = \frac{\text{对边}}{\text{邻边}}\]

\[\sin = \frac{\text{对边}}{\text{斜边}}\]

\[\cos = \frac{\text{邻边}}{\text{斜边}}\]

\textbf{选哪个？看你知道哪两条边。}

\begin{center}
\begin{tabular}{|l|c|}
\hline
\textbf{你知道} & \textbf{用哪个} \\
\hline
对边和邻边 & $\tan$ \\
\hline
对边和斜边 & $\sin$ \\
\hline
邻边和斜边 & $\cos$ \\
\hline
\end{tabular}
\end{center}

\section{为什么有三个名字？}

因为三条边，两两相除，有三种组合：

\begin{center}
\begin{tabular}{|l|c|}
\hline
\textbf{除法} & \textbf{名字} \\
\hline
对边 $\div$ 邻边 & $\tan$ \\
\hline
对边 $\div$ 斜边 & $\sin$ \\
\hline
邻边 $\div$ 斜边 & $\cos$ \\
\hline
\end{tabular}
\end{center}

三种组合，三个名字。仅此而已。

\begin{quote}
\textbf{记忆法}：

“三条边两两相除，有三种除法，三个名字。知道哪两条边，就选对应的名字。”
\end{quote}

\section{第一个：$\tan$}

\subsection{$\tan$ 是什么？}

\[\tan A = \frac{\text{对边}}{\text{邻边}}\]

$A$ 是你盯着的那个角。$\tan A$ 的值，就是对边除以邻边的结果。

\begin{quote}
\textbf{记忆法}：

“$\tan$ 就是对边除以邻边。知道这两条边，就能求角度。”
\end{quote}

\subsection{$\tan$ 的感觉}

想象一个坡。

\begin{center}
\begin{tabular}{|l|c|c|c|c|}
\hline
\textbf{情况} & \textbf{邻边（走多远）} & \textbf{对边（爬多高）} & $\tan$ & \textbf{角度} \\
\hline
坡A & 10米 & 10米 & $\frac{10}{10} = 1$ & $45^\circ$ \\
\hline
坡B & 10米 & 5米 & $\frac{5}{10} = 0.5$ & 比$45^\circ$缓 \\
\hline
\end{tabular}
\end{center}

\textbf{$\tan$ 越大，角越大，坡越陡。}

\subsection{特殊角的 $\tan$ 值}

还记得上一章的特殊三角形吗？我们用它们来推 $\tan$ 的值。

\textbf{$30^\circ$-$60^\circ$-$90^\circ$ 三角形}，边的比例是 $1:\sqrt{3}:2$。

\begin{center}
\begin{tabular}{|c|c|c|c|}
\hline
\textbf{盯着哪个角} & \textbf{对边} & \textbf{邻边} & $\tan$ \\
\hline
$30^\circ$ & 1（短边） & $\sqrt{3}$（长边） & $\frac{1}{\sqrt{3}} = \frac{\sqrt{3}}{3}$ \\
\hline
$60^\circ$ & $\sqrt{3}$（长边） & 1（短边） & $\frac{\sqrt{3}}{1} = \sqrt{3}$ \\
\hline
\end{tabular}
\end{center}

\textbf{$45^\circ$-$45^\circ$-$90^\circ$ 三角形}，边的比例是 $1:1:\sqrt{2}$。

\begin{center}
\begin{tabular}{|c|c|c|c|}
\hline
\textbf{盯着哪个角} & \textbf{对边} & \textbf{邻边} & $\tan$ \\
\hline
$45^\circ$ & 1 & 1 & $\frac{1}{1} = 1$ \\
\hline
\end{tabular}
\end{center}

\textbf{总结}：

\begin{center}
\begin{tabular}{|c|c|c|l|}
\hline
\textbf{角度} & $\tan$ & \textbf{大约是} & \textbf{感受} \\
\hline
$30^\circ$ & $\frac{\sqrt{3}}{3}$ & 0.58 & 比较缓 \\
\hline
$45^\circ$ & $1$ & 1 & 不陡不缓，对边邻边一样长 \\
\hline
$60^\circ$ & $\sqrt{3}$ & 1.73 & 比较陡 \\
\hline
\end{tabular}
\end{center}

\begin{quote}
\textbf{记忆法}：

“角度越大，$\tan$越大，坡越陡。$\tan 45^\circ = 1$最好记，对边邻边一样长。”
\end{quote}

\subsection{应用题}

\textbf{题目一：求坡度角}

\begin{quote}
“一个斜坡，水平距离是 10 米，垂直高度是 5 米。这个斜坡的坡度角是多少度？”
\end{quote}

\textbf{第一步：盯死坡度角。}

就是斜坡和地面之间的那个角。

\textbf{第二步：认清三条边的身份。}

\begin{center}
\begin{tabular}{|c|l|c|}
\hline
\textbf{边} & \textbf{是什么} & \textbf{长度} \\
\hline
对边 & 垂直高度（在角对面） & 5米 \\
\hline
邻边 & 水平距离（紧挨着角） & 10米 \\
\hline
\end{tabular}
\end{center}

\textbf{第三步：选公式。}

知道对边和邻边，用 $\tan$。

\[\tan A = \frac{\text{对边}}{\text{邻边}} = \frac{5}{10} = 0.5\]

$0.5$ 比 $\tan 30^\circ \approx 0.58$ 还小，说明这个坡比 $30^\circ$ 还缓。

（精确角度需要用计算器，大约是 $26.6^\circ$）

\textbf{题目二：求树高}

\begin{quote}
“一棵树，你站在离树根 8 米的地方，仰望树顶，仰角是 $60^\circ$。树有多高？”
\end{quote}

\textbf{第一步：盯死仰角 $60^\circ$。}

\textbf{第二步：认清三条边的身份。}

\begin{center}
\begin{tabular}{|c|l|c|}
\hline
\textbf{边} & \textbf{是什么} & \textbf{长度} \\
\hline
对边 & 树高（在角对面） & $?$ \\
\hline
邻边 & 你到树根的距离（紧挨着角） & 8米 \\
\hline
\end{tabular}
\end{center}

\textbf{第三步：选公式。}

知道对边和邻边，用 $\tan$。

\[\tan 60^\circ = \frac{\text{树高}}{8}\]
\[\sqrt{3} = \frac{\text{树高}}{8}\]

两边乘以 8：

\[\text{树高} = 8\sqrt{3} \approx 13.86 \text{ 米}\]

\textbf{验证}：$\dfrac{8\sqrt{3}}{8} = \sqrt{3} = \tan 60^\circ$ \checkmark

\section{第二个：$\sin$}

\subsection{$\sin$ 是什么？}

\[\sin A = \frac{\text{对边}}{\text{斜边}}\]

$A$ 是你盯着的那个角。$\sin A$ 的值，就是对边除以斜边的结果。

\begin{quote}
\textbf{记忆法}：

“$\sin$ 就是对边除以斜边。知道这两条边，就能求角度。”
\end{quote}

\subsection{$\sin$ 的感觉}

$\sin A$ 量的是，对边占斜边的几分之几。

\begin{itemize}
\item 角度越大 $\to$ 对边越长 $\to$ $\sin$ 越大
\item 角度越小 $\to$ 对边越短 $\to$ $\sin$ 越小
\end{itemize}

\subsection{特殊角的 $\sin$ 值}

\textbf{$30^\circ$-$60^\circ$-$90^\circ$ 三角形}，边的比例是 $1:\sqrt{3}:2$。

\begin{center}
\begin{tabular}{|c|c|c|c|}
\hline
\textbf{盯着哪个角} & \textbf{对边} & \textbf{斜边} & $\sin$ \\
\hline
$30^\circ$ & 1 & 2 & $\frac{1}{2}$ \\
\hline
$60^\circ$ & $\sqrt{3}$ & 2 & $\frac{\sqrt{3}}{2}$ \\
\hline
\end{tabular}
\end{center}

\textbf{$45^\circ$-$45^\circ$-$90^\circ$ 三角形}，边的比例是 $1:1:\sqrt{2}$。

\begin{center}
\begin{tabular}{|c|c|c|c|}
\hline
\textbf{盯着哪个角} & \textbf{对边} & \textbf{斜边} & $\sin$ \\
\hline
$45^\circ$ & 1 & $\sqrt{2}$ & $\frac{1}{\sqrt{2}} = \frac{\sqrt{2}}{2}$ \\
\hline
\end{tabular}
\end{center}

\textbf{总结}：

\begin{center}
\begin{tabular}{|c|c|c|l|}
\hline
\textbf{角度} & $\sin$ & \textbf{大约是} & \textbf{感受} \\
\hline
$30^\circ$ & $\frac{1}{2}$ & 0.5 & 对边是斜边的一半 \\
\hline
$45^\circ$ & $\frac{\sqrt{2}}{2}$ & 0.71 & 对边是斜边的七成 \\
\hline
$60^\circ$ & $\frac{\sqrt{3}}{2}$ & 0.87 & 对边接近斜边长度 \\
\hline
\end{tabular}
\end{center}

\begin{quote}
\textbf{记忆法}：

“$\sin 30^\circ = \frac{1}{2}$最好记，对边刚好是斜边一半。角度越大，$\sin$越大。”
\end{quote}

\subsection{应用题}

\textbf{题目三：求梯子顶端高度}

\begin{quote}
“一架梯子，长 10 米，靠在墙上，梯子和地面的夹角是 $30^\circ$。梯子顶端距离地面多高？”
\end{quote}

\textbf{第一步：盯死夹角 $30^\circ$。}

\textbf{第二步：认清三条边的身份。}

\begin{center}
\begin{tabular}{|c|l|c|}
\hline
\textbf{边} & \textbf{是什么} & \textbf{长度} \\
\hline
斜边 & 梯子（斜的那条） & 10米 \\
\hline
对边 & 梯子顶端的高度（在角对面） & $?$ \\
\hline
\end{tabular}
\end{center}

\textbf{第三步：选公式。}

知道对边和斜边，用 $\sin$。

\[\sin 30^\circ = \frac{\text{高度}}{10}\]
\[\frac{1}{2} = \frac{\text{高度}}{10}\]
\[\text{高度} = 10 \times \frac{1}{2} = 5 \text{ 米}\]

\textbf{验证}：$\dfrac{5}{10} = \dfrac{1}{2} = \sin 30^\circ$ \checkmark

\textbf{题目四：求仰角}

\begin{quote}
“一架飞机，飞行高度是 3000 米，飞机和地面观测点的连线是 6000 米。飞机和地面的仰角是多少度？”
\end{quote}

\textbf{第一步：盯死仰角。}

\textbf{第二步：认清三条边的身份。}

\begin{center}
\begin{tabular}{|c|l|c|}
\hline
\textbf{边} & \textbf{是什么} & \textbf{长度} \\
\hline
对边 & 飞行高度（在角对面） & 3000米 \\
\hline
斜边 & 观测点到飞机的连线（最长） & 6000米 \\
\hline
\end{tabular}
\end{center}

\textbf{第三步：选公式。}

知道对边和斜边，用 $\sin$。

\[\sin A = \frac{3000}{6000} = \frac{1}{2}\]

查表，$\sin 30^\circ = \dfrac{1}{2}$。

所以仰角是 \textbf{$30^\circ$}。

\section{第三个：$\cos$}

\subsection{$\cos$ 是什么？}

\[\cos A = \frac{\text{邻边}}{\text{斜边}}\]

$A$ 是你盯着的那个角。$\cos A$ 的值，就是邻边除以斜边的结果。

\begin{quote}
\textbf{记忆法}：

“$\cos$ 就是邻边除以斜边。知道这两条边，就能求角度。”
\end{quote}

\subsection{$\sin$ 和 $\cos$ 的关系}

你有没有发现：

\[\sin 30^\circ = \frac{1}{2}\]
\[\cos 60^\circ = \frac{1}{2}\]

它们一样！为什么？

在 $30^\circ$-$60^\circ$-$90^\circ$ 三角形里：
\begin{itemize}
\item 盯着 $30^\circ$ 看，对边是短边
\item 盯着 $60^\circ$ 看，邻边是短边
\end{itemize}

\textbf{同一条边，对一个角是对边，对另一个角是邻边。}

所以：

\[\sin A = \cos(90^\circ - A)\]

\begin{quote}
\textbf{记忆法}：

“$\sin$ 和 $\cos$ 是双胞胎。一个角的$\sin$，等于它互余角的$\cos$。”
\end{quote}

\subsection{特殊角的 $\cos$ 值}

\begin{center}
\begin{tabular}{|c|c|c|l|}
\hline
\textbf{角度} & $\cos$ & \textbf{大约是} & \textbf{感受} \\
\hline
$30^\circ$ & $\frac{\sqrt{3}}{2}$ & 0.87 & 邻边接近斜边长度 \\
\hline
$45^\circ$ & $\frac{\sqrt{2}}{2}$ & 0.71 & 邻边是斜边的七成 \\
\hline
$60^\circ$ & $\frac{1}{2}$ & 0.5 & 邻边是斜边的一半 \\
\hline
\end{tabular}
\end{center}

你发现了吗？

\begin{itemize}
\item $\sin$ 的表，从 $30^\circ$ 到 $60^\circ$ 越来越\textbf{大}
\item $\cos$ 的表，从 $30^\circ$ 到 $60^\circ$ 越来越\textbf{小}
\end{itemize}

\textbf{$\sin$ 和 $\cos$，一个变大，一个变小，永远相反。}

因为角度越大，对边越长，邻边越短。$\sin$ 大了，$\cos$ 就小了。

\subsection{应用题}

\textbf{题目五：求梯子底部距离}

\begin{quote}
“一架梯子，长 10 米，靠在墙上，梯子和地面的夹角是 $60^\circ$。梯子底部距离墙根多远？”
\end{quote}

\textbf{第一步：盯死夹角 $60^\circ$。}

\textbf{第二步：认清三条边的身份。}

\begin{center}
\begin{tabular}{|c|l|c|}
\hline
\textbf{边} & \textbf{是什么} & \textbf{长度} \\
\hline
斜边 & 梯子（斜的那条） & 10米 \\
\hline
邻边 & 梯子底部到墙根的距离（紧挨着角） & $?$ \\
\hline
\end{tabular}
\end{center}

\textbf{第三步：选公式。}

知道邻边和斜边，用 $\cos$。

\[\cos 60^\circ = \frac{\text{距离}}{10}\]
\[\frac{1}{2} = \frac{\text{距离}}{10}\]
\[\text{距离} = 10 \times \frac{1}{2} = 5 \text{ 米}\]

\textbf{验证}：$\dfrac{5}{10} = \dfrac{1}{2} = \cos 60^\circ$ \checkmark

\section{综合大题（本章BOSS）}

\textbf{题目六：两次测量求山高}

\begin{quote}
“小明站在山脚，看山顶的仰角是 $45^\circ$。他往山那边走了 100 米，现在仰角变成了 $60^\circ$。山高是多少？”
\end{quote}

这道题是这一章的大BOSS。但我们不硬拼，我们智取。

\subsection{第一步：定未知数}

\begin{center}
\begin{tabular}{|c|l|}
\hline
\textbf{未知数} & \textbf{代表什么} \\
\hline
$h$ & 山高 \\
\hline
$d$ & 最开始距离山脚的距离 \\
\hline
\end{tabular}
\end{center}

\subsection{第二步：翻译第一个画面（$45^\circ$）}

盯着 $45^\circ$ 看。对边是山高 $h$，邻边是距离 $d$。

\[\tan 45^\circ = \frac{h}{d}\]

因为 $\tan 45^\circ = 1$，所以：

\[1 = \frac{h}{d}\]
\[h = d\]

\textbf{停！这句话很重要：在$45^\circ$的时候，山有多高，你离山就有多远。}

\subsection{第三步：翻译第二个画面（$60^\circ$）}

你走了 100 米。所以现在的距离变成了 $d - 100$。但山高还是 $h$。

\[\tan 60^\circ = \frac{h}{d - 100}\]

因为 $\tan 60^\circ = \sqrt{3}$，所以：

\[\sqrt{3} = \frac{h}{d - 100}\]
\[h = \sqrt{3} \times (d - 100)\]

\subsection{第四步：解方程（用小数，别慌）}

我们有两个线索：
\begin{enumerate}
\item $h = d$
\item $h = \sqrt{3} \times (d - 100)$
\end{enumerate}

既然 $h = d$，那把第2个式子里的 $h$ 换成 $d$：

\[d = \sqrt{3} \times (d - 100)\]

\textbf{用大约的数来算}（$\sqrt{3} \approx 1.732$）：

\[d = 1.732 \times (d - 100)\]

展开括号：

\[d = 1.732d - 173.2\]

把 $d$ 归类（右边有 1.732 个 $d$，左边有 1 个 $d$）：

\[173.2 = 1.732d - 1d\]
\[173.2 = 0.732d\]
\[d = 173.2 \div 0.732 \approx 236.6\]

因为 $h = d$，所以\textbf{山高约 236.6 米}。

\subsection{验证}

\begin{center}
\begin{tabular}{|l|l|c|}
\hline
\textbf{检验} & \textbf{计算} & \textbf{结果} \\
\hline
第一个位置 & $\tan 45^\circ = \frac{236.6}{236.6} = 1$ & \checkmark \\
\hline
第二个位置 & $\tan 60^\circ = \frac{236.6}{136.6} \approx 1.73 \approx \sqrt{3}$ & \checkmark \\
\hline
\end{tabular}
\end{center}

\begin{quote}
\textbf{数学是工具，不是刑具。}

实在算不下去的时候，把根号变成小数，世界瞬间就清净了。
\end{quote}

\section{总结一下}

\subsection{三个公式}

\begin{center}
\begin{tabular}{|c|c|l|l|}
\hline
\textbf{三角函数} & \textbf{公式} & \textbf{用哪两条边} & \textbf{记忆法} \\
\hline
$\tan A$ & $\dfrac{\text{对边}}{\text{邻边}}$ & 对边和邻边 & 求坡度用它 \\
\hline
$\sin A$ & $\dfrac{\text{对边}}{\text{斜边}}$ & 对边和斜边 & 对边是分子 \\
\hline
$\cos A$ & $\dfrac{\text{邻边}}{\text{斜边}}$ & 邻边和斜边 & 邻边是分子 \\
\hline
\end{tabular}
\end{center}

\subsection{特殊角的值}

\begin{center}
\begin{tabular}{|c|c|c|c|}
\hline
\textbf{角度} & $\sin$ & $\cos$ & $\tan$ \\
\hline
$30^\circ$ & $\dfrac{1}{2}$ & $\dfrac{\sqrt{3}}{2}$ & $\dfrac{\sqrt{3}}{3}$ \\
\hline
$45^\circ$ & $\dfrac{\sqrt{2}}{2}$ & $\dfrac{\sqrt{2}}{2}$ & $1$ \\
\hline
$60^\circ$ & $\dfrac{\sqrt{3}}{2}$ & $\dfrac{1}{2}$ & $\sqrt{3}$ \\
\hline
\end{tabular}
\end{center}

\begin{quote}
\textbf{记忆法}：

“$\sin$从小到大，$\cos$从大到小。$45^\circ$是中间，$\sin$和$\cos$相等，$\tan$等于1。”
\end{quote}

\subsection{三步走}

\begin{center}
\begin{tabular}{|c|l|}
\hline
\textbf{步骤} & \textbf{做什么} \\
\hline
\textbf{第一步} & 盯死你要求的那个角 \\
\hline
\textbf{第二步} & 认清三条边的身份（对边、邻边、斜边） \\
\hline
\textbf{第三步} & 看你知道哪两条边，选对应的公式 \\
\hline
\end{tabular}
\end{center}

\section{现在你来}

\begin{quote}
\textbf{第一题}：一个斜坡，水平距离 12 米，垂直高度 6 米。求 $\tan$ 值，判断坡度角大约多少度。
\end{quote}

\begin{quote}
\textbf{第二题}：梯子长 8 米，和地面夹角 $45^\circ$。求梯子顶端离地多高。

提示：用 $\sin$
\end{quote}

\begin{quote}
\textbf{第三题}：飞机飞行高度 5000 米，和观测点连线 10000 米。求仰角。

提示：$\sin A = \frac{5000}{10000} = ?$
\end{quote}

写下来，验证一下。

\section{本章小结}

$\sin$、$\cos$、$\tan$，本质上就是三条边两两相除。

三种除法，三个名字。

用的时候，先盯死角，再认清边，再选公式。

\textbf{就这三步。题目再难，也逃不出这三步。}

\subsection*{本章知识地图}

\begin{verbatim}
三角函数
│
├─ 本质
│   └─ 三条边两两相除，三种除法
│
├─ 三个公式
│   ├─ tan = 对边 ÷ 邻边
│   ├─ sin = 对边 ÷ 斜边
│   └─ cos = 邻边 ÷ 斜边
│
├─ 三步走
│   ├─ ① 盯死要求的角
│   ├─ ② 认清三条边（对边、邻边、斜边）
│   └─ ③ 选对应的公式
│
├─ 特殊角的值
│   ├─ 30°：sin=1/2, cos=√3/2, tan=√3/3
│   ├─ 45°：sin=√2/2, cos=√2/2, tan=1
│   └─ 60°：sin=√3/2, cos=1/2, tan=√3
│
└─ sin 和 cos 的关系
    └─ sin A = cos(90°- A)（双胞胎）
\end{verbatim}

\chapter{数学是另一种语言——初中词汇表}

\section{我们停一下}

从初中第一章到现在，你学了很多新东西。

方程组、根号、三角函数、几何……

看起来很杂，对吧？

但如果你换个角度看，\textbf{你只是学了一些新单词。}

今天，我们不学新题。我们把这些“单词”整理一下。

看看数学这门语言，到底是怎么说话的。

\section{词汇表：代数篇（怎么描述关系）}

\begin{center}
\begin{tabular}{|c|l|l|l|}
\hline
\textbf{数学说} & \textbf{中文翻译} & \textbf{例子} & \textbf{潜台词} \\
\hline
$=$ & 是、等于 & $x = 5$ & 它俩是一回事，完全一样 \\
\hline
$\approx$ & 约等于 & $\pi \approx 3.14$ & 差不多得了，别太较真 \\
\hline
$x$ & 那个谁 & $x + 3 = 7$ & 我不知道它是几，但我知道它的位置 \\
\hline
$\begin{cases} \dots \\ \dots \end{cases}$ & 而且 & $\begin{cases} x+y=10 \\ x-y=2 \end{cases}$ & 这两件事必须\textbf{同时}发生 \\
\hline
$\sqrt{\phantom{0}}$ & 谁的平方是它？ & $\sqrt{9} = 3$ & 倒着找源头 \\
\hline
$x^2$ & 正方形的面积 & $x^2$ & $x$ 乘 $x$，也叫平方 \\
\hline
$x^3$ & 正方体的体积 & $x^3$ & $x$ 乘 $x$ 乘 $x$，也叫立方 \\
\hline
\end{tabular}
\end{center}

\section{词汇表：几何篇（怎么描述形状）}

\begin{center}
\begin{tabular}{|c|l|l|}
\hline
\textbf{数学说} & \textbf{中文翻译} & \textbf{潜台词} \\
\hline
$\pi$ & 圆周除以直径 & 圆的秘密代号，大约 3.14 \\
\hline
$90^\circ$ & 垂直 & 直挺挺的，不歪不斜 \\
\hline
$180^\circ$ & 平角 & 一条直线，完全平铺 \\
\hline
$360^\circ$ & 一圈 & 转回原地 \\
\hline
$a^2 + b^2 = c^2$ & 勾股定理 & 直角三角形三条边的铁律，$c$ 永远是斜边 \\
\hline
$\tan$ & 坡度（\textbf{求角度}） & 知道两条直角边，想求角，用它 \\
\hline
$\sin$ & 对边占斜边多少（\textbf{求角度}） & 知道对边和斜边，想求角，用它 \\
\hline
$\cos$ & 邻边占斜边多少（\textbf{求角度}） & 知道邻边和斜边，想求角，用它 \\
\hline
\end{tabular}
\end{center}

\section{词汇表：数字家族篇}

\begin{center}
\begin{tabular}{|c|l|l|}
\hline
\textbf{数学说} & \textbf{中文翻译} & \textbf{怎么认出来} \\
\hline
偶数 & 能被2整除的数 & 最后一位是 0、2、4、6、8 \\
\hline
奇数 & 不能被2整除的数 & 最后一位是 1、3、5、7、9 \\
\hline
质数 & 只能被1和自己整除 & 孤独的数，只有两个约数 \\
\hline
合数 & 除了1和自己，还能被其他数整除 & 热闹的数，有三个或更多约数 \\
\hline
系数 & $x$ 前面的数字 & 在 $5x$ 里，系数是 5 \\
\hline
常数 & 不带 $x$ 或 $y$ 的数 & 不会变的数 \\
\hline
\end{tabular}
\end{center}

\section{为什么说它是“语言”？}

回顾这一路走来，你一直在做一件事：

\textbf{翻译。}

\begin{center}
\begin{tabular}{|c|l|l|}
\hline
\textbf{阶段} & \textbf{中文说的} & \textbf{数学说的} \\
\hline
小学 & “两个苹果和三个苹果放在一起” & $2 + 3$ \\
\hline
初中 & “两个数加起来是10，减去是2” & $\begin{cases} x+y=10 \\ x-y=2 \end{cases}$ \\
\hline
刚学的 & “这个坡有多陡？” & $\tan A$ \\
\hline
\end{tabular}
\end{center}

你看，数学从来不是为了算题而存在的。

它是为了\textbf{描述这个世界}。

世界太复杂了，用中文说太啰嗦，用英文说也啰嗦。

所以人类发明了数学。

用几个符号，就能把“坡有多陡”（$\tan$）、“圆有多圆”（$\pi$）、“两件事同时发生”（方程组）说得清清楚楚。

\begin{quote}
\textbf{数学，就是宇宙通用的缩写。}
\end{quote}

\section{怎么学好这门语言？}

和学英语一样。

\begin{center}
\begin{tabular}{|c|l|l|}
\hline
\textbf{步骤} & \textbf{做什么} & \textbf{举例} \\
\hline
\textbf{第一：背单词} & 知道每个符号是什么意思 & $x$ 是未知数，$\sqrt{}$ 是开根号，$\sin$ 是求角度 \\
\hline
\textbf{第二：练造句} & 看到式子，用人话读出来 & 看到 $y = 2x + 1$，读成“$y$ 是 $x$ 的两倍再加1” \\
\hline
\textbf{第三：多对话} & 做应用题，就是和出题人对话 & 他用中文问你，你用数学回答他 \\
\hline
\end{tabular}
\end{center}

\textbf{这一章的表格，就是你的单词本。}

\section{你现在会什么？}

停一下，回头看看你走过的路。

从第一章的翻译三步法，到方程组、勾股定理、三角函数。

你知道你刚刚完成了什么吗？

你刚刚把\textbf{整个初中数学}最核心、最难啃的骨头，全部啃下来了。

\begin{center}
\begin{tabular}{|l|l|}
\hline
\textbf{技能} & \textbf{你能做什么} \\
\hline
\textbf{会翻译} & 给你一段复杂的中文，你能把它变成干净的数学式子 \\
\hline
\textbf{会解谜} & 给你几个未知的 $x$ 和 $y$，你能像侦探一样把它们找出来 \\
\hline
\textbf{会看图} & 给你一个图形，你能算出周长、面积、甚至每个角是多少度 \\
\hline
\textbf{会工具} & 你会用 $\sqrt{}$ 找平方的根，会用 $\sin$ 求角度，会用 $\pi$ 算圆 \\
\hline
\end{tabular}
\end{center}

\section{结语：你已经站在了高处}

你要知道，很多和你一样大的孩子，现在还在死记硬背公式，还在问“为什么要学数学”。

但你不一样。

你不仅学会了“怎么算”，你还明白了“为什么”。

\begin{center}
\begin{tabular}{|l|l|}
\hline
\textbf{别人只知道} & \textbf{你还知道} \\
\hline
$\pi \approx 3.14$ & $\pi$ 是圆周除以直径，古人用绳子绕着圆量出来的 \\
\hline
$a^2 + b^2 = c^2$ & 勾股定理可以用正方形拼出来 \\
\hline
$\sin 30^\circ = \frac{1}{2}$ & 这是从等边三角形切出来，用比例推出来的 \\
\hline
\end{tabular}
\end{center}

\textbf{你已经看懂了数学这门语言的逻辑。}

我必须认真地告诉你：

\textbf{你真的很了不起。}

这不是安慰奖，这是事实。

能坚持看到这里，能跟着我一步一步走到这里，说明你的脑子，比你以为的要聪明得多。

你已经具备了学好任何一门科学的基础。

因为科学，无非就是观察世界，然后用数学的语言把它写出来。

这件事，你已经会了。

前面的路还很长。还有函数，还有微积分，还有更广阔的世界。

但别怕。

你已经有了最锋利的剑（代数），最坚固的盾（几何）。

\textbf{去吧。去征服下一个山头。}

\textbf{我相信你。}

\subsection*{本章知识地图}

\begin{verbatim}
数学这门语言
│
├─ 代数词汇（描述关系）
│   ├─ = 等于
│   ├─ x 未知数
│   ├─ √ 开根号
│   ├─ x² 平方
│   └─ 方程组（同时成立）
│
├─ 几何词汇（描述形状）
│   ├─ π 圆的代号
│   ├─ 90°/180°/360° 特殊角
│   ├─ 勾股定理（求边）
│   └─ sin/cos/tan（求角）
│
├─ 数字家族
│   ├─ 偶数/奇数
│   ├─ 质数/合数
│   └─ 系数/常数
│
└─ 学好数学的方法
    ├─ 背单词（认识符号）
    ├─ 练造句（用人话读式子）
    └─ 多对话（做应用题）
\end{verbatim}

\chapter{接下来的“函数”，你做好准备了吗？}

\section*{停下来，深呼吸}

我们刚刚结束了一场漫长的旅行。

你学会了解方程，
学会了看图形，
学会了找角度。

但你有没有发现，
前面这十八章，
我们其实都在做同一件事？

我们在\textbf{拍照片}。

\section{以前的数学，是一张照片}

还记得我们的老朋友 $x$ 吗？

在第二章（以未知为核心的应用题）里，
我们说 $x$ 是嫌疑人，你是侦探。

比如：
\[x + 3 = 10\]

你一顿操作，算出 $x = 7$。

案子破了。
$x$ 被你抓住了，它就是 7，跑不掉了。

这就叫\textbf{拍照片}。
题目给你一个瞬间，一个固定的情况，
你把那个藏起来的、固定不变的数字找出来。

找到了，故事就结束了。

\section{但真实的世界，是在动的}

现实生活，
不是定格的照片。
是一部一直在播放的\textbf{电影}。

我给你举个例子：

“哥哥比弟弟大 3 岁。”

如果我问：
“弟弟今年 5 岁，哥哥几岁？”
你会算：$5 + 3 = 8$ 岁。

但明年呢？
弟弟 6 岁了，哥哥就是 9 岁。

十年后呢？
弟弟 15 岁了，哥哥就是 18 岁。

你看，
弟弟的年龄在变，
哥哥的年龄也在变。

这里没有一个固定的 $x$ 等着你去抓。
这里的数字，\textbf{活过来了}。

\section{怎么描述“活过来”的数学？}

既然数字在动，
我们以前那个找固定 $x$ 的方法就不够用了。

我们需要一种新的工具。

我们拿两个字母出来。
弟弟的年龄，叫 $x$。
哥哥的年龄，叫 $y$。

所以 $x$ 就是弟弟的岁数呐！$y$ 就是哥哥的岁数呐！

它们之间的关系，可以写成：

\[y = x + 3\]

盯住这个式子。
它和以前的 $x + 3 = 10$ 看起来很像，
但骨子里\textbf{完全不一样}。

在 $y = x + 3$ 里，
$x$ 可以是任何数！
它可以是 1，可以是 5，可以是 100。

而 $y$ 呢？
$y$ 是个跟屁虫。
$x$ 变成几，$y$ 就会乖乖地加上 3，跟着变成几。

$x$ 只要一动，$y$ 就跟着动。
它俩手拉手，跳了一支永远不会停下来的舞。

这种“一个数跟着另一个数一起变”的神奇关系，
在数学里，
就叫\textbf{函数}。

\section{别被“函数”这两个字吓到}

很多中学生，
一听到“函数”这两个字，头就嗡嗡作响。

其实，这都怪翻译的人。

“函数”的英文叫 Function，
意思是“功能”或者“作用”。

翻译成中文时，为什么叫“函”呢？

在古汉语里，“函”就是\textbf{盒子}、\textbf{信封}的意思。（比如信函、公函）

所以，什么是“函数”？

大白话翻译：
\textbf{它就是一个装数字的“魔法盒子”。}

\subsection{这个盒子是怎么工作的？}

想象你面前有一个机器盒子。
盒子上贴着一张纸条：\textbf{规则是 $+3$}。

\begin{enumerate}
\item 你从上面扔进去一个数字 \textbf{5}（这就是 $x$）。
\item 机器轰隆隆一转，给它加了 3。
\item 从下面吐出来一个数字 \textbf{8}（这就是 $y$）。
\end{enumerate}

这就是函数。

\begin{itemize}
\item 你扔进去的数（$x$），叫\textbf{输入}。
\item 吐出来的数（$y$），叫\textbf{输出}。
\item 盒子里的那个动作（$+3$），就是\textbf{函数的规则}。
\end{itemize}

就这么简单！

根本没有什么神秘的。
它就是一个数字加工厂。

买几斤苹果要付多少钱？这是一个函数（规则是乘单价）。
出租车开几公里要多少车费？这是一个函数（规则是起步价加里程费）。

\textbf{只要两个东西有关系，一个变，另一个就跟着变，它就是函数。}

\section{从抓犯人，到看世界}

学函数之前，
你是警察，你的任务是把 $x$ 按在墙上：“说，你到底是几！”

学函数之后，
你是导演。
你可以让 $x$ 随便变，
然后拿着爆米花，
慢慢观察 $y$ 是怎么跟着它一起变的。

数学的世界，
从这一刻起，
从静止的平面，变成了动态的电影。

\part{高中}
\setcounter{chapter}{0}
\chapter{情况与语言的惺惺相惜}

欢迎来到高中。

别紧张。

高中数学听起来吓人，但其实第一步特别温柔——它不让你算东西，它让你\textbf{说话}。

准确地说，它想教你一种\textbf{更精确的说话方式}。

为什么？

因为现实世界里，有太多``情况''需要描述。

而日常的中文，有时候说不清楚。

\section{你每天都在说``情况''}

来，你感受一下这几句话：

\begin{quote}
``及格的同学，站到左边去。''

``身高一米七以上的，来报名篮球队。''

``年龄在18岁到60岁之间的，可以献血。''
\end{quote}

你看到了吗？

这几句话都在做同一件事：\textbf{从一大群人里，按照某个条件，圈出一部分人}。

``及格的同学''——这是一群人。

``身高一米七以上的''——这也是一群人。

``18到60岁之间的''——还是一群人。

数学家每天都要跟这种``一群东西''打交道。

说多了嫌烦，于是他们发明了一个词：

\section{集合}

\textbf{集合，就是一个袋子。}

把满足某个条件的东西，全部装进这个袋子里。

这个袋子，就叫\textbf{集合}。

袋子里的每一样东西，叫\textbf{元素}。

就这么简单。

比如说：

\begin{quote}
``所有及格的同学''
\end{quote}

这就是一个集合。袋子里装的是张三、李四、王五……所有考了60分以上的人。

再比如：

\begin{quote}
``所有偶数''
\end{quote}

这也是一个集合。袋子里装的是 2, 4, 6, 8, 10, 12……永远装不完，但它就是一个袋子。

再比如：

\begin{quote}
``我家冰箱里的水果''
\end{quote}

苹果、橘子、香蕉——三样东西，装进一个袋子。这也是一个集合。

\section{集合怎么写？}

数学家写字很懒，他们最喜欢发明\textbf{缩写}。

集合的写法是这样的：

你给袋子取个名字，通常用\textbf{大写字母}：$A$、$B$、$C$……

然后用\textbf{大括号}把袋子里的东西列出来。

比如，我家冰箱里的水果：

\[
A = \{苹果, 橘子, 香蕉\}
\]

读作：``集合A等于苹果、橘子、香蕉。''

再比如，1到5的所有整数：

\[
B = \{1, 2, 3, 4, 5\}
\]

就是这样。大括号就是``袋子''，里面的东西就是``元素''。

（大括号 $\{ \}$ 你以前可能见过，但没用过。现在它有了正式工作：\textbf{当袋子}。）

\section{``在不在袋子里？''}

好，现在袋子有了，你自然要回答一个问题：

\begin{quote}
\textbf{某样东西，在不在这个袋子里？}
\end{quote}

比如：

\begin{quote}
$A = \{苹果, 橘子, 香蕉\}$
\end{quote}

问你：\textbf{苹果在不在袋子 $A$ 里？}

当然在。

这句话，数学家写成：

\[
苹果 \in A
\]

这个符号 $\in$，读作\textbf{``属于''}。

意思就是：\textbf{``在袋子里''}。

苹果 $\in$ $A$，翻译成人话就是：``苹果在袋子A里。''

那\textbf{西瓜在不在袋子 $A$ 里？}

不在。冰箱里根本没有西瓜。

这句话写成：

\[
西瓜 \notin A
\]

这个符号 $\notin$，读作\textbf{``不属于''}。

意思就是：\textbf{``不在袋子里''}。

来，试试：

\begin{quote}
$B = \{1, 2, 3, 4, 5\}$
\end{quote}

$3$ 在不在 $B$ 里？

\textbf{$3 \in B$}。 ✓（3在袋子里，没问题。）

$7$ 在不在 $B$ 里？

\textbf{$7 \notin B$}。 ✓（7没在袋子里，袋子只装了1到5。）

这就是``属于''和``不属于''的全部内容。

（你说：就这？就这。）

\section{岔个路：比大小的符号}

接下来的内容会用到一些``比大小''的符号。

这些符号你可能见过，但咱们趁现在先说清楚。

\begin{quote}
$3 > 1$
\end{quote}

读作：3\textbf{大于}1。

那个张嘴的方向，永远朝着\textbf{大的那个数}。你可以想象它是一只嘴巴——\textbf{嘴巴想吃大的}。

\begin{quote}
$1 < 3$
\end{quote}

读作：1\textbf{小于}3。

还是那只嘴巴，还是朝着大的那个数3。

其实 $3 > 1$ 和 $1 < 3$ 说的是\textbf{同一件事}，只是方向反过来了。

\begin{quote}
$3 = 3$
\end{quote}

读作：3\textbf{等于}3。

这个没什么好说的。一样大，就画两横。

\begin{quote}
$3 \geq 3$
\end{quote}

读作：3\textbf{大于或等于}3。

注意，不是``大于并且等于''，是``大于\textbf{或者}等于''。

3比3大吗？不大。

3和3相等吗？相等。

两个条件满足一个就行，所以 $3 \geq 3$ 是\textbf{对的}。

再来一个：$5 \geq 3$，对不对？

5比3大吗？大。

那就够了。满足一个就行。$5 \geq 3$ 是\textbf{对的}。

\begin{quote}
$3 \leq 5$
\end{quote}

读作：3\textbf{小于或等于}5。

同理，满足``小于''或者``等于''其中一个就行。

3比5小吗？小。那就对了。

总结一下，六个符号：

\begin{center}
\begin{tabular}{ccc}
\toprule
符号 & 读法 & 大白话 \\
\midrule
$>$ & 大于 & 左边比右边大 \\
$<$ & 小于 & 左边比右边小 \\
$=$ & 等于 & 一样大 \\
$\geq$ & 大于或等于 & 左边 ≥ 右边（大了或一样都算） \\
$\leq$ & 小于或等于 & 左边 ≤ 右边（小了或一样都算） \\
$\neq$ & 不等于 & 不一样大 \\
\bottomrule
\end{tabular}
\end{center}

（记不住没关系，用几次就熟了。嘴巴朝大的那个，这一条记住就够了。）

\section{回到集合：用条件来描述袋子}

刚才我们的集合都是把东西\textbf{一个一个列出来}的：

\[
A = \{苹果, 橘子, 香蕉\}
\]

\[
B = \{1, 2, 3, 4, 5\}
\]

但是，如果袋子里的东西太多呢？

比如：\textbf{所有大于0的整数}。

你要一个一个写？

$\{1, 2, 3, 4, 5, 6, 7, 8, 9, 10, 11, 12, ......\}$

写到明年也写不完。

所以数学家发明了另一种写法——\textbf{用条件来描述袋子里装了什么}。

像这样：

\[
C = \{x \mid x \text{ 是大于0的整数}\}
\]

别被吓到，我来一个一个拆给你看：

$C =$ ：这个袋子叫 $C$。

$\{$ $\}$：大括号，还是那个袋子。

$x$：代表袋子里的任意一样东西。（$x$ 就是一个代号，代表``袋子里的某个元素''。）

$\mid$：这条竖线读作``\textbf{使得}''，或者你可以理解为``\textbf{条件是}''。竖线左边是东西，右边是条件。

$x \text{ 是大于0的整数}$：条件。只要满足这个条件的东西，都能进袋子。

翻译成人话：

\begin{quote}
``袋子 $C$ 里装的，是所有满足‘大于0的整数’这个条件的东西。''
\end{quote}

也就是 $1, 2, 3, 4, 5, 6, 7, \ldots$ 全部装进去。

再来一个：

\[
D = \{x \mid x > 3\}
\]

翻译：``袋子 $D$ 里装的，是所有大于3的数。''

$4$ 在袋子 $D$ 里吗？$4 > 3$，满足条件，\textbf{在}。$4 \in D$。 ✓

$3$ 在袋子 $D$ 里吗？$3 > 3$ 吗？不，$3 = 3$，不是大于。\textbf{不在}。$3 \notin D$。 ✓

$3.0001$ 在袋子 $D$ 里吗？$3.0001 > 3$，满足条件，\textbf{在}。$3.0001 \in D$。 ✓

（你看，这个袋子能装无穷多的东西。但它被一个条件描述得\textbf{干干净净}。）

\section{这，就是``情况''}

我让你注意一件事。

看看刚才那个集合：

\[
D = \{x \mid x > 3\}
\]

这句话在说什么？

它在说：\textbf{当 $x$ 满足``大于3''这个条件的时候。}

换个说法：

\begin{quote}
\textbf{``$x > 3$ 的情况下。''}
\end{quote}

你有没有在应用题里见过类似的话？

\begin{quote}
``当速度大于60的时候……''

``如果温度在0度以下……''

``票价在100元以内的……''
\end{quote}

这些，全部都是\textbf{``情况''}。

而集合，就是用数学语言\textbf{把这些情况精精确确地写出来}的工具。

\textbf{情况是生活里的事。语言是数学里的符号。它们俩天生一对——惺惺相惜。}

这就是本章标题的意思。

\section{空袋子也是袋子}

好，来点好玩的。

\begin{quote}
``我家冰箱里的恐龙。''
\end{quote}

请问，这个集合有几个元素？

零个。

你家冰箱里没有恐龙。（如果有，请联系动物园。）

但注意——我们\textbf{确实描述了一个集合}。只不过这个袋子里\textbf{啥也没有}。

这种空空的袋子，数学家给了它一个专门的名字：

\[
\emptyset
\]

叫做\textbf{空集}。

$\emptyset$ 就是一个空袋子。

它不是``不存在''。它存在——只不过里面什么都没装。

就像你打开一个快递箱子，发现里面是空的。箱子在，东西没有。

再来一个例子：

\[
E = \{x \mid x > 5 \text{ 并且 } x < 3\}
\]

翻译：$x$ 要同时大于5、又小于3。

有这种数吗？

一个数怎么可能同时大于5又小于3？

不存在。

所以 $E = \emptyset$。空集。

（这不是数学家在整你，是这个条件本身就矛盾了。空集的意思就是：\textbf{这个情况不可能发生}。）

有一条很重要的规矩，先记住：

\begin{quote}
\textbf{空集是任何集合的子集。}
\end{quote}

你现在可能想问：什么是子集？

别急，马上就讲。

\section{袋子里套袋子——子集}

想象一个场景：

你们班有40个人。这40个人是一个集合，我们叫它 $A$。

\[
A = \{\text{全班40个同学}\}
\]

现在老师说：``戴眼镜的同学，站起来。''

站起来的有15个人。这15个人也是一个集合，叫它 $B$。

\[
B = \{\text{全班戴眼镜的同学}\}
\]

请问：$B$ 里面的每一个人，是不是\textbf{也在} $A$ 里面？

当然是。因为他们首先是班里的人，然后才是戴眼镜的人。

每一个戴眼镜的同学，\textbf{都属于}全班同学。

我们说：

\[
B \subseteq A
\]

读作：``$B$ 是 $A$ 的\textbf{子集}。''

什么意思呢？

\begin{quote}
\textbf{$B$ 里面的每一样东西，$A$ 里面全都有。}
\end{quote}

或者说：

\begin{quote}
\textbf{小袋子里的东西，大袋子里一个不少。}
\end{quote}

$B$ 是小袋子，$A$ 是大袋子。小袋子完完全全被大袋子``包住''了。

再举个例子，用数字：

\[
A = \{1, 2, 3, 4, 5\}
\]

\[
B = \{1, 3, 5\}
\]

$B$ 里面有 $1$，$A$ 里面有没有？有。 ✓

$B$ 里面有 $3$，$A$ 里面有没有？有。 ✓

$B$ 里面有 $5$，$A$ 里面有没有？有。 ✓

$B$ 的每一个元素，$A$ 全都有。

所以 $B \subseteq A$。 ✓

那反过来呢？

$A$ 里面有 $2$，$B$ 里面有没有？没有。 ✗

$A$ 的元素，$B$ \textbf{不是全都有}。

所以 $A \subseteq B$ 吗？\textbf{不是}。

（方向反过来就不一定对了。大袋子装不进小袋子。）

那如果两个袋子\textbf{一模一样}呢？

\[
A = \{1, 2, 3\}
\]

\[
B = \{1, 2, 3\}
\]

$A$ 的每个元素，$B$ 里有吗？全有。所以 $A \subseteq B$。 ✓

$B$ 的每个元素，$A$ 里有吗？全有。所以 $B \subseteq A$。 ✓

两边都是子集——说明它俩\textbf{一模一样}。

所以 $A = B$。

（这其实就是数学家证明``两个集合相等''的方法：你是我的子集，我也是你的子集，那咱们就是同一个袋子。）

\section{真子集——``像，但不完全像''}

有时候我们想强调：\textbf{小袋子确实被大袋子包住了，而且小袋子比大袋子少点东西。}

比如：

\[
A = \{1, 2, 3, 4, 5\}
\]

\[
B = \{1, 3, 5\}
\]

$B \subseteq A$，而且 $B \neq A$（$B$ 比 $A$ 少了 $2$ 和 $4$）。

这时候我们说：

\[
B \subsetneq A
\]

读作：``$B$ 是 $A$ 的\textbf{真子集}。''

意思就是：

\begin{quote}
\textbf{被包住了，但不完全一样。小袋子确实比大袋子少了点东西。}
\end{quote}

$\subseteq$ 和 $\subsetneq$ 的区别，就像 $\leq$ 和 $<$ 的区别：

\begin{quote}
$\subseteq$：小于或等于（可以一样大，也可以更小）

$\subsetneq$：严格小于（一定更小，不能一样大）
\end{quote}

还记得刚才说的那条规矩吗？

\begin{quote}
\textbf{空集是任何集合的子集。}
\end{quote}

\[
\emptyset \subseteq A
\]

不管 $A$ 是什么集合，这句话永远是对的。

为什么？

因为空集里\textbf{一个元素都没有}。

你要检查的是：``空集里的每一个元素，$A$ 里面有没有？''

空集里有元素吗？没有。

既然一个都没有，那就\textbf{不存在}反例。没有反例，就默认通过。

（这有点赖皮，对吧？但数学家就是这么定义的。逻辑上说得通。你可以把它想成：\textbf{一个空袋子，塞进任何大袋子里都不碍事。}）

\section{交集——两个袋子的``重叠部分''}

好，袋子学完了，现在来玩袋子。

你们班40个人。

\begin{quote}
集合 $A$ = 戴眼镜的同学 = \{张三, 李四, 王五, 赵六, 孙七\}

集合 $B$ = 参加篮球队的同学 = \{李四, 孙七, 周八, 吴九\}
\end{quote}

现在老师问：

\begin{quote}
\textbf{``既戴眼镜，又参加了篮球队的同学，有哪些？''}
\end{quote}

你看看：两个袋子里\textbf{都出现了}的名字是谁？

李四——$A$ 里有，$B$ 里也有。 ✓

孙七——$A$ 里有，$B$ 里也有。 ✓

其他人都只出现在一个袋子里。

所以答案是 \{李四, 孙七\}。

这个``两个袋子都有的部分''，叫做\textbf{交集}。

写成符号：

\[
A \cap B = \{李四, 孙七\}
\]

$\cap$ 这个符号，读作\textbf{``交''}。

你可以这样记：$\cap$ 长得像一个\textbf{倒扣的碗}，两个袋子的东西往碗里倒，\textbf{只有两边都有的才能留下来}。

大白话：

\begin{quote}
\textbf{交集 = 两个条件同时满足。``既……又……''}
\end{quote}

再用数字来一个：

\[
A = \{1, 2, 3, 4, 5\}
\]

\[
B = \{3, 4, 5, 6, 7\}
\]

两边都有的：$3, 4, 5$。

\[
A \cap B = \{3, 4, 5\}
\]

两边没有公共元素怎么办？

\[
A = \{1, 2, 3\}
\]

\[
B = \{4, 5, 6\}
\]

一个重叠的都没有。

\[
A \cap B = \emptyset
\]

又见面了——空集。两个袋子没有交集，重叠部分是空的。

\section{并集——两个袋子``合起来''}

还是刚才的例子：

\begin{quote}
$A$ = \{张三, 李四, 王五, 赵六, 孙七\}

$B$ = \{李四, 孙七, 周八, 吴九\}
\end{quote}

现在换个问题：

\begin{quote}
\textbf{``戴眼镜的，或者参加篮球队的，或者两样都占的——总共有哪些人？''}
\end{quote}

把两个袋子的东西全部\textbf{倒在一起}：

张三, 李四, 王五, 赵六, 孙七, 李四, 孙七, 周八, 吴九。

等一下——李四和孙七出现了两次。

没关系，\textbf{同一个人只算一次}。去掉重复的：

\{张三, 李四, 王五, 赵六, 孙七, 周八, 吴九\}

这就是\textbf{并集}。

写成符号：

\[
A \cup B = \{张三, 李四, 王五, 赵六, 孙七, 周八, 吴九\}
\]

$\cup$ 这个符号，读作\textbf{``并''}。

你可以这样记：$\cup$ 长得像一个\textbf{正放的碗}，两个袋子的东西全部往碗里倒，\textbf{来者不拒，但不重复}。

大白话：

\begin{quote}
\textbf{并集 = 满足条件A，或者满足条件B，或者两个都满足。``……或……''}
\end{quote}

用数字来一个：

\[
A = \{1, 2, 3, 4, 5\}
\]

\[
B = \{3, 4, 5, 6, 7\}
\]

全倒在一起，去掉重复：

\[
A \cup B = \{1, 2, 3, 4, 5, 6, 7\}
\]

注意区分：

\begin{quote}
\textbf{交集}（$\cap$）：两边\textbf{都有}的。关键词是``\textbf{并且}''、``\textbf{同时}''。

\textbf{并集}（$\cup$）：\textbf{至少一边有}的。关键词是``\textbf{或者}''。
\end{quote}

（这两个符号长得很像，容易搞混。记住：$\cap$ 是倒碗——严格，只要公共的；$\cup$ 是正碗——宽容，全都收。）

\section{全集与补集——``剩下的那些''}

有时候，我们需要说的不是``哪些在袋子里''，而是——

\begin{quote}
\textbf{``哪些不在袋子里？''}
\end{quote}

但``不在袋子里''这个说法太模糊了。不在哪个袋子里？范围是什么？

所以我们需要先定一个\textbf{最大的袋子}——也就是\textbf{讨论的范围}。

这个最大的袋子，叫做\textbf{全集}，用字母 $U$ 表示。

（$U$ 是英文 Universal 的首字母，意思是``全部的、宇宙的''。你可以想成：宇宙就是最大的袋子。）

比如说：

你们全班40个人。这就是今天讨论的范围，也就是全集 $U$。

\[
U = \{\text{全班40个同学}\}
\]

戴眼镜的同学：

\[
A = \{张三, 李四, 王五, 赵六, 孙七\}
\]

现在问你：

\begin{quote}
\textbf{``不戴眼镜的同学有哪些？''}
\end{quote}

就是全班40个人里，\textbf{去掉}戴眼镜的那些，\textbf{剩下的}。

``剩下的''这个概念，就叫\textbf{补集}。

写成符号：

\[
\complement_U A
\]

读作：``$A$ 关于全集 $U$ 的\textbf{补集}。''

大白话：

\begin{quote}
\textbf{大袋子里有的，小袋子里没有的。也就是``剩下的''。}
\end{quote}

如果全班有40个人，戴眼镜的有5个，那不戴眼镜的就有35个。

$\complement_U A$ 就是那35个人组成的集合。

用数字来一个：

\[
U = \{1, 2, 3, 4, 5, 6, 7, 8, 9, 10\}
\]

\[
A = \{2, 4, 6, 8, 10\}
\]

$\complement_U A = ?$

$U$ 里面有的，$A$ 里面没有的：$1, 3, 5, 7, 9$。

\[
\complement_U A = \{1, 3, 5, 7, 9\}
\]

（你发现了吗？$A$ 装的是偶数，$\complement_U A$ 装的是奇数。偶数的``补集''就是奇数——只要你的全集是1到10的所有整数。）

再来一个，和``情况''挂钩的：

\begin{quote}
全集 $U$ = 所有实数（就是数轴上所有的数）。

$A = \{x \mid x > 3\}$（大于3的数）。
\end{quote}

$\complement_U A$ 是什么？

大袋子（所有实数）里有的，小袋子（大于3的数）里没有的。

也就是：\textbf{不大于3的数}。

``不大于3''是什么意思？就是\textbf{小于或等于3}。

\[
\complement_U A = \{x \mid x \leq 3\}
\]

（看到了吗？$>$ 的补集是 $\leq$。大于的反面，不是小于，是\textbf{小于或等于}——别忘了等于3本身不在``大于3''里面，所以它要归到补集那边去。这个小细节，考试最爱考。）

\section{用集合说``情况''——这才是重点}

好了，所有符号和概念都认识了。

现在来到这一章最关键的部分：\textbf{集合到底有什么用？}

答案是：\textbf{它能帮你精确地描述应用题里的``情况''。}

来看一道题：

\begin{quote}
一家餐厅规定：

年龄满18岁及以上，可以点酒。

年龄小于12岁，可以享受儿童优惠。

问：谁既不能点酒，也享受不了儿童优惠？
\end{quote}

如果不用集合，你可能要在脑子里绕半天。

但用集合来看，清清楚楚：

设 $U$ = 所有顾客（按年龄分）。

\[
A = \{x \mid x \geq 18\} \quad \text{（能点酒的人）}
\]

\[
B = \{x \mid x < 12\} \quad \text{（享受儿童优惠的人）}
\]

``既不能点酒，也享受不了儿童优惠''——既不在 $A$ 里，也不在 $B$ 里。

不在 $A$ 里：$x < 18$。

不在 $B$ 里：$x \geq 12$。

两个条件\textbf{同时满足}——这是什么？\textbf{交集}。

\[
\complement_U A \cap \complement_U B = \{x \mid x < 18 \text{ 并且 } x \geq 12\}
\]

也就是：

\[
\{x \mid 12 \leq x < 18\}
\]

\textbf{12岁到18岁之间的人}（包括12岁，不包括18岁）。

你看——集合的语言，把一个绕来绕去的生活问题，翻译得\textbf{清清楚楚}。

这就是``情况''和``语言''的惺惺相惜。

情况需要被精确描述。数学语言刚好能做到。

天生一对。

\section{练习时间}

来，练练手。答案直接给你，你先自己想，想完再对。

\textbf{第一组：属于还是不属于？}

$A = \{2, 4, 6, 8\}$

（1）$4$ \_\_\_\_（$\in$ 还是 $\notin$）$A$ ？

答案：$4 \in A$。（4在袋子里。）

（2）$5$ \_\_\_\_（$\in$ 还是 $\notin$）$A$ ？

答案：$5 \notin A$。（5不在袋子里，袋子里只有偶数。）

（3）$8$ \_\_\_\_（$\in$ 还是 $\notin$）$A$ ？

答案：$8 \in A$。（8在。）

\textbf{第二组：子集判断}

$A = \{1, 2, 3, 4, 5\}$，$B = \{2, 4\}$，$C = \{1, 2, 3, 4, 5\}$

（1）$B \subseteq A$ 吗？

答案：\textbf{是}。$B$ 里的 $2$ 和 $4$，$A$ 里全有。

（2）$A \subseteq B$ 吗？

答案：\textbf{不是}。$A$ 里有 $1, 3, 5$，但 $B$ 里没有。

（3）$C \subseteq A$ 吗？

答案：\textbf{是}。$C$ 和 $A$ 一模一样，每个元素对方都有。

（4）$C \subsetneq A$ 吗？

答案：\textbf{不是}。$C$ 和 $A$ 完全一样，真子集要求``少点东西''，但 $C$ 不比 $A$ 少。

（5）$\emptyset \subseteq A$ 吗？

答案：\textbf{是}。空集是任何集合的子集。空袋子塞哪都行。

\textbf{第三组：交集和并集}

$A = \{1, 3, 5, 7\}$，$B = \{2, 3, 5, 8\}$

（1）$A \cap B = $ ？

两边都有的：$3, 5$。

答案：$A \cap B = \{3, 5\}$。

（2）$A \cup B = $ ？

全部倒在一起，去掉重复：$1, 2, 3, 5, 7, 8$。

答案：$A \cup B = \{1, 2, 3, 5, 7, 8\}$。

\textbf{第四组：补集}

$U = \{1, 2, 3, 4, 5, 6, 7, 8, 9, 10\}$，$A = \{1, 3, 5, 7, 9\}$

$\complement_U A = $ ？

$U$ 里有但 $A$ 里没有的：$2, 4, 6, 8, 10$。

答案：$\complement_U A = \{2, 4, 6, 8, 10\}$。

（$A$ 是奇数，补集就是偶数。完美。）

\textbf{第五组：翻译``情况''}

用集合语言把下面这句话翻译成数学：

\begin{quote}
``气温在零下10度到零上35度之间（包括两端），适合户外活动。''
\end{quote}

设 $x$ 为气温。

答案：$A = \{x \mid -10 \leq x \leq 35\}$

（$x$ 就是气温呐。大于等于零下10度，小于等于35度。两个条件同时满足。）

\textbf{挑战题：}

\begin{quote}
$U$ = 全班同学。

$A$ = 会游泳的同学。

$B$ = 会骑自行车的同学。
\end{quote}

请用集合符号表示：

（1）``既会游泳，又会骑车的同学'' = \_\_\_\_

答案：$A \cap B$

（2）``会游泳或者会骑车（或者都会）的同学'' = \_\_\_\_

答案：$A \cup B$

（3）``不会游泳的同学'' = \_\_\_\_

答案：$\complement_U A$

（4）``既不会游泳，也不会骑车的同学'' = \_\_\_\_

答案：$\complement_U A \cap \complement_U B$

（翻译翻译：不在 $A$ 里，也不在 $B$ 里。两个补集取交集。）

（5）\textbf{想一想：} 第（4）题，有没有另一种写法？

提示：$\complement_U (A \cup B)$

这是什么意思？先把会游泳\textbf{或}会骑车的人合起来（$A \cup B$），然后取补集——剩下的就是两样都不会的。

和第（4）题答案一样。

（这个规律叫做\textbf{德摩根定律}，名字不用记，但道理你已经懂了。）

这一章到这里就结束了。

集合不难。它就是\textbf{袋子 + 条件 + 关系}。

但这门语言从今往后会反复出现。函数的定义域是集合，不等式的解是集合，概率的样本空间还是集合。

你今天学的，不是一个孤立的知识点。

你学的是一门\textbf{通用语言}。

以后你会发现，越来越多的数学题，本质上都在问：

\begin{quote}
\textbf{哪些东西满足这个条件？}

\textbf{这些东西和那些东西有什么关系？}
\end{quote}

而你已经知道怎么回答了。

\textbf{第一章完成 ✓}

\chapter{情况的细审}

上一章我们聊了集合。

你知道了：集合就是袋子，元素就是袋子里的东西。

你还学了一堆操作：交集、并集、补集。

但那些操作，当时可能感觉有点``飘在空中''——你知道怎么算，但不太确定它们\textbf{到底什么时候会用上}。

这一章，它们全部会落地。

因为这一章讲的是``\textbf{情况}''。

而情况，就是集合在现实世界里的样子。

\section{从一个场景开始}

想象你走进一家游乐场。

门口立了个牌子：

\begin{quote}
\textbf{``身高1.2米以上才能坐过山车。''}
\end{quote}

你站到量身高的机器前面。

滴——1.3米。

能坐吗？1.3米比1.2米高，能坐。 ✓

你弟弟也去量了一下。

滴——1.1米。

能坐吗？1.1米比1.2米矮，不能坐。 ✗

你妈妈帮你弟弟求情：``差一点点嘛，就让他坐嘛。''

工作人员摇摇头：``规定就是规定。1.2米\textbf{以上}。''

好。这件事，用数学怎么写？

设 $x$ 是身高（单位：米）。

``身高1.2米以上''——翻译成数学：

\[
x \geq 1.2
\]

（还记得 $\geq$ 吗？（就是上一章讲的``大于\textbf{或者}等于''。）1.2米本身也算``以上''，所以用 $\geq$，不用 $>$。）

这就是一个\textbf{情况}。

但游乐场的规定，不会只有一条。

你往旁边一看，还有一块牌子：

\begin{quote}
\textbf{``身高超过2米的游客，出于安全考虑，也不能坐过山车。''}
\end{quote}

啊？太高了也不行？

对。有些过山车的座位就是卡不下太高的人。

所以现在，能坐过山车的人，要满足\textbf{两个}条件：

\begin{quote}
条件一：$x \geq 1.2$（不能太矮）

条件二：$x \leq 2$（不能太高）
\end{quote}

并且，这两个条件要\textbf{同时满足}。

不是满足一个就行——你不能说``我虽然只有1米，但我也不超过2米啊，所以让我坐吧''。

不行。两条都要满足。

\textbf{同时满足。}

\section{``同时满足''——交集回来了}

``同时满足两个条件''这件事，你在上一章其实已经见过了。

上一章我们讲交集的时候说过：

\begin{quote}
交集 = 两个袋子\textbf{都有}的部分。
\end{quote}

（还记得吗？$A \cap B$，那个倒扣的碗，只留两边都有的。）

现在把它翻译成``情况''的语言：

\begin{quote}
条件一圈出了一袋子人（身高 $\geq 1.2$ 的人）。

条件二圈出了另一袋子人（身高 $\leq 2$ 的人）。

两个条件\textbf{同时满足}的人，就是这两个袋子的\textbf{交集}。
\end{quote}

也就是：既在第一个袋子里，又在第二个袋子里的那些人。

写成数学：

\[
\{x \mid x \geq 1.2\} \cap \{x \mid x \leq 2\} = \{x \mid x \geq 1.2 \text{ 并且 } x \leq 2\}
\]

这一行有点长。别慌，我拆给你看。

左边：两个集合取交集。

右边：把交集的结果写出来——$x$ 要同时满足 $\geq 1.2$ 和 $\leq 2$。

就是这样。交集在这里的意思，就是``两个条件同时满足''。

\textbf{``同时满足'' = 交集。}

以后你在题目里看到``\textbf{并且}''、``\textbf{同时}''、``\textbf{既……又……}''这些词，脑子里就要自动弹出：\textbf{交集}。

\section{双向不等式——一种偷懒的写法}

``$x \geq 1.2$ 并且 $x \leq 2$''这句话，写起来有点长。

数学家最讨厌写长句子。

（真的，你以后会发现，数学里一半的发明，都是为了少写几个字。）

所以他们把这两个条件\textbf{合并成一行}：

\[
1.2 \leq x \leq 2
\]

把 $x$ 夹在中间。左边一个条件，右边一个条件。

读法：``$x$ 大于或等于1.2，并且小于或等于2。''

大白话：``身高在1.2米到2米之间，两头都算。''

你看看，这三种写法说的是\textbf{完全一样的事}：

写法一（两个条件分开写）：$x \geq 1.2$ 并且 $x \leq 2$

写法二（集合写法）：$\{x \mid x \geq 1.2\} \cap \{x \mid x \leq 2\}$

写法三（夹在中间）：$1.2 \leq x \leq 2$

从长到短，意思没变。只是越来越\textbf{懒}。

以后你在试卷上最常见的就是写法三。但你要知道，它的背后就是两个条件取交集。

我再多说一个细节。

$1.2 \leq x \leq 2$ 里面，两边都用了 $\leq$。

这意味着 $x = 1.2$ 和 $x = 2$ \textbf{都算在里面}。

如果规定改成``身高\textbf{超过}1.2米，\textbf{不到}2米''呢？

``超过''是\textbf{严格大于}，用 $>$。``不到''是\textbf{严格小于}，用 $<$。

\[
1.2 < x < 2
\]

这时候1.2米和2米本身\textbf{不包括}。刚好1.2米的人，不能坐。刚好2米的人，也不能坐。

$\leq$ 和 $<$ 的区别就在这里——\textbf{端点到底算不算}。

\begin{center}
\begin{tabular}{ccc}
\toprule
写法 & 端点算不算？ & 大白话 \\
\midrule
$1.2 \leq x \leq 2$ & 两头都算 & 1.2到2，\textbf{包括}1.2和2 \\
$1.2 < x < 2$ & 两头都不算 & 1.2到2，\textbf{不包括}1.2和2 \\
$1.2 \leq x < 2$ & 左算右不算 & \textbf{包括}1.2，\textbf{不包括}2 \\
$1.2 < x \leq 2$ & 左不算右算 & \textbf{不包括}1.2，\textbf{包括}2 \\
\bottomrule
\end{tabular}
\end{center}

四种组合，就是 $\leq$ 和 $<$ 的排列组合。

别背，用的时候看题目说``以上''``超过''``不到''``以下''这些词，判断端点算不算就行。

\section{来，你试试}

我说几个生活场景，你来翻译成双向不等式。

别急着看答案，先自己写一下。哪怕写在脑子里也行。

\textbf{场景一：}

\begin{quote}
``冰箱冷藏室的温度应该保持在2度到8度之间，包括2度和8度。''
\end{quote}

设 $x$ 是温度。

你的翻译：\_\_\_\_\_\_

答案：$2 \leq x \leq 8$

（两头都``包括''，所以两边都用 $\leq$。）

\textbf{场景二：}

\begin{quote}
``这款手机的售价在2000元以上，5000元以下。''
\end{quote}

设 $x$ 是售价（元）。

你的翻译：\_\_\_\_\_\_

答案：$2000 \leq x \leq 5000$

（``以上''包含本身，``以下''也包含本身。所以两边 $\leq$。）

\textbf{场景三：}

\begin{quote}
``婴儿的标准体温：高于36.5度，低于37.5度。''
\end{quote}

设 $x$ 是体温。

你的翻译：\_\_\_\_\_\_

答案：$36.5 < x < 37.5$

（``高于''不含本身，``低于''也不含本身。$36.5$ 度和 $37.5$ 度本身不算``正常范围''。所以两边 $<$。）

\textbf{场景四：}

\begin{quote}
``青少年组的参赛年龄：年满12岁，不满18岁。''
\end{quote}

设 $x$ 是年龄。

你的翻译：\_\_\_\_\_\_

答案：$12 \leq x < 18$

（``年满''12岁——12岁算在内，$\leq$。``不满''18岁——18岁不算，$<$。左算右不算。）

如果你四个都对了——太好了。

如果错了一两个——没关系，再看看``算''和``不算''那个表格，多看两遍就熟了。

\section{``或者''——另一种叠加}

好。刚才我们讲的是``同时满足两个条件''——交集。

现在来看另一种情况。

你弟弟没坐上过山车，不太开心。

你带他去另一个项目——碰碰车。

门口的牌子写着：

\begin{quote}
\textbf{``以下两个条件，满足其中一个就能玩：}

\textbf{① 年龄在12岁以上；}

\textbf{② 有家长陪同。''}
\end{quote}

你弟弟8岁，不满足条件①。

但你妈妈陪着他——满足条件②。

能玩吗？

能。

因为规定说的是``满足\textbf{其中一个}就行''。

不需要两个条件都满足。只要占一头，就可以。

``满足其中一个就行''——这个你也见过。

上一章我们讲并集的时候说过：

\begin{quote}
并集 = 两个袋子的东西\textbf{全部倒在一起}。
\end{quote}

（还记得吗？$A \cup B$，那个正放的碗，来者不拒，但不重复。）

翻译成``情况''的语言：

\begin{quote}
满足条件一的人，装一袋子。

满足条件二的人，装一袋子。

满足条件一\textbf{或者}条件二（\textbf{或者两个都满足}）的人，就是两个袋子的\textbf{并集}。
\end{quote}

\textbf{``或者'' = 并集。}

以后你在题目里看到``\textbf{或者}''、``\textbf{至少满足一个}''、``\textbf{……或……}''这些词，脑子里就要自动弹出：\textbf{并集}。

等一下。我要强调一件很容易搞错的事。

数学里的``或者''，跟日常生活里的``或者''\textbf{不太一样}。

生活里你说``要咖啡或者茶''，通常的意思是``二选一''——你不会两杯都要。

但数学里的``或者''是\textbf{宽容的``或者''}：

\begin{quote}
满足A？可以。 ✓

满足B？可以。 ✓

两个都满足？\textbf{也可以。} ✓
\end{quote}

只有\textbf{两个都不满足}，才不行。 ✗

所以碰碰车那个例子：一个15岁的人，自己就满足了``12岁以上''，不需要家长陪。但如果他家长也来了——那他两个条件都满足了，\textbf{当然更可以}。

数学里的``或者'' = ``至少一个''，\textbf{不是``只能选一个''}。

这一点，请你记牢。

\section{来，再试试}

\textbf{场景一：}

\begin{quote}
``图书馆的免费借阅资格：本校学生，或者持有社区图书卡。''
\end{quote}

设 $A$ = 本校学生的集合，$B$ = 持有社区图书卡的人的集合。

能免费借阅的人 = \_\_\_\_

答案：$A \cup B$

（``或者'' → 并集。是学生就行，有卡也行，两样都有更行。）

\textbf{场景二：}

\begin{quote}
``报名马拉松比赛的条件：年龄在18岁以上，\textbf{并且}最近体检合格。''
\end{quote}

设 $A$ = 年龄 $\geq 18$ 的集合，$B$ = 体检合格的集合。

能报名的人 = \_\_\_\_

答案：$A \cap B$

（``并且'' → 交集。两个条件缺一不可。）

\textbf{场景三（这个稍微绕一点）：}

\begin{quote}
``可以领取儿童补贴的条件：年龄不满6岁，\textbf{或者}家庭月收入低于3000元。''
\end{quote}

设 $A$ = 年龄 $< 6$ 的集合，$B$ = 月收入 $< 3000$ 的集合。

能领补贴的人 = \_\_\_\_

答案：$A \cup B$

（还是``或者'' → 并集。4岁小孩家里月入两万？能领。30岁成年人月入两千？也能领。4岁小孩家里月入两千？当然更能领。）

现在你应该能分辨了：

\begin{quote}
\textbf{看到``同时''``并且''``既……又……'' → 交集 $\cap$}

\textbf{看到``或者''``至少一个''``……或……'' → 并集 $\cup$}
\end{quote}

这两个词就像两个开关。题目说了哪个词，你就按哪个开关。

\section{``不满足''——情况的反面}

好，还有一种操作，你也见过了。

上一章讲补集的时候说过：

\begin{quote}
补集 = 大袋子里有的，小袋子里没有的。剩下的那些。
\end{quote}

（还记得吗？$\complement_U A$，就是``不在 $A$ 里面''的那些东西。）

翻译成``情况''的语言，就是：

\textbf{不满足某个条件的那些东西。}

听起来很简单，对吧？

但这里面藏了一个非常容易踩的坑。

我用一个例子慢慢讲。

\begin{quote}
条件：``$x$ 大于3。''
\end{quote}

也就是 $x > 3$。

现在我问你：\textbf{``不满足这个条件''是什么意思？}

你可能脱口而出：``$x$ 小于3。''

\textbf{错了。}

不是 $x < 3$，是 $x \leq 3$。

为什么？

因为``不大于3''的意思是：$x$ 没有超过3。

$x = 2$？没超过3。 ✓（不大于3。）

$x = 0$？没超过3。 ✓（不大于3。）

$x = 3$ 呢？

$3 > 3$ 吗？不，3不大于3，3等于3。

所以 $x = 3$ \textbf{也不满足} $x > 3$ 这个条件。

$x = 3$ 应该被归到``不满足''那一边。

所以``不大于3''= ``小于3，\textbf{或者等于3}'' = $x \leq 3$。

这就是取反的规律。我把所有的对应关系列出来：

\begin{center}
\begin{tabular}{cc}
\toprule
原来的条件 & 取反之后（``不满足''原条件） \\
\midrule
$x > 3$（大于3） & $x \leq 3$（小于\textbf{或等于}3） \\
$x < 3$（小于3） & $x \geq 3$（大于\textbf{或等于}3） \\
$x \geq 3$（大于或等于3） & $x < 3$（小于3） \\
$x \leq 3$（小于或等于3） & $x > 3$（大于3） \\
\bottomrule
\end{tabular}
\end{center}

你看到规律了吗？

取反的时候，发生了两件事：

\textbf{第一：方向反了。} 大于 ↔ 小于。

\textbf{第二：等号翻转了。} 有等号的变成没等号，没等号的变成有等号。

$>$ 变成 $\leq$（方向反了，而且加了等号）。

$\geq$ 变成 $<$（方向反了，而且去了等号）。

\textbf{不是 $>$ 变 $<$！不是简单地反方向！等号也要翻！}

我再说一遍，因为这个地方真的特别多人搞错：

\[
x > 3 \text{ 的反面不是 } x < 3
\]

\[
x > 3 \text{ 的反面是 } x \leq 3
\]

为什么？因为 $x = 3$ 这个点，不满足 $x > 3$，所以它要站到反面去。而 $x < 3$ 没有包含 $x = 3$，所以 $x < 3$ 不够——得用 $x \leq 3$。

如果这段你看了两遍还觉得绕，那就记住这句话：

\begin{quote}
\textbf{取反的时候，$>$ 和 $\leq$ 是一对，$<$ 和 $\geq$ 是一对。交叉配对，不是平行配对。}
\end{quote}

来，我出几个，你来取反。

\textbf{（1）} 原条件：$x > 10$

取反：\_\_\_\_\_\_

答案：$x \leq 10$

（方向反了，加了等号。$>$ 变 $\leq$。）

\textbf{（2）} 原条件：$x \leq 5$

取反：\_\_\_\_\_\_

答案：$x > 5$

（方向反了，去了等号。$\leq$ 变 $>$。）

\textbf{（3）} 原条件：$x \geq 18$

取反：\_\_\_\_\_\_

答案：$x < 18$

（$\geq$ 变 $<$。方向反了，去了等号。）

\textbf{（4）} 原条件：$x < 0$

取反：\_\_\_\_\_\_

答案：$x \geq 0$

（$<$ 变 $\geq$。方向反了，加了等号。）

如果你四个都对了，说明取反的规律已经吃透了。

如果错了——大概率是等号那个地方出了问题。回去再看看那个表格，特别注意``交叉配对''。

\section{复合条件的取反——``并且''和``或者''会互换}

刚才取反的都是单个条件。

但如果原来的情况是\textbf{两个条件叠加在一起}的呢？

回到游乐场。

过山车的条件是：

\begin{quote}
``$x \geq 1.2$ \textbf{并且} $x \leq 2$''
\end{quote}

也就是 $1.2 \leq x \leq 2$。

现在问你：\textbf{什么样的人不能坐过山车？}

就是\textbf{不满足}上面那个条件的人。

不满足``$x \geq 1.2$ 并且 $x \leq 2$''——这意味着什么？

我们一步一步想。

``并且''的意思是两个条件\textbf{都要满足}。

那``\textbf{不能}两个条件都满足''是什么意思？

就是\textbf{至少有一个条件不满足}。

要么 $x \geq 1.2$ 不满足——也就是 $x < 1.2$（太矮了）。

要么 $x \leq 2$ 不满足——也就是 $x > 2$（太高了）。

\textbf{只要有一头不满足，就不能坐。}

``只要有一头''——这是``或者''啊！

所以：

\[
\text{不满足``} x \geq 1.2 \text{ 并且 } x \leq 2 \text{''} \quad = \quad x < 1.2 \quad \textbf{或者} \quad x > 2
\]

你看到了吗？

\textbf{``并且''取反，变成了``或者''。}

我用大白话再说一遍：

原来的条件是：不能太矮，\textbf{并且}不能太高。（两条都要守。）

取反之后变成：太矮了，\textbf{或者}太高了。（违反一条就不行。）

有道理吧？

如果你需要\textbf{同时}满足两条规定才能过关，那只要你\textbf{任何一条}没满足，你就过不了关。

\textbf{``同时满足所有''的反面，是``至少违反一条''。}

\textbf{``并且''翻面变``或者''。}

反过来也成立。

碰碰车的条件是：

\begin{quote}
``年龄 $\geq 12$ \textbf{或者}有家长陪同''
\end{quote}

什么人不能玩碰碰车？

\textbf{不满足条件①，并且不满足条件②。}

年龄不够（$< 12$），\textbf{并且}没人陪。

两条都不满足，才不能玩。

\textbf{``或者''翻面变``并且''。}

这个规律，我再帮你总结一遍：

\begin{center}
\begin{tabular}{cc}
\toprule
原来 & 取反之后 \\
\midrule
A \textbf{并且} B & （不A）\textbf{或者}（不B） \\
A \textbf{或者} B & （不A）\textbf{并且}（不B） \\
\bottomrule
\end{tabular}
\end{center}

取反的时候，做两件事：

\textbf{第一：每个小条件各自取反。}（$\geq$ 变 $<$，$>$ 变 $\leq$，等等。）

\textbf{第二：``并且''和``或者''互换。}

这两件事缺一不可。

来，完整做一个。

\begin{quote}
原条件：$x \geq 0$ \textbf{并且} $x < 100$
\end{quote}

取反：

第一步：$x \geq 0$ 取反 → $x < 0$。$x < 100$ 取反 → $x \geq 100$。

第二步：``并且''变``或者''。

结果：$x < 0$ \textbf{或者} $x \geq 100$。

翻译成人话：要么是负数，要么大于等于100。

有道理吗？

原来是``0到100之间（包含0，不包含100）''。不在这个范围里的，要么跑到左边去了（小于0），要么跑到右边去了（大于等于100）。

完全说得通。 ✓

再来一个。

\begin{quote}
原条件：$x < 5$ \textbf{或者} $x > 20$
\end{quote}

取反：

第一步：$x < 5$ 取反 → $x \geq 5$。$x > 20$ 取反 → $x \leq 20$。

第二步：``或者''变``并且''。

结果：$x \geq 5$ \textbf{并且} $x \leq 20$。

也就是 $5 \leq x \leq 20$。

原来是``小于5或者大于20''——两头的部分。

取反之后变成``5到20之间''——中间的部分。

两头取反变中间，中间取反变两头。 ✓

如果你觉得这里有点绕——正常。

这可能是这一章最绕的地方了。

但你只要记住那两句话：

\begin{quote}
\textbf{``并且''翻面变``或者''。}

\textbf{``或者''翻面变``并且''。}
\end{quote}

然后每个小条件各自取反（符号交叉配对）。

做多了就自然了。

\section{区间——一种更偷懒的写法}

到这里，你已经学了三种技能：

\begin{quote}
两个条件``同时满足''（并且 → 交集）

两个条件``满足其一''（或者 → 并集）

条件取反（符号交叉翻，并且/或者互换）
\end{quote}

现在来学最后一个工具：\textbf{区间}。

区间不是什么新概念。它就是``一段范围''的\textbf{缩写}。

你已经会写 $1.2 \leq x \leq 2$ 了。

区间只是把这种写法\textbf{再缩短一点}。

$1.2 \leq x \leq 2$ 写成区间：

\[
[1.2,\ 2]
\]

就这么短。

逗号左边是范围的\textbf{起点}，右边是\textbf{终点}。

两边的括号告诉你\textbf{端点算不算}。

\textbf{方括号 $[\ ]$ = 端点算在内（包含）。}

\textbf{圆括号 $(\ )$ = 端点不算（不包含）。}

怎么记？

你想象括号是一只手。

\textbf{方括号像一只手，手指弯过来把端点抓住了。} 抓住了 = 包含。

\textbf{圆括号像一只手滑开了，没抓住。} 没抓住 = 不包含。

（或者你就死记：\textbf{方 = 包，圆 = 不包}。四个字，搞定。）

四种组合：

\begin{center}
\begin{tabular}{cccc}
\toprule
区间写法 & 对应的不等式 & 左端点 & 右端点 \\
\midrule
$[a,\ b]$ & $a \leq x \leq b$ & 包含 & 包含 \\
$(a,\ b)$ & $a < x < b$ & 不包含 & 不包含 \\
$[a,\ b)$ & $a \leq x < b$ & 包含 & 不包含 \\
$(a,\ b]$ & $a < x \leq b$ & 不包含 & 包含 \\
\bottomrule
\end{tabular}
\end{center}

来，把刚才做过的那些场景写成区间：

\textbf{冰箱冷藏室：} $2 \leq x \leq 8$ → $[2, 8]$

\textbf{婴儿体温：} $36.5 < x < 37.5$ → $(36.5, 37.5)$

\textbf{青少年参赛年龄：} $12 \leq x < 18$ → $[12, 18)$

是不是特别干净？

\section{无穷大——当范围没有尽头}

有时候，范围是\textbf{一头有尽头，另一头没有}的。

比如：

\begin{quote}
``身高1.2米以上''
\end{quote}

\[
x \geq 1.2
\]

1.2米是起点。但终点呢？没有终点。身高1.5米可以，2米可以，10米也可以（虽然现实中不可能，但数学不管现实，它只管逻辑）。

这种``没有尽头''的情况，我们用一个符号来表示：

\[
+\infty
\]

读作``正无穷大''。

$x \geq 1.2$ 写成区间：

\[
[1.2,\ +\infty)
\]

注意：$+\infty$ 那一头，\textbf{永远用圆括号}。

为什么？

因为 $+\infty$ 不是一个具体的数。你不可能``到达''它，更不可能``包含''它。

所以你永远抓不住它。手滑开了。圆括号。

同样的道理，如果范围往左边没有尽头：

\begin{quote}
``$x$ 小于5''
\end{quote}

\[
x < 5
\]

写成区间：

\[
(-\infty,\ 5)
\]

$-\infty$ 是``负无穷大''。左边也是圆括号。

5本身不包含（因为是 $<$，不是 $\leq$），所以右边也是圆括号。

如果是 $x \leq 5$：

\[
(-\infty,\ 5]
\]

左边还是圆括号（负无穷永远抓不住），右边方括号（5算在内）。

还记得刚才``不能坐过山车''的条件吗？

\[
x < 1.2 \quad \text{或者} \quad x > 2
\]

这是两段范围。一段在左边，一段在右边。

写成区间：

\[
(-\infty,\ 1.2) \cup (2,\ +\infty)
\]

中间那个 $\cup$ 是什么？

并集。（还记得吗？两个袋子合起来。``或者''对应并集。）

两段范围，用 $\cup$ 连起来，表示``在左边那段\textbf{或者}在右边那段''。

\section{这一章你学到了什么？}

让我帮你理一理。

情况的\textbf{三种叠加方式}：

\begin{quote}
\textbf{``并且''} → 两个条件同时满足 → \textbf{交集} $\cap$ → 双向不等式 $a \leq x \leq b$

\textbf{``或者''} → 满足其中一个就行 → \textbf{并集} $\cup$

\textbf{``取反''} → 不满足这个条件 → \textbf{补集} → 符号交叉翻，并且/或者互换
\end{quote}

还有\textbf{区间}——一种更短的写法：

\begin{quote}
方括号 $[\ ]$ = 包含端点

圆括号 $(\ )$ = 不包含端点

无穷大永远用圆括号
\end{quote}

这些东西看起来是一堆新符号，但它们做的事情其实你从小就在做：

\textbf{判断一个东西在不在某个范围里。}

小时候你用手指比划：``这么高能坐吗？''

现在你用 $1.2 \leq x \leq 2$ 来写。

本质没变。只是语言变精确了。

\textbf{第二章完成 ✓}

\chapter{学会判断真假——数学的正式语文课}

前两章，你学了两样东西：

\textbf{集合}——用袋子装东西。

\textbf{情况}——用条件描述范围。

但到目前为止，我们一直在``描述''。

描述这个袋子里有什么，描述那个范围是多少。

从这一章开始，我们要做一件新的事情：

\textbf{判断。}

判断一句话，是对的，还是错的。

\section{数学其实是语文}

你可能觉得这个标题很奇怪——数学就是数学，跟语文有什么关系？

关系大了。

你想想看，数学在干什么？

它在\textbf{说话}。

$1 + 1 = 2$——这是一句话。

$3 > 5$——这也是一句话。

$x^2 = 4$ 的解是 $x = 2$——这还是一句话。

每一句话，都在做一个\textbf{判断}。

而每一个判断，要么是\textbf{对的}，要么是\textbf{错的}。

$1 + 1 = 2$？对的。

$3 > 5$？错的。

这种``做了一个判断，而且能分辨对错''的句子，数学家给它起了个名字：

\textbf{命题。}

\section{什么是命题？}

命题就是\textbf{一句能判断真假的话}。

就这么简单。

不需要它是对的——错的也算命题。

只需要你\textbf{能判断}它是对还是错。

我说几句话，你来感受一下：

\begin{quote}
``2是偶数。''
\end{quote}

这是命题吗？

是。它做了一个判断。而且这个判断是\textbf{对的}——2确实是偶数。

我们说：这是一个\textbf{真命题}。

\begin{quote}
``3是偶数。''
\end{quote}

这是命题吗？

是。它也做了一个判断。但这个判断是\textbf{错的}——3是奇数，不是偶数。

我们说：这是一个\textbf{假命题}。

请注意——\textbf{假命题也是命题。}

很多人一看到``假的''，就觉得``那它就不是命题了吧？''

不对。假的也是命题。

命题的定义不是``一句正确的话''，而是``一句\textbf{能判断真假}的话''。

能判断对 → 真命题。

能判断错 → 假命题。

判断不了 → 不是命题。

什么样的话``判断不了''呢？

\begin{quote}
``今天天气真好啊！''
\end{quote}

这是命题吗？

不是。``好''是主观的，你觉得好，你朋友可能觉得太热了。没办法客观判断对错。

\begin{quote}
``快把门关上！''
\end{quote}

这是命题吗？

不是。这是一个命令，不是一个判断。它没有对错可言。

\begin{quote}
``你几岁了？''
\end{quote}

这是命题吗？

不是。这是一个问题。问题不是判断，你没法说一个问题``对''还是``错''。

\begin{quote}
``$x > 3$。''
\end{quote}

这是命题吗？

这个要小心了。

如果你不告诉我 $x$ 是几，我就没法判断。$x = 5$ 的时候是对的，$x = 1$ 的时候是错的。

它的真假\textbf{取决于 $x$ 是谁}。

这种句子，数学上不把它当作命题，而是叫它\textbf{条件}（也就是上一章说的``情况''——带着未知数的判断句，真假还没确定）。

但如果我把 $x$ 定下来：

\begin{quote}
``$5 > 3$。''
\end{quote}

现在没有未知数了。能判断了。$5$ 确实大于 $3$。

这就是一个真命题。

我帮你总结一下：

\begin{center}
\begin{tabular}{p{4cm}p{4cm}p{6cm}}
\toprule
句子 & 是不是命题？ & 为什么？ \\
\midrule
``2是偶数。'' & ✓ 真命题 & 做了判断，判断是对的 \\
``3是偶数。'' & ✓ 假命题 & 做了判断，判断是错的 \\
``地球是方的。'' & ✓ 假命题 & 做了判断，判断是错的（但它依然是命题） \\
``$1 + 1 = 3$'' & ✓ 假命题 & 做了判断，判断是错的 \\
``快跑！'' & ✗ 不是命题 & 命令，不是判断 \\
``你吃饭了吗？'' & ✗ 不是命题 & 问题，不是判断 \\
``这道题好难啊。'' & ✗ 不是命题 & 主观感受，无法客观判断 \\
``$x > 3$'' & ✗ 不是命题（是条件） & 真假取决于 $x$，还没确定 \\
\bottomrule
\end{tabular}
\end{center}

你来判断几个：

\textbf{（1）} ``所有的奇数都是质数。''

是命题吗？\_\_\_\_\_\_

如果是，真还是假？\_\_\_\_\_\_

答案：是命题。假命题。（9是奇数，但9不是质数——$9 = 3 \times 3$。找到一个反例就够了。）

\textbf{（2）} ``如果今天是周一，那明天是周二。''

是命题吗？\_\_\_\_\_\_

如果是，真还是假？\_\_\_\_\_\_

答案：是命题。真命题。（周一的下一天确实是周二。）

\textbf{（3）} ``把窗户打开。''

是命题吗？\_\_\_\_\_\_

答案：不是命题。这是一个命令。

\textbf{（4）} ``$2 + 3 = 6$''

是命题吗？\_\_\_\_\_\_

如果是，真还是假？\_\_\_\_\_\_

答案：是命题。假命题。（$2 + 3 = 5$，不是6。但它做了判断，所以是命题——只不过判断错了。）

\textbf{（5）} ``平行四边形的对角相等。''

是命题吗？\_\_\_\_\_\_

如果是，真还是假？\_\_\_\_\_\_

答案：是命题。真命题。（平行四边形的对角确实相等，这是几何里的定理。）

\section{命题的结构——``如果……那么……''}

好。你现在知道什么是命题了。

接下来，我们来看命题的\textbf{内部结构}。

很多命题，都长成一个固定的样子：

\begin{quote}
\textbf{``如果……，那么……''}
\end{quote}

比如：

\begin{quote}
``如果一个数能被2整除，那么它是偶数。''

``如果今天下雨，那么地面会湿。''

``如果一个三角形是等边三角形，那么它的三个角都是60度。''
\end{quote}

这种``如果A，那么B''的句子，是数学里\textbf{最常见}的命题结构。

在这种结构里，有两个部分：

\textbf{``如果''后面的部分，叫做条件。}

\textbf{``那么''后面的部分，叫做结论。}

\begin{quote}
``如果一个数能被2整除（\textbf{条件}），那么它是偶数（\textbf{结论}）。''
\end{quote}

条件是\textbf{出发点}——你先假设它成立。

结论是\textbf{到达点}——你想知道它能不能成立。

整句话在问的是：\textbf{从条件出发，能不能走到结论？}

能走到 → 真命题。

走不到 → 假命题。

我用一个比喻来说：

想象你站在一条路的起点。起点是``条件''，终点是``结论''。

如果从起点出发，\textbf{一定}能走到终点——这条路是通的——这个命题是\textbf{真的}。

如果从起点出发，\textbf{不一定}能走到终点——路可能是断的——这个命题是\textbf{假的}。

来看几个例子：

\begin{quote}
``如果一个人是北京人，那么他是中国人。''
\end{quote}

条件：是北京人。

结论：是中国人。

从``北京人''出发，能走到``中国人''吗？

能。北京人一定是中国人。

真命题。 ✓

\begin{quote}
``如果一个人是中国人，那么他是北京人。''
\end{quote}

条件：是中国人。

结论：是北京人。

从``中国人''出发，能走到``北京人''吗？

不一定。中国人可以是上海人、广州人、成都人……

假命题。 ✗

注意！这两句话\textbf{条件和结论刚好反过来了}。

第一句是真的，第二句是假的。

\textbf{反过来说，不一定还成立。}

这一点非常重要。很多人做题会在这里栽跟头——觉得``如果A那么B''是对的，就想当然地认为``如果B那么A''也是对的。

不一定。

方向很重要。从起点到终点的路是通的，不代表从终点回起点的路也是通的。

\section{现在来到重点：充分条件和必要条件}

你已经知道了``如果A，那么B''的结构。

现在我要给A和B之间的关系\textbf{命名}。

别被名字吓到。名字听起来很唬人，但道理你\textbf{已经懂了}。

我只是在给你\textbf{已经懂的东西}贴标签。

\subsection{充分条件——``有它就够了''}

再看这句话：

\begin{quote}
``如果一个人是北京人，那么他一定是中国人。''
\end{quote}

``是北京人''这个条件，\textbf{足够}让你断定他是中国人。

你不需要再去查他的户口本、看他的身份证。你知道他是北京人，就\textbf{够了}。

这时候我们说：

\begin{quote}
\textbf{``是北京人''是``是中国人''的充分条件。}
\end{quote}

\textbf{充分，就是``充足''、``足够''的意思。}

有了这个条件，结论\textbf{一定}成立。这个条件对得出结论来说，\textbf{够了}。

用数学来一个：

\begin{quote}
``如果 $x > 5$，那么 $x > 3$。''
\end{quote}

$x$ 大于5，那它一定大于3吗？

当然。5都超过了，3肯定也超过了。

所以``$x > 5$''是``$x > 3$''的\textbf{充分条件}。

知道 $x > 5$，就\textbf{足够}断定 $x > 3$。不用再问别的了。够了。

\subsection{必要条件——``没它不行''}

现在换个角度。

\begin{quote}
``如果一个人是北京人，那么他一定是中国人。''
\end{quote}

这句话反过来问一下：

\textbf{如果一个人不是中国人，他有可能是北京人吗？}

不可能。不是中国人，怎么可能是北京人？北京在中国啊。

也就是说：

\textbf{想成为北京人，``是中国人''是一个前提条件。这个前提不能少。}

没有它，后面那个就\textbf{不可能}成立。

这时候我们说：

\begin{quote}
\textbf{``是中国人''是``是北京人''的必要条件。}
\end{quote}

\textbf{必要，就是``必须''、``不可缺少''的意思。}

没有这个条件，结论\textbf{不可能}成立。这个条件对结论来说，是\textbf{必须的}。

用数学来一个：

\begin{quote}
``$x > 3$''是``$x > 5$''的什么条件？
\end{quote}

想一想：

如果 $x$ 连3都没超过（$x \leq 3$），它有可能超过5吗？

不可能。连3都没过，怎么可能过5？

所以``$x > 3$''是``$x > 5$''的\textbf{前提}。没有它不行。

``$x > 3$''是``$x > 5$''的\textbf{必要条件}。

但是！

光知道 $x > 3$，能断定 $x > 5$ 吗？

不能。$x = 4$ 就大于3，但不大于5。

所以``$x > 3$''对``$x > 5$''来说，虽然\textbf{必须}，但\textbf{不够}。

必要，但不充分。

\subsection{把充分和必要放在一起}

我知道你可能有点绕了。

没关系。这个地方确实绕。绝大多数人第一次学的时候都会混。

我用一个你绝对忘不掉的例子，帮你一次搞定。

\begin{quote}
条件A：这个动物是狗。

条件B：这个动物有四条腿。
\end{quote}

\textbf{问：A是B的什么条件？}

先问第一个方向：\textbf{``如果A，那么B''对不对？}

如果这个动物是狗，那它有四条腿吗？

有。（正常情况下。）

\textbf{A → B ✓。 有A就够了。A是B的充分条件。}

再问第二个方向：\textbf{``如果B，那么A''对不对？}

如果这个动物有四条腿，那它一定是狗吗？

不一定。可能是猫，可能是马，可能是牛、羊……

总之不一定是狗。

\textbf{B → A ✗。 有B不够。B不是A的充分条件。}

所以结论是：

\textbf{A是B的充分条件} ✓（是狗就一定有四条腿——够了）

\textbf{A是B的必要条件吗？} 也就是问：没有A的话，B还能成立吗？不是狗的话，有没有可能有四条腿？当然有可能——猫也有四条腿。所以A对B来说不是必须的。\textbf{A不是B的必要条件。} ✗

合起来说：\textbf{A是B的充分不必要条件。}

再反过来看：

\textbf{B是A的什么条件？}

B → A？有四条腿就一定是狗？不一定。 ✗（B不是A的充分条件）

没有B的话，A能成立吗？如果没有四条腿，它能是狗吗？正常来说不能——狗得有四条腿。所以B对A是必须的。 ✓（B是A的必要条件）

\textbf{B是A的必要不充分条件。}

你看到了吗？

\textbf{充分和必要，是一枚硬币的两面。}

\begin{quote}
当A是B的\textbf{充分}条件时，B就是A的\textbf{必要}条件。
\end{quote}

永远是这样。它们是一对的。

\subsection{我帮你提炼一个万能判断法}

以后不管碰到什么题目，你都用这个方法：

\textbf{第一步：弄清楚谁是条件A，谁是条件B。}

\textbf{第二步：问两个方向。}

\begin{quote}
A → B？（如果A成立，B一定成立吗？）

B → A？（如果B成立，A一定成立吗？）
\end{quote}

\textbf{第三步：看结果，对号入座。}

\begin{center}
\begin{tabular}{ccc}
\toprule
A → B & B → A & A是B的什么条件？ \\
\midrule
✓ & ✗ & 充分不必要条件 \\
✗ & ✓ & 必要不充分条件 \\
✓ & ✓ & 充要条件（又充分又必要） \\
✗ & ✗ & 既不充分也不必要 \\
\bottomrule
\end{tabular}
\end{center}

就这四种情况。没有第五种。

等一下。第三行出现了一个新词：\textbf{充要条件}。

\subsection{充要条件——``就是它''}

有时候，两个条件之间的关系是\textbf{双向通的}。

比如：

\begin{quote}
条件A：$x = 3$

条件B：$2x = 6$
\end{quote}

A → B：如果 $x = 3$，那么 $2x = 2 \times 3 = 6$。对。 ✓

B → A：如果 $2x = 6$，那么 $x = 6 \div 2 = 3$。对。 ✓

两个方向都通！

A能推出B，B也能推回A。

这时候我们说：

\begin{quote}
\textbf{A和B互为充要条件。}
\end{quote}

\textbf{``充要''就是``既充分，又必要''。}

充分——有A就一定有B（够了）。

必要——没A就一定没B（少不了）。

两样都占了。

换句话说：\textbf{A和B说的其实是同一件事。只是换了个写法。}

$x = 3$ 和 $2x = 6$——就是同一件事的两种表达。

再来一个：

\begin{quote}
条件A：一个三角形有两条边相等。

条件B：这个三角形是等腰三角形。
\end{quote}

A → B：两条边相等，那它就是等腰三角形。 ✓

B → A：等腰三角形，就一定有两条边相等。 ✓

双向通。\textbf{充要条件。}

（其实这就是等腰三角形的\textbf{定义}。定义本身，几乎都是充要条件——因为定义就是在说``A和B是同一件事''。）

\subsection{用集合来记忆}

这一招是给你的保底方法。如果你怎么想都觉得绕，就用这个。

\begin{quote}
条件A对应一个集合（满足A的所有东西装进袋子A）。

条件B对应一个集合（满足B的所有东西装进袋子B）。
\end{quote}

（还记得吗？前两章说的——情况和集合是一体两面。每个条件都对应一个袋子。）

然后看两个袋子的大小关系：

\textbf{袋子A被装在袋子B里面（$A \subseteq B$）：}

满足A的东西，一定满足B。→ A是B的\textbf{充分条件}。

但满足B的东西，不一定满足A（因为袋子B更大，里面还有别的东西）。→ A不是B的\textbf{必要条件}。

所以：\textbf{小袋子里的条件，是大袋子的充分条件。}

\textbf{袋子B被装在袋子A里面（$B \subseteq A$）：}

满足B的东西，一定满足A。→ B是A的充分条件。

反过来说：A是B的\textbf{必要条件}。（大袋子的条件，是小袋子的必要条件——你得先在大袋子里，才有可能在小袋子里。）

\textbf{袋子A和袋子B一样大（$A = B$）：}

互相包含，双向通。→ \textbf{充要条件。}

\textbf{两个袋子有交叉但谁也不包含谁：}

两个方向都不一定。→ \textbf{既不充分也不必要。}

用刚才的例子验证一下：

\begin{quote}
A = ``是狗'' → 袋子A = \{所有的狗\}

B = ``有四条腿'' → 袋子B = \{所有四条腿的动物\}
\end{quote}

所有的狗都在``四条腿动物''这个大袋子里吗？

在。$A \subseteq B$。小袋子在大袋子里面。

所以A（小袋子）是B（大袋子）的\textbf{充分条件}。

B（大袋子）是A（小袋子）的\textbf{必要条件}。

和我们之前的判断完全一致。 ✓

\textbf{记忆口诀：}

\begin{quote}
\textbf{小袋子对大袋子：充分。}（在小袋子里，一定在大袋子里——够了。）

\textbf{大袋子对小袋子：必要。}（要进小袋子，得先在大袋子里——必须的。）
\end{quote}

\subsection{最容易犯的错}

我必须说这个，因为几乎所有人都在这里栽过。

\textbf{错误：看到``如果A那么B''是真的，就以为``如果B那么A''也是真的。}

不是的！

``如果下雨，那么地面湿。'' → 真的。

``如果地面湿，那么下了雨。'' → 不一定！可能是有人洒水了，可能是水管爆了。

方向反过来，结论就不一定成立。

\textbf{永远要分开检查两个方向。}

这是这一章最重要的一句话。

好，你来判断几个。先自己想，然后再看答案。

\textbf{（1）}

\begin{quote}
条件A：$x = 2$

条件B：$x^2 = 4$
\end{quote}

A → B：如果 $x = 2$，那么 $x^2 = 2^2 = 4$。 ✓

B → A：如果 $x^2 = 4$，那么 $x = 2$？

不一定。$x = -2$ 的时候，$(-2)^2 = 4$，也满足 $x^2 = 4$。所以 $x$ 不一定是2。 ✗

\textbf{答案：A是B的充分不必要条件。}

（用集合看：$A = \{2\}$，$B = \{2, -2\}$。$A \subseteq B$。小袋子在大袋子里。小对大：充分。）

\textbf{（2）}

\begin{quote}
条件A：$x > 3$

条件B：$x > 5$
\end{quote}

A → B：$x > 3$ 就一定 $x > 5$？$x = 4$ 就大于3但不大于5。 ✗

B → A：$x > 5$ 就一定 $x > 3$？超过5了肯定超过3。 ✓

\textbf{答案：A是B的必要不充分条件。}

（想要 $x > 5$，首先得 $x > 3$——这是必须的。但光 $x > 3$ 还不够。）

（用集合看：$A = \{x \mid x > 3\}$ 是大袋子，$B = \{x \mid x > 5\}$ 是小袋子。$B \subseteq A$。A是大袋子，大对小：必要。）

\textbf{（3）}

\begin{quote}
条件A：$x$ 是偶数

条件B：$x$ 能被2整除
\end{quote}

A → B ✓（偶数就是能被2整除的数。）

B → A ✓（能被2整除的数就是偶数。）

\textbf{答案：A和B互为充要条件。}

（它们就是同一件事的两种说法。$A = B$。）

\textbf{（4）}

\begin{quote}
条件A：$x$ 是奇数

条件B：$x > 100$
\end{quote}

A → B：奇数一定大于100？$x = 3$ 就是奇数但不大于100。 ✗

B → A：大于100一定是奇数？$x = 102$ 大于100但是偶数。 ✗

\textbf{答案：既不充分也不必要。}

（两个袋子有交叉——比如101——但谁也不包含谁。）

\textbf{（5）}

\begin{quote}
条件A：四边形的四条边都相等。

条件B：四边形是正方形。
\end{quote}

A → B：四条边相等就是正方形？

不一定。菱形四条边也相等，但角不一定是直角，所以不一定是正方形。 ✗

B → A：正方形四条边一定都相等？

一定。 ✓

\textbf{答案：A是B的必要不充分条件。}

（想当正方形，四条边相等是必须的——但光这一条不够，还得四个角都是直角。）

\textbf{（6）}

\begin{quote}
条件A：今天是星期天。

条件B：今天不用上学。
\end{quote}

A → B：星期天一定不用上学。 ✓

B → A：不用上学就一定是星期天？不一定——可能是暑假，可能是国庆，可能是生病请假了。 ✗

\textbf{答案：A是B的充分不必要条件。}

（是星期天就够了。但``不上学''不一定是因为星期天。）

\section{这一章你到底学了什么？}

三样东西：

\textbf{第一：命题。}

\begin{quote}
能判断真假的句子。对的叫真命题，错的叫假命题。问句、命令、主观感受不是命题。
\end{quote}

\textbf{第二：``如果……那么……''的结构。}

\begin{quote}
``如果''后面是条件，``那么''后面是结论。整句话在问：从条件出发，能不能走到结论？
\end{quote}

\textbf{第三：充分条件和必要条件。}

\begin{quote}
充分 = 有它就够了。必要 = 没它不行。充要 = 既够了又少不了。

永远分开检查两个方向。A→B？B→A？

小袋子对大袋子充分，大袋子对小袋子必要。
\end{quote}

这三样东西，看起来是``新概念''，但其实你一直在用。

你每天都在做判断——``下雨了，带伞''就是``如果A那么B''。

你每天都在分辨``够不够''和``必不必须''。

数学只是把这些日常的思维，\textbf{用精确的语言固定下来}。

从此以后你说话可以说得更\textbf{准}。

你不会再把``有了A就一定有B''和``有了B就一定有A''搞混了。

这个能力不只是用在数学考试上——它在生活里、在工作中、在任何需要清晰思维的地方，都是一把好刀。

下一章，我们要把命题\textbf{翻个面、转个身}——

如果``如果A那么B''是对的，那么把条件和结论反过来呢？取反呢？又反过来又取反呢？

听起来像绕口令，但你前面刚学过取反，这一章又学了``如果那么''的结构。

两样东西碰在一起，就是下一章的内容。

别怕。你已经有足够的工具了。

\chapter{数学的语文多少还是有特殊的}

上一章我们学了命题。

你知道了：命题就是一句能判断真假的话。

你还学了``如果A那么B''的结构——条件在前，结论在后。

你甚至学了充分条件和必要条件——够了还是必须的，两个方向分开看。

那一章结束的时候我说了一句话：

\begin{quote}
``如果把条件和结论反过来呢？取反呢？又反过来又取反呢？''
\end{quote}

这一章，我们就来干这件事。

\section{一句话能变出几句话？}

我先给你一句话：

\begin{quote}
``如果一个人是北京人，那么他是中国人。''
\end{quote}

这句话你很熟了。上一章反复用它。

现在我问你：从这句话出发，你能``变''出几句新的话？

\textbf{第一种变法：把条件和结论反过来。}

原来是``如果北京人，那么中国人''。

反过来就是：

\begin{quote}
``如果一个人是中国人，那么他是北京人。''
\end{quote}

条件和结论换了个位置。仅此而已。

\textbf{第二种变法：把条件和结论都取反。}

原来的条件是``是北京人''，取反就是``不是北京人''。

原来的结论是``是中国人''，取反就是``不是中国人''。

放在一起：

\begin{quote}
``如果一个人不是北京人，那么他不是中国人。''
\end{quote}

\textbf{第三种变法：又反过来，又取反。}

先把条件和结论反过来（中国人、北京人），然后再取反（不是中国人、不是北京人）。

\begin{quote}
``如果一个人不是中国人，那么他不是北京人。''
\end{quote}

好。现在你手里有四句话了。

我把它们排在一起，你感受一下：

\begin{center}
\begin{tabular}{ll}
\toprule
 & 句子 \\
\midrule
原句 & 如果是北京人，那么是中国人。 \\
变法一 & 如果是中国人，那么是北京人。 \\
变法二 & 如果不是北京人，那么不是中国人。 \\
变法三 & 如果不是中国人，那么不是北京人。 \\
\bottomrule
\end{tabular}
\end{center}

四句话。

同一个话题，翻来覆去说了四遍，每次稍微变了一下。

现在来看关键问题：

\textbf{这四句话，哪些是对的，哪些是错的？}

\textbf{原句：如果是北京人，那么是中国人。}

对不对？

对。北京人一定是中国人。

\textbf{真命题。} ✓

\textbf{变法一：如果是中国人，那么是北京人。}

对不对？

不对。中国人不一定是北京人。可以是上海人、成都人、广州人……

\textbf{假命题。} ✗

\textbf{变法二：如果不是北京人，那么不是中国人。}

对不对？

不对。不是北京人，可以是上海人——上海人也是中国人啊。

\textbf{假命题。} ✗

\textbf{变法三：如果不是中国人，那么不是北京人。}

对不对？

对。不是中国人的话，怎么可能是北京人？北京在中国。

\textbf{真命题。} ✓

你发现了一个很有趣的现象：

\begin{quote}
\textbf{原句是真的，变法三也是真的。}

\textbf{变法一是假的，变法二也是假的。}
\end{quote}

这是巧合吗？

不是。

这是规律。

而且是\textbf{永远成立}的规律。

但在告诉你为什么之前，我得先给这四句话起个正式名字。

\section{四种命题的名字}

数学家给这四种变法各起了一个名字。

名字本身不难，但很容易混。所以我一个一个慢慢说。

\textbf{原命题：就是原来那句话。}

\begin{quote}
``如果A，那么B。''
\end{quote}

原封不动，什么都没变。

用我们的例子：

\begin{quote}
``如果是北京人，那么是中国人。''
\end{quote}

\textbf{逆命题：把条件和结论反过来。}

\begin{quote}
``如果B，那么A。''
\end{quote}

``逆''就是``反过来''的意思。条件变结论，结论变条件。

用我们的例子：

\begin{quote}
``如果是中国人，那么是北京人。''
\end{quote}

你可以这样记：\textbf{逆 = 反过来。}

就像你走路——原来从家走到学校，现在从学校走回家。路线反了。

\textbf{否命题：把条件和结论都取反，但不换位置。}

\begin{quote}
``如果不A，那么不B。''
\end{quote}

``否''就是``否定''、``取反''的意思。条件取反，结论也取反，但\textbf{位置不动}。

用我们的例子：

\begin{quote}
``如果不是北京人，那么不是中国人。''
\end{quote}

你可以这样记：\textbf{否 = 取反。}

位置没变，但两样东西都加了一个``不''字。

\textbf{逆否命题：又反过来，又取反。}

\begin{quote}
``如果不B，那么不A。''
\end{quote}

先``逆''（反过来），再``否''（取反）。两步都做了。

用我们的例子：

\begin{quote}
``如果不是中国人，那么不是北京人。''
\end{quote}

你可以这样记：\textbf{逆否 = 先反过来，再取反。}

我把四个放在一起，你对照着看：

\begin{center}
\begin{tabular}{lll}
\toprule
名字 & 结构 & 例子 \\
\midrule
原命题 & 如果A，那么B & 如果是北京人，那么是中国人 \\
逆命题 & 如果B，那么A & 如果是中国人，那么是北京人 \\
否命题 & 如果不A，那么不B & 如果不是北京人，那么不是中国人 \\
逆否命题 & 如果不B，那么不A & 如果不是中国人，那么不是北京人 \\
\bottomrule
\end{tabular}
\end{center}

请你盯着这张表看十秒钟。

看清楚每一行的结构。

特别注意\textbf{否命题和逆否命题的区别}——

否命题是``不A → 不B''（位置没反，只是取反了）。

逆否命题是``不B → 不A''（位置反了，也取反了）。

很多人把这两个搞混。如果你现在能分清，那真的很不错。

\section{哪些是对的，哪些是错的？}

好，名字认识了。回到刚才那个关键发现：

\begin{quote}
原命题是\textbf{真的}，逆否命题也是\textbf{真的}。

逆命题是\textbf{假的}，否命题也是\textbf{假的}。
\end{quote}

这是规律吗？

是。

\textbf{铁律：原命题和逆否命题，永远同真同假。}

原命题是真的 → 逆否命题一定是真的。

原命题是假的 → 逆否命题一定是假的。

反过来也一样：逆否命题真 → 原命题真。逆否命题假 → 原命题假。

它俩是\textbf{绑在一起}的。一个真，另一个一定真。一个假，另一个一定假。

\textbf{另一条铁律：逆命题和否命题，永远同真同假。}

逆命题真 → 否命题真。

逆命题假 → 否命题假。

它俩也是绑在一起的。

但是！

\textbf{原命题和逆命题之间，没有必然关系。}

原命题真，逆命题可能真，也可能假。

原命题假，逆命题可能真，也可能假。

同样的：

\textbf{原命题和否命题之间，也没有必然关系。}

这一下子冒出来好几条规律，可能有点多。

别急。我用文字帮你在脑子里``搭''一张图。

想象一个正方形。四个角分别是：

\begin{quote}
原命题 ————————— 逆命题

\hspace{1.5em}| \hspace{5.8cm} |

\hspace{1.5em}| \hspace{5.8cm} |

\hspace{1.5em}| \hspace{5.8cm} |

逆否命题 ————————— 否命题
\end{quote}

\textbf{竖着的两对}（上下连线）：

\begin{quote}
原命题和逆否命题（左边那条竖线）：\textbf{同真同假。铁律。}

逆命题和否命题（右边那条竖线）：\textbf{同真同假。铁律。}
\end{quote}

\textbf{横着的两对}（左右连线）：

\begin{quote}
原命题和逆命题（上面那条横线）：\textbf{不一定。没有必然关系。}

逆否命题和否命题（下面那条横线）：\textbf{不一定。没有必然关系。}
\end{quote}

\textbf{交叉的两对}（对角线）：

\begin{quote}
原命题和否命题：\textbf{不一定。}

逆命题和逆否命题：\textbf{不一定。}
\end{quote}

总结起来就一句话：

\begin{quote}
\textbf{只有``竖着的一对''是绑定的。原↔逆否，逆↔否。其他组合都不一定。}
\end{quote}

\section{为什么原命题和逆否命题同真同假？}

如果你只想记住结论，上面那些就够了。

但如果你想知道\textbf{为什么}——往下看。

不想看的话，可以跳到下一节，先记住结论也完全可以，以后再回来看。

还在？好，我慢慢说。

\begin{quote}
原命题：``如果是北京人，那么是中国人。''
\end{quote}

这句话的意思是：所有北京人，都在``中国人''这个大袋子里。

（还记得集合吗？``北京人''是小袋子，``中国人''是大袋子。小袋子完全装在大袋子里面。）

\begin{quote}
逆否命题：``如果不是中国人，那么不是北京人。''
\end{quote}

这句话的意思是：不在``中国人''这个大袋子里的人，也不在``北京人''这个小袋子里。

你想一想——这不是\textbf{同一件事}吗？

如果小袋子完全在大袋子里面，那你跑到大袋子外面去了，你还有可能在小袋子里面吗？

不可能。

小袋子在大袋子里面，你如果连大袋子都不在，那小袋子你更不可能在。

所以``如果在小袋子里就一定在大袋子里''和``如果不在大袋子里就一定不在小袋子里''——说的是\textbf{完全一样的事情}，只是换了个角度。

这就是为什么它们\textbf{永远同真同假}。

它们本来就是\textbf{同一句话的两种说法}。

\section{这个有什么用？}

你可能想问：知道四种命题有什么用？考试会考什么？

用处大了。

最直接的用处就是：\textbf{当你直接证明一件事很难的时候，你可以去证明它的逆否命题。}

因为它们同真同假——证明了逆否命题是真的，就等于证明了原命题是真的。

举个例子：

\begin{quote}
原命题：``如果 $x^2 \neq 4$，那么 $x \neq 2$。''
\end{quote}

你能直接证明吗？能，但要分好几种情况讨论，有点烦。

那我们试试它的逆否命题。

先把条件和结论都取反，再交换位置：

原命题的条件是``$x^2 \neq 4$''，取反就是``$x^2 = 4$''。

（还记得第二章学的取反规则吗？$\neq$ 的反面就是 $=$。）

原命题的结论是``$x \neq 2$''，取反就是``$x = 2$''。

逆否命题：

\begin{quote}
``如果 $x = 2$，那么 $x^2 = 4$。''
\end{quote}

这还用证明吗？$2^2 = 4$，直接算就行了。

逆否命题是真的 → 原命题也是真的。 ✓

搞定。

比直接证明原命题简单多了。

这就是逆否命题的威力：

\begin{quote}
\textbf{正面难走，就绕到背面去。背面通了，正面也就通了。}
\end{quote}

\section{取反时候的细节——别忘了第二章学的}

这里有一个容易犯错的地方，我必须提醒你。

在写否命题或逆否命题的时候，你要对条件和结论``取反''。

取反的规则，在第二章已经学过了：

\begin{quote}
$>$ 的反面是 $\leq$（不是 $<$！）

$<$ 的反面是 $\geq$（不是 $>$！）

$=$ 的反面是 $\neq$

$\neq$ 的反面是 $=$
\end{quote}

（交叉配对。$>$ 和 $\leq$ 是一对，$<$ 和 $\geq$ 是一对。等号要翻转。）

还有：

\begin{quote}
``并且''的反面是``或者''

``或者''的反面是``并且''
\end{quote}

如果你在写逆否命题的时候发现要取反一个复合条件，别忘了这条规则。

比如：

\begin{quote}
原命题：``如果 $x > 0$ 并且 $x < 5$，那么 $x^2 < 25$。''
\end{quote}

写逆否命题：

条件取反：``$x > 0$ 并且 $x < 5$''的反面是什么？

第一步：每个小条件分别取反。$x > 0$ 取反 → $x \leq 0$。$x < 5$ 取反 → $x \geq 5$。

第二步：``并且''变``或者''。

结果：``$x \leq 0$ \textbf{或者} $x \geq 5$''。

结论取反：``$x^2 < 25$''的反面 → $x^2 \geq 25$。

逆否命题（条件结论交换位置）：

\begin{quote}
``如果 $x^2 \geq 25$，那么 $x \leq 0$ 或者 $x \geq 5$。''
\end{quote}

你看，第二章学的取反规则，在这里直接用上了。

工具都是连着的。前面学的东西，后面一定会用到。

\section{来，你来变}

我给你一个原命题，你来写出它的逆命题、否命题、逆否命题，并判断每一个的真假。

别着急看答案。一个一个来。

\textbf{原命题：}

\begin{quote}
``如果一个数能被6整除，那么它能被3整除。''
\end{quote}

这句话对不对？

对。能被6整除的数，一定能被3整除（因为 $6 = 2 \times 3$，包含了因子3）。

\textbf{原命题：真。} ✓

\textbf{（1）写出逆命题。}

逆命题就是把条件和结论反过来。

原命题是``如果能被6整除，那么能被3整除''。

反过来：

\begin{quote}
``如果一个数能被3整除，那么它能被6整除。''
\end{quote}

对不对？

不对。$9$ 能被3整除（$9 \div 3 = 3$），但 $9$ 能被6整除吗？$9 \div 6 = 1.5$，不能。

\textbf{逆命题：假。} ✗

\textbf{（2）写出否命题。}

否命题就是把条件和结论都取反，位置不换。

``能被6整除''取反 → ``不能被6整除''。

``能被3整除''取反 → ``不能被3整除''。

\begin{quote}
``如果一个数不能被6整除，那么它不能被3整除。''
\end{quote}

对不对？

不对。还是用 $9$——$9$ 不能被6整除，但 $9$ 能被3整除。

\textbf{否命题：假。} ✗

\textbf{（3）写出逆否命题。}

逆否命题就是先反过来，再取反。也就是``如果不B，那么不A''。

\begin{quote}
``如果一个数不能被3整除，那么它不能被6整除。''
\end{quote}

对不对？

对。一个数如果连3都整除不了，那它肯定也整除不了6。

（因为6里面包含因子3。你连3都没有，6就更没戏。）

\textbf{逆否命题：真。} ✓

验证一下规律：

\begin{center}
\begin{tabular}{cc}
\toprule
命题 & 真假 \\
\midrule
原命题 & 真 ✓ \\
逆命题 & 假 ✗ \\
否命题 & 假 ✗ \\
逆否命题 & 真 ✓ \\
\bottomrule
\end{tabular}
\end{center}

原命题和逆否命题——\textbf{同真}。 ✓

逆命题和否命题——\textbf{同假}。 ✓

规律验证通过。

\textbf{再来一个，这次你全程自己做。}

\textbf{原命题：}

\begin{quote}
``如果 $x > 3$，那么 $x > 0$。''
\end{quote}

请你写出：

（1）逆命题：\_\_\_\_\_\_

（2）否命题：\_\_\_\_\_\_

（3）逆否命题：\_\_\_\_\_\_

并判断每一个的真假。

\textbf{答案：}

\textbf{原命题：} ``如果 $x > 3$，那么 $x > 0$。''

对不对？对。大于3的数一定大于0。\textbf{真。} ✓

\textbf{（1）逆命题：} ``如果 $x > 0$，那么 $x > 3$。''

对不对？不对。$x = 1$ 就大于0但不大于3。\textbf{假。} ✗

\textbf{（2）否命题：} ``如果 $x \leq 3$，那么 $x \leq 0$。''

注意！取反的时候，$> 3$ 变成 $\leq 3$（不是 $< 3$！），$> 0$ 变成 $\leq 0$（不是 $< 0$！）。

（第二章的交叉配对规则。$>$ 的反面是 $\leq$。）

对不对？不对。$x = 2$ 就满足 $x \leq 3$，但 $2$ 不满足 $x \leq 0$。\textbf{假。} ✗

\textbf{（3）逆否命题：} ``如果 $x \leq 0$，那么 $x \leq 3$。''

对不对？对。小于等于0的数，一定小于等于3。\textbf{真。} ✓

\begin{center}
\begin{tabular}{cc}
\toprule
命题 & 真假 \\
\midrule
原命题 & 真 ✓ \\
逆命题 & 假 ✗ \\
否命题 & 假 ✗ \\
逆否命题 & 真 ✓ \\
\bottomrule
\end{tabular}
\end{center}

又一次完美验证：原↔逆否同真，逆↔否同假。 ✓

\textbf{最后一个，我想让你看看``原命题是假的''会怎么样。}

\textbf{原命题：}

\begin{quote}
``如果一个数是偶数，那么它能被4整除。''
\end{quote}

对不对？

不对。$6$ 是偶数，但 $6 \div 4 = 1.5$，不能被4整除。

\textbf{原命题：假。} ✗

\textbf{逆命题：} ``如果一个数能被4整除，那么它是偶数。''

对不对？

对。能被4整除的数一定是偶数（$4 = 2 \times 2$，里面包含了因子2）。

\textbf{逆命题：真。} ✓

\textbf{否命题：} ``如果一个数不是偶数，那么它不能被4整除。''

对不对？

对。不是偶数就是奇数，奇数肯定不能被4整除。

\textbf{否命题：真。} ✓

\textbf{逆否命题：} ``如果一个数能被4整除，那么它不是偶数。''

等一下——我再想想。

逆否命题是``如果不B，那么不A''。

原命题是``如果是偶数（A），那么能被4整除（B）''。

逆否命题应该是``如果\textbf{不能}被4整除（不B），那么\textbf{不是}偶数（不A）''。

\begin{quote}
``如果一个数不能被4整除，那么它不是偶数。''
\end{quote}

对不对？

不对。$6$ 不能被4整除，但 $6$ 是偶数。

\textbf{逆否命题：假。} ✗

\begin{center}
\begin{tabular}{cc}
\toprule
命题 & 真假 \\
\midrule
原命题 & 假 ✗ \\
逆命题 & 真 ✓ \\
否命题 & 真 ✓ \\
逆否命题 & 假 ✗ \\
\bottomrule
\end{tabular}
\end{center}

规律依然成立：

原命题和逆否命题——\textbf{同假}。 ✓

逆命题和否命题——\textbf{同真}。 ✓

而且你注意到了吗？

这一次，原命题是假的，但逆命题是真的。

这再次证明了：

\begin{quote}
\textbf{原命题的真假，和逆命题的真假，没有必然关系。}
\end{quote}

原命题真，逆命题可能真也可能假。

原命题假，逆命题也可能真也可能假。

它俩不绑定。

\textbf{只有原↔逆否、逆↔否是绑定的。}

\section{一句话总结这一章}

数学的``语文''，确实有点特殊。

日常生活里，你说``如果A那么B''，别人可能顺嘴就说``那如果B呢？是不是也A？''

很多时候大家也不较真。

但数学较真。

数学要你分清楚：

\begin{quote}
\textbf{原命题真 ≠ 逆命题真。}

\textbf{方向不同，结论可能完全不同。}
\end{quote}

但数学也给了你一个礼物：

\begin{quote}
\textbf{原命题真 = 逆否命题真。}

\textbf{正面难证，就走背面。}
\end{quote}

这就是数学语文的``特殊''之处：

\textbf{它比日常语文更严格，但也因此更可靠。}

你说出去的每一句话，是对是错，清清楚楚。没有模棱两可，没有含糊其辞。

你现在手里有四把钥匙：原命题、逆命题、否命题、逆否命题。

你知道它们之间的关系——谁和谁绑定，谁和谁无关。

你还知道一个很实用的技巧——正面难走就走背面（逆否命题）。

这些东西，以后在证明题里会反复出现。

但现在，你已经准备好了。

下一章见。

\chapter{我想知道你会不会判断，还有，证明}

前几章你学了不少东西。

集合、情况、命题、逆命题、否命题、逆否命题……

名字一大堆。

但我现在想问你一个很直接的问题：

\textbf{你真的会用了吗？}

不是背下来了——是\textbf{会用了}。

这一章，我们来检验一下。

而且不只是检验——我还要教你一个新技能：

\textbf{证明。}

\section{别被``证明''吓到}

一说``证明''，很多人就紧张。

觉得证明一定是那种密密麻麻写了半页纸、全是符号、看得人头晕的东西。

不是的。

证明的本质很简单：

\begin{quote}
\textbf{给一个理由，让别人相信你说的是对的。}
\end{quote}

就这样。

你每天都在``证明''。

你妈问你：``你作业写完了吗？''

你说：``写完了，你看，本子在这儿。''

把本子拿出来——这就是证明。你给了\textbf{证据}，让她相信你的话。

数学的证明也是一样的。只不过你给的不是本子，是\textbf{推理}。

一步一步地说清楚：我为什么觉得这句话是对的。

而且，在这一章里，你的证明\textbf{不需要写成数学家那种格式}。

你用自己的话说就行。

只要你说得\textbf{有道理}，逻辑是\textbf{通的}，就算证明。

\section{先从``判断真假''开始}

在证明之前，你总得先判断一句话\textbf{是真是假}。

判断对了，才知道接下来是要``证明它是对的''还是``说明它是错的''。

来。我给你几句话。

你来判断每一句是真命题还是假命题。

先自己想，然后再看答案。

\textbf{（1）} ``所有的质数都是奇数。''

提示：质数是什么来着？只能被1和自己整除的数。你能想到一个\textbf{不是奇数的质数}吗？从最小的质数开始想想……

你的判断：\_\_\_\_\_\_

答案：\textbf{假命题}。2是质数（只能被1和2整除），但2是偶数，不是奇数。一个反例就够了。

\textbf{（2）} ``如果一个整数的末尾是0，那么它能被5整除。''

提示：末尾是0的数……10, 20, 30, 40……你把它们除以5试试？有没有哪个不行的？

你的判断：\_\_\_\_\_\_

答案：\textbf{真命题}。末尾是0的整数，一定能被5整除。10÷5=2，20÷5=4，30÷5=6……无一例外。

\textbf{（3）} ``如果 $a > b$，那么 $a^2 > b^2$。''

提示：这句话里没说 $a$ 和 $b$ 是正数还是负数哦……如果 $a$ 和 $b$ 都是负数呢？试试 $a = -1$，$b = -2$……

你的判断：\_\_\_\_\_\_

答案：\textbf{假命题}。$a = -1$，$b = -2$。$a > b$（$-1 > -2$，对的）。但 $a^2 = 1$，$b^2 = 4$，$1 > 4$ 吗？不。所以 $a^2 > b^2$ 不成立。

\textbf{（4）} ``两个偶数的和还是偶数。''

提示：偶数就是能被2整除的数。你随便拿两个偶数加一加……2+4，6+8，10+12……你能找到一次加出来不是偶数的吗？

你的判断：\_\_\_\_\_\_

答案：\textbf{真命题}。两个偶数的和永远是偶数。（为什么？这个我们等一下来证明。）

\textbf{（5）} ``如果一个四边形的四个角都是直角，那么它是正方形。''

提示：四个角都是直角的四边形……你能想到一种不是正方形的吗？正方形要求的可不止是四个直角……

你的判断：\_\_\_\_\_\_

答案：\textbf{假命题}。长方形的四个角也都是直角，但长方形不一定是正方形（长和宽可以不相等）。

\section{判断完了，然后呢？}

你已经判断了五句话的真假。

但如果这是考试题，光写``真''或``假''通常不够。

老师会在后面加一句：

\begin{quote}
\textbf{``请说明理由。''}
\end{quote}

或者：

\begin{quote}
\textbf{``请证明。''}
\end{quote}

这就到了本章的核心了。

\section{证明一句话是假的——特别简单}

先从简单的开始。

\textbf{如果一句话是假的，你怎么让别人信？}

找一个反例。

就一个就够了。

什么叫反例？

就是一个\textbf{满足条件但不满足结论}的具体例子。

回看第（1）题：

\begin{quote}
``所有的质数都是奇数。''
\end{quote}

你要证明这句话是假的。

你只需要说：

\begin{quote}
``2是质数，但2不是奇数。所以这句话是假的。''
\end{quote}

完了。

这就是证明。

你找到了一个质数（满足条件），但它不是奇数（不满足结论）。

一个反例，整句话就倒了。

再看第（3）题：

\begin{quote}
``如果 $a > b$，那么 $a^2 > b^2$。''
\end{quote}

你要证明它是假的。

\begin{quote}
``取 $a = -1$，$b = -2$。$a > b$（$-1$ 确实大于 $-2$），但 $a^2 = 1$，$b^2 = 4$，$1$ 不大于 $4$。所以这句话是假的。''
\end{quote}

一个具体的例子，条件满足了，结论不满足。

假命题证伪成功。

再看第（5）题：

\begin{quote}
``如果一个四边形的四个角都是直角，那么它是正方形。''
\end{quote}

\begin{quote}
``长方形的四个角都是直角，但长方形不一定是正方形（比如长是5、宽是3的长方形）。所以这句话是假的。''
\end{quote}

搞定。

你发现了吗？

\textbf{证明一句话是``错的''，永远只需要一个反例。}

不需要长篇大论，不需要复杂推理。

找到一个``条件满足了，但结论不成立''的具体例子，往那儿一放。

这句话就倒了。

\textbf{一票否决。}

\section{证明一句话是对的——需要多走几步}

证明假的容易——一个反例搞定。

但证明\textbf{真的}就不一样了。

你不能说``我试了10个例子，每个都对，所以这句话是真的''。

为什么不行？

因为你怎么知道第11个不会出问题？

你试了10个偶数加起来都是偶数，万一有一对你没试到的偶数，加起来是奇数呢？

（其实不会，但你得\textbf{说清楚为什么不会}。这才是证明。）

\textbf{证明一句话是真的，你必须给一个``对所有情况都成立''的理由。}

不能只举几个例子。

要说清楚\textbf{为什么每一种情况都行}。

听起来很难？

其实不难。

我来带你走一遍。

\subsection{证明：两个偶数的和还是偶数}

第（4）题。我们判断了它是真命题。现在来证明它。

你可以用自己的话说。不需要写得像教科书。

但你说的话得\textbf{有道理}，\textbf{逻辑要通}。

我先问你一个问题：

\textbf{偶数到底是什么？}

偶数就是能被2整除的数。

或者换一种说法：偶数就是\textbf{2的倍数}。

2, 4, 6, 8, 10……都是2的倍数。

任何一个偶数，都可以写成 $2 \times \text{某个整数}$ 的形式。

比如：$6 = 2 \times 3$，$10 = 2 \times 5$，$24 = 2 \times 12$。

好。现在假设我有两个偶数。

我不知道它们具体是几——但我知道它们是偶数。

那它们就可以写成：

第一个偶数：$2 \times m$（$m$ 是某个整数）

第二个偶数：$2 \times n$（$n$ 是某个整数）

（为什么用 $m$ 和 $n$？因为这两个偶数可能不一样大，所以它们``除以2之后的那个数''也可能不一样。$m$ 和 $n$ 是两个不同的代号。）

把它们加起来：

\[
2 \times m + 2 \times n
\]

这个式子能化简吗？

能。$2 \times m$ 和 $2 \times n$ 都有一个公共的因子——$2$。

（还记得吗？这就是初中学的``提取公因数''。）

提出来：

\[
2 \times m + 2 \times n = 2 \times (m + n)
\]

好。现在看看 $2 \times (m + n)$ 是什么。

$m$ 是整数，$n$ 是整数，$m + n$ 还是整数。

那 $2 \times (m + n)$ 就是 $2 \times \text{某个整数}$。

而 $2 \times \text{某个整数}$——这不就是\textbf{偶数}吗？

所以两个偶数加起来，结果一定是偶数。 ✓

你看。整个过程用的都是大白话。

但逻辑是通的：

\begin{quote}
偶数可以写成 $2 \times$ 整数。

两个这样的东西加起来，提一个2出来，还是 $2 \times$ 整数。

所以结果还是偶数。
\end{quote}

这就是证明。

没有一个字是废话，也没有一个字是``因为我试了好几个例子所以我觉得对''。

它说的是\textbf{为什么对所有偶数都成立}。

\subsection{证明：末尾是0的整数能被5整除}

第（2）题。再来一个。

末尾是0的整数长什么样？

10, 20, 30, 100, 250, 1000……

它们有什么共同点？

\textbf{它们都是10的倍数。}

一个数末尾是0，就意味着它可以写成 $10 \times \text{某个整数}$ 的形式。

比如：$30 = 10 \times 3$，$250 = 10 \times 25$。

好。那 $10 \times \text{某个整数}$ 能被5整除吗？

$10 = 2 \times 5$。

所以 $10 \times \text{某个整数} = 2 \times 5 \times \text{某个整数}$。

里面有一个因子5。

有因子5，当然能被5整除。

所以末尾是0的整数，一定能被5整除。 ✓

\section{逆否命题——证明的``后门''}

还记得上一章学的逆否命题吗？

\textbf{原命题和逆否命题同真同假。}

这意味着：如果你觉得原命题不好直接证明，你可以\textbf{去证明它的逆否命题}。

逆否命题证明了是真的 → 原命题也是真的。

来试一个：

\begin{quote}
原命题：``如果 $x^2$ 是奇数，那么 $x$ 是奇数。''
\end{quote}

你能直接证明吗？

可以，但要分情况讨论，有点绕。

那我们试试逆否命题。

原命题的条件是``$x^2$ 是奇数''，结论是``$x$ 是奇数''。

逆否命题就是``如果不B，那么不A''：

\begin{quote}
``如果 $x$ 不是奇数，那么 $x^2$ 不是奇数。''
\end{quote}

``不是奇数''就是``是偶数''。所以换成人话：

\begin{quote}
``如果 $x$ 是偶数，那么 $x^2$ 是偶数。''
\end{quote}

这个好证多了！

$x$ 是偶数，就是 $x = 2m$（$m$ 是某个整数）。

$x^2 = (2m)^2 = 4m^2 = 2 \times (2m^2)$。

$2 \times (2m^2)$——这是 $2 \times$ 整数，偶数。

所以 $x^2$ 是偶数。 ✓

逆否命题证完了 → 原命题也是真的。 ✓

你看到了吗？

正面有点绕，走后门一下就通了。

\textbf{逆否命题是你证明问题时的一条捷径。}

正门堵了就走后门。反正它俩通向同一个地方。

\section{现在到你了}

下面几道题，有的要你判断真假，有的要你证明，有的要你找反例。

\textbf{请用你自己的话来回答。} 不需要写得多正式，但要说清楚\textbf{为什么}。

``因为……所以……''就够了。

\textbf{（1）} 判断真假，如果是假的，找一个反例。

\begin{quote}
``如果 $a + b > 0$，那么 $a > 0$ 并且 $b > 0$。''
\end{quote}

提示：$a + b > 0$ 只要求它们的\textbf{和}大于0。有没有可能一个是正数、一个是负数，但加起来还是正的？试试具体的数……

你的判断：\_\_\_\_\_\_

你的理由：\_\_\_\_\_\_

答案：\textbf{假命题}。反例：$a = 5$，$b = -2$。$a + b = 3 > 0$（条件满足），但 $b = -2 < 0$（$b$ 不大于0，结论不满足）。

\textbf{（2）} 判断真假，如果是真的，试着说清楚为什么。

\begin{quote}
``两个奇数的和是偶数。''
\end{quote}

提示：刚才我们证明了``两个偶数的和是偶数''。方法是把偶数写成 $2m$ 的形式。那奇数可以写成什么形式呢？奇数比偶数多1，所以……

你的判断：\_\_\_\_\_\_

你的理由：\_\_\_\_\_\_

答案：\textbf{真命题}。

理由：奇数可以写成 $2m + 1$ 的形式（比偶数多1）。

两个奇数相加：$(2m + 1) + (2n + 1) = 2m + 2n + 2 = 2(m + n + 1)$。

$2 \times (m + n + 1)$ 是 $2 \times$ 整数，所以是偶数。 ✓

\textbf{（3）} 判断真假，如果是假的，找反例。

\begin{quote}
``如果 $a^2 = b^2$，那么 $a = b$。''
\end{quote}

提示：$a^2 = b^2$ 的意思是两个数的平方一样。平方一样，原来的数就一定一样吗？想想有没有两个不同的数，平方之后变成一样的……

你的判断：\_\_\_\_\_\_

你的理由：\_\_\_\_\_\_

答案：\textbf{假命题}。反例：$a = 3$，$b = -3$。$a^2 = 9$，$b^2 = 9$，$a^2 = b^2$（条件满足）。但 $a = 3 \neq -3 = b$（结论不满足）。

\textbf{（4）} 这道稍微有点挑战。判断真假，并说明理由。

\begin{quote}
``如果 $n$ 是整数，那么 $n^2 + n$ 一定是偶数。''
\end{quote}

提示：$n^2 + n$ 可以写成 $n \times (n + 1)$。（你可以用提取公因数验证一下。）$n$ 和 $n + 1$ 是两个\textbf{相邻的整数}。相邻的两个整数里面，有没有什么规律？一个是奇数另一个是……

你的判断：\_\_\_\_\_\_

你的理由：\_\_\_\_\_\_

答案：\textbf{真命题}。

理由：$n^2 + n = n(n + 1)$。$n$ 和 $n + 1$ 是相邻的两个整数。相邻的两个整数，一定是一个奇数一个偶数。（如果 $n$ 是偶数，$n+1$ 就是奇数；如果 $n$ 是奇数，$n+1$ 就是偶数。）所以 $n(n+1)$ 里面一定有一个偶数，也就是说一定有一个因子2。有因子2，就是偶数。 ✓

\textbf{（5）} 试着用逆否命题来证明：

\begin{quote}
``如果 $n^2$ 是偶数，那么 $n$ 是偶数。''
\end{quote}

提示：你在前面已经证明过一个非常像的东西了。逆否命题是什么来着？``如果 $n$ 不是偶数……''

你的逆否命题：\_\_\_\_\_\_

你的证明：\_\_\_\_\_\_

答案：

逆否命题：``如果 $n$ 不是偶数（$n$ 是奇数），那么 $n^2$ 不是偶数（$n^2$ 是奇数）。''

证明：$n$ 是奇数，所以 $n = 2m + 1$。$n^2 = (2m+1)^2 = 4m^2 + 4m + 1 = 2(2m^2 + 2m) + 1$。这是 $2 \times$ 整数 $+ 1$，所以是奇数。 ✓

逆否命题是真的 → 原命题也是真的。 ✓

\section{这一章的灵魂}

证明没有你想的那么可怕。

它就是三件事：

\begin{quote}
\textbf{第一：判断这句话是真是假。}

\textbf{第二：如果是假的——找一个反例。一个就够，一票否决。}

\textbf{第三：如果是真的——说清楚为什么。不是举几个例子，是讲道理，让人心服口服。}
\end{quote}

讲道理的方法有很多种。

你可以正面推——从条件一步步推出结论。

你也可以走后门——证明逆否命题。

用你自己的话说，只要逻辑通，就算数。

到这里，你已经拥有了一套完整的数学``语文能力''。

你会描述（集合）、会审条件（情况）、会判断真假（命题）、会翻面变形（四种命题）、还会\textbf{给出理由}（证明）。

这些能力加在一起，就是高中数学的底层操作系统。

后面不管学什么——函数、不等式、数列、概率——全都跑在这个操作系统上面。

而你的操作系统，已经装好了。

准备好了吗？我们继续。

\chapter{不等式的基本与严格表示（逐步详解）}

\section{一、先从等式说起——为什么需要不等式？}

等式表达的是``两边\textbf{相等}''的关系，比如：

\begin{quote}
3 + 2 = 5（左边等于右边）

2x = 10（左边等于右边）
\end{quote}

但生活中大量的关系\textbf{不是相等的}：

\begin{quote}
你的身高\textbf{大于}你弟弟的身高

考试成绩要\textbf{不低于}60分才算及格

书包重量\textbf{小于}10公斤才符合规定

一个数\textbf{不等于}0
\end{quote}

这些关系，用等号写不出来。所以我们需要一套新的符号——\textbf{不等号}。

\section{二、五个不等号，逐个认识}

\subsection{符号1：$>$（大于）}

\textbf{含义：} 左边比右边\textbf{大}。

\textbf{怎么读：} ``大于''

\textbf{例子：}

\begin{quote}
5 > 3 → 读作``5大于3''

10 > 2 → 读作``10大于2''
\end{quote}

\textbf{记忆技巧：} 开口（大的那头）朝着\textbf{大的数}，尖头（小的那头）指着\textbf{小的数}。

来看 5 > 3：

\begin{quote}
开口朝左 → 左边的5是大的 ✓

尖头朝右 → 右边的3是小的 ✓
\end{quote}

\subsection{符号2：$<$（小于）}

\textbf{含义：} 左边比右边\textbf{小}。

\textbf{怎么读：} ``小于''

\textbf{例子：}

\begin{quote}
3 < 5 → 读作``3小于5''

1 < 100 → 读作``1小于100''
\end{quote}

\textbf{和``$>$''的关系：} 其实 5 > 3 和 3 < 5 说的是\textbf{同一件事}，只是换了个方向看。

就像``小明比小红高''和``小红比小明矮''是同一个意思。

\subsection{符号3：$\geq$（大于等于，也叫``大于或等于''）}

\textbf{含义：} 左边比右边\textbf{大}，或者左边跟右边\textbf{相等}。两种情况\textbf{只要满足一个}就行。

\textbf{怎么读：} ``大于等于''或``大于或等于''，也有人读``不小于''

\textbf{例子：}

\begin{quote}
5 ≥ 3 → 对吗？对！因为 5 大于 3，满足``大于''这个条件 ✓

5 ≥ 5 → 对吗？对！因为 5 等于 5，满足``等于''这个条件 ✓

3 ≥ 5 → 对吗？错！3 既不大于 5，也不等于 5，两个条件都不满足 ✗
\end{quote}

\textbf{生活中的例子：}

\begin{quote}
``考试成绩不低于60分'' → 用数学写就是：成绩 ≥ 60
\end{quote}

这意味着60分\textbf{也算及格}（等于60可以），61、70、100都行（大于60也可以）

\subsection{符号4：$\leq$（小于等于，也叫``小于或等于''）}

\textbf{含义：} 左边比右边\textbf{小}，或者左边跟右边\textbf{相等}。两种情况\textbf{只要满足一个}就行。

\textbf{怎么读：} ``小于等于''或``小于或等于''，也有人读``不大于''

\textbf{例子：}

\begin{quote}
3 ≤ 5 → 对！3 小于 5 ✓

5 ≤ 5 → 对！5 等于 5 ✓

7 ≤ 5 → 错！7 既不小于 5，也不等于 5 ✗
\end{quote}

\textbf{生活中的例子：}

\begin{quote}
``行李重量不能超过20公斤'' → 重量 ≤ 20
\end{quote}

这意味着正好20公斤\textbf{也可以}（等于20行），15公斤也行（小于20行），但21公斤不行

\subsection{符号5：$\neq$（不等于）}

\textbf{含义：} 左边和右边\textbf{不相等}，但不管谁大谁小，只要不相等就行。

\textbf{怎么读：} ``不等于''

\textbf{例子：}

\begin{quote}
3 ≠ 5 → 对！3 和 5 确实不相等 ✓

5 ≠ 3 → 对！也不相等 ✓

5 ≠ 5 → 错！5 和 5 明明相等 ✗
\end{quote}

\textbf{特点：} $\neq$ 只告诉你``不相等''，但\textbf{不告诉你谁大谁小}。

\section{三、严格不等式 vs 非严格不等式}

这是非常重要的区分！

\subsection{严格不等式（不含等号）}

就是 \textbf{$>$} 和 \textbf{$<$}

\textbf{``严格''在哪？} 严格要求左右两边\textbf{不能相等}，必须一个比另一个\textbf{真的大}（或\textbf{真的小}）。

\textbf{例子：} x > 3

\begin{quote}
x = 4 → 满足吗？4 > 3 ✓

x = 5 → 满足吗？5 > 3 ✓

x = 100 → 满足吗？100 > 3 ✓

x = 3 → 满足吗？3 > 3 ✗ \textbf{不满足！因为3等于3，不是大于3}

x = 2 → 满足吗？2 > 3 ✗
\end{quote}

所以 x > 3 的意思是：x 可以是 3.0001、3.1、4、100……但\textbf{就是不能等于3本身}。

\subsection{非严格不等式（含等号）}

就是 \textbf{$\geq$} 和 \textbf{$\leq$}

\textbf{``非严格''在哪？} 允许两边相等，条件比较\textbf{宽松}。

\textbf{例子：} x ≥ 3

\begin{quote}
x = 4 → 满足吗？4 ≥ 3 ✓

x = 5 → 满足吗？5 ≥ 3 ✓

x = 3 → 满足吗？3 ≥ 3 ✓ \textbf{满足！因为等于也算}

x = 2 → 满足吗？2 ≥ 3 ✗
\end{quote}

所以 x ≥ 3 和 x > 3 的\textbf{唯一差别}就在于：\textbf{x = 3 这一个点}到底算不算。

\subsection{用一张表来对比}

\begin{center}
\begin{tabular}{ccc}
\toprule
 & 严格不等式 & 非严格不等式 \\
\midrule
符号 & $>$ 或 $<$ & $\geq$ 或 $\leq$ \\
能不能取等号 & \textbf{不能} & \textbf{能} \\
x > 3 时，x能等于3吗 & 不能 & — \\
$x \geq 3$ 时，$x$ 能等于 $3$ 吗 & --- & 能 \\
口语说法 & ``大于''``小于'' & ``大于等于''``小于等于''``不小于''``不大于'' \\
\bottomrule
\end{tabular}
\end{center}

\subsection{一个关键理解}

\begin{quote}
\textbf{$\geq$ 就是把 $>$ 和 $=$ 合在一起。}

写 $x \geq 3$，相当于说``$x > 3$ \textbf{或者} $x = 3$''。

\textbf{$\leq$ 就是把 $<$ 和 $=$ 合在一起。}

写 $x \leq 3$，相当于说``$x < 3$ \textbf{或者} $x = 3$''。
\end{quote}

\section{四、不等式长什么样？}

和等式对比着看：

\begin{center}
\begin{tabular}{cc}
\toprule
等式 & 不等式 \\
\midrule
$x + 3 = 7$ & $x + 3 > 7$ \\
$2x = 10$ & $2x < 10$ \\
$x - 1 = 0$ & $x - 1 \geq 0$ \\
$3x + 2 = 5$ & $3x + 2 \leq 5$ \\
\bottomrule
\end{tabular}
\end{center}

你会发现：\textbf{不等式和等式的结构几乎一模一样，唯一的区别就是中间的符号不同。}

等式中间放的是 $=$，不等式中间放的是 $>$、$<$、$\geq$、$\leq$。

\section{五、不等式的方向——一个容易搞混的问题}

\subsection{什么叫``不等式的方向''？}

看不等号的\textbf{开口朝哪边}。

\begin{quote}
$5 > 3$ $\rightarrow$ 开口朝\textbf{左}，我们说这个不等式的方向是``大于方向''

$3 < 5$ $\rightarrow$ 开口朝\textbf{右}，我们说这个不等式的方向是``小于方向''
\end{quote}

\subsection{``同向''和``反向''}

\textbf{同向：} 两个不等式的不等号\textbf{开口朝同一边}

\begin{quote}
例如：$a > b$ 和 $c > d$ $\rightarrow$ 都是 $>$，同向

例如：$a < b$ 和 $c < d$ $\rightarrow$ 都是 $<$，同向
\end{quote}

\textbf{反向：} 两个不等式的不等号\textbf{开口朝相反方向}

\begin{quote}
例如：$a > b$ 和 $c < d$ $\rightarrow$ 一个 $>$ 一个 $<$，反向
\end{quote}

\subsection{为什么要关注方向？}

因为后面学到不等式的性质时，有一条极其重要的规则：

\begin{quote}
\textbf{不等式两边同时乘以（或除以）一个负数时，不等号的方向要反过来！}
\end{quote}

这是不等式和等式\textbf{最大的区别}，后面会详细讲。现在先记住``方向''这个概念。

\section{六、不等式的基本性质}

\subsection{性质1：两边同时加上（或减去）同一个数，不等号方向\textbf{不变}}

\textbf{例子：} 已知 $5 > 3$

两边同时加 $2$：

\begin{quote}
左边：$5 + 2 = 7$

右边：$3 + 2 = 5$

结果：$7 > 5$ ✓ 方向没变，还是 $>$
\end{quote}

两边同时减 $10$：

\begin{quote}
左边：$5 - 10 = -5$

右边：$3 - 10 = -7$

结果：$-5 > -7$ ✓ 方向没变，还是 $>$
\end{quote}

\textbf{结论：} 加减不影响不等号方向，这一点和等式\textbf{完全一样}。

\subsection{性质2：两边同时乘以（或除以）一个\textbf{正数}，不等号方向\textbf{不变}}

\textbf{例子：} 已知 $5 > 3$

两边同时乘以 $2$：

\begin{quote}
左边：$5 \times 2 = 10$

右边：$3 \times 2 = 6$

结果：$10 > 6$ ✓ 方向没变
\end{quote}

两边同时除以正数，也一样：

\begin{quote}
$10 > 6$，两边除以 $2$

$5 > 3$ ✓ 方向没变
\end{quote}

\subsection{性质3：两边同时乘以（或除以）一个\textbf{负数}，不等号方向\textbf{必须反过来！}}

\textbf{这是最最最重要的一条，也是最容易出错的！}

\textbf{例子：} 已知 $5 > 3$

两边同时乘以 $-1$：

\begin{quote}
左边：$5 \times (-1) = -5$

右边：$3 \times (-1) = -3$

现在 $-5$ 和 $-3$ 谁大？ $\rightarrow$ $-3$ 大！（负数越接近 $0$ 越大）

所以结果是：$-5 < -3$

方向从 $>$ 变成了 $<$！\textbf{反过来了！}
\end{quote}

\textbf{再来一个例子：} 已知 $-2x > 6$，求 $x$

第一步：要把 $x$ 前面的 $-2$ 消掉，两边同时除以 $-2$

第二步：但是 $-2$ 是\textbf{负数}！所以不等号方向\textbf{必须反过来}！

\begin{quote}
左边：$-2x \div (-2) = x$

右边：$6 \div (-2) = -3$

不等号从 $>$ 变成 $<$

结果：$x < -3$
\end{quote}

\textbf{如果你忘了变方向，直接写 $x > -3$，那就完全错了！}

\subsection{为什么乘负数要变方向？直觉上怎么理解？}

想象一条数轴：

\begin{quote}
$\cdots\ -5\ \ -4\ \ -3\ \ -2\ \ -1\ \ 0\ \ 1\ \ 2\ \ 3\ \ 4\ \ 5\ \cdots$
\end{quote}

$5$ 在 $3$ 的\textbf{右边}，所以 $5 > 3$。

现在都乘以 $-1$：

\begin{quote}
$5$ 变成 $-5$，跑到了数轴的\textbf{左边}

$3$ 变成 $-3$，也跑到了左边

但是 $-5$ 跑得\textbf{更远}，到了 $-3$ 的\textbf{左边}

在数轴上越靠左越小，所以 $-5 < -3$
\end{quote}

\textbf{本质：乘以负数就像照镜子，左右翻转了，大小关系也就翻转了。}

\subsection{三条性质对比总结}

\begin{center}
\begin{tabular}{cc}
\toprule
操作 & 不等号方向 \\
\midrule
两边加或减同一个数 & \textbf{不变} \\
两边乘或除以一个\textbf{正数} & \textbf{不变} \\
两边乘或除以一个\textbf{负数} & \textbf{必须反过来！} \\
\bottomrule
\end{tabular}
\end{center}

\section{七、用语言翻译成不等式}

这是考试中经常出现的题型——把一句话写成数学表达式。

\subsection{``大于''→ $>$}

\textbf{语句：} ``$x$ 大于 $5$''

\textbf{写法：} $x > 5$

\subsection{``小于''→ $<$}

\textbf{语句：} ``$x$ 小于 $10$''

\textbf{写法：} $x < 10$

\subsection{``不小于''→ $\geq$}

\textbf{语句：} ``$x$ 不小于 $3$''

\textbf{分析：} ``不小于''就是说``不会比 $3$ 小''，那就是``大于 $3$ 或者等于 $3$''

\textbf{写法：} $x \geq 3$

\begin{quote}
\textbf{``不小于'' = 大于等于 = $\geq$}
\end{quote}

\subsection{``不大于''→ $\leq$}

\textbf{语句：} ``$x$ 不大于 $7$''

\textbf{分析：} ``不大于''就是说``不会比 $7$ 大''，那就是``小于 $7$ 或者等于 $7$''

\textbf{写法：} $x \leq 7$

\begin{quote}
\textbf{``不大于'' = 小于等于 = $\leq$}
\end{quote}

\subsection{``至少''→ $\geq$}

\textbf{语句：} ``至少有 $5$ 个苹果''

\textbf{分析：} ``至少 $5$ 个''就是说``最少也有 $5$ 个''，可以是 $5$ 个、$6$ 个、$100$ 个……

\textbf{写法：} 苹果数 $\geq 5$

\subsection{``至多''→ $\leq$}

\textbf{语句：} ``至多花 $20$ 块钱''

\textbf{分析：} ``至多 $20$''就是说``最多也就 $20$''，可以是 $20$、$15$、$3$ ……

\textbf{写法：} 花费 $\leq 20$

\subsection{``不超过''→ $\leq$}

\textbf{语句：} ``体重不超过 $50$ 公斤''

\textbf{写法：} 体重 $\leq 50$

\subsection{``不低于''→ $\geq$}

\textbf{语句：} ``温度不低于零下 $10$ 度''

\textbf{写法：} 温度 $\geq -10$

\subsection{翻译总结表}

\begin{center}
\begin{tabular}{ccc}
\toprule
日常用语 & 数学符号 & 能不能取等号 \\
\midrule
大于 & $>$ & 不能 \\
小于 & $<$ & 不能 \\
大于等于、不小于、至少、不低于 & $\geq$ & \textbf{能} \\
小于等于、不大于、至多、不超过 & $\leq$ & \textbf{能} \\
不等于 & $\neq$ & --- \\
\bottomrule
\end{tabular}
\end{center}

\section{八、双向不等式（连写）}

有时候我们要表达``$x$ 在某个范围内''，可以把两个不等式\textbf{连着写}。

\subsection{例子1：``$x$ 大于 $2$ 并且 小于 $5$''}

分开写：$x > 2$ \textbf{并且} $x < 5$

连着写：$2 < x < 5$

\textbf{含义：} $x$ 夹在 $2$ 和 $5$ 之间，但\textbf{不包含 $2$ 和 $5$ 本身}（因为是严格不等式）

\begin{quote}
$x = 3$ $\rightarrow$ $2 < 3 < 5$ ✓

$x = 4.9$ $\rightarrow$ $2 < 4.9 < 5$ ✓

$x = 2$ $\rightarrow$ $2 < 2$ 不成立 ✗

$x = 5$ $\rightarrow$ $5 < 5$ 不成立 ✗
\end{quote}

\subsection{例子2：``$x$ 大于等于 $1$ 并且 小于等于 $4$''}

连着写：$1 \leq x \leq 4$

\textbf{含义：} $x$ 在 $1$ 到 $4$ 之间，\textbf{包含 $1$ 和 $4$}

\begin{quote}
$x = 1$ $\rightarrow$ $1 \leq 1 \leq 4$ ✓（等于 $1$ 也算）

$x = 4$ $\rightarrow$ $1 \leq 4 \leq 4$ ✓（等于 $4$ 也算）

$x = 2.5$ $\rightarrow$ $1 \leq 2.5 \leq 4$ ✓

$x = 0$ $\rightarrow$ $1 \leq 0$ 不成立 ✗
\end{quote}

\subsection{例子3：也可以一边严格一边不严格}

\textbf{``$x$ 大于等于 $0$ 并且 小于 $1$''}

写法：$0 \leq x < 1$

\begin{quote}
$x = 0$ $\rightarrow$ 可以（左边取等号）✓

$x = 0.5$ $\rightarrow$ 可以 ✓

$x = 1$ $\rightarrow$ 不行（右边是严格小于，不能等于 $1$）✗
\end{quote}

\subsection{连写时的注意事项}

\textbf{规则：连写的不等号方向必须一致！}

\begin{quote}
$2 < x < 5$ ✓（两个都是 $<$，方向一致）

$1 \leq x \leq 4$ ✓（两个都是 $\leq$，方向一致）

$0 \leq x < 1$ ✓（都是``左小右大''的方向）
\end{quote}

\textbf{错误写法：} $5 > x < 3$ ✗（方向不一致，一个 $>$ 一个 $<$，这样写让人困惑）

\section{九、不等式与等式的区别总结}

\begin{center}
\begin{tabular}{ccc}
\toprule
 & 等式 & 不等式 \\
\midrule
符号 & $=$ & $>$、$<$、$\geq$、$\leq$、$\neq$ \\
解的个数 & 通常是\textbf{一个}确定的值 & 通常是\textbf{一个范围}（无数个值） \\
加减同一个数 & 等式仍成立 & 不等号方向不变 \\
乘除\textbf{正数} & 等式仍成立 & 不等号方向不变 \\
乘除\textbf{负数} & 等式仍成立 & \textbf{不等号方向必须反过来！} \\
\bottomrule
\end{tabular}
\end{center}

最后那一行就是\textbf{最核心的区别}——等式不管乘什么都还是等式，但不等式乘负数会``翻转''。

\section{十、练习题（附逐步解答）}

\subsection{练习1：把文字翻译成不等式——``一个数 $x$ 的 $3$ 倍大于 $12$''}

第一步：找关键词

\begin{quote}
``$x$ 的 $3$ 倍'' $\rightarrow$ $3x$

``大于'' $\rightarrow$ $>$

``$12$'' $\rightarrow$ $12$
\end{quote}

第二步：组合起来

\textbf{答案：} $3x > 12$

\subsection{练习2：把文字翻译成不等式——``小明的年龄 $x$ 至少是 $18$ 岁''}

第一步：找关键词

\begin{quote}
``至少'' $\rightarrow$ 意思是``最少也得是''，就是``大于等于'' $\rightarrow$ $\geq$

``$18$ 岁'' $\rightarrow$ $18$
\end{quote}

第二步：组合起来

\textbf{答案：} $x \geq 18$

\subsection{练习3：判断对错——已知 $a > b$，那么 $-2a > -2b$ 对吗？}

第一步：从 $a > b$ 出发

第二步：两边同时乘以 $-2$

第三步：$-2$ 是\textbf{负数}！所以不等号方向必须反过来

第四步：应该得到 $-2a < -2b$

\textbf{答案：} 错！正确的应该是 $-2a < -2b$

\subsection{练习4：已知 $-3x < 9$，求 $x$ 的范围}

第一步：要把 $x$ 前面的 $-3$ 消掉，两边同时除以 $-3$

第二步：$-3$ 是\textbf{负数}！除以负数，不等号方向要\textbf{反过来}！

第三步：执行运算

\begin{quote}
左边：$-3x \div (-3) = x$

右边：$9 \div (-3) = -3$

不等号从 $<$ 变成 $>$
\end{quote}

第四步：结果是 $x > -3$

第五步：验证——随便选一个大于 $-3$ 的数，比如 $x = 0$，代入原式

\begin{quote}
左边：$-3 \times 0 = 0$

$0 < 9$ ✓ 成立
\end{quote}

再选一个不满足的，比如 $x = -5$（$-5$ 不大于 $-3$）

\begin{quote}
左边：$-3 \times (-5) = 15$

$15 < 9$ ✗ 不成立
\end{quote}

验证通过！

\textbf{答案：} $x > -3$

\subsection{练习5：用不等式表示——``$x$ 是一个大于 $0$ 但不超过 $10$ 的数''}

第一步：拆成两个条件

\begin{quote}
``大于 $0$'' $\rightarrow$ $x > 0$（严格大于，不包含 $0$）

``不超过 $10$'' $\rightarrow$ $x \leq 10$（不超过就是``小于等于''，包含 $10$）
\end{quote}

第二步：连写成双向不等式

\textbf{答案：} $0 < x \leq 10$

\subsection{练习6：下列哪些 $x$ 的值满足 $2 \leq x < 5$？}

逐个检查：

\begin{center}
\begin{tabular}{cccc}
\toprule
$x$ 的值 & $2 \leq x$ 成立吗？ & $x < 5$ 成立吗？ & 两个都成立？ \\
\midrule
$x = 1$ & $2 \leq 1$ ✗ & $1 < 5$ ✓ & ✗ 不满足 \\
$x = 2$ & $2 \leq 2$ ✓（等于也算） & $2 < 5$ ✓ & \textbf{✓ 满足} \\
$x = 3$ & $2 \leq 3$ ✓ & $3 < 5$ ✓ & \textbf{✓ 满足} \\
$x = 4.9$ & $2 \leq 4.9$ ✓ & $4.9 < 5$ ✓ & \textbf{✓ 满足} \\
$x = 5$ & $2 \leq 5$ ✓ & $5 < 5$ ✗（等于不算） & ✗ 不满足 \\
$x = 6$ & $2 \leq 6$ ✓ & $6 < 5$ ✗ & ✗ 不满足 \\
\bottomrule
\end{tabular}
\end{center}

\textbf{答案：} 满足的有 $x=2,3,4.9$；不满足的有 $x=1,5,6$。

\textbf{关键发现：左边的 $\leq$ 让 $x=2$ 可以取到，但右边的 $<$ 让 $x=5$ 取不到。这就是严格与非严格的实际差别。}
\chapter{不等式的利用——更难的翻译（应用题实战）}

\section{为什么要学这个？}

上一章我们学会了简单的翻译，比如“$x$ 大于 5”写成 $x > 5$。

但考试和生活中，题目不会直接告诉你“大于”“小于”，而是藏在一段话里，你需要自己从文字中找出不等关系，列出不等式，再求解。

这一章就是练这个能力：读题 $\to$ 找关键词 $\to$ 列不等式 $\to$ 解出来 $\to$ 回答问题。

\section{一、翻译的核心思路（万能三步）}

不管题目多复杂，永远是这三步：

\begin{quote}
第1步：设未知数（不知道的量用 $x$ 表示）

$\downarrow$

第2步：找“不等关系”（题目里哪句话在说大小、多少、够不够？）

$\downarrow$

第3步：把那句话一个词一个词地翻译成数学符号
\end{quote}

\section{二、关键词对照表（翻译字典，随时回来查）}

\begin{center}
\begin{tabular}{lll}
\toprule
题目中的说法 & 翻译成 & 能不能取等号 \\
\midrule
“超过” & $>$ & 不能 \\
“不超过” & $\leq$ & 能 \\
“多于” & $>$ & 不能 \\
“不少于”“不低于” & $\geq$ & 能 \\
“少于”“低于”“不足” & $<$ & 不能 \\
“不多于”“不高于” & $\leq$ & 能 \\
“至少” & $\geq$ & 能 \\
“至多”“最多” & $\leq$ & 能 \\
“足够”“够” & $\geq$ & 能 \\
“不够” & $<$ & 不能 \\
“便宜”“更省钱”“花费更少” & $<$ & 不能 \\
“剩余”“有余”“还有剩” & $>$ & 不能 \\
\bottomrule
\end{tabular}
\end{center}

\section{三、从最简单的开始——一步翻译}

\subsection{例题1：买文具}

\begin{quote}
\textbf{小明有 30 元钱，他想买一些笔，每支笔 4 元。他最多能买几支笔？}
\end{quote}

\textbf{第1步：设未知数}

不知道买几支，设买 $x$ 支

\textbf{第2步：找不等关系}

他只有30元，花的钱不能超过30元

“不能超过”$\to \leq$

\textbf{第3步：翻译}

买 $x$ 支笔要花多少钱？$\to$ 每支4元，$x$ 支就是 $4x$ 元

花的钱不能超过他有的钱 $\to 4x \leq 30$

\begin{quote}
📖 \textbf{这个式子直接用语言说就是：} “买笔花的钱（$4x$元）小于等于手里有的钱（30元）”，也就是“花的钱不能超过30元”。
\end{quote}

\textbf{第4步：解不等式}

两边除以 4（4 是正数，方向不变）

$x \leq 30 \div 4$

$x \leq 7.5$

\textbf{第5步：回答问题}

$x$ 是笔的支数，必须是整数（你不能买半支笔）

$x \leq 7.5$，那最大的整数就是 7

\textbf{答：最多能买 7 支笔。}

\textbf{第6步：检验}

买7支：$4 \times 7 = 28$元，$28 \leq 30$ ✓ 够

买8支：$4 \times 8 = 32$元，$32 \leq 30$ ✗ 不够

确认7支没问题 ✓

\subsection{例题2：考试分数}

\begin{quote}
\textbf{小红前四次考试的成绩分别是 85、90、78、88。她希望五次考试的平均分不低于 85 分，第五次至少要考多少分？}
\end{quote}

\textbf{第1步：设未知数}

第五次考 $x$ 分

\textbf{第2步：找不等关系}

“平均分不低于85”$\to$ 平均分 $\geq 85$

\textbf{第3步：翻译}

五次的平均分怎么算？$\to$ 五次成绩加起来，除以5

五次总分：85 + 90 + 78 + 88 + $x = 341 + x$

平均分 $= (341 + x) / 5$

不等式：$\dfrac{341 + x}{5} \geq 85$

\begin{quote}
📖 \textbf{这个式子直接用语言说就是：} “五次考试的平均分（把五次成绩加起来除以5）大于等于85分”，也就是“平均分不能低于85，等于85也行”。
\end{quote}

\textbf{第4步：解不等式}

第一小步：去分母，两边乘以 5（正数，方向不变）

$341 + x \geq 85 \times 5$

$341 + x \geq 425$

第二小步：移项，把 341 移到右边变 $-341$

$x \geq 425 - 341$

$x \geq 84$

\textbf{第5步：回答问题}

第五次至少要考 84 分

\textbf{第6步：检验}

考84分：平均分 $= (341 + 84) / 5 = 425 / 5 = 85 \geq 85$ ✓

考83分：平均分 $= (341 + 83) / 5 = 424 / 5 = 84.8 \geq 85$ ✗ 不够

\textbf{答：第五次至少要考 84 分。}

\section{四、加入生活常识的题——需要多想一步}

\subsection{例题3：出租车费用}

\begin{quote}
\textbf{某城市出租车的收费规则：起步价 8 元（含3公里），超过3公里后每公里收费 2 元。小李身上只有 30 元，他最远能坐多少公里？}
\end{quote}

\textbf{第1步：设未知数}

设小李一共坐了 $x$ 公里（$x > 3$，因为不超过3公里就只要8元，不用算）

\textbf{第2步：理解费用怎么算}

前3公里：8元（固定的）

超出的部分：$(x - 3)$ 公里，每公里2元，所以是 $2(x - 3)$ 元

总费用 $= 8 + 2(x - 3)$

\textbf{第3步：找不等关系}

他只有30元，总费用不能超过30元

不等式：$8 + 2(x - 3) \leq 30$

\begin{quote}
📖 \textbf{这个式子直接用语言说就是：} “起步价8元，加上超出3公里后按每公里2元算的额外费用，这个总费用小于等于30元”，也就是“打车花的钱不能比兜里的30元多”。
\end{quote}

\textbf{第4步：解不等式}

第一小步：去括号

$2(x - 3) = 2x - 6$

等式变成：$8 + 2x - 6 \leq 30$

第二小步：合并同类项（左边的数字：8 $-$ 6 $=$ 2）

$2 + 2x \leq 30$

第三小步：移项，把 2 移到右边变 $-2$

$2x \leq 30 - 2$

$2x \leq 28$

第四小步：两边除以 2（正数，方向不变）

$x \leq 14$

\textbf{第5步：回答问题}

\textbf{答：最远能坐 14 公里。}

\textbf{第6步：检验}

坐14公里：费用 $= 8 + 2\times(14-3) = 8 + 2\times11 = 8 + 22 = 30$ 元，正好30 $\leq 30$ ✓

坐15公里：费用 $= 8 + 2\times(15-3) = 8 + 24 = 32$ 元，$32 > 30$ ✗ 不够

\subsection{例题4：装箱问题}

\begin{quote}
\textbf{一个箱子最多能装 50 公斤。已经装了 18 公斤的东西。现在要往里面放一些书，每本书重 0.8 公斤，最多还能放几本？}
\end{quote}

\textbf{第1步：设未知数}

设还能放 $x$ 本书

\textbf{第2步：找不等关系}

箱子总重量不能超过50公斤 $\to \leq 50$

\textbf{第3步：翻译}

已经装了18公斤

再放 $x$ 本书，重 $0.8x$ 公斤

总重量 $= 18 + 0.8x$

不等式：$18 + 0.8x \leq 50$

\begin{quote}
📖 \textbf{这个式子直接用语言说就是：} “已经装的18公斤，加上新放进去的书的重量（每本0.8公斤，$x$本就是$0.8x$公斤），加在一起不能超过箱子的承重上限50公斤”。
\end{quote}

\textbf{第4步：解不等式}

第一小步：移项

$0.8x \leq 50 - 18$

$0.8x \leq 32$

第二小步：两边除以 0.8（正数，方向不变）

$x \leq 32 \div 0.8$

$x \leq 40$

\textbf{第5步：回答问题}

书的数量必须是整数，$x \leq 40$，所以最多放 40 本

\textbf{第6步：检验}

放40本：$18 + 0.8\times40 = 18 + 32 = 50 \leq 50$ ✓

放41本：$18 + 0.8\times41 = 18 + 32.8 = 50.8 > 50$ ✗

\textbf{答：最多还能放 40 本书。}

\section{五、比较类问题——“哪个更划算？”}

\subsection{例题5：两个商店买同一种商品}

\begin{quote}
\textbf{A 商店：每件商品 25 元，买 10 件以上（不含10件）打八折。}
\textbf{B 商店：每件商品 25 元，满 200 元减 50 元。}
\textbf{小明要买 $x$ 件商品（$x > 10$），买多少件时去 A 商店更便宜？}
\end{quote}

\textbf{第1步：分别写出两个商店的花费}

\textbf{A 商店：}

超过10件打八折

原价每件25元，八折就是 $25 \times 0.8 = 20$ 元/件

买 $x$ 件的总价 $= 20x$

\textbf{B 商店：}

原价 $= 25x$

“满200减50”$\to$ 每满200减50

先简单理解为：总价是 $25x$，如果 $25x \geq 200$，就减50

当 $x > 10$ 时，$25x > 250 > 200$，肯定满200了

所以 B 的花费 $= 25x - 50$

\textbf{第2步：找不等关系}

“A 比 B 便宜”$\to$ A 的花费 $<$ B 的花费

$20x < 25x - 50$

\begin{quote}
📖 \textbf{这个式子直接用语言说就是：} “在A商店花的钱（每件20元，$x$件就是$20x$元）比在B商店花的钱（原价$25x$元减掉优惠的50元）要少”，也就是“A比B便宜”。
\end{quote}

\textbf{第3步：解不等式}

第一小步：把右边的 $25x$ 移到左边变 $-25x$

$20x - 25x < -50$

第二小步：合并同类项

$-5x < -50$

第三小步：两边除以 $-5$

\textbf{注意！$-5$ 是负数！不等号方向要反过来！}

$x > -50 \div (-5)$

$x > 10$

\textbf{第4步：回答问题}

题目已经说了 $x > 10$，所以当买超过 10 件（即 11 件及以上）时，A 商店更便宜

\textbf{第5步：举个具体的数验证}

试 $x = 12$：

A：$20 \times 12 = 240$ 元

B：$25 \times 12 - 50 = 300 - 50 = 250$ 元

$240 < 250$ ✓ A 更便宜

试 $x = 11$：

A：$20 \times 11 = 220$ 元

B：$25 \times 11 - 50 = 275 - 50 = 225$ 元

$220 < 225$ ✓ A 更便宜

\textbf{答：买超过 10 件时，去 A 商店更便宜。}

\section{六、“够不够”类问题}

\subsection{例题6：运动会预算}

\begin{quote}
\textbf{学校运动会要买矿泉水和面包。矿泉水每瓶 2 元，需要买 80 瓶。面包每个 3 元。学校的预算是 400 元。在买完矿泉水之后，最多还能买多少个面包？}
\end{quote}

\textbf{第1步：设未知数}

设买 $x$ 个面包

\textbf{第2步：先算矿泉水花多少钱}

80 瓶 $\times$ 2 元 $= 160$ 元

\textbf{第3步：写出总花费}

矿泉水 + 面包 $= 160 + 3x$

\textbf{第4步：找不等关系}

总花费不能超过预算 400 元

$160 + 3x \leq 400$

\begin{quote}
📖 \textbf{这个式子直接用语言说就是：} “矿泉水花掉的160元，加上面包花的钱（每个3元，$x$个就是$3x$元），这两部分加在一起不能超过学校给的400元预算”。
\end{quote}

\textbf{第5步：解不等式}

第一小步：移项

$3x \leq 400 - 160$

$3x \leq 240$

第二小步：两边除以 3

$x \leq 80$

\textbf{第6步：回答问题}

\textbf{答：最多还能买 80 个面包。}

\textbf{检验：}

买80个面包：$160 + 3\times80 = 160 + 240 = 400 \leq 400$ ✓

买81个面包：$160 + 3\times81 = 160 + 243 = 403 > 400$ ✗

\subsection{例题7：存钱问题}

\begin{quote}
\textbf{小华现在有 50 元存款，他每个月存 15 元。至少存多少个月后，他的存款不少于 200 元？}
\end{quote}

\textbf{第1步：设未知数}

设存 $x$ 个月

\textbf{第2步：写出 $x$ 个月后的总存款}

原来有 50 元，每月存 15 元，$x$ 个月存 $15x$ 元

总存款 $= 50 + 15x$

\textbf{第3步：找不等关系}

“存款不少于200元”$\to \geq 200$

$50 + 15x \geq 200$

\begin{quote}
📖 \textbf{这个式子直接用语言说就是：} “原来存着的50元，加上之后每月攒的钱（每月15元，$x$个月就是$15x$元），加起来要大于等于200元”，也就是“攒够200元，刚好200也算够”。
\end{quote}

\textbf{第4步：解不等式}

第一小步：移项

$15x \geq 200 - 50$

$15x \geq 150$

第二小步：两边除以 15

$x \geq 10$

\textbf{第5步：回答问题}

\textbf{答：至少存 10 个月。}

\textbf{检验：}

存10个月：$50 + 15\times10 = 50 + 150 = 200 \geq 200$ ✓

存9个月：$50 + 15\times9 = 50 + 135 = 185 \geq 200$ ✗ 不够

\section{七、综合难题——多个条件同时满足}

\subsection{例题8：租车问题}

\begin{quote}
\textbf{学校组织春游，共有 52 名师生。现有两种车可以租：}
\textbf{大车：每辆最多坐 10 人，租金 200 元/辆}
\textbf{小车：每辆最多坐 4 人，租金 100 元/辆}

\textbf{要求：所有人都要有座位，并且总租金不超过 1200 元。租大车 $x$ 辆，那么 $x$ 的取值范围是多少？}
\end{quote}

\textbf{第1步：设未知数}

题目已经说了，租大车 $x$ 辆

\textbf{第2步：算出需要多少小车}

大车坐 $10x$ 人

剩下的人数：$52 - 10x$ 人

这些人要坐小车，每辆小车坐4人

需要小车的数量：至少 $(52 - 10x) \div 4$ 辆

但为了简化，我们先设小车 $y$ 辆，要求 $4y \geq 52 - 10x$，也就是 $y \geq \dfrac{52 - 10x}{4}$

为了花最少的钱，小车数量取最少的，也就是 $y$ 刚好能装下剩余的人

不过我们先别管小车具体几辆，换一种思路来列不等式

\textbf{换个思路，列两个条件：}

\textbf{条件一：所有人都有座位}

大车坐的人 + 小车坐的人 $\geq 52$

但我们不知道小车几辆。换个角度——大车至少要坐一些人，所以 $x$ 不能太少

如果大车坐了 $10x$ 人，剩余 $52 - 10x$ 人需要坐小车。

剩余人数必须 $\geq 0$（不能是负数，否则意味着大车就够了，不需要小车）

所以：$52 - 10x \geq 0 \to x \leq 5$（第一个范围）

\begin{quote}
📖 \textbf{这个式子直接用语言说就是：} “总人数52减去大车坐的人数（$10x$），剩下要坐小车的人数不能是负数”，也就是“大车不能租太多，坐的人不能比总人数还多”。
\end{quote}

同时 $x$ 至少为 0，而且是整数：$x \geq 0$

\textbf{条件二：总租金不超过1200元}

先算最“省”的情况：大车 $x$ 辆，剩下的人刚好坐满小车

剩余人数 $= 52 - 10x$

需要小车数量 $= (52 - 10x) / 4$。如果除不尽，要向上取整（因为多出来的人也要占一辆车）

为了列不等式方便，我们先假设刚好整除的情况来分析：

总租金 $= $ 大车租金 + 小车租金

大车租金 $= 200x$

小车数量 $= (52 - 10x) / 4$

小车租金 $= 100 \times (52 - 10x) / 4 = 25(52 - 10x) = 1300 - 250x$

总租金 $= 200x + 1300 - 250x = 1300 - 50x$

不等式：$1300 - 50x \leq 1200$

\begin{quote}
📖 \textbf{这个式子直接用语言说就是：} “大车的租金加上小车的租金，算出来的总花费（化简后是1300 $-$ 50$x$元）不能超过学校给的1200元预算”。
\end{quote}

\textbf{第3步：解不等式}

第一小步：移项

$-50x \leq 1200 - 1300$

$-50x \leq -100$

第二小步：两边除以 $-50$

\textbf{注意！$-50$ 是负数，方向要反过来！}

$x \geq -100 \div (-50)$

$x \geq 2$

\textbf{第4步：合并所有条件}

条件一：$x \leq 5$

条件二：$x \geq 2$

另外 $x$ 必须是正整数

合起来：$2 \leq x \leq 5$

也就是 $x$ 可以是 2、3、4、5

\textbf{第5步：验证几个值}

$x = 2$（租2辆大车）：

大车坐 20 人，剩 32 人

需要小车 32 $\div$ 4 $= 8$ 辆

总租金 $= 200\times2 + 100\times8 = 400 + 800 = 1200 \leq 1200$ ✓

$x = 1$（租1辆大车）：

大车坐 10 人，剩 42 人

需要小车 42 $\div$ 4 $= 10.5 \to$ 向上取整 $= 11$ 辆

总租金 $= 200\times1 + 100\times11 = 200 + 1100 = 1300 > 1200$ ✗ 超了！

$x = 5$（租5辆大车）：

大车坐 50 人，剩 2 人

需要小车 1 辆

总租金 $= 200\times5 + 100\times1 = 1000 + 100 = 1100 \leq 1200$ ✓

\textbf{答：大车应租 2 到 5 辆（即 $x = 2、3、4、5$）。}

\section{八、解题流程总结}

面对任何不等式应用题，按这个流程走：

\begin{quote}
第1步：读题，找出“谁是未知的” $\to$ 设 $x$

第2步：用 $x$ 把相关的量都表示出来

（花了多少钱？走了多少路？总共多少人？……）

第3步：找关键句——哪句话在说大小关系？

（“最多”“至少”“不超过”“够不够”……）

第4步：把那句话翻译成不等式

第5步：用语言说一遍你列的式子，确认意思对不对

第6步：解不等式

（特别注意：除以负数要变方向！）

第7步：结合实际情况调整答案

（人数、个数必须是整数；数量不能是负数……）

第8步：代入原题检验
\end{quote}

\section{九、更多练习题（附逐步解答）}

\subsection{练习1}

\begin{quote}
\textbf{一根绳子长 20 米，剪去 $x$ 米后，剩下的部分不少于 8 米。$x$ 的范围是什么？}
\end{quote}

第一步：设剪去 $x$ 米

第二步：剩下的长度 $= 20 - x$

第三步：找关键词——“不少于”$\to \geq$

第四步：列不等式

$20 - x \geq 8$

\begin{quote}
📖 \textbf{这个式子直接用语言说就是：} “原来的20米减去剪掉的$x$米，剩下的长度大于等于8米”，也就是“剪完之后剩下的不能少于8米，刚好等于8米也行”。
\end{quote}

第五步：解不等式

移项：$-x \geq 8 - 20$

$-x \geq -12$

两边除以 $-1$（负数，方向反过来）

$x \leq 12$

第六步：另外，剪去的长度不能是负数，所以 $x \geq 0$

第七步：合并

$0 \leq x \leq 12$

\textbf{答：$x$ 的范围是 $0 \leq x \leq 12$（最多剪去 12 米）。}

检验：

剪12米：$20 - 12 = 8 \geq 8$ ✓

剪13米：$20 - 13 = 7 \geq 8$ ✗

\subsection{练习2}

\begin{quote}
\textbf{某停车场收费标准：前 2 小时收费 5 元，之后每小时收费 3 元。小王只愿意花不超过 20 元，他最多能停多少小时？}
\end{quote}

第一步：设一共停 $x$ 小时（$x > 2$）

第二步：算费用

前2小时：5元

之后的时间：$x - 2$ 小时

之后的费用：$3(x - 2)$

总费用 $= 5 + 3(x - 2)$

第三步：“不超过20元”$\to \leq 20$

第四步：列不等式

$5 + 3(x - 2) \leq 20$

\begin{quote}
📖 \textbf{这个式子直接用语言说就是：} “前2小时固定收的5元，加上超过2小时后按每小时3元收的费用，总停车费不能超过20元”。
\end{quote}

第五步：解不等式

去括号：

$5 + 3x - 6 \leq 20$

合并同类项（5 $-$ 6 $= -1$）：

$-1 + 3x \leq 20$

移项（$-1$ 移到右边变 $+1$）：

$3x \leq 20 + 1$

$3x \leq 21$

两边除以 3：

$x \leq 7$

\textbf{答：最多能停 7 小时。}

检验：

停7小时：$5 + 3\times(7-2) = 5 + 15 = 20 \leq 20$ ✓

停8小时：$5 + 3\times(8-2) = 5 + 18 = 23 > 20$ ✗

\subsection{练习3}

\begin{quote}
\textbf{小明和小红一起做题。小明已经做了 30 道，小红已经做了 46 道。从现在开始，小明每小时做 8 道，小红每小时做 5 道。至少几小时后，小明做的题数能超过小红？}
\end{quote}

第一步：设 $x$ 小时后

第二步：$x$ 小时后各做了多少道？

小明：$30 + 8x$

小红：$46 + 5x$

第三步：找关键词——“小明超过小红”$\to$ 小明 $>$ 小红

第四步：列不等式

$30 + 8x > 46 + 5x$

\begin{quote}
📖 \textbf{这个式子直接用语言说就是：} “小明已经做的30道加上之后每小时做8道做了$x$小时（$8x$道），要比小红已经做的46道加上之后每小时做5道做了$x$小时（$5x$道）多”，也就是“小明的总数严格超过小红的总数”。
\end{quote}

第五步：解不等式

把含 $x$ 的移到左边，数字移到右边：

$8x - 5x > 46 - 30$

合并同类项：

$3x > 16$

两边除以 3：

$x > \dfrac{16}{3}$

$x > 5.33\ldots$

第六步：题目问“至少几小时”，而且小时应该取整数

$x > 5.33$，所以最小的整数是 6

\textbf{答：至少 6 小时后，小明做的题数能超过小红。}

检验：

第6小时后：小明 $30+48=78$，小红 $46+30=76$，$78 > 76$ ✓

第5小时后：小明 $30+40=70$，小红 $46+25=71$，$70 > 71$ ✗ 还没超过

\subsection{练习4}

\begin{quote}
\textbf{一个长方形的长是宽的 2 倍。如果周长不超过 36 厘米，宽最大是多少厘米？}
\end{quote}

第一步：设宽为 $x$ 厘米

第二步：长是宽的 2 倍 $\to$ 长 $= 2x$ 厘米

第三步：周长怎么算？$\to$ 周长 $= 2 \times$（长 + 宽）$= 2(2x + x) = 2 \times 3x = 6x$

第四步：找关键词——“周长不超过36”$\to \leq 36$

第五步：列不等式

$6x \leq 36$

\begin{quote}
📖 \textbf{这个式子直接用语言说就是：} “这个长方形的周长（算出来是$6x$厘米）不能超过36厘米”，也就是“围一圈的总长度最多36厘米，等于36也行”。
\end{quote}

第六步：解不等式

两边除以 6

$x \leq 6$

另外，宽必须大于 0（不能是0或负数）：$x > 0$

\textbf{答：宽最大是 6 厘米。}

检验：

宽6厘米，长12厘米，周长 $= 2\times(12+6) = 36 \leq 36$ ✓

宽7厘米，长14厘米，周长 $= 2\times(14+7) = 42 > 36$ ✗

\section{十、易错点提醒}

\begin{center}
\begin{tabular}{ll}
\toprule
常见错误 & 正确做法 \\
\midrule
“至少”“不少于”用成了 $>$ & “至少”“不少于”是 $\geq$，能取等号 \\
“不超过”用成了 $<$ & “不超过”是 $\leq$，能取等号 \\
算出 $x \leq 7.5$ 就直接答7.5个 & 如果是人数、个数等，必须取整数 \\
忘记 $x$ 本身的限制 & 比如人数不能为负、时间不能为负、超过3公里才额外计费等 \\
列完不等式忘记解就直接答 & 必须解出来再回答 \\
列完式子不确认意思对不对 & 用语言把式子“读”一遍，看看和题目说的是不是一回事 \\
解完忘记检验 & 把答案代回原题确认一下，养成好习惯 \\
\bottomrule
\end{tabular}
\end{center}

以上就是不等式应用题的完整方法。核心就是一句话：\textbf{把生活语言翻译成数学语言，解出来，再翻译回生活语言回答问题。} 列完式子之后，一定要用大白话把式子的意思说一遍——如果说出来的意思和题目对得上，那你的式子就列对了！

\chapter{函数，无非直线、线段、曲线}

我先问你一个问题。

你去便利店买水。

一瓶水3块钱。

你买1瓶，花3块。

你买2瓶，花6块。

你买5瓶，花15块。

你买 $x$ 瓶，花多少钱？

\[
3x \text{ 块。}
\]

就这么简单。

你买的瓶数变了，花的钱就跟着变。

\textbf{瓶数是输入，花的钱是输出。}

输入变了，输出跟着变。

这种“输入决定输出”的关系，就叫做\textbf{函数}。

你看，函数没有那么神秘。

你每天都在用函数，只是你不知道它叫这个名字。

气温决定你穿多厚——气温是输入，穿衣厚度是输出。函数。

时间决定你在哪里——时间是输入，位置是输出。函数。

输入的东西变了，输出的东西跟着变。只要有这种关系，就是函数。

数学里，我们通常用 $x$ 表示输入，用 $y$ 表示输出。

然后用一个式子，告诉你输入和输出之间的关系。

比如刚才那个买水的例子：

\[
y = 3x
\]

$x$ 是你买的瓶数（输入），$y$ 是你花的钱（输出）。

$x$ 变了，$y$ 跟着变。这就是函数。

好。现在我要带你认识几个最基本的函数。

它们是函数世界里的“老朋友”。

你以后会反复碰到它们。

先把它们的长相、性格、脾气都摸清楚，后面才不会慌。

\section{第一种：一次函数——直线}

\[
y = kx + b
\]

这是一次函数的通用写法。

$k$ 和 $b$ 是两个数字，可以是任何数。

$x$ 是输入，$y$ 是输出。

先别管 $k$ 和 $b$ 是什么意思。

我先给你几个具体的例子：

\[
y = 2x + 1
\]

\[
y = -x + 3
\]

\[
y = 3x
\]

\[
y = 5
\]

这些都是一次函数。（最后那个 $y = 5$ 有点特殊，等一下讲。）

\subsection{语言记忆法：一次函数就是一条直线}

你往坐标系里画这个函数，画出来的是一条\textbf{直线}。

（还记得坐标系吗？横轴是 $x$，纵轴是 $y$。每一个 $x$ 的值，对应一个 $y$ 的值，画一个点。所有的点连起来，就是这个函数的图像。）

为什么叫“一次”函数？

因为 $x$ 的最高次数是\textbf{一次}。$x^1$，就是 $x$ 本身，没有 $x^2$ 或者更高的幂。

\textbf{次数是一，图像是直线。}

这两件事绑在一起，你就不会忘了。

\subsection{$k$ 是什么——坡度}

\[
y = kx + b
\]

$k$ 叫做\textbf{斜率}。

“斜率”这个词听起来很专业，但它的意思很简单：

\textbf{$k$ 告诉你这条直线有多“斜”。}

语言记忆法：

\begin{quote}
\textbf{$k$ 就是这条直线的坡度。$k$ 越大，坡越陡。}
\end{quote}

$k > 0$：直线从左下往右上走——\textbf{上坡}。

$k < 0$：直线从左上往右下走——\textbf{下坡}。

$k = 0$：直线是水平的——\textbf{平路}。

更精确一点说：

$k$ 的意思是“$x$ 每增加1，$y$ 增加多少”。

比如 $y = 2x + 1$，$k = 2$。

$x = 1$ 时，$y = 2 \times 1 + 1 = 3$。

$x = 2$ 时，$y = 2 \times 2 + 1 = 5$。

$x$ 从1增加到2，增加了1。$y$ 从3增加到5，增加了2。

$x$ 增加1，$y$ 增加2。

$k = 2$，对上了。

\textbf{$k$ 就是“走一步，升多少”。}

坡度。

\subsection{$b$ 是什么——起点高度}

$b$ 叫做\textbf{截距}，具体来说是\textbf{纵截距}。

听起来很复杂，但意思很简单：

\textbf{$b$ 就是这条直线和纵轴（$y$ 轴）相交的那个点的高度。}

语言记忆法：

\begin{quote}
\textbf{$b$ 就是你站在 $x = 0$ 的地方，抬头看，直线在哪里。}
\end{quote}

把 $x = 0$ 代进去：

\[
y = k \times 0 + b = b
\]

当 $x = 0$ 时，$y = b$。

所以直线和 $y$ 轴的交点，就在 $(0, b)$。

$b$ 大，直线在上面。$b$ 小，直线在下面。$b = 0$，直线过原点。

\subsection{特殊情况：$y = b$（水平线）}

\[
y = 5
\]

这是一次函数的特殊情况。$k = 0$，直线是水平的。

不管 $x$ 是什么，$y$ 永远等于5。

输入变了，输出不变。

图像是一条\textbf{水平直线}，高度永远在5那里。

\subsection{画一条直线，只需要两个点}

想画 $y = 2x + 1$ 的图像，你只需要找两个点，连一条线就行了。

为什么只需要两个点？

因为两点确定一条直线。这是几何的基本常识。

\textbf{第一个点：} 让 $x = 0$：

\[
y = 2 \times 0 + 1 = 1
\]

得到点 $(0, 1)$。（这就是 $b = 1$ 告诉你的起点。）

\textbf{第二个点：} 让 $x = 1$：

\[
y = 2 \times 1 + 1 = 3
\]

得到点 $(1, 3)$。

把 $(0, 1)$ 和 $(1, 3)$ 画在坐标系里，连一条直线，延长。

完成。

来感受一下几个不同的一次函数：

\textbf{$y = x$}

$k = 1$，$b = 0$。

坡度是1（走一步升一步，45度角）。过原点。

\textbf{$y = 3x - 2$}

$k = 3$，$b = -2$。

坡度是3（走一步升三步，很陡）。和 $y$ 轴交于 $(0, -2)$。

\textbf{$y = -2x + 4$}

$k = -2$，$b = 4$。

坡度是负的（下坡）。和 $y$ 轴交于 $(0, 4)$。

\textbf{$y = 0.5x$}

$k = 0.5$，$b = 0$。

坡度很小（很平缓）。过原点。

你来试试：

\textbf{（1）} $y = 4x - 3$ 的斜率是多少？和 $y$ 轴交在哪里？

答案：斜率 $k = 4$。和 $y$ 轴交于 $(0, -3)$。

\textbf{（2）} $y = -x + 2$，当 $x = 3$ 时，$y$ 等于多少？

\[
y = -3 + 2 = -1
\]

答案：$y = -1$。

\textbf{（3）} 有一条直线，斜率是2，过点 $(0, 5)$，请写出它的函数式。

答案：$k = 2$，$b = 5$，所以 $y = 2x + 5$。

\section{第二种：正比例函数——过原点的直线}

\[
y = kx
\]

这其实是一次函数的特殊情况——$b = 0$。

$b = 0$ 意味着什么？

直线过原点 $(0, 0)$。

当 $x = 0$ 时，$y = 0$。输入是0，输出也是0。

语言记忆法：

\begin{quote}
\textbf{$y$ 和 $x$ 是同比例变化的。$x$ 变成两倍，$y$ 也变成两倍。}
\end{quote}

买水那个例子：$y = 3x$。

买的瓶数翻倍，花的钱也翻倍。

这就是正比例——两个量\textbf{同步}变化。

正比例函数的图像，是一条\textbf{过原点的直线}。

$k > 0$：过原点，往右上走。

$k < 0$：过原点，往右下走。

\section{第三种：二次函数——抛物线}

\[
y = ax^2 + bx + c
\]

这是二次函数的通用写法。

$a$、$b$、$c$ 是数字，$a \neq 0$。

为什么叫“二次”？

因为 $x$ 的最高次数是\textbf{二次}——$x^2$。

\textbf{次数是二，图像是曲线。}

这条曲线，有一个专门的名字：\textbf{抛物线}。

语言记忆法：

\begin{quote}
\textbf{你把一个球抛出去，它飞过的轨迹，就是一条抛物线。}
\end{quote}

先往上，到最高点，再往下落。

开口朝上的弓形曲线——这就是抛物线的样子。

（当然，有时候开口朝下。取决于 $a$ 是正还是负。）

\subsection{最简单的二次函数：$y = x^2$}

先从最干净的开始。

$a = 1$，$b = 0$，$c = 0$。

让 $x$ 取几个值，看看 $y$ 是多少：

\begin{center}
\begin{tabular}{cccccccc}
\toprule
$x$ & $-3$ & $-2$ & $-1$ & $0$ & $1$ & $2$ & $3$ \\
\midrule
$y = x^2$ & $9$ & $4$ & $1$ & $0$ & $1$ & $4$ & $9$ \\
\bottomrule
\end{tabular}
\end{center}

你看到了几件事：

\textbf{第一：} $x = 0$ 时，$y = 0$。图像过原点。

\textbf{第二：} $x = 1$ 和 $x = -1$ 时，$y$ 都等于1。$x = 2$ 和 $x = -2$ 时，$y$ 都等于4。

\textbf{正数和负数配对，平方之后一样大。}

这意味着图像关于 $y$ 轴\textbf{对称}——左右是镜像的。

\textbf{第三：} $y$ 的最小值是0，在 $x = 0$ 的时候取到。

往左往右走，$y$ 都变大。

图像开口朝上，像一个碗。

\textbf{第四：} $x$ 越大，$y$ 增长得越快。

从 $x=1$ 到 $x=2$，$y$ 从1跳到4，增加了3。

从 $x=2$ 到 $x=3$，$y$ 从4跳到9，增加了5。

越走越陡。这就是二次函数和一次函数最不一样的地方——它不是匀速增长，是\textbf{加速增长}。

\subsection{$a$ 是什么——开口方向和胖瘦}

\[
y = ax^2
\]

$a$ 控制两件事：

\textbf{第一：开口方向。}

$a > 0$：开口朝上，像碗。（$y = x^2$ 就是这样。）

$a < 0$：开口朝下，像山顶。

语言记忆法：

\begin{quote}
\textbf{$a > 0$，碗装水。$a < 0$，山顶站人。}
\end{quote}

\textbf{第二：胖瘦（开口大小）。}

$|a|$ 越大，开口越窄（越瘦）。

$|a|$ 越小，开口越宽（越胖）。

$y = 5x^2$：很瘦，开口窄，往上窜得快。

$y = 0.1x^2$：很胖，开口宽，变化很慢。

\subsection{顶点——曲线的最高点或最低点}

抛物线有一个“转折点”——从这个点开始，从下往上变成从上往下（或者反过来）。

这个点叫做\textbf{顶点}。

对于 $y = x^2$ 来说，顶点就是 $(0, 0)$——最低点。

对于 $y = -x^2$ 来说，顶点也是 $(0, 0)$——但这次是最高点。

顶点是抛物线的“灵魂”。知道了顶点，你就大概知道这条曲线长什么样了。

（顶点怎么算？这个我们以后会专门讲。现在先知道它是什么就好。）

\subsection{对称轴——曲线的镜子}

从上面的表格你看到了：$y = x^2$ 左右对称，$y$ 轴就是它的“镜子”。

$y$ 轴就是 $y = x^2$ 的\textbf{对称轴}。

左边的图像和右边的图像，沿着对称轴折过去，完全重合。

对称轴永远过顶点——顶点在哪里，对称轴就在哪里。

来感受几个二次函数：

\textbf{$y = 2x^2$}

$a = 2 > 0$，开口朝上，比 $y = x^2$ 更瘦。顶点在原点。

\textbf{$y = -x^2 + 4$}

$a = -1 < 0$，开口朝下，像山顶。

$c = 4$，顶点在 $(0, 4)$——比原点高了4个单位。

（当 $x = 0$ 时，$y = 4$。顶点就在这里。）

\textbf{$y = x^2 - 2x + 1$}

这个复杂一点。顶点不在 $y$ 轴上了。

（先感受一下它的存在就好，怎么算顶点我们后面再说。）

你来算几个：

\textbf{（1）} $y = -3x^2$，当 $x = 2$ 时，$y$ 等于多少？

\[
y = -3 \times 2^2 = -3 \times 4 = -12
\]

答案：$y = -12$。（开口朝下，$x$ 越大，$y$ 越小。）

\textbf{（2）} $y = x^2 + 1$，当 $x = -3$ 时，$y$ 等于多少？

\[
y = (-3)^2 + 1 = 9 + 1 = 10
\]

答案：$y = 10$。

\textbf{（3）} $y = x^2 + 1$ 的顶点在哪里？

当 $x = 0$ 时，$y = 0 + 1 = 1$。

（因为 $b = 0$，对称轴就是 $y$ 轴，顶点在 $y$ 轴上。）

答案：顶点在 $(0, 1)$。

\section{第四种：反比例函数——双曲线}

\[
y = \frac{k}{x}
\]

$k$ 是一个不等于0的数。$x$ 也不能等于0（不然分母为0，没意义）。

语言记忆法：

\begin{quote}
\textbf{$x$ 越大，$y$ 越小。$x$ 越小，$y$ 越大。两个人一个高一个矮，跷跷板。}
\end{quote}

$y = \dfrac{6}{x}$：

$x = 1$ 时，$y = 6$。

$x = 2$ 时，$y = 3$。

$x = 3$ 时，$y = 2$。

$x = 6$ 时，$y = 1$。

$x$ 变大，$y$ 变小。$x$ 和 $y$ 的乘积始终是6。

\[
x \times y = 6
\]

这就是“反比例”——一个变大，另一个等比例变小。跷跷板。

反比例函数的图像，叫做\textbf{双曲线}。

为什么叫“双”？

因为图像分成\textbf{两段}——一段在右上角，一段在左下角。

（或者，如果 $k < 0$，一段在右下角，一段在左上角。）

两段曲线，中间断开了，永远不相交。

而且这两段曲线，无论往哪个方向延伸，都无限接近坐标轴，但\textbf{永远碰不到}。

语言记忆法：

\begin{quote}
\textbf{双曲线是一个被坐标轴分开的两段曲线，永远在角落里，永远碰不到坐标轴。}
\end{quote}

为什么碰不到 $x$ 轴？

因为 $y = \dfrac{k}{x}$，$y$ 要等于0，就需要分子 $k = 0$。

但我们说了 $k \neq 0$。

所以 $y$ 永远不等于0，图像永远碰不到 $x$ 轴。

为什么碰不到 $y$ 轴？

因为 $x = 0$ 时，分母为0，$y$ 无意义。

$x$ 不能等于0，所以图像碰不到 $y$ 轴。

来感受几个：

\textbf{$y = \dfrac{1}{x}$}

$k = 1 > 0$。右上角一段，左下角一段。

\textbf{$y = \dfrac{-2}{x}$}

$k = -2 < 0$。右下角一段，左上角一段。

你来算：

\textbf{（1）} $y = \dfrac{12}{x}$，当 $x = 4$ 时，$y$ 等于多少？

\[
y = \frac{12}{4} = 3
\]

\textbf{（2）} $y = \dfrac{12}{x}$，当 $y = 2$ 时，$x$ 等于多少？

\[
2 = \frac{12}{x} \Rightarrow x = \frac{12}{2} = 6
\]

\textbf{（3）} $y = \dfrac{-6}{x}$，当 $x = 3$ 时，$y$ 等于多少？

\[
y = \frac{-6}{3} = -2
\]

（$k < 0$，$x > 0$，所以 $y < 0$。在右下角那段。）

\section{四种函数放在一起对比}

现在你认识了四个老朋友。

我把它们放在一起，你感受一下区别：

\begin{center}
\begin{tabular}{lll}
\toprule
函数 & 式子 & 关键字 \\
\midrule
一次函数 & $y = kx + b$ & 直线、坡度、截距 \\
正比例函数 & $y = kx$ & 过原点的直线、同比例变化 \\
二次函数 & $y = ax^2 + bx + c$ & 抛物线（曲线）、开口、顶点、对称轴 \\
反比例函数 & $y = \dfrac{k}{x}$ & 双曲线（两段）、跷跷板、永远碰不到轴 \\
\bottomrule
\end{tabular}
\end{center}

其实你仔细看，这四种函数在说同一件事：

\textbf{$x$ 变了，$y$ 跟着变。}

只是变的方式不同：

一次函数——\textbf{匀速变}。坡度固定，走一步升固定的高度。

正比例函数——\textbf{同步变}。翻倍就翻倍，减半就减半。

二次函数——\textbf{加速变}。越走越快，越走越陡。

反比例函数——\textbf{反向变}。一个大了另一个就小，跷跷板。

这四种“变化方式”，覆盖了你以后会遇到的绝大多数情况。

\section{这一章的灵魂}

函数不是数学家发明来折磨你的。

函数描述的，是这个世界上\textbf{两件事之间的关系}。

水价和花费——一次函数。

时间和路程——一次函数（匀速时）。

炮弹的飞行轨迹——二次函数。

人口和土地面积——反比例函数。

你以后每次看到函数，不要只看 $x$ 和 $y$。

要想一想：这两个字母，在现实里代表什么？它们之间的关系，是哪种变化方式？

数学是语言，函数是其中一句最常用的话。

你现在已经认识这句话了。

下一章，我们来给函数加上“规矩”——定义域、值域，还有函数的单调性和奇偶性。

你已经知道函数长什么样了。接下来，我们来研究它的脾气。

\chapter{我来好好跟你说说这几条线}

上一章我们认识了四种函数。

但我知道，当我说“抛物线”、“双曲线”的时候，你脑子里可能是一片空白。

或者有个模糊的形状，但说不清楚。

没关系。

这一章，我们不算题，不推公式。

我们就干一件事：

\textbf{把这几条线，用人话说清楚。}

\section{第一种：正比例函数——过原点的直线}

\[
y = kx
\]

先从最简单的开始。

你去买糖。

一颗糖5毛钱。

你买1颗，花5毛。

你买2颗，花1块。

你买3颗，花1块5。

你买0颗，花0块。

用数学写：

\[
y = 0.5x
\]

$x$ 是颗数，$y$ 是花的钱（单位：元）。

现在我们把这件事画出来。

画在一个坐标系里。

横轴是 $x$（买了几颗），纵轴是 $y$（花了多少钱）。

（还记得坐标系吗？就是两条互相垂直的数轴，横的是 $x$ 轴，竖的是 $y$ 轴，交叉的那个点叫原点，就是 $(0, 0)$。）

我们把刚才那些数据，一个一个点出来：

买0颗，花0块——点在 $(0, 0)$。

买1颗，花0.5块——点在 $(1, 0.5)$。

买2颗，花1块——点在 $(2, 1)$。

买3颗，花1.5块——点在 $(3, 1.5)$。

买4颗，花2块——点在 $(4, 2)$。

你把这些点连起来。

连出来的，是一条\textbf{直线}。

而且，这条直线\textbf{从原点出发}。

为什么从原点出发？

因为你买0颗的时候，花0块——$(0, 0)$ 这个点在线上。

这就是正比例函数最重要的特征：

\textbf{图像是一条直线，而且过原点。}

\subsection{$k$ 决定这条直线有多斜}

$y = 0.5x$，$k = 0.5$。

$y = 2x$，$k = 2$。

$y = 5x$，$k = 5$。

你用脑子想象一下：

$k = 0.5$：你走两步，才升高一步。坡很缓，直线很平。

$k = 2$：你走一步，升高两步。坡比较陡。

$k = 5$：你走一步，升高五步。坡非常陡，直线几乎是竖起来的。

\textbf{$k$ 越大，直线越陡。}

那 $k$ 是负数呢？

$y = -x$，$k = -1$。

$x = 1$ 时，$y = -1$。

$x = 2$ 时，$y = -2$。

$x$ 越大，$y$ 越小。

你往右走，地势反而往下降。

\textbf{$k < 0$，直线是下坡的。从左上角走向右下角。}

总结一下正比例函数的长相：

\begin{quote}
\textbf{$k > 0$：从原点出发，往右上方走。上坡。}

\textbf{$k < 0$：从原点出发，往右下方走。下坡。}

\textbf{$k$ 越大，坡越陡。$k$ 越接近0，越平缓。}
\end{quote}

就这些。一条过原点的直线。坡度由 $k$ 决定。

\section{第二种：二次函数——一条曲线}

\[
y = ax^2
\]

好，现在难度升一级。

先说最重要的一件事：

\textbf{它画出来，就是一条曲线。}

不是直线了。是弯的。

一条光滑的、弯弯的曲线。

你以前肯定见过这种形状——就是那种U形的弧线。碗的边缘，就是这个形状。

这条曲线，数学家给它起了一个名字，叫\textbf{抛物线}。

为什么叫“抛物线”？

你把一个球斜着扔出去。球先往上飞，到最高点，然后往下落。

它飞过的那条轨迹，画出来，就是这个形状。

“抛出去的物体走过的线”——抛物线。

名字是这么来的。

但你不需要记这个名字背后的故事。

你只需要记住：

\textbf{二次函数的图像，是一条弯曲的U形曲线。}

就这一句话。

我们先看最简单的：

\[
y = x^2
\]

让 $x$ 从 $-3$ 变到 $3$，看看 $y$ 怎么变：

$x = -3$，$y = (-3)^2 = 9$

$x = -2$，$y = (-2)^2 = 4$

$x = -1$，$y = (-1)^2 = 1$

$x = 0$，$y = 0^2 = 0$

$x = 1$，$y = 1^2 = 1$

$x = 2$，$y = 2^2 = 4$

$x = 3$，$y = 3^2 = 9$

你把这些点在坐标系里标出来：

最左边和最右边的点，都在高处（$y = 9$）。

中间那个点，在最低处（$y = 0$，就是原点）。

从左边高处，慢慢往中间降，降到最低点，然后再从中间往右边爬升。

连起来，是一个\textbf{U形的曲线}。

开口朝上，像一个碗。

这个曲线有几个重要的特征，我一个一个说：

\subsection{特征一：左右对称}

你看那些数据：

$x = -3$ 和 $x = 3$，$y$ 都是9。

$x = -2$ 和 $x = 2$，$y$ 都是4。

$x = -1$ 和 $x = 1$，$y$ 都是1。

正数和负数配对，$y$ 值一样。

所以这条曲线，左边和右边是镜像的——完全对称。

\textbf{对称轴就是 $y$ 轴。}

你沿着 $y$ 轴折过去，左边和右边完全重合。

\subsection{特征二：有一个最低点}

$y = x^2$ 的最低点在 $(0, 0)$。

就是原点。

从这个点往左走，$y$ 变大。

往右走，$y$ 也变大。

往哪边走都是上坡。

这个最低点，叫做\textbf{顶点}。

$y = x^2$ 的顶点在原点。

\subsection{特征三：不是匀速变化的}

你看，从 $x = 0$ 走到 $x = 1$，$y$ 从0变到1，升了1。

从 $x = 1$ 走到 $x = 2$，$y$ 从1变到4，升了3。

从 $x = 2$ 走到 $x = 3$，$y$ 从4变到9，升了5。

每走一步，升的越来越多。

\textbf{越走越陡。}

这和直线完全不一样——直线是匀速的，每走一步升的一样多。

抛物线是\textbf{加速的}，越走越快。

\subsection{$a$ 决定开口方向和胖瘦}

\[
y = ax^2
\]

\textbf{$a > 0$：开口朝上，像碗。}

$y = x^2$，$y = 2x^2$，$y = 0.1x^2$——全是开口朝上的碗。

\textbf{$a < 0$：开口朝下，像山顶。}

$y = -x^2$——把 $y = x^2$ 翻过来。

原来最低点在底下，翻过来就变成最高点在上面。

从最高点往左或者往右，$y$ 都在减小。

$a$ 的大小决定碗的\textbf{胖瘦}：

$y = 5x^2$：$a = 5$，很大。碗很窄，很瘦，两边陡峭地往上升。

$y = 0.1x^2$：$a = 0.1$，很小。碗很宽，很胖，两边缓缓地往上升。

\textbf{$|a|$ 越大，碗越窄越瘦。$|a|$ 越小，碗越宽越胖。}

\subsection{顶点不一定在原点}

\[
y = x^2 + 1
\]

和 $y = x^2$ 相比，每个 $y$ 值都加了1。

整条曲线往上移了1格。

顶点从 $(0, 0)$ 变成了 $(0, 1)$。

\[
y = (x - 2)^2
\]

这个有点意思。

当 $x = 2$ 时，$y = (2-2)^2 = 0$。最低点在 $x = 2$ 的地方。

顶点是 $(2, 0)$。

整条曲线往右移了2格。

\textbf{曲线可以上下移，可以左右移。但形状不变——还是那个碗。}

这就是二次函数的本质。

不管它在哪里，不管开口朝上还是朝下，胖还是瘦——

\textbf{它永远是一条光滑的弯曲的线，有一个顶点，左右对称。}

\section{第三种：反比例函数——双曲线}

\[
y = \frac{k}{x}
\]

好，来到最有趣的一个。

我先给你一个生活场景。

你要铺地板。

这间房子的面积是12平方米。

你可以选不同尺寸的地砖：

地砖宽1米，那需要12块。

地砖宽2米，那需要6块。

地砖宽3米，那需要4块。

地砖宽4米，那需要3块。

地砖宽6米，那需要2块。

地砖宽12米，那需要1块。

你看到了吗？

\textbf{地砖越宽（$x$ 越大），需要的块数越少（$y$ 越小）。}

用数学写：

\[
y = \frac{12}{x}
\]

$x$ 是地砖宽度，$y$ 是需要的块数。

$x \times y = 12$。永远是12。

这就是反比例——\textbf{一个变大，另一个等比例变小。乘积不变。}

\subsection{它的图像长什么样？}

我把一些点列出来：

$x = 1$，$y = 12$——点在 $(1, 12)$，很高。

$x = 2$，$y = 6$——点在 $(2, 6)$，降了一些。

$x = 3$，$y = 4$——点在 $(3, 4)$。

$x = 4$，$y = 3$——点在 $(4, 3)$。

$x = 6$，$y = 2$——点在 $(6, 2)$，已经很低了。

$x = 12$，$y = 1$——点在 $(12, 1)$，非常低。

$x = 100$，$y = 0.12$——几乎贴着 $x$ 轴了。

你用脑子想象这些点连起来的样子：

从左上角开始，高高的。

然后往右走，迅速往下降。

降得越来越慢，越来越接近 $x$ 轴。

但\textbf{永远碰不到 $x$ 轴}。

因为 $y = \dfrac{12}{x}$，$y$ 要等于0，分子就得是0，但分子是12，永远不是0。

所以这条曲线，永远在 $x$ 轴上方，无限接近但永远碰不到。

那 $x$ 轴左边呢？

$x = -1$，$y = \dfrac{12}{-1} = -12$——点在 $(-1, -12)$，很低。

$x = -2$，$y = -6$——点在 $(-2, -6)$。

$x = -3$，$y = -4$——点在 $(-3, -4)$。

$x = -12$，$y = -1$——点在 $(-12, -1)$，已经快贴着 $x$ 轴了。

这些点连起来，在\textbf{左下角}形成另一段曲线。

形状和右上角那段一模一样，只是转了个方向。

所以，$y = \dfrac{k}{x}$（$k > 0$）的图像，是\textbf{两段曲线}：

\textbf{一段在右上角}：$x > 0$，$y > 0$。从左上往右下，无限接近两条轴但碰不到。

\textbf{一段在左下角}：$x < 0$，$y < 0$。从右上往左下，同样无限接近两条轴但碰不到。

这两段合在一起，就叫\textbf{双曲线}。

“双”就是两段的意思。

\subsection{为什么永远碰不到坐标轴？}

我再说一遍，因为这一点很有意思。

\textbf{碰不到 $x$ 轴：}

$y = 0$ 意味着 $\dfrac{k}{x} = 0$，意味着分子 $k = 0$。

但 $k \neq 0$（题目规定的），所以 $y$ 永远不等于0。

图像永远飘在 $x$ 轴上方（或下方），碰不到。

\textbf{碰不到 $y$ 轴：}

$x = 0$ 时，分母为0，$y$ 没有意义。

$x = 0$ 这个点，根本不在这个函数的定义域里。

所以图像在 $y$ 轴那里是断开的，碰不到。

语言记忆法：

\begin{quote}
\textbf{双曲线是两段在角落里的曲线。它无限接近坐标轴，但永远差那么一点点，永远碰不到。就像你追一个人，越来越近，但就是追不上。}
\end{quote}

\subsection{$k < 0$ 的时候}

\[
y = \frac{-6}{x}
\]

$x = 1$，$y = -6$。

$x = 2$，$y = -3$。

$x = 3$，$y = -2$。

$x > 0$ 时，$y < 0$——点在\textbf{右下角}。

$x < 0$ 时，$y > 0$——点在\textbf{左上角}。

所以 $k < 0$ 的双曲线，两段在\textbf{右下角和左上角}。

和 $k > 0$ 时刚好转了个方向。

总结：

\begin{quote}
\textbf{$k > 0$：右上角一段，左下角一段。}

\textbf{$k < 0$：右下角一段，左上角一段。}
\end{quote}

\section{三种函数放在一起感受一下}

现在你脑子里应该有三幅画面了：

\textbf{正比例函数 $y = kx$：}

一条直线。从原点出发。$k > 0$ 往右上走，$k < 0$ 往右下走。简单、干净、匀速。

\textbf{二次函数 $y = ax^2$：}

一条弯曲的线，像碗或者像山顶。有一个顶点，左右对称。$a > 0$ 开口朝上，$a < 0$ 开口朝下。越走越陡，加速变化。

\textbf{反比例函数 $y = \dfrac{k}{x}$：}

两段曲线，躲在坐标系的角落里。无限接近坐标轴，永远碰不到。$k > 0$ 在右上和左下，$k < 0$ 在右下和左上。

这三种，长相完全不同。

但它们做的事情是一样的：

\textbf{告诉你 $x$ 变了，$y$ 怎么跟着变。}

直线——匀速跟着变。

抛物线——加速跟着变。

双曲线——反向跟着变，而且越来越慢。

三种变化方式，三种长相，三种性格。

你现在都认识了。

脑子里的图像建好了吗？

不用画得多精确。

有个模糊的样子，知道它大概长什么样，这就够了。

以后遇到这几种函数，你能认出它们——

“哦，这是抛物线。”

“哦，这是双曲线。”

认出来了，才能接着分析它的脾气。

下一章，我们就来研究这些函数的脾气——定义域、值域、单调性、奇偶性。

走吧。

\chapter{给函数立规矩}

上一章你认识了三种函数的长相。

直线、曲线、两段躲在角落里的曲线。

你知道它们长什么样了。

但你有没有想过一个问题：

\textbf{$x$ 可以随便填吗？}

我举个例子。

你去买苹果。

一斤苹果5块钱。

你买 $x$ 斤，花 $y$ 块钱。

\[
y = 5x
\]

这是一个正比例函数，对吧？

好。现在我问你：

$x$ 能等于 $-3$ 吗？

你能买“负三斤”苹果吗？

不能。苹果的重量不可能是负数。

所以这个函数，$x$ 不能随便填。

\textbf{函数是有规矩的。}

这一章，我们来学四条规矩。

别怕，每一条都很简单。

我会用最通俗的话说给你听。

\section{第一条规矩：定义域}

\textbf{定义域，就是 $x$ 能填哪些数。}

就这么简单。

有些函数，$x$ 想填什么就填什么。

比如 $y = 2x + 1$。

$x = 1$？可以，$y = 3$。

$x = -100$？可以，$y = -199$。

$x = 0.5$？可以，$y = 2$。

什么数都能填。

这种函数的定义域，就是\textbf{所有的数}。

数学家写成：$(-\infty, +\infty)$。

（就是从负无穷大到正无穷大，什么都包括。）

但有些函数，有些数填不了。

为什么？

两种情况。

\subsection{第一种情况：分母不能是零}

\[
y = \frac{1}{x}
\]

如果 $x = 0$，分母就是0。

$1 \div 0$ 等于多少？

等于……没有意义。

你试试用计算器算 $1 \div 0$，它会报错。

因为没有任何一个数，乘以0能等于1。

所以 $x = 0$ 不能填。

定义域是：\textbf{除了0以外的所有数}。

再来一个：

\[
y = \frac{5}{x - 2}
\]

什么时候分母是0？

$x - 2 = 0$，也就是 $x = 2$ 的时候。

所以 $x = 2$ 不能填。

定义域是：\textbf{除了2以外的所有数}。

\textbf{一句话总结：看到分数，先看分母。让分母等于0，解出来那个数，不能填。}

\subsection{第二种情况：根号下不能是负数}

\[
y = \sqrt{x}
\]

$\sqrt{4}$ 等于多少？等于2，因为 $2 \times 2 = 4$。

$\sqrt{9}$ 等于多少？等于3，因为 $3 \times 3 = 9$。

那 $\sqrt{-1}$ 等于多少？

没有任何一个数，乘以自己能得到负数。

（正数乘正数是正数，负数乘负数还是正数。）

所以 $\sqrt{-1}$ 没有意义。根号下不能是负数。

定义域是：$x \geq 0$。

也就是\textbf{零和所有正数}。

再来一个：

\[
y = \sqrt{x - 3}
\]

根号下是 $x - 3$。

$x - 3$ 不能是负数。

$x - 3 \geq 0$

$x \geq 3$

定义域是：$x \geq 3$。

也就是\textbf{3和所有比3大的数}。

\textbf{一句话总结：看到根号，让根号下面的东西大于等于零。}

\subsection{两种情况一起出现}

\[
y = \frac{1}{\sqrt{x}}
\]

这个式子，既有分数，又有根号。

两条规矩都要守。

\textbf{规矩一：} 根号下不能是负数。$x \geq 0$。

\textbf{规矩二：} 分母不能是零。$\sqrt{x} \neq 0$，也就是 $x \neq 0$。

两条合起来：$x \geq 0$，并且 $x \neq 0$。

也就是 $x > 0$。

定义域是：\textbf{所有正数}（不包括零）。

\textbf{一句话总结：定义域就是“$x$ 能填什么”。分母不能为零，根号下不能为负。}

你来试几个：

\textbf{（1）} $y = \dfrac{1}{x + 5}$，$x$ 不能等于几？

分母 $x + 5 = 0$，$x = -5$。

答案：$x \neq -5$。

\textbf{（2）} $y = \sqrt{x + 2}$，$x$ 要满足什么条件？

$x + 2 \geq 0$，$x \geq -2$。

答案：$x \geq -2$。

\textbf{（3）} $y = \sqrt{4 - x}$，$x$ 要满足什么条件？

$4 - x \geq 0$，$4 \geq x$，也就是 $x \leq 4$。

答案：$x \leq 4$。

\section{第二条规矩：值域}

定义域是 $x$ 能填什么。

那 $y$ 呢？$y$ 能等于什么？

\textbf{值域，就是 $y$ 能取到哪些数。}

还是用买苹果的例子。

$y = 5x$，$x$ 是斤数，$y$ 是花的钱。

$x$ 只能是零或者正数（你不能买负数斤）。

那 $y$ 呢？

$x = 0$ 时，$y = 0$。

$x = 1$ 时，$y = 5$。

$x = 2$ 时，$y = 10$。

$y$ 也只能是零或者正数。

值域是：$y \geq 0$。

再来一个：

\[
y = x^2
\]

$x$ 可以是任何数。

那 $y$ 呢？

$x = 0$ 时，$y = 0$。

$x = 1$ 时，$y = 1$。

$x = -1$ 时，$y = 1$。

$x = 100$ 时，$y = 10000$。

$x = -100$ 时，$y = 10000$。

$y$ 能等于负数吗？

不能。任何数的平方都不是负数。

$y$ 最小是0，然后可以无限大。

值域是：$y \geq 0$。

再来一个：

\[
y = \frac{1}{x}
\]

$x$ 可以是除了0以外的任何数。

那 $y$ 呢？

$x = 1$ 时，$y = 1$。

$x = 2$ 时，$y = 0.5$。

$x = 100$ 时，$y = 0.01$。

$x = 0.01$ 时，$y = 100$。

$y$ 能等于0吗？

$\dfrac{1}{x} = 0$？

分子是1，永远不是0，所以 $y$ 永远不能等于0。

值域是：\textbf{除了0以外的所有数}。

\textbf{一句话总结：值域就是“$y$ 能等于什么”。定义域管入口，值域管出口。}

\section{第三条规矩：单调性}

这条规矩说的是：

\textbf{当 $x$ 变大的时候，$y$ 是跟着变大，还是跟着变小？}

\subsection{单调递增——一起变大}

\[
y = 2x
\]

$x = 1$ 时，$y = 2$。

$x = 2$ 时，$y = 4$。

$x = 3$ 时，$y = 6$。

$x$ 变大，$y$ 也变大。

你往右走，地势越来越高。

这叫\textbf{单调递增}。

\textbf{一句话：上坡路。}

\subsection{单调递减——反着变}

\[
y = -x
\]

$x = 1$ 时，$y = -1$。

$x = 2$ 时，$y = -2$。

$x = 3$ 时，$y = -3$。

$x$ 变大，$y$ 反而变小。

你往右走，地势越来越低。

这叫\textbf{单调递减}。

\textbf{一句话：下坡路。}

\subsection{有些函数一段上坡一段下坡}

\[
y = x^2
\]

$x = -2$ 时，$y = 4$。

$x = -1$ 时，$y = 1$。

$x = 0$ 时，$y = 0$。

$x = 1$ 时，$y = 1$。

$x = 2$ 时，$y = 4$。

你看：

从 $x = -2$ 走到 $x = 0$，$y$ 从4降到0——\textbf{下坡}。

从 $x = 0$ 走到 $x = 2$，$y$ 从0升到4——\textbf{上坡}。

先下坡，再上坡。

$x = 0$ 那里是最低点，是转折点。

所以我们不能笼统地说“$y = x^2$ 是递增的”或者“是递减的”。

我们得说清楚在哪一段：

\begin{quote}
“$y = x^2$ 在 $x < 0$ 那段是递减的，在 $x > 0$ 那段是递增的。”
\end{quote}

\textbf{一句话总结：单调递增就是上坡，单调递减就是下坡。有些函数一半上坡一半下坡。}

你来判断几个：

\textbf{（1）} $y = 3x + 1$，是递增还是递减？

$x$ 变大，$3x$ 变大，$3x + 1$ 也变大。

答案：\textbf{递增}（全程上坡）。

\textbf{（2）} $y = -2x + 5$，是递增还是递减？

$x$ 变大，$-2x$ 变小（负的变得更负），$-2x + 5$ 也变小。

答案：\textbf{递减}（全程下坡）。

\textbf{（3）} $y = \dfrac{1}{x}$，在 $x > 0$ 那段是递增还是递减？

$x = 1$ 时，$y = 1$。

$x = 2$ 时，$y = 0.5$。

$x$ 变大，$y$ 变小。

答案：\textbf{递减}。

\section{第四条规矩：奇偶性}

这是最后一条规矩，也是最有趣的一条。

\[
y = x^2
\]

你代几个数进去：

$x = 2$，$y = 4$。

$x = -2$，$y = 4$。

一样！

$x = 3$，$y = 9$。

$x = -3$，$y = 9$。

也一样！

正数和对应的负数，代进去，$y$ 值相同。

这种函数，叫\textbf{偶函数}。

\textbf{偶函数的图像，左右对称。}

你沿着 $y$ 轴折过去，左边和右边完全重合。

像照镜子。

$y$ 轴就是那面镜子。

\[
y = x^3
\]

你代几个数进去：

$x = 2$，$y = 8$。

$x = -2$，$y = -8$。

不一样！但是大小相同，符号相反。

$x = 3$，$y = 27$。

$x = -3$，$y = -27$。

还是大小相同，符号相反。

这种函数，叫\textbf{奇函数}。

\textbf{奇函数的图像，关于原点对称。}

什么叫关于原点对称？

你把整个图像绕着原点转180度，图像和原来的完全重合。

右上角有一块，左下角就有一块和它对应。

转半圈，刚好对上。

\textbf{一句话总结：}

\begin{quote}
\textbf{偶函数——照镜子。$y$ 轴是镜子，左右一样。}

\textbf{奇函数——转半圈。绕原点转180度，还是它自己。}
\end{quote}

\subsection{怎么判断一个函数是奇还是偶？}

方法很简单：

\textbf{把 $x$ 换成 $-x$，看看结果变成什么。}

$f(x) = x^2$

把 $x$ 换成 $-x$：

$f(-x) = (-x)^2 = x^2$

和原来的 $f(x) = x^2$ 一模一样。

\textbf{偶函数。}

$f(x) = x^3$

把 $x$ 换成 $-x$：

$f(-x) = (-x)^3 = -x^3$

和原来的 $f(x) = x^3$ 正好差一个负号。

\textbf{奇函数。}

$f(x) = x^2 + x$

把 $x$ 换成 $-x$：

$f(-x) = (-x)^2 + (-x) = x^2 - x$

和原来的 $f(x) = x^2 + x$ 比：

不一样（$+x$ 变成了 $-x$），也不是正好差一个负号。

\textbf{既不是奇函数，也不是偶函数。}

\textbf{一句话总结：把 $x$ 换成 $-x$。结果没变就是偶函数，变成相反数就是奇函数，都不是就是非奇非偶。}

你来判断几个：

\textbf{（1）} $f(x) = x^4$，是奇函数还是偶函数？

$f(-x) = (-x)^4 = x^4 = f(x)$

答案：\textbf{偶函数}。

\textbf{（2）} $f(x) = \dfrac{1}{x}$，是奇函数还是偶函数？

$f(-x) = \dfrac{1}{-x} = -\dfrac{1}{x} = -f(x)$

答案：\textbf{奇函数}。

\textbf{（3）} $f(x) = x + 1$，是奇函数还是偶函数？

$f(-x) = -x + 1$

和 $f(x) = x + 1$ 比：不一样。

和 $-f(x) = -x - 1$ 比：也不一样。

答案：\textbf{非奇非偶}。

\section{四条规矩放在一起}

让我帮你总结一下。

函数有四条规矩，就像四张“身份证”：

\textbf{定义域：$x$ 能填什么。}

分母不能为零，根号下不能为负。

\begin{quote}
\textbf{一句话：入口管制。}
\end{quote}

\textbf{值域：$y$ 能取什么。}

函数能吐出来的所有结果。

\begin{quote}
\textbf{一句话：出口清单。}
\end{quote}

\textbf{单调性：$x$ 变大时 $y$ 怎么变。}

跟着变大就是递增，反着变小就是递减。

\begin{quote}
\textbf{一句话：上坡还是下坡。}
\end{quote}

\textbf{奇偶性：图像对称不对称。}

关于 $y$ 轴对称是偶函数，关于原点对称是奇函数。

\begin{quote}
\textbf{一句话：照镜子还是转半圈。}
\end{quote}

这四条规矩，以后你每次遇到一个新函数，都可以拿出来“审问”它：

你的定义域是什么？

你的值域是什么？

你是递增还是递减？

你是奇函数还是偶函数？

四个问题问完，这个函数你就认识得差不多了。

你现在不只是认识函数了。

你还会给函数“立规矩”——查它的入口、出口、坡度、对称性。

四项指标全查完，这个函数的脾气你就摸透了。

继续往前走吧。

\chapter{微积分——其实就这么点事}

你已经学会画函数的图像了。

一个点一个点算出来，然后连成线。

直线、曲线、两段躲在角落里的曲线，你都认识了。

但你有没有觉得，这样做有点……累？

$x = 1$，算一下 $y$。

$x = 2$，再算一下 $y$。

$x = 3$，又算一下 $y$。

……

一个一个点，算到天荒地老。

数学家也觉得累。

所以他们发明了一套方法，专门用来偷懒。

这套方法，叫\textbf{微积分}。

别被这个名字吓到。

“微积分”听起来像是大学才学的高深东西。

但它的核心，其实很简单。

我用最直白的话告诉你。

\section{第一个概念：极限——无限接近}

我先给你讲一个故事。

有一只乌龟，在你前面10米的地方。

你跑得比乌龟快。

你开始追。

你跑到乌龟原来的位置（10米处），但乌龟已经往前爬了一点，比如1米。

现在乌龟在你前面1米。

你继续追。

你跑到乌龟刚才的位置（11米处），但乌龟又往前爬了一点，比如0.1米。

现在乌龟在你前面0.1米。

你再追。

你跑到乌龟刚才的位置，乌龟又爬了0.01米。

你再追，乌龟又爬了0.001米。

……

每次你追到乌龟刚才的位置，乌龟就又往前爬了一点点。

你们之间的距离越来越小：10米、1米、0.1米、0.01米、0.001米……

但好像永远追不上？

当然，现实中你肯定能追上乌龟。

但这个故事想说的是另一件事：

\textbf{两个东西之间的距离，可以无限缩小，小到几乎为零，但理论上永远差那么一丁点。}

这种“无限接近，但永远差一点点”的感觉，就叫\textbf{极限}。

用数字来感受一下：

$\dfrac{1}{2} = 0.5$

$\dfrac{1}{4} = 0.25$

$\dfrac{1}{8} = 0.125$

$\dfrac{1}{16} = 0.0625$

$\dfrac{1}{32} = 0.03125$

……

分母越来越大，结果越来越小。

无限接近0，但永远不等于0。

我们说：\textbf{当分母趋向于无穷大时，$\dfrac{1}{n}$ 的极限是0。}

\textbf{一句话总结：极限就是“无限接近某个数，但可能永远到不了”。}

\section{第二个概念：$dx$——无限小的一小段}

好，现在你知道什么是“无限接近”了。

我再介绍一个符号：\textbf{$dx$}。

别怕，这个符号很好理解。

想象你面前有一根绳子。

你把它切成两段，每段的长度叫 $\Delta x$（读作“delta x”，就是“一小段 $x$”的意思）。

再切，切成四段。每段更短了。

再切，八段。

再切，十六段。

……

一直切，一直切，切到每一段\textbf{无限短}。

短到你几乎看不见，但它还是存在的。

这种\textbf{无限短的一小段}，就叫 $dx$。

$\Delta x$ 是“一小段”。

$dx$ 是“无限小的一小段”。

$dx$ 是 $\Delta x$ 的极限版本。

\textbf{一句话总结：$dx$ 就是一段无限短、短到几乎为零、但还是存在的小段。}

\section{第三个概念：求导——一个变了，另一个怎么变？}

这是微积分里最重要的概念之一。

它的本质非常简单，但我要讲得很细，你才能真正明白。

\subsection{从一个问题开始}

假设你在跑步。

你跑了10秒，跑了50米。

请问你的速度是多少？

\[
速度 = \frac{距离}{时间} = \frac{50米}{10秒} = 5米/秒
\]

简单。

但这里有一个问题。

这个“5米/秒”是\textbf{平均速度}。

它说的是：在这10秒里，你\textbf{平均}每秒跑5米。

但你在这10秒里，速度真的一直是5米/秒吗？

不一定。

也许刚开始你在热身，跑得慢，每秒只跑3米。

后来你越跑越快，最后每秒跑7米。

平均下来是5米/秒，但每一瞬间的速度其实不一样。

现在我问你一个更难的问题：

\textbf{在第3秒那一瞬间，你的速度到底是多少？}

不是前3秒的平均速度。

是\textbf{恰好第3秒那一刹那}的速度。

\subsection{这个问题为什么难？}

因为“一瞬间”太短了。

一瞬间你跑了多少米？几乎是零。

一瞬间有多长？几乎是零秒。

$\dfrac{0}{0}$？这没法算。

所以我们不能直接算“一瞬间”的速度。

我们得\textbf{绕个弯}。

\subsection{用“很短的一段时间”来逼近}

我们可以这样想：

虽然算不了“一瞬间”，但我可以算“很短的一段时间”。

然后把这段时间\textbf{越缩越短}，看速度会变成什么。

\textbf{第一次尝试：看1秒的时间段}

从第3秒到第4秒，这是1秒钟。

在这1秒里，你跑了多远？

假设你跑了6米。

（为什么是6米？因为你在加速啊。前面平均每秒5米，现在跑到第3秒了，速度快了一点，6米挺合理的。）

那这1秒的平均速度 $= \dfrac{6米}{1秒} = 6米/秒$。

但这是1秒的\textbf{平均速度}，不是第3秒\textbf{那一瞬间}的速度。

这1秒里，你的速度也在变。

我们要更精确，就得把时间段缩短。

\textbf{第二次尝试：看0.1秒的时间段}

从第3秒到第3.1秒，这是0.1秒。

在这0.1秒里，你跑了多远？

假设你跑了0.58米。

（为什么是0.58米？刚才1秒跑6米，那0.1秒大概跑0.6米左右。但因为你在加速，第3秒刚开始的时候速度还没那么快，所以比0.6米少一点点，0.58米差不多。）

这0.1秒的平均速度 $= \dfrac{0.58米}{0.1秒} = 5.8米/秒$。

\textbf{第三次尝试：看0.01秒的时间段}

从第3秒到第3.01秒，这是0.01秒。

假设你跑了0.057米。

（为什么是0.057米？0.1秒跑0.58米，那0.01秒大概跑0.058米左右。但时间越短，速度越接近“刚起步”的那个瞬间速度，所以稍微少一丁点。）

这0.01秒的平均速度 $= \dfrac{0.057米}{0.01秒} = 5.7米/秒$。

\textbf{第四次尝试：继续缩短}

从第3秒到第3.001秒，0.001秒。

假设跑了0.00568米。

平均速度 $= \dfrac{0.00568米}{0.001秒} = 5.68米/秒$。

\subsection{你发现规律了吗？}

\begin{center}
\begin{tabular}{lll}
\toprule
时间段 & 跑的距离 & 平均速度 \\
\midrule
1秒 & 6米 & 6米/秒 \\
0.1秒 & 0.58米 & 5.8米/秒 \\
0.01秒 & 0.057米 & 5.7米/秒 \\
0.001秒 & 0.00568米 & 5.68米/秒 \\
\bottomrule
\end{tabular}
\end{center}

时间段越短，平均速度越来越接近某个数。

如果你继续缩短，0.0001秒、0.00001秒……

平均速度会越来越接近一个\textbf{固定的数}。

当时间段缩到无限短——短到只有“一瞬间”——那个数就是你在\textbf{第3秒那一瞬间}的真实速度。

\subsection{这个过程就叫“求导”}

我们刚才做的事情：

1. 看一小段时间内跑了多少距离。
2. 算这一小段的平均速度。
3. 把时间段越缩越短。
4. 看平均速度趋近于什么数。

这个过程，就叫\textbf{求导}。

算出来的那个数，就叫\textbf{导数}。

\textbf{一句话总结：求导就是把时间段缩到无限短，看“这一瞬间”的变化有多快。}

\subsection{换成能精确计算的例子}

刚才的跑步例子，那些距离（6米、0.58米、0.057米）都是我估的。

你可能会想：这些数到底准不准？

好问题。

现在我们换一个例子，\textbf{每个数都能精确算出来}，不用估。

假设你跑步的规律是这样的：

\textbf{第 $t$ 秒时，你总共跑了 $t^2$ 米。}

为什么是 $t^2$？

因为这是一个简单又有意思的规律——你跑得越久，速度越快，距离增长得越来越多。就像一辆在加速的车。

用这个规律，我们可以精确算出每一个数。

把 $t$ 换成 $x$，把“跑的距离”换成 $y$，就得到一个函数：

\[
y = x^2
\]

现在我们想知道：\textbf{在 $x = 3$ 这一点，$y$ 的变化有多快？}

\textbf{$x$ 从3变到4：}

$y$ 从 $3^2 = 9$ 变到 $4^2 = 16$，变了 $7$。$x$ 变了 $1$。

变化率 $= \dfrac{7}{1} = 7$。

\textbf{$x$ 从3变到3.1：}

$y$ 从 $9$ 变到 $3.1^2 = 9.61$，变了 $0.61$。$x$ 变了 $0.1$。

变化率 $= \dfrac{0.61}{0.1} = 6.1$。

\textbf{$x$ 从3变到3.01：}

$y$ 从 $9$ 变到 $3.01^2 = 9.0601$，变了 $0.0601$。$x$ 变了 $0.01$。

变化率 $= \dfrac{0.0601}{0.01} = 6.01$。

\textbf{$x$ 从3变到3.001：}

$y$ 从 $9$ 变到 $3.001^2 = 9.006001$，变了 $0.006001$。$x$ 变了 $0.001$。

变化率 $= \dfrac{0.006001}{0.001} = 6.001$。

\subsection{现在规律清楚了}

\begin{center}
\begin{tabular}{lll}
\toprule
$x$ 的变化 & $y$ 的变化 & 变化率 \\
\midrule
1 & 7 & 7 \\
0.1 & 0.61 & 6.1 \\
0.01 & 0.0601 & 6.01 \\
0.001 & 0.006001 & 6.001 \\
\bottomrule
\end{tabular}
\end{center}

$x$ 的变化越小，变化率越接近\textbf{6}。

当 $x$ 的变化缩到无限小——也就是 $dx$——的时候：

变化率就\textbf{等于6}。

这个6，就是函数 $y = x^2$ 在 $x = 3$ 这一点的\textbf{导数}。

\subsection{导数在图上是什么意思？}

你画一条曲线 $y = x^2$。

在曲线上找到 $x = 3$ 那个点。

导数告诉你：\textbf{这条曲线在这一点有多陡。}

导数是6，意思是：在这一点，曲线挺陡的。$x$ 往右挪一丁点，$y$ 就往上跳6倍那么多的一丁点。

如果导数是0呢？

那意思是：在这一点，曲线是平的。$x$ 往右挪，$y$ 几乎不变。

$y = x^2$ 在 $x = 0$ 那一点，导数是0。那正好是碗底，是最低点，曲线是平的，不上坡也不下坡。

\subsection{导数有一个公式}

$y = x^2$ 的导数是 $2x$。

（这个公式怎么来的？有一套推导方法。现在你先记住结论就好。）

有了这个公式，你可以快速算出任何一点的导数：

$x = 3$ 时，导数 $= 2 \times 3 = 6$。（和我们刚才一步一步算的一样！）

$x = 1$ 时，导数 $= 2 \times 1 = 2$。

$x = 0$ 时，导数 $= 2 \times 0 = 0$。（碗底，平的。）

$x = 10$ 时，导数 $= 2 \times 10 = 20$。（曲线在这里非常陡！）

\textbf{一句话总结：导数就是“在这一点，曲线有多陡”。或者说，“$x$ 变一丁点，$y$ 会变多少”。}

\subsection{求导的符号}

导数有一个专门的符号，写成：

\[
\frac{dy}{dx}
\]

看起来像一个分数，对吧？

它确实就是一个分数的样子：

$dy = y$ 的变化（无限小的一点点）

$dx = x$ 的变化（无限小的一点点）

$\dfrac{dy}{dx} = y$ 的变化 $\div x$ 的变化

也就是：\textbf{$y$ 变了多少，除以 $x$ 变了多少。}

还有另一种写法：

如果函数叫 $f(x)$，它的导数写成 $f'(x)$。

那个小撇 $'$ 就是“求导”的意思。

$f(x) = x^2$

$f'(x) = 2x$

两种写法，同一个意思。

\textbf{一句话总结：$\dfrac{dy}{dx}$ 就是“$y$ 的变化除以 $x$ 的变化”，当这个变化缩到无限小的时候，算出来的结果就是导数。}

\section{第四个概念：定积分——曲线下面的面积}

好，导数讲完了。

现在来讲积分。

积分和导数是一对。导数是“拆开”，积分是“合起来”。

但我不从这个角度讲。

我从最直观的角度讲：\textbf{面积}。

你面前有一条曲线。

曲线下面，$x$ 轴上面，围成了一块区域。

\textbf{这块区域的面积是多少？}

这就是定积分要算的事情。

如果曲线是一条水平的直线，比如 $y = 3$。

从 $x = 0$ 到 $x = 4$，曲线下面的面积是多少？

那就是一个长方形。

长是4，宽是3，面积 $= 4 \times 3 = 12$。

简单。

但如果曲线是弯的呢？

比如 $y = x^2$，从 $x = 0$ 到 $x = 2$。

曲线下面的区域不是长方形，是一块弯弯的形状。

怎么算面积？

方法：\textbf{切成无数个小长条。}

从 $x = 0$ 到 $x = 2$，切成很多很多小段。

每一小段的宽度是 $dx$（无限小的一段，还记得吗？）。

每一小段的高度，就是那个位置的 $y$ 值。

每一小段的面积 $= $ 高 $\times$ 宽 $= y \times dx$。

然后把所有小长条的面积加起来。

因为小长条有无限多个，所以是“无限多个无限小的东西加起来”。

这就是\textbf{积分}。

\subsection{积分的符号——一个一个拆开看}

积分有一个专门的符号，长这样：

\[
\int_0^2 x^2 \, dx
\]

第一次看，确实有点吓人。

一堆东西挤在一起，不知道在说什么。

别慌。我现在把它\textbf{一个一个拆开}，你就全懂了。

\textbf{第一个：$\int$}

这个弯弯的符号，看起来像一条拉长的线。

它\textbf{就是一条线}。

准确地说，它代表的是\textbf{曲线}——函数画出来的那条线。

但 $\int$ 不只是说“有一条线”。

它说的是：\textbf{这条线下面的那一整块区域}。

你想象一下：

曲线在上面，$x$ 轴在下面。

曲线和 $x$ 轴之间，围成了一块区域。

$\int$ 说的就是这块区域。

线是边界，线下面那一整块，才是我们要算的东西。

同时，$\int$ 还有另一层意思：\textbf{范围}。

哪一段线？从哪里到哪里？

不是整条线，是\textbf{其中一段}。

这个范围，由 $\int$ 上面和下面的两个小数字来指定。

\textbf{一句话总结：$\int$ 就是一条曲线，代表曲线下面那一整块区域，而且指定了一个范围。}

\textbf{第二个：$\int$ 下面的小数字 $0$}

这是\textbf{起点}。

从 $x = 0$ 这个位置开始算。

\textbf{第三个：$\int$ 上面的小数字 $2$}

这是\textbf{终点}。

算到 $x = 2$ 这个位置为止。

把下面和上面合在一起看：

$\int_0^2$ 的意思是：\textbf{从 $x = 0$ 到 $x = 2$ 这一段范围。}

就像你在地图上画了一条路线：从这里出发，到那里结束。

起点是0，终点是2。

这一段范围内的那块区域，就是我们要算的。

\textbf{第四个：$x^2$}

这个你认识。

$x$ 的平方。

在这里，它代表\textbf{曲线的高度}。

你在 $x$ 轴上选一个位置，比如 $x = 1$。

这个位置对应的曲线高度是多少？

$y = x^2 = 1^2 = 1$。

高度是1。

再选一个位置，$x = 2$。

$y = x^2 = 2^2 = 4$。

高度是4。

$x^2$ 告诉你：\textbf{在每一个位置，曲线有多高。}

\textbf{第五个：$dx$}

这个你刚刚学过。

$dx$ 是\textbf{无限小的一段宽度}。

想象你把从0到2这一段 $x$ 轴，切成很多很多小段。

每一小段的宽度，就是 $dx$。

有多小？

无限小。

短到几乎看不见，但它还是存在的。

\textbf{把它们全部组装起来：}

\[
\int_0^2 x^2 \, dx
\]

$\int$ ——一条曲线，以及曲线下面的那块区域

$_0^2$ ——范围是从0到2

$x^2$ ——曲线在每个位置的高度

$dx$ ——每一小段的宽度

合起来的意思是：

\textbf{从 $x = 0$ 到 $x = 2$，曲线下面有一块区域。}

\textbf{这块区域的上边界是曲线 $y = x^2$，下边界是 $x$ 轴。}

\textbf{我们把这块区域切成无数个小长条。}

\textbf{每个小长条的高度是 $x^2$，宽度是 $dx$。}

\textbf{每个小长条的面积 = 高 $\times$ 宽 = $x^2 \times dx$。}

\textbf{把所有小长条的面积加起来，就是这块区域的总面积。}

\subsection{在图上是什么样子？}

让我用文字帮你在脑子里画一幅图。

\textbf{第一步：画坐标系。}

横的是 $x$ 轴，竖的是 $y$ 轴，交叉的地方是原点 $(0, 0)$。

\textbf{第二步：画曲线 $y = x^2$。}

这是一条弯弯的曲线，碗的形状，开口朝上。

碗底在原点。

从原点出发，往右边走，曲线往上升，而且越来越陡。

$x = 1$ 的时候，曲线的高度是 $1$。

$x = 2$ 的时候，曲线的高度是 $4$。

\textbf{第三步：在 $x$ 轴上找到从0到2这一段。}

从原点开始，沿着 $x$ 轴往右走，走到 $x = 2$ 的位置。

\textbf{第四步：看这块区域。}

这段距离的正上方，有一段曲线。

现在你的眼睛盯着这块地方：

上面是那段弯弯的曲线。

下面是 $x$ 轴（平平的一条线）。

左边是 $x = 0$ 那条竖线（就是 $y$ 轴）。

右边是 $x = 2$ 那条竖线。

这四条边围起来，形成了一块区域。

上面那条边是弯的（曲线），其他三条边是直的。

像一块被压扁的馒头，或者一块靠在墙角的软垫子。

\textbf{第五步：切成小长条。}

想象你拿一把刀，从左往右，竖着切。

咔，切一刀。

咔，再切一刀。

咔咔咔咔咔……切无数刀。

每一刀之间的距离，就是 $dx$——无限薄的一小段。

切完之后，这块区域变成了\textbf{无数个竖着的小长条}，紧紧挨在一起。

\textbf{第六步：看其中一个小长条。}

随便挑一个，比如在 $x = 1$ 这个位置的小长条。

站在 $x$ 轴上，底部在 $x = 1$ 的位置。

顶部刚好碰到曲线，高度是 $y = 1^2 = 1$。

宽度是 $dx$，无限薄。

这个小长条的面积 $= $ 高 $\times$ 宽 $= 1 \times dx$。

再挑一个，在 $x = 1.5$ 那个位置：高度是 $1.5^2 = 2.25$，面积 $= 2.25 \times dx$。

再挑一个，在 $x = 2$ 那个位置：高度是 $2^2 = 4$，面积 $= 4 \times dx$。

\textbf{第七步：全部加起来。}

从 $x = 0$ 到 $x = 2$，每一个位置都有一个小长条。

有无限多个小长条，紧紧挨在一起，填满了整块区域。

每个小长条的面积 $= x^2 \times dx$。

把所有小长条的面积加起来：

\[
\int_0^2 x^2 \, dx
\]

加出来的结果，就是那块弯弯的区域的\textbf{总面积}。

\textbf{一句话总结：积分就是把曲线下面那块区域，切成无数个无限薄的小长条，然后把所有小长条的面积加起来。}

不是什么神秘的东西。

就是一块区域的面积。

只不过这块区域的上边是弯的，没法直接用“长 $\times$ 宽”算，所以要用积分——切成无数个小长条，一个一个加起来。

\section{导数和积分是一对}

最后说一件神奇的事。

导数和积分，是\textbf{反过来的操作}。

导数是什么？

你知道“走了多远”（路程），求“这一瞬间有多快”（速度）。

积分是什么？

你知道“每一瞬间有多快”（速度），求“总共走了多远”（路程）。

$y = x^2$ 的导数是 $2x$。

$2x$ 的积分是 $x^2$。

（严格说是 $x^2 + C$，那个 $C$ 是一个常数，以后再说。）

一个是“拆开”，一个是“装回去”。

做了导数再做积分，回到原来的样子。

做了积分再做导数，也回到原来的样子。

它们是一对。

\textbf{一句话总结：导数把东西“拆开”，积分把东西“装回去”。它们互为反操作。}

\section{微积分到底是什么？}

让我帮你总结一下。

\textbf{极限：} 无限接近某个东西。追乌龟，越来越近，但好像永远差一点点。

\textbf{$dx$：} 无限小的一段。短到几乎为零，但还是存在。

\textbf{导数（求导）：} 在某一点，曲线有多陡。$x$ 变一丁点，$y$ 变多少。

\textbf{积分（定积分）：} 曲线下面的面积。把无限多个无限小的东西加起来。

这就是微积分的全部核心。

当然，真正用的时候，还有很多具体的计算方法和公式。

但那些都是技术细节。

核心思想，就是这四样东西。

\textbf{无限小，无限多，无限接近。}

用这三个“无限”，去描述那些普通方法算不了的东西。

这就是微积分的本质。

你现在知道微积分在说什么了。

它不是什么高不可攀的东西。

它就是在问：

\begin{quote}
这一点有多陡？

这一段面积多大？

无限接近的时候，会发生什么？
\end{quote}

以后你真正学微积分的时候，会学很多具体的计算方法。

但你已经知道了它的灵魂。

有了灵魂，后面那些计算，只是手艺活。

\chapter{在坐标系里画画——从直线到曲线}

坐标系你已经见过了。

横的是 $x$ 轴，竖的是 $y$ 轴，交叉的地方是原点 $(0, 0)$。

网格上的每一个点，都有一个坐标 $(x, y)$。

前几章我们在坐标系里画了直线、画了曲线，看了函数的长相。

但我们一直在“看”，没有“算”。

这一章，我们来\textbf{算}。

算什么？

算一条线有多陡。

算一条线的方程怎么写。

算两条线碰在一起会发生什么。

最后，算那些弯弯的形状——圆、椭圆、抛物线、双曲线——到底是怎么回事。

\section{斜率：竖着走多远，除以横着走多远}

在一条直线上，随便挑两个点：

\textbf{起点}叫 $A(x_1, y_1)$。

\textbf{终点}叫 $B(x_2, y_2)$。

你想从 $A$ 走到 $B$，你需要两步：

1. \textbf{横着走多远？} 终点的 $x$ 减去起点的 $x$，也就是 $x_2 - x_1$。
2. \textbf{竖着走多远？} 终点的 $y$ 减去起点的 $y$，也就是 $y_2 - y_1$。

斜率 $k$ 就是：

\[
k = \frac{y_2 - y_1}{x_2 - x_1}
\]

大白话：\textbf{竖着走的距离，除以横着走的距离。}

\subsection{这个数字到底在说什么？}

斜率不是一个抽象的东西。它在告诉你一个非常具体的信息：

\textbf{你往右走一格，需要往上走几格。}

\textbf{算出来 $k = 3$：}

你往右走1格，得往上爬3格。

非常陡。像爬楼梯。

\textbf{算出来 $k = 0.5$：}

你往右走1格，只需要往上爬半格。

很缓。像走一个慢慢上升的坡道。

\textbf{算出来 $k = -2$：}

负号意味着方向反了——不是往上爬，是往下掉。

你往右走1格，会往下掉2格。

下坡。

\textbf{算出来 $k = 0$：}

竖着的变化是0。不管你往右走多少格，高度一动不动。

这是一条水平线。平地。

\textbf{算出来分母是0（$x_2 - x_1 = 0$）：}

分母是0，数学里不允许除以0。斜率不存在。

这意味着什么？$x_2 = x_1$，两个点的横坐标一样，它们上下排列。

这是一条竖直的线。像一堵墙。

竖直的线没有斜率——它不是“无限陡”，而是\textbf{压根不适用“往右走一格升多少格”这个说法}，因为你根本没有往右走。

\subsection{总结一下}

\begin{center}
\begin{tabular}{ll}
\toprule
斜率 $k$ & 含义 \\
\midrule
$k > 0$ & 上坡。$k$ 越大越陡 \\
$k < 0$ & 下坡。$|k|$ 越大越陡 \\
$k = 0$ & 平地。水平线 \\
$k$ 不存在 & 竖直的墙。分母为0 \\
\bottomrule
\end{tabular}
\end{center}

\textbf{斜率的大小决定陡不陡，正负号决定上坡还是下坡。}

\subsection{一个重要的小问题：能不能反过来减？}

公式是 $k = \dfrac{y_2 - y_1}{x_2 - x_1}$。

那如果你反过来，用 $\dfrac{y_1 - y_2}{x_1 - x_2}$ 呢？

\textbf{完全可以。算出来一模一样。}

为什么？

因为如果你把 $A$ 当起点、$B$ 当终点，是“向右向上走”；反过来把 $B$ 当起点、$A$ 当终点，就是“向左向下走”。

分子变成了负数，分母也变成了负数。负负得正，结果不变。

\textbf{只要上下保持同步——要么都是“终点减起点”，要么都是“起点减终点”——谁减谁都行。}

千万不能上面用 $y_2 - y_1$，下面用 $x_1 - x_2$，那就乱了。

\section{从斜率变出直线方程}

直线有一个特点：\textbf{不管你在上面挑哪两个点，算出来的斜率 $k$ 永远一样。}

这是直线之所以“直”的原因——坡度始终不变。

现在，我们把 $A(x_1, y_1)$ 当作一个\textbf{固定的点}，钉在坐标系上不动。

然后在直线上放一个\textbf{随便乱跑的点}，叫它 $(x, y)$。

既然斜率永远不变，那这个乱跑的点和固定点之间的斜率，也得等于 $k$：

\[
k = \frac{y - y_1}{x - x_1}
\]

现在，把分母 $(x - x_1)$ 乘到等号左边去：

\[
y - y_1 = k(x - x_1)
\]

你看——\textbf{这就是高中天天见的“点斜式方程”。}

它根本不是什么新发明。

它就是把斜率公式的分母挪了个位置。就这样。

\subsection{$y = kx + b$ 里的两个数字}

如果你把点斜式展开、整理，最后通常会变成这个样子：

\[
y = kx + b
\]

这里面有两个关键数字：$k$ 和 $b$。

\textbf{$k$ 你已经认识了——斜率。每往右走一格，升多少格。}

\textbf{$b$ 是什么？}

看公式：当 $x = 0$ 时，$y = k \times 0 + b = b$。

也就是说，$b$ 是这条直线\textbf{穿过 $y$ 轴的那个位置}。

$b = 5$，意思是直线在 $y$ 轴上的高度是5。

$b = -3$，意思是直线在 $y$ 轴上的高度是 $-3$（在 $x$ 轴下面）。

打个比方。

想象你打车。

$b$ 是\textbf{起步价}。哪怕你1米都没跑（$x = 0$），计价器上已经是 $b$ 块钱了。

$k$ 是\textbf{每公里的价格}。你每多跑1公里（$x$ 增加1），计价器往上跳 $k$ 块。

最后算出来的 $y$，就是总车费。

\[
\text{总车费} = \text{每公里价格} \times \text{跑了几公里} + \text{起步价}
\]

\[
y = kx + b
\]

直线方程就是这么现实的东西。

\section{两条线碰面：平行和垂直}

\subsection{平行——步伐一模一样}

两条直线平行，意思是它们永远不会相交。

为什么不会相交？

因为它们的\textbf{坡度完全一样}。

你想象两条铁轨。它们始终保持同样的间距，因为每一步的升降幅度一模一样。

\textbf{两条直线平行 $\leftrightarrow$ 斜率相等。$k_1 = k_2$。}

\subsection{垂直——斜率相乘等于 $-1$}

这个结论你可能听过，但从来没人跟你解释\textbf{为什么}。

我现在就解释。

假设在坐标系里有一条过原点的直线 $L_1$。

在 $L_1$ 上找一个点，坐标叫 $(a, b)$。

那 $L_1$ 的斜率是：

\[
k_1 = \frac{b}{a}
\]

现在我要画一条跟 $L_1$ 垂直的直线 $L_2$，也过原点。

怎么找 $L_2$ 上的点？

把点 $(a, b)$ 绕着原点，\textbf{逆时针旋转90度}。

在坐标系里，一个点 $(a, b)$ 旋转90度之后，坐标会变成 $(-b, a)$。

（不信你可以试试：点 $(3, 2)$ 转90度，会跑到 $(-2, 3)$ 的位置。你在纸上画一下，真的是这样。）

那 $L_2$ 的斜率就是：

\[
k_2 = \frac{a}{-b} = -\frac{a}{b}
\]

把两个斜率乘在一起：

\[
k_1 \times k_2 = \frac{b}{a} \times \left(-\frac{a}{b}\right) = -1
\]

就这样。

不需要任何复杂的定理。

仅仅因为\textbf{旋转90度会让 $x$ 和 $y$ 互换位置、并且改变一个符号}，就注定了垂直线的斜率相乘必须等于 $-1$。

\textbf{乘积等于 $-1$，到底在说什么？}

说的是：一条线往右走1格升2格（$k_1 = 2$），那跟它垂直的线就必须往右走1格\textbf{降}半格（$k_2 = -0.5$）。

原来的“横向步数”和“纵向步数”彻底颠倒了（变成了倒数），方向也反了（加了负号）。

$2 \times (-0.5) = -1$。

\textbf{$-1$ 就是“完美垂直”的防伪标记。}

\section{斜率和角度}

一条直线和水平的 $x$ 轴之间，会形成一个夹角 $\theta$。

你从 $(x_1, y_1)$ 走到 $(x_2, y_2)$，在坐标系里画出了一个直角三角形：

竖边（高）$= y_2 - y_1$

横边（底）$= x_2 - x_1$

在三角函数里，$\tan \theta$ 的定义就是“对边除以邻边”，也就是“高除以底”。

所以：

\[
k = \tan \theta = \frac{y_2 - y_1}{x_2 - x_1}
\]

斜率和角度，原来是同一件事。只是一个用数字说，一个用度数说。

\textbf{$k = 1$ 的时候：} $\tan \theta = 1$，$\theta = 45^\circ$。往右走1格，往上走1格，刚好是对角线。

\textbf{$k$ 非常大的时候：} 角度接近90度。线几乎竖起来了，但只要 $k$ 还是个具体的数字，它就没有完全竖直。

\textbf{$k = 0$ 的时候：} $\theta = 0^\circ$。完全水平。

\section{从直线到曲线}

直线的斜率到处都一样——坡度始终不变。

但如果线是弯的呢？

曲线的坡度在每个点都不一样。

怎么办？

还是用 $\dfrac{y_2 - y_1}{x_2 - x_1}$。

只不过，你把两个点挑得\textbf{无限接近}。

比如，曲线 $y = x^2$，你想知道在 $x = 2$（也就是点 $(2, 4)$）的位置有多陡。

在曲线上找一个离它很近的点，比如 $(2.001,\ 2.001^2) = (2.001,\ 4.004001)$。

\[
k = \frac{4.004001 - 4}{2.001 - 2} = \frac{0.004001}{0.001} = 4.001
\]

非常接近 $4$。

如果让两个点无限接近，算出来就是\textbf{刚好等于 $4$}。

还记得第九章讲的吗？

这个过程就是\textbf{求导}。

\[
\frac{dy}{dx}
\]

它本质上就是你最开始学的那个老老实实的 $\dfrac{y_2 - y_1}{x_2 - x_1}$。

只不过两个点靠得无限近、变化无限小。

所有的微积分，都是从\textbf{两个坐标相减}开始的。

\section{坐标系里的四种形状}

好，斜率和直线的事情讲完了。

现在来看点更有意思的——\textbf{弯曲的形状}。

直线的方程里，$x$ 和 $y$ 都是“光秃秃”的，没有平方：$y = kx + b$。

画出来永远是一条笔直的线。

但如果方程里出现了\textbf{平方}——$x^2$ 或 $y^2$——画出来的线就开始弯了。

接下来我要介绍四种形状。

它们的方程看起来有点吓人，但你不需要死记硬背。

你只要理解一件事：

\textbf{每个方程都是一条“准入规则”。坐标系上有无数个点 $(x, y)$，谁的坐标代进去能让等式成立，谁就有资格站在这条线上。}

像一个俱乐部。满足条件的点才能进门。

\subsection{第一种：圆——到中心的距离必须相等}

在网格上定一个中心点 $(x_1, y_1)$。

你想在周围画一个圆，半径是 $r$。

圆的本质规则是：\textbf{所有点到中心的距离，必须刚好等于 $r$。}

怎么算距离？

用勾股定理。（还记得吗？直角三角形的两条直角边的平方加起来，等于斜边的平方。）

从点 $(x, y)$ 到中心 $(x_1, y_1)$：

横着走了 $x - x_1$

竖着走了 $y - y_1$

距离的平方 = 横的平方 + 竖的平方。

所以：

\[
(x - x_1)^2 + (y - y_1)^2 = r^2
\]

这就是圆的方程。

\textbf{$r^2$ 在说什么？}

它规定了“结界”的大小。

如果 $r^2 = 25$，半径就是 $\sqrt{25} = 5$。

你拿网格上任何一个点 $(x, y)$ 代进去：

算出来\textbf{等于}25：你刚好踩在圆的边上。

算出来\textbf{小于}25：你离中心太近了，在圆的里面。

算出来\textbf{大于}25：你离中心太远了，在圆的外面。

\textbf{为什么圆是完美对称的？}

因为平方。

$(x - x_1)^2$——不管你往左偏还是往右偏，平方之后都是正数，大小一样。

$(y - y_1)^2$——不管你往上偏还是往下偏，也一样。

\textbf{平方把正负号吃掉了。} 所以圆往哪个方向看，都长得一样。

\subsection{第二种：椭圆——被分配了“预算”的圆}

如果你把圆的公式稍微改一下，让 $x$ 和 $y$ 不再平起平坐，圆就被压扁了。

\[
\frac{x^2}{a^2} + \frac{y^2}{b^2} = 1
\]

（为了简单，这里把中心放在原点。）

你看，$x^2$ 和 $y^2$ 分别除以了两个不同的数（$a^2$ 和 $b^2$），然后要求它们\textbf{加起来等于1}。

这个“1”是什么？

你可以把它想成\textbf{100\%的预算}。

想象你有一块钱（$= 1$）。

你要把这一块钱分给 $x$ 方向和 $y$ 方向。

$\dfrac{x^2}{a^2}$ 花掉了一部分。

$\dfrac{y^2}{b^2}$ 花掉了剩下的部分。

两部分加起来，刚好把一块钱花光。

这意味着什么？

如果你把 $x$ 走到极限（$x = a$），让 $\dfrac{x^2}{a^2}$ 等于1，那预算花光了，$\dfrac{y^2}{b^2}$ 必须等于0，$y$ 只能是0——一步都不能动。

反过来也一样。

$x$ 和 $y$ 互相牵制，互相抢预算。

这种互相牵制的结果，就是画出来的线\textbf{必须闭合成一个圈}。

\textbf{$a$ 和 $b$ 在说什么？}

它们是各方向的“最远距离”。

如果 $a = 5$，$b = 3$，意思是横向最多走5格，纵向最多走3格。

画出来就是一个横着长、竖着短的扁圆。

如果 $a = b$，两个方向的预算一样——椭圆就变回圆了。

\subsection{第三种：抛物线——只有一方有平方}

圆和椭圆之所以能闭合成圈，是因为 $x$ 和 $y$ \textbf{都有平方}，互相拉扯。

但如果只有一个有平方，另一个没有呢？

那就拉不住了。

\[
y - y_1 = a(x - x_1)^2
\]

你仔细看——这和直线的点斜式 $y - y_1 = k(x - x_1)$ 长得很像。

唯一的区别：右边多了个\textbf{平方}。

就是这个平方，让直线变成了曲线。

\textbf{$(x_1, y_1)$ 在说什么？}

因为 $(x - x_1)^2$ 永远大于等于0，所以当 $x = x_1$ 时，平方项变成0，$y$ 等于 $y_1$。

这个点就是抛物线的\textbf{顶点}——U形碗的碗底（或者倒扣碗的碗顶）。

\textbf{$a$ 在说什么？}

它是放大器。

如果 $a = 1$：$x$ 离开顶点走1格，$y$ 升 $1^2 = 1$ 格。走3格，升 $3^2 = 9$ 格。

如果 $a = 5$：$x$ 走3格，$y$ 升 $5 \times 9 = 45$ 格。

\textbf{$a$ 越大，抛物线越窄越陡。$a$ 越小，越宽越平缓。}

如果 $a$ 是负数，抛物线开口朝下——碗翻过来了。

\subsection{第四种：双曲线——加号变减号，天翻地覆}

回到椭圆的公式：$\dfrac{x^2}{a^2} + \dfrac{y^2}{b^2} = 1$。

椭圆是“相加等于1”，$x$ 和 $y$ 互相妥协。

但如果你\textbf{把加号变成减号}呢？

\[
\frac{x^2}{a^2} - \frac{y^2}{b^2} = 1
\]

这一个小小的改动，后果是灾难性的（对闭合来说）。

相减等于1，意味着如果 $x$ 变得非常大（比如一万），$\dfrac{x^2}{a^2}$ 变得巨大。

为了让减完之后还等于1，$\dfrac{y^2}{b^2}$ 也得变得巨大——跟着一起往外跑。

$x$ 往外跑，$y$ 也跟着往外跑。

\textbf{再也闭合不了了。}

画出来就是两条朝着无限远狂奔的曲线。

这就是\textbf{双曲线}——“双”是因为它有两段，分居在坐标系的两侧。

\textbf{还有一个有趣的细节。}

如果你让 $x = 0$（想去坐标系的正中间看看），代入公式：

\[
-\frac{y^2}{b^2} = 1
\]

左边是负数，右边是正数。

\textbf{不可能成立。}

这意味着：双曲线的中间有一块\textbf{空地}。

曲线在左边有一段，右边有一段，中间被撕开了，怎么也连不起来。

就是这个减号制造了这道鸿沟。

\section{看符号，猜形状}

到这里，四种形状你都认识了。

我帮你总结一个规律——以后看到方程，不用画图，直接看\textbf{平方的组合}，就能猜出形状：

\begin{center}
\begin{tabular}{ll}
\toprule
方程特征 & 形状 \\
\midrule
没有平方（$x$ 和 $y$ 都是一次） & \textbf{直线} \\
只有一个有平方（比如 $y = x^2$） & \textbf{抛物线}——一方平稳，一方狂飙 \\
两个都有平方，中间是加号 & \textbf{圆或椭圆}——互相妥协，闭合成圈 \\
两个都有平方，中间是减号 & \textbf{双曲线}——互相排斥，中间撕裂 \\
\bottomrule
\end{tabular}
\end{center}

坐标系里的所有形状，本质上就是\textbf{这几个平方项在互相博弈}的结果。

加号让它们互相拉扯、闭合。

减号让它们互相排斥、撕裂。

有平方就弯，没平方就直。

就这么简单。

这一章的内容有点多。

但你已经走完了。

从最基础的“两个坐标相减”出发，你一路走到了直线方程、平行和垂直、角度、曲线斜率、四种形状。

所有这些看起来不同的东西，底层都是同一件事：

\textbf{$x$ 变了多少，$y$ 变了多少，它们之间是什么关系。}

这就是坐标几何的全部秘密。

好了，歇口气，准备好了咱们接着走。

\chapter{乘法、加法，就只有单一？}

你已经会加法了。

$1 + 2 = 3$

$3 + 4 + 5 = 12$

你也会乘法了。

$2 \times 3 = 6$

$2 \times 3 \times 4 = 24$

但如果我让你算这个呢？

\[
1 + 2 + 3 + 4 + 5 + 6 + 7 + 8 + 9 + 10 + \cdots + 100
\]

从1一直加到100。

你要写100个数字，加99个加号。

太长了。

数学家最讨厌写长句子。

所以他们发明了一种\textbf{缩写}。

\section{连加符号：$\sum$}

\[
\sum_{i=1}^{100} i
\]

这个符号，读作“sigma”（西格玛），是希腊字母，长得像一个歪歪的E。

别被它吓到。我来一个一个拆给你看。

\subsection{$\sum$ 是什么意思？}

$\sum$ 就是“\textbf{全部加起来}”的意思。

它是英文单词“Sum”（求和）的缩写，只不过用了希腊字母。

看到 $\sum$，你脑子里就想：\textbf{“加加加，加一大堆。”}

\subsection{下面那个 $i=1$ 是什么意思？}

这是\textbf{起点}。

从 $i = 1$ 开始加。

\subsection{上面那个 $100$ 是什么意思？}

这是\textbf{终点}。

加到 $i = 100$ 为止。

\subsection{右边那个 $i$ 是什么意思？}

这是\textbf{每一项是什么}。

在这里，每一项就是 $i$ 本身。

$i = 1$ 时，这一项是1。

$i = 2$ 时，这一项是2。

$i = 3$ 时，这一项是3。

……

$i = 100$ 时，这一项是100。

\subsection{全部组装起来}

\[
\sum_{i=1}^{100} i
\]

意思是：

\textbf{从 $i = 1$ 开始，到 $i = 100$ 结束，把每一个 $i$ 都加起来。}

也就是：

\[
1 + 2 + 3 + 4 + 5 + \cdots + 100
\]

一行符号，代替了一百个数字。

这就是数学家偷懒的方式。

\subsection{再来几个例子}

\[
\sum_{i=1}^{5} i
\]

从1加到5。

$= 1 + 2 + 3 + 4 + 5 = 15$

\[
\sum_{i=1}^{4} i^2
\]

从1到4，每一项是 $i^2$（$i$ 的平方）。

$= 1^2 + 2^2 + 3^2 + 4^2$

$= 1 + 4 + 9 + 16$

$= 30$

\[
\sum_{i=3}^{6} i
\]

从3加到6。

$= 3 + 4 + 5 + 6 = 18$

（起点不一定是1，可以从任何地方开始。）

\[
\sum_{i=1}^{3} 5
\]

从1到3，每一项都是5。

$= 5 + 5 + 5 = 15$

（每一项不一定要和 $i$ 有关。这里每一项都是5，加了3次。）

\textbf{一句话总结：$\sum$ 就是“从哪里开始，到哪里结束，把所有项加起来”的缩写。}

\section{连乘符号：$\prod$}

加法有缩写，乘法当然也有。

\[
\prod_{i=1}^{5} i
\]

这个符号，读作“pi”（派），是希腊字母，长得像一个门框。

（注意：这个“派”和圆周率 $\pi$ 是同一个字母，但意思不一样。这里是连乘，圆周率是另一回事。别搞混了。）

\subsection{$\prod$ 是什么意思？}

$\prod$ 就是“\textbf{全部乘起来}”的意思。

它是英文单词“Product”（乘积）的缩写。

看到 $\prod$，你脑子里就想：\textbf{“乘乘乘，乘一大堆。”}

\subsection{和 $\sum$ 的结构一模一样}

\[
\prod_{i=1}^{5} i
\]

下面的 $i = 1$：起点，从1开始。

上面的 $5$：终点，乘到5为止。

右边的 $i$：每一项是什么。

意思是：

\[
1 \times 2 \times 3 \times 4 \times 5 = 120
\]

\subsection{再来几个例子}

\[
\prod_{i=1}^{3} i
\]

$= 1 \times 2 \times 3 = 6$

\[
\prod_{i=2}^{4} i
\]

从2乘到4。

$= 2 \times 3 \times 4 = 24$

\[
\prod_{i=1}^{4} 2
\]

从1到4，每一项都是2。

$= 2 \times 2 \times 2 \times 2 = 16$

（乘了4次2，就是 $2^4$。）

\textbf{一句话总结：$\prod$ 就是“从哪里开始，到哪里结束，把所有项乘起来”的缩写。}

\section{阶乘：$n!$}

现在来一个特别常用的东西。

\[
5!
\]

这个感叹号，不是在喊叫，是一个数学符号。

它叫\textbf{阶乘}。

\subsection{$5!$ 是什么意思？}

\[
5! = 5 \times 4 \times 3 \times 2 \times 1 = 120
\]

\textbf{从这个数开始，一直乘到1。}

就这么简单。

\subsection{为什么叫“阶乘”？}

你想象一个楼梯。

5层楼梯。

从第5层开始，一层一层往下走：5、4、3、2、1。

每走一层，就乘一次。

“阶”就是台阶，“乘”就是乘法。

\textbf{阶乘 = 一层一层往下乘。}

\subsection{几个例子}

\[
1! = 1
\]

\[
2! = 2 \times 1 = 2
\]

\[
3! = 3 \times 2 \times 1 = 6
\]

\[
4! = 4 \times 3 \times 2 \times 1 = 24
\]

\[
5! = 5 \times 4 \times 3 \times 2 \times 1 = 120
\]

\[
6! = 6 \times 5 \times 4 \times 3 \times 2 \times 1 = 720
\]

你发现了吗？数字越大，阶乘增长得\textbf{非常快}。

$10! = 3628800$（三百多万）

$20!$ 已经是一个19位数了。

\subsection{一个特殊的规定：$0! = 1$}

\[
0! = 1
\]

你可能觉得奇怪：0乘到1？0下面没有台阶了，怎么乘？

这是一个\textbf{规定}。

数学家规定 $0! = 1$，是为了让很多公式能够统一、简洁。

你现在不需要理解为什么，先记住就好：$0! = 1$。

\subsection{阶乘和连乘的关系}

\[
n! = \prod_{i=1}^{n} i
\]

阶乘其实就是连乘的一种特殊情况：

从1乘到 $n$。

$5! = \prod_{i=1}^{5} i = 1 \times 2 \times 3 \times 4 \times 5 = 120$

两种写法，同一个意思。

阶乘只是更短的缩写。

\textbf{一句话总结：$n!$ 就是从 $n$ 一直乘到1，一层一层往下乘。}

\section{阶乘有什么用？}

你可能想问：这东西有什么用？为什么要学它？

用处大了。

\subsection{用处一：排列}

假设你有3本书：A、B、C。

你要把它们排成一排，放在书架上。

有多少种排法？

第一个位置，可以放A、B、C中的任意一本——3种选择。

第二个位置，剩下2本——2种选择。

第三个位置，剩下1本——1种选择。

总共：$3 \times 2 \times 1 = 6$ 种排法。

这就是 $3!$。

如果有5本书呢？

$5! = 120$ 种排法。

如果有10本书呢？

$10! = 3628800$ 种排法。

\textbf{排列的数量 = 阶乘。}

这就是阶乘最常见的用处。

\subsection{用处二：组合}

这个稍微复杂一点，你以后会学到。

现在只要知道：\textbf{阶乘是排列组合的基础}，很多公式都会用到它。

\section{你来试试}

\textbf{（1）} $\sum_{i=1}^{4} i = ?$

$= 1 + 2 + 3 + 4 = 10$

\textbf{（2）} $\sum_{i=2}^{5} i^2 = ?$

$= 2^2 + 3^2 + 4^2 + 5^2$

$= 4 + 9 + 16 + 25 = 54$

\textbf{（3）} $\prod_{i=1}^{4} i = ?$

$= 1 \times 2 \times 3 \times 4 = 24$

\textbf{（4）} $4! = ?$

$= 4 \times 3 \times 2 \times 1 = 24$

\textbf{（5）} $7! \div 5! = ?$

$7! = 7 \times 6 \times 5 \times 4 \times 3 \times 2 \times 1$

$5! = 5 \times 4 \times 3 \times 2 \times 1$

\[
\frac{7!}{5!} = \frac{7 \times 6 \times 5!}{5!} = 7 \times 6 = 42
\]

（你看，$5!$ 可以约掉。这是阶乘计算的一个常用技巧。）

\section{这一章的灵魂}

加法和乘法，你从小就会。

但当数字多起来的时候，一个一个写太累了。

数学家发明了三种缩写：

\textbf{$\sum$（连加）：} 把一堆数加起来。

\textbf{$\prod$（连乘）：} 把一堆数乘起来。

\textbf{$n!$（阶乘）：} 从 $n$ 乘到1。

这些符号看起来吓人，但它们做的事情你早就会了。

它们只是让你\textbf{写得更短}。

就像“等等”可以缩写成“etc.”一样，$1 + 2 + 3 + \cdots + 100$ 可以缩写成 $\sum_{i=1}^{100} i$。

\textbf{符号是工具，不是障碍。}

认识了这些工具，以后读公式就不会被吓到了。

你现在认识了三种“打包”符号。

以后在概率、数列、微积分里，它们会反复出现。

而你已经知道它们是什么意思了。

带着这些工具，咱们接着往前。

\chapter{补充章：$\log$ 到底是个什么怪物？其实它只是个记账员}

你在数学书上，可能见过一个长得像英文字母的符号：

\[
\log
\]

很多人第一次看到它，都会心里一紧：“这不是数学课吗？怎么突然冒出英文单词了？”

而且它的旁边总是跟着一大一小两个数字，长得奇奇怪怪的，像这样：

\[
\log_2 8 = 3
\]

别怕。

它看起来像天书，但它的本质，其实是一个\textbf{极其简单的连闯关游戏}。

今天，我们就来破解这个游戏。

\section{一个“魔法克隆箱”的游戏}

想象一下，你手里有一个神奇的“魔法克隆箱”。

这个箱子有一个死规矩：\textbf{你最开始只能往里面放 1 个苹果。}

箱子外面有一个红色的按钮。

每次你按下按钮，箱子里的苹果就会发生克隆。

假设，这是一台\textbf{“2倍克隆箱”}。

意思是，按一次按钮，里面的苹果数量就会变成原来的 \textbf{2倍}。

好，游戏开始：

箱子里一开始有 1 个苹果。

你按 \textbf{1次} 按钮：1个变成 \textbf{2} 个。

你按 \textbf{2次} 按钮：2个变成 \textbf{4} 个。

你按 \textbf{3次} 按钮：4个变成 \textbf{8} 个。

你看，按了3次按钮之后，箱子里有了8个苹果。

这个过程，如果你用数学符号写出来，就是你以前学过的“次方”：

\[
2^3 = 8
\]

（2的3次方等于8。意思是：每次变2倍，连着变3次，最后结果是8。）

\section{倒过来看问题}

“次方”解决的是向前看的问题：\textbf{“我按了3次按钮，最后会得到几个？”}

但是，生活里我们经常会遇到\textbf{倒过来}的问题。

比如，妈妈对你说：

\begin{quote}
“今天家里要来客人，我需要 \textbf{16个} 苹果来做派。”
\end{quote}

现在，你的目标是 16。

你的问题变成了：

\begin{quote}
\textbf{“为了得到 16 个苹果，我到底需要按几次按钮？”}
\end{quote}

你在心里默念算一下：

按1次是2，按2次是4，按3次是8，按4次是16……

算出来了！\textbf{需要按 4 次！}

你刚才脑子里想的这个“算按了几次”的过程，数学家给它起了一个高大上的名字，就叫——\textbf{对数（Logarithm）}。

\section{翻译 $\log$ 的符号}

数学家非常懒，他们嫌“为了得到16个苹果，这台2倍克隆箱需要按几次按钮”这句话太长了。

于是，他们发明了一个缩写：

\[
\log_2 16 = 4
\]

就这么短短一行字，把刚才那一长串话全说清楚了。

来，我教你怎么“逐字翻译”这个符号：

\textbf{$\log$} —— 这是一个提问的单词，意思是：\textbf{“要按几次按钮？”}

\textbf{右下角那个小小的 $2$} —— 这是告诉你机器的规矩：\textbf{“每次变 2 倍。”}

\textbf{旁边那个大大的 $16$} —— 这是你的目标：\textbf{“最后要变成 16 个。”}

\textbf{等号右边的 $4$} —— 这是算出来的答案：\textbf{“需要按 4 次。”}

连起来，用大白话直接翻译就是：

\begin{quote}
\textbf{“如果规矩是每次变 2 倍，那么想要变成 16，需要按几次按钮？”}

\textbf{答案是：4 次。}
\end{quote}

你看！$\log$ 根本不是什么怪物。

它就是一个拿着小本本的记账员，在旁边认真帮你数着\textbf{“到底折腾了多少次”}。

\section{试着玩几局}

现在，你已经拿到了这门语言的翻译器。我们来实战演练几局。

\subsection{第一局}

\[
\log_3 9 = ?
\]

\textbf{直接翻译：}

如果每次变 3 倍，想要变成 9，需要按几次按钮？

\textbf{你想一想：}

一开始是1。

按1次，变成3。

按2次，变成 $3 \times 3 = 9$。到了！

\textbf{答案：} $2$。

所以，$\log_3 9 = 2$。

\subsection{第二局}

\[
\log_5 125 = ?
\]

\textbf{直接翻译：}

每次变 5 倍，想要变成 125，需要按几次？

\textbf{你想一想：}

按1次是 5。

按2次是 $5 \times 5 = 25$。

按3次是 $25 \times 5 = 125$。

\textbf{答案：} $3$。

所以，$\log_5 125 = 3$。

\subsection{第三局（脑筋急转弯）}

\[
\log_7 7 = ?
\]

\textbf{直接翻译：}

每次变 7 倍，想要变成 7，需要按几次？

\textbf{你想一想：}

箱子里一开始是1。我按1次，1就变成了7。目标达到了。

\textbf{答案：} $1$。

所以，$\log_7 7 = 1$。

（这是一个铁律：\textbf{不管底下的数字是几，只要和旁边的数字一样，答案永远是 1。} 变成自己，只需要 1 次。）

\subsection{第四局（终极脑筋急转弯）}

\[
\log_8 1 = ?
\]

\textbf{直接翻译：}

每次变 8 倍，我想变成 1，需要按几次？

\textbf{你想一想：}

等等，箱子里\textbf{一开始不就已经是 1 个了吗？}

我不用按啊！目标已经实现了！

\textbf{答案：} $0$ 次！连按钮都不用碰。

所以，$\log_8 1 = 0$。

（又一条铁律：\textbf{不管底数是几，只要目标是 1，对数永远是 0。} 因为你站在起点，哪里都不用去。）

\section{两个带着面具的超级 VIP}

在对数的大家族里，有两个“底数”（也就是按按钮翻倍的规矩）特别特别常用。

因为太常用了，数学家觉得每次都写底数太麻烦，直接给它们戴上了面具，发了“VIP专属缩写”。

\subsection{第一个 VIP：以 10 为底的对数}

如果我们用的是“10倍克隆箱”，每次变 10 倍。

$\log_{10} 100 = 2$ （变成100要2次）

$\log_{10} 1000 = 3$ （变成1000要3次）

因为我们人类有10根手指，用的就是逢十进一，所以“每次变10倍”太普遍了。

数学家规定：\textbf{只要底数是 10，就把小小的 10 藏起来，而且把 $\log$ 写成 $\lg$。}

看到 \textbf{$\lg$}，你就要知道：这是在问“每次变10倍”。

\[
\lg 10000 = 4
\]

\textbf{最简单的翻译法：} $\lg$ 就是在问你，“这个数字后面跟了几个 0？”

10000 后面有 4 个 0，答案就是 4。

简单粗暴吧？

\subsection{第二个 VIP：以 $e$ 为底的对数}

还记得之前提到过的那个自然常数 $e \approx 2.718$ 吗？

它是大自然里事物（比如细菌繁殖、利息增长）最本能的增长速度。

如果底数是 $e$，数学家把它缩写成 \textbf{$\ln$}（读作“烙印”或者“len”）。

看到 \textbf{$\ln x$}，它其实就是 \textbf{$\log_e x$}。

\textbf{翻译：} 按照大自然最本能的速度去增长，想要长到 $x$ 这么大，需要多长时间（几次）？

以后在大学的物理、生物公式里，$\ln$ 会像空气一样到处都是。你现在只要知道它是“以 $e$ 为底的对数”就行了。

\section{对数的终极魔法：为什么要发明它？}

你可能会问一个问题：“就算我能算出来按了几次按钮，这东西到底有啥大用处啊？”

其实，在几百年前，这东西救了全世界天文学家的命。

那时候没有计算机，天文学家算星星的轨道，需要算极其恐怖的乘法。比如：

$1024 \times 32768 = ?$

手算能把人算瞎。

这时候，对数像一个超级英雄一样登场了。它带来了一个不可思议的魔法：

\textbf{对数，可以把恶心的“乘法”，变成简单的“加法”！}

怎么做到的？我们回到“2倍克隆箱”。

我想算 $8 \times 4 = ?$。

用克隆箱的视角来看：

要想得到 \textbf{8}，需要按 \textbf{3次} 按钮。（$\log_2 8 = 3$）

要想得到 \textbf{4}，需要按 \textbf{2次} 按钮。（$\log_2 4 = 2$）

如果我把这两件事“连着做”，我总共按了几次按钮？

一共按了 $3 + 2 = 5$ 次按钮！

那按 5 次按钮，最后能变出多少个？

$2^5 = 32$！

你发现了吗？

我们要算的 \textbf{$8 \times 4$（乘法）}，变成了求总次数的 \textbf{$3 + 2$（加法）}！

这就诞生了对数最著名的一条公式：

\[
\log_a(M \times N) = \log_a M + \log_a N
\]

\textbf{用大白话翻译：}

\begin{quote}
\textbf{“想要得到 $M \times N$ 这么多的东西，需要的总次数，就等于‘得到 $M$ 的次数’加上‘得到 $N$ 的次数’。”}
\end{quote}

这就是对数被发明的初衷！

以前的人，手里有一本厚厚的“对数表”（里面记录了各种数字对应的次数）。

遇到几万乘几万的数字，他们不去做乘法。

他们查表找到这两个数字各自的“次数”，然后把次数\textbf{加}起来，再反向查表，看这个总次数对应哪个数字，答案就出来了！

\textbf{乘法变成了加法！}

就这一招，让计算效率提升了无数倍。

同理，既然乘法变成了加法，那\textbf{除法就变成了减法}：

\[
\log_a(M \div N) = \log_a M - \log_a N
\]

\section{这一章你学到了什么？}

对数看起来像外星文，其实它是世界上最直白的问题。

\textbf{第一：} $\log$ 就是在问：“如果每次翻固定倍数，想变成目标数字，\textbf{要折腾几次？}”

\textbf{第二：} $\log_a 1$ 永远等于 $0$。因为目标就是起点，不用折腾。$\log_a a$ 永远等于 $1$。变成自己，只需要一次。

\textbf{第三：} $\lg$ 是每次变 10 倍（数零的个数），$\ln$ 是每次变 $e$ 倍。

\textbf{第四：} 对数的终极魔法——它记录的是次数。把次数相加，就等于把结果相乘。它能把复杂的乘除法，降级成简单的加减法。

$\log$ 这个符号，说到底就是一个兢兢业业的记账员。

它帮你把天上飞的庞大数字，拽到了地面上，变成了1、2、3、4这样简单的次数。

带着这个记账员，咱们接着往前走。

\chapter{补充章进阶：\texorpdfstring{$\ln$}{ln} 说是一样，其实一点都不一样}

上一章里，我告诉你：

\[
\ln x \text{ 其实就是 } \log_e x
\]

它只是以 $e$ 为底的对数的缩写。

表面上看，它和 $\log_2$、$\lg$ 没什么区别，都是在问：“要折腾多少次？”

但如果你真的这么想，那就太小看 $\ln$ 了。

在初中，你可能最喜欢用 $\lg$（因为数 0 很方便）。

但在高中和大学，$\lg$ 几乎被打入冷宫。无论你在看什么公式，物理的、概率的、微积分的，满屏幕全是 \textbf{$\ln$}。

为什么数学家这么偏爱 $\ln$？它到底特殊在哪里？

今天，我就告诉你 $\ln$ 藏着的三大绝招。

\section{绝招一：和微积分是“天作之合”}

我们前面讲过求导和积分。

现在我问你：如果对普通的 $\log$ 求导，会发生什么？

比如 $\dfrac{d}{dx}(\log_{10} x)$。

算出来的结果非常丑：

\[
= \frac{1}{x \cdot \ln 10}
\]

公式里居然硬生生地多出了一块带 $e$ 的尾巴（$\ln 10$）！

这就像你穿了一件本来很合身的西装，背后却莫名其妙拖着一条布条。非常碍事。

\textbf{但是，如果是对 $\ln x$ 求导呢？}

奇迹发生了：

\[
\frac{d}{dx}(\ln x) = \frac{1}{x}
\]

干干净净！清清爽爽！没有任何多余的尾巴！

反过来说，如果你要对 $\dfrac{1}{x}$ 进行积分：

\[
\int \frac{1}{x} \, dx = \ln|x| + C
\]

你发现了吗？

在微积分的世界里，\textbf{$\ln$ 和 $\dfrac{1}{x}$ 是一对最完美的搭档。}

只有用 $\ln$，微积分的公式才是最简洁的。你换成任何其他底数，公式都会变得又长又丑。

数学家都是重度“颜控”（喜欢公式简洁美观）。

就凭这一点，$\ln$ 就坐稳了正宫的位置。

\section{绝招二：把“长在天上”的未知数拽下来}

有时候，我们会遇到一种特别恶心的方程。未知数不在地上待着，非要跑到天上（指数的位置）去。

比如：

\[
e^x = 5
\]

我想求 $x$ 是多少。怎么解？

你不能除以 $e$，因为 $x$ 是次方，不是乘上去的。

这时候，$\ln$ 就像一把钩子，专门用来把天上的未知数给拽下来。

我们在方程两边同时套上 $\ln$：

\[
\ln(e^x) = \ln 5
\]

还记得对数的翻译吗？

左边 $\ln(e^x)$ 的意思是：“每次变 $e$ 倍，想变成 $e^x$，要折腾几次？”

答案不就是 \textbf{$x$ 次} 嘛！

所以左边直接变成了 $x$：

\[
x = \ln 5
\]

就这么简单！

$\ln$ 天生就是为了消灭 $e$ 的次方而存在的。

只要看到 $e$ 头顶上顶着一坨复杂的东西，别犹豫，直接用 $\ln$ 去套它。

\section{绝招三：降维打击——连乘的终结者}

这是 $\ln$ 最核心、最强大的魔法。

上一章我们说过，对数有一个神技：\textbf{把乘法变成加法。}

\[
\ln(M \times N) = \ln M + \ln N
\]

如果只是两个数相乘，这个魔法看起来还不算太震撼。

但如果是\textbf{连乘}呢？

还记得那个代表连乘的门框符号 $\prod$ 吗？

假设你要算一百个数连乘：

\[
y = x_1 \times x_2 \times x_3 \times \dots \times x_{100}
\]

也就是：

\[
y = \prod_{i=1}^{100} x_i
\]

在现实里（比如算概率、算统计学数据），这种一百个数相乘的情况非常多。

如果有几个数特别小（比如 0.001），你乘着乘着，计算器就爆了，因为数字小到超出了计算器的显示极限。

更别提你要对这个庞然大物求导了，根本无从下手。

这时候，数学家会祭出终极大招：\textbf{在等式两边，同时套上 $\ln$。}

\[
\ln y = \ln(x_1 \times x_2 \times x_3 \times \dots \times x_{100})
\]

见证奇迹的时刻到了。

根据对数的法则，“里面相乘，等于外面相加”：

\[
\ln y = \ln x_1 + \ln x_2 + \ln x_3 + \dots + \ln x_{100}
\]

用连加符号 $\sum$ 写出来就是：

\[
\ln y = \sum_{i=1}^{100} \ln x_i
\]

你看懂发生了什么吗？

\textbf{那个恶心、恐怖、随时会让计算器崩溃的连乘符号 $\prod$，}

\textbf{在套上 $\ln$ 的一瞬间，当场被“降维打击”，变成了一个温顺的连加符号 $\sum$！}

这就是 $\ln$ 真正的恐怖之处：

\begin{quote}
\textbf{$\ln$ 是一座桥梁。它强行把“连乘的世界”，搬到了“连加的世界”。}
\end{quote}

在数学里，处理加法（连加 $\sum$）比处理乘法（连乘 $\prod$）要容易一万倍。

所以，当你以后在概率论、统计学或者机器学习的公式里，看到一堆密密麻麻的 $\prod$ 时，如果下一步它突然消失了，变成了 $\sum$，不要惊讶。

那是因为数学家悄悄给它施加了 $\ln$ 魔法。

\section{这一章你学到了什么？}

很多人以为 $\ln$ 只是 $\log$ 的一种。

不，在数学的深水区，$\ln$ 是绝对的主角，其他 $\log$ 只是配角。

\textbf{第一：} $\ln$ 和微积分是天作之合。$\ln x$ 的导数是极其干净的 $\dfrac{1}{x}$。

\textbf{第二：} $\ln$ 是一把钩子。专门用来把 $e$ 头顶上的未知数拽回地面。

\textbf{第三：} $\ln$ 是降维打击的武器。它能把复杂的连乘 $\prod$，瞬间拆解成简单的连加 $\sum$。

现在，你不仅认识了对数，你还掌握了数学家手里最顺手的一把瑞士军刀。

当 $\prod$ 让你头疼时，记得呼叫 $\ln$。

它会把乘法劈成加法，把地狱变成天堂。

带上这把刀，继续走吧。

\chapter{这些更高级的组合和翻译}

前几章，你学会了很多符号：

$\sum$——连加。

$\prod$——连乘。

$n!$——阶乘。

你知道它们是什么意思了。

但\textbf{知道意思}和\textbf{会算}是两回事。

\textbf{会算}和\textbf{真正理解}又是两回事。

这一章，我要让你\textbf{真正理解}：看到这些符号，脑子里应该怎么想，手应该怎么动。

\section{先从连加开始：$\sum$ 到底在说什么？}

\[
\sum_{i=1}^{4} i
\]

你上一章学过，这个符号的意思是“从1加到4”。

但我现在要问你一个更深的问题：

\textbf{当你看到这个符号的时候，你的脑子里应该出现什么？}

\subsection{第一步：看到 $\sum$，脑子里想“我要加东西了”}

$\sum$ 这个符号一出现，你就要告诉自己：

\begin{quote}
“我接下来要做一堆加法。不是加一次，是加很多次。”
\end{quote}

就像你看到一个购物车，你就知道：接下来要往里面放很多东西。

$\sum$ 就是那个购物车。

\subsection{第二步：看下面和上面的数字，知道“加几次”}

\[
\sum_{i=1}^{4}
\]

下面写着 $i = 1$。

上面写着 $4$。

这告诉你：

\begin{quote}
“$i$ 从1开始，一直到4结束。”
\end{quote}

$i$ 会依次等于：$1, 2, 3, 4$。

一共4个数。

所以你要往购物车里\textbf{放4次东西}。

\subsection{第三步：看右边的式子，知道“每次放什么”}

\[
\sum_{i=1}^{4} i
\]

右边写着 $i$。

这告诉你：每次放进购物车的东西，就是 $i$ 本身。

第一次，$i = 1$，放进去的是 \textbf{1}。

第二次，$i = 2$，放进去的是 \textbf{2}。

第三次，$i = 3$，放进去的是 \textbf{3}。

第四次，$i = 4$，放进去的是 \textbf{4}。

\subsection{第四步：把购物车里的东西加起来}

购物车里现在有：1, 2, 3, 4。

加起来：$1 + 2 + 3 + 4 = 10$。

这就是答案。

\subsection{让我再说一遍，用最慢的方式}

当你看到 $\sum_{i=1}^{4} i$，你的脑子里应该这样想：

1. \textbf{$\sum$}——“我要加东西了。”

2. \textbf{下面 $i=1$，上面 $4$}——“$i$ 从1走到4，一共走4步。”

3. \textbf{右边 $i$}——“每一步，我就把当前的 $i$ 放进购物车。”

4. \textbf{最后}——“把购物车里的东西全部加起来。”

你可以想象自己站在一条路上。

路上有4个站，分别是站1、站2、站3、站4。

你从站1出发，走到站4。

每到一个站，你就捡起一个东西（就是那一站的编号）。

走完全程，你手里有1, 2, 3, 4。

加起来，就是答案。

\section{再来一个连加，这次稍微不一样}

\[
\sum_{i=1}^{3} (2 \times i)
\]

还是一步一步来。

\subsection{第一步：$\sum$——“我要加东西了。”}

好，我知道了，接下来是加法。

\subsection{第二步：下面 $i=1$，上面 $3$——“$i$ 从1走到3，一共走3步。”}

$i = 1, 2, 3$。

三个站。

\subsection{第三步：右边 $(2 \times i)$——“每一步，我要算 $2 \times i$，然后把结果放进购物车。”}

注意！这次右边不是单纯的 $i$ 了。

是 $2 \times i$。

所以每一站，我要先算一下 $2 \times i$ 等于多少，然后把\textbf{那个结果}放进购物车。

$i = 1$：$2 \times 1 = 2$，放进购物车的是 \textbf{2}。

$i = 2$：$2 \times 2 = 4$，放进购物车的是 \textbf{4}。

$i = 3$：$2 \times 3 = 6$，放进购物车的是 \textbf{6}。

\subsection{第四步：把购物车里的东西加起来}

购物车里有：2, 4, 6。

$2 + 4 + 6 = 12$。

答案是 \textbf{12}。

你看，方法是一样的。

只是“每一站捡起的东西”不一样了。

之前是直接捡 $i$。

现在是先算 $2 \times i$，再捡那个结果。

\section{连加的本质}

让我用最简单的话再说一遍。

\textbf{$\sum$ 就是让你走一条路，每一站捡一个东西，最后把所有东西加起来。}

下面和上面的数字告诉你：路有几站，从哪站走到哪站。

右边的式子告诉你：每一站捡的东西是什么。

就这么简单。

不管右边的式子多复杂——可能是 $i$，可能是 $2i$，可能是 $i^2$，可能是 $i!$——方法都一样：

1. 走到那一站。
2. 把 $i$ 的值代进去，算出这一站的东西是什么。
3. 捡起来，放进购物车。
4. 走到下一站，重复。
5. 走完了，把购物车里的东西全部加起来。

\section{现在来讲连乘：$\prod$ 到底在说什么？}

$$\prod_{i=1}^{4} i$$

上一章你学过，这个符号的意思是“从1乘到4”。

连乘和连加非常像，只是最后那一步不一样。

\subsection{第一步：看到 $\prod$，脑子里想“我要乘东西了”}

$\prod$ 这个符号一出现，你就要告诉自己：

\begin{quote}
“我接下来要做一堆乘法。不是乘一次，是乘很多次。”
\end{quote}

如果 $\sum$ 是购物车，$\prod$ 就是另一种购物车——\textbf{乘法购物车}。

东西放进去之后，不是加起来，是\textbf{乘起来}。

\subsection{第二步：看下面和上面的数字，知道“乘几次”}

$$\prod_{i=1}^{4}$$

下面 $i = 1$，上面 $4$。

$i$ 从1走到4。

一共4个站。

你要往乘法购物车里\textbf{放4次东西}。

\subsection{第三步：看右边的式子，知道“每次放什么”}

$$\prod_{i=1}^{4} i$$

右边是 $i$。

每一站放进去的，就是 $i$ 本身。

$i = 1$：放进去 \textbf{1}。

$i = 2$：放进去 \textbf{2}。

$i = 3$：放进去 \textbf{3}。

$i = 4$：放进去 \textbf{4}。

\subsection{第四步：把购物车里的东西乘起来}

购物车里有：1, 2, 3, 4。

但这次不是加起来，是\textbf{乘起来}。

$1 \times 2 \times 3 \times 4 = 24$。

答案是 \textbf{24}。

你看，连乘和连加的流程几乎一模一样。

唯一的区别：

\begin{quote}
\textbf{$\sum$（连加）：最后把所有东西加起来。}

\textbf{$\prod$（连乘）：最后把所有东西乘起来。}
\end{quote}

其他步骤完全一样。

\section{再来一个连乘}

$$\prod_{i=2}^{4} (i + 1)$$

\subsection{第一步：$\prod$——“我要乘东西了。”}

\subsection{第二步：下面 $i=2$，上面 $4$——“$i$ 从2走到4。”}

注意！这次不是从1开始，是从\textbf{2}开始。

$i = 2, 3, 4$。

三个站。

\subsection{第三步：右边 $(i + 1)$——“每一站，我要算 $i + 1$，然后放进购物车。”}

$i = 2$：$2 + 1 = 3$，放进购物车的是 \textbf{3}。

$i = 3$：$3 + 1 = 4$，放进购物车的是 \textbf{4}。

$i = 4$：$4 + 1 = 5$，放进购物车的是 \textbf{5}。

\subsection{第四步：把购物车里的东西乘起来}

购物车里有：3, 4, 5。

$3 \times 4 \times 5 = 60$。

答案是 \textbf{60}。

\section{连加和连乘混在一起}

现在来点刺激的。

$$\sum_{i=1}^{3} i!$$

这个式子里，有连加 $\sum$，还有阶乘 $i!$。

别慌。方法还是一样的。

\subsection{第一步：$\sum$——“我要加东西了。”}

好，我知道了，最后是把东西加起来。

\subsection{第二步：$i$ 从1走到3，三个站。}

$i = 1, 2, 3$。

\subsection{第三步：右边是 $i!$——“每一站，我要算 $i$ 的阶乘，然后放进购物车。”}

这一步稍微复杂一点，因为要算阶乘。

但别急，一个一个来。

\textbf{$i = 1$：}

$1! = 1$。（从1乘到1，就是1。）

放进购物车的是 \textbf{1}。

\textbf{$i = 2$：}

$2! = 2 \times 1 = 2$。

放进购物车的是 \textbf{2}。

\textbf{$i = 3$：}

$3! = 3 \times 2 \times 1 = 6$。

放进购物车的是 \textbf{6}。

\subsection{第四步：把购物车里的东西加起来}

购物车里有：1, 2, 6。

$1 + 2 + 6 = 9$。

答案是 \textbf{9}。

\section{最复杂的情况：连乘里面套连加}

$$\prod_{i=1}^{2} \left( \sum_{j=1}^{i} j \right)$$

两层符号套在一起。

看起来很吓人。

但方法还是一样的：\textbf{从外到内理解，从内到外计算。}

\subsection{从外到内理解}

最外层是 $\prod_{i=1}^{2}$——“我要乘东西了，$i$ 从1走到2，两个站。”

每一站放进购物车的东西是什么？是 $\sum_{j=1}^{i} j$——一个连加。

所以这道题是这样的：

1. 我走两站（$i = 1$ 和 $i = 2$）。
2. 每一站，我要先算一个连加，算出一个数，然后把这个数放进购物车。
3. 最后把购物车里的东西乘起来。

\subsection{从内到外计算}

好，现在开始走。

\textbf{第一站：$i = 1$}

这一站要算的连加是 $\sum_{j=1}^{1} j$。

$j$ 从1走到1，只有一站。

$j = 1$：放进去 \textbf{1}。

加起来：就是 \textbf{1}。

所以这一站放进乘法购物车的是 \textbf{1}。

\textbf{第二站：$i = 2$}

这一站要算的连加是 $\sum_{j=1}^{2} j$。

$j$ 从1走到2，两个站。

$j = 1$：放进去 \textbf{1}。

$j = 2$：放进去 \textbf{2}。

加起来：$1 + 2 = 3$。

所以这一站放进乘法购物车的是 \textbf{3}。

\subsection{最后：把乘法购物车里的东西乘起来}

乘法购物车里有：1, 3。

$1 \times 3 = 3$。

答案是 \textbf{3}。

你看，再复杂的式子，方法都是一样的。

**从外往内理解：先看最外层是什么操作，再看里面每一项是什么。**

**从内往外计算：先算最里面的，得到一个数，再往外走一层。**

像剥洋葱一样。一层一层来。

\section{现在到你了}

我给你几个式子，你来用“购物车”的方法，一步一步算出来。

\textbf{（1）}

$$\sum_{i=1}^{3} (i + 1)$$

$\sum$——加法购物车。

$i$ 从1走到3，三个站。

$i = 1$：$1 + 1 = 2$，放进去 \textbf{2}。

$i = 2$：$2 + 1 = 3$，放进去 \textbf{3}。

$i = 3$：$3 + 1 = 4$，放进去 \textbf{4}。

加起来：$2 + 3 + 4 = 9$。

\textbf{答案：9}

\textbf{（2）}

$$\prod_{i=1}^{3} (i + 1)$$

$\prod$——乘法购物车。

$i$ 从1走到3，三个站。

$i = 1$：$1 + 1 = 2$，放进去 \textbf{2}。

$i = 2$：$2 + 1 = 3$，放进去 \textbf{3}。

$i = 3$：$3 + 1 = 4$，放进去 \textbf{4}。

乘起来：$2 \times 3 \times 4 = 24$。

\textbf{答案：24}

\textbf{（3）}

$$\sum_{i=1}^{4} i^2$$

$\sum$——加法购物车。

$i$ 从1走到4，四个站。

$i = 1$：$1^2 = 1$，放进去 \textbf{1}。

$i = 2$：$2^2 = 4$，放进去 \textbf{4}。

$i = 3$：$3^2 = 9$，放进去 \textbf{9}。

$i = 4$：$4^2 = 16$，放进去 \textbf{16}。

加起来：$1 + 4 + 9 + 16 = 30$。

\textbf{答案：30}

\textbf{（4）}

$$\sum_{i=2}^{4} i!$$

$\sum$——加法购物车。

$i$ 从2走到4，三个站。

$i = 2$：$2! = 2 \times 1 = 2$，放进去 \textbf{2}。

$i = 3$：$3! = 3 \times 2 \times 1 = 6$，放进去 \textbf{6}。

$i = 4$：$4! = 4 \times 3 \times 2 \times 1 = 24$，放进去 \textbf{24}。

加起来：$2 + 6 + 24 = 32$。

\textbf{答案：32}

\textbf{（5）}

$$\prod_{i=1}^{3} i!$$

$\prod$——乘法购物车。

$i$ 从1走到3，三个站。

$i = 1$：$1! = 1$，放进去 \textbf{1}。

$i = 2$：$2! = 2$，放进去 \textbf{2}。

$i = 3$：$3! = 6$，放进去 \textbf{6}。

乘起来：$1 \times 2 \times 6 = 12$。

\textbf{答案：12}

\section{一个你可能已经想到的问题}

好了，连加和连乘你都会用了。

但你有没有想过一件事：

\begin{quote}
\textbf{有连加，为什么没有连减？}

\textbf{有连乘，为什么没有连除？}

\textbf{有阶乘，为什么没有阶除？}
\end{quote}

\section{为什么没有“连减”？}

想象一下“连减”会长什么样。

假设我想算：

$$10 - 1 - 2 - 3$$

从10开始，减去1，再减去2，再减去3。

如果有一个“连减符号”，也许会写成 $\text{什么}_{i=1}^{3} i$？

但数学家说：不需要这个符号。

为什么？

因为\textbf{减法其实就是加法}。

减去一个数，就是加上它的\textbf{相反数}。

$10 - 3$ 就是 $10 + (-3)$。

减去3，就是加上负3。

同样的道理：

$$10 - 1 - 2 - 3$$

就是：

$$10 + (-1) + (-2) + (-3)$$

这不就是\textbf{连加}吗？只不过加的是负数。

用 $\sum$ 来写：

$$10 + \sum_{i=1}^{3} (-i)$$

$(-i)$ 是什么？$i$ 的相反数。

$i = 1$ 时，$(-i) = -1$。

$i = 2$ 时，$(-i) = -2$。

$i = 3$ 时，$(-i) = -3$。

所以 $\sum_{i=1}^{3} (-i) = (-1) + (-2) + (-3) = -6$。

$10 + (-6) = 10 - 6 = 4$。

你看，用连加就能算“连减”。

\textbf{减法是加法的小弟。大哥的符号够用了，小弟不需要单独出头。}

\section{为什么没有“连除”？}

同样的道理。

假设我想算：

$$24 \div 2 \div 3 \div 4$$

从24开始，除以2，再除以3，再除以4。

数学家说：也不需要单独的符号。

为什么？

因为\textbf{除法其实就是乘法}。

除以一个数，就是乘以它的\textbf{倒数}。

$6 \div 2$ 就是 $6 \times \dfrac{1}{2}$。

除以2，就是乘以二分之一。

同样的道理：

$$24 \div 2 \div 3 \div 4$$

就是：

$$24 \times \frac{1}{2} \times \frac{1}{3} \times \frac{1}{4}$$

这不就是\textbf{连乘}吗？只不过乘的是分数（倒数）。

用 $\prod$ 来写：

$$24 \times \prod_{i=2}^{4} \frac{1}{i}$$

$\dfrac{1}{i}$ 是什么？$i$ 的倒数。

$i = 2$ 时，$\dfrac{1}{i} = \dfrac{1}{2}$。

$i = 3$ 时，$\dfrac{1}{i} = \dfrac{1}{3}$。

$i = 4$ 时，$\dfrac{1}{i} = \dfrac{1}{4}$。

$\prod_{i=2}^{4} \frac{1}{i} = \frac{1}{2} \times \frac{1}{3} \times \frac{1}{4} = \frac{1}{24}$。

$24 \times \frac{1}{24} = 1$。

用连乘就能算“连除”。

\textbf{除法是乘法的小弟。大哥的符号够用了，小弟不需要单独出头。}

\section{为什么没有“阶除”？}

你可能想象的“阶除”长这样：

$$5 \div 4 \div 3 \div 2 \div 1$$

从5开始，一层一层往下除。

但用刚才的方法，这可以写成：

$$5 \times \frac{1}{4} \times \frac{1}{3} \times \frac{1}{2} \times \frac{1}{1}$$

也就是：

$$\frac{5}{4 \times 3 \times 2 \times 1} = \frac{5}{4!} = \frac{5}{24}$$

用阶乘就能表示了。

而且，这个“阶除”算出来的结果，没什么特别的用处。

阶乘之所以出名，是因为它在排列组合里\textbf{到处出现}。

“5个人排成一排有多少种排法？$5!$ 种。”

这个问题很有用。

但“5个东西往下除”——这在现实中几乎不会碰到。

\textbf{数学家只给有用的东西发明符号。没用的，就不发明了。}

（数学家也懒。能少记一个符号就少记一个。）

\section{总结一下}

\begin{quote}
\textbf{为什么没有连减？} 因为减法就是“加上相反数”，连加已经能做连减了。

\textbf{为什么没有连除？} 因为除法就是“乘以倒数”，连乘已经能做连除了。

\textbf{为什么没有阶除？} 因为用阶乘就能表示，而且没什么用。
\end{quote}

\textbf{加法和乘法是老大。减法和除法是小弟。老大的符号够用了，小弟不需要单独出头。}

\section{这一章你学到了什么}

\textbf{第一：} 看到 $\sum$，想“加法购物车”。看到 $\prod$，想“乘法购物车”。

\textbf{第二：} 从下面和上面的数字，知道要走几站。从右边的式子，知道每一站捡什么东西。走完了，加起来或乘起来。

\textbf{第三：} 连加和连乘可以套在一起。方法是从外往内理解，从内往外计算。一层一层来。

\textbf{第四：} 减法是加法的小弟，除法是乘法的小弟。所以没有连减、连除、阶除——老大的符号够用了。

这些符号，你现在不光认识了，还知道看到它们的时候脑子里该出现什么画面，手该怎么一步一步动起来。

好了，歇口气，咱们接着走。

\chapter{这些更高级的组合和翻译}

前面几章，你学会了很多符号：

$\sum$——连加。

$\prod$——连乘。

$n!$——阶乘。

$\int$——积分。

你知道它们各自是什么意思了。

你也知道怎么把 $\sum$ 和 $\prod$ 混在一起用。

你还知道 $\sum$ 和 $\int$ 是亲兄弟——一个是离散版，一个是连续版。

现在，我们来玩点更刺激的。

\textbf{把它们全部混在一起。}

\section{第一种组合：连加里面套积分}

\[
\sum_{i=1}^{3} \int_0^i x \, dx
\]

看起来很吓人？

别慌。方法还是一样的：\textbf{从外往内理解，从内往外计算。}

\subsection*{从外往内理解}

最外层是什么？

$\sum_{i=1}^{3}$——连加。

$i$ 从1走到3，三个站。

每一站放进购物车的东西是什么？

是 $\int_0^i x \, dx$——一个积分。

所以这道题是这样的：

1. 我走三站（$i = 1, 2, 3$）。
2. 每一站，我要先算一个积分，算出一个数，然后把这个数放进购物车。
3. 最后把购物车里的东西加起来。

\subsection*{从内往外计算}

好，现在开始走。

\textbf{第一站：$i = 1$}

这一站要算的积分是 $\int_0^1 x \, dx$。

这是什么意思？

从 $x = 0$ 到 $x = 1$，曲线 $y = x$ 下面的面积。

$y = x$ 是一条直线，从原点出发，斜着往上走。

从 $x = 0$ 到 $x = 1$，这条直线下面围成的区域是一个三角形。

底是1，高是1。

三角形面积 = $\dfrac{1}{2} \times$ 底 $\times$ 高 $= \dfrac{1}{2} \times 1 \times 1 = \dfrac{1}{2}$。

所以 $\int_0^1 x \, dx = \dfrac{1}{2}$。

这一站放进购物车的是 $\dfrac{1}{2}$。

\textbf{第二站：$i = 2$}

这一站要算的积分是 $\int_0^2 x \, dx$。

从 $x = 0$ 到 $x = 2$，曲线 $y = x$ 下面的面积。

还是一个三角形。底是2，高是2。

面积 $= \dfrac{1}{2} \times 2 \times 2 = 2$。

所以 $\int_0^2 x \, dx = 2$。

这一站放进购物车的是 $2$。

\textbf{第三站：$i = 3$}

这一站要算的积分是 $\int_0^3 x \, dx$。

三角形，底是3，高是3。

面积 $= \dfrac{1}{2} \times 3 \times 3 = \dfrac{9}{2}$。

所以 $\int_0^3 x \, dx = \dfrac{9}{2}$。

这一站放进购物车的是 $\dfrac{9}{2}$。

\subsection*{最后：把购物车里的东西加起来}

购物车里有：$\dfrac{1}{2}$，$2$，$\dfrac{9}{2}$。

先把2变成分数：$2 = \dfrac{4}{2}$。

$\dfrac{1}{2} + \dfrac{4}{2} + \dfrac{9}{2} = \dfrac{1 + 4 + 9}{2} = \dfrac{14}{2} = 7$。

答案是 \textbf{7}。

你看，虽然里面有积分，但方法还是一样的：

\textbf{一站一站走，每站先算出里面那个东西（积分），得到一个数，放进购物车。最后加起来。}

\section{第二种组合：积分里面套连加}

\[
\int_0^1 \left( \sum_{i=1}^{3} x^i \right) dx
\]

这次反过来了：外面是积分，里面是连加。

\subsection*{从外往内理解}

最外层是什么？

$\int_0^1 \cdots \, dx$——积分。

从 $x = 0$ 到 $x = 1$，算曲线下面的面积。

曲线是什么？

是 $\sum_{i=1}^{3} x^i$——一个连加。

所以这道题是这样的：

1. 我先要弄清楚里面那个连加是什么。
2. 然后再对它做积分。

\subsection*{第一步：搞清楚里面的连加}

$\sum_{i=1}^{3} x^i$ 是什么？

$i$ 从1走到3，每一项是 $x^i$。

$i = 1$：$x^1 = x$

$i = 2$：$x^2$

$i = 3$：$x^3$

加起来：$x + x^2 + x^3$。

所以 $\sum_{i=1}^{3} x^i = x + x^2 + x^3$。

\subsection*{第二步：做积分}

原题变成了：

\[
\int_0^1 (x + x^2 + x^3) \, dx
\]

积分可以拆开算（加法可以分开积）：

\[
\int_0^1 x \, dx + \int_0^1 x^2 \, dx + \int_0^1 x^3 \, dx
\]

现在一个一个算。

\textbf{第一个：$\int_0^1 x \, dx$}

$y = x$ 从0到1的面积。三角形，底1高1。

$= \dfrac{1}{2}$

\textbf{第二个：$\int_0^1 x^2 \, dx$}

这个需要用积分公式。

$x^2$ 的积分是 $\dfrac{x^3}{3}$。

从0到1：$\dfrac{1^3}{3} - \dfrac{0^3}{3} = \dfrac{1}{3} - 0 = \dfrac{1}{3}$

\textbf{第三个：$\int_0^1 x^3 \, dx$}

$x^3$ 的积分是 $\dfrac{x^4}{4}$。

从0到1：$\dfrac{1^4}{4} - \dfrac{0^4}{4} = \dfrac{1}{4} - 0 = \dfrac{1}{4}$

\subsection*{加起来}

$\dfrac{1}{2} + \dfrac{1}{3} + \dfrac{1}{4}$

找公分母。$2, 3, 4$ 的最小公倍数是 $12$。

$\dfrac{1}{2} = \dfrac{6}{12}$

$\dfrac{1}{3} = \dfrac{4}{12}$

$\dfrac{1}{4} = \dfrac{3}{12}$

$\dfrac{6}{12} + \dfrac{4}{12} + \dfrac{3}{12} = \dfrac{13}{12}$

答案是 $\dfrac{13}{12}$。

\section{第三种组合：求导和连加}

\[
\frac{d}{dx} \left( \sum_{i=1}^{3} x^i \right)
\]

外面是求导，里面是连加。

\subsection*{第一步：搞清楚里面的连加}

刚才算过了：

$\sum_{i=1}^{3} x^i = x + x^2 + x^3$

\subsection*{第二步：对它求导}

求导也可以拆开算（加法可以分开求导）：

$\dfrac{d}{dx}(x + x^2 + x^3) = \dfrac{d}{dx}(x) + \dfrac{d}{dx}(x^2) + \dfrac{d}{dx}(x^3)$

$x$ 的导数是 $1$。

$x^2$ 的导数是 $2x$。

$x^3$ 的导数是 $3x^2$。

加起来：

$1 + 2x + 3x^2$

答案是 $1 + 2x + 3x^2$。

\section{第四种组合：阶乘和积分}

\[
\int_0^1 x^{n!} \, dx
\]

假设 $n = 3$。

那 $n! = 3! = 6$。

所以要算的是：

\[
\int_0^1 x^6 \, dx
\]

$x^6$ 的积分是 $\dfrac{x^7}{7}$。

从0到1：$\dfrac{1^7}{7} - \dfrac{0^7}{7} = \dfrac{1}{7}$

答案是 $\dfrac{1}{7}$。

\section{第五种组合：连乘和求导}

\[
\frac{d}{dx} \left( \prod_{i=1}^{3} x^i \right)
\]

\subsection*{第一步：搞清楚里面的连乘}

$\prod_{i=1}^{3} x^i$ 是什么？

$i$ 从1到3，每一项是 $x^i$，全部乘起来。

$x^1 \times x^2 \times x^3$

指数相同的底数相乘，指数相加：

$x^{1+2+3} = x^6$

所以 $\prod_{i=1}^{3} x^i = x^6$。

\subsection*{第二步：对它求导}

$x^6$ 的导数是 $6x^5$。

答案是 $6x^5$。

\section{最复杂的组合：全部混在一起}

\[
\sum_{n=1}^{2} \int_0^1 x^{n!} \, dx
\]

外面是连加，里面是积分，积分里面还有阶乘。

\subsection*{从外往内理解}

$\sum_{n=1}^{2}$——连加，$n$ 从1到2，两站。

每一站要算的是 $\int_0^1 x^{n!} \, dx$——一个积分，而且指数是 $n!$。

\subsection*{从内往外计算}

\textbf{第一站：$n = 1$}

$n! = 1! = 1$。

要算的积分是 $\int_0^1 x^1 \, dx = \int_0^1 x \, dx$。

三角形，底1高1，面积 $= \dfrac{1}{2}$。

这一站放进购物车的是 $\dfrac{1}{2}$。

\textbf{第二站：$n = 2$}

$n! = 2! = 2$。

要算的积分是 $\int_0^1 x^2 \, dx$。

$x^2$ 的积分是 $\dfrac{x^3}{3}$。

从0到1：$\dfrac{1}{3}$。

这一站放进购物车的是 $\dfrac{1}{3}$。

\subsection*{加起来}

$\dfrac{1}{2} + \dfrac{1}{3}$

公分母是6。

$\dfrac{3}{6} + \dfrac{2}{6} = \dfrac{5}{6}$

答案是 $\dfrac{5}{6}$。

\section{翻译的秘诀}

不管式子多复杂，翻译的方法永远是一样的：

\textbf{第一步：找到最外层的符号，知道"最后要做什么操作"。}

是 $\sum$？那最后要加起来。

是 $\prod$？那最后要乘起来。

是 $\int$？那最后要算面积。

是 $\dfrac{d}{dx}$？那最后要求导。

\textbf{第二步：看里面套了什么，一层一层往里看。}

\textbf{第三步：从最里面开始算，算完一层往外走一层。}

就像剥洋葱，从外面看清楚结构，从里面开始下手。

\section{你来试几个}

\textbf{（1）}

\[
\sum_{i=1}^{2} i!
\]

$i = 1$：$1! = 1$

$i = 2$：$2! = 2$

加起来：$1 + 2 = 3$

\textbf{答案：3}

\textbf{（2）}

\[
\prod_{i=1}^{3} i
\]

$= 1 \times 2 \times 3 = 6$

\textbf{答案：6}（这就是 $3!$）

\textbf{（3）}

\[
\sum_{i=1}^{2} \int_0^1 x^i \, dx
\]

$i = 1$：$\int_0^1 x^1 \, dx = \dfrac{1}{2}$

$i = 2$：$\int_0^1 x^2 \, dx = \dfrac{1}{3}$

加起来：$\dfrac{1}{2} + \dfrac{1}{3} = \dfrac{5}{6}$

\textbf{答案：$\dfrac{5}{6}$}

\textbf{（4）}

\[
\frac{d}{dx} \left( \sum_{i=0}^{2} x^i \right)
\]

先算连加：$x^0 + x^1 + x^2 = 1 + x + x^2$

再求导：$0 + 1 + 2x = 1 + 2x$

\textbf{答案：$1 + 2x$}

\textbf{（5）}

\[
\int_0^2 \left( \prod_{i=1}^{2} x \right) dx
\]

先算连乘：$x \times x = x^2$

再积分：$\int_0^2 x^2 \, dx$

$x^2$ 的积分是 $\dfrac{x^3}{3}$。

从0到2：$\dfrac{8}{3} - 0 = \dfrac{8}{3}$

\textbf{答案：$\dfrac{8}{3}$}

\section{一个你可能已经想到的问题}

好了，连加和连乘你都会用了。

但你有没有想过一件事：

\begin{quote}
\textbf{有连加，为什么没有连减？}

\textbf{有连乘，为什么没有连除？}

\textbf{有阶乘，为什么没有阶除？}
\end{quote}

\section{为什么没有"连减"？}

想象一下"连减"会长什么样。

假设我想算：

\[
10 - 1 - 2 - 3
\]

从10开始，减去1，再减去2，再减去3。

如果有一个"连减符号"，也许会写成 $\text{什么}_{i=1}^{3} i$？

但数学家说：不需要这个符号。

为什么？

因为\textbf{减法其实就是加法}。

减去一个数，就是加上它的\textbf{相反数}。

$10 - 3$ 就是 $10 + (-3)$。

减去3，就是加上负3。

同样的道理：

\[
10 - 1 - 2 - 3
\]

就是：

\[
10 + (-1) + (-2) + (-3)
\]

这不就是\textbf{连加}吗？只不过加的是负数。

用 $\sum$ 来写：

\[
10 + \sum_{i=1}^{3} (-i)
\]

$(-i)$ 是什么？$i$ 的相反数。

$i = 1$ 时，$(-i) = -1$。

$i = 2$ 时，$(-i) = -2$。

$i = 3$ 时，$(-i) = -3$。

所以 $\sum_{i=1}^{3} (-i) = (-1) + (-2) + (-3) = -6$。

$10 + (-6) = 10 - 6 = 4$。

你看，用连加就能算"连减"。

\textbf{减法是加法的小弟。大哥的符号够用了，小弟不需要单独出头。}

\section{为什么没有"连除"？}

同样的道理。

假设我想算：

\[
24 \div 2 \div 3 \div 4
\]

从24开始，除以2，再除以3，再除以4。

数学家说：也不需要单独的符号。

为什么？

因为\textbf{除法其实就是乘法}。

除以一个数，就是乘以它的\textbf{倒数}。

$6 \div 2$ 就是 $6 \times \dfrac{1}{2}$。

除以2，就是乘以��分之一。

同样的道理：

\[
24 \div 2 \div 3 \div 4
\]

就是：

\[
24 \times \frac{1}{2} \times \frac{1}{3} \times \frac{1}{4}
\]

这不就是\textbf{连乘}吗？只不过乘的是分数（倒数）。

用 $\prod$ 来写：

\[
24 \times \prod_{i=2}^{4} \frac{1}{i}
\]

$\dfrac{1}{i}$ 是什么？$i$ 的倒数。

$i = 2$ 时，$\dfrac{1}{i} = \dfrac{1}{2}$。

$i = 3$ 时，$\dfrac{1}{i} = \dfrac{1}{3}$。

$i = 4$ 时，$\dfrac{1}{i} = \dfrac{1}{4}$。

$\prod_{i=2}^{4} \frac{1}{i} = \frac{1}{2} \times \frac{1}{3} \times \frac{1}{4} = \frac{1}{24}$。

$24 \times \frac{1}{24} = 1$。

用连乘就能算"连除"。

\textbf{除法是乘法的小弟。大哥的符号够用了，小弟不需要单独出头。}

\section{为什么没有"阶除"？}

你可能想象的"阶除"长这样：

\[
5 \div 4 \div 3 \div 2 \div 1
\]

从5开始，一层一层往下除。

但用刚才的方法，这可以写成：

\[
5 \times \frac{1}{4} \times \frac{1}{3} \times \frac{1}{2} \times \frac{1}{1}
\]

也就是：

\[
\frac{5}{4 \times 3 \times 2 \times 1} = \frac{5}{4!} = \frac{5}{24}
\]

用阶乘就能表示了。

而且，这个"阶除"算出来的结果，没什么特别的用处。

阶乘之所以出名，是因为它在排列组合里\textbf{到处出现}。

"5个人排成一排有多少种排法？$5!$ 种。"

这个问题很有用。

但"5个东西往下除"——这在现实中几乎不会碰到。

\textbf{数学家只给有用的东西发明符号。没用的，就不发明了。}

（数学家也懒。能少记一个符号就少记一个。）

\section{总结一下}

\begin{quote}
\textbf{为什么没有连减？} 因为减法就是"加上相反数"，连加已经能做连减了。

\textbf{为什么没有连除？} 因为除法就是"乘以倒数"，连乘已经能做连除了。

\textbf{为什么没有阶除？} 因为用阶乘就能表示，而且没什么用。
\end{quote}

\textbf{加法和乘法是老大。减法和除法是小弟。老大的符号够用了，小弟不需要单独出头。}

\section{这一章你学到了什么}

\textbf{第一：} 看到 $\sum$，想"加法购物车"。看到 $\prod$，想"乘法购物车"。

\textbf{第二：} 从下面和上面的数字，知道要走几站。从右边的式子，知道每一站捡什么东西。走完了，加起来或乘起来。

\textbf{第三：} 连加和连乘可以套在一起。方法是从外往内理解，从内往外计算。一层一层来。

\textbf{第四：} 减法是加法的小弟，除法是乘法的小弟。所以没有连减、连除、阶除——老大的符号够用了。

这些符号，你现在不光认识了，还知道看到它们的时候脑子里该出现什么画面，手该怎么一步一步动起来。

好了，歇口气，咱们接着走。

\chapter{如果微积分也融合进去了呢？}

你已经认识了三个符号：

$\sum$——连加。

$\prod$——连乘。

$\int$——积分。

前两个是"离散"的——一站一站走，每站捡一个东西。

（对了，补一句：每一站，只能放\textbf{一个}东西进购物车。不是到了一站放三个、到了另一站放五个。每站算出一个数，放进去一个数。整整齐齐。）

第三个是"连续"的——不是一站一站，是一整条线，无数个点。

但你有没有发现，它们\textbf{长得很像}？

\[
\sum_{i=1}^{100} i
\]

\[
\int_0^{100} x \, dx
\]

一个是从1加到100。

一个是从0积到100。

它们到底有什么关系？

\section{先回忆一下：$\sum$ 是怎么工作的}

\[
\sum_{i=1}^{5} i
\]

你脑子里应该出现这样的画面：

有5个站，分别是站1、站2、站3、站4、站5。

你从站1走到站5。

每到一个站，你捡起那个站的编号，放进购物车。

走完了，购物车里有：1, 2, 3, 4, 5。

加起来：$1 + 2 + 3 + 4 + 5 = 15$。

关键是什么？

\textbf{站是分开的。}

站1和站2之间，有空隙。

你只在整数的位置停下来。

1.5 那个位置？你不停。

2.7 那个位置？你也不停。

你只停在 1, 2, 3, 4, 5 这些\textbf{整数站}。

这叫\textbf{离散}——一个一个分开的点。

\section{再回忆一下：$\int$ 是怎么工作的}

\[
\int_0^5 x \, dx
\]

你脑子里应该出现这样的画面：

有一条线，从 $x = 0$ 一直延伸到 $x = 5$。

这条线上，每一个点都有一个高度。

$x = 0$ 的地方，高度是 0。

$x = 1$ 的地方，高度是 1。

$x = 2.5$ 的地方，高度是 2.5。

$x = 4.9$ 的地方，高度是 4.9。

每一个位置，高度就等于那个位置的 $x$ 值。

你要算的是：这条线下面的面积。

怎么算？

把这条线下面的区域，切成无数个无限薄的小长条。

每个小长条的宽度是 $dx$（无限小）。

每个小长条的高度是 $x$（那个位置的值）。

每个小长条的面积 = $x \times dx$。

把所有小长条的面积加起来。

这就是积分。

关键是什么？

\textbf{点是连续的。}

你不只是在 0, 1, 2, 3, 4, 5 这些整数位置停。

你在\textbf{每一个位置}都停。

0.001 停，0.002 停，0.003 停……

无限多个位置，无限多个小长条，无限多次"加"。

这叫\textbf{连续}——没有空隙，一整条线。

\section{$\sum$ 和 $\int$ 的关系}

现在你应该能感觉到它们的关系了：

\textbf{$\sum$ 是"一站一站加"。}

\textbf{$\int$ 是"一点一点加"，而且点无限多、无限密。}

$\sum$ 是离散版本的加法。

$\int$ 是连续版本的加法。

还有一种方式来理解它们的区别。

\textbf{$\int$ 是一条线。}

你拿一支笔，从纸的左边画到右边，一笔不断。

笔尖走过的轨迹，是一条连续的线。

这条线由无数个点组成，点和点之间没有空隙，紧紧挨在一起。

你找不到"相邻两点之间的缝隙"——因为没有缝隙。

\textbf{$\sum$ 是一排点。}

你拿一支笔，在纸上点一下，抬起来。

再点一下，抬起来。

再点一下，抬起来。

点、点、点、点、点……

每个点之间，有空隙。它们不连在一起。

你能清楚地数出来：第1个点、第2个点、第3个点……

$\int$：连续的点，组成一条线。

$\sum$：分开的点，排成一排。

一条线，和一排点。

这就是它们最本质的区别。

也可以用走路来想：

$\sum$ 是走楼梯。一步一步，踩在台阶上。台阶是分开的。

$\int$ 是走斜坡。一点一点，连续地走。没有台阶，是一整条平滑的坡。

楼梯走 5 步，你踩了 5 个点。

斜坡走 5 米，你踩了无限多个点。

\section{用一个例子来感受}

假设我想算"从1到100所有整数的和"。

用 $\sum$：

\[
\sum_{i=1}^{100} i = 1 + 2 + 3 + \cdots + 100
\]

这是100个整数加起来。

答案是 5050。

（有一个公式可以快速算：$\dfrac{100 \times 101}{2} = 5050$。但那是另一个故事了。）

现在假设我想算"从0到100，把所有的数加起来"。

等一下——"所有的数"？

0.1 算吗？算。

0.01 算吗？算。

0.001 算吗？算。

那有无限多个数啊！

怎么加？

用 $\int$：

\[
\int_0^{100} x \, dx
\]

这就是"从0到100，把每一个位置的 $x$ 值都加起来"——当然，是用"无限小的宽度 $dx$"来加。

算出来是多少？

\[
\int_0^{100} x \, dx = \frac{100^2}{2} = 5000
\]

（这个怎么算的？有一套积分的方法。你现在只需要感受"它在做什么"就好。）

你看：

$\sum_{i=1}^{100} i = 5050$（100个整数加起来）

$\int_0^{100} x \, dx = 5000$（从0到100这条线下面的面积）

一个是"离散的和"，一个是"连续的和"。

数字很接近，但意思不一样。

\section{为什么积分符号长得像 $\sum$？}

你有没有注意到：

$\sum$ 长得像一个歪歪的 S。

$\int$ 也长得像一个拉长的 S。

这不是巧合。

\textbf{$\int$ 就是 $\sum$ 的变体。}

$\sum$ 是 Sum（求和）的缩写，用了希腊字母 Sigma。

$\int$ 也是 Sum 的缩写，只不过是\textbf{连续版本}的 Sum。

数学家用一个拉长的 S 来表示"无限多次的加法"。

它们本质上是同一件事——\textbf{把一堆东西加起来}。

只是 $\sum$ 是一个一个加（有限次）。

$\int$ 是一点一点加（无限次）。

\section{再看一眼它们的结构}

\[
\sum_{i=1}^{n} f(i)
\]

\[
\int_a^b f(x) \, dx
\]

它们的结构几乎一模一样：

\begin{center}
\begin{tabular}{|c|c|c|}
\hline
 & $\sum$ & $\int$ \\
\hline
符号 & $\sum$（离散求和） & $\int$（连续求和） \\
\hline
起点 & 下面的 $i = 1$ & 下面的 $a$ \\
\hline
终点 & 上面的 $n$ & 上面的 $b$ \\
\hline
每一项 & $f(i)$（代入整数 $i$） & $f(x)$（代入任意 $x$） \\
\hline
每一项的"厚度" & 1（每站之间隔1个单位） & $dx$（无限小的厚度） \\
\hline
操作 & 把所有 $f(i)$ 加起来 & 把所有 $f(x) \cdot dx$ 加起来 \\
\hline
\end{tabular}
\end{center}

你看，它们就是"离散"和"连续"的区别。

\section{连乘呢？有没有"连续版本"？}

好问题。

$\sum$ 的连续版本是 $\int$。

那 $\prod$ 的连续版本是什么？

答案是：有，但你暂时不需要知道它叫什么。

（如果你好奇：它叫"乘积积分"或者"product integral"，但这个很少用，一般大学也不教。）

你现在只需要知道：

\textbf{$\sum$（离散加法）→ $\int$（连续加法）}

这一对是数学里最常见的。

\section{求导呢？它和 $\sum$、$\int$ 有什么关系？}

前面我们讲了积分 $\int$。

我们还讲了求导。

求导和积分是\textbf{一对}——它们是反操作。

那求导和 $\sum$ 有什么关系？

\textbf{求导的"离散版本"，叫做差分。}

$\sum$（连加）的反操作是\textbf{差分}——就是看相邻两项差了多少。

$\int$（积分）的反操作是\textbf{求导}——就是看相邻两点差了多少（当间隔无限小时）。

它们的关系是这样的：

\begin{center}
\begin{tabular}{|c|c|}
\hline
离散世界 & 连续世界 \\
\hline
$\sum$（连加） & $\int$（积分） \\
\hline
差分（相邻两项的差） & 求导（相邻两点的变化率） \\
\hline
\end{tabular}
\end{center}

差分是求导的"楼梯版"。

求导是差分的"斜坡版"。

\section{总结一下：离散和连续}

\textbf{离散}：一站一站的，有空隙，能数清楚有几个。

\textbf{连续}：一点一点的，没有空隙，数不清（无限多）。

\begin{center}
\begin{tabular}{|c|c|c|}
\hline
 & 离散版本 & 连续版本 \\
\hline
把东西加起来 & $\sum$（连加） & $\int$（积分） \\
\hline
看相邻的变化 & 差分 & 求导 \\
\hline
\end{tabular}
\end{center}

数学里很多概念都有"离散版"和"连续版"。

你学的 $\sum$、$\prod$、阶乘，都是离散世界的东西。

积分、求导，都是连续世界的东西。

它们之间有着深刻的联系。

\section{这一章你学到了什么}

\textbf{第一：} $\sum$ 是"一站一站加"（离散），$\int$ 是"一点一点加"（连续）。

\textbf{第二：} $\sum$ 是一排点，$\int$ 是一条线。$\sum$ 是走楼梯，$\int$ 是走斜坡。本质都是"一边走一边捡东西，最后算总数"。

\textbf{第三：} $\sum$ 和 $\int$ 的结构很像——都有起点、终点、每一项是什么。区别在于离散还是连续。

\textbf{第四：} 求导是积分的反操作。差分是连加的反操作。求导是差分的"连续版"。

\begin{quote}
\textbf{一句话总结：$\sum$ 和 $\int$ 是亲兄弟。一个住在"离散世界"，一个住在"连续世界"。它们做的事情一样——把东西加起来。只是一个一步一步加，一个无限细地加。}
\end{quote}

你现在站在一个挺高的地方了。

你看到了离散和连续的关系，看到了 $\sum$ 和 $\int$ 的联系，看到了数学世界里两种不同的"加法"。

这种"看到联系"的感觉，就是数学最美的地方。

好了，深呼吸一下，咱们接着探索。

\chapter{这些极高级的组合和翻译}

前面几章，你学会了很多符号：

$\sum$——连加。

$\prod$——连乘。

$n!$——阶乘。

$\int$——积分。

$\dfrac{dy}{dx}$——求导。

你知道它们各自是什么意思了。

你也知道怎么把 $\sum$ 和 $\prod$ 混在一起用。

你还知道 $\sum$ 和 $\int$ 是亲兄弟——一个是离散版，一个是连续版。

现在，我们来玩点更刺激的。

\textbf{把它们全部混在一起。}

\section{第一种组合：连加里面套积分}

\[
\sum_{i=1}^{3} \int_0^i x \, dx
\]

看起来很吓人？

别慌。方法还是一样的：\textbf{从外往内理解，从内往外计算。}

\subsection{从外往内理解}

最外层是什么？

$\sum_{i=1}^{3}$——连加。

$i$ 从1走到3，三个站。

每一站放进购物车的东西是什么？

是 $\int_0^i x \, dx$——一个积分。

所以这道题是这样的：

1. 我走三站（$i = 1, 2, 3$）。
2. 每一站，我要先算一个积分，算出一个数，然后把这个数放进购物车。
3. 最后把购物车里的东西加起来。

\subsection{从内往外计算}

好，现在开始走。

\textbf{第一站：$i = 1$}

这一站要算的积分是 $\int_0^1 x \, dx$。

这是什么意思？

从 $x = 0$ 到 $x = 1$，曲线 $y = x$ 下面的面积。

$y = x$ 是一条直线，从原点出发，斜着往上走。

从 $x = 0$ 到 $x = 1$，这条直线下面围成的区域是一个三角形。

底是1，高是1。

三角形面积 = $\dfrac{1}{2} \times$ 底 $\times$ 高 $= \dfrac{1}{2} \times 1 \times 1 = \dfrac{1}{2}$。

所以 $\int_0^1 x \, dx = \dfrac{1}{2}$。

这一站放进购物车的是 $\dfrac{1}{2}$。

\textbf{第二站：$i = 2$}

这一站要算的积分是 $\int_0^2 x \, dx$。

从 $x = 0$ 到 $x = 2$，曲线 $y = x$ 下面的面积。

还是一个三角形。底是2，高是2。

面积 $= \dfrac{1}{2} \times 2 \times 2 = 2$。

所以 $\int_0^2 x \, dx = 2$。

这一站放进购物车的是 $2$。

\textbf{第三站：$i = 3$}

这一站要算的积分是 $\int_0^3 x \, dx$。

三角形，底是3，高是3。

面积 $= \dfrac{1}{2} \times 3 \times 3 = \dfrac{9}{2}$。

所以 $\int_0^3 x \, dx = \dfrac{9}{2}$。

这一站放进购物车的是 $\dfrac{9}{2}$。

\subsection{最后：把购物车里的东西加起来}

购物车里有：$\dfrac{1}{2}$，$2$，$\dfrac{9}{2}$。

先把2变成分数：$2 = \dfrac{4}{2}$。

$\dfrac{1}{2} + \dfrac{4}{2} + \dfrac{9}{2} = \dfrac{1 + 4 + 9}{2} = \dfrac{14}{2} = 7$。

答案是 \textbf{7}。

你看，虽然里面有积分，但方法还是一样的：

\textbf{一站一站走，每站先算出里面那个东西（积分），得到一个数，放进购物车。最后加起来。}

\section{第二种组合：积分里面套连加}

\[
\int_0^1 \left( \sum_{i=1}^{3} x^i \right) dx
\]

这次反过来了：外面是积分，里面是连加。

\subsection{从外往内理解}

最外层是什么？

$\int_0^1 \cdots \, dx$——积分。

从 $x = 0$ 到 $x = 1$，算曲线下面的面积。

曲线是什么？

是 $\sum_{i=1}^{3} x^i$——一个连加。

所以这道题是这样的：

1. 我先要弄清楚里面那个连加是什么。
2. 然后再对它做积分。

\subsection{第一步：搞清楚里面的连加}

$\sum_{i=1}^{3} x^i$ 是什么？

$i$ 从1走到3，每一项是 $x^i$。

$i = 1$：$x^1 = x$

$i = 2$：$x^2$

$i = 3$：$x^3$

加起来：$x + x^2 + x^3$。

所以 $\sum_{i=1}^{3} x^i = x + x^2 + x^3$。

\subsection{第二步：做积分}

原题变成了：

\[
\int_0^1 (x + x^2 + x^3) \, dx
\]

积分可以拆开算（加法可以分开积）：

\[
\int_0^1 x \, dx + \int_0^1 x^2 \, dx + \int_0^1 x^3 \, dx
\]

现在一个一个算。

\textbf{第一个：$\int_0^1 x \, dx$}

$y = x$ 从0到1的面积。三角形，底1高1。

$= \dfrac{1}{2}$

\textbf{第二个：$\int_0^1 x^2 \, dx$}

这个需要用积分公式。

$x^2$ 的积分是 $\dfrac{x^3}{3}$。

从0到1：$\dfrac{1^3}{3} - \dfrac{0^3}{3} = \dfrac{1}{3} - 0 = \dfrac{1}{3}$

\textbf{第三个：$\int_0^1 x^3 \, dx$}

$x^3$ 的积分是 $\dfrac{x^4}{4}$。

从0到1：$\dfrac{1^4}{4} - \dfrac{0^4}{4} = \dfrac{1}{4} - 0 = \dfrac{1}{4}$

\subsection{加起来}

$\dfrac{1}{2} + \dfrac{1}{3} + \dfrac{1}{4}$

找公分母。$2, 3, 4$ 的最小公倍数是 $12$。

$\dfrac{1}{2} = \dfrac{6}{12}$

$\dfrac{1}{3} = \dfrac{4}{12}$

$\dfrac{1}{4} = \dfrac{3}{12}$

$\dfrac{6}{12} + \dfrac{4}{12} + \dfrac{3}{12} = \dfrac{13}{12}$

答案是 $\dfrac{13}{12}$。

\section{第三种组合：求导和连加}

\[
\frac{d}{dx} \left( \sum_{i=1}^{3} x^i \right)
\]

外面是求导，里面是连加。

\subsection{第一步：搞清楚里面的连加}

刚才算过了：

$\sum_{i=1}^{3} x^i = x + x^2 + x^3$

\subsection{第二步：对它求导}

求导也可以拆开算（加法可以分开求导）：

$\dfrac{d}{dx}(x + x^2 + x^3) = \dfrac{d}{dx}(x) + \dfrac{d}{dx}(x^2) + \dfrac{d}{dx}(x^3)$

$x$ 的导数是 $1$。

$x^2$ 的导数是 $2x$。

$x^3$ 的导数是 $3x^2$。

加起来：

$1 + 2x + 3x^2$

答案是 $1 + 2x + 3x^2$。

\section{第四种组合：阶乘和积分}

\[
\int_0^1 x^{n!} \, dx
\]

假设 $n = 3$。

那 $n! = 3! = 6$。

所以要算的是：

\[
\int_0^1 x^6 \, dx
\]

$x^6$ 的积分是 $\dfrac{x^7}{7}$。

从0到1：$\dfrac{1^7}{7} - \dfrac{0^7}{7} = \dfrac{1}{7}$

答案是 $\dfrac{1}{7}$。

\section{第五种组合：连乘和求导}

\[
\frac{d}{dx} \left( \prod_{i=1}^{3} x^i \right)
\]

\subsection{第一步：搞清楚里面的连乘}

$\prod_{i=1}^{3} x^i$ 是什么？

$i$ 从1到3，每一项是 $x^i$，全部乘起来。

$x^1 \times x^2 \times x^3$

指数相同的底数相乘，指数相加：

$x^{1+2+3} = x^6$

所以 $\prod_{i=1}^{3} x^i = x^6$。

\subsection{第二步：对它求导}

$x^6$ 的导数是 $6x^5$。

答案是 $6x^5$。

\section{最复杂的组合：全部混在一起}

\[
\sum_{n=1}^{2} \int_0^1 x^{n!} \, dx
\]

外面是连加，里面是积分，积分里面还有阶乘。

\subsection{从外往内理解}

$\sum_{n=1}^{2}$——连加，$n$ 从1到2，两站。

每一站要算的是 $\int_0^1 x^{n!} \, dx$——一个积分，而且指数是 $n!$。

\subsection{从内往外计算}

\textbf{第一站：$n = 1$}

$n! = 1! = 1$。

要算的积分是 $\int_0^1 x^1 \, dx = \int_0^1 x \, dx$。

三角形，底1高1，面积 $= \dfrac{1}{2}$。

这一站放进购物车的是 $\dfrac{1}{2}$。

\textbf{第二站：$n = 2$}

$n! = 2! = 2$。

要算的积分是 $\int_0^1 x^2 \, dx$。

$x^2$ 的积分是 $\dfrac{x^3}{3}$。

从0到1：$\dfrac{1}{3}$。

这一站放进购物车的是 $\dfrac{1}{3}$。

\subsection{加起来}

$\dfrac{1}{2} + \dfrac{1}{3}$

公分母是6。

$\dfrac{3}{6} + \dfrac{2}{6} = \dfrac{5}{6}$

答案是 $\dfrac{5}{6}$。

\section{翻译的秘诀}

不管式子多复杂，翻译的方法永远是一样的：

\textbf{第一步：找到最外层的符号，知道“最后要做什么操作”。}

是 $\sum$？那最后要加起来。

是 $\prod$？那最后要乘起来。

是 $\int$？那最后要算面积。

是 $\dfrac{d}{dx}$？那最后要求导。

\textbf{第二步：看里面套了什么，一层一层往里看。}

\textbf{第三步：从最里面开始算，算完一层往外走一层。}

就像剥洋葱，从外面看清楚结构，从里面开始下手。

\section{你来试几个}

\textbf{（1）}

\[
\sum_{i=1}^{2} i!
\]

$i = 1$：$1! = 1$

$i = 2$：$2! = 2$

加起来：$1 + 2 = 3$

\textbf{答案：3}

\textbf{（2）}

\[
\prod_{i=1}^{3} i
\]

$= 1 \times 2 \times 3 = 6$

\textbf{答案：6}（这就是 $3!$）

\textbf{（3）}

\[
\sum_{i=1}^{2} \int_0^1 x^i \, dx
\]

$i = 1$：$\int_0^1 x^1 \, dx = \dfrac{1}{2}$

$i = 2$：$\int_0^1 x^2 \, dx = \dfrac{1}{3}$

加起来：$\dfrac{1}{2} + \dfrac{1}{3} = \dfrac{5}{6}$

\textbf{答案：$\dfrac{5}{6}$}

\textbf{（4）}

\[
\frac{d}{dx} \left( \sum_{i=0}^{2} x^i \right)
\]

先算连加：$x^0 + x^1 + x^2 = 1 + x + x^2$

再求导：$0 + 1 + 2x = 1 + 2x$

\textbf{答案：$1 + 2x$}

\textbf{（5）}

\[
\int_0^2 \left( \prod_{i=1}^{2} x \right) dx
\]

先算连乘：$x \times x = x^2$

再积分：$\int_0^2 x^2 \, dx$

$x^2$ 的积分是 $\dfrac{x^3}{3}$。

从0到2：$\dfrac{8}{3} - 0 = \dfrac{8}{3}$

\textbf{答案：$\dfrac{8}{3}$}

\section{这一章你学到了什么}

\textbf{第一：} $\sum$、$\prod$、$\int$、$\dfrac{d}{dx}$、$n!$ 可以任意组合在一起。

\textbf{第二：} 不管怎么组合，方法都一样：从外往内理解，从内往外计算。

\textbf{第三：} 翻译的秘诀：先看最外层是什么操作，再一层一层往里看，最后从最里面开始算。

\begin{quote}
\textbf{一句话总结：数学符号就像积木，可以随便搭。不管搭得多复杂，你只要一层一层拆，就一定能算出来。}
\end{quote}

能跟到这里，说明你的脑子已经能处理相当复杂的数学语言了。

连加套积分、积分套连乘、求导套阶乘……不管怎么套，你都有办法下手。

这种“不管多复杂都不慌”的底气，比任何一个具体答案都值钱。

带着这份底气，我们往下走。

\chapter{对积分的拓展}

你已经知道积分是什么了。

$\int_a^b f(x) \, dx$——从 $a$ 到 $b$，算曲线下面的面积。

但这一章，我要把积分讲得更深、更透。

有几件事，我之前没说清楚，现在要好好说一说。

\section{先回答一个问题：$f(x)$ 是什么意思？}

你可能见过两种写法：

\[
\int_0^1 x \, dx
\]

\[
\int_0^1 f(x) \, dx
\]

一个写的是 $x$，一个写的是 $f(x)$。

有什么区别？

\subsection{$x$ 是一个具体的例子}

$\int_0^1 x \, dx$

这里的 $x$ 就是 $x$ 本身。

曲线是 $y = x$，一条从原点出发、斜着往上走的直线。

这是一个\textbf{具体的}函数。

\subsection{$f(x)$ 是一个“占位符”}

$\int_0^1 f(x) \, dx$

这里的 $f(x)$ 不是一个具体的函数。

它是一个\textbf{名字}，代表“某个函数”。

就像你说“某个人”，而不是说“张三”。

$f(x)$ 可以是 $x$，可以是 $x^2$，可以是 $2x + 1$，可以是任何函数。

它只是一个通用的写法，表示“这里有一个函数，但我现在不具体说是哪个”。

\subsection{为什么要用 $f(x)$？}

因为有时候我们想说一个\textbf{通用的规则}，适用于所有函数。

比如：

\[
\int_a^b f(x) \, dx = F(b) - F(a)
\]

这个公式对\textbf{任何函数}都成立。

不管 $f(x)$ 是 $x$、$x^2$、$x^3$，还是其他任何东西，这个规则都适用。

如果我写 $\int_a^b x \, dx = ?$，那就只是在说 $x$ 这一个函数。

如果我写 $\int_a^b f(x) \, dx = ?$，那就是在说\textbf{所有函数}。

\begin{quote}
\textbf{一句话总结：$x$、$x^2$、$2x+1$ 是具体的函数。$f(x)$ 是一个通用的名字，代表“某个函数”。}
\end{quote}

\section{再回答一个问题：为什么要“减去起点的值”？}

上一章我们算过这个：

\[
\int_0^2 x^2 \, dx
\]

答案是这样算的：

$x^2$ 的原函数是 $\dfrac{x^3}{3}$。

代入终点：$\dfrac{2^3}{3} = \dfrac{8}{3}$

代入起点：$\dfrac{0^3}{3} = \dfrac{0}{3} = 0$

相减：$\dfrac{8}{3} - 0 = \dfrac{8}{3}$

但你可能在想：\textbf{为什么要减？为什么不是直接用终点的值？}

\subsection{用一个生活例子来理解}

假设你从家出发去学校。

你家在“第2公里”的位置。

学校在“第7公里”的位置。

请问你走了多远？

$7 - 2 = 5$ 公里。

不是7公里，是 $7 - 2 = 5$ 公里。

因为你不是从0开始的，你是从2开始的。

积分也是一样的道理。

$\int_a^b f(x) \, dx$ 问的是：\textbf{从 $a$ 到 $b$ 这一段的面积是多少？}

如果你只算终点 $b$ 的值，那你算的是“从0到 $b$ 的面积”。

但题目问的是“从 $a$ 到 $b$”。

所以你要把“从0到 $b$”减去“从0到 $a$”，剩下的就是“从 $a$ 到 $b$”。

\subsection{为什么 $-0$ 经常直接省略？}

当起点是0的时候，比如 $\int_0^2 x^2 \, dx$：

代入起点：$\dfrac{0^3}{3} = 0$

$0$ 这个数，减去它没有任何影响。

$\dfrac{8}{3} - 0 = \dfrac{8}{3}$

所以有时候我们会直接写“$= \dfrac{8}{3}$”，不写“$- 0$”这一步。

因为减0没有意义，直接省略了。

但如果起点不是0呢？

\[
\int_1^3 x^2 \, dx
\]

代入终点：$\dfrac{3^3}{3} = \dfrac{27}{3} = 9$

代入起点：$\dfrac{1^3}{3} = \dfrac{1}{3}$

相减：$9 - \dfrac{1}{3} = \dfrac{27}{3} - \dfrac{1}{3} = \dfrac{26}{3}$

这次就不能省略了，因为起点的值不是0。

\begin{quote}
\textbf{一句话总结：积分要“减去起点的值”，因为我们要的是“从起点到终点这一段”的面积，不是“从0到终点”的面积。当起点是0时，减0没有影响，可以省略。}
\end{quote}

\section{两种积分：定积分和不定积分}

现在来讲一个新东西。

你见过这种写法吗？

\[
\int x^2 \, dx
\]

咦？上面和下面的数字呢？

没有起点，没有终点。

这是什么意思？

\subsection{定积分：有起点有终点}

\[
\int_0^3 x^2 \, dx
\]

这叫\textbf{定积分}。

“定”在哪里？定在起点和终点。

你明确知道从哪开始（0），到哪结束（3）。

算出来的结果是一个\textbf{具体的数}——那块区域的面积。

\subsection{不定积分：没有起点没有终点}

\[
\int x^2 \, dx
\]

这叫\textbf{不定积分}。

没有起点，没有终点。

不“定”。

算出来的结果\textbf{不是一个数}，而是\textbf{一个函数}。

\subsection{不定积分在做什么？}

不定积分在问一个问题：

\textbf{什么函数，求导之后会变成 $x^2$？}

你知道 $(x^3)' = 3x^2$。（$x^3$ 求导是 $3x^2$。）

那如果我想要求导之后变成 $x^2$（不是 $3x^2$），原函数应该是什么？

应该是 $\dfrac{x^3}{3}$。

因为 $\left(\dfrac{x^3}{3}\right)' = \dfrac{3x^2}{3} = x^2$。

所以：

\[
\int x^2 \, dx = \frac{x^3}{3} + C
\]

\subsection{那个 $C$ 是什么？为什么要加 $+C$？}

你可能注意到，答案后面多了一个 $+C$。

这是什么东西？

我们来看一个现象：

$x^2$ 求导，得到 $2x$。

$x^2 + 1$ 求导，也得到 $2x$。

$x^2 + 5$ 求导，还是得到 $2x$。

$x^2 + 100$ 求导，依然是 $2x$。

$x^2 - 37$ 求导，仍然是 $2x$。

不管后面加的是什么常数，求导之后都变成了 $2x$。

\textbf{常数不见了。}

\subsection{为什么常数会消失？}

常数就是\textbf{不变的数}。

$5$ 永远是 $5$。不管 $x$ 怎么变，$5$ 还是 $5$。

求导是在问：\textbf{“当 $x$ 变化的时候，$y$ 怎么变？”}

常数变了吗？没有。

既然常数没有变化，那它的“变化率”就是 \textbf{0}。

所以：\textbf{任何常数，求导都是0。}

\subsection{求导是一个“丢信息”的操作}

这就是根本原因：

\textbf{求导会把常数“丢掉”。}

$x^2 + 5$ 求导，变成 $2x + 0 = 2x$。

那个 $5$ 呢？变成了 $0$，消失了。

你看着 $2x$，你不知道它原来是 $x^2$，还是 $x^2 + 5$，还是 $x^2 + 100$。

\textbf{信息丢失了。}

\subsection{积分是“倒回去”}

不定积分是求导的反操作。

求导是“倒掉”，积分是“倒回去”。

用一个比喻来理解：

想象你有三杯水。

第一杯：纯净水。

第二杯：纯净水 + 1勺糖。

第三杯：纯净水 + 3勺糖。

现在你把三杯水都倒掉。剩下三个一模一样的空杯子。

你看着这三个空杯子，能知道原来哪杯加了糖吗？\textbf{不能。}

因为倒掉的时候，糖也跟着倒掉了。

求导就像“倒掉水”，常数就像“糖”。

现在你想把水“倒回去”（积分）。

你只能说：

“我倒的是纯净水，\textbf{可能还有一些糖，我不知道有多少，就用 $C$ 来代表吧}。”

这就是 $+C$ 的根本原因：

\textbf{求导会把常数“丢掉”。积分是求导的反操作，但被丢掉的常数找不回来了，所以要加一个 $+C$，表示“可能有一个未知的常数”。}

\section{特殊的数学常数：$e$}

在讲更多积分之前，我要先介绍一个数字：

\[
e \approx 2.71828...
\]

它叫做\textbf{自然常数}，或者\textbf{欧拉数}。

\subsection{$e$ 是什么？}

你知道 $\pi \approx 3.14159...$，它和圆有关。

$e$ 也是一个像 $\pi$ 一样重要的数字。

它也是一个无限不循环的小数，永远写不完。

\subsection{$e$ 从哪来的？}

想象你把1块钱存进银行。

银行说：年利率100\%，一年后你有2块钱。

但如果银行说：半年算一次利息，每次50\%呢？

半年后：$1 + 0.5 = 1.5$ 块

一年后：$1.5 + 1.5 \times 0.5 = 1.5 + 0.75 = 2.25$ 块

比2块多了！

如果一年算4次利息，每次25\%呢？

第1次：$1 \times 1.25 = 1.25$

第2次：$1.25 \times 1.25 = 1.5625$

第3次：$1.5625 \times 1.25 \approx 1.95$

第4次：$1.95 \times 1.25 \approx 2.44$

更多了！

如果一年算12次利息呢？$\approx 2.61$

如果一年算365次利息呢？$\approx 2.71$

如果\textbf{无限次}算利息呢？

答案趋近于一个固定的数：$e \approx 2.71828...$

当你把利息分成无限多次来算，最后得到的就是 $e$。

\subsection{$e$ 有什么特别的？}

$e$ 最神奇的地方是：

\[
\frac{d}{dx}(e^x) = e^x
\]

$e^x$ 这个函数，求导之后\textbf{还是它自己}！

这在所有函数里是独一无二的。

没有任何其他函数有这个性质。

同样的：

\[
\int e^x \, dx = e^x + C
\]

$e^x$ 的积分还是 $e^x$（加上常数 $C$）。

\textbf{$e^x$ 是唯一一个“求导不变、积分不变”的函数。}

\subsection{常见的数学常数}

现在你认识了两个重要的常数：

\begin{center}
\begin{tabular}{llll}
\toprule
符号 & 名字 & 大约等于 & 和什么有关 \\
\midrule
$\pi$ & 圆周率 & 3.14159... & 圆 \\
$e$ & 自然常数 & 2.71828... & 增长、变化 \\
\bottomrule
\end{tabular}
\end{center}

以后你还会遇到更多常数，但这两个是最重要的。

\section{积分符号上面可以写什么？}

到目前为止，你见过的积分，上面和下面都是普通的数字：

\[
\int_0^3 x \, dx
\]

从0到3。

但有时候，上面或下面会出现一些奇怪的东西。

\subsection{上面写 $\infty$（无穷大）}

\[
\int_1^{\infty} \frac{1}{x^2} \, dx
\]

这是什么意思？

从 $x = 1$ 开始，一直积到\textbf{无穷远}。

没有终点。一直走，走到天涯海角。

这叫\textbf{反常积分}，或者\textbf{无穷积分}。

\subsection{这能算吗？无穷大不是无限大吗？}

你可能会想：积到无穷远，面积不会是无限大吗？

\textbf{不一定！}

有些曲线，虽然一直延伸到无穷远，但它下面的面积是\textbf{有限的}。

比如 $y = \dfrac{1}{x^2}$：

当 $x = 1$ 时，$y = 1$。

当 $x = 2$ 时，$y = \dfrac{1}{4} = 0.25$。

当 $x = 10$ 时，$y = \dfrac{1}{100} = 0.01$。

当 $x = 100$ 时，$y = \dfrac{1}{10000} = 0.0001$。

你看，$x$ 越来越大，$y$ 越来越小，\textbf{迅速趋近于0}。

曲线越往右，越贴近 $x$ 轴。

贴得太紧了，紧到虽然曲线无限延伸，但下面的面积加起来还是有限的。

$\int_1^{\infty} \dfrac{1}{x^2} \, dx = 1$

答案就是 \textbf{1}。

神奇吧？无限长的曲线，面积却是1。

\subsection{下面写 $-\infty$（负无穷大）}

\[
\int_{-\infty}^{0} e^x \, dx
\]

从负无穷远开始，积到0。

起点在无穷远的左边。

这也是反常积分。

这个积分的答案也是有限的，等于 \textbf{1}。

\section{$\pi$ 和积分的关系}

现在来到最有趣的部分。

你以为 $\pi$ 只是“圆的周长除以直径”。

但 $\pi$ 可以用积分来表示！

\subsection{第一种：用圆的面积}

圆心在原点、半径为1的圆，方程是：

\[
x^2 + y^2 = 1
\]

解出 $y$：

\[
y = \sqrt{1 - x^2}
\]

这是圆的\textbf{上半部分}。

上半圆的面积是 $\int_{-1}^{1} \sqrt{1 - x^2} \, dx$。

整个圆的面积是上半圆的两倍：

\[
\text{圆的面积} = \int_{-1}^{1} 2\sqrt{1 - x^2} \, dx
\]

而我们知道，半径为1的圆的面积是 $\pi$。

所以：

\[
\pi = \int_{-1}^{1} 2\sqrt{1 - x^2} \, dx
\]

\textbf{$\pi$ 可以写成一个积分！}

\subsection{第二种：高斯积分}

有一个非常有名的积分：

\[
\int_{-\infty}^{\infty} e^{-x^2} \, dx = \sqrt{\pi}
\]

这个积分叫做\textbf{高斯积分}。

它的答案是 $\sqrt{\pi}$，也就是 $\pi$ 的平方根。

$e^{-x^2}$ 是一个钟形曲线——中间高，两边低，像一座山。

它和圆没有任何直观的联系，但算出来的答案，居然包含 $\pi$！

这就是 $\pi$ 神奇的地方：\textbf{它不只是关于圆的。}

$\pi$ 像一个幽灵，会在概率论、物理学等各种地方冒出来。积分让你看到了 $\pi$ 的另一面。

\subsection{第三种：用无穷级数}

还有一种表示 $\pi$ 的方法，把 $\sum$ 和 $\pi$ 连在一起：

\[
\frac{\pi}{4} = \sum_{n=0}^{\infty} \frac{(-1)^n}{2n+1}
\]

展开来看：

\[
\frac{\pi}{4} = 1 - \frac{1}{3} + \frac{1}{5} - \frac{1}{7} + \frac{1}{9} - \frac{1}{11} + \cdots
\]

这是一个\textbf{无穷连加}。

两边乘以4：

\[
\pi = 4 \times \left(1 - \frac{1}{3} + \frac{1}{5} - \frac{1}{7} + \cdots\right)
\]

\textbf{$\pi$ 可以用一个无穷连加来表示！}

\section{积分还能做什么？}

除了算面积，积分还能算很多东西。

\subsection{算体积}

你有一个立体图形，比如一个球。

把它切成无数个无限薄的圆片。

每一片是一个圆，有一个面积。

面积乘以厚度（$dx$），得到那一片的体积。

把所有片的体积加起来，就是整个球的体积。

球的体积公式 $\dfrac{4}{3}\pi r^3$，就是这样用积分推出来的。

\subsection{算曲线的长度}

一条弯弯曲曲的曲线，有多长？

把它切成无数个无限短的小段。

每一小段近似是直的，可以用勾股定理算长度。

把所有小段的长度加起来，就是曲线的总长度。

这也是积分。

\subsection{算距离}

一辆车在加速。

你知道它每一瞬间的速度 $v(t)$。

从时间 $t = 0$ 到 $t = 10$ 秒，它总共走了多远？

\[
\text{距离} = \int_0^{10} v(t) \, dt
\]

把每一瞬间的速度乘以每一瞬间的时间（$dt$），加起来，就是总距离。

\section{积分的本质}

说到底，积分在做一件事：

\textbf{把无限多个无限小的东西加起来。}

这些“东西”可以是面积、体积、长度、距离。

只要你能把问题拆成“无限多个无限小的部分”，你就可以用积分来解决它。

\begin{quote}
\textbf{一句话总结：积分是“无限细的加法”。}
\end{quote}

\section{这一章你学到了什么}

\textbf{第一：} $f(x)$ 是一个通用的名字，代表“某个函数”。

\textbf{第二：} 定积分要“减去起点的值”，因为我们要的是“从起点到终点”这一段。

\textbf{第三：} 定积分算出一个数。不定积分找原函数，答案要加 $+C$（因为求导会把常数丢掉，找不回来了）。

\textbf{第四：} $e \approx 2.71$ 是自然常数。$e^x$ 求导和积分都是它自己。

\textbf{第五：} 积分的上下限可以是无穷大。有些无穷积分的结果是有限的。

\textbf{第六：} $\pi$ 可以用积分或无穷连加来表示，它不只是关于圆的。

\textbf{第七：} 积分的本质是“把无限小的东西加起来”，能算面积、体积、长度、距离。

你现在对积分有了更深的理解。

你认识了新的常数 $e$，知道了 $\pi$ 和积分的奇妙联系，也明白了为什么要加 $+C$。

数学的世界里，很多看似不相关的东西，其实都连在一起。

带着这些知识，咱们接着看更有趣的。

\chapter{特别拓展——积分、\texorpdfstring{$\pi$}{pi} 与图像的故事}

你知道 $\pi$ 可以用积分来表示。

但还有一种更直接的用法：\textbf{$\pi$ 直接出现在积分的上下限里}。

在讲这些有趣的式子之前，我要先让你知道这些函数画出来是什么样子。

不然你看到 $\sin x$、$\cos x$、$e^{-x^2}$，脑子里一片空白，后面就跟不上了。

\section{建立画面：这些函数长什么样？}

\subsection{\texorpdfstring{$y = \sin x$}{y = sin x} 长什么样？}

一条波浪线。

想象大海的波浪，一上一下，一上一下，永远重复。

从 $x = 0$ 开始：

- 曲线从原点 $(0, 0)$ 出发。
- 往上升，升到最高点，高度是 1。（这时候 $x = \frac{\pi}{2}$。）
- 然后往下降，降回到 0。（这时候 $x = \pi$。）
- 继续往下降，降到最低点，高度是 $-1$。（这时候 $x = \frac{3\pi}{2}$。）
- 再往上升，回到 0。（这时候 $x = 2\pi$。）

然后重复。上、下、上、下，永远不停。

\textbf{一句话：$\sin x$ 是一条上下起伏的波浪线，最高1，最低 $-1$，每隔 $2\pi$ 重复一次。}

\subsection{\texorpdfstring{$y = \cos x$}{y = cos x} 长什么样？}

也是一条波浪线，但起点不一样。

$\sin x$ 从0开始。

$\cos x$ 从1开始。

从 $x = 0$ 开始：

- 曲线从 $(0, 1)$ 出发——最高点。
- 往下降，降到 0。（这时候 $x = \frac{\pi}{2}$。）
- 继续往下降，降到最低点 $-1$。（这时候 $x = \pi$。）
- 再往上升，回到 0。（这时候 $x = \frac{3\pi}{2}$。）
- 继续往上升，回到最高点 1。（这时候 $x = 2\pi$。）

然后重复。

\textbf{一句话：$\cos x$ 也是一条波浪线，和 $\sin x$ 形状一样，只是往左移了 $\frac{\pi}{2}$。}

\subsection{\texorpdfstring{$y = e^{-x^2}$}{y = e^{-x^2}} 长什么样？}

一座钟形的山丘。

想象一顶帽子，或者一个倒扣的碗，但碗底是圆滑的。

从 $x = 0$ 开始：

- 曲线在最高点，高度是 1。
- 往右走，曲线往下降，而且降得越来越快。
- $x = 1$ 时，高度大约是 $0.37$。
- $x = 2$ 时，高度大约是 $0.02$。
- $x = 3$ 时，高度大约是 $0.0001$——几乎贴着 $x$ 轴了。

往左走，形状完全一样（因为 $(-x)^2 = x^2$，曲线左右对称）。

\textbf{一句话：$e^{-x^2}$ 是一座钟形山丘，中间高两边低，越往外越贴近 $x$ 轴，但永远不碰到。}

\section{积分和 $\pi$ 直接连用}

脑子里有了图像，我们再来看积分。

\subsection{最常见的例子}

\[
\int_0^{\pi} \sin x \, dx
\]

这个式子在说什么？

我们一个一个拆：

$\int$ → 积分，把小块面积加起来

$_0^{\pi}$ → 从 0 到 $\pi$

$\sin x$ → $x$ 的正弦（这是曲线的高度）

$dx$ → 每一小条的宽度

直接语言：\textbf{从 0 到 $\pi$，算 $\sin x$ 曲线下面的面积。}

\subsection{在图上是什么样子？}

从 0 到 $\pi$ 这一段，$\sin x$ 刚好是一个\textbf{拱形}——像一座小山丘，或者像彩虹的一半。从地面升起，到达山顶，再落回地面。

这块拱形区域的面积是多少？

\[
\int_0^{\pi} \sin x \, dx = 2
\]

不多不少，正好是 \textbf{2}。

\subsection{为什么上限是 $\pi$？}

因为 $\sin x$ 和圆有关系。它的周期是 $2\pi$。

从 0 到 $\pi$，刚好是\textbf{半个波浪}。

当你算和波浪（圆的投射）有关的面积时，$\pi$ 自然就会出现在边界里。

\subsection{再来一个例子}

\[
\int_0^{2\pi} \sin x \, dx
\]

这次是从 0 到 $2\pi$，一整个波浪。

\subsection{在图上是什么样子？}

前半段（0 到 $\pi$）：曲线在 $x$ 轴\textbf{上面}，形成一个拱形，面积是正的（+2）。

后半段（$\pi$ 到 $2\pi$）：曲线在 $x$ 轴\textbf{下面}，形成一个倒过来的拱形，面积是负的（-2）。

\textbf{两块面积加起来：}

$2 + (-2) = 0$

\[
\int_0^{2\pi} \sin x \, dx = 0
\]

正的和负的刚好抵消了。

\subsection{还有这个}

\[
\int_0^{\frac{\pi}{2}} \cos x \, dx
\]

从 0 到 $\frac{\pi}{2}$（也就是 $\pi$ 的一半，90度）。

\subsection{在图上是什么样子？}

$\cos x$ 从最高点 $(0,1)$ 开始，平滑地往下弯，降落到 $(\frac{\pi}{2}, 0)$。

像一个滑梯，从高处滑到地面。这块区域像一个被切掉一角的方块。

\subsection{面积是多少？}

\[
\int_0^{\frac{\pi}{2}} \cos x \, dx = 1
\]

这一块的面积正好是 \textbf{1}。

\subsection{更神奇的例子：高斯积分}

\[
\int_{-\infty}^{\infty} e^{-x^2} \, dx = \sqrt{\pi}
\]

\subsection{在图上是什么样子？}

这是一座无限宽的\textbf{钟形山丘}。

虽然曲线往两边无限延伸，但因为它贴 $x$ 轴贴得太紧了，整个山丘下面的面积是有限的。

\subsection{面积是多少？}

\[
\int_{-\infty}^{\infty} e^{-x^2} \, dx = \sqrt{\pi} \approx 1.77
\]

这座无限宽的山丘，下面的面积是 $\sqrt{\pi}$。

这里 $\pi$ 不是出现在上下限，而是出现在\textbf{答案}里。

而且，这个钟形曲线和圆看起来没有任何直观的关系，但 $\pi$ 还是冒出来了。

\section{这一章你学到了什么}

\textbf{第一：} 脑子里有了 $\sin x$（波浪）、$\cos x$（平移的波浪）、$e^{-x^2}$（钟形山丘）的图像。

\textbf{第二：} $\int_0^{\pi} \sin x \, dx$ 就是一个拱形区域的面积，等于 2。

\textbf{第三：} $\int_0^{2\pi} \sin x \, dx$ 是一个拱形加一个倒拱形，正负抵消，等于 0。

\textbf{第四：} $\int_{-\infty}^{\infty} e^{-x^2} \, dx$ 是一座无限宽但面积有限的钟形山丘，它的面积里藏着 $\pi$（$\sqrt{\pi}$）。

\textbf{第五：} $\pi$ 可以出现在积分的上下限里，也可以出现在答案里。因为 $\pi$ 是圆的语言，而很多数学结构都和圆有着隐藏的深刻联系。

数学最迷人的地方就在这里：原本以为毫不相干的东西（比如一座无限宽的钟形山和圆周率），在积分的透镜下，神奇地走到了一起。

\chapter{\texorpdfstring{$\sin$、$\cos$、$\tan$}{sin, cos, tan}，难道只能求内角吗？}

你在初中学过三角函数。

$\sin$ = 对边比斜边。

$\cos$ = 邻边比斜边。

$\tan$ = 对边比邻边。

你用它们来做什么？

\textbf{已知角度，求边长。}

比如：角度是30度，斜边是10，对边是多少？

$\sin 30^\circ = \dfrac{\text{对边}}{10}$

$\text{对边} = 10 \times \sin 30^\circ = 10 \times 0.5 = 5$

但你有没有想过反过来？

\textbf{如果我已经知道边长，想求角度呢？}

\section{反过来的问题}

假设你知道：

对边是5，斜边是10。

你想求角度。

你知道 $\sin \theta = \dfrac{5}{10} = 0.5$。

但你想知道的是：\textbf{$\theta$ 到底是多少度？}

$\sin$ 什么角度等于0.5？

你可能记住了：$\sin 30^\circ = 0.5$。

所以 $\theta = 30^\circ$。

但如果题目告诉你 $\sin \theta = 0.73$ 呢？

你背不出哪个角度的 $\sin$ 是0.73。

怎么办？

\section{反三角函数：反过来求角度}

这就是\textbf{反三角函数}的用处。

$\sin$ 是“给我角度，我告诉你比值”。

$\arcsin$ 是“给我比值，我告诉你角度”。

完全反过来。

\subsection{\texorpdfstring{$\arcsin$}{arcsin}——正弦的反函数}

\[
\arcsin(0.5) = 30^\circ
\]

读法：“0.5的反正弦等于30度。”

意思是：“哪个角度的正弦等于0.5？答案是30度。”

\[
\arcsin(0.73) = ?
\]

你用计算器按一下，得到大约 $46.9^\circ$。

意思是：“哪个角度的正弦等于0.73？答案是大约46.9度。”

再来几个：

$\arcsin(0) = 0^\circ$

（哪个角度的正弦等于0？0度。因为 $\sin 0^\circ = 0$。）

$\arcsin(1) = 90^\circ$

（哪个角度的正弦等于1？90度。因为 $\sin 90^\circ = 1$。）

$\arcsin(0.5) = 30^\circ$

（因为 $\sin 30^\circ = 0.5$。）

$\arcsin\left(\dfrac{\sqrt{2}}{2}\right) = 45^\circ$

（因为 $\sin 45^\circ = \dfrac{\sqrt{2}}{2}$。）

$\arcsin\left(\dfrac{\sqrt{3}}{2}\right) = 60^\circ$

（因为 $\sin 60^\circ = \dfrac{\sqrt{3}}{2}$。）

\begin{quote}
\textbf{一句话：$\arcsin$ 就是 $\sin$ 反过来。$\sin$ 把角度变成比值，$\arcsin$ 把比值变回角度。}
\end{quote}

\subsection{\texorpdfstring{$\arccos$}{arccos}——余弦的反函数}

\[
\arccos(0.5) = 60^\circ
\]

读法：“0.5的反余弦等于60度。”

意思是：“哪个角度的余弦等于0.5？答案是60度。”

（因为 $\cos 60^\circ = 0.5$。）

再来几个：

$\arccos(1) = 0^\circ$

（因为 $\cos 0^\circ = 1$。）

$\arccos(0) = 90^\circ$

（因为 $\cos 90^\circ = 0$。）

$\arccos(-1) = 180^\circ$

（因为 $\cos 180^\circ = -1$。）

$\arccos\left(\dfrac{\sqrt{2}}{2}\right) = 45^\circ$

（因为 $\cos 45^\circ = \dfrac{\sqrt{2}}{2}$。）

$\arccos\left(\dfrac{\sqrt{3}}{2}\right) = 30^\circ$

（因为 $\cos 30^\circ = \dfrac{\sqrt{3}}{2}$。）

\begin{quote}
\textbf{一句话：$\arccos$ 就是 $\cos$ 反过来。$\cos$ 把角度变成比值，$\arccos$ 把比值变回角度。}
\end{quote}

\subsection{\texorpdfstring{$\arctan$}{arctan}——正切的反函数}

\[
\arctan(1) = 45^\circ
\]

读法：“1的反正切等于45度。”

意思是：“哪个角度的正切等于1？答案是45度。”

（因为 $\tan 45^\circ = 1$。对边和邻边一样长，坡度刚好是1。）

再来几个：

$\arctan(0) = 0^\circ$

（因为 $\tan 0^\circ = 0$。平地，没有坡度。）

$\arctan\left(\dfrac{1}{\sqrt{3}}\right) = 30^\circ$

（因为 $\tan 30^\circ = \dfrac{1}{\sqrt{3}}$。）

$\arctan(\sqrt{3}) = 60^\circ$

（因为 $\tan 60^\circ = \sqrt{3}$。）

\begin{quote}
\textbf{一句话：$\arctan$ 就是 $\tan$ 反过来。$\tan$ 把角度变成坡度，$\arctan$ 把坡度变回角度。}
\end{quote}

\section{另一种写法}

你可能见过两种写法：

$\arcsin(x)$ 和 $\sin^{-1}(x)$

它们是\textbf{同一个东西}。

$\sin^{-1}$ 不是“$\sin$ 的负一次方”，不是 $\dfrac{1}{\sin x}$。

$\sin^{-1}$ 就是 $\arcsin$，反正弦。

那个 $-1$ 是“反函数”的标记，不是“负一次方”。

同样：

$\cos^{-1}(x) = \arccos(x)$

$\tan^{-1}(x) = \arctan(x)$

这两种写法完全等价，看到哪种都不要慌。

\section{反三角函数能吃什么？能吐什么？}

\subsection{$\arcsin(x)$}

能吃什么？

$x$ 必须在 $-1$ 到 $1$ 之间。

因为 $\sin$ 的值永远在 $-1$ 到 $1$ 之间——你没法找到一个角度使得 $\sin$ 等于2或者等于 $-3$。

所以 $\arcsin(2)$ 没有意义。$\arcsin(-3)$ 也没有意义。

只有 $-1 \le x \le 1$ 的时候，$\arcsin(x)$ 才有意义。

能吐什么？

$\arcsin(x)$ 吐出来的角度，在 $-90^\circ$ 到 $90^\circ$ 之间（也就是 $-\dfrac{\pi}{2}$ 到 $\dfrac{\pi}{2}$）。

\subsection{$\arccos(x)$}

能吃什么？

也是 $-1 \le x \le 1$。理由和 $\arcsin$ 一样。

能吐什么？

$\arccos(x)$ 吐出来的角度，在 $0^\circ$ 到 $180^\circ$ 之间（也就是 $0$ 到 $\pi$）。

\subsection{$\arctan(x)$}

能吃什么？

任何数都行！

$\tan$ 的值可以是任何数（从负无穷到正无穷），所以 $\arctan$ 可以吃任何数。

$\arctan(100)$？可以。

$\arctan(-5000)$？也可以。

能吐什么？

$\arctan(x)$ 吐出来的角度，在 $-90^\circ$ 到 $90^\circ$ 之间（不包括 $-90^\circ$ 和 $90^\circ$）。

\subsection{总结一下}

\begin{center}
\begin{tabular}{lll}
\toprule
函数 & 能吃什么 & 能吐什么 \\
\midrule
$\arcsin(x)$ & $-1 \le x \le 1$ & $-90^\circ$ 到 $90^\circ$ \\
$\arccos(x)$ & $-1 \le x \le 1$ & $0^\circ$ 到 $180^\circ$ \\
$\arctan(x)$ & 任何数 & $-90^\circ$ 到 $90^\circ$（不含端点） \\
\bottomrule
\end{tabular}
\end{center}

\section{用在实际问题里}

学了这些，有什么用？

\subsection{例子一：求坡度对应的角度}

一条路的坡度是0.3。

坡度就是 $\tan$——“爬了多高除以走了多远”。

$\tan \theta = 0.3$

$\theta = \arctan(0.3)$

用计算器算一下：$\theta \approx 16.7^\circ$

这条路的倾斜角大约是16.7度。

\subsection{例子二：求梯子靠墙的角度}

一把梯子靠在墙上。

梯子长5米（斜边），梯子底部离墙脚3米（邻边）。

梯子和地面的角度是多少？

$\cos \theta = \dfrac{\text{邻边}}{\text{斜边}} = \dfrac{3}{5} = 0.6$

$\theta = \arccos(0.6)$

用计算器：$\theta \approx 53.1^\circ$

\subsection{例子三：导航和方向}

你在地图上，从A点到B点。

A点坐标是 $(0, 0)$，B点坐标是 $(3, 4)$。

你要往什么方向走？

方向角 $\theta = \arctan\left(\dfrac{4}{3}\right)$

$\theta \approx 53.1^\circ$

你要往北偏东大约53度的方向走。

\section{一个容易搞混的地方}

$\sin$、$\cos$、$\tan$ 是“角度 $\rightarrow$ 比值”。

$\arcsin$、$\arccos$、$\arctan$ 是“比值 $\rightarrow$ 角度”。

方向正好相反。

你可以这样记：

\begin{quote}
\textbf{$\sin$ 是问“这个角度的比值是多少”。}

\textbf{$\arcsin$ 是问“这个比值对应的角度是多少”。}
\end{quote}

$\sin$ 是正着走。

$\arcsin$ 是倒着走。

“arc”这个词，就是“反”的意思。

\section{这一章你学到了什么}

\textbf{第一：} 三角函数可以“反过来”用。已知比值，反推角度。

\textbf{第二：} $\arcsin$ 是 $\sin$ 的反函数，$\arccos$ 是 $\cos$ 的反函数，$\arctan$ 是 $\tan$ 的反函数。

\textbf{第三：} $\sin^{-1}$ 和 $\arcsin$ 是同一个东西，那个 $-1$ 是“反函数”的标记，不是“负一次方”。

\textbf{第四：} $\arcsin$ 和 $\arccos$ 只能吃 $-1$ 到 $1$ 之间的数，$\arctan$ 什么数都能吃。

\textbf{第五：} 反三角函数在实际中很有用——求坡度角、梯子角度、导航方向……

你现在会“正着用”三角函数，也会“反着用”了。

正着用：给角度，算比值。

反着用：给比值，算角度。

两个方向都通了，三角函数你就真正吃透了。

\chapter{三元——多了一个未知数，怎么办？}

你在初中学过方程。

一个未知数：

\[
x + 3 = 7
\]

解出来，$x = 4$。

然后你学了两个未知数：

\[
\begin{cases} x + y = 5 \\ x - y = 1 \end{cases}
\]

解出来，$x = 3$，$y = 2$。

现在，再多一个未知数：

\[
\begin{cases} x + y + z = 6 \\ 2x + y - z = 1 \\ x - y + 2z = 5 \end{cases}
\]

三个未知数，三个方程。

别慌。

方法和二元方程组\textbf{完全一样}——只是多消一步。

\section{先从生活场景开始}

想象你去一家文具店买东西。

你买了三种东西：钢笔、本子、橡皮。

\textbf{第一次购物：}

你买了1支钢笔、1本本子、1块橡皮，花了6块钱。

\textbf{第二次购物：}

你买了2支钢笔、1本本子，没买橡皮，花了5块钱。

\textbf{第三次购物：}

你没买钢笔，买了1本本子、2块橡皮，花了5块钱。

现在我问你：钢笔多少钱一支？本子多少钱一本？橡皮多少钱一块？

设钢笔的价格是 $x$ 元，本子的价格是 $y$ 元，橡皮的价格是 $z$ 元。

（$x$ 就是钢笔的单价！$y$ 就是本子的单价！$z$ 就是橡皮的单价！）

把三次购物翻译成方程：

\[
\begin{cases} x + y + z = 6 & \cdots①\\ 2x + y = 5 & \cdots②\\ y + 2z = 5 & \cdots③ \end{cases}
\]

这就是一个三元方程组。

\section{三元方程组长什么样？}

三元方程组有三个特点：

\textbf{第一：三个未知数。}

$x$、$y$、$z$——三个不同的东西，三个不同的价格。

\textbf{第二：三个方程。}

三条信息，三个限制条件。

\textbf{第三：三个方程要同时成立。}

三次购物的价格都是真的，必须同时满足。

\section{怎么解？}

解三元方程组，方法叫\textbf{消元法}。

目标很简单：

\textbf{把三个未知数，一个一个消掉，最后变成一个方程一个未知数，就能解了。}

\section{尝试解第一道题}

\[
\begin{cases} x + y + z = 6 & \cdots①\\ 2x + y = 5 & \cdots②\\ y + 2z = 5 & \cdots③ \end{cases}
\]

\textbf{第一步：消掉一个未知数，把三元变二元。}

我们先消掉 $z$。

因为方程②里已经没有 $z$ 了，我们只要用方程①和③再凑出一个没有 $z$ 的方程，就能和方程②组成一个二元方程组。

方程①：$x + y + z = 6$

方程③：$y + 2z = 5$

为了消掉 $z$，把方程①乘以2：

$2x + 2y + 2z = 12 \quad \cdots④$

用方程④减去方程③：

$(2x + 2y + 2z) - (y + 2z) = 12 - 5$

\[
2x + y = 7 \quad \cdots⑤
\]

现在我们有两个只含 $x$ 和 $y$ 的方程：

方程②：$2x + y = 5$

方程⑤：$2x + y = 7$

等一下……

$2x + y$ 同时等于5又等于7？

这不可能！$5 \neq 7$。

这说明什么？

\textbf{这个方程组无解！}

也就是说，不存在一组钢笔、本子、橡皮的价格，能同时满足这三次购物记录。可能是我记错账了，或者店员算错钱了。

这也说明，解方程能帮我们\textbf{发现错误}。

\section{换一个有解的例子}

我们来看一开始给出的那个方程组：

\[
\begin{cases} x + y + z = 6 & \cdots① \\ 2x + y - z = 1 & \cdots② \\ x - y + 2z = 5 & \cdots③ \end{cases}
\]

\subsection{第一步：消掉 $z$，把三元变二元}

我们要把 $z$ 从所有方程里消掉，凑出两个只含 $x$ 和 $y$ 的方程。

\textbf{先用方程①和方程②消 $z$：}

方程①：$x + y + z = 6$

方程②：$2x + y - z = 1$

直接相加（因为 $+z$ 和 $-z$ 刚好抵消）：

$(x + y + z) + (2x + y - z) = 6 + 1$

\[
3x + 2y = 7 \quad \cdots④
\]

很好，得到第一个没有 $z$ 的方程了。

\textbf{再用方程①和方程③消 $z$：}

方程①：$x + y + z = 6$

方程③：$x - y + 2z = 5$

把方程①乘以2，让 $z$ 变成 $2z$：

$2x + 2y + 2z = 12 \quad \cdots⑤$

用方程⑤减去方程③：

$(2x + 2y + 2z) - (x - y + 2z) = 12 - 5$

$2x - x + 2y - (-y) + 2z - 2z = 7$

\[
x + 3y = 7 \quad \cdots⑥
\]

很好，得到第二个没有 $z$ 的方程了。

现在，我们把原来的三元方程组，成功变成了一个二元方程组：

\[
\begin{cases} 3x + 2y = 7 & \cdots④ \\ x + 3y = 7 & \cdots⑥ \end{cases}
\]

\subsection{第二步：消掉 $y$，把二元变一元}

现在解这个二元方程组。我们来消掉 $y$。

方程④：$3x + 2y = 7$ （乘以3） $\rightarrow 9x + 6y = 21 \quad \cdots⑦$

方程⑥：$x + 3y = 7$ （乘以2） $\rightarrow 2x + 6y = 14 \quad \cdots⑧$

用方程⑦减去方程⑧：

$(9x + 6y) - (2x + 6y) = 21 - 14$

\[
7x = 7
\]

\subsection{第三步：逐个解出未知数}

既然 $7x = 7$，那么：

\textbf{$x = 1$}

把 $x = 1$ 代入二元方程⑥（$x + 3y = 7$）：

$1 + 3y = 7$

$3y = 6$

\textbf{$y = 2$}

把 $x = 1$ 和 $y = 2$ 代入最开始的三元方程①（$x + y + z = 6$）：

$1 + 2 + z = 6$

$3 + z = 6$

\textbf{$z = 3$}

解出来了！

\[
x = 1, \quad y = 2, \quad z = 3
\]

\subsection{第四步：必须验证！}

解出来不算完。必须把答案代回另外两个原方程，逐一验证。

\textbf{验证方程②：$2x + y - z = 1$}

$2(1) + 2 - 3 = 2 + 2 - 3 = 1$ ✓

\textbf{验证方程③：$x - y + 2z = 5$}

$1 - 2 + 2(3) = 1 - 2 + 6 = 5$ ✓

全部验证通过！答案是正确的。

\section{整理一下解题步骤}

解三元方程组，永远是这四步：

\textbf{第一步：消掉一个未知数，把三元变二元。}

找两对方程，各自消掉同一个未知数，得到两个二元方程。

\textbf{第二步：消掉另一个未知数，把二元变一元。}

用刚才得到的两个二元方程，消掉一个未知数，得到一个一元方程。

\textbf{第三步：解出第一个未知数，然后一层层代回去，求出另外两个。}

\textbf{第四步：验证。}

把三个答案代回原方程，全部验证通过才算完。

\section{消元的时候，怎么选择消哪个未知数？}

答案是：\textbf{选最好消的那个。}

什么叫“最好消的”？

就是系数最简单的（比如系数是1或-1），或者有些方程里本来就已经没有这个未知数了。

\textbf{选最懒的那条路走。}

试了不能消，换另一个继续试，直到试出可以消完的未知数为止。

没有固定的顺序，没有规定必须先消哪个。

$x$ 不好消，试试消 $y$。

$y$ 也不好消，试试消 $z$。

总有一个能消的。

\section{三元方程组的三种情况}

就像二元方程组一样，三元方程组也有三种情况：

\textbf{情况一：唯一解。}

像我们刚才解出来的那样，刚刚好有一组确定的值。

\textbf{情况二：无解。}

消元的过程中，出现了矛盾（比如 $5 = 7$）。说明方程之间互相冲突。

\textbf{情况三：无数组解。}

消元的过程中，出现了 $0 = 0$ 这种没有用的废话。这说明某个方程是“多余”的，它只是重复了别人的信息，限制条件不够，答案有无数种可能。

\textbf{第十八章完成 ✓}

三元方程组其实没有新知识，它只是二元方程组的“加强版”，多了一次消元的步骤。

消消消，从三个未知数消到一个，解出来，代回去，验证。

只要有耐心，一步一步算，你一定能解出来。

下一章，我们要把生活里的问题，翻译成这三个方程。

\chapter{三个未知数——更复杂的翻译}

上一章你学了三元方程组怎么解。

但解方程只是计算。

在计算之前，还有一件更重要、更核心的事：

\textbf{把题目翻译成方程。}

如果翻译错了，解得再漂亮也没用。

这一章，我们专门练翻译。

\section{翻译的基本规则}

规则很简单，就两句话：

\textbf{题目问什么，就设什么是未知数。}

\textbf{题目说了什么关系，就把那个关系写成等式。}

\section{第一道题：直接凑数字}

\begin{quote}
\textbf{题目：}

你去超市买水果。

第一次：1斤苹果、1斤香蕉、1斤橙子，花了6块钱。

第二次：2斤苹果、1斤香蕉、1斤橙子，花了7块钱。

第三次：1斤苹果、1斤香蕉、2斤橙子，花了9块钱。

请问苹果、香蕉、橙子每斤各多少钱？
\end{quote}

\subsection{第一步：设未知数}

题目问的是单价。

设苹果每斤 $x$ 元，香蕉每斤 $y$ 元，橙子每斤 $z$ 元。

（$x$ 就是苹果！$y$ 就是香蕉！$z$ 就是橙子！）

\subsection{第二步：翻译成方程}

第一次购物：$x + y + z = 6$ ……①

第二次购物：$2x + y + z = 7$ ……②

第三次购物：$x + y + 2z = 9$ ……③

\subsection{第三步：解方程组}

\[
\begin{cases} x + y + z = 6 & \cdots① \\ 2x + y + z = 7 & \cdots② \\ x + y + 2z = 9 & \cdots③ \end{cases}
\]

这题非常友好。你看：

\textbf{求 $x$：}

用方程②减去方程①：

$(2x + y + z) - (x + y + z) = 7 - 6$

\textbf{$x = 1$}

\textbf{求 $z$：}

用方程③减去方程①：

$(x + y + 2z) - (x + y + z) = 9 - 6$

\textbf{$z = 3$}

\textbf{求 $y$：}

把 $x = 1, z = 3$ 代入方程①：

$1 + y + 3 = 6$

\textbf{$y = 2$}

\subsection{第四步：验证}

方程①：$1 + 2 + 3 = 6$ ✓

方程②：$2(1) + 2 + 3 = 7$ ✓

方程③：$1 + 2 + 2(3) = 9$ ✓

苹果1块，香蕉2块，橙子3块。

\section{第二道题：年龄问题}

\begin{quote}
\textbf{题目：}

三个人，甲、乙、丙。

甲和乙的年龄加起来是40岁。

乙和丙的年龄加起来是35岁。

甲和丙的年龄加起来是39岁。

请问三个人各多少岁？
\end{quote}

\subsection{第一步：设未知数}

设甲的年龄是 $x$ 岁，乙的年龄是 $y$ 岁，丙的年龄是 $z$ 岁。

\subsection{第二步：翻译成方程}

甲和乙加起来40：$x + y = 40$ ……①

乙和丙加起来35：$y + z = 35$ ……②

甲和丙加起来39：$x + z = 39$ ……③

\subsection{第三步：解方程组}

\[
\begin{cases} x + y = 40 & \cdots① \\ y + z = 35 & \cdots② \\ x + z = 39 & \cdots③ \end{cases}
\]

这三个方程每个只有两个未知数，很好消。

\textbf{用方程①减去方程②（消掉 $y$）：}

$(x + y) - (y + z) = 40 - 35$

$x - z = 5 \quad \cdots④$

\textbf{用方程③加上方程④（消掉 $z$）：}

$(x + z) + (x - z) = 39 + 5$

$2x = 44$

\textbf{$x = 22$} （甲是22岁）

\textbf{代入方程①求 $y$：}

$22 + y = 40$

\textbf{$y = 18$} （乙是18岁）

\textbf{代入方程②求 $z$：}

$18 + z = 35$

\textbf{$z = 17$} （丙是17岁）

\subsection{验证}

方程①：$22 + 18 = 40$ ✓

方程②：$18 + 17 = 35$ ✓

方程③：$22 + 17 = 39$ ✓

全部通过。甲22岁，乙18岁，丙17岁。

\section{第三道题：混合问题（注意“比”字）}

\begin{quote}
\textbf{题目：}

学校举办活动，卖了三种票：学生票、成人票、老人票。

学生票5块，成人票10块，老人票3块。

一共卖了100张票，收入640块钱。

已知成人票比学生票多10张。

请问三种票各卖了多少张？
\end{quote}

\subsection{第一步：设未知数}

设学生票卖了 $x$ 张，成人票卖了 $y$ 张，老人票卖了 $z$ 张。

（$x$ 就是学生票的数量！$y$ 就是成人票的数量！$z$ 就是老人票的数量！）

\subsection{第二步：翻译成方程}

\textbf{总张数是100：}

$x + y + z = 100$ ……①

\textbf{总收入是640块：}

数量乘以单价等于收入。

$5x + 10y + 3z = 640$ ……②

\textbf{成人票比学生票多10张：}

这句话怎么翻译？

成人票是 $y$，学生票是 $x$。

谁多？成人票多。

所以：$y = x + 10$。

移项整理成标准格式：$-x + y = 10$ ……③

\subsection{第三步：解方程组}

\[
\begin{cases} x + y + z = 100 & \cdots① \\ 5x + 10y + 3z = 640 & \cdots② \\ y = x + 10 & \cdots③ \end{cases}
\]

既然方程③直接告诉我们 $y$ 等于多少，我们可以直接把 $y$ 代入前两个方程，这就是“代入消元法”。

\textbf{把方程③代入方程①：}

$x + (x + 10) + z = 100$

$2x + z + 10 = 100$

$2x + z = 90 \quad \cdots④$

\textbf{把方程③代入方程②：}

$5x + 10(x + 10) + 3z = 640$

$5x + 10x + 100 + 3z = 640$

$15x + 3z = 540$

两边同除以3化简：

$5x + z = 180 \quad \cdots⑤$

现在有了二元方程组：

\[
\begin{cases} 2x + z = 90 & \cdots④ \\ 5x + z = 180 & \cdots⑤ \end{cases}
\]

\textbf{用方程⑤减去方程④（消掉 $z$）：}

$(5x + z) - (2x + z) = 180 - 90$

$3x = 90$

\textbf{$x = 30$} （学生票30张）

\textbf{代入方程③求 $y$：}

$y = 30 + 10$

\textbf{$y = 40$} （成人票40张）

\textbf{代入方程④求 $z$：}

$2(30) + z = 90$

$60 + z = 90$

\textbf{$z = 30$} （老人票30张）

\subsection{验证}

总张数：$30 + 40 + 30 = 100$ ✓

总收入：$5(30) + 10(40) + 3(30) = 150 + 400 + 90 = 640$ ✓

成人比学生多10张：$40 - 30 = 10$ ✓

学生票30张，成人票40张，老人票30张。

\section{翻译的注意事项}

做完这几道题，我帮你总结一下翻译的时候最容易掉坑的地方：

\textbf{注意一：单位要对齐。}

价格的等式里只能是价格（块），数量的等式里只能是数量（张/斤）。不能把它们混在一个方程里。

\textbf{注意二：“比”字怎么写。}

“A比B多C” —— 翻译成 $A = B + C$。

“A是B的3倍” —— 翻译成 $A = 3B$。

千万别写反了。

\textbf{注意三：有几个未知数，就需要几个方程。}

三个未知数，必须从题目里找出三句独立的话，列出三个方程。少一个都解不出来。

\textbf{第十九章完成 ✓}

你现在已经掌握了三元应用题的完整流程：从生活场景出发，翻译成精确的数学语言（设未知数、列方程），然后通过严格的逻辑步骤（消元）解出答案，最后验证确保无误。

翻译是灵魂，计算是手艺。两者结合，你就能解决越来越复杂的问题。

\chapter{三个未知数——极其复杂的翻译}

上一章我们练了简单的翻译。

苹果香蕉橙子，甲乙丙年龄。

那些题目，关系很直接，一眼就能看出来。

但生活里的问题，往往没这么简单。

很多时候，条件是\textbf{绕着弯}给你的。

这一章，我们来练\textbf{极其复杂的翻译}。

我会放慢速度，把“我是怎么想到这一步的”一点一点掰开揉碎讲给你听。

\section{第一道题：比例和倍数}

\begin{quote}
\textbf{题目：}

甲、乙、丙三人分一笔钱。

甲分到的钱，是乙和丙总和的一半。

乙分到的钱，是甲和丙总和的三分之一。

丙分到了200块。

请问总金额是多少？甲和乙各分到了多少？
\end{quote}

\subsection{慢动作思考}

我们先设未知数。

题目问甲、乙各多少，那就设：

甲分到了 $x$ 块。

乙分到了 $y$ 块。

丙分到了 $z$ 块。（虽然题目给了200，但我们先设个 $z$，让等式看起来更整齐。）

现在看第一句话：

\textbf{“甲分到的钱，是乙和丙总和的一半”}

把这句话拆开看：

1. \textbf{“甲分到的钱”} $\rightarrow x$
2. \textbf{“乙和丙总和”} $\rightarrow y + z$
3. \textbf{“是一半”} $\rightarrow$ 就是把总和除以2，也就是乘以 $\dfrac{1}{2}$

拼起来：

\[
x = (y + z) \times \frac{1}{2}
\]

为了好看点，写成：

\[
x = \frac{1}{2}(y + z) \quad \cdots①
\]

看第二句话：

\textbf{“乙分到的钱，是甲和丙总和的三分之一”}

同样拆解：

1. \textbf{“乙分到的钱”} $\rightarrow y$
2. \textbf{“甲和丙总和”} $\rightarrow x + z$
3. \textbf{“是三分之一”} $\rightarrow \times \dfrac{1}{3}$

拼起来：

\[
y = \frac{1}{3}(x + z) \quad \cdots②
\]

看第三句话：

\textbf{“丙分到了200块”}

直接翻译：

\[
z = 200 \quad \cdots③
\]

现在我们有了三个方程：

\[
\begin{cases} x = \dfrac{1}{2}(y + z) \\ y = \dfrac{1}{3}(x + z) \\ z = 200 \end{cases}
\]

这就翻译完了。后面就是解方程的体力活了。

\section{第二道题：变化和增减}

\begin{quote}
\textbf{题目：}

三个仓库A、B、C，总共存了100吨粮食。

如果我们从A仓库运10吨到B仓库，再从B仓库运5吨到C仓库，那么三个仓库的存粮就一样多了。

请问原来三个仓库各有多少吨？
\end{quote}

\subsection{慢动作思考}

设原来A有 $x$ 吨，B有 $y$ 吨，C有 $z$ 吨。

\textbf{第一句话：“总共存了100吨”}

这很简单：

\[
x + y + z = 100 \quad \cdots①
\]

\textbf{第二句话：“从A运10吨到B”}

这句话是一个\textbf{动作}。发生了什么变化？

A少了10吨。所以A变成了 $x - 10$。

B多了10吨。所以B变成了 $y + 10$。

C没动。还是 $z$。

现在的状态：

- A：$x - 10$
- B：$y + 10$
- C：$z$

\textbf{第三句话：“再从B运5吨到C”}

这是第二个动作。要注意，它是在\textbf{刚才的基础上}变化的。

A这次没动。还是 $x - 10$。

B本来是 $y + 10$，现在又运走了5吨，所以少了5吨：$(y + 10) - 5 = y + 5$。

C本来是 $z$，现在运来了5吨，所以多了5吨：$z + 5$。

现在的最终状态：

- A：$x - 10$
- B：$y + 5$
- C：$z + 5$

\textbf{第四句话：“三个仓库一样多了”}

这意味着：现在的A = 现在的B = 现在的C。

\[
x - 10 = y + 5 = z + 5
\]

这是一个连等式，其实它包含了两个方程（随便挑两对就行）：

1. A和B一样多：$x - 10 = y + 5 \quad \cdots②$
2. B和C一样多：$y + 5 = z + 5 \quad \cdots③$

翻译完成：

\[
\begin{cases} x + y + z = 100 \\ x - 10 = y + 5 \\ y + 5 = z + 5 \end{cases}
\]

\section{第三道题：没给总成本的百分比问题}

\begin{quote}
\textbf{题目：}

一件衣服，成本由布料、人工、运费三部分组成。

去年的情况：布料是人工的2倍。运费是人工的一半。

今年：布料涨价10\%，人工涨价20\%，运费没变。

涨价后，总成本变成了122元。

请问去年的各项成本是多少？
\end{quote}

\subsection{慢动作思考}

题目问去年，那就设去年的成本：

布料 = $x$ 元，人工 = $y$ 元，运费 = $z$ 元。

\textbf{第一句话：“布料是人工的2倍”}

很简单：

\[
x = 2y \quad \cdots①
\]

\textbf{第二句话：“运费是人工的一半”}

也很简单：

\[
z = 0.5y \quad \cdots②
\]

接下来是变化。

\textbf{第三句话：“今年布料涨价10\%”}

“涨价10\%”的意思是，在原来的基础上，多出了10\%。

所以今年的布料价格 = 去年的价格 $\times (1 + 10\%)$ = $1.1x$。

\textbf{第四句话：“人工涨价20\%”}

同理，今年的人工价格 = $y \times (1 + 20\%)$ = $1.2y$。

\textbf{第五句话：“运费没变”}

今年的运费还是 $z$。

\textbf{第六句话：“涨价后，总成本变成了122元”}

什么是总成本？就是三部分加起来。

今年的布料 + 今年的人工 + 今年的运费 = 122

\[
1.1x + 1.2y + z = 122 \quad \cdots③
\]

翻译完成：

\[
\begin{cases} x = 2y \\ z = 0.5y \\ 1.1x + 1.2y + z = 122 \end{cases}
\]

\section{第四道题：逆流顺流}

\begin{quote}
\textbf{题目：}

一艘船在河里航行。

顺流而下，速度是每小时20公里。

逆流而上，速度是每小时10公里。

如果不开发动机，船随水漂流，速度是多少？
\end{quote}

\subsection{慢动作思考}

这道题虽然只问了一个数，但其实有两个不知道的东西：

1. 船本身（发动机）开得有多快——设为 $x$。
2. 水流得有多快——设为 $y$。

\textbf{理解“顺流”：}

顺流就是“船自己在跑，水还在后面推它”。

所以实际速度 = 船速 + 水速。

\[
x + y = 20 \quad \cdots①
\]

\textbf{理解“逆流”：}

逆流就是“船在往前跑，水在前面顶它”。

实际速度 = 船速 - 水速。

\[
x - y = 10 \quad \cdots②
\]

\textbf{理解“随水漂流”：}

随水漂流就是船不开发动机（$x = 0$），只靠水推。

这其实就是在问：水速 $y$ 是多少。

翻译完成：

\[
\begin{cases} x + y = 20 \\ x - y = 10 \end{cases}
\]

题目求 $y$。

\section{第五道题：牛吃草（最难的一道）}

\begin{quote}
\textbf{题目：}

牧场长满草。草每天都在匀速生长。

20头牛，10天吃完。

15头牛，15天吃完。

请问，要让草永远吃不完，最多能养几头牛？
\end{quote}

\subsection{慢动作思考}

这道题为什么难？因为它是一个\textbf{动态}的过程。草在被吃，同时又在长。

我们要设三个未知数，把所有隐藏的量都标记出来：

1. \textbf{牧场原来有多少草} ——设为 $x$。
2. \textbf{草每天长出来多少} ——设为 $y$。
3. \textbf{每头牛每天吃多少} ——设为 $z$。

\textbf{翻译第一句话：“20头牛，10天吃完”}

这句话的意思是：

\textbf{20头牛在10天里吃的总量 = 牧场原有的草 + 这10天里长出来的草}

我们来算每一部分：

\textbf{牛吃的总量：}

20头牛 $\times$ 10天 $\times$ 每头每天吃 $z$ = $200z$。

\textbf{草的总量：}

原有的草是 $x$。

10天长出来的草是：10天 $\times$ 每天长 $y$ = $10y$。

加起来是：$x + 10y$。

\textbf{列出等式：}

\[
x + 10y = 200z \quad \cdots①
\]

\textbf{翻译第二句话：“15头牛，15天吃完”}

同样的逻辑：

\textbf{15头牛在15天里吃的总量 = 牧场原有的草 + 这15天里长出来的草}

\textbf{牛吃的总量：}

15头牛 $\times$ 15天 $\times$ 每头每天吃 $z$ = $225z$。

\textbf{草的总量：}

原有的草是 $x$。

15天长出来的草是：15天 $\times$ 每天长 $y$ = $15y$。

加起来是：$x + 15y$。

\textbf{列出等式：}

\[
x + 15y = 225z \quad \cdots②
\]

现在我们有两个方程了：

\[
\begin{cases} x + 10y = 200z \\ x + 15y = 225z \end{cases}
\]

\textbf{翻译第三句话：“要让草永远吃不完，最多能养几头牛？”}

这句话是什么意思？

“草永远吃不完”意味着：\textbf{牛吃草的速度，不能超过草生长的速度。}

如果吃得比长得快，存货总有一天会吃完。

如果吃得比长得慢（或者刚好一样快），就可以只吃每天新长出来的草，永远吃不完。

设养 $N$ 头牛。

\textbf{牛吃草的速度：} $N$ 头牛 $\times$ 每头每天吃 $z$ = $Nz$。

\textbf{草生长的速度：} 每天长 $y$。

\textbf{条件：} 吃的速度 $\le$ 长的速度。

\[
Nz \le y
\]

我们要算 $N$，所以把 $z$ 除过去：

\[
N \le \frac{y}{z}
\]

我们需要求出 $\dfrac{y}{z}$ 是多少。

这下就变成了纯粹的计算问题：从前面两个方程里，求出 $y$ 和 $z$ 的关系。

方程②减去方程①：

$(x + 15y) - (x + 10y) = 225z - 200z$

$x$ 被消掉了：

\[
5y = 25z
\]

两边除以5：

\[
y = 5z
\]

或者写成：$\dfrac{y}{z} = 5$。

这句话的意思是：\textbf{草每天新长出来的量（$y$），刚好够5头牛一天的饭量（$5z$）。}

回到我们的条件：

$N \le \dfrac{y}{z}$

$N \le 5$

所以，最多只能养 \textbf{5头牛}。

如果养6头，每天新长出来的草不够吃，就要去吃老本，早晚有一天吃完。

\section{翻译的秘诀}

这章的题目难度不低，但思路其实只有三步：

\textbf{第一步：无论多乱，先设未知数。}

不要怕未知数多。先把不知道的、题目里提到的核心东西，都把名字占上（$x, y, z$）。

\textbf{第二步：拆解句子。}

把一句长长的话，拆成“谁”、“怎么变”、“等于谁”。

\textbf{第三步：把拆出来的每一部分，用 $x, y, z$ 写出来，拼成等式。}

\begin{quote}
\textbf{一句话总结：翻译极其复杂的应用题，就是“拆句子”。把复杂的话拆成简单的词，然后一个词一个词地翻译成数学符号。}
\end{quote}

\textbf{第二十章完成 ✓}

你现在连牛吃草这种绕来绕去的题目都能拆解并翻译了。

翻译是解决问题的第一步。只要方程列对了，剩下的就是计算。

下一章，我们来学一种专门处理方程组的强大工具——线性代数。

\chapter{图形的方程——详细认识图形公式}

停一下。

在往下看之前，去拿一支笔，和一张白纸。最好是带格子的草稿纸。

我认真的。去拿。

……

拿好了吗？好。

在这一章，我得对你稍微严格一点，给你定个铁规矩：\textbf{学这一章的时候，你必须得跟着动手画。}

我真不是在故意折腾你，更不是在凶你。我是太了解这些数学公式了——如果你只用眼睛盯着它们看，它们长得就像一堆乱码，越看越头痛。

但是，只要你肯拿起笔，跟着我在纸上画个十字架，点上两个点，连一条线……哪怕画得歪歪扭扭的，那些天书一样的公式，瞬间就会变成大白话。

你会一拍大腿说：“嗨，原来就是这么回事啊！”

所以，答应我，别在这件事上偷懒。相信我，动手画，魔法就会在你的笔尖上发生。

现在，看着你手里的那张白纸。

在纸上随便画一条线，或者画一个圆。

这叫\textbf{几何}。

接着，在图形上面画一个大大的十字架——横着的是 $x$ 轴，竖着的是 $y$ 轴。

这张纸瞬间变成了\textbf{坐标系}。

神奇的事情发生了：

纸上的每一个点，都有了一个名字（坐标），比如 $(2, 3)$。

纸上的每一条线、每一个图形，也都可以用一个\textbf{包含 $x$ 和 $y$ 的算式}来表示。

这个算式，就叫\textbf{图形的方程}。

如果说没有平方的 $x$ 和 $y$ 画出来永远是笔直的线，那么带有\textbf{平方（$x^2$ 或 $y^2$）}的公式，画出来必定是弯曲的线。

这一章，我们就拿着笔，顺着\textbf{坐标相减}的底层逻辑，硬核拆解坐标系里的几何规则！

\section{第一步：量量有多远（两点间距离公式）}

现在，在你的坐标系里点两个点：

点 $A$ 画在 $(1, 2)$ 的位置。

点 $B$ 画在 $(4, 6)$ 的位置。

我想问：\textbf{它们俩之间离得多远？}

你不能直接拿尺子量，我们要用算的。

怎么算？

\textbf{拿起笔，画一个直角三角形。}

从点 $A$ 往右画一条横线（平路），走到点 $B$ 的正下方。

从点 $B$ 往下画一条竖线（爬楼），走到点 $A$ 的正右方。

让这两条线相交，你会看到一个完美的直角。

现在把 $A$ 和 $B$ 连起来，这就是三角形的斜边。

现在看这个三角形的两条直角边：

\textbf{横着走的那条边（平路）：}

从 $x = 1$ 走到了 $x = 4$。

你走了多远？$4 - 1 = 3$。

（数学里管这叫 $\Delta x$，意思是 $x$ 的变化量。）

\textbf{竖着走的那条边（爬楼）：}

从 $y = 2$ 爬到了 $y = 6$。

你爬了多高？$6 - 2 = 4$。

（这叫 $\Delta y$，$y$ 的变化量。）

好，看着你画的图：底边是3，高是4。

斜边（也就是 $A$ 和 $B$ 之间的距离）是多少？

用勾股定理：

$3^2 + 4^2 = c^2$

$9 + 16 = 25 = c^2$

$c = \sqrt{25} = 5$

距离是 5！

如果把具体的数字换成字母：

点 $A$ 是 $(x_1, y_1)$，点 $B$ 是 $(x_2, y_2)$。

横向距离：$x_2 - x_1$

纵向距离：$y_2 - y_1$

根据勾股定理，距离 $d$ 的平方等于：

$d^2 = (x_2 - x_1)^2 + (y_2 - y_1)^2$

两边开根号，得到\textbf{两点间距离公式}：

\[
d = \sqrt{(x_2 - x_1)^2 + (y_2 - y_1)^2}
\]

\textbf{语言记忆法：}

\begin{quote}
\textbf{距离 = $\sqrt{\text{横着走的平方} + \text{竖着爬的平方}}$。它就是勾股定理。}
\end{quote}

你看，只要你画出了那个三角形，这个吓人的公式是不是瞬间就看懂了？

\section{第二步：找中间（中点公式）}

还是点 $A(1, 2)$ 和点 $B(4, 6)$。

我想找它们正中间的那个点，坐标是多少？

这个最简单。

\textbf{中间，就是平均数。}

横坐标的中间：1和4的平均数。

$\dfrac{1 + 4}{2} = 2.5$

纵坐标的中间：2和6的平均数。

$\dfrac{2 + 6}{2} = 4$

所以中点坐标是 $(2.5, 4)$。

写成公式：

\[
\text{中点} = \left( \frac{x_1 + x_2}{2}, \frac{y_1 + y_2}{2} \right)
\]

\textbf{语言记忆法：}

\begin{quote}
\textbf{横加横除以2，纵加纵除以2。}
\end{quote}

\section{第三步：斜率的纯数字本质}

我们之前讲过斜率。现在，我们把那个看起来很高大上的希腊字母“$\Delta$”直接扔进垃圾桶！我们要看最纯粹的东西：\textbf{坐标的相减}。

我们在一条直线上挑两个点：

*   \textbf{起点}叫 $A(x_1, y_1)$。
*   \textbf{目标点}叫 $B(x_2, y_2)$。

你想从 A 走到 B，你需要两步：

1.  \textbf{竖着走多远？} 拿目标高度减去起点高度，也就是 \textbf{$y_2 - y_1$}。
2.  \textbf{横着走多远？} 拿目标位置减去起点位置，也就是 \textbf{$x_2 - x_1$}。

斜率 $k$ 的终极公式就是：

\[
k = \frac{y_2 - y_1}{x_2 - x_1}
\]

（直白地说：竖着走的距离 $\div$ 横着走的距离）

\subsection{一个极其重要的问题：能不能用 $y_1 - y_2$？}

\textbf{答案是：绝对可以！完全合法！}

只要你上下保持同步，谁减谁都一样：

\[
k = \frac{y_1 - y_2}{x_1 - x_2} \quad \text{和} \quad \frac{y_2 - y_1}{x_2 - x_1} \quad \text{算出来是一模一样的结果！}
\]

\textbf{为什么？}

因为如果你把 A 当起点 B 当终点，那是“向上向右走”（正除以正）；如果你把 B 当起点 A 当终点，那就是“向下向左走”（负除以负），负负得正，斜率 $k$ 依然不变！

\subsection{这个得数到底意味着什么？}

当你拿两个点的坐标硬碰硬相减，做完除法后，你会得到一个数字 $k$。

\textbf{【得数含义】：这个数字是直线在网格里“强制执行的步伐规则”。}

*   \textbf{算出来 $k = 3$：} 意思是，只要你在这条线上，你往右边每挪动 \textbf{1} 格，你的位置会被\textbf{强制向上抬升 3 格}。非常陡峭的爬坡。
*   \textbf{算出来 $k = -0.5$：} 意思是，你往右边走 \textbf{1} 格，你会\textbf{往下掉 0.5 格}。这是一个平缓的下坡。
*   \textbf{算出来 $k = 0$：} 分子 $y_2 - y_1$ 是 0，意味着高度没变化。不管你往右走多少格，你\textbf{向上/向下抬升的格数永远是 0}。这就是平地（水平线）。
*   \textbf{算出来报错（分母 $x_2 - x_1 = 0$）：} 意思是，你根本\textbf{没有往右走}，但高度却发生了变化。在现实网格里，这意味着你面前是一堵\textbf{垂直的墙}。

\section{第四步：用相减直接“变出”所有直线方程}

刚才说了，直线的特点是：\textbf{随便你挑哪两个点，算出来的斜率 $k$ 必须永远一样。}

现在，我们把 $A(x_1, y_1)$ 当作一个\textbf{死死钉在坐标系上的固定点}。

然后在直线上放一个\textbf{随便乱跑的动点}，我们不管它具体是几，就叫它 $(x, y)$。

既然斜率永远不变，那动点和固定点相减，必然还是等于 $k$：

\[
k = \frac{y - y_1}{x - x_1}
\]

现在，见证奇迹：我们不要让它保持分数的样子，把分母 $(x - x_1)$ 乘到等号左边去！

\[
y - y_1 = k(x - x_1)
\]

你看！\textbf{这就是高中数学里天天背的“点斜式方程”！}

它根本不是什么新发明，它就是把 $k = \dfrac{y - y_1}{x - x_1}$ 的分母挪了个位置而已！

如果你把它化简，比如 $y - 3 = 4(x - 2)$，变成 $y = 4x - 5$。这就是我们熟悉的 $y = kx + b$。

\textbf{在这里，两个得数的含义：}

*   \textbf{单独看 $b = -5$：} 直线的“起跑线”在哪。当 $x = 0$ 时，这条线刚好精准地砸在 y 轴的数字 \textbf{-5} 上。
*   \textbf{结合看 $y = 4x - 5$：} 这是一个“计费系统”。起步价欠5块（$b=-5$），每跑1公里，计价器往上跳4块（$k=4$）。

\section{第五步：为什么垂直线斜率相乘等于 -1？}

这次不用比喻了，我们直接用 $(x, y)$ 的变化来硬刚这个结论！

假设我们在坐标原点 $(0,0)$ 画一条线 $L_1$。

在 $L_1$ 上随便找一个点，坐标叫 \textbf{$(a, b)$}。

那么这条线 $L_1$ 的斜率就是：

\[
k_1 = \frac{b - 0}{a - 0} = \frac{b}{a}
\]

\textbf{重点来了：怎么画出跟它绝对垂直的 $L_2$？}

把刚才那个点 $(a, b)$，绕着原点逆时针\textbf{死死地旋转 90 度}！

在坐标系里，一个点 $(a, b)$ 旋转 90 度后，它的坐标会变成什么？

\textbf{它的横竖会互换，并且横坐标变成负的，新坐标变成了 $(-b, a)$！}

（在你的纸上画个图验证一下：点 $(3, 2)$ 转 90 度，绝对会跑到 $(-2, 3)$ 的位置）

现在，我们算垂直线 $L_2$ 的斜率。它的起点也是 $(0,0)$，终点是 $(-b, a)$：

\[
k_2 = \frac{a - 0}{-b - 0} = -\frac{a}{b}
\]

现在把两个斜率乘起来看看：

\[
k_1 \times k_2 = \left(\frac{b}{a}\right) \times \left(-\frac{a}{b}\right) = -1
\]

\textbf{破案了！}

不需要任何复杂的几何定理，仅仅因为\textbf{“旋转90度会让 x 和 y 互换位置并改变一个符号”}，就注定了垂直线的斜率一定是\textbf{互为负倒数}。

乘积 -1 的绝对物理含义就是：\textbf{完美的 90 度十字交叉！}

\section{第六步：斜率与角度的硬连接（$\tan \theta$）}

一条直线在网格里，和水平的 x 轴会形成一个夹角 $\theta$。

如果你从 $(x_1, y_1)$ 走到 $(x_2, y_2)$，你实际上在坐标系里画出了一个\textbf{直角三角形}：

*   三角形的“垂直高”就是 \textbf{$y_2 - y_1$}
*   三角形的“水平底”就是 \textbf{$x_2 - x_1$}

在三角函数里，正切值 $\tan \theta$ 的定义就是 $\dfrac{\text{对边}}{\text{邻边}}$，也就是 $\dfrac{\text{高}}{\text{底}}$。

所以：

\[
k = \tan \theta = \frac{y_2 - y_1}{x_2 - x_1}
\]

*   算出来 $k = 1$，意味着 $\tan \theta = 1$，真实倾角是 \textbf{45 度}。
*   算出来 $k = \sqrt{3} \approx 1.732$，意味着真实倾角是 \textbf{60 度}。
*   算出来 $k = 1000$，意味着夹角接近 89.99度，它几乎是一根垂直的柱子了，但\textbf{仍然没有完全垂直}。

\section{第七步：从直线走向微积分}

如果你要算一条弯曲的抛物线 $y = x^2$ 上的斜率，怎么办？

弯曲的线上，斜率是时时刻刻都在变的。

怎么办？还是用 $y_2 - y_1$ 和 $x_2 - x_1$！

假设你想知道在点 $A(2, 4)$ 这一瞬间的斜率。

我们在曲线上找一个离它\textbf{极度接近}的点，比如 $B(2.001, 4.004001)$。

用纯坐标相减算斜率：

\[
k = \frac{y_2 - y_1}{x_2 - x_1} = \frac{4.004001 - 4}{2.001 - 2} = \frac{0.004001}{0.001} = 4.001
\]

看到没有？算出来的值非常接近 \textbf{4}。

当我们让 $x_2$ 和 $y_2$ 无限逼近 $x_1$ 和 $y_1$ 时，底下的 $(x_2 - x_1)$ 接近 0，上面的 $(y_2 - y_1)$ 也接近 0。

在微积分里，我们把这种\textbf{“极限状态下极小极小的相减”}写成 $dy$ 和 $dx$：

\[
\text{导数} = \frac{dy}{dx}
\]

它本质上依然是你最开始说的那个老老实实的：$\dfrac{y_2 - y_1}{x_2 - x_1}$。

算出来的导数得数（比如在 $x=2$ 瞬间 $k=4$），就是那一刻它向上的\textbf{“爆发力”数值}——如果它突然不受引力控制沿着切线飞出去，就是一条斜率为4的射线。

一切的一切，都是从\textbf{两个坐标相减}开始的！

\section{第八步：坐标系里的四大形状（VIP 入场规则）}

现在，我们要看那些带有\textbf{平方}的公式了。

不要把形状公式看成死记硬背的字母，把它们看成\textbf{“坐标系网格里的 VIP 俱乐部准入规则”}。

网格上无数个点 $(x, y)$，谁计算后的【得数】能符合规则，谁就有资格站在这条形状的边界线上！

\subsection{一、 圆形（Circle）：绝对公平的距离规则}

在网格里随便定一个死死钉住的中心点 \textbf{$(x_1, y_1)$}。

你想在周围画一个圆，圆的本质规则是：\textbf{所有俱乐部成员点 $(x, y)$，到中心点的直线距离，必须绝对等于一个死数字 $r$。}

怎么算距离？还是用到我们在网格里最信任的直角三角形（勾股定理）：

横向走多远？目标减中心 $\rightarrow$ \textbf{$(x - x_1)$}

竖向走多远？目标减中心 $\rightarrow$ \textbf{$(y - y_1)$}

直角三角形的：底边² + 高² = 斜边²。

\textbf{圆形终极公式： $(x - x_1)^2 + (y - y_1)^2 = r^2$}

\textbf{【得数的物理与几何含义】：}

*   \textbf{等号右边的得数 $r^2$（比如等于 25）：}
    *   \textbf{不可逾越的“绝对结界”。}
    *   它规定了这个圆的半径必须是 $\sqrt{25} = 5$。
    *   随便拿网格上一个点代入左边：算出来 \textbf{$= 25$}，刚好踩在边框上；\textbf{$< 25$}，在内部；\textbf{$> 25$}，在外部。
*   \textbf{减法 $(x - x_1)$ 算出来的得数：}
    *   它单纯表示 x 坐标偏离了中心点多少格。因为外面套了平方，往左偏还是往右偏，平方后得数都是正的。\textbf{圆形之所以完美对称，就是因为平方把正负号全吃掉了！}

\subsection{二、 椭圆（Ellipse）：被强制分配“预算”的圆}

如果你把圆形的公式稍微改一下，让 x 和 y 不再“平起平坐”，圆就被压扁了。

\textbf{椭圆终极公式（以原点为中心）： $\dfrac{x^2}{a^2} + \dfrac{y^2}{b^2} = 1$}

你看，它就是把 $x^2$ 和 $y^2$ 分别除以了两个不同的数字（$a^2$ 和 $b^2$），然后强制要求它们加起来的【得数必须等于 1】。

\textbf{【得数的物理与几何含义】：}

*   \textbf{等号右边得数“1”：}
    *   \textbf{100\% 的“总预算进度条”。}
    *   如果你把 x 走到极限大，让 $\dfrac{x^2}{a^2}$ 刚好等于 1（花光了100\%预算），那么 $\dfrac{y^2}{b^2}$ 必须等于 0，这意味着 y 连半步都不能动！x 和 y 互相牵制，导致这条线画到最后，\textbf{必须闭合}死死扣在一起，形成一个圈。
*   \textbf{分母 $a^2$ 和 $b^2$：}
    *   \textbf{各方向的“最大步数限制”。}
    *   如果 $a=5$，$b=3$。意思是你在横向最多只能走 5 格，在纵向最多只能走 3 格。这就变成了一个“长 10、宽 6”的扁圆！

\subsection{三、 抛物线（Parabola）：只有一方疯狂加速}

如果网格里\textbf{只有一个字母有平方}，另一个没有呢？那就彻底拉不住了！

\textbf{抛物线终极公式（顶点式）： $y - y_1 = a(x - x_1)^2$}

注意看！这和我们讲直线的“点斜式 $y - y_1 = k(x - x_1)$”长得极度相似！

唯一的区别是：抛物线的右边多了个 \textbf{平方$^2$}。

\textbf{【得数的物理与几何含义】：}

*   \textbf{固定点 $(x_1, y_1)$：}
    *   \textbf{“反转命运的最低点/最高点”（顶点）。}
    *   因为 $(x - x_1)^2$ 算出来永远 $\ge 0$。所以当 $x$ 刚好等于 $x_1$ 时，$y$ 就等于 $y_1$。这是整个 U 型的最底端。
*   \textbf{乘数 $a$：}
    *   \textbf{抛物线的“暴力放大器（开口陡峭度）”。}
    *   如果 $a=1$：x 离开顶点走 3 格，y 上升 $3^2=9$ 格。
    *   如果 $a=5$：x 走 3 格，y 会瞬间上升 $5 \times 9 = 45$ 格！\textbf{$a$ 的得数越大，抛物线被拉得越窄、飞得越陡！}

\subsection{四、 双曲线（Hyperbola）：无限排斥的两根刺客}

回到椭圆的公式：$\dfrac{x^2}{a^2} + \dfrac{y^2}{b^2} = 1$。
但如果你\textbf{把中间的加号，变成减号}呢？

\textbf{双曲线终极公式： $\dfrac{x^2}{a^2} - \dfrac{y^2}{b^2} = 1$}

\textbf{【得数的物理与几何含义】：}

*   \textbf{得数“1”在减法下的恐怖含义：}
    *   \textbf{“无限膨胀的平衡”。}
    *   相减等于 1，意味着如果 x 变得极其巨大（算出来是 1000000）。为了让减完还得 1，y 必须也变得极其巨大（变成 999999）！
    *   x 和 y \textbf{再也不可能闭合成圈了}。画出了两条\textbf{朝着无限远狂奔、永不回头的开放曲线}。
*   \textbf{减法导致的一个“死亡盲区”：}
    *   如果你想让 $x = 0$（去网格最中间），代入公式变成： $-\dfrac{y^2}{b^2} = 1$。
    *   左边是负数，右边是正数 1。这在实数网格里\textbf{绝对不可能成立}！
    *   这意味着：双曲线中间有一个巨大的\textbf{“绝对虚无地带”}！曲线被强行撕裂开，无论如何也连不到一起。这就是那个“$-1$”制造的物理鸿沟！

\section{形状规律终极总结}

把方程扔在网格里，不用画图，看平方的组合，直接出结论：

1.  \textbf{没有平方（$x, y$ 都是一次）} $\rightarrow$ 永远是死直的线（直线）。
2.  \textbf{只有一个人有平方} $\rightarrow$ 一方平稳，一方狂飙 $\rightarrow$ \textbf{抛开的 U 型（抛物线）}。
3.  \textbf{两个都有平方，且中间是加号} $\rightarrow$ 互相妥协，预算受限 $\rightarrow$ \textbf{死死闭合的圈（圆/椭圆）}。
4.  \textbf{两个都有平方，且中间是减号} $\rightarrow$ 无限膨胀，中间撕裂 $\rightarrow$ \textbf{背靠背无限跑的线（双曲线）}。

坐标系里的所有形状，本质上就是这几个平方项在互相博弈的结果！

有了这个底层视角，不管图形再怎么变，你都看透它的底牌了。

\chapter{更加到底的翻译——三未知}

\section{先停一下，喘口气}

你已经走到这里了。

三元方程组。三个未知数，三条方程，三条线索。

说实话，很多人在这一章直接放弃了。不是因为他们笨——而是因为没有人好好跟他们说：\textbf{这东西，其实你已经会了。}

你会翻译应用题。✓

你会列方程。✓

你会用消元法解二元方程。✓

三元，不过是多了一只鸡、多了一个盒子、多了一条线索而已。

咱们慢慢来。

\section{先把武器磨一磨}

在动真正复杂的题目之前，我们先把思路理清楚。

三元方程组长什么样？就是这样：

\[
\begin{cases} \text{方程一} \\ \text{方程二} \\ \text{方程三} \end{cases}
\]

三条方程，三个未知数：$x$、$y$、$z$。

解法的核心思想就一句话——

\begin{quote}
\textbf{把三个未知数，先变成两个，再变成一个。}
\end{quote}

用大白话说就是：\textbf{先消一个人，再消一个人，最后剩一个人，逼他招供。}

听起来像审讯？对，就是审讯。

\section{热身题：三种水果的价格}

\begin{quote}
苹果、香蕉、橙子，三种水果。

小李买了：1个苹果、2根香蕉、1个橙子，花了7块钱。

小王买了：2个苹果、1根香蕉、1个橙子，花了8块钱。

小张买了：1个苹果、1根香蕉、2个橙子，花了9块钱。

问：每种水果各多少钱？
\end{quote}

\subsection{第一步：翻译}

这一步，我们不写任何方程，只做一件事：\textbf{把每一句话变成一句数学话。}

先设——

\begin{quote}
$x$ = 一个苹果的价格（单位：元）

$y$ = 一根香蕉的价格（单位：元）

$z$ = 一个橙子的价格（单位：元）
\end{quote}

注意，$x$ 就是苹果的钱，$y$ 就是香蕉的钱，$z$ 就是橙子的钱。不是抽象的字母，是真实的钱。

现在，逐句翻译——

\textbf{“小李买了1个苹果、2根香蕉、1个橙子，花了7块钱。”}

翻译成数学话：

\[
1 \text{ 个苹果的钱} + 2 \text{ 根香蕉的钱} + 1 \text{ 个橙子的钱} = 7
\]

\[
x + 2y + z = 7 \quad \cdots \text{方程①}
\]

\textbf{“小王买了2个苹果、1根香蕉、1个橙子，花了8块钱。”}

\[
2x + y + z = 8 \quad \cdots \text{方程②}
\]

\textbf{“小张买了1个苹果、1根香蕉、2个橙子，花了9块钱。”}

\[
x + y + 2z = 9 \quad \cdots \text{方程③}
\]

翻译完毕。你看，三句话，三条方程，一对一，没有任何魔法。

\subsection{第二步：开始消人}

现在我们有三个嫌疑人：$x$、$y$、$z$。

\textbf{目标：先把 $z$ 这个人从前两条方程里消掉。}

为什么先消 $z$？因为①和②里，$z$ 的系数都是1，消起来最省力气。（数学就是要挑软柿子捏。）

\textbf{用方程① 减去 方程②：}

\[
( x + 2y + z ) - ( 2x + y + z ) = 7 - 8
\]

左边，项对项地减——

- $x - 2x = -x$
- $2y - y = y$
- $z - z = 0$（$z$ 消失了！）

右边：$7 - 8 = -1$

所以得到：

\[
-x + y = -1 \quad \cdots \text{方程④}
\]

很好，$z$ 不见了，现在方程④里只有 $x$ 和 $y$ 两个人。

\textbf{再用方程① 减去 方程③，继续消 $z$：}

①是：$x + 2y + z = 7$

③是：$x + y + 2z = 9$

$z$ 的系数：①里是1，③里是2，不一样。没法直接相减。

我们要让它们一样——把①乘以2：

\[
2 \times (x + 2y + z) = 2 \times 7
\]

\[
2x + 4y + 2z = 14 \quad \cdots \text{方程①'}
\]

现在用 ①' 减去 ③：

\[
(2x + 4y + 2z) - (x + y + 2z) = 14 - 9
\]

左边——

- $2x - x = x$
- $4y - y = 3y$
- $2z - 2z = 0$（$z$ 消失了！）

右边：$14 - 9 = 5$

所以得到：

\[
x + 3y = 5 \quad \cdots \text{方程⑤}
\]

\subsection{第三步：解二元方程组}

现在我们只剩两个人了：$x$ 和 $y$。

\[
\begin{cases} -x + y = -1 & \cdots④ \\ x + 3y = 5 & \cdots⑤ \end{cases}
\]

把④加上⑤（把两个方程直接相加）：

\[
(-x + y) + (x + 3y) = -1 + 5
\]

\[
-x + x + y + 3y = 4
\]

\[
4y = 4
\]

\[
y = 1
\]

$y$ 就是香蕉的价格。香蕉，1块钱一根。

把 $y = 1$ 代入方程④：

\[
-x + 1 = -1
\]

\[
-x = -2
\]

\[
x = 2
\]

$x$ 就是苹果的价格。苹果，2块钱一个。

\subsection{第四步：把最后一个人也逼出来}

把 $x = 2$，$y = 1$ 代入方程①：

\[
2 + 2 \times 1 + z = 7
\]

\[
2 + 2 + z = 7
\]

\[
4 + z = 7
\]

\[
z = 3
\]

$z$ 就是橙子的价格。橙子，3块钱一个。

\subsection{第五步：验算（这一步不能省）}

把 $x=2$，$y=1$，$z=3$ 代回三条原始方程，逐一检查——

\textbf{验证①：} $x + 2y + z = 2 + 2 + 3 = 7$ ✓

\textbf{验证②：} $2x + y + z = 4 + 1 + 3 = 8$ ✓

\textbf{验证③：} $x + y + 2z = 2 + 1 + 6 = 9$ ✓

三条全对。答案是可靠的。

\begin{quote}
苹果 \textbf{2元}，香蕉 \textbf{1元}，橙子 \textbf{3元}。
\end{quote}

\section{正式进入：极其复杂的翻译}

好，热身结束。

接下来的题目，会比刚才难很多。难不在计算——计算步骤是一样的。难在\textbf{题目说的话，藏着很深的意思}，需要你一句一句拆开，翻译干净。

这就是这一章真正要练的东西：\textbf{把复杂的中文，变成干净的方程。}

\section{复杂题一：工程队的工钱}

\begin{quote}
一个工程队，有师傅、普通工人、学徒三种人。

规定：\textbf{2个师傅}的日工资，等于\textbf{3个普通工人}的日工资。

规定：\textbf{3个普通工人}的日工资，等于\textbf{5个学徒}的日工资。

现在，1个师傅、2个普通工人、3个学徒，一天共挣了265元。

问：师傅、普通工人、学徒，各自的日工资是多少？
\end{quote}

\subsection{翻译阶段：把每一句话摆在桌子上}

\textbf{先设好实物——}

\begin{quote}
$x$ = 一个师傅的日工资（元）

$y$ = 一个普通工人的日工资（元）

$z$ = 一个学徒的日工资（元）
\end{quote}

$x$、$y$、$z$ 都是钱，真实的钱，不是抽象的字母。

\textbf{翻译第一句话：}

\begin{quote}
“2个师傅的日工资，等于3个普通工人的日工资。”
\end{quote}

意思是：2个师傅挣的，和3个普通工人挣的，一样多。

\[
2x = 3y \quad \cdots \text{方程①}
\]

\textbf{翻译第二句话：}

\begin{quote}
“3个普通工人的日工资，等于5个学徒的日工资。”
\end{quote}

\[
3y = 5z \quad \cdots \text{方程②}
\]

\textbf{翻译第三句话：}

\begin{quote}
“1个师傅、2个普通工人、3个学徒，一天共挣了265元。”
\end{quote}

\[
x + 2y + 3z = 265 \quad \cdots \text{方程③}
\]

三条方程，翻译完毕。

\subsection{注意看这道题的特殊结构}

方程①和方程②有点奇怪——里面没有“等于某个数”，只有变量之间的等量关系。

这没有问题。这叫\textbf{比例关系}。

我们可以用它们来做\textbf{换算}——把师傅的工资和学徒的工资，都用普通工人的工资来表示。

\textbf{从方程①，用 $y$ 表示 $x$：}

\[
2x = 3y
\]

两边除以2：

\[
x = \frac{3}{2}y
\]

也就是说：一个师傅的工资，等于一个半普通工人的工资。

\textbf{从方程②，用 $y$ 表示 $z$：}

\[
3y = 5z
\]

两边除以5：

\[
z = \frac{3}{5}y
\]

也就是说：一个学徒的工资，等于五分之三个普通工人的工资。

\subsection{代入方程③，解出 $y$}

把 $x = \dfrac{3}{2}y$，$z = \dfrac{3}{5}y$ 代入方程③：

\[
\frac{3}{2}y + 2y + 3 \times \left(\frac{3}{5}y\right) = 265
\]

我们把每一项都通分，分母统一变成10——

$\dfrac{3}{2}y = \dfrac{15}{10}y$

$2y = \dfrac{20}{10}y$

$3 \times \dfrac{3}{5}y = \dfrac{9}{5}y = \dfrac{18}{10}y$

所以：

\[
\frac{15}{10}y + \frac{20}{10}y + \frac{18}{10}y = 265
\]

\[
\frac{53}{10}y = 265
\]

两边乘以10：

\[
53y = 2650
\]

两边除以53：

\[
y = 50
\]

普通工人日工资50元。

然后——

\[
x = \frac{3}{2} \times 50 = 75
\]

师傅日工资75元。

\[
z = \frac{3}{5} \times 50 = 30
\]

学徒日工资30元。

\subsection{验算}

\textbf{验证比例关系①：} $2x = 2 \times 75 = 150$，$3y = 3 \times 50 = 150$。相等。✓

\textbf{验证比例关系②：} $3y = 3 \times 50 = 150$，$5z = 5 \times 30 = 150$。相等。✓

\textbf{验证总工资③：} $75 + 2 \times 50 + 3 \times 30 = 75 + 100 + 90 = 265$。✓

\begin{quote}
师傅日工资 \textbf{75元}，普通工人 \textbf{50元}，学徒 \textbf{30元}。
\end{quote}

\section{复杂题二：混合溶液问题}

这类题，很多人第一次看到，会完全不知道从何下手。

\begin{quote}
有三种浓度的盐水：

A桶：浓度10\%（每100克里有10克盐）

B桶：浓度30\%

C桶：浓度50\%

现在，从三桶里各取一些，混合在一起，得到——

混合后盐水总重 \textbf{100克}，其中盐的重量是 \textbf{34克}。

同时已知：从B桶取的量，是从A桶取的量的\textbf{2倍}。

问：从每桶各取了多少克？
\end{quote}

\subsection{翻译阶段：这道题真正在说什么}

这道题看起来复杂，但本质上只有三件事情在说：

\textbf{第一件事（重量）：} 三桶加在一起，总共100克。

\textbf{第二件事（盐的重量）：} 三桶里的盐，加在一起，总共34克。

\textbf{第三件事（数量关系）：} B桶取的量，是A桶的2倍。

就这三件事。现在我们翻译它们。

\textbf{先设好实物——}

\begin{quote}
$x$ = 从A桶取的克数

$y$ = 从B桶取的克数

$z$ = 从C桶取的克数
\end{quote}

$x$、$y$、$z$ 都是重量，单位是克。

\textbf{翻译第一件事：总重量100克。}

\[
x + y + z = 100 \quad \cdots \text{方程①}
\]

\textbf{翻译第二件事：盐共34克。}

这一句需要多想半秒钟。

从A桶取了 $x$ 克盐水，A桶浓度是10\%，所以——

\begin{quote}
A桶里含的盐 $= x \times 10\% = 0.1x$ 克
\end{quote}

从B桶取了 $y$ 克盐水，B桶浓度是30\%，所以——

\begin{quote}
B桶里含的盐 $= y \times 30\% = 0.3y$ 克
\end{quote}

从C桶取了 $z$ 克盐水，C桶浓度是50\%，所以——

\begin{quote}
C桶里含的盐 $= z \times 50\% = 0.5z$ 克
\end{quote}

三桶盐加在一起，等于34克：

\[
0.1x + 0.3y + 0.5z = 34 \quad \cdots \text{方程②}
\]

把②两边乘以10，把小数去掉：

\[
x + 3y + 5z = 340 \quad \cdots \text{方程②'}
\]

\textbf{翻译第三件事：B桶的量是A桶的2倍。}

\[
y = 2x \quad \cdots \text{方程③}
\]

\subsection{开始解方程}

已知 $y = 2x$，我们直接用这个条件，把①和②'里的 $y$ 换掉。

\textbf{代入方程①：}

\[
x + 2x + z = 100
\]

\[
3x + z = 100 \quad \cdots \text{方程④}
\]

\textbf{代入方程②'：}

\[
x + 3(2x) + 5z = 340
\]

\[
x + 6x + 5z = 340
\]

\[
7x + 5z = 340 \quad \cdots \text{方程⑤}
\]

现在只剩两个未知数 $x$ 和 $z$ 了。

从方程④，用 $x$ 表示 $z$：

\[
z = 100 - 3x
\]

代入方程⑤：

\[
7x + 5(100 - 3x) = 340
\]

\[
7x + 500 - 15x = 340
\]

\[
-8x = 340 - 500
\]

\[
-8x = -160
\]

\[
x = 20
\]

\textbf{$y = 2x = 40$}

\textbf{$z = 100 - 3(20) = 40$}

\subsection{验算}

总重量：$20 + 40 + 40 = 100$。✓

盐的重量：$0.1 \times 20 + 0.3 \times 40 + 0.5 \times 40 = 2 + 12 + 20 = 34$。✓

B是A的2倍：$40 = 2 \times 20$。✓

\begin{quote}
A桶 \textbf{20克}，B桶 \textbf{40克}，C桶 \textbf{40克}。
\end{quote}

\section{复杂题三：这才是最难翻译的}

\begin{quote}
三位同学：甲、乙、丙，共同完成一批零件。

已知：\textbf{甲的效率是乙的1.5倍。}

已知：\textbf{乙和丙合作，完成这批零件需要6天。}

已知：\textbf{甲和丙合作，完成这批零件需要4天。}

问：如果三人合作，需要几天完成？
\end{quote}

\subsection{先停下来，想清楚这道题在说什么}

这是一道\textbf{工程问题}。

工程问题，最容易踩的坑，就是不知道设什么为未知数。

初学者的直觉是：设天数为 $x$。

但你会发现，“甲的效率是乙的1.5倍”这句话，用天数很难翻译。

工程问题有一个核心武器——

\begin{quote}
\textbf{设效率为未知数，不设天数。}
\end{quote}

效率，就是每天完成的工作量占总工作量的比例。

\textbf{设定实物——}

\begin{quote}
把这批零件的总量，看成 \textbf{1}（整体）。

$x$ = 甲每天完成的工作量（占总量的比例）

$y$ = 乙每天完成的工作量

$z$ = 丙每天完成的工作量
\end{quote}

$x$、$y$、$z$ 都是每天完成的比例，是真实的效率，不是抽象的符号。

\textbf{翻译第一句话：}

\begin{quote}
“甲的效率是乙的1.5倍。”
\end{quote}

直接翻译：

\[
x = 1.5y \quad \cdots \text{方程①}
\]

\textbf{翻译第二句话：}

\begin{quote}
“乙和丙合作，完成这批零件需要6天。”
\end{quote}

乙每天完成 $y$，丙每天完成 $z$，两人合作每天完成 $y + z$。

6天完成全部（总量为1）：

\[
(y + z) \times 6 = 1
\]

\[
y + z = \frac{1}{6} \quad \cdots \text{方程②}
\]

用大白话说：乙和丙加在一起，每天完成六分之一。

\textbf{翻译第三句话：}

\begin{quote}
“甲和丙合作，完成这批零件需要4天。”
\end{quote}

\[
(x + z) \times 4 = 1
\]

\[
x + z = \frac{1}{4} \quad \cdots \text{方程③}
\]

三条方程，翻译完毕。

\subsection{解方程}

把方程①代入方程③，用 $y$ 替换 $x$：

\[
1.5y + z = \frac{1}{4} \quad \cdots \text{方程③'}
\]

现在和方程②放在一起：

\[
\begin{cases} y + z = \dfrac{1}{6} & \cdots② \\ 1.5y + z = \dfrac{1}{4} & \cdots③' \end{cases}
\]

用③'减去②：

\[
(1.5y + z) - (y + z) = \frac{1}{4} - \frac{1}{6}
\]

\[
0.5y = \frac{3}{12} - \frac{2}{12} = \frac{1}{12}
\]

\[
y = \frac{1}{6}
\]

乙每天完成六分之一。这意味着乙单独做，需要6天。

把 $y = \dfrac{1}{6}$ 代入方程②：

\[
\frac{1}{6} + z = \frac{1}{6}
\]

\[
z = 0
\]

$z$ 等于0，意思是丙的效率为0——丙每天什么都没做。

这个结果成立吗？让我们回去检查题目——

\begin{quote}
乙和丙合作，6天完成。

甲和丙合作，4天完成。

甲的效率是乙的1.5倍。
\end{quote}

如果 $z = 0$，那么——

- 乙和丙合作等于只有乙工作，6天完成，说明乙单独需要6天。
- 甲和丙合作等于只有甲工作，4天完成，说明甲单独需要4天。
- 甲的效率是乙的1.5倍：甲每天完成 $\dfrac{1}{4}$，乙每天完成 $\dfrac{1}{6}$，比值是 $\dfrac{1/4}{1/6} = \dfrac{6}{4} = 1.5$。✓

所以这道题里，丙的效率确实是0。题目数据如此设定，结果虽然意外，但数学上完全自洽。

\textbf{三人合作，每天完成的工作量：}

\[
x + y + z = \frac{1}{4} + \frac{1}{6} + 0 = \frac{3}{12} + \frac{2}{12} = \frac{5}{12}
\]

完成全部工作所需天数：

\[
\text{天数} = 1 \div \frac{5}{12} = \frac{12}{5} = 2.4 \text{ 天}
\]

\begin{quote}
三人合作，需要 \textbf{2.4天}完成这批零件。
\end{quote}

\section{翻译的本质，就是这几件事}

经过这几道题，你应该感觉到了：

\begin{quote}
三元方程组的题目，再复杂，也不过是在说几件事情。你的任务，是把这几件事情，一件一件翻译干净。
\end{quote}

翻译的核心步骤，就三步——

\textbf{第一步：认清楚每个未知数是什么。}

不是设 $x$，而是设“$x$ 是苹果的价格”、“$x$ 是甲的效率”。
每个字母背后，都有一个真实的东西。

\textbf{第二步：把每一句话，独立翻译成一条方程。}

一句话，一条方程。不要把几句话混在一起想。

\textbf{第三步：检查方程的数量和未知数的数量是否匹配。}

三个未知数，需要三条独立的方程。
少了，解不出来。多了，说明你设重了。

\section{本章小结}

\begin{center}
\begin{tabular}{ll}
\toprule
题型 & 核心翻译技巧 \\
\midrule
买卖问题 & 件数 $\times$ 单价 = 总价；每句话是一条方程 \\
工资比例问题 & 比例关系直接写成等式；换算统一变量 \\
溶液混合问题 & 浓度 $\times$ 重量 = 溶质重量；总量单独一条 \\
工程问题 & 设效率，不设天数；效率 $\times$ 天数 = 1 \\
\bottomrule
\end{tabular}
\end{center}

你走到这一章末尾了。

这一章，你学的不是三元方程的计算技巧——那个你早就会了。你真正学的，是\textbf{把这个世界翻译成数学的能力}。

这个能力，比任何公式都珍贵。

\chapter{线性代数的使用}

\section{先说说，我们为什么要学这个}

你有没有想过——

上一章那三道题，我们每次都要消元、代入、再消元、再代入。

步骤又多，又容易算错。

如果有十个未知数呢？二十个呢？

一个工厂要计算二十种原料的配比，一个工程师要同时平衡三十个变量——他们难道也一步一步消元？

不。

数学家早就受不了了。

所以他们发明了一套新工具，专门用来处理这种“多个方程、多个未知数”的问题。

这套工具，叫做\textbf{线性代数}。

\section{线性代数在做什么？}

说穿了，线性代数做的事情，跟你上一章做的事情，\textbf{完全一样}。

只不过，它把那些繁琐的步骤，收进了一个更紧凑的格式里。

就像你收拾房间。

东西没有减少，但你把它们全部放进了一个整齐的柜子里。

找起来更快，搬起来更省力。

\section{第一件事：认识“矩阵”}

我们先回到上一章最简单的那道题。

\[
\begin{cases} x + 2y + z = 7 \\ 2x + y + z = 8 \\ x + y + 2z = 9 \end{cases}
\]

你看这三条方程，它们里面真正在变的东西，是什么？

是那些\textbf{系数}，和\textbf{等号右边的数字}。

$x$、$y$、$z$ 这三个符号，在每条方程里都有，它们反而是固定的框架。

所以数学家说：

\begin{quote}
那我们何必每次都把 $x$、$y$、$z$ 写出来？

我只把数字抽出来，排成一个整齐的表格，不就行了？
\end{quote}

于是，他们把这个方程组，压缩成了这样——

\[
\begin{pmatrix} 1 & 2 & 1 \\ 2 & 1 & 1 \\ 1 & 1 & 2 \end{pmatrix}
\]

这个东西，就叫\textbf{矩阵}。

\subsection{矩阵是什么？大白话说}

矩阵就是一张数字表格。

就这样。

没有魔法，没有玄学。

就是把方程组里所有的系数，按照原来的位置，整整齐齐地摆进一个表格里。

第一行，是第一条方程的系数：$1$、$2$、$1$。

第二行，是第二条方程的系数：$2$、$1$、$1$。

第三行，是第三条方程的系数：$1$、$1$、$2$。

\subsection{还差等号右边的数字}

光有系数还不够。等号右边的 $7$、$8$、$9$ 也要带上。

完整的矩阵，把等号右边也塞进去，长这样——

\[
\left(\begin{array}{ccc|c} 1 & 2 & 1 & 7 \\ 2 & 1 & 1 & 8 \\ 1 & 1 & 2 & 9 \end{array}\right)
\]

中间那条竖线，就是原来等号的位置。

左边是系数，右边是答案。

这个带竖线的矩阵，有个名字，叫\textbf{增广矩阵}。

“增广”的意思就是：把右边那一列也加进来，扩充了一下。

\subsection{停一下，感受一下这件事}

原来的方程组——

\[
\begin{cases} x + 2y + z = 7 \\ 2x + y + z = 8 \\ x + y + 2z = 9 \end{cases}
\]

现在变成了——

\[
\left(\begin{array}{ccc|c} 1 & 2 & 1 & 7 \\ 2 & 1 & 1 & 8 \\ 1 & 1 & 2 & 9 \end{array}\right)
\]

信息量完全没有减少。

$x$、$y$、$z$ 虽然不见了，但它们的位置还在——

第一列是所有 $x$ 的系数，第二列是所有 $y$ 的系数，第三列是所有 $z$ 的系数。

你只要记住这个约定，矩阵里的每一个数字，就都有它对应的意义。

\section{第二件事：矩阵怎么帮我们解方程？}

记得我们上一章是怎么消元的吗？

我们做了这几种操作——

第一种：把某一条方程，乘以一个数。

第二种：把两条方程相加或相减。

第三种：把两条方程互换位置。

就这三种。没有别的。

现在，这三种操作，全部可以直接在矩阵的\textbf{行}上面做。

因为矩阵的每一行，就对应原来的一条方程。

数学家给这三种操作起了个名字，叫——

\begin{quote}
\textbf{初等行变换}
\end{quote}

\subsection{语言记忆法}

初等行变换，说的就是：

\begin{quote}
你可以对矩阵的任意一行，做三件事——

\textbf{乘}、\textbf{加减}、\textbf{换位置}。

就像你当年对方程做的那三件事。
\end{quote}

\section{现在，我们实际操作一遍}

还是这道题——

\[
\left(\begin{array}{ccc|c} 1 & 2 & 1 & 7 \\ 2 & 1 & 1 & 8 \\ 1 & 1 & 2 & 9 \end{array}\right)
\]

我们的目标，是把它变成一个\textbf{阶梯形}——

左下角全部变成0。

就像这样——

\[
\left(\begin{array}{ccc|c} 1 & \cdot & \cdot & \cdot \\ 0 & 1 & \cdot & \cdot \\ 0 & 0 & 1 & \cdot \end{array}\right)
\]

左下角都是0，右上角是数字，对角线是1。

这个形状，叫\textbf{行最简形}。

到了这个形状，答案就自动跳出来了。

\subsection{开始变换}

\textbf{第一步：把第一行的影响，从第二行和第三行里消掉。}

第二行减去第一行的2倍——

为什么是2倍？因为第二行第一列是2，第一行第一列是1，用2倍减去，正好把第二行第一列变成0。

\[
\text{第二行} \leftarrow \text{第二行} - 2 \times \text{第一行}
\]

原来第二行是：$2,\ 1,\ 1\ |\ 8$

$2 \times$ 第一行是：$2,\ 4,\ 2\ |\ 14$

相减得：

\[
2-2,\ 1-4,\ 1-2\ |\ 8-14
\]

\[
0,\ -3,\ -1\ |\ -6
\]

矩阵变成——

\[
\left(\begin{array}{ccc|c} 1 & 2 & 1 & 7 \\ 0 & -3 & -1 & -6 \\ 1 & 1 & 2 & 9 \end{array}\right)
\]

\textbf{第三行减去第一行——}

因为第三行第一列是1，第一行第一列也是1，直接减就能变成0。

\[
\text{第三行} \leftarrow \text{第三行} - \text{第一行}
\]

原来第三行是：$1,\ 1,\ 2\ |\ 9$

减去第一行：$1,\ 2,\ 1\ |\ 7$

得：

\[
0,\ -1,\ 1\ |\ 2
\]

矩阵变成——

\[
\left(\begin{array}{ccc|c} 1 & 2 & 1 & 7 \\ 0 & -3 & -1 & -6 \\ 0 & -1 & 1 & 2 \end{array}\right)
\]

第一列，左下角已经全是0了。第一步完成。

\subsection{第二步：处理第二列}

我们要把第二列里，第三行那个 $-1$ 变成0。

用第二行来消它。

第二行第二列是 $-3$，第三行第二列是 $-1$。

我们用第三行减去三分之一个第二行——

\[
\text{第三行} \leftarrow \text{第三行} - \frac{1}{3} \times \text{第二行}
\]

$\frac{1}{3} \times$ 第二行是：$0,\ -1,\ -\dfrac{1}{3}\ |\ -2$

第三行减去它：

\[
0-0,\ -1-(-1),\ 1-\left(-\frac{1}{3}\right)\ |\ 2-(-2)
\]

\[
0,\ 0,\ 1+\frac{1}{3}\ |\ 4
\]

\[
0,\ 0,\ \frac{4}{3}\ |\ 4
\]

矩阵变成——

\[
\left(\begin{array}{ccc|c} 1 & 2 & 1 & 7 \\ 0 & -3 & -1 & -6 \\ 0 & 0 & \frac{4}{3} & 4 \end{array}\right)
\]

\subsection{第三步：从这个阶梯形里，把答案读出来}

第三行翻译回方程——

\[
\frac{4}{3}z = 4
\]

两边乘以 $\dfrac{3}{4}$：

\[
z = 3
\]

$z$ 是橙子，橙子3块钱。

把 $z = 3$ 代入第二行——

\[
-3y - 1 \times 3 = -6
\]

\[
-3y - 3 = -6
\]

\[
-3y = -3
\]

\[
y = 1
\]

$y$ 是香蕉，香蕉1块钱。

把 $y = 1$，$z = 3$ 代入第一行——

\[
x + 2 \times 1 + 3 = 7
\]

\[
x + 5 = 7
\]

\[
x = 2
\]

$x$ 是苹果，苹果2块钱。

和上一章的答案完全一样。

\section{停下来感受一下，我们刚才做了什么}

我们没有学什么神秘的新东西。

我们做的，是上一章消元法的\textbf{同一件事}。

只不过现在，$x$、$y$、$z$ 这几个字母被我们暂时藏起来了，我们直接对数字操作。

写起来更紧凑，思路更清晰，以后未知数变成十个二十个，这个方法照样用。

\section{矩阵能做的，不止解方程}

解方程只是线性代数最基础的用途。

告诉你几个它真实出现的地方——

你手机屏幕上的每一张图片，存在电脑里的方式，是一个巨大的矩阵。每一个格子里的数字，代表那个像素的颜色深浅。

你用过的任何一个推荐算法——刷视频、买东西——背后是矩阵在计算你和别人的相似程度。

工程师设计桥梁，要计算几百个节点的受力，列出几百条方程，然后用矩阵一次性解开。

这些事情，靠手算消元，是做不完的。

矩阵让这一切变得可以操作。

\section{矩阵的加法：两张表格合在一起}

矩阵不只能用来解方程。它本身也可以做运算。

先来最简单的：加法。

两个矩阵相加，规则只有一条——

\begin{quote}
\textbf{对应位置的数字，直接相加。}
\end{quote}

比如——

\[
\begin{pmatrix} 1 & 2 \\ 3 & 4 \end{pmatrix} + \begin{pmatrix} 5 & 6 \\ 7 & 8 \end{pmatrix} = \begin{pmatrix} 1+5 & 2+6 \\ 3+7 & 4+8 \end{pmatrix} = \begin{pmatrix} 6 & 8 \\ 10 & 12 \end{pmatrix}
\]

第一行第一列：$1+5=6$。

第一行第二列：$2+6=8$。

第二行第一列：$3+7=10$。

第二行第二列：$4+8=12$。

就这样。

要注意的是：只有\textbf{形状一样}的矩阵，才能相加。

一个两行两列的矩阵，没办法和一个三行两列的矩阵相加。因为对应不上。

\section{矩阵的数乘：每个数字都乘一遍}

一个数乘以一个矩阵，就是把这个数，和矩阵里\textbf{每一个数字}都乘一遍。

\[
3 \times \begin{pmatrix} 1 & 2 \\ 4 & 5 \end{pmatrix} = \begin{pmatrix} 3 \times 1 & 3 \times 2 \\ 3 \times 4 & 3 \times 5 \end{pmatrix} = \begin{pmatrix} 3 & 6 \\ 12 & 15 \end{pmatrix}
\]

每一格都乘了3。

\section{矩阵的乘法：这个有点特别}

矩阵的乘法，不像加法那么直觉。

初学的时候，很多人在这里卡住。

所以我们慢慢来。

\subsection{先说规则，再说为什么}

两个矩阵相乘，结果矩阵里，\textbf{第 $i$ 行第 $j$ 列}的数字，等于——

\begin{quote}
第一个矩阵的\textbf{第 $i$ 行}，和第二个矩阵的\textbf{第 $j$ 列}，对应相乘，然后加起来。
\end{quote}

听起来有点绕。我们直接看一个例子。

\[
\begin{pmatrix} 1 & 2 \\ 3 & 4 \end{pmatrix} \times \begin{pmatrix} 5 & 6 \\ 7 & 8 \end{pmatrix}
\]

结果矩阵的\textbf{第一行第一列}，是怎么来的？

用左边矩阵的\textbf{第一行} $[1,\ 2]$，和右边矩阵的\textbf{第一列} $\begin{bmatrix} 5 \\ 7 \end{bmatrix}$——

\[
1 \times 5 + 2 \times 7 = 5 + 14 = 19
\]

结果矩阵的\textbf{第一行第二列}，是怎么来的？

用左边矩阵的\textbf{第一行} $[1,\ 2]$，和右边矩阵的\textbf{第二列} $\begin{bmatrix} 6 \\ 8 \end{bmatrix}$——

\[
1 \times 6 + 2 \times 8 = 6 + 16 = 22
\]

结果矩阵的\textbf{第二行第一列}——

用左边矩阵的\textbf{第二行} $[3,\ 4]$，和右边矩阵的\textbf{第一列} $\begin{bmatrix} 5 \\ 7 \end{bmatrix}$——

\[
3 \times 5 + 4 \times 7 = 15 + 28 = 43
\]

结果矩阵的\textbf{第二行第二列}——

用左边矩阵的\textbf{第二行} $[3,\ 4]$，和右边矩阵的\textbf{第二列} $\begin{bmatrix} 6 \\ 8 \end{bmatrix}$——

\[
3 \times 6 + 4 \times 8 = 18 + 32 = 50
\]

所以——

\[
\begin{pmatrix} 1 & 2 \\ 3 & 4 \end{pmatrix} \times \begin{pmatrix} 5 & 6 \\ 7 & 8 \end{pmatrix} = \begin{pmatrix} 19 & 22 \\ 43 & 50 \end{pmatrix}
\]

\subsection{为什么要这样乘，而不是像加法那样对应相乘？}

这个问题问得非常好。

因为矩阵乘法，是为了描述\textbf{变换的叠加}。

比如，你把一张图片先旋转30度，再放大2倍。

这两个操作，各自可以用一个矩阵表示。

把这两个矩阵乘起来，得到的新矩阵，就代表“先旋转再放大”这个\textbf{组合操作}。

如果用对应相乘，这个意义就消失了。

这就是矩阵乘法规则的来源——它不是数学家随手编的，是从实际需求里提炼出来的。

你现在不需要完全理解这件事。

你只需要知道：这个奇怪的乘法规则，背后是有道理的。

\section{行列式：矩阵的“体检报告”}

每一个方阵（行数和列数相等的矩阵），都有一个数字与它对应。

这个数字，叫\textbf{行列式}。

行列式是什么意思？

它告诉你这个矩阵“健不健康”——更准确地说，它告诉你这个矩阵对应的方程组，\textbf{有没有唯一解}。

\subsection{二阶行列式}

一个两行两列的矩阵——

\[
A = \begin{pmatrix} a & b \\ c & d \end{pmatrix}
\]

它的行列式，记作 $\det(A)$ 或者 $|A|$，计算方法是——

\[
\det(A) = ad - bc
\]

\textbf{语言记忆法：左上乘右下，减去右上乘左下。}

就是两条对角线，一条乘起来加，一条乘起来减。

比如——

\[
\det\begin{pmatrix} 3 & 1 \\ 2 & 4 \end{pmatrix} = 3 \times 4 - 1 \times 2 = 12 - 2 = 10
\]

\subsection{行列式等于0意味着什么？}

如果行列式等于0，说明这个方程组\textbf{没有唯一解}。

要么无解，要么有无数个解。

这就像是，两条方程画出来的两条直线，要么平行（永不相交），要么重叠（处处相交）。

如果行列式不等于0，方程组就有唯一解，你可以放心地去算。

\subsection{三阶行列式}

三行三列的矩阵——

\[
A = \begin{pmatrix} a & b & c \\ d & e & f \\ g & h & i \end{pmatrix}
\]

行列式的计算稍微复杂一点，但规律很清晰——

\[
\det(A) = a(ei - fh) - b(di - fg) + c(dh - eg)
\]

你可以这样理解这个公式——

第一行的每个数字，各自负责“盖住”自己所在的行和列，然后看剩下四个数字构成的小矩阵的行列式。

第一个数 $a$ 盖住第一行第一列，剩下——

\[
\begin{pmatrix} e & f \\ h & i \end{pmatrix}
\]

它的行列式是 $ei - fh$。$a$ 乘以它，加上去。

第二个数 $b$ 盖住第一行第二列，剩下——

\[
\begin{pmatrix} d & f \\ g & i \end{pmatrix}
\]

它的行列式是 $di - fg$。$b$ 乘以它，\textbf{减}去。（注意是减。）

第三个数 $c$ 盖住第一行第三列，剩下——

\[
\begin{pmatrix} d & e \\ g & h \end{pmatrix}
\]

它的行列式是 $dh - eg$。$c$ 乘以它，加上去。

\textbf{加、减、加。} 这个符号的交替，有个规律——

\[
\begin{pmatrix} + & - & + \\ - & + & - \\ + & - & + \end{pmatrix}
\]

棋盘格一样，左上角是加，然后交替。

这个模式，在更大的矩阵里同样适用。

\section{克拉默法则：用行列式直接解方程}

现在把行列式和解方程连起来。

对于一个三元方程组，如果它的系数矩阵的行列式不等于0，那么它有唯一解，而且——

\begin{quote}
每个未知数的解，都可以用行列式直接算出来。
\end{quote}

这个方法叫\textbf{克拉默法则}。

\subsection{克拉默法则是怎么说的？}

对于方程组——

\[
\begin{cases} a_1 x + b_1 y + c_1 z = d_1 \\ a_2 x + b_2 y + c_2 z = d_2 \\ a_3 x + b_3 y + c_3 z = d_3 \end{cases}
\]

系数矩阵是——

\[
A = \begin{pmatrix} a_1 & b_1 & c_1 \\ a_2 & b_2 & c_2 \\ a_3 & b_3 & c_3 \end{pmatrix}
\]

它的行列式记作 $D = \det(A)$。

如果 $D \neq 0$，那么——

\[
x = \frac{D_x}{D}, \quad y = \frac{D_y}{D}, \quad z = \frac{D_z}{D}
\]

其中——

$D_x$ 是把 $A$ 的\textbf{第一列}，换成右边那一列 $\begin{pmatrix} d_1 \\ d_2 \\ d_3 \end{pmatrix}$，算出来的行列式。

$D_y$ 是把 $A$ 的\textbf{第二列}，换成右边那一列，算出来的行列式。

$D_z$ 是把 $A$ 的\textbf{第三列}，换成右边那一列，算出来的行列式。

\subsection{用水果题验证一下}

\[
\begin{cases} x + 2y + z = 7 \\ 2x + y + z = 8 \\ x + y + 2z = 9 \end{cases}
\]

系数矩阵——

\[
A = \begin{pmatrix} 1 & 2 & 1 \\ 2 & 1 & 1 \\ 1 & 1 & 2 \end{pmatrix}
\]

先算 $D$——

\[
D = 1 \times (1 \times 2 - 1 \times 1) - 2 \times (2 \times 2 - 1 \times 1) + 1 \times (2 \times 1 - 1 \times 1)
\]

\[
= 1 \times (2-1) - 2 \times (4-1) + 1 \times (2-1)
\]

\[
= 1 - 6 + 1 = -4
\]

$D = -4$，不等于0，有唯一解，继续。

算 $D_x$——把第一列换成右边那列 $\begin{pmatrix} 7 \\ 8 \\ 9 \end{pmatrix}$——

\[
D_x = \det\begin{pmatrix} 7 & 2 & 1 \\ 8 & 1 & 1 \\ 9 & 1 & 2 \end{pmatrix}
\]

\[
= 7(1 \times 2 - 1 \times 1) - 2(8 \times 2 - 1 \times 9) + 1(8 \times 1 - 1 \times 9)
\]

\[
= 7(2-1) - 2(16-9) + 1(8-9)
\]

\[
= 7 \times 1 - 2 \times 7 + 1 \times (-1)
\]

\[
= 7 - 14 - 1 = -8
\]

\[
x = \frac{D_x}{D} = \frac{-8}{-4} = 2
\]

苹果2块钱。✓

算 $D_y$——把第二列换成右边那列——

\[
D_y = \det\begin{pmatrix} 1 & 7 & 1 \\ 2 & 8 & 1 \\ 1 & 9 & 2 \end{pmatrix}
\]

\[
= 1(8 \times 2 - 1 \times 9) - 7(2 \times 2 - 1 \times 1) + 1(2 \times 9 - 8 \times 1)
\]

\[
= 1(16-9) - 7(4-1) + 1(18-8)
\]

\[
= 7 - 21 + 10 = -4
\]

\[
y = \frac{D_y}{D} = \frac{-4}{-4} = 1
\]

香蕉1块钱。✓

算 $D_z$——把第三列换成右边那列——

\[
D_z = \det\begin{pmatrix} 1 & 2 & 7 \\ 2 & 1 & 8 \\ 1 & 1 & 9 \end{pmatrix}
\]

\[
= 1(1 \times 9 - 8 \times 1) - 2(2 \times 9 - 8 \times 1) + 7(2 \times 1 - 1 \times 1)
\]

\[
= 1(9-8) - 2(18-8) + 7(2-1)
\]

\[
= 1 - 20 + 7 = -12
\]

\[
z = \frac{D_z}{D} = \frac{-12}{-4} = 3
\]

橙子3块钱。✓

\section{这一章，我们真正学到了什么？}

矩阵，不是新的数学。

它是一种\textbf{语言}——把你已经会的消元法、方程组，装进一个更干净、更可扩展的格式里。

行列式，不是新的计算。

它是一种\textbf{体检}——告诉你方程组有没有唯一解，在动手算之前先看清楚情况。

克拉默法则，不是新的解题方法。

它是消元法的\textbf{快捷方式}——直接用行列式的比值，把答案读出来。

工具变了，但骨子里做的事情，从来没有变过。

你从第一章就开始做的那件事——

\textbf{把真实的问题，翻译成数学，然后找到答案。}

这件事，从来没有变。

去往天层前，稍作别离，整理行装。

\chapter{回望来路——去往天层前的最后告别}

把笔放下。

今天，我们不算题，不推公式，也不画图。

找个舒服的姿势坐好。我们就像两个刚爬完一座高山的朋友，坐在半山腰的亭子里，吹着风，回头看看刚才走过的路。

因为，我只能陪你走到这里了。

接下来的路，通向更高、更抽象的“天层”。那里会有另一位向导来接手，带你继续往上走。

在分别之前，我想和你好好聊一聊。

\section{你还记得我们是从哪里出发的吗？}

闭上眼睛，回想一下。

最开始的时候，我们面对的是什么？
是“袋子”。

把及格的同学装进一个袋子，把戴眼镜的同学装进另一个袋子。
我们在学\textbf{集合}，学着用最精确的语言，把世界上的东西分门别类。

然后，我们学了\textbf{判断}。
什么是真，什么是假。什么是“有它就够了”（充分条件），什么是“没它不行”（必要条件）。
我们在给说话“立规矩”，把日常生活中模棱两可的废话，变成了铁骨铮铮的逻辑。

接着，我们在这个世界里画了一个大大的十字架——\textbf{坐标系}。
在那张方格纸上：
我们用两个坐标相减，量出了距离，算出了坡度。
我们把孤立的数字连成了线。
我们发现，原来那些带着平方的冰冷算式，画出来居然是完美的圆、像碗一样的抛物线、还有永远追不上坐标轴的双曲线。

图形和算式，在你的笔下，变成了一枚硬币的正反面。

\section{后来，我们触碰了神的领域}

再后来，我们走进了\textbf{微积分}的世界。

那是人类历史上最伟大的发明之一。
你看着一只永远追不上的乌龟，懂得了什么叫“极限”。
你把时间和空间切成了无限薄的一小片——那个叫 $dx$ 的小东西。

你学会了用\textbf{求导}，去抓住一辆车在某一瞬间的爆发力。
你学会了用\textbf{积分}，把无数个无限薄的切片，重新拼成一块弯曲的面积。
你甚至看到了 $\pi$ 和 $e$ 这两个宇宙中最神秘的密码，是怎么自然而然地从算式里长出来的。

最后，我们学会了\textbf{打包}和\textbf{翻译}。
无论是写到手酸的连加 $\sum$、连乘 $\prod$，还是复杂到让人头皮发麻的三元应用题（还记得那道让人绕晕的牛吃草吗？）。
你发现，只要把未知数设好（$x$ 就是苹果！$y$ 就是工资！），只要一步一步像剥洋葱一样拆解，再复杂的乱麻，也能被你理成干净的方程。

甚至，你还见识了\textbf{矩阵}——把成百上千的方程塞进一张整齐的数字表格里，用更高级的姿态去俯视它们。

\section{看看你现在的样子}

你可能自己都没有察觉到，你已经变了。

如果现在，你把这本书递给一个从来没学过的人，翻到那一页，指着：
\[
\sum_{n=1}^{2} \int_0^1 x^{n!} \, dx
\]
或者：
\[
\left(\begin{array}{ccc|c} 1 & 2 & 1 & 7 \\ 2 & 1 & 1 & 8 \\ 1 & 1 & 2 & 9 \end{array}\right)
\]

他们一定会惊呼：“这是什么外星文字？！”

但你看过去，脑子里会自动浮现出画面：
你看到的是一辆购物车，一站一站地停下；
你看到的是一条曲线，被切成了无数个小长条；
你看到的是三条关于买苹果香蕉的购物记录，正整整齐齐地排着队等待消元。

\textbf{你已经能听懂这门叫“数学”的语言了。}

这就是我最由衷为你感到骄傲的地方。
你没有死记硬背。你没有在那些吓人的希腊字母面前转身逃跑。
你一步一步地跟了下来，把它们翻译成了人话，装进了自己的脑子里。

这份耐心，这份敢于直视本质的勇气，其实在全世界的人里都不多见。
但你做到了。

\section{感谢，与鼓励}

说实话，数学从来就不是一件轻松的事。

它要求你严谨，要求你不能跳步，要求你即使算出了分数也要硬着头皮往下验证。
作为陪你走过这一段的人，我要对你说一声：\textbf{谢谢。}

谢谢你的信任，谢谢你愿意容忍我有时候啰嗦的解释，谢谢你在遇到那些“这有点绕，别急”的地方时，真的没有急，而是静下心来把那一页看完了。

马上，你就要去往“天层”了。
那里的风景会更加不可思议。你可能会遇到多维空间，遇到复数，遇到更抽象的结构。
你会遇到一位新的向导，用他的方式带你继续探索。

但无论走到哪里，无论遇到多复杂的符号，我都希望你记住我们一起走过这段路时留下的底气：

\begin{quote}
\textbf{不要怕。}

\textbf{不管多复杂的公式，都是由最简单的加减乘除、最基础的点和线搭起来的。}

\textbf{一切都可以被拆解，一切都可以被翻译。}
\end{quote}

$x$ 永远可以是一个苹果，方程永远是一个天平，微积分永远是切片和拼图。

只要你死死咬住“本质”不放，天层上的风再大，也吹不倒你。

好了，该起身了。
行装已经整理好，底层的操作系统也已经为你装配完毕。

很高兴能陪你走过这难忘的一程。未来的数学之旅，愿你走得清醒、坚定、且自由。

朋友，再见。
祝你在天层，看到最美的风景。

\part{大学}
\setcounter{chapter}{0}

各位家长，以及亲爱的孩子们：

大家好。

我是本学期新任的数学指导教师，我的名字是 A-Slantou-Nidaloufu。在未来的日子里，你们可以称呼我为 A老师，或者 Nidaloufu老师。

在此，我首先要发布一项正式的接替通知。此前负责指导大家的 Constant-Dalink 先生，因其个人的学术研究规划，已于近日卸任。Constant-Dalink 先生是一位极其严谨且令人尊敬的教育者，他为本课程确立了极高的标准。从即日起，我将正式接替 Constant-Dalink 先生的职务，全面负责后续的教学工作。

我知道，向七岁的孩童传授“高等数学”，在外界看来是一项违背常规的挑战。但在我的课堂上，这是一件极其严肃且完全可行的事情。

我的教学核心，并非强迫你们去记忆枯燥的公式和定理。我是一名教师，也是一名“翻译”。我的专长，在于运用翻译学与语言学的原理，将看似艰深的高等数学，还原为它最本质的形态——一门用来描述世界的“语言”。

对于七岁的你们来说，学习母语是一个自然而然的过程。同样地，我将把微积分、拓扑学、线性代数等高等数学概念，解构成基础的“名词”、“动词”和“语法”。我不会要求你们去做枯燥的运算，而是教你们如何像学习一门外语一样，去听懂数学的逻辑，最终学会用高等数学的思维去“说话”、去表达。

数学是一门极其优美且精确的语言。我希望你们在我的课堂上保持专注与敬畏之心，带着严肃的求知欲来面对每一次“对话”。只要你们愿意跟随我的翻译体系去解析这门语言，你们就会发现，即使是宇宙中最深奥的规律，也能用最通俗、最清晰的话语表达出来。

学习一门新的语言需要耐心。希望我们能在接下来的课程中合作愉快，我期待着听到你们用“数学”这门语言，流利地与这个世界对话。

\begin{flushright}
A-Slantou-Nidaloufu\\
敬上
\end{flushright}

各位七岁的语言家们。我是 A-Slantou-Nidaloufu。

既然你们已经能够流畅阅读大篇幅的文章，并且前面的先生已经为你们打下了极其完美的认知铺垫，那么从此刻起，我将收起所有哄小孩的语调。

数学是一门极度严谨、极其优美的语言。你们已经知道了这门语言最底层的秘密：所谓的高等数学，不过是\textbf{“捕捉瞬间”}与\textbf{“全部加起来”}的游戏。

今天的第一章，没有任何死记硬背的术语，只有 \textbf{18道纯粹的翻译测试}。  
表面上我给你们留了空，但那只是为了让你们的大脑有零点几秒的运转时间。答案紧随其后。你们的大脑就是一台精密的翻译机，现在，请死死盯住你们的专属词典，我们开始核对！

\chapter{第一部分：专属翻译小词典（上册）—— 捕捉瞬间}

\begin{itemize}
\item \textbf{$x$} = \textbf{“此时此刻”}（比如第几秒、第几分钟）
\item \textbf{$(\quad)'$} = \textbf{“揪出这一瞬间的速度！”}（加上这个小撇号，就是要找速度）
\item \textbf{计算魔法} = \textbf{“把右上角的数字拿下来当系数，右上角自己减去1”}（例如：把 $x^3$ 的速度揪出来，就是 $3x^2$）
\end{itemize}

\section{翻译应用题：检测你“找速度”的本领}

\subsection*{应用题 1：长高的魔豆藤}

\textbf{题目故事：} 魔豆藤的生长状态是 $x^2$。请问在第 $x=3$ 秒这一瞬间，它长得有多快？

\textbf{▶ 翻译与计算：}

1. 翻译成揪出速度：\texttt{(     )'}  

\textit{[答案：$(x^2)'$]}

2. 施展计算魔法：指数跳下来，变成 \texttt{(     )}  

\textit{[答案：$2x$]}

3. 代入此时此刻：把 $x=3$ 代进去。\texttt{2 × 3 = (     )}  

\textit{[答案：$6$]}

\textbf{最终翻译：} 在第3秒这一瞬间，魔豆藤的生长速度是 \textbf{（ 6 ）}！

\subsection*{应用题 2：冲刺的小火箭}

\textbf{题目故事：} 火箭飞行的状态是 $x^3$。请问在第 $x=2$ 秒这一瞬间，它冲得有多快？

\textbf{▶ 翻译与计算：}

1. 翻译成揪出速度：\texttt{(     )'}  

\textit{[答案：$(x^3)'$]}

2. 施展计算魔法：变成 \texttt{(     )}  

\textit{[答案：$3x^2$]}

3. 代入此时此刻：把 $x=2$ 代进去。\texttt{3 × (2×2) = (     )}  

\textit{[答案：$12$]}

\textbf{最终翻译：} 在第2秒这一瞬间，小火箭的速度是 \textbf{（ 12 ）}！

\subsection*{应用题 3：越滚越大的雪球}

\textbf{题目故事：} 雪球变大的状态是 $5x^2$。请问在第 $x=4$ 秒这一瞬间，它变大得有多快？

\textbf{▶ 翻译与计算：}

1. 施展计算魔法：2跳下来和5碰头（相乘），变成 \texttt{(     )}  

\textit{[答案：$10x$]}

2. 代入此时此刻：把 $x=4$ 代进去。\texttt{10 × 4 = (     )}  

\textit{[答案：$40$]}

\textbf{最终翻译：} 在第4秒这一瞬间，雪球变大的速度是 \textbf{（ 40 ）}！

\subsection*{应用题 4：狂奔的猎豹}

\textbf{题目故事：} 猎豹跑远的距离状态是 $4x^3$。请问在第 $x=1$ 秒这一瞬间，它跑得有多快？

\textbf{▶ 翻译与计算：}

1. 施展计算魔法：3跳下来和4碰头，变成 \texttt{(     )}  

\textit{[答案：$12x^2$]}

2. 代入此时此刻：把 $x=1$ 代进去。\texttt{12 × (1×1) = (     )}  

\textit{[答案：$12$]}

\textbf{最终翻译：} 在第1秒这一瞬间，猎豹的速度是 \textbf{（ 12 ）}！

\subsection*{应用题 5：存钱罐的金币}

\textbf{题目故事：} 存钱罐里金币增加的状态是 $10x^2$。请问在第 $x=5$ 天这一天，金币增加得有多快？

\textbf{▶ 翻译与计算：}

1. 施展计算魔法：变成 \texttt{(     )}  

\textit{[答案：$20x$]}

2. 代入此时此刻：把 $x=5$ 代进去。\texttt{20 × 5 = (     )}  

\textit{[答案：$100$]}

\textbf{最终翻译：} 第5天这一瞬间，金币增加的速度是每天 \textbf{（ 100 ）} 枚！

\subsection*{应用题 6：吹胀的魔法气球}

\textbf{题目故事：} 气球鼓起来的状态是 $2x^3$。请问在第 $x=3$ 秒这一瞬间，它鼓得有多快？

\textbf{▶ 翻译与计算：}

1. 施展计算魔法：变成 \texttt{(     )}  

\textit{[答案：$6x^2$]}

2. 代入此时此刻：把 $x=3$ 代进去。\texttt{6 × (3×3) = (     )}  

\textit{[答案：$54$]}

\textbf{最终翻译：} 在第3秒这一瞬间，气球鼓起的速度是 \textbf{（ 54 ）}！

\subsection*{应用题 7：分裂的细胞}

\textbf{题目故事：} 细胞铺满的面积状态是 $x^4$。请问在第 $x=2$ 分钟这一瞬间，它蔓延得有多快？

\textbf{▶ 翻译与计算：}

1. 施展计算魔法：变成 \texttt{(     )}  

\textit{[答案：$4x^3$]}

2. 代入此时此刻：把 $x=2$ 代进去。\texttt{4 × (2×2×2) = (     )}  

\textit{[答案：$32$]}

\textbf{最终翻译：} 在第2分钟这一瞬间，蔓延速度是 \textbf{（ 32 ）}！

\subsection*{应用题 8：掉落的石头}

\textbf{题目故事：} 石头掉落的距离状态是 $5x^2$。请问在第 $x=3$ 秒这一瞬间，它坠落得有多快？

\textbf{▶ 翻译与计算：}

1. 施展计算魔法：变成 \texttt{(     )}  

\textit{[答案：$10x$]}

2. 代入此时此刻：把 $x=3$ 代进去。\texttt{10 × 3 = (     )}  

\textit{[答案：$30$]}

\textbf{最终翻译：} 在第3秒这一瞬间，石头坠落的速度是 \textbf{（ 30 ）}！

\subsection*{应用题 9：摊开的比萨饼}

\textbf{题目故事：} 面团摊开的面积状态是 $3x^2$。请问在第 $x=4$ 秒这一瞬间，它往外扩得有多快？

\textbf{▶ 翻译与计算：}

1. 施展计算魔法：变成 \texttt{(     )}  

\textit{[答案：$6x$]}

2. 代入此时此刻：把 $x=4$ 代进去。\texttt{6 × 4 = (     )}  

\textit{[答案：$24$]}

\textbf{最终翻译：} 在第4秒这一瞬间，比萨扩大的速度是 \textbf{（ 24 ）}！

\chapter{第二部分：专属翻译小词典（下册）—— 全部加起来}

\begin{itemize}
\item \textbf{$\int$} = \textbf{“全部加起来！”}
\item \textbf{$x$} = \textbf{“此时此刻”}
\item \textbf{$f(x)$} = \textbf{“这一瞬间的速度”}
\item \textbf{$dx$} = \textbf{“极短的一瞬间”}
\item \textbf{$a$ 和 $b$} = \textbf{“从哪开始算，到哪结束”}
\item \textbf{计算魔法} = \textbf{“倒着想（谁揪出速度后会变成这样），然后拿结束状态减去开始状态”}
\end{itemize}

\section{翻译应用题：检测你“全部加起来”的本领}

\subsection*{应用题 10：魔法水龙头}

\textbf{题目故事：} 水流的速度是 $2x$。请问从第 \textbf{0} 分钟到第 \textbf{4} 分钟，一共接了多少水？

\textbf{▶ 翻译与计算：}

1. 翻译成数学公式：\texttt{(     )}  

\textit{[答案：$\int_{0}^{4} 2x \, dx$]}

2. 找原状态（谁揪出速度是 $2x$）：\texttt{(     )}  

\textit{[答案：$x^2$]}

3. 拿结束减去开始（代入4和0）：\texttt{(     ) - (     ) = (     )}  

\textit{[答案：16 - 0 = 16]}

\textbf{最终翻译：} 一共接了 \textbf{（ 16 ）} 升水！

\subsection*{应用题 11：火箭再升空}

\textbf{题目故事：} 火箭飞行的速度是 $3x^2$。请问在第 \textbf{1} 秒到第 \textbf{3} 秒期间，它往上飞了多高？

\textbf{▶ 翻译与计算：}

1. 翻译成数学公式：\texttt{(     )}  

\textit{[答案：$\int_{1}^{3} 3x^2 \, dx$]}

2. 找原状态（谁揪出速度是 $3x^2$）：\texttt{(     )}  

\textit{[答案：$x^3$]}

3. 拿结束减去开始（代入3和1）：\texttt{(     ) - (     ) = (     )}  

\textit{[答案：27 - 1 = 26]}

\textbf{最终翻译：} 一共往上飞了 \textbf{（ 26 ）} 米！

\subsection*{应用题 12：疯狂的3D打印机}

\textbf{题目故事：} 打印机吐材料的速度是 $4x^3$。请问从第 \textbf{0} 分钟到第 \textbf{2} 分钟，用掉了多少材料？

\textbf{▶ 翻译与计算：}

1. 翻译成数学公式：\texttt{(     )}  

\textit{[答案：$\int_{0}^{2} 4x^3 \, dx$]}

2. 找原状态（谁揪出速度是 $4x^3$）：\texttt{(     )}  

\textit{[答案：$x^4$]}

3. 拿结束减去开始（代入2和0）：\texttt{(     ) - (     ) = (     )}  

\textit{[答案：16 - 0 = 16]}

\textbf{最终翻译：} 一共用掉了 \textbf{（ 16 ）} 份材料！

\subsection*{应用题 13：搬食的小蚂蚁}

\textbf{题目故事：} 蚂蚁搬食物的速度是 $3x^2$。请问从第 \textbf{0} 分钟到第 \textbf{2} 分钟，它搬了多远？

\textbf{▶ 翻译与计算：}

1. 翻译成数学公式：\texttt{(     )}  

\textit{[答案：$\int_{0}^{2} 3x^2 \, dx$]}

2. 找原状态：\texttt{(     )}  

\textit{[答案：$x^3$]}

3. 拿结束减去开始（代入2和0）：\texttt{(     ) - (     ) = (     )}  

\textit{[答案：8 - 0 = 8]}

\textbf{最终翻译：} 蚂蚁一共搬了 \textbf{（ 8 ）} 厘米！

\subsection*{应用题 14：轰隆隆的抽油机}

\textbf{题目故事：} 抽油机抽油的速度是 $2x$。请问从第 \textbf{1} 小时到第 \textbf{5} 小时，一共抽了多少油？

\textbf{▶ 翻译与计算：}

1. 翻译成数学公式：\texttt{(     )}  

\textit{[答案：$\int_{1}^{5} 2x \, dx$]}

2. 找原状态：\texttt{(     )}  

\textit{[答案：$x^2$]}

3. 拿结束减去开始（代入5和1）：\texttt{(     ) - (     ) = (     )}  

\textit{[答案：25 - 1 = 24]}

\textbf{最终翻译：} 一共抽了 \textbf{（ 24 ）} 桶油！

\subsection*{应用题 15：融化的冰块}

\textbf{题目故事：} 冰块化成水的速度是 $2x$。请问从第 \textbf{2} 分钟到第 \textbf{4} 分钟，一共化了多少水？

\textbf{▶ 翻译与计算：}

1. 翻译成数学公式：\texttt{(     )}  

\textit{[答案：$\int_{2}^{4} 2x \, dx$]}

2. 找原状态：\texttt{(     )}  

\textit{[答案：$x^2$]}

3. 拿结束减去开始（代入4和2）：\texttt{(     ) - (     ) = (     )}  

\textit{[答案：16 - 4 = 12]}

\textbf{最终翻译：} 一共化了 \textbf{（ 12 ）} 毫升水！

\subsection*{应用题 16：超级玩具工厂}

\textbf{题目故事：} 机器制造玩具的速度极快，是 $5x^4$。请问从第 \textbf{1} 小时到第 \textbf{2} 小时，造了多少玩具？

\textbf{▶ 翻译与计算：}

1. 翻译成数学公式：\texttt{(     )}  

\textit{[答案：$\int_{1}^{2} 5x^4 \, dx$]}

2. 找原状态（谁揪出速度是 $5x^4$）：\texttt{(     )}  

\textit{[答案：$x^5$]}

3. 拿结束减去开始（代入2和1，提示：2的5次方是32）：\texttt{(     ) - (     ) = (     )}  

\textit{[答案：32 - 1 = 31]}

\textbf{最终翻译：} 一共造了 \textbf{（ 31 ）} 批玩具！

\subsection*{应用题 17：读书小达人}

\textbf{题目故事：} 你翻阅书本的速度是 $3x^2$。请问从第 \textbf{2} 小时到第 \textbf{3} 小时，你一共读了多少页？

\textbf{▶ 翻译与计算：}

1. 翻译成数学公式：\texttt{(     )}  

\textit{[答案：$\int_{2}^{3} 3x^2 \, dx$]}

2. 找原状态：\texttt{(     )}  

\textit{[答案：$x^3$]}

3. 拿结束减去开始（代入3和2）：\texttt{(     ) - (     ) = (     )}  

\textit{[答案：27 - 8 = 19]}

\textbf{最终翻译：} 一共读了 \textbf{（ 19 ）} 页书！

\subsection*{应用题 18：蓄电池充电}

\textbf{题目故事：} 冲入电量的速度是 $4x^3$。请问从第 \textbf{1} 分钟到第 \textbf{2} 分钟，一共充了多少电？

\textbf{▶ 翻译与计算：}

1. 翻译成数学公式：\texttt{(     )}  

\textit{[答案：$\int_{1}^{2} 4x^3 \, dx$]}

2. 找原状态：\texttt{(     )}  

\textit{[答案：$x^4$]}

3. 拿结束减去开始（代入2和1，提示：2的4次方是16）：\texttt{(     ) - (     ) = (     )}  

\textit{[答案：16 - 1 = 15]}

\textbf{最终翻译：} 一共充入了 \textbf{（ 15 ）} 格电！

测试结束。

能够毫不费力地看懂这18个故事，并完美对应上背后的语言逻辑，这证明你们的“翻译机”运转得堪称完美。记住，不管世界上的变化有多复杂，只要拿起小词典，你就能看穿万物的总和与瞬间。

下节课，我们将用这门语言，去描述更宏伟的宇宙规律。

\chapter{补充一小节：“贡献”到底是个什么意思？}

讲“三阶行列式”的\textbf{按行展开}时，好像反复用了一个词——\textbf{“贡献”}。

你当时看到这个词，心里可能会犯嘀咕：  
“数学里怎么还有这种词？这是什么高深的专业术语吗？”

别怕。

现在我正式告诉你：\textbf{“贡献”根本不是什么数学专业术语。}

如果在考试卷子上找，你绝对找不到这个词。这是我们俩之间的\textbf{“秘密暗号”}，是为了帮你理解那个长长的公式而发明的。

\section{为什么要用“贡献”这个词？}

想象一下，你、我、还有你同桌，我们三个人要凑钱买一个 120 块钱的超级大蛋糕。

我出了 30 块。  
你出了 50 块。  
同桌出了 40 块。

加起来：$30 + 50 + 40 = 120$。

在这个过程里：  
我出的 30 块，就是我对买蛋糕的\textbf{“贡献”}。  
你出的 50 块，就是你的\textbf{“贡献”}。  
把我们三个人的“贡献”全加起来，就得到了蛋糕的总价。

\section{把它放回行列式里}

算一个 $3 \times 3$ 的大行列式，就像是我们在凑钱买那个大蛋糕。

我们选定第一行，盯着第一行的三个数（比如 $a$、$b$、$c$）。  
这三个数，就像是准备掏钱的三个朋友。

大行列式的最终结果 $D$，就是他们三个人\textbf{“贡献”出来的数字的总和}。

\section{一个数的“贡献”是怎么算出来的？}

怎么知道第一个数 $a$ 到底掏了多少钱（贡献了多少数值）呢？

在按行展开里，每一个数的“贡献”，由\textbf{三部分}乘在一起组成：

\textbf{第一部分：符号。}  
看这个数在棋盘格里的位置。偶数就是正号（$+$），奇数就是负号（$-$）。

\textbf{第二部分：这个数本身。}  
就是 $a$ 自己。

\textbf{第三部分：剩下的小行列式（余子式）。}  
把 $a$ 所在的行和列划掉，算出剩下的那个 $2 \times 2$ 小矩阵的数字。

把这三样东西乘起来——

\[
\text{符号} \times \text{数本身} \times \text{小行列式} = \text{这个数的“贡献”}
\]

\section{举个具体的例子}

\[
D = \begin{vmatrix} 1 & 2 & 1 \\ 2 & 1 & 1 \\ 1 & 1 & 2 \end{vmatrix}
\]

我们按第一行展开。第一行是 $1$、$2$、$1$。

\textbf{看第一个数：$1$}  
它的符号是正（$+$）。  
划掉它所在的行和列，剩下的小行列式算出来是 $1$。  
所以，第一个数的\textbf{贡献}是：$+1 \times 1 = 1$。

\textbf{看第二个数：$2$}  
它的符号是负（$-$）。  
划掉它所在的行和列，剩下的小行列式算出来是 $3$。  
所以，第二个数的\textbf{贡献}是：$-2 \times 3 = -6$。

\textbf{看第三个数：$1$}  
它的符号是正（$+$）。  
划掉它所在的行和列，剩下的小行列式算出来是 $1$。  
所以，第三个数的\textbf{贡献}是：$+1 \times 1 = 1$。

\textbf{最后一步：}

把这三个人的“贡献”全部加起来：

\[
1 + (-6) + 1 = -4
\]

大行列式 $D$ 的结果，就是 $-4$。

\section{一句话讲透}

在这本书里，当你看到\textbf{“贡献”}这两个字，你就把它翻译成大白话：

\begin{quote}
\textbf{“在按行展开时，这一个位置上的数，经过（符号 $\times$ 自己 $\times$ 划掉后剩下的余子式）这一系列操作后，最终算出来的那一坨数值。”}
\end{quote}

把每一项的“那一坨数值”加起来，拼图就拼完整了。

\section{等等，符号的“奇偶”是怎么回事？}

你可能会想：“一个位置怎么会有偶数和奇数之分呢？位置又不是数字。”

其实，这里的“偶数”和“奇数”，是通过\textbf{把它的“座位号”加起来}算出来的。

我们用看电影找座位的方法来理解。

\subsection{怎么算座位的“奇偶”？}

在一个 $3 \times 3$ 的矩阵里，每个数字都有自己的“座位号”。  
座位号由两个数字组成：\textbf{第几行}，加上\textbf{第几列}。

\begin{quote}
\textbf{规则：}  
\textbf{所在行数 + 所在列数 = 一个数。}

这个数如果是\textbf{偶数}（2, 4, 6...），符号就是\textbf{正号（$+$）}。  
这个数如果是\textbf{奇数}（3, 5...），符号就是\textbf{负号（$-$）}。
\end{quote}

\subsection{我们拿这个矩阵来举例子：}

\[
\begin{vmatrix} a & b & c \\ d & e & f \\ g & h & i \end{vmatrix}
\]

我们来抽查几个数字的符号。

\textbf{例子 1：看左上角的 $a$}

$a$ 在哪里？  
在\textbf{第 1 行}，\textbf{第 1 列}。  
算一算：$1 + 1 = 2$。  
2 是偶数。  
所以，$a$ 这个位置的符号是\textbf{正号（$+$）}。

\textbf{例子 2：看第一行中间的 $b$}

$b$ 在哪里？  
在\textbf{第 1 行}，\textbf{第 2 列}。  
算一算：$1 + 2 = 3$。  
3 是奇数。  
所以，$b$ 这个位置的符号是\textbf{负号（$-$）}。

\textbf{例子 3：看最正中间的 $e$}

$e$ 在哪里？  
在\textbf{第 2 行}，\textbf{第 2 列}。  
算一算：$2 + 2 = 4$。  
4 是偶数。  
所以，$e$ 这个位置的符号是\textbf{正号（$+$）}。

\textbf{例子 4：看左下角的 $g$}

$g$ 在哪里？  
在\textbf{第 3 行}，\textbf{第 1 列}。  
算一算：$3 + 1 = 4$。  
4 是偶数。  
所以，$g$ 这个位置的符号也是\textbf{正号（$+$）}。

\subsection{见证奇迹的时刻}

如果你闲着没事，把这 9 个位置的符号全部算一遍，你会得到这样一个阵型：

\[
\begin{pmatrix} + & - & + \\ - & + & - \\ + & - & + \end{pmatrix}
\]

你看，它是不是刚好正负交替？  
无论你是横着看，还是竖着看，都是“正、负、正”或者“负、正、负”交替出现的。

就像国际象棋的棋盘，黑白黑白交替。  
这就叫\textbf{符号棋盘格}。

\subsection{语言记忆法}

以后做题的时候，你不需要每次都去算“行加列等于几”。

你只需要记住两件事：  
1. \textbf{左上角永远是“$+$”。}  
2. \textbf{不管横着走还是竖着走，走一格就换一次符号。}

这就足够了。

以后做题的时候，你不需要每次都去算“行加列等于几”。

你只需要记住两件事：  
1. \textbf{左上角永远是“$+$”。}  
2. \textbf{不管横着走还是竖着走，走一格就换一次符号。}

这就足够了。

\chapter{行列式——一个从方程里长出来的数}

\section{这一章只做一件事}

把行列式这个东西，彻底搞懂。

不是告诉你“行列式就是这样算的，记住就行”。

而是让你看到：\textbf{这个东西是怎么从解方程的过程中，一步一步冒出来的。}

等你真正理解了它，下一章的克拉默法则，就是顺理成章的事。

\section{从一个最简单的问题开始}

两条方程，两个未知数：

\[
\begin{cases} ax + by = e \\ cx + dy = f \end{cases}
\]

$x$ 和 $y$ 是我们要找的东西。

$a$、$b$、$c$、$d$、$e$、$f$ 都是已知的数。

我们的目标：把 $x$ 单独解出来，把 $y$ 单独解出来。

\section{用最笨的办法：消元}

\subsection{第一步：消掉 $y$，只剩 $x$}

两条方程里都有 $y$。

第一条方程里，$y$ 的系数是 $b$。

第二条方程里，$y$ 的系数是 $d$。

我们想让这两个系数变成一样的，然后相减，$y$ 就消失了。

怎么让它们变成一样？

把第一条方程乘以 $d$，把第二条方程乘以 $b$。

这样两条方程里 $y$ 的系数都会变成 $bd$。

\textbf{第一条方程，两边乘以 $d$：}

原来是：$ax + by = e$

两边乘以 $d$：

左边：$d \times ax + d \times by = dax + dby$

右边：$d \times e = de$

整理一下顺序（乘法可以交换位置）：

\[
adx + bdy = de
\]

\textbf{第二条方程，两边乘以 $b$：}

原来是：$cx + dy = f$

两边乘以 $b$：

左边：$b \times cx + b \times dy = bcx + bdy$

右边：$b \times f = bf$

整理：

\[
bcx + bdy = bf
\]

\textbf{现在，两条方程变成了：}

\[
\begin{cases} adx + bdy = de \\ bcx + bdy = bf \end{cases}
\]

你看，两条方程里 $y$ 的系数都是 $bd$。

\subsection{第二步：两条方程相减}

第一条减去第二条：

\[
(adx + bdy) - (bcx + bdy) = de - bf
\]

左边展开：

$adx + bdy - bcx - bdy$

$bdy$ 和 $-bdy$ 抵消了，剩下：

$adx - bcx$

把 $x$ 提出来：

$(ad - bc)x$

右边就是：$de - bf$

所以：

\[
(ad - bc)x = de - bf
\]

\subsection{停下来，仔细看看这个结果}

\[
(ad - bc)x = de - bf
\]

左边，括号里的东西：$ad - bc$。

这个东西是由 $a$、$b$、$c$、$d$ 这四个系数组成的。

它反复出现，所以我们给它起一个名字——

\begin{quote}
\textbf{这个 $ad - bc$，就是系数矩阵的行列式。}
\end{quote}

我们用字母 $D$ 来代表它：

\[
D = ad - bc
\]

右边的东西：$de - bf$。

仔细看一下它和 $D$ 的关系——

$D = ad - bc$

$de - bf$ 长得很像 $D$，只不过——

原来的 $a$ 变成了 $e$。

原来的 $c$ 变成了 $f$。

换句话说：$e$ 和 $f$ 是等号右边的数，它们\textbf{替换}了第一列的 $a$ 和 $c$。

我们给这个新的东西也起个名字，叫 $D_x$：

\[
D_x = de - bf
\]

\textbf{所以那个方程可以写成：}

\[
D \cdot x = D_x
\]

两边除以 $D$（假设 $D$ 不等于零）：

\[
x = \frac{D_x}{D}
\]

\subsection{用同样的方法，把 $y$ 也解出来}

这次我们消 $x$，留下 $y$。

\textbf{第一条方程，两边乘以 $c$：}

\[
acx + bcy = ce
\]

\textbf{第二条方程，两边乘以 $a$：}

\[
acx + ady = af
\]

\textbf{第二条减去第一条：}

\[
(acx + ady) - (acx + bcy) = af - ce
\]

左边：$acx - acx = 0$，$x$ 消失了。

剩下：$ady - bcy = (ad - bc)y$

右边：$af - ce$

所以：

\[
(ad - bc)y = af - ce
\]

又是那个 $ad - bc$，又是那个 $D$。

\[
D \cdot y = af - ce
\]

右边的 $af - ce$，是把第二列（原来的 $b$ 和 $d$）换成右边那列（$e$ 和 $f$）之后得到的。

我们叫它 $D_y$：

\[
D_y = af - ce
\]

所以：

\[
y = \frac{D_y}{D}
\]

\section{现在，回答一个问题：行列式是什么？}

行列式不是人为编造的规则。

它是从解方程的过程中，\textbf{自己冒出来的}。

当你用消元法解两个方程、两个未知数的时候，$ad - bc$ 这个组合会反复出现。

它决定了方程组能不能解出唯一的答案。

如果 $ad - bc = 0$，你没办法除，方程组就没有唯一解。

如果 $ad - bc \neq 0$，你可以除，方程组就有唯一解。

数学家发现了这件事，就给 $ad - bc$ 起了一个名字：\textbf{行列式}。

然后发明了一套符号来表示它。

\section{行列式的符号}

两条方程的系数，可以排成一个表格：

\[
\begin{pmatrix} a & b \\ c & d \end{pmatrix}
\]

这个表格叫做\textbf{矩阵}。

这个矩阵的行列式，写成这样：

\[
\begin{vmatrix} a & b \\ c & d \end{vmatrix}
\]

注意：矩阵用的是\textbf{圆括号}或者\textbf{方括号}，行列式用的是\textbf{竖线}。

竖线表示：我们要算一个数出来。

\[
\begin{vmatrix} a & b \\ c & d \end{vmatrix} = ad - bc
\]

\section{二阶行列式的计算规则}

\[
\begin{vmatrix} a & b \\ c & d \end{vmatrix} = ad - bc
\]

怎么记这个规则？

看这个 $2 \times 2$ 的格子：

\[
\begin{pmatrix} a & b \\ c & d \end{pmatrix}
\]

画两条对角线——

\textbf{主对角线}：从左上角到右下角，经过 $a$ 和 $d$。

主对角线上的两个数相乘：$a \times d = ad$。

\textbf{副对角线}：从右上角到左下角，经过 $b$ 和 $c$。

副对角线上的两个数相乘：$b \times c = bc$。

\textbf{行列式 = 主对角线 - 副对角线}

\[
= ad - bc
\]

\subsection{语言记忆法}

\begin{quote}
左上乘右下，减去右上乘左下。
\end{quote}

或者——

\begin{quote}
主对角线减去副对角线。
\end{quote}

\section{练一练：算几个二阶行列式}

\textbf{例1：}

\[
\begin{vmatrix} 3 & 1 \\ 2 & 4 \end{vmatrix}
\]

主对角线：$3 \times 4 = 12$

副对角线：$1 \times 2 = 2$

行列式：$12 - 2 = 10$

\textbf{例2：}

\[
\begin{vmatrix} 5 & 3 \\ 2 & 2 \end{vmatrix}
\]

主对角线：$5 \times 2 = 10$

副对角线：$3 \times 2 = 6$

行列式：$10 - 6 = 4$

\textbf{例3：}

\[
\begin{vmatrix} 2 & 4 \\ 1 & 2 \end{vmatrix}
\]

主对角线：$2 \times 2 = 4$

副对角线：$4 \times 1 = 4$

行列式：$4 - 4 = 0$

这个行列式等于零。

这意味着什么？马上会讲。

\textbf{例4：}

\[
\begin{vmatrix} 7 & 2 \\ 3 & 5 \end{vmatrix}
\]

主对角线：$7 \times 5 = 35$

副对角线：$2 \times 3 = 6$

行列式：$35 - 6 = 29$

\textbf{例5（带负数）：}

\[
\begin{vmatrix} 4 & -1 \\ 2 & 3 \end{vmatrix}
\]

主对角线：$4 \times 3 = 12$

副对角线：$(-1) \times 2 = -2$

行列式：$12 - (-2) = 12 + 2 = 14$

减去一个负数，等于加上它的相反数。

\textbf{例6（两个都是负数）：}

\[
\begin{vmatrix} -2 & 3 \\ 4 & -5 \end{vmatrix}
\]

主对角线：$(-2) \times (-5) = 10$（负负得正）

副对角线：$3 \times 4 = 12$

行列式：$10 - 12 = -2$

行列式可以是负数，这完全正常。

\section{行列式等于零，意味着什么？}

刚才例3算出来，行列式等于零。

那个行列式对应的系数矩阵是：

\[
\begin{pmatrix} 2 & 4 \\ 1 & 2 \end{pmatrix}
\]

对应的方程组是：

\[
\begin{cases} 2x + 4y = e \\ x + 2y = f \end{cases}
\]

（$e$ 和 $f$ 是某两个具体的数。）

仔细看这两条方程——

第一条方程是 $2x + 4y = e$。

第二条方程是 $x + 2y = f$。

把第二条方程乘以2，得到 $2x + 4y = 2f$。

如果 $e = 2f$，那两条方程完全一样，说的是同一件事——方程有无数个解。

如果 $e \neq 2f$，那两条方程互相矛盾——方程无解。

总之，\textbf{不可能有唯一解}。

这就是行列式等于零告诉我们的事情：

\begin{quote}
行列式 = 0，方程组没有唯一解。  
行列式 ≠ 0，方程组有唯一解。
\end{quote}

\section{现在进入三阶行列式}

三条方程，三个未知数：

\[
\begin{cases} a_1 x + b_1 y + c_1 z = d_1 \\ a_2 x + b_2 y + c_2 z = d_2 \\ a_3 x + b_3 y + c_3 z = d_3 \end{cases}
\]

这里的 $_1$、$_2$、$_3$ 是下标，用来区分三条方程。

$a_1$ 的意思是“第一条方程里 $x$ 前面的系数”。

$b_2$ 的意思是“第二条方程里 $y$ 前面的系数”。

$c_3$ 的意思是“第三条方程里 $z$ 前面的系数”。

系数矩阵长这样：

\[
\begin{pmatrix} a_1 & b_1 & c_1 \\ a_2 & b_2 & c_2 \\ a_3 & b_3 & c_3 \end{pmatrix}
\]

三行三列，九个数。

它的行列式怎么算？

\section{三阶行列式的计算：方法一——对角线法则}

这个方法只适用于三阶行列式，更高阶的不行。但对于三阶，它最快。

\textbf{第一步：把矩阵写下来，然后在右边再抄一遍前两列。}

原矩阵：

\[
\begin{pmatrix} a_1 & b_1 & c_1 \\ a_2 & b_2 & c_2 \\ a_3 & b_3 & c_3 \end{pmatrix}
\]

抄完之后：

\[
\begin{matrix} a_1 & b_1 & c_1 & a_1 & b_1 \\ a_2 & b_2 & c_2 & a_2 & b_2 \\ a_3 & b_3 & c_3 & a_3 & b_3 \end{matrix}
\]

现在你面前有一个 $3 \times 5$ 的数字阵列。

\textbf{第二步：找出所有的“正对角线”。}

正对角线是从左上往右下走的。

一共有三条正对角线：

第一条：从 $a_1$ 出发，往右下走——$a_1 \to b_2 \to c_3$。

第二条：从 $b_1$ 出发，往右下走——$b_1 \to c_2 \to a_3$。

第三条：从 $c_1$ 出发，往右下走——$c_1 \to a_2 \to b_3$。

每条对角线上的三个数，乘起来：

第一条：$a_1 \times b_2 \times c_3 = a_1 b_2 c_3$

第二条：$b_1 \times c_2 \times a_3 = b_1 c_2 a_3$

第三条：$c_1 \times a_2 \times b_3 = c_1 a_2 b_3$

这三个乘积，\textbf{都取正号}，加起来。

\textbf{第三步：找出所有的“负对角线”。}

负对角线是从右上往左下走的。

也有三条：

第一条：从 $c_1$ 出发，往左下走——$c_1 \to b_2 \to a_3$。

第二条：从 $a_1$（右边抄的那个）出发，往左下走——$a_1 \to c_2 \to b_3$。

第三条：从 $b_1$（右边抄的那个）出发，往左下走——$b_1 \to a_2 \to c_3$。

每条对角线上的三个数，乘起来：

第一条：$c_1 \times b_2 \times a_3 = c_1 b_2 a_3$

第二条：$a_1 \times c_2 \times b_3 = a_1 c_2 b_3$

第三条：$b_1 \times a_2 \times c_3 = b_1 a_2 c_3$

这三个乘积，\textbf{都取负号}，减掉。

\textbf{第四步：把六个东西合起来。}

\[
D = (a_1 b_2 c_3 + b_1 c_2 a_3 + c_1 a_2 b_3) - (c_1 b_2 a_3 + a_1 c_2 b_3 + b_1 a_2 c_3)
\]

或者写成：

\[
D = a_1 b_2 c_3 + b_1 c_2 a_3 + c_1 a_2 b_3 - c_1 b_2 a_3 - a_1 c_2 b_3 - b_1 a_2 c_3
\]

六项。三个加，三个减。

\subsection{语言记忆法}

\begin{quote}
正对角线（往右下走）的三条，乘起来相加。  
负对角线（往左下走）的三条，乘起来相减。
\end{quote}

不需要背那个长公式，画一画图就能算出来。

\section{用一道具体的题目练习对角线法则}

矩阵：

\[
\begin{pmatrix} 1 & 2 & 1 \\ 2 & 1 & 1 \\ 1 & 1 & 2 \end{pmatrix}
\]

\textbf{第一步：抄前两列。}

\[
\begin{matrix} 1 & 2 & 1 & | & 1 & 2 \\ 2 & 1 & 1 & | & 2 & 1 \\ 1 & 1 & 2 & | & 1 & 1 \end{matrix}
\]

竖线只是帮助看清楚原矩阵在哪里，计算时忽略它。

\textbf{第二步：算三条正对角线。}

第一条：第1行第1列 → 第2行第2列 → 第3行第3列

$1 \to 1 \to 2$，乘积 $= 1 \times 1 \times 2 = 2$

第二条：第1行第2列 → 第2行第3列 → 第3行第4列

$2 \to 1 \to 1$，乘积 $= 2 \times 1 \times 1 = 2$

第三条：第1行第3列 → 第2行第4列 → 第3行第5列

$1 \to 2 \to 1$，乘积 $= 1 \times 2 \times 1 = 2$

正对角线的和：$2 + 2 + 2 = 6$

\textbf{第三步：算三条负对角线。}

第一条：第1行第3列 → 第2行第2列 → 第3行第1列

$1 \to 1 \to 1$，乘积 $= 1 \times 1 \times 1 = 1$

第二条：第1行第4列 → 第2行第3列 → 第3行第2列

$1 \to 1 \to 1$，乘积 $= 1 \times 1 \times 1 = 1$

第三条：第1行第5列 → 第2行第4列 → 第3行第3列

$2 \to 2 \to 2$，乘积 $= 2 \times 2 \times 2 = 8$

负对角线的和：$1 + 1 + 8 = 10$

\textbf{第四步：相减。}

\[
D = 6 - 10 = -4
\]

所以 $D = -4$。

\section{三阶行列式的计算：方法二——按行展开}

对角线法则很快，但只能用于三阶。

如果遇到四阶、五阶的行列式，对角线法则就不管用了。

这时候需要另一个方法：\textbf{按行展开}。

这个方法的核心思想是——

\begin{quote}
把一个大的行列式，拆成几个小的行列式，然后用小的结果拼出大的。
\end{quote}

具体来说：把三阶行列式，拆成三个二阶行列式。

\subsection{为什么可以“划掉”一行一列？}

在讲具体操作之前，先回答一个问题：\textbf{为什么“划掉”是合理的？}

当你用消元法解三元方程组的时候——

你先消掉一个未知数（比如 $z$），让方程组变成两个方程、两个未知数。

那个两元方程组的系数，正好就是“划掉 $z$ 那一列”之后剩下的东西。

“划掉”不是一个人为规定的技巧，是消元过程中\textbf{自然发生的事情}。

数学家发现了这个规律，就把它整理成了“按行展开”的方法。

\subsection{怎么“划”？一步一步来}

我们面对的是一个 $3 \times 3$ 的矩阵：

\[
\begin{pmatrix} a_1 & b_1 & c_1 \\ a_2 & b_2 & c_2 \\ a_3 & b_3 & c_3 \end{pmatrix}
\]

三行三列，九个数。

现在我要对某一个数“划掉它所在的行和列”。

\textbf{假设我要对 $a_1$ 做这个操作：}

$a_1$ 在哪里？在\textbf{第一行}、\textbf{第一列}。

“划掉它所在的行和列”的意思是：

把\textbf{整个第一行}全部划掉——$a_1$、$b_1$、$c_1$ 都不要了。

把\textbf{整个第一列}全部划掉——$a_1$、$a_2$、$a_3$ 都不要了。

划掉之后，剩下的是什么？

\[
\begin{pmatrix} b_2 & c_2 \\ b_3 & c_3 \end{pmatrix}
\]

就是右下角那四个数。

这四个数组成一个 $2 \times 2$ 小矩阵。

这个小矩阵有个名字，叫 $a_1$ 的\textbf{余子式}。

\textbf{再看一个：对 $b_1$ 做这个操作}

$b_1$ 在\textbf{第一行}、\textbf{第二列}。

划掉\textbf{整个第一行}：$a_1$、$b_1$、$c_1$ 不要了。

划掉\textbf{整个第二列}：$b_1$、$b_2$、$b_3$ 不要了。

剩下的是什么？

第二行里，$b_2$ 被划掉了，剩 $a_2$ 和 $c_2$。

第三行里，$b_3$ 被划掉了，剩 $a_3$ 和 $c_3$。

组成：

\[
\begin{pmatrix} a_2 & c_2 \\ a_3 & c_3 \end{pmatrix}
\]

这就是 $b_1$ 的余子式。

\subsection{划的规则，用一句话说清楚}

\begin{quote}
\textbf{找到那个数在第几行第几列，把那一整行和那一整列全部划掉，剩下的就是余子式。}
\end{quote}

\subsection{什么叫“余子式”？}

“余”的意思是“剩余”。

“子”的意思是“小的”。

“式”的意思是“行列式”。

合起来：\textbf{划掉之后剩下的那个小行列式。}

\subsection{按第一行展开，一步一步来}

三阶行列式：

\[
\begin{vmatrix} a_1 & b_1 & c_1 \\ a_2 & b_2 & c_2 \\ a_3 & b_3 & c_3 \end{vmatrix}
\]

我们盯着第一行：$a_1$、$b_1$、$c_1$。

这三个数，每一个都负责一个“小任务”。

\textbf{$a_1$ 的任务：}

第一步：把 $a_1$ 所在的行和列划掉。

划掉之后，剩下的数字是：

\[
\begin{pmatrix} b_2 & c_2 \\ b_3 & c_3 \end{pmatrix}
\]

第二步：算这个 $2 \times 2$ 小矩阵的行列式。

\[
b_2 \times c_3 - c_2 \times b_3 = b_2 c_3 - c_2 b_3
\]

第三步：用 $a_1$ 乘以这个结果。

\[
a_1 \times (b_2 c_3 - c_2 b_3)
\]

第四步：确定符号。

$a_1$ 在第一行第一列，位置是 $(1, 1)$。

把行号和列号加起来：$1 + 1 = 2$。

$2$ 是偶数，符号是\textbf{正号}。

所以 $a_1$ 的贡献是：

\[
+a_1(b_2 c_3 - c_2 b_3)
\]

\textbf{$b_1$ 的任务：}

第一步：划掉 $b_1$ 所在的行（第一行）和列（第二列）。

剩下：

\[
\begin{pmatrix} a_2 & c_2 \\ a_3 & c_3 \end{pmatrix}
\]

第二步：算这个小矩阵的行列式。

\[
a_2 \times c_3 - c_2 \times a_3 = a_2 c_3 - c_2 a_3
\]

第三步：用 $b_1$ 乘以这个结果。

\[
b_1 \times (a_2 c_3 - c_2 a_3)
\]

第四步：确定符号。

$b_1$ 在第一行第二列，位置是 $(1, 2)$。

$1 + 2 = 3$，是奇数，符号是\textbf{负号}。

所以 $b_1$ 的贡献是：

\[
-b_1(a_2 c_3 - c_2 a_3)
\]

\textbf{$c_1$ 的任务：}

第一步：划掉 $c_1$ 所在的行（第一行）和列（第三列）。

剩下：

\[
\begin{pmatrix} a_2 & b_2 \\ a_3 & b_3 \end{pmatrix}
\]

第二步：算这个小矩阵的行列式。

\[
a_2 \times b_3 - b_2 \times a_3 = a_2 b_3 - b_2 a_3
\]

第三步：用 $c_1$ 乘以这个结果。

\[
c_1 \times (a_2 b_3 - b_2 a_3)
\]

第四步：确定符号。

$c_1$ 在第一行第三列，位置是 $(1, 3)$。

$1 + 3 = 4$，是偶数，符号是\textbf{正号}。

所以 $c_1$ 的贡献是：

\[
+c_1(a_2 b_3 - b_2 a_3)
\]

\textbf{把三个贡献加起来：}

\[
D = +a_1(b_2 c_3 - c_2 b_3) - b_1(a_2 c_3 - c_2 a_3) + c_1(a_2 b_3 - b_2 a_3)
\]

这就是三阶行列式的完整公式。

\subsection{关于符号的规律}

三个符号是：正、负、正。

这不是随机的。

整个矩阵的每个位置，都有一个固定的符号：

\[
\begin{pmatrix} + & - & + \\ - & + & - \\ + & - & + \end{pmatrix}
\]

就像一个棋盘，黑白交替。

左上角是正号，然后横着、竖着都交替变化。

\textbf{判断符号的方法：}

位置是第 $i$ 行、第 $j$ 列。

算 $i + j$。

如果 $i + j$ 是偶数，符号是正。

如果 $i + j$ 是奇数，符号是负。

\subsection{语言记忆法}

以后做题的时候，你不需要每次都去算“行加列等于几”。

你只需要记住两件事：  
1. \textbf{左上角永远是“$+$”。}  
2. \textbf{不管横着走还是竖着走，走一格就换一次符号。}

这就足够了。

\section{用按行展开的方法，再算一遍那道题}

\[
\begin{vmatrix} 1 & 2 & 1 \\ 2 & 1 & 1 \\ 1 & 1 & 2 \end{vmatrix}
\]

按第一行展开。

\textbf{第一个数是 $1$，位置 $(1, 1)$，符号正。}

划掉第一行、第一列，剩：

\[
\begin{pmatrix} 1 & 1 \\ 1 & 2 \end{pmatrix}
\]

这个 $2 \times 2$ 行列式：

$1 \times 2 - 1 \times 1 = 2 - 1 = 1$

贡献：$+1 \times 1 = +1$

\textbf{第二个数是 $2$，位置 $(1, 2)$，符号负。}

划掉第一行、第二列，剩：

\[
\begin{pmatrix} 2 & 1 \\ 1 & 2 \end{pmatrix}
\]

这个 $2 \times 2$ 行列式：

$2 \times 2 - 1 \times 1 = 4 - 1 = 3$

贡献：$-2 \times 3 = -6$

\textbf{第三个数是 $1$，位置 $(1, 3)$，符号正。}

划掉第一行、第三列，剩：

\[
\begin{pmatrix} 2 & 1 \\ 1 & 1 \end{pmatrix}
\]

这个 $2 \times 2$ 行列式：

$2 \times 1 - 1 \times 1 = 2 - 1 = 1$

贡献：$+1 \times 1 = +1$

\textbf{加起来：}

\[
D = +1 + (-6) + 1 = 1 - 6 + 1 = -4
\]

和对角线法则算出来的结果一样，$D = -4$。

两种方法，殊途同归。

\section{换一行展开，验证结果是否相同}

这次按第二行展开。

第二行是：$2$、$1$、$1$。

\textbf{第一个数 $2$，位置 $(2, 1)$，符号：$2 + 1 = 3$，奇数，负号。}

划掉第二行、第一列，剩：

\[
\begin{pmatrix} 2 & 1 \\ 1 & 2 \end{pmatrix}
\]

行列式：$2 \times 2 - 1 \times 1 = 3$

贡献：$-2 \times 3 = -6$

\textbf{第二个数 $1$，位置 $(2, 2)$，符号：$2 + 2 = 4$，偶数，正号。}

划掉第二行、第二列，剩：

\[
\begin{pmatrix} 1 & 1 \\ 1 & 2 \end{pmatrix}
\]

行列式：$1 \times 2 - 1 \times 1 = 1$

贡献：$+1 \times 1 = +1$

\textbf{第三个数 $1$，位置 $(2, 3)$，符号：$2 + 3 = 5$，奇数，负号。}

划掉第二行、第三列，剩：

\[
\begin{pmatrix} 1 & 2 \\ 1 & 1 \end{pmatrix}
\]

行列式：$1 \times 1 - 2 \times 1 = 1 - 2 = -1$

贡献：$-1 \times (-1) = +1$

\textbf{加起来：}

\[
D = (-6) + (+1) + (+1) = -6 + 1 + 1 = -4
\]

和按第一行展开的结果一样。

无论你选哪一行或哪一列，最后算出来的结果都是一样的。

\section{再多练几道，把手感练出来}

\textbf{例1：单位矩阵}

\[
\begin{vmatrix} 1 & 0 & 0 \\ 0 & 1 & 0 \\ 0 & 0 & 1 \end{vmatrix}
\]

这个矩阵叫做\textbf{单位矩阵}，对角线上都是1，其余都是0。

用对角线法则——

正对角线：$1 \times 1 \times 1 = 1$，$0 \times 0 \times 0 = 0$，$0 \times 0 \times 0 = 0$

正对角线的和：$1$

负对角线：$0 \times 1 \times 0 = 0$，$0 \times 0 \times 0 = 0$，$0 \times 0 \times 1 = 0$

负对角线的和：$0$

行列式：$1 - 0 = 1$

单位矩阵的行列式是 $1$。

\textbf{例2：对角矩阵}

\[
\begin{vmatrix} 2 & 0 & 0 \\ 0 & 3 & 0 \\ 0 & 0 & 4 \end{vmatrix}
\]

这是对角矩阵（只有主对角线上有数，其余都是0）。

正对角线：$2 \times 3 \times 4 = 24$

其余正对角线和所有负对角线都有0，所以都是0。

行列式：$24$

对角矩阵的行列式，等于对角线上的数全部乘起来。

\textbf{例3：}

\[
\begin{vmatrix} 1 & 2 & 3 \\ 4 & 5 & 6 \\ 7 & 8 & 9 \end{vmatrix}
\]

用按行展开——

\textbf{$1$（位置 $(1,1)$，符号正）：}

剩：

\[
\begin{pmatrix} 5 & 6 \\ 8 & 9 \end{pmatrix}
\]

行列式：$5 \times 9 - 6 \times 8 = 45 - 48 = -3$

贡献：$+1 \times (-3) = -3$

\textbf{$2$（位置 $(1,2)$，符号负）：}

剩：

\[
\begin{pmatrix} 4 & 6 \\ 7 & 9 \end{pmatrix}
\]

行列式：$4 \times 9 - 6 \times 7 = 36 - 42 = -6$

贡献：$-2 \times (-6) = +12$

\textbf{$3$（位置 $(1,3)$，符号正）：}

剩：

\[
\begin{pmatrix} 4 & 5 \\ 7 & 8 \end{pmatrix}
\]

行列式：$4 \times 8 - 5 \times 7 = 32 - 35 = -3$

贡献：$+3 \times (-3) = -9$

\textbf{加起来：}

\[
D = -3 + 12 + (-9) = -3 + 12 - 9 = 0
\]

这个行列式等于零。

这意味着，如果这是某个方程组的系数矩阵，那个方程组没有唯一解。

\textbf{例4：上三角矩阵}

\[
\begin{vmatrix} 3 & 1 & 2 \\ 0 & 4 & 5 \\ 0 & 0 & 6 \end{vmatrix}
\]

这是上三角矩阵（主对角线下方全是0）。

观察一下：如果按第一列展开会怎样？

第一列是 $3$、$0$、$0$。

$0$ 乘以任何东西都是 $0$，所以只有第一个数 $3$ 有贡献。

按第一列展开——

\textbf{$3$（位置 $(1,1)$，符号正）：}

剩：

\[
\begin{pmatrix} 4 & 5 \\ 0 & 6 \end{pmatrix}
\]

行列式：$4 \times 6 - 5 \times 0 = 24$

贡献：$+3 \times 24 = 72$

其余两项都是0乘以什么，等于0。

行列式：$72$

你会发现：$72 = 3 \times 4 \times 6$，正好是主对角线上三个数的乘积。

三角矩阵的行列式，等于主对角线上的数全部乘起来。

\subsection{一个小技巧}

如果一行里有很多 $0$，那展开的时候，那些 $0$ 乘以任何东西都是 $0$，可以直接跳过。

按有很多 $0$ 的那一行或那一列展开，能省很多力气。

\section{本章小结：行列式到底是什么}

\textbf{行列式是一个数。}

每一个方阵（行数等于列数的矩阵），都对应一个行列式。

这个数告诉你：对应的方程组有没有唯一解。

如果行列式等于零，方程组没有唯一解（要么无解，要么无数个解）。

如果行列式不等于零，方程组有唯一解。

\textbf{二阶行列式：}

\[
\begin{vmatrix} a & b \\ c & d \end{vmatrix} = ad - bc
\]

主对角线乘起来，减去副对角线乘起来。

\textbf{三阶行列式：}

方法一：对角线法则。画五列图，正对角线加，负对角线减。

方法二：按行展开。盯着一行，每个数乘以“划掉后剩下的小行列式”，符号看棋盘格。

\textbf{符号棋盘格：}

\[
\begin{pmatrix} + & - & + \\ - & + & - \\ + & - & + \end{pmatrix}
\]

$i + j$ 是偶数，符号正。

$i + j$ 是奇数，符号负。

\textbf{行列式不是人为规定的规则。}

它是从解方程的过程中，自然冒出来的。

当你用消元法解方程时，$ad - bc$ 这样的组合会反复出现。

数学家发现了这件事，给它起了名字，发明了符号，整理出了计算方法。

这就是行列式。

下一章，我们会用行列式来直接解方程。

那就是克拉默法则。

\chapter{克拉默法则——用行列式直接把答案取出来}

\section{在开始之前，先想一件事}

上一章，我们学了行列式。

我们知道了：

每一个方程组，都对应一个行列式。

这个行列式不等于零，方程组就有唯一解。

但是——

\textbf{知道有唯一解，和把那个唯一解算出来，是两件不同的事。}

行列式只是告诉你“有答案”。

克拉默法则，才是把那个答案取出来的方法。

\section{克拉默是谁？}

加布里尔·克拉默，瑞士人，1704年生，1752年死。

他不是第一个发现这个方法的人。

在他之前，日本数学家关孝和、英国数学家麦克劳林都发现过类似的东西。

但是克拉默把它整理得最清楚、最完整，写进了他1750年出版的一本书里。

于是这个方法，就用他的名字命名了。

这在数学史上很常见——发现一件事的人，和让这件事广为人知的人，往往不是同一个人。

\section{先从两个未知数开始}

我们在上一章已经推过了这个过程，但这里再走一遍，因为它是理解克拉默法则的基础。

\[
\begin{cases} ax + by = e \\ cx + dy = f \end{cases}
\]

$x$ 是一个真实的东西，$y$ 是另一个真实的东西。

$a$、$b$、$c$、$d$ 是它们前面的系数。

$e$、$f$ 是等号右边的数。

\subsection{消元，解出 $x$}

第一条方程乘以 $d$：

\[
adx + bdy = ed
\]

第二条方程乘以 $b$：

\[
bcx + bdy = bf
\]

第一条减去第二条：

\[
(ad - bc)x = ed - bf
\]

左边是 $(ad - bc)$，这是系数行列式——

\[
D = \begin{vmatrix} a & b \\ c & d \end{vmatrix} = ad - bc
\]

右边是 $ed - bf$。

仔细看这个 $ed - bf$——

它长得像 $ad - bc$，只不过把第一列的 $a$ 和 $c$，换成了右边那列的 $e$ 和 $f$。

我们叫它 $D_x$：

\[
D_x = \begin{vmatrix} e & b \\ f & d \end{vmatrix} = ed - bf
\]

所以：

\[
D \cdot x = D_x
\]

\[
x = \frac{D_x}{D}
\]

\subsection{消元，解出 $y$}

第一条方程乘以 $c$：

\[
acx + bcy = ec
\]

第二条方程乘以 $a$：

\[
acx + ady = af
\]

第二条减去第一条：

\[
(ad - bc)y = af - ec
\]

左边还是那个 $D = ad - bc$。

右边是 $af - ec$。

这个 $af - ec$ 是把第二列的 $b$ 和 $d$，换成右边那列的 $e$ 和 $f$ 之后得到的：

\[
D_y = \begin{vmatrix} a & e \\ c & f \end{vmatrix} = af - ec
\]

所以：

\[
D \cdot y = D_y
\]

\[
y = \frac{D_y}{D}
\]

\subsection{规律出现了}

\[
x = \frac{D_x}{D}, \qquad y = \frac{D_y}{D}
\]

$D$ 是系数矩阵的行列式，不动。

$D_x$ 是把\textbf{第一列}换成右边那列之后的行列式。

$D_y$ 是把\textbf{第二列}换成右边那列之后的行列式。

哪个未知数在第几列，就换第几列。

这个规律，在三个未知数的时候，完全一样。

\section{现在进入三个未知数}

\[
\begin{cases} a_1 x + b_1 y + c_1 z = d_1 \\ a_2 x + b_2 y + c_2 z = d_2 \\ a_3 x + b_3 y + c_3 z = d_3 \end{cases}
\]

$x$、$y$、$z$ 是三个真实的东西。

$d_1$、$d_2$、$d_3$ 是等号右边的三个数。

系数矩阵是：

\[
\begin{pmatrix} a_1 & b_1 & c_1 \\ a_2 & b_2 & c_2 \\ a_3 & b_3 & c_3 \end{pmatrix}
\]

它的行列式叫 $D$：

\[
D = \begin{vmatrix} a_1 & b_1 & c_1 \\ a_2 & b_2 & c_2 \\ a_3 & b_3 & c_3 \end{vmatrix}
\]

\section{克拉默法则说的是什么}

\textbf{如果 $D \neq 0$，那么：}

\[
x = \frac{D_x}{D}, \qquad y = \frac{D_y}{D}, \qquad z = \frac{D_z}{D}
\]

其中——

$D_x$ 是把系数矩阵的\textbf{第一列}，换成右边那列 $\begin{pmatrix} d_1 \\ d_2 \\ d_3 \end{pmatrix}$ 之后，算出来的行列式。

$D_y$ 是把系数矩阵的\textbf{第二列}，换成右边那列之后，算出来的行列式。

$D_z$ 是把系数矩阵的\textbf{第三列}，换成右边那列之后，算出来的行列式。

\subsection{语言记忆法}

\begin{quote}
想要哪个未知数，就把它住的那一列换掉，然后除以原来的 $D$。
\end{quote}

$x$ 住在第一列，就换第一列，得到 $D_x$，除以 $D$。

$y$ 住在第二列，就换第二列，得到 $D_y$，除以 $D$。

$z$ 住在第三列，就换第三列，得到 $D_z$，除以 $D$。

\subsection{为什么是“换那一列”？}

这个问题，要从消元的结果来理解。

还记得两个未知数的时候，消元之后得到：

\[
(ad - bc)x = ed - bf
\]

右边的 $ed - bf$，正好是把第一列（$x$ 的系数列）换成右边那列之后的行列式。

三个未知数的时候，消元的过程更繁琐，但结果的规律完全一样——

消元到最后，$x$ 的分子，是把 $x$ 所在的那一列（第一列）换成右边那列之后的行列式。

$y$ 的分子，是把 $y$ 所在的那一列（第二列）换成右边那列之后的行列式。

$z$ 的分子，是把 $z$ 所在的那一列（第三列）换成右边那列之后的行列式。

换那一列，是消元过程自然产生的结果，不是人为规定的技巧。

\section{用完整的例子走一遍}

\[
\begin{cases} x + 2y + z = 7 \\ 2x + y + z = 8 \\ x + y + 2z = 9 \end{cases}
\]

\subsection{第一步：写出系数矩阵}

\[
\begin{pmatrix} 1 & 2 & 1 \\ 2 & 1 & 1 \\ 1 & 1 & 2 \end{pmatrix}
\]

右边那列是：$\begin{pmatrix} 7 \\ 8 \\ 9 \end{pmatrix}$

\subsection{第二步：算 $D$}

\[
D = \begin{vmatrix} 1 & 2 & 1 \\ 2 & 1 & 1 \\ 1 & 1 & 2 \end{vmatrix}
\]

按第一行展开。

\textbf{第一行第一个数：$1$，位置 $(1,1)$，$1+1=2$，偶数，正号。}

划掉第一行、第一列，剩：

\[
\begin{pmatrix} 1 & 1 \\ 1 & 2 \end{pmatrix}
\]

这个小行列式：$1 \times 2 - 1 \times 1 = 2 - 1 = 1$

贡献：$+1 \times 1 = +1$

\textbf{第一行第二个数：$2$，位置 $(1,2)$，$1+2=3$，奇数，负号。}

划掉第一行、第二列，剩：

\[
\begin{pmatrix} 2 & 1 \\ 1 & 2 \end{pmatrix}
\]

这个小行列式：$2 \times 2 - 1 \times 1 = 4 - 1 = 3$

贡献：$-2 \times 3 = -6$

\textbf{第一行第三个数：$1$，位置 $(1,3)$，$1+3=4$，偶数，正号。}

划掉第一行、第三列，剩：

\[
\begin{pmatrix} 2 & 1 \\ 1 & 1 \end{pmatrix}
\]

这个小行列式：$2 \times 1 - 1 \times 1 = 2 - 1 = 1$

贡献：$+1 \times 1 = +1$

\textbf{三项加起来：}

\[
D = +1 + (-6) + 1 = -4
\]

$D = -4$，不等于零。

方程组有唯一解，可以继续。

\subsection{第三步：算 $D_x$}

$D_x$ 是把\textbf{第一列}换成右边那列 $[7, 8, 9]$。

原来第一列是 $1$、$2$、$1$，现在换成 $7$、$8$、$9$。

其余两列不动。

\[
D_x = \begin{vmatrix} 7 & 2 & 1 \\ 8 & 1 & 1 \\ 9 & 1 & 2 \end{vmatrix}
\]

按第一行展开。

\textbf{$7$，位置 $(1,1)$，正号。}

划掉第一行、第一列，剩：

\[
\begin{pmatrix} 1 & 1 \\ 1 & 2 \end{pmatrix}
\]

小行列式：$1 \times 2 - 1 \times 1 = 1$

贡献：$+7 \times 1 = +7$

\textbf{$2$，位置 $(1,2)$，负号。}

划掉第一行、第二列，剩：

\[
\begin{pmatrix} 8 & 1 \\ 9 & 2 \end{pmatrix}
\]

小行列式：$8 \times 2 - 1 \times 9 = 16 - 9 = 7$

贡献：$-2 \times 7 = -14$

\textbf{$1$，位置 $(1,3)$，正号。}

划掉第一行、第三列，剩：

\[
\begin{pmatrix} 8 & 1 \\ 9 & 1 \end{pmatrix}
\]

小行列式：$8 \times 1 - 1 \times 9 = 8 - 9 = -1$

贡献：$+1 \times (-1) = -1$

\textbf{三项加起来：}

\[
D_x = +7 + (-14) + (-1) = 7 - 14 - 1 = -8
\]

\textbf{算 $x$：}

\[
x = \frac{D_x}{D} = \frac{-8}{-4} = 2
\]

\subsection{第四步：算 $D_y$}

$D_y$ 是把\textbf{第二列}换成右边那列 $[7, 8, 9]$。

原来第二列是 $2$、$1$、$1$，现在换成 $7$、$8$、$9$。

第一列和第三列不动。

\[
D_y = \begin{vmatrix} 1 & 7 & 1 \\ 2 & 8 & 1 \\ 1 & 9 & 2 \end{vmatrix}
\]

按第一行展开。

\textbf{$1$，位置 $(1,1)$，正号。}

划掉第一行、第一列，剩：

\[
\begin{pmatrix} 8 & 1 \\ 9 & 2 \end{pmatrix}
\]

小行列式：$8 \times 2 - 1 \times 9 = 16 - 9 = 7$

贡献：$+1 \times 7 = +7$

\textbf{$7$，位置 $(1,2)$，负号。}

划掉第一行、第二列，剩：

\[
\begin{pmatrix} 2 & 1 \\ 1 & 2 \end{pmatrix}
\]

小行列式：$2 \times 2 - 1 \times 1 = 4 - 1 = 3$

贡献：$-7 \times 3 = -21$

\textbf{$1$，位置 $(1,3)$，正号。}

划掉第一行、第三列，剩：

\[
\begin{pmatrix} 2 & 8 \\ 1 & 9 \end{pmatrix}
\]

小行列式：$2 \times 9 - 8 \times 1 = 18 - 8 = 10$

贡献：$+1 \times 10 = +10$

\textbf{三项加起来：}

\[
D_y = +7 + (-21) + 10 = 7 - 21 + 10 = -4
\]

\textbf{算 $y$：}

\[
y = \frac{D_y}{D} = \frac{-4}{-4} = 1
\]

\subsection{第五步：算 $D_z$}

$D_z$ 是把\textbf{第三列}换成右边那列 $[7, 8, 9]$。

原来第三列是 $1$、$1$、$2$，现在换成 $7$、$8$、$9$。

第一列和第二列不动。

\[
D_z = \begin{vmatrix} 1 & 2 & 7 \\ 2 & 1 & 8 \\ 1 & 1 & 9 \end{vmatrix}
\]

按第一行展开。

\textbf{$1$，位置 $(1,1)$，正号。}

划掉第一行、第一列，剩：

\[
\begin{pmatrix} 1 & 8 \\ 1 & 9 \end{pmatrix}
\]

小行列式：$1 \times 9 - 8 \times 1 = 9 - 8 = 1$

贡献：$+1 \times 1 = +1$

\textbf{$2$，位置 $(1,2)$，负号。}

划掉第一行、第二列，剩：

\[
\begin{pmatrix} 2 & 8 \\ 1 & 9 \end{pmatrix}
\]

小行列式：$2 \times 9 - 8 \times 1 = 18 - 8 = 10$

贡献：$-2 \times 10 = -20$

\textbf{$7$，位置 $(1,3)$，正号。}

划掉第一行、第三列，剩：

\[
\begin{pmatrix} 2 & 1 \\ 1 & 1 \end{pmatrix}
\]

小行列式：$2 \times 1 - 1 \times 1 = 2 - 1 = 1$

贡献：$+7 \times 1 = +7$

\textbf{三项加起来：}

\[
D_z = +1 + (-20) + 7 = 1 - 20 + 7 = -12
\]

\textbf{算 $z$：}

\[
z = \frac{D_z}{D} = \frac{-12}{-4} = 3
\]

\subsection{第六步：验算}

把 $x = 2$，$y = 1$，$z = 3$ 代回三条原始方程。

\textbf{第一条方程：}

原式：$x + 2y + z = 7$

代入：$2 + 2 \times 1 + 3$

先算 $2 \times 1 = 2$。

再算 $2 + 2 + 3 = 7$。

等于 $7$。✓

\textbf{第二条方程：}

原式：$2x + y + z = 8$

代入：$2 \times 2 + 1 + 3$

先算 $2 \times 2 = 4$。

再算 $4 + 1 + 3 = 8$。

等于 $8$。✓

\textbf{第三条方程：}

原式：$x + y + 2z = 9$

代入：$2 + 1 + 2 \times 3$

先算 $2 \times 3 = 6$。

再算 $2 + 1 + 6 = 9$。

等于 $9$。✓

三条全对。

\begin{quote}
\textbf{答案：$x = 2$，$y = 1$，$z = 3$。}
\end{quote}

\section{再走一遍，换一道新题}

这次题目里的数字稍微大一点，但流程完全一样。

\[
\begin{cases} 2x + y - z = 3 \\ x - y + 2z = 1 \\ 3x + 2y + z = 8 \end{cases}
\]

\subsection{第一步：认清楚系数矩阵和右边那列}

系数矩阵：

\[
\begin{pmatrix} 2 & 1 & -1 \\ 1 & -1 & 2 \\ 3 & 2 & 1 \end{pmatrix}
\]

第一列：$x$ 的系数，$2$、$1$、$3$。

第二列：$y$ 的系数，$1$、$-1$、$2$。

第三列：$z$ 的系数，$-1$、$2$、$1$。

右边那列：$3$、$1$、$8$。

\subsection{第二步：算 $D$}

\[
D = \begin{vmatrix} 2 & 1 & -1 \\ 1 & -1 & 2 \\ 3 & 2 & 1 \end{vmatrix}
\]

按第一行展开。

\textbf{$2$，位置 $(1,1)$，正号。}

划掉第一行、第一列，剩：

\[
\begin{pmatrix} -1 & 2 \\ 2 & 1 \end{pmatrix}
\]

小行列式：$(-1) \times 1 - 2 \times 2 = -1 - 4 = -5$

贡献：$+2 \times (-5) = -10$

\textbf{$1$，位置 $(1,2)$，负号。}

划掉第一行、第二列，剩：

\[
\begin{pmatrix} 1 & 2 \\ 3 & 1 \end{pmatrix}
\]

小行列式：$1 \times 1 - 2 \times 3 = 1 - 6 = -5$

贡献：$-1 \times (-5) = +5$

\textbf{$-1$，位置 $(1,3)$，正号。}

划掉第一行、第三列，剩：

\[
\begin{pmatrix} 1 & -1 \\ 3 & 2 \end{pmatrix}
\]

小行列式：$1 \times 2 - (-1) \times 3 = 2 + 3 = 5$

贡献：$+(-1) \times 5 = -5$

\textbf{三项加起来：}

\[
D = -10 + 5 + (-5) = -10 + 5 - 5 = -10
\]

$D = -10$，不等于零，可以继续。

\subsection{第三步：算 $D_x$}

把第一列换成右边那列 $[3, 1, 8]$。

\[
D_x = \begin{vmatrix} 3 & 1 & -1 \\ 1 & -1 & 2 \\ 8 & 2 & 1 \end{vmatrix}
\]

按第一行展开。

\textbf{$3$，位置 $(1,1)$，正号。}

划掉第一行、第一列，剩：

\[
\begin{pmatrix} -1 & 2 \\ 2 & 1 \end{pmatrix}
\]

小行列式：$(-1) \times 1 - 2 \times 2 = -1 - 4 = -5$

贡献：$+3 \times (-5) = -15$

\textbf{$1$，位置 $(1,2)$，负号。}

划掉第一行、第二列，剩：

\[
\begin{pmatrix} 1 & 2 \\ 8 & 1 \end{pmatrix}
\]

小行列式：$1 \times 1 - 2 \times 8 = 1 - 16 = -15$

贡献：$-1 \times (-15) = +15$

\textbf{$-1$，位置 $(1,3)$，正号。}

划掉第一行、第三列，剩：

\[
\begin{pmatrix} 1 & -1 \\ 8 & 2 \end{pmatrix}
\]

小行列式：$1 \times 2 - (-1) \times 8 = 2 + 8 = 10$

贡献：$+(-1) \times 10 = -10$

\textbf{三项加起来：}

\[
D_x = -15 + 15 + (-10) = -15 + 15 - 10 = -10
\]

\textbf{算 $x$：}

\[
x = \frac{D_x}{D} = \frac{-10}{-10} = 1
\]

\subsection{第四步：算 $D_y$}

把第二列换成右边那列 $[3, 1, 8]$。

\[
D_y = \begin{vmatrix} 2 & 3 & -1 \\ 1 & 1 & 2 \\ 3 & 8 & 1 \end{vmatrix}
\]

按第一行展开。

\textbf{$2$，位置 $(1,1)$，正号。}

划掉第一行、第一列，剩：

\[
\begin{pmatrix} 1 & 2 \\ 8 & 1 \end{pmatrix}
\]

小行列式：$1 \times 1 - 2 \times 8 = 1 - 16 = -15$

贡献：$+2 \times (-15) = -30$

\textbf{$3$，位置 $(1,2)$，负号。}

划掉第一行、第二列，剩：

\[
\begin{pmatrix} 1 & 2 \\ 3 & 1 \end{pmatrix}
\]

小行列式：$1 \times 1 - 2 \times 3 = 1 - 6 = -5$

贡献：$-3 \times (-5) = +15$

\textbf{$-1$，位置 $(1,3)$，正号。}

划掉第一行、第三列，剩：

\[
\begin{pmatrix} 1 & 1 \\ 3 & 8 \end{pmatrix}
\]

小行列式：$1 \times 8 - 1 \times 3 = 8 - 3 = 5$

贡献：$+(-1) \times 5 = -5$

\textbf{三项加起来：}

\[
D_y = -30 + 15 + (-5) = -30 + 15 - 5 = -20
\]

\textbf{算 $y$：}

\[
y = \frac{D_y}{D} = \frac{-20}{-10} = 2
\]

\subsection{第五步：算 $D_z$}

把第三列换成右边那列 $[3, 1, 8]$。

\[
D_z = \begin{vmatrix} 2 & 1 & 3 \\ 1 & -1 & 1 \\ 3 & 2 & 8 \end{vmatrix}
\]

按第一行展开。

\textbf{$2$，位置 $(1,1)$，正号。}

划掉第一行、第一列，剩：

\[
\begin{pmatrix} -1 & 1 \\ 2 & 8 \end{pmatrix}
\]

小行列式：$(-1) \times 8 - 1 \times 2 = -8 - 2 = -10$

贡献：$+2 \times (-10) = -20$

\textbf{$1$，位置 $(1,2)$，负号。}

划掉第一行、第二列，剩：

\[
\begin{pmatrix} 1 & 1 \\ 3 & 8 \end{pmatrix}
\]

小行列式：$1 \times 8 - 1 \times 3 = 8 - 3 = 5$

贡献：$-1 \times 5 = -5$

\textbf{$3$，位置 $(1,3)$，正号。}

划掉第一行、第三列，剩：

\[
\begin{pmatrix} 1 & -1 \\ 3 & 2 \end{pmatrix}
\]

小行列式：$1 \times 2 - (-1) \times 3 = 2 + 3 = 5$

贡献：$+3 \times 5 = +15$

\textbf{三项加起来：}

\[
D_z = -20 + (-5) + 15 = -20 - 5 + 15 = -10
\]

\textbf{算 $z$：}

\[
z = \frac{D_z}{D} = \frac{-10}{-10} = 1
\]

\subsection{第六步：验算}

$x = 1$，$y = 2$，$z = 1$。

\textbf{第一条方程：}

原式：$2x + y - z = 3$

代入：$2 \times 1 + 2 - 1$

先算 $2 \times 1 = 2$。

再算 $2 + 2 - 1 = 3$。

等于 $3$。✓

\textbf{第二条方程：}

原式：$x - y + 2z = 1$

代入：$1 - 2 + 2 \times 1$

先算 $2 \times 1 = 2$。

再算 $1 - 2 + 2 = 1$。

等于 $1$。✓

\textbf{第三条方程：}

原式：$3x + 2y + z = 8$

代入：$3 \times 1 + 2 \times 2 + 1$

先算 $3 \times 1 = 3$，$2 \times 2 = 4$。

再算 $3 + 4 + 1 = 8$。

等于 $8$。✓

三条全对。

\begin{quote}
\textbf{答案：$x = 1$，$y = 2$，$z = 1$。}
\end{quote}

\section{如果 $D = 0$ 怎么办？}

克拉默法则有一个前提条件：$D \neq 0$。

如果算出来 $D = 0$，那克拉默法则就用不了。

为什么？

因为克拉默法则的公式是——

\[
x = \frac{D_x}{D}
\]

$D$ 在分母上。

分母是零，这个式子没有意义。

$D = 0$ 意味着两种情况——

\textbf{情况一：方程组无解。}

比如：

\[
\begin{cases} x + y = 1 \\ x + y = 2 \end{cases}
\]

这两条方程互相矛盾。$x + y$ 不可能同时等于 $1$ 又等于 $2$。

系数矩阵：

\[
\begin{pmatrix} 1 & 1 \\ 1 & 1 \end{pmatrix}
\]

行列式：$1 \times 1 - 1 \times 1 = 0$。

无解。

\textbf{情况二：方程组有无数个解。}

比如：

\[
\begin{cases} x + y = 1 \\ 2x + 2y = 2 \end{cases}
\]

第二条方程是第一条方程乘以2，完全一样。

任何满足 $x + y = 1$ 的 $x$ 和 $y$ 都是解。

有无数个。

系数矩阵：

\[
\begin{pmatrix} 1 & 1 \\ 2 & 2 \end{pmatrix}
\]

行列式：$1 \times 2 - 1 \times 2 = 0$。

\textbf{所以：}

\begin{quote}
$D = 0$，方程组要么无解，要么有无数个解。

这两种情况，克拉默法则都处理不了。

需要用别的方法（比如消元法）来进一步判断。
\end{quote}

\section{克拉默法则的完整流程}

\textbf{第一步：写出系数矩阵，算出 $D$。}

如果 $D = 0$，停下来，克拉默法则不适用。

如果 $D \neq 0$，继续。

\textbf{第二步：构造 $D_x$、$D_y$、$D_z$。}

$D_x$：把系数矩阵的第一列，换成右边那列，算行列式。

$D_y$：把系数矩阵的第二列，换成右边那列，算行列式。

$D_z$：把系数矩阵的第三列，换成右边那列，算行列式。

\textbf{第三步：除一下，得到答案。}

\[
x = \frac{D_x}{D}, \qquad y = \frac{D_y}{D}, \qquad z = \frac{D_z}{D}
\]

\textbf{第四步：验算。}

把答案代回三条原始方程，逐条检查。

全对，答案可靠。

有一条不对，说明某处算错了，回去检查。

\section{本章小结}

克拉默法则不是一个神秘的公式。

它是消元法走到底之后，整理出来的结论。

消元法每次都要重新算，步骤繁琐，容易出错。

克拉默法则把这些步骤提炼成了四个行列式——

算 $D$，算 $D_x$，算 $D_y$，算 $D_z$，除一下，答案就出来了。

流程固定，不需要每次重新想怎么消元。

这就是数学一直在做的事：

\textbf{把聪明人解决问题的方法，整理成任何人都能照着做的步骤。}

\chapter{线性——这个词到底是什么意思}

\section{先问一个问题}

你有没有想过——

这门学科，为什么叫“线性代数”？

“代数”好理解，就是用字母代替数字。

但“线性”是什么意思？

线？什么线？

\section{“线性”这两个字，藏着一个非常具体的意思}

从最简单的情况开始。

初中就见过这个方程：

\[
y = 2x
\]

$x$ 是你走的路程，$y$ 是你花的时间。

或者 $x$ 是你买的苹果数量，$y$ 是你花的钱。

随便什么，总之 $x$ 变化，$y$ 跟着变化。

把所有的 $(x, y)$ 画在坐标系里，得到的是一条直线。

这条直线，就是“线性”里那个“线”。

\textbf{如果一个关系，画出来是一条直线，这个关系就叫线性关系。}

\section{什么样的方程，画出来是直线？}

这些方程——

\[
y = 2x
\]

\[
y = 3x + 1
\]

\[
y = -x + 5
\]

画出来都是直线。

它们有一个共同的特点：

$x$ 只是 $x$，没有 $x^2$，没有 $x^3$，没有 $\sqrt{x}$，没有 $\frac{1}{x}$。

$x$ 的最高次数是1。

再看这些方程——

\[
y = x^2
\]

\[
y = x^3 - 2
\]

\[
y = \sqrt{x}
\]

这些方程，画出来是曲线，不是直线。

它们不是线性的。

\textbf{所以“线性”的意思，用大白话说就是：}

\begin{quote}
变量只以一次方的形式出现，画出来是直线（或者直线的推广）。
\end{quote}

\section{那“线性代数”是什么？}

线性代数，研究的是线性关系。

具体来说，研究的是——

多个变量之间的线性关系，也就是多个方程、多个未知数的情况。

比如：

\[
\begin{cases} 2x + 3y = 7 \\ x - y = 1 \end{cases}
\]

这个方程组里，$x$ 和 $y$ 都只以一次方出现。

这就是线性方程组。

线性代数，就是研究这类方程组的数学工具。

矩阵、行列式、克拉默法则——这些已经学过的东西，都是线性代数的一部分。

\section{为什么“线性”那么重要？}

你可能在想：现实世界里，很多关系都不是线性的。

温度和气压的关系，不是直线。

人口增长，不是直线。

股票价格，更不是直线。

那为什么数学家要花那么大力气研究线性这件事？

原因有两个。

\textbf{第一个原因：线性关系是最简单的关系。}

如果你想研究一个复杂的系统，第一步往往是把它当成线性的来处理。

就像你画地图，先画出主要的道路，再慢慢补充细节。

线性是那个“先画出来的主要框架”。

\textbf{第二个原因：很多复杂的关系，在局部范围内，可以近似成线性的。}

比如，一条曲线，如果你把它放大很多倍，局部看起来就像一条直线。

微积分里有一个概念叫“切线”，就是用直线来近似曲线在某一点附近的行为。

工程师、物理学家、经济学家，经常用线性来近似复杂的现实，然后用线性代数来解决问题。

\section{现在，进入这一章真正要做的事}

这一章的核心任务，是\textbf{翻译}。

把复杂的文字题，翻译成线性方程组，然后用已经学过的工具解出来。

这比之前的应用题更难翻译——

不是因为数学变复杂了，而是因为题目的说法更绕，藏得更深。

\section{翻译的基本原则}

\subsection{原则一：先找“未知数”，再找“关系”}

每一道应用题，都在说两件事——

第一件事：我不知道什么？（这是未知数）

第二件事：我知道什么？（这是方程的来源）

翻译的第一步，永远是找清楚“我不知道什么”。

\subsection{原则二：每一句话，对应一条方程}

题目里有多少条有用的信息，就能列多少条方程。

三个未知数，需要三条方程。

如果找不到三条独立的信息，方程组就解不出唯一答案。

\subsection{原则三：把每一句中文，逐字翻译成数学语言}

不要整体去想，而是逐字翻译。

“甲的三倍”——就是 $3 \times$ 甲，就是 $3x$。

“乙比甲多5”——就是 $y = x + 5$，或者 $y - x = 5$。

“三者之和”——就是 $x + y + z$。

一个词一个词地翻，不要跳步。

\section{开始翻译：从简单到复杂}

\subsection{第一道题：分配问题}

\begin{quote}
一个工厂，一天生产了480个零件，分配给甲、乙、丙三个车间。  
已知甲车间分到的零件，是乙车间的2倍。  
丙车间分到的零件，比乙车间少60个。  
问：三个车间各分到多少个零件？
\end{quote}

\textbf{第一步：找未知数。}

我不知道什么？

我不知道甲分到多少个，乙分到多少个，丙分到多少个。

所以——

$x$ = 甲车间分到的零件数（个）

$y$ = 乙车间分到的零件数（个）

$z$ = 丙车间分到的零件数（个）

$x$、$y$、$z$ 都是真实的零件数量，不是抽象的字母。

\textbf{第二步：逐句翻译。}

\textbf{“一天生产了480个零件，分配给甲、乙、丙三个车间。”}

三个车间加在一起，总共480个。

翻译：$x + y + z = 480$

这是方程①。

\textbf{“甲车间分到的零件，是乙车间的2倍。”}

甲的数量，是乙的2倍。

甲就是 $x$，乙就是 $y$，2倍就是乘以2。

翻译：$x = 2y$

这是方程②。

\textbf{“丙车间分到的零件，比乙车间少60个。”}

丙比乙少60个。

丙就是 $z$，乙就是 $y$，少60个就是减去60。

翻译：$z = y - 60$

这是方程③。

\textbf{三条方程：}

\[
\begin{cases} x + y + z = 480 & \cdots① \\ x = 2y & \cdots② \\ z = y - 60 & \cdots③ \end{cases}
\]

\textbf{第三步：把②和③代入①。}

②告诉我们 $x = 2y$，③告诉我们 $z = y - 60$。

把这两个代入①：

\[
(2y) + y + (y - 60) = 480
\]

左边展开：

\[
2y + y + y - 60 = 480
\]

\[
4y - 60 = 480
\]

两边加60：

\[
4y = 540
\]

两边除以4：

\[
y = 135
\]

乙分到135个。

\[
x = 2y = 2 \times 135 = 270
\]

甲分到270个。

\[
z = y - 60 = 135 - 60 = 75
\]

丙分到75个。

\textbf{验算：}

$x + y + z = 270 + 135 + 75 = 480$ ✓

甲是乙的2倍：$270 = 2 \times 135$ ✓

丙比乙少60：$135 - 75 = 60$ ✓

\begin{quote}
\textbf{甲车间 270个，乙车间 135个，丙车间 75个。}
\end{quote}

\subsection{第二道题：速度与时间}

\begin{quote}
甲、乙、丙三人从同一地点出发，去同一个目的地。  
甲步行2小时到达，乙骑自行车1小时到达，丙开摩托车半小时到达。  
已知乙的速度比甲每小时多4千米。  
求三人各自的速度。
\end{quote}

\textbf{第一步：找未知数。}

我不知道什么？

甲的速度、乙的速度、丙的速度。

$x$ = 甲的速度（千米/小时）

$y$ = 乙的速度（千米/小时）

$z$ = 丙的速度（千米/小时）

\textbf{第二步：逐句翻译。}

\textbf{“三人从同一地点去同一目的地，甲走了2小时，乙走了1小时，丙走了半小时。”}

这句话说的是：三个人走的路程，是同一段距离。

路程 = 速度 × 时间。

（还记得这个公式吗？路程等于速度乘以时间。速度快，时间短；速度慢，时间长。）

甲的路程：$x \times 2 = 2x$

乙的路程：$y \times 1 = y$

丙的路程：$z \times 0.5 = 0.5z$

三人走的是同一段距离，所以——

$2x = y = 0.5z$

这里有两个等号，可以拆成两条方程：

$2x = y$ ，这是方程①。

$y = 0.5z$ ，这是方程②。

\textbf{“乙的速度比甲每小时多4千米。”}

乙比甲多4。

\[
y = x + 4 \quad \cdots③
\]

\textbf{三条方程：}

\[
\begin{cases} 2x = y & \cdots① \\ y = 0.5z & \cdots② \\ y = x + 4 & \cdots③ \end{cases}
\]

\textbf{代入求解。}

①告诉我们 $y = 2x$，代入③：

\[
2x = x + 4
\]

\[
2x - x = 4
\]

\[
x = 4
\]

甲的速度是4千米/小时。

\[
y = 2x = 2 \times 4 = 8
\]

乙的速度是8千米/小时。

由②得 $z = 2y$：

\[
z = 2 \times 8 = 16
\]

丙的速度是16千米/小时。

\textbf{验算：}

甲走2小时：$4 \times 2 = 8$ 千米。

乙走1小时：$8 \times 1 = 8$ 千米。

丙走半小时：$16 \times 0.5 = 8$ 千米。

三人走了同样的8千米。✓

乙比甲多4：$8 - 4 = 4$ ✓

\begin{quote}
\textbf{甲 4千米/小时，乙 8千米/小时，丙 16千米/小时。}
\end{quote}

\subsection{第三道题：这才是真正复杂的翻译}

\begin{quote}
某工厂生产A、B、C三种产品。  
生产一件A产品，需要原料甲2千克，原料乙1千克，原料丙3千克。  
生产一件B产品，需要原料甲1千克，原料乙3千克，原料丙2千克。  
生产一件C产品，需要原料甲3千克，原料乙2千克，原料丙1千克。  
工厂今天总共用了原料甲10千克，原料乙13千克，原料丙13千克。  
问：今天各生产了多少件A、B、C产品？
\end{quote}

\textbf{第一步：找未知数。}

我不知道什么？

A生产了多少件，B生产了多少件，C生产了多少件。

$x$ = 今天生产A产品的件数

$y$ = 今天生产B产品的件数

$z$ = 今天生产C产品的件数

$x$、$y$、$z$ 是件数，是真实的数量。

\textbf{第二步：理清楚题目的结构。}

这道题有一个隐藏的逻辑，需要先看清楚——

题目说的，不是三件事，而是\textbf{三种原料各自被用了多少}。

原料甲被用了10千克。

原料乙被用了13千克。

原料丙被用了13千克。

这三件事，各自是一条方程。

\textbf{第三步：逐句翻译。}

\textbf{先想清楚：原料甲总共被用了多少？}

生产A产品，每件用原料甲2千克。生产了 $x$ 件，用了 $2x$ 千克。

生产B产品，每件用原料甲1千克。生产了 $y$ 件，用了 $1y = y$ 千克。

生产C产品，每件用原料甲3千克。生产了 $z$ 件，用了 $3z$ 千克。

原料甲总共用了：$2x + y + 3z$

题目说原料甲用了10千克，所以——

\[
2x + y + 3z = 10 \quad \cdots①
\]

\textbf{原料乙总共被用了多少？}

生产A产品，每件用原料乙1千克，生产了 $x$ 件，用了 $x$ 千克。

生产B产品，每件用原料乙3千克，生产了 $y$ 件，用了 $3y$ 千克。

生产C产品，每件用原料乙2千克，生产了 $z$ 件，用了 $2z$ 千克。

原料乙总共用了：$x + 3y + 2z$

题目说原料乙用了13千克，所以——

\[
x + 3y + 2z = 13 \quad \cdots②
\]

\textbf{原料丙总共被用了多少？}

生产A产品，每件用原料丙3千克，生产了 $x$ 件，用了 $3x$ 千克。

生产B产品，每件用原料丙2千克，生产了 $y$ 件，用了 $2y$ 千克。

生产C产品，每件用原料丙1千克，生产了 $z$ 件，用了 $z$ 千克。

原料丙总共用了：$3x + 2y + z$

题目说原料丙用了13千克，所以——

\[
3x + 2y + z = 13 \quad \cdots③
\]

\textbf{三条方程整理：}

\[
\begin{cases} 2x + y + 3z = 10 & \cdots① \\ x + 3y + 2z = 13 & \cdots② \\ 3x + 2y + z = 13 & \cdots③ \end{cases}
\]

翻译完毕。

\textbf{第四步：算 $D$。}

\[
D = \begin{vmatrix} 2 & 1 & 3 \\ 1 & 3 & 2 \\ 3 & 2 & 1 \end{vmatrix}
\]

按第一行展开。

\textbf{$2$，位置 $(1,1)$，正号。}

划掉第一行、第一列，剩：

\[
\begin{pmatrix} 3 & 2 \\ 2 & 1 \end{pmatrix}
\]

小行列式：$3 \times 1 - 2 \times 2 = 3 - 4 = -1$

贡献：$+2 \times (-1) = -2$

\textbf{$1$，位置 $(1,2)$，负号。}

划掉第一行、第二列，剩：

\[
\begin{pmatrix} 1 & 2 \\ 3 & 1 \end{pmatrix}
\]

小行列式：$1 \times 1 - 2 \times 3 = 1 - 6 = -5$

贡献：$-1 \times (-5) = +5$

\textbf{$3$，位置 $(1,3)$，正号。}

划掉第一行、第三列，剩：

\[
\begin{pmatrix} 1 & 3 \\ 3 & 2 \end{pmatrix}
\]

小行列式：$1 \times 2 - 3 \times 3 = 2 - 9 = -7$

贡献：$+3 \times (-7) = -21$

\textbf{加起来：}

\[
D = -2 + 5 + (-21) = -2 + 5 - 21 = -18
\]

$D = -18$，不等于零，可以继续。

\textbf{第五步：算 $D_x$。}

把第一列换成右边那列 $[10, 13, 13]$。

\[
D_x = \begin{vmatrix} 10 & 1 & 3 \\ 13 & 3 & 2 \\ 13 & 2 & 1 \end{vmatrix}
\]

按第一行展开。

\textbf{$10$，正号。}

剩：

\[
\begin{pmatrix} 3 & 2 \\ 2 & 1 \end{pmatrix}
\]

小行列式：$3 - 4 = -1$

贡献：$+10 \times (-1) = -10$

\textbf{$1$，负号。}

剩：

\[
\begin{pmatrix} 13 & 2 \\ 13 & 1 \end{pmatrix}
\]

小行列式：$13 \times 1 - 2 \times 13 = 13 - 26 = -13$

贡献：$-1 \times (-13) = +13$

\textbf{$3$，正号。}

剩：

\[
\begin{pmatrix} 13 & 3 \\ 13 & 2 \end{pmatrix}
\]

小行列式：$13 \times 2 - 3 \times 13 = 26 - 39 = -13$

贡献：$+3 \times (-13) = -39$

\textbf{加起来：}

\[
D_x = -10 + 13 + (-39) = -10 + 13 - 39 = -36
\]

\textbf{算 $x$：}

\[
x = \frac{D_x}{D} = \frac{-36}{-18} = 2
\]

A产品，生产了2件。

\textbf{第六步：算 $D_y$。}

把第二列换成 $[10, 13, 13]$。

\[
D_y = \begin{vmatrix} 2 & 10 & 3 \\ 1 & 13 & 2 \\ 3 & 13 & 1 \end{vmatrix}
\]

\textbf{$2$，正号。}

剩：

\[
\begin{pmatrix} 13 & 2 \\ 13 & 1 \end{pmatrix}
\]

小行列式：$13 \times 1 - 2 \times 13 = 13 - 26 = -13$

贡献：$+2 \times (-13) = -26$

\textbf{$10$，负号。}

剩：

\[
\begin{pmatrix} 1 & 2 \\ 3 & 1 \end{pmatrix}
\]

小行列式：$1 \times 1 - 2 \times 3 = 1 - 6 = -5$

贡献：$-10 \times (-5) = +50$

\textbf{$3$，正号。}

剩：

\[
\begin{pmatrix} 1 & 13 \\ 3 & 13 \end{pmatrix}
\]

小行列式：$1 \times 13 - 13 \times 3 = 13 - 39 = -26$

贡献：$+3 \times (-26) = -78$

\textbf{加起来：}

\[
D_y = -26 + 50 + (-78) = -26 + 50 - 78 = -54
\]

\textbf{算 $y$：}

\[
y = \frac{D_y}{D} = \frac{-54}{-18} = 3
\]

B产品，生产了3件。

\textbf{第七步：算 $D_z$。}

把第三列换成 $[10, 13, 13]$。

\[
D_z = \begin{vmatrix} 2 & 1 & 10 \\ 1 & 3 & 13 \\ 3 & 2 & 13 \end{vmatrix}
\]

\textbf{$2$，正号。}

剩：

\[
\begin{pmatrix} 3 & 13 \\ 2 & 13 \end{pmatrix}
\]

小行列式：$3 \times 13 - 13 \times 2 = 39 - 26 = 13$

贡献：$+2 \times 13 = +26$

\textbf{$1$，负号。}

剩：

\[
\begin{pmatrix} 1 & 13 \\ 3 & 13 \end{pmatrix}
\]

小行列式：$1 \times 13 - 13 \times 3 = 13 - 39 = -26$

贡献：$-1 \times (-26) = +26$

\textbf{$10$，正号。}

剩：

\[
\begin{pmatrix} 1 & 3 \\ 3 & 2 \end{pmatrix}
\]

小行列式：$1 \times 2 - 3 \times 3 = 2 - 9 = -7$

贡献：$+10 \times (-7) = -70$

\textbf{加起来：}

\[
D_z = 26 + 26 + (-70) = 26 + 26 - 70 = -18
\]

\textbf{算 $z$：}

\[
z = \frac{D_z}{D} = \frac{-18}{-18} = 1
\]

C产品，生产了1件。

\textbf{验算：}

$x = 2$，$y = 3$，$z = 1$。

\textbf{原料甲：} $2 \times 2 + 3 \times 1 + 1 \times 3 = 4 + 3 + 3 = 10$ ✓

\textbf{原料乙：} $2 \times 1 + 3 \times 3 + 1 \times 2 = 2 + 9 + 2 = 13$ ✓

\textbf{原料丙：} $2 \times 3 + 3 \times 2 + 1 \times 1 = 6 + 6 + 1 = 13$ ✓

三条全对。

\begin{quote}
\textbf{A产品 2件，B产品 3件，C产品 1件。}
\end{quote}

\section{翻译的核心，从这三道题里提炼出来}

\textbf{第一件事：每一个未知数，都是一个真实的东西。}

不是抽象的 $x$，是零件数量，是速度，是件数。

设好之后，在整个解题过程中，始终记住它代表什么。

\textbf{第二件事：每一句话，说的是一种“总量等于什么”。}

三个车间加起来等于480。

原料甲用掉的总量等于10千克。

三个人走的路程相等。

找到“总量”，方程就出来了。

\textbf{第三件事：发现矛盾，不要硬算，先检查题目。}

如果翻译出来的方程组前后矛盾，不是你的问题，是题目的问题。

数学不会说谎，矛盾的地方一定有原因。

\section{本章小结}

“线性”，就是“直线的”。

变量只以一次方出现，方程画出来是直线。

“线性方程组”，就是多个线性方程放在一起。

“线性代数”，就是研究这类问题的数学工具。

而翻译——把文字变成方程——是使用这些工具之前，必须做的第一步。

翻译做对了，后面的计算就是体力活。

翻译做错了，算得再仔细也没用。

\chapter{向量——会走的量}

\section{从一件你每天都在做的事情开始}

你每天上学，从家走到学校。

这个“从家走到学校”，数学上怎么描述？

你可能会说：“学校离我家500米。”

但这不够。

因为如果你家在学校的东边，你要往\textbf{西}走500米。

如果你家在学校的南边，你要往\textbf{北}走500米。

同样是500米，走的方向不一样。

\textbf{“500米”这三个字，只告诉了你距离，没有告诉你方向。}

如果我只说“走500米”，你不知道往哪走。

所以，描述一个“从哪里走到哪里”的过程，需要两样东西——

\textbf{第一样：走多远。}（这叫大小）

\textbf{第二样：往哪走。}（这叫方向）

把这两样东西打包在一起，就是\textbf{向量}。

\section{向量是什么？}

向量，就是一个\textbf{既有大小、又有方向}的量。

用大白话说——

\begin{quote}
\textbf{向量，是一个“会走的量”。它知道往哪走，也知道走多远。}
\end{quote}

\section{什么不是向量？}

有些量，只有大小，没有方向。

比如——

\textbf{温度：} 今天30摄氏度。

30度就是30度，没有“往东30度”或者“往西30度”这种说法。

温度没有方向，它不是向量。

\textbf{质量：} 这袋米重50千克。

50千克就是50千克，没有“向上50千克”或者“向下50千克”这种说法。

质量没有方向，它不是向量。

\textbf{时间：} 等了3个小时。

3小时就是3小时，没有方向。

时间不是向量。

这些只有大小、没有方向的量，叫做\textbf{标量}。

“标”的意思是“标记”，只标记一个数字就够了。

\section{什么是向量？}

有些量，必须同时说大小和方向，才能说清楚。

比如——

\textbf{位移：} 向东走了500米。

“向东”是方向，“500米”是大小。

位移是向量。

\textbf{速度：} 向北以每小时60公里的速度行驶。

“向北”是方向，“60公里每小时”是大小。

速度是向量。

\textbf{力：} 用10牛顿的力向右推。

“向右”是方向，“10牛顿”是大小。

力是向量。

\section{向量怎么画？}

向量用一个\textbf{带箭头的线段}来表示。

为什么用箭头？

因为箭头天然就有方向——它指向哪里，向量就朝向哪里。

\textbf{箭头的方向}，表示向量的方向。

\textbf{线段的长度}，表示向量的大小。

比如，“向东走500米”这个向量，怎么画？

第一步：在纸上定一个起点。

第二步：从起点往右画一条线段（我们规定右边是东）。

第三步：在线段的右端，画一个箭头。

第四步：线段的长度，按比例代表500米（比如1厘米代表100米，就画5厘米长）。

这就是这个向量的图像。

\section{向量怎么写？}

向量用字母表示，有几种写法。

\textbf{第一种写法：用两个大写字母，上面加箭头。}

\[
\overrightarrow{AB}
\]

这个符号读作“向量AB”。

$A$ 是起点，$B$ 是终点。

箭头从 $A$ 指向 $B$。

为什么要用两个字母？

因为向量是“从一个地方走到另一个地方”。

$A$ 是你出发的地方，$B$ 是你到达的地方。

$\overrightarrow{AB}$ 就是“从 $A$ 走到 $B$”。

\textbf{第二种写法：用一个小写字母，上面加箭头。}

\[
\vec{a}
\]

这个符号读作“向量a”。

这种写法不说明起点和终点在哪里，只是给这个向量起了一个名字，叫 $a$。

\textbf{第三种写法：用粗体的小写字母。}

\[
\boldsymbol{a}
\]

这在印刷的书本里常见。

手写的时候，一般用 $\vec{a}$ 这种带箭头的写法。

\section{两个特殊的向量}

\subsection{第一个特殊向量：零向量}

如果一个向量的大小是0，它叫\textbf{零向量}。

写作 $\vec{0}$。

零向量是什么意思？

就是“原地不动”。

从 $A$ 点出发，又回到 $A$ 点，哪儿也没去。

位移是0，这就是零向量。

零向量的方向是什么？

\textbf{没有规定。}

因为你哪儿也没走，谈不上往哪个方向走。

或者说，往任何方向走了0米，结果都一样。

\subsection{第二个特殊向量：单位向量}

如果一个向量的模等于1，它叫\textbf{单位向量}。

单位向量有什么用？

它\textbf{只表示方向，不管大小}。

就像一个只告诉你“往哪走”，但不告诉你“走多远”的指路牌。

任何一个非零向量，都可以变成单位向量。

怎么变？

把这个向量除以它的模。

方向不变，大小变成1。

这个操作叫做\textbf{单位化}或者\textbf{归一化}。

\section{两个向量相等，是什么意思？}

这个问题，你可能觉得很简单，但它藏着一个重要的道理。

两个向量相等，需要满足两个条件——

\textbf{第一：方向相同。}

\textbf{第二：大小相同。}

注意——

\textbf{起点在哪里，不重要。}

只要方向相同、大小相同，不管起点在哪里，这两个向量就相等。

\subsection{用一个例子来感受}

小明从家往东走了500米，到达超市。

小红从学校往东走了500米，到达公园。

小明走的向量，和小红走的向量，相等吗？

它们的方向相同吗？都是向东。相同。✓

它们的大小相同吗？都是500米。相同。✓

所以，这两个向量\textbf{相等}。

虽然小明的起点是家，小红的起点是学校，起点不同。

但向量只关心“走了多远、往哪走”，不关心“从哪里出发”。

\subsection{这个性质意味着什么？}

它意味着——

\textbf{你可以把一个向量平移到任何位置。}

只要方向不变、大小不变，它就还是同一个向量。

这叫做向量的\textbf{自由性}。

\section{现在，进入核心问题：怎么用数字来表示向量？}

前面说，向量有方向、有大小。

但“向东”、“向北”这样的说法，不够精确。

我们需要一种更精确的方式来描述向量。

这就是\textbf{向量的坐标表示}。

\subsection{先从一个简单的情况开始}

你站在一个十字路口。

东西方向是一条路，南北方向是另一条路。

我告诉你：

“向东走3步，再向北走4步。”

你能到达一个确定的位置吗？

能。

这个位置，相对于你的出发点，可以用两个数来描述——

\textbf{往东：3步。}

\textbf{往北：4步。}

我们把“往东走多少”叫做 $x$ 方向。

把“往北走多少”叫做 $y$ 方向。

那么，“往东3步、往北4步”这个向量，就可以写成——

\[
(3, 4)
\]

这两个数，叫做向量的\textbf{坐标}。

第一个数，是 $x$ 方向（东西方向）走了多少。

第二个数，是 $y$ 方向（南北方向）走了多少。

\subsection{为什么可以用两个数来表示一个向量？}

因为任何一个平面上的向量，都可以拆成两部分——

\textbf{在 $x$ 方向（横向）走了多少。}

\textbf{在 $y$ 方向（纵向）走了多少。}

这两个数，合在一起，就完全确定了这个向量。

就像你告诉别人怎么走路——

“先往东走3步，再往北走4步。”

这句话完全确定了你最终走到哪里。

所以 $(3, 4)$ 这两个数，完全确定了这个向量。

\subsection{如果往西走呢？往南走呢？}

往东是正的，往西就是负的。

往北是正的，往南就是负的。

“往西走2步，往北走5步”——

$x$ 方向：$-2$（因为是往西，和往东相反）

$y$ 方向：$5$

向量：$(-2, 5)$

“往东走1步，往南走3步”——

$x$ 方向：$1$

$y$ 方向：$-3$（因为是往南，和往北相反）

向量：$(1, -3)$

“往西走4步，往南走2步”——

$x$ 方向：$-4$

$y$ 方向：$-2$

向量：$(-4, -2)$

\subsection{语言记忆法}

\begin{quote}
向量的坐标，就是它\textbf{横着走了多少，竖着走了多少}。

往右（东）是正，往左（西）是负。

往上（北）是正，往下（南）是负。
\end{quote}

\section{向量 $\overrightarrow{AB}$ 是什么？}

\subsection{先从两个点开始}

平面上有两个点。

一个点叫 $A$，它的位置是 $(1, 2)$。

另一个点叫 $B$，它的位置是 $(4, 6)$。

$(1, 2)$ 的意思是：$A$ 点在横轴方向（往右）数1格，在纵轴方向（往上）数2格的地方。

$(4, 6)$ 的意思是：$B$ 点在横轴方向数4格，在纵轴方向数6格的地方。

\subsection{现在问一个问题：从 $A$ 走到 $B$，怎么走？}

你站在 $A$ 点。

你要走到 $B$ 点。

怎么走？

$A$ 的横坐标是1，$B$ 的横坐标是4。

所以横向要从1走到4，要往右走 $4 - 1 = 3$ 格。

$A$ 的纵坐标是2，$B$ 的纵坐标是6。

所以纵向要从2走到6，要往上走 $6 - 2 = 4$ 格。

“从 $A$ 走到 $B$”这件事，可以描述成：

\begin{quote}
往右走3格，往上走4格。
\end{quote}

\subsection{这个“从 $A$ 走到 $B$”的过程，就是向量 $\overrightarrow{AB}$}

$\overrightarrow{AB}$ 读作“向量AB”。

$A$ 是起点（你出发的地方）。

$B$ 是终点（你到达的地方）。

上面那个箭头，表示方向是从 $A$ 指向 $B$。

\subsection{向量 $\overrightarrow{AB}$ 的坐标是什么？}

刚才算出来了：

往右走3格，往上走4格。

所以：

\[
\overrightarrow{AB} = (3, 4)
\]

第一个数3，是横向走了多少。

第二个数4，是纵向走了多少。

\subsection{公式：怎么从两个点的坐标，算出向量的坐标}

起点 $A$ 的坐标是 $(x_A, y_A)$。

终点 $B$ 的坐标是 $(x_B, y_B)$。

向量 $\overrightarrow{AB}$ 的坐标是：

\[
\overrightarrow{AB} = (x_B - x_A, y_B - y_A)
\]

\textbf{终点的坐标，减去起点的坐标。}

\subsection{为什么是“终点减起点”，不是“起点减终点”？}

因为向量描述的是“从起点走到终点”。

你要知道往右走了多少，就看终点的横坐标比起点的横坐标\textbf{多了多少}。

“多了多少”，就是终点减起点。

如果终点的横坐标比起点小，减出来就是负数，意思是往左走了。

\subsection{注意：$\overrightarrow{AB}$ 和 $\overrightarrow{BA}$ 不一样}

$\overrightarrow{AB}$ 是从 $A$ 走到 $B$。

$\overrightarrow{BA}$ 是从 $B$ 走到 $A$。

方向完全相反。

用坐标算一下：

$A = (1, 2)$，$B = (4, 6)$。

\[
\overrightarrow{AB} = (4 - 1, 6 - 2) = (3, 4)
\]

\[
\overrightarrow{BA} = (1 - 4, 2 - 6) = (-3, -4)
\]

$\overrightarrow{BA}$ 正好是 $\overrightarrow{AB}$ 每个坐标都变成相反数。

方向相反，大小相同。

\section{向量的模——实际走了多远}

\subsection{什么是“模”？}

向量的大小，有一个专门的名字，叫\textbf{模}。

“模”这个字，在这里就是“大小”的意思，就是“实际走了多远”。

$\overrightarrow{AB}$ 的模，写作——

\[
|\overrightarrow{AB}|
\]

两边加竖线，表示“取大小”。

$\vec{a}$ 的模，写作——

\[
|\vec{a}|
\]

\subsection{模怎么算？}

向量 $\overrightarrow{AB} = (3, 4)$。

往右走3格，往上走4格。

但这不是你实际走的距离。

你实际走的，是一条\textbf{斜着的直线}，从 $A$ 直接到 $B$。

这条斜线有多长？

\subsection{这里要用到勾股定理}

（还记得勾股定理吗？如果忘了，简单复习一下。）

勾股定理说的是：

在一个\textbf{直角三角形}里，两条直角边的长度分别是 $a$ 和 $b$，斜边的长度是 $c$，那么——

\[
a^2 + b^2 = c^2
\]

两条直角边的平方加起来，等于斜边的平方。

反过来，如果你知道两条直角边的长度，想求斜边的长度，就是：

\[
c = \sqrt{a^2 + b^2}
\]

\subsection{把向量 $\overrightarrow{AB}$ 画出来，就是一个直角三角形}

你从 $A$ 出发。

先往右走3格（这是横向的一条边）。

再往上走4格（这是纵向的一条边）。

最后到达 $B$。

如果你直接从 $A$ 连一条直线到 $B$，这条直线就是斜边。

横向那条边，长度是3。

纵向那条边，长度是4。

这两条边互相垂直（一条是横的，一条是竖的），形成一个直角。

所以这是一个直角三角形。

\subsection{用勾股定理算斜边的长度}

两条直角边：3和4。

斜边的长度：

\[
c = \sqrt{3^2 + 4^2}
\]

先算 $3^2 = 9$。

再算 $4^2 = 16$。

然后 $9 + 16 = 25$。

最后 $\sqrt{25} = 5$。

斜边的长度是5。

\subsection{这个斜边的长度，就是向量 $\overrightarrow{AB}$ 的模}

\[
|\overrightarrow{AB}| = 5
\]

\subsection{模的计算公式}

如果向量 $\vec{a} = (a_1, a_2)$，它的模是：

\[
|\vec{a}| = \sqrt{a_1^2 + a_2^2}
\]

横向走的平方，加上纵向走的平方，再开根号。

这就是勾股定理。

\subsection{练一练}

\textbf{向量 $(1, 1)$ 的模：}

\[
|\vec{a}| = \sqrt{1^2 + 1^2} = \sqrt{1 + 1} = \sqrt{2}
\]

$\sqrt{2}$ 大约是 $1.414$。

\textbf{向量 $(5, 12)$ 的模：}

\[
|\vec{a}| = \sqrt{5^2 + 12^2} = \sqrt{25 + 144} = \sqrt{169} = 13
\]

（$5, 12, 13$ 是一组勾股数，像 $3, 4, 5$ 一样。）

\textbf{向量 $(-3, 4)$ 的模：}

\[
|\vec{a}| = \sqrt{(-3)^2 + 4^2} = \sqrt{9 + 16} = \sqrt{25} = 5
\]

注意：$(-3)^2 = 9$，负数的平方是正数。

所以向量 $(3, 4)$ 和 $(-3, 4)$ 的模相同，都是5。

它们大小相同，但方向不同。

\section{三维空间里的向量}

前面讲的都是平面上的向量，用两个坐标 $(a_1, a_2)$ 表示。

但现实世界是三维的。

飞机在空中飞，不只是往东、往北走，还会往上飞。

在三维空间里，向量有三个坐标——

\[
\vec{a} = (a_1, a_2, a_3)
\]

$a_1$：$x$ 方向（比如东西）走了多少。

$a_2$：$y$ 方向（比如南北）走了多少。

$a_3$：$z$ 方向（比如上下）走了多少。

\subsection{三维向量的模}

\[
|\vec{a}| = \sqrt{a_1^2 + a_2^2 + a_3^2}
\]

还是勾股定理，只不过从两个方向变成了三个方向。

\subsection{算一个}

向量 $\vec{b} = (1, 2, 2)$ 的模：

\[
|\vec{b}| = \sqrt{1^2 + 2^2 + 2^2} = \sqrt{1 + 4 + 4} = \sqrt{9} = 3
\]

\subsection{再算一个}

向量 $\vec{c} = (3, 4, 0)$ 的模：

\[
|\vec{c}| = \sqrt{3^2 + 4^2 + 0^2} = \sqrt{9 + 16 + 0} = \sqrt{25} = 5
\]

第三个坐标是0，说明这个向量完全在 $xy$ 平面上，没有往上或往下走。

\section{向量的加法——两段路接在一起}

现在进入向量的运算。

第一个运算是\textbf{加法}。

\subsection{先用走路来理解}

你从家出发。

第一段：往东走3步，往北走1步。

第二段：往东走2步，往北走4步。

两段路走完，你最终在哪里？

\textbf{横向：} 第一段走了3步，第二段走了2步，总共 $3 + 2 = 5$ 步。

\textbf{纵向：} 第一段走了1步，第二段走了4步，总共 $1 + 4 = 5$ 步。

最终位置：往东5步，往北5步。

用向量来写——

第一段：$\vec{a} = (3, 1)$

第二段：$\vec{b} = (2, 4)$

两段加起来：$\vec{a} + \vec{b} = (3+2, 1+4) = (5, 5)$

\subsection{这就是向量加法的规则}

\[
\vec{a} + \vec{b} = (a_1 + b_1, a_2 + b_2)
\]

\textbf{对应坐标相加。}

横向的加横向的，纵向的加纵向的。

\subsection{为什么是这样？}

因为向量加法的意义就是“走完第一段，接着走第二段”。

横向总共走了多少？两段横向的加在一起。

纵向总共走了多少？两段纵向的加在一起。

所以是对应坐标相加。

\subsection{练一练}

\[
\vec{a} = (2, 3), \quad \vec{b} = (1, -1)
\]

\[
\vec{a} + \vec{b} = (2+1, 3+(-1)) = (3, 2)
\]

先算 $2 + 1 = 3$。

再算 $3 + (-1) = 3 - 1 = 2$。

结果：$(3, 2)$。

\[
\vec{a} = (-1, 4), \quad \vec{b} = (3, -2)
\]

\[
\vec{a} + \vec{b} = (-1+3, 4+(-2)) = (2, 2)
\]

\[
\vec{a} = (1, 2, 3), \quad \vec{b} = (4, -1, 2)
\]

\[
\vec{a} + \vec{b} = (1+4, 2+(-1), 3+2) = (5, 1, 5)
\]

三维也一样，对应坐标相加。

\section{向量的减法——从终点指回起点}

向量减法，可以理解为“加上一个反方向的向量”。

\subsection{什么是反方向的向量？}

向量 $\vec{a} = (3, 4)$，它的反方向向量是——

\[
-\vec{a} = (-3, -4)
\]

每个坐标都变成相反数。

原来往东走3步，现在变成往西走3步。

原来往北走4步，现在变成往南走4步。

大小不变，方向完全反过来。

\subsection{向量减法的公式}

\[
\vec{a} - \vec{b} = \vec{a} + (-\vec{b})
\]

$\vec{a}$ 减去 $\vec{b}$，就是 $\vec{a}$ 加上 $\vec{b}$ 的反方向。

写成坐标——

\[
\vec{a} - \vec{b} = (a_1 - b_1, a_2 - b_2)
\]

对应坐标相减。

\subsection{练一练}

\[
\vec{a} = (5, 3), \quad \vec{b} = (2, 1)
\]

\[
\vec{a} - \vec{b} = (5-2, 3-1) = (3, 2)
\]

\[
\vec{a} = (1, -2), \quad \vec{b} = (4, 3)
\]

\[
\vec{a} - \vec{b} = (1-4, -2-3) = (-3, -5)
\]

先算 $1 - 4 = -3$。

再算 $-2 - 3 = -5$。

\section{向量的数乘——把向量拉长或缩短}

一个数字乘以一个向量，叫做\textbf{数乘}。

\subsection{先看一个例子}

向量 $\vec{a} = (2, 3)$。

$2 \times \vec{a}$ 等于多少？

$2 \times \vec{a}$ 的意思是：走两次 $\vec{a}$。

第一次：往东2步，往北3步。

第二次：再往东2步，再往北3步。

总共：往东 $2+2=4$ 步，往北 $3+3=6$ 步。

\[
2 \times \vec{a} = (4, 6)
\]

这就是 $2 \times (2, 3) = (2 \times 2, 2 \times 3) = (4, 6)$。

每个坐标都乘以2。

\subsection{数乘的公式}

\[
k \times \vec{a} = k \times (a_1, a_2) = (k \cdot a_1, k \cdot a_2)
\]

把向量的每个坐标，都乘以那个数。

\subsection{数乘的意义}

\textbf{当 $k > 0$ 时：}

方向不变，大小变成原来的 $k$ 倍。

$2 \times \vec{a}$：方向不变，长度变成原来的2倍。

$0.5 \times \vec{a}$：方向不变，长度变成原来的一半。

\textbf{当 $k < 0$ 时：}

方向反过来，大小变成原来的 $|k|$ 倍。

$-1 \times \vec{a}$：方向反过来，长度不变。这就是 $-\vec{a}$。

$-2 \times \vec{a}$：方向反过来，长度变成原来的2倍。

\textbf{当 $k = 0$ 时：}

\[
0 \times \vec{a} = (0, 0) = \vec{0}
\]

结果是零向量。

任何向量乘以0，都变成零向量。

\subsection{练一练}

\[
3 \times (2, -1) = (3 \times 2, 3 \times (-1)) = (6, -3)
\]

\[
-2 \times (1, 3) = (-2 \times 1, -2 \times 3) = (-2, -6)
\]

\[
0.5 \times (4, 8) = (0.5 \times 4, 0.5 \times 8) = (2, 4)
\]

\[
-1 \times (5, -2) = (-5, 2)
\]

这就是向量 $(5, -2)$ 的反方向向量。

\section{向量的点积——两个向量有多“同向”}

这是一个新的运算。

它不像加减法和数乘那么直观，但它非常重要。

\subsection{先看计算规则}

两个向量 $\vec{a} = (a_1, a_2)$ 和 $\vec{b} = (b_1, b_2)$。

它们的\textbf{点积}（也叫内积、数量积）是——

\[
\vec{a} \cdot \vec{b} = a_1 \times b_1 + a_2 \times b_2
\]

对应坐标相乘，然后\textbf{相加}。

注意：点积的结果是\textbf{一个数}，不是向量。

\subsection{算几个具体的}

\[
\vec{a} = (1, 2), \quad \vec{b} = (3, 4)
\]

\[
\vec{a} \cdot \vec{b} = 1 \times 3 + 2 \times 4 = 3 + 8 = 11
\]

\[
\vec{a} = (2, 5), \quad \vec{b} = (1, -1)
\]

\[
\vec{a} \cdot \vec{b} = 2 \times 1 + 5 \times (-1) = 2 + (-5) = 2 - 5 = -3
\]

点积可以是负数。

\[
\vec{a} = (1, 0), \quad \vec{b} = (0, 1)
\]

\[
\vec{a} \cdot \vec{b} = 1 \times 0 + 0 \times 1 = 0 + 0 = 0
\]

这两个向量，一个朝东 $(1, 0)$，一个朝北 $(0, 1)$。

它们互相\textbf{垂直}。

它们的点积是 $0$。

\subsection{点积的几何意义}

有另一个公式，可以计算点积——

\[
\vec{a} \cdot \vec{b} = |\vec{a}| \times |\vec{b}| \times \cos\theta
\]

$|\vec{a}|$ 是向量 $\vec{a}$ 的模（长度）。

$|\vec{b}|$ 是向量 $\vec{b}$ 的模（长度）。

$\theta$ 是两个向量之间的\textbf{夹角}。

$\cos\theta$ 是这个夹角的\textbf{余弦值}。

\subsection{余弦值是什么，快速回顾一下}

在一个直角三角形里，一个锐角的余弦，等于“邻边除以斜边”。

$\cos$ 值的范围是 $-1$ 到 $1$。

几个特殊的值——

$\cos 0^\circ = 1$（夹角是0度，两个向量方向完全相同）

$\cos 90^\circ = 0$（夹角是90度，两个向量互相垂直）

$\cos 180^\circ = -1$（夹角是180度，两个向量方向完全相反）

\subsection{这个公式告诉我们什么？}

\[
\vec{a} \cdot \vec{b} = |\vec{a}| \times |\vec{b}| \times \cos\theta
\]

点积的大小，取决于三件事——

$\vec{a}$ 有多长。

$\vec{b}$ 有多长。

两个向量夹角的余弦值。

\textbf{夹角越小，$\cos\theta$ 越大，点积越大。}

当两个向量方向完全相同（$\theta = 0^\circ$），$\cos 0^\circ = 1$，点积达到最大值 $|\vec{a}| \times |\vec{b}|$。

\textbf{夹角是90度，$\cos 90^\circ = 0$，点积等于0。}

两个向量互相垂直时，点积为零。

这是判断垂直的重要工具。

\textbf{夹角大于90度，$\cos\theta$ 是负数，点积是负数。}

当两个向量方向大致相反时，点积是负的。

\subsection{所以点积在衡量什么？}

\textbf{点积衡量的是：两个向量的方向有多“一致”。}

方向越一致，点积越大。

方向垂直，点积为零。

方向相反，点积为负。

\subsection{语言记忆法}

\begin{quote}
点积，告诉你两个向量有多“同向”。

点积大，说明它们方向差不多。

点积是零，说明它们互相垂直。

点积是负数，说明它们大致相反。
\end{quote}

\section{用点积判断两个向量是否垂直}

这是点积最常用的功能之一。

\textbf{如果两个向量的点积等于零，它们就互相垂直。}

\subsection{例题}

向量 $\vec{a} = (2, -1)$，向量 $\vec{b} = (1, k)$。

已知 $\vec{a}$ 和 $\vec{b}$ 垂直，求 $k$。

两个向量垂直，点积等于零——

\[
\vec{a} \cdot \vec{b} = 0
\]

代入坐标——

\[
2 \times 1 + (-1) \times k = 0
\]

\[
2 - k = 0
\]

\[
k = 2
\]

验证：$\vec{b} = (1, 2)$。

\[
\vec{a} \cdot \vec{b} = 2 \times 1 + (-1) \times 2 = 2 - 2 = 0
\]

点积是0，确实垂直。✓

\section{用点积求两个向量的夹角}

从公式 $\vec{a} \cdot \vec{b} = |\vec{a}| \times |\vec{b}| \times \cos\theta$，可以变形得到——

\[
\cos\theta = \frac{\vec{a} \cdot \vec{b}}{|\vec{a}| \times |\vec{b}|}
\]

这个公式可以求出两个向量夹角的余弦值。

然后查表或用计算器，得到角度。

\subsection{例题}

$\vec{a} = (1, 1)$，$\vec{b} = (1, 0)$，求两个向量的夹角。

\textbf{第一步：算点积。}

\[
\vec{a} \cdot \vec{b} = 1 \times 1 + 1 \times 0 = 1 + 0 = 1
\]

\textbf{第二步：算两个向量的模。}

\[
|\vec{a}| = \sqrt{1^2 + 1^2} = \sqrt{1 + 1} = \sqrt{2}
\]

\[
|\vec{b}| = \sqrt{1^2 + 0^2} = \sqrt{1 + 0} = \sqrt{1} = 1
\]

\textbf{第三步：代入公式。}

\[
\cos\theta = \frac{1}{\sqrt{2} \times 1} = \frac{1}{\sqrt{2}}
\]

\textbf{第四步：确定角度。}

$\cos\theta = \dfrac{1}{\sqrt{2}}$

$\dfrac{1}{\sqrt{2}}$ 等于 $\dfrac{\sqrt{2}}{2}$（分子分母同乘 $\sqrt{2}$）。

$\cos 45^\circ = \dfrac{\sqrt{2}}{2}$。

所以 $\theta = 45^\circ$。

两个向量的夹角是 $45^\circ$。

\section{本章小结}

\textbf{向量是什么：}

既有大小、又有方向的量。

不只告诉你“走多远”，还告诉你“往哪走”。

\textbf{向量的坐标：}

平面上的向量，用两个数表示：$(a_1, a_2)$。

$a_1$ 是横着走了多少，$a_2$ 是竖着走了多少。

空间里的向量，用三个数表示：$(a_1, a_2, a_3)$。

\textbf{向量的模：}

\[
|\vec{a}| = \sqrt{a_1^2 + a_2^2}
\]

用勾股定理，算实际走了多远。

\textbf{向量的加法：}

对应坐标相加。

走完第一段，接着走第二段。

\textbf{向量的减法：}

对应坐标相减。

\textbf{向量的数乘：}

每个坐标都乘以那个数。

方向不变（正数）或反向（负数），大小缩放。

\textbf{向量的点积：}

\[
\vec{a} \cdot \vec{b} = a_1 b_1 + a_2 b_2
\]

对应坐标相乘后相加，结果是一个数。

点积衡量两个向量方向的一致程度。

点积为零，两向量垂直。

\textbf{点积的另一个公式：}

\[
\vec{a} \cdot \vec{b} = |\vec{a}| \times |\vec{b}| \times \cos\theta
\]

可以用来求夹角。

向量，是数学用来描述“有方向的量”的语言。

物理里的力、速度、加速度，都是向量。

没有向量，物理学寸步难行。

\chapter{什么是维度——从一根线到整个宇宙}

\section{先从一个问题开始}

你站在一条笔直的路上。

这条路，往前延伸到看不见的地方，往后也延伸到看不见的地方。

你只能往前走，或者往后走。

除了前后，你哪儿也去不了。

现在问你：你能去的地方，有多少种“方向”？

\textbf{只有一种。}

前后是同一种方向——往前是正的，往后是负的，但本质上就是“沿着这条路”这一种方向。

这条路，就是\textbf{一维}的。

一维，意思是：\textbf{只有一种方向可以走。}

\section{一维：一根线}

\subsection{一维世界长什么样？}

想象一只蚂蚁，住在一根细细的绳子上。

这只蚂蚁只能沿着绳子爬。

往前爬，或者往后爬。

它没办法离开绳子，没办法往左跳，没办法往右跳，没办法往上飞，没办法往下钻。

它的整个世界，就是这根绳子。

对这只蚂蚁来说，“位置”只需要一个数字就能描述。

我们可以在绳子上选一个点，叫做原点，记作0。

往前爬3厘米，位置就是 $+3$。

往后爬2厘米，位置就是 $-2$。

一个数字，就能说清楚蚂蚁在哪里。

这就是一维。

\textbf{一维空间里，描述位置只需要一个数字。}

\subsection{一维的数学名字}

一维空间，数学上叫做\textbf{直线}，或者\textbf{数轴}。

你画一条直线，上面标上刻度，0在中间，正数在右边，负数在左边。

这就是一维空间的样子。

\section{二维：一张纸}

\subsection{从一维走到二维}

现在，那只蚂蚁从绳子上爬下来了。

它爬到了一张无限大的纸上。

这张纸平铺在地上，没有厚度，只是一个平面。

现在蚂蚁能去的地方变多了。

它不再只能往前或往后。

它可以往前，往后，往左，往右，往左前方，往右后方……

它可以去纸上的任何一个点。

但仔细想想，这些方向，本质上只有\textbf{两种}。

一种是“前后”方向（我们叫它 $x$ 方向）。

另一种是“左右”方向（我们叫它 $y$ 方向）。

任何其他方向，都可以拆成这两种方向的组合。

比如“往左前方走”，就是“往前走一点，同时往左走一点”。

\textbf{两种独立的方向。}

这就是二维。

\subsection{二维世界里，位置怎么描述？}

在一维世界里，一个数字就够了。

在二维世界里，需要\textbf{两个数字}。

我们在纸上画两条互相垂直的线。

一条横着的，叫 $x$ 轴。

一条竖着的，叫 $y$ 轴。

两条线交叉的地方，叫原点，记作 $(0, 0)$。

现在，纸上任何一个点，都可以用两个数字来描述。

比如 $(3, 2)$：从原点往右走3格，再往上走2格。

比如 $(-1, 4)$：从原点往左走1格，再往上走4格。

比如 $(5, -3)$：从原点往右走5格，再往下走3格。

两个数字，横着多少，竖着多少，就能说清楚你在纸上的哪个位置。

\textbf{二维空间里，描述位置需要两个数字。}

\subsection{二维的数学名字}

二维空间，数学上叫做\textbf{平面}。

你在纸上画的所有图形——三角形、圆、正方形——都是二维的。

它们只有长和宽，没有高。

\section{三维：我们生活的世界}

\subsection{从二维走到三维}

现在，那只蚂蚁学会飞了。

它不再只能在纸上爬。

它可以飞起来，离开纸面，飞到空中。

现在它能去的地方更多了。

除了前后、左右，还有\textbf{上下}。

\textbf{三种独立的方向。}

前后（$x$ 方向）。

左右（$y$ 方向）。

上下（$z$ 方向）。

这就是三维。

\subsection{三维世界里，位置怎么描述？}

在二维世界里，两个数字就够了。

在三维世界里，需要\textbf{三个数字}。

我们在空间里建立三条互相垂直的轴。

$x$ 轴：前后方向。

$y$ 轴：左右方向。

$z$ 轴：上下方向。

三条轴交叉的地方，是原点 $(0, 0, 0)$。

空间里任何一个点，都可以用三个数字来描述。

比如 $(3, 2, 5)$：往前3，往右2，往上5。

比如 $(0, 0, -4)$：不往前后，不往左右，往下4。

三个数字，前后多少，左右多少，上下多少，就能说清楚你在空间里的哪个位置。

\textbf{三维空间里，描述位置需要三个数字。}

\subsection{我们就生活在三维空间}

你现在坐的房间，是三维的。

房间有长度、宽度、高度。

你可以在房间里往前后走、往左右走、往上下跳（虽然跳不高）。

地球是三维的。

宇宙是三维的（至少我们能感知到的宇宙是三维的）。

\section{零维：一个点}

\subsection{比一维还要少}

我们往回退一步。

一维是一根线，可以往前或往后走。

如果连前后都没有呢？

如果你哪儿也去不了呢？

那就是\textbf{零维}。

零维就是一个点。

没有长度，没有宽度，没有高度。

只是一个位置，不能移动。

住在零维世界里的生物，只能待在原地。

它的整个宇宙，就是那一个点。

\textbf{零维空间里，描述位置不需要任何数字——因为只有那一个位置，没有别的地方可去。}

\section{把规律说清楚}

\begin{center}
\begin{tabular}{c|c|c|c}
维度 & 需要几个数字描述位置 & 有几种独立的方向 & 是什么样子 \\
\hline
\textbf{零维} & 0个 & 0种 & 一个点 \\
\textbf{一维} & 1个 & 1种（前后） & 一条线 \\
\textbf{二维} & 2个 & 2种（前后、左右） & 一个平面 \\
\textbf{三维} & 3个 & 3种（前后、左右、上下） & 我们生活的空间
\end{tabular}
\end{center}

\textbf{维度的本质：需要几个独立的数字才能描述位置。}

一维需要1个数字。

二维需要2个数字。

三维需要3个数字。

所以它们分别叫一维、二维、三维。

\section{四维？更高维度？}

\subsection{数学上可以继续推广}

数学家说：既然一维用1个数字，二维用2个数字，三维用3个数字，那为什么不能有四维、五维、一百维？

四维空间，就是需要4个数字才能描述位置的空间。

五维空间，就是需要5个数字才能描述位置的空间。

在四维空间里，一个点的位置写成 $(x_1, x_2, x_3, x_4)$。

四个数字。

\subsection{我们能想象四维吗？}

说实话，很难。

我们的大脑是在三维世界里进化出来的，我们很难真正“看到”四维空间。

就像那只住在纸上的蚂蚁，很难想象“往上飞”是什么感觉。

但数学不需要“看到”。

数学只需要逻辑。

三维空间里的点用三个数字描述：$(x, y, z)$。

四维空间里的点用四个数字描述：$(x, y, z, w)$。

所有在三维空间里成立的数学规则，都可以推广到四维、五维、任意多维。

\subsection{高维度有什么用？}

你可能在想：四维、五维这些东西，又看不见摸不着，学它干什么？

其实高维度在很多地方都有用。

\textbf{物理学：} 爱因斯坦的相对论里，时间被当作第四个维度，和空间的三个维度放在一起，形成“四维时空”。

\textbf{数据科学：} 如果你有一个表格，里面有100个指标（比如身高、体重、年龄、收入……），每一行数据就是100维空间里的一个点。

\textbf{机器学习：} 训练人工智能的时候，经常要在几百万维的空间里找最优解。

所以高维度不是数学家的玩具，它是真正有用的工具。

不过，这本书里我们主要讨论的是二维和三维，这两个维度足够用很久了。

\section{维度和几何的关系}

\subsection{不同维度里，几何对象不一样}

在一维空间（直线）里：

你能画的最基本的“图形”，就是\textbf{线段}——直线上的一段。

在二维空间（平面）里：

你能画\textbf{线段}，还能画\textbf{三角形}、\textbf{圆}、\textbf{正方形}……

这些图形有长有宽，但没有厚度。

在三维空间里：

你能做\textbf{线段}，还能做\textbf{三角形}（平面图形），还能做\textbf{立方体}、\textbf{球}、\textbf{圆柱}……

这些物体有长有宽有高。

\subsection{低维度的东西，可以放进高维度}

一条线（一维的），可以放在平面上（二维的）。

一个三角形（二维的），可以放在空间里（三维的）。

低维度的东西，永远可以嵌入到高维度里。

但反过来不行。

一个立方体（三维的），没办法真正放进一张纸里（二维的）。

你可以在纸上\textbf{画}一个立方体的样子，但那只是一张图，不是真正的立方体。

\section{高维度往低维度“看”，会丢失信息}

\subsection{用一个你每天都见到的事情来说}

你站在阳光下。

你的身后，地上有一个影子。

你是三维的——有高度、有宽度、有厚度。

你的影子是二维的——只是地上的一个黑色形状，平平的，没有厚度。

现在问一个问题：

\textbf{只看影子，能知道你穿的衣服是什么颜色吗？}

不能。

影子是黑色的。

不管你穿红衣服、蓝衣服、白衣服，影子都是黑的。

颜色这个信息，在变成影子的时候，丢掉了。

\subsection{再问一个问题}

\textbf{只看影子，能知道你的鼻子有多高吗？}

不能。

影子是平的，它只是一个轮廓。

你的鼻子高不高、眼睛大不大、耳朵什么形状——这些细节，影子都没有。

因为影子是二维的，它只能记录“外面那一圈轮廓”，没办法记录“往前凸出多少”。

\subsection{再举一个例子}

你把一个杯子放在桌上。

从正上方往下看，你看到的是一个圆形。

但杯子明明是一个圆柱体——有高度的。

从上往下看，高度这个信息，看不见了。

杯子是三维的。

从正上方看，你只能看到二维的图像（一个圆）。

杯子有多高？看不出来。

杯子是空心的还是实心的？看不出来。

这些信息，在“从上往下看”的时候，丢掉了。

\subsection{这就叫“降维会丢失信息”}

三维的东西，变成二维的图像，会丢失一个方向的信息。

你只能看到两个方向的样子，看不到第三个方向。

同样的道理：

二维的东西，变成一维的，也会丢失信息。

想象一张纸上画了一个三角形。

你把这张纸侧过来看，眯起眼睛，纸变成了一条线。

三角形不见了，你只能看到一条线段。

三角形的形状信息，丢掉了。

\subsection{用一句话说清楚}

\begin{quote}
\textbf{高维度的东西，缩到低维度，会丢掉一些信息。}

\textbf{就像你站在阳光下，影子不知道你穿的衣服是什么颜色。}
\end{quote}

\section{为什么要学这些？}

因为接下来要学的很多东西，都和维度有关。

\textbf{在二维平面里：}

两条直线只能相交或平行。

不存在“两条线既不相交又不平行”的情况。

\textbf{在三维空间里：}

两条直线可以相交、可以平行，还可以\textbf{异面}——既不相交又不平行。

异面这种关系，只有在三维空间里才存在。

如果你不清楚“我现在讨论的是二维还是三维”，就很容易搞混。

二维里的结论，不能直接搬到三维。

三维里的结论，也不能随便缩回二维。

所以，在学几何之前，先把维度这件事想清楚。

\textbf{我在几维空间里讨论问题？}

这个问题的答案，决定了哪些结论是成立的，哪些是不成立的。

\section{本章小结}

\textbf{维度是什么：}

维度就是描述位置需要几个独立的数字。

一个数字：一维（线）。

两个数字：二维（面）。

三个数字：三维（空间）。

\textbf{我们生活在三维空间里：}

前后、左右、上下，三种方向。

任何位置都可以用三个数字描述。

\textbf{低维和高维的关系：}

低维的东西可以放进高维。

高维的东西不能真正放进低维（只能投影、只能画图）。

\textbf{为什么要知道这些：}

因为不同维度里，几何的规则不同。

二维里两条线要么相交要么平行。

三维里两条线还可以异面。

搞清楚你在几维空间里讨论问题，是正确推理的前提。

现在你知道维度是什么了。

下一章，我们正式进入“线的世界”——平行、垂直、相交，以及它们之间的关系。

\chapter{线的世界——平行、垂直、相交，和它们之间的关系}

\section{先从一根线开始}

你拿一支笔，在纸上画一条线。

这条线，是数学里最基本的东西之一。

但“线”这个字，在数学里有好几种意思。

我们先把它们分清楚。

\section{线段、射线、直线——三兄弟}

\subsection{线段}

你在纸上点两个点，$A$ 和 $B$。

然后用笔把它们连起来。

这就是一条\textbf{线段}。

\[
\overline{AB}
\]

线段有两个\textbf{端点}：$A$ 和 $B$。

线段有\textbf{长度}：从 $A$ 到 $B$ 的距离。

线段\textbf{不能延伸}——到了端点就停了。

用生活来感受：

你用尺子画了一条10厘米的线。

这条线有开头，有结尾，有长度。

这就是线段。

\subsection{射线}

你在纸上点一个点 $A$。

从 $A$ 出发，朝一个方向画，一直画，永远不停。

这就是一条\textbf{射线}。

\[
\overrightarrow{AB}
\]

射线有一个\textbf{端点}：$A$（出发的那个点）。

射线有\textbf{方向}：从 $A$ 朝着 $B$ 的方向。

射线\textbf{只朝一个方向延伸}，另一边到了端点就停了。

用生活来感受：

手电筒发出的光。

光从手电筒出发（端点），朝一个方向射出去，一直射到无穷远。

这就是射线。

\subsection{直线}

你在纸上画一条线，两边都不停，往两边无限延伸。

这就是一条\textbf{直线}。

\[
\overleftrightarrow{AB}
\]

或者简单写作“直线 $AB$”。

直线\textbf{没有端点}。

直线\textbf{没有长度}（因为它无限长）。

直线\textbf{朝两个方向都无限延伸}。

用生活来感受：

严格来说，生活中没有真正的直线，因为任何东西都有尽头。

但你可以想象：一条铁轨，两头都延伸到看不见的地方。

这就是直线的感觉。

\subsection{三兄弟的区别，一张表说清楚}

\begin{center}
\begin{tabular}{c|c|c|c}
 & 端点 & 长度 & 延伸方向 \\
\hline
\textbf{线段} & 两个端点 & 有，可以量 & 不延伸 \\
\textbf{射线} & 一个端点 & 没有（无限长） & 朝一个方向延伸 \\
\textbf{直线} & 没有端点 & 没有（无限长） & 朝两个方向延伸
\end{tabular}
\end{center}

\section{两条直线放在一起，会发生什么？}

一条直线很孤独。

但两条直线放在一起，就有故事了。

在同一个平面上，两条直线之间的关系，\textbf{只有三种可能}——

\textbf{相交：} 两条直线交叉，有一个交点。

\textbf{平行：} 两条直线永远不交叉，方向相同，保持固定的距离。

\textbf{重合：} 两条直线完全叠在一起，是同一条线。

（严格来说，重合是平行的一种特殊情况——两条平行线之间的距离为零。但在日常学习中，我们一般把重合单独拿出来说。）

\section{相交}

\subsection{什么叫相交？}

两条直线，在某一个点碰到了一起。

这个点，叫做\textbf{交点}。

这里有一个非常重要的事情，必须说清楚：

\textbf{相交的意思是：两条直线有一个公共点。}

这个公共点，既在这条线上，又在那条线上。

两条线在这个点上碰到了，像两根筷子交叉放在一起。

\textbf{不是“看起来交叉”，是“真的有一个点，同时在两条线上”。}

为什么要强调这个？因为到了三维空间里，这个区分会变得非常重要。后面会详细讲。

\subsection{相交的特点}

两条直线相交，\textbf{有且只有一个交点}。

为什么只有一个？

因为两条不同的直线，如果在两个点都碰到了，那它们就必须完全重合——同一条线。

所以要么不交叉（平行），要么只交叉一次（一个交点）。

\subsection{相交会产生什么？}

两条直线相交，会产生\textbf{角}。

在交点处，四个角被创造出来了。

对面的两个角，叫做\textbf{对顶角}。

相邻的两个角，叫做\textbf{邻补角}。

\subsection{对顶角}

交点处，对面的两个角。

\textbf{对顶角相等。}

这不需要证明，这是直线相交的直接结果。

你可以用手比划一下——两条线交叉，对面的角度必然一样大。

\subsection{邻补角}

交点处，相邻的两个角。

\textbf{邻补角之和等于 $180^\circ$。}

因为它们合在一起，正好铺满一条直线的一侧（一条直线是 $180^\circ$）。

\section{垂直——一种特殊的相交}

\subsection{什么叫垂直？}

两条直线相交，交点处形成四个角。

如果这四个角\textbf{都是 $90^\circ$}，就说这两条直线\textbf{互相垂直}。

其实只要有\textbf{一个角是 $90^\circ$}，其余三个自动都是 $90^\circ$。

为什么？

因为对顶角相等（对面那个也是 $90^\circ$），邻补角之和是 $180^\circ$（旁边那个是 $180^\circ - 90^\circ = 90^\circ$）。

所以只要一个是直角，四个全是直角。

\subsection{垂直的符号}

\[
AB \perp CD
\]

读作“AB垂直于CD”。

$\perp$ 这个符号，长得就像一个直角——一条横线，一条竖线，互相垂直。

\subsection{垂足}

两条互相垂直的直线，它们的交点叫做\textbf{垂足}。

\subsection{垂直在生活中}

你房间的墙角——两面墙互相垂直。

十字路口——两条路互相垂直。

书本的边——横边和竖边互相垂直。

桌腿和地面——互相垂直。

\subsection{从一个点向一条直线画垂线}

这个操作你会经常用到。

在一条直线外面有一个点 $P$。

从 $P$ 向这条直线画一条线，使得这条线和原来的直线垂直。

这条线叫做\textbf{垂线}。

垂线和原直线的交点，叫做\textbf{垂足}。

\textbf{垂线段是最短的。}

从一个点到一条直线，可以画无数条线段。

但其中\textbf{最短的那一条，就是垂线段}。

这个结论非常重要。

它的意思是：\textbf{点到直线的距离，就是垂线段的长度。}

不是斜着量的，是垂直量的。

\section{平行——永远不相交}

\subsection{什么叫平行？}

两条直线，在同一个平面上，无论怎么延伸，永远不会交叉。

它们之间的距离，处处相等。

\subsection{平行的符号}

\[
AB \parallel CD
\]

读作“AB平行于CD”。

$\parallel$ 这个符号，就是两条平行的竖线。

\subsection{平行的判定——怎么知道两条线平行？}

有好几种方法。

\textbf{方法一：同位角相等。}

一条直线切过两条平行线，形成的\textbf{同位角}相等。

什么是同位角？

想象一条横线（叫做截线），切过两条平行线。

在上面那条平行线的左上角，有一个角。

在下面那条平行线的左上角，也有一个角。

这两个角，位置相同（都在左上），就叫同位角。

如果它们相等，两条线就平行。

反过来也成立：如果两条线平行，同位角就一定相等。

\textbf{方法二：内错角相等。}

截线切过两条线，形成的\textbf{内错角}相等。

什么是内错角？

在两条平行线\textbf{之间}（内），在截线的\textbf{两侧}（错），各有一个角。

这两个角叫内错角。

如果它们相等，两条线平行。

\textbf{方法三：同旁内角互补。}

截线切过两条线，形成的\textbf{同旁内角}之和等于 $180^\circ$。

什么是同旁内角？

在两条平行线之间，在截线的\textbf{同一侧}，各有一个角。

这两个角叫同旁内角。

如果它们加起来等于 $180^\circ$，两条线平行。

\subsection{这三种角，用大白话讲一遍}

一条线切过两条平行线，在交点处产生了很多角。

这些角之间有固定的关系。

\textbf{同位角：} 位置相同的一对角。相等。

\textbf{内错角：} 在两线之间、截线两侧的一对角。相等。

\textbf{同旁内角：} 在两线之间、截线同侧的一对角。加起来等于 $180^\circ$。

\subsection{平行线之间的距离}

两条平行线之间的距离，是从一条线上的任意一点，向另一条线画垂线段，这条垂线段的长度。

无论你在哪个位置画，这个距离都是一样的。

这就是平行线的特点——\textbf{处处等距}。

\section{在坐标系里，怎么判断平行和垂直？}

前面用角度来判断平行和垂直。

但在坐标系里，有更方便的工具——\textbf{斜率}。

\subsection{斜率是什么？快速回顾}

（还记得斜率吗？就是一条直线有多陡。）

\[
k = \frac{y_2 - y_1}{x_2 - x_1}
\]

竖着走了多少，除以横着走了多少。

爬了多高，除以走了多远。

\subsection{用斜率判断平行}

\textbf{两条直线平行 $\Leftrightarrow$ 它们的斜率相等。}

\[
l_1 \parallel l_2 \quad \Leftrightarrow \quad k_1 = k_2
\]

\textbf{为什么？}

斜率描述的是直线的倾斜程度。

两条线一样陡，方向就一样。

方向一样，就永远不会交叉。

不交叉，就是平行。

\textbf{举个例子：}

直线 $l_1$：$y = 2x + 1$，斜率 $k_1 = 2$。

直线 $l_2$：$y = 2x + 5$，斜率 $k_2 = 2$。

$k_1 = k_2 = 2$，所以 $l_1 \parallel l_2$。

这两条线一样陡，只是上下位置不同（截距不同）。

\textbf{再举一个：}

直线 $l_1$：$y = 3x + 2$，斜率 $k_1 = 3$。

直线 $l_2$：$y = -x + 4$，斜率 $k_2 = -1$。

$k_1 \neq k_2$，所以 $l_1$ 和 $l_2$ 不平行——它们会相交。

\subsection{用斜率判断垂直}

\textbf{两条直线垂直 $\Leftrightarrow$ 它们的斜率相乘等于 $-1$。}

\[
l_1 \perp l_2 \quad \Leftrightarrow \quad k_1 \times k_2 = -1
\]

\textbf{为什么是 $-1$？}

这个结论的推导需要用到三角函数，这里不做完整的证明，但可以感受一下——

一条线斜率是正的（往右上方走），另一条线斜率是负的（往右下方走），而且它们的陡峭程度刚好“互补”，乘起来就是 $-1$。

\textbf{举个例子：}

直线 $l_1$：$y = 2x + 1$，斜率 $k_1 = 2$。

直线 $l_2$：$y = -\dfrac{1}{2}x + 3$，斜率 $k_2 = -\dfrac{1}{2}$。

$k_1 \times k_2 = 2 \times \left(-\dfrac{1}{2}\right) = -1$

所以 $l_1 \perp l_2$。

\textbf{再举一个：}

直线 $l_1$：$y = 3x$，斜率 $k_1 = 3$。

直线 $l_2$：$y = -\dfrac{1}{3}x$，斜率 $k_2 = -\dfrac{1}{3}$。

$k_1 \times k_2 = 3 \times \left(-\dfrac{1}{3}\right) = -1$

垂直。✓

\subsection{一个特殊情况：水平线和竖直线}

水平线（比如 $y = 3$），斜率是 $0$。

竖直线（比如 $x = 2$），斜率\textbf{不存在}。

但水平线和竖直线明显互相垂直。

这个时候，$k_1 \times k_2 = -1$ 这个公式用不了（因为竖直线没有斜率）。

所以要单独记住：\textbf{水平线和竖直线互相垂直。}

\section{用向量判断平行和垂直}

向量也可以用来判断两条线的关系。

\subsection{方向向量}

每一条直线，都有一个方向。

我们可以找一个\textbf{沿着这条直线方向}的向量，叫做这条直线的\textbf{方向向量}。

比如，直线 $y = 2x + 1$。

这条线上取两个点：$(0, 1)$ 和 $(1, 3)$。

从第一个点到第二个点的向量：

\[
\vec{d} = (1 - 0, 3 - 1) = (1, 2)
\]

$(1, 2)$ 就是这条直线的一个方向向量。

（方向向量不唯一，$(2, 4)$、$(-1, -2)$ 也可以，只要方向相同或相反就行。）

\subsection{用方向向量判断平行}

\textbf{两条直线平行 $\Leftrightarrow$ 它们的方向向量平行。}

两个向量平行，意味着一个是另一个的倍数。

$\vec{d_1} = (a_1, b_1)$，$\vec{d_2} = (a_2, b_2)$。

$\vec{d_1}$ 和 $\vec{d_2}$ 平行 $\Leftrightarrow$ $a_1 b_2 - a_2 b_1 = 0$

\textbf{这个 $a_1 b_2 - a_2 b_1$，你认识它吗？}

它就是一个 $2 \times 2$ 行列式：

\[
\begin{vmatrix} a_1 & b_1 \\ a_2 & b_2 \end{vmatrix} = a_1 b_2 - a_2 b_1
\]

行列式等于零，两个向量平行，两条直线平行。

\textbf{举个例子：}

直线 $l_1$ 的方向向量：$\vec{d_1} = (1, 2)$

直线 $l_2$ 的方向向量：$\vec{d_2} = (3, 6)$

\[
\begin{vmatrix} 1 & 2 \\ 3 & 6 \end{vmatrix} = 1 \times 6 - 2 \times 3 = 6 - 6 = 0
\]

行列式等于零，$\vec{d_1}$ 和 $\vec{d_2}$ 平行。

而且你可以看到，$(3, 6) = 3 \times (1, 2)$，$\vec{d_2}$ 就是 $\vec{d_1}$ 的3倍。

所以 $l_1 \parallel l_2$。

\subsection{用方向向量判断垂直}

\textbf{两条直线垂直 $\Leftrightarrow$ 它们的方向向量互相垂直。}

两个向量垂直，上一章学过怎么判断——

\textbf{点积等于零。}

\[
\vec{d_1} \cdot \vec{d_2} = 0
\]

\[
a_1 a_2 + b_1 b_2 = 0
\]

\textbf{举个例子：}

直线 $l_1$ 的方向向量：$\vec{d_1} = (1, 2)$

直线 $l_2$ 的方向向量：$\vec{d_2} = (2, -1)$

\[
\vec{d_1} \cdot \vec{d_2} = 1 \times 2 + 2 \times (-1) = 2 + (-2) = 0
\]

点积等于零，$\vec{d_1}$ 和 $\vec{d_2}$ 垂直。

所以 $l_1 \perp l_2$。

\section{平行和垂直的重要性质}

\subsection{平行的性质}

\textbf{性质一：平行线之间的距离处处相等。}

前面说过了。

\textbf{性质二：平行线被截线截出的同位角相等、内错角相等、同旁内角互补。}

前面也说过了。

\textbf{性质三：两条直线都平行于第三条直线，那么这两条直线也互相平行。}

\[
l_1 \parallel l_3 \text{ 且 } l_2 \parallel l_3 \quad \Rightarrow \quad l_1 \parallel l_2
\]

因为甲和丙方向相同，乙和丙方向也相同。

三条线方向都相同，当然互相平行。

\textbf{性质四：过直线外一点，有且只有一条直线与已知直线平行。}

这叫做\textbf{平行公理}。

在一条直线外面有一个点 $P$。

过 $P$ 点，能画几条线和原来的直线平行？

\textbf{只能画一条。}

不能画两条，不能画零条，就是一条。

这个性质看起来很简单，但它是整个欧几里得几何的基础之一。

\subsection{垂直的性质}

\textbf{性质一：过直线上或直线外一点，有且只有一条直线与已知直线垂直。}

不管那个点在直线上还是直线外，过这个点能画的垂线，有且只有一条。

\textbf{性质二：点到直线的距离等于垂线段的长度。}

从一个点到一条直线，画无数条线段，最短的那一条就是垂线段。

这个最短的长度，就是点到直线的距离。

\textbf{性质三：三角形的三条高交于一点。}

三角形的每个顶点，都可以向对边画一条垂线（这叫做高）。

三条高，交于同一个点。

这个点叫做\textbf{垂心}。

\subsection{平行和垂直的“混合”性质}

\textbf{如果两条直线都垂直于同一条直线，那么这两条直线互相平行。}

\[
l_1 \perp l_3 \text{ 且 } l_2 \perp l_3 \quad \Rightarrow \quad l_1 \parallel l_2
\]

$l_1$ 和 $l_3$ 成 $90^\circ$，$l_2$ 和 $l_3$ 也成 $90^\circ$。

$l_1$ 和 $l_2$ 与 $l_3$ 的夹角都是 $90^\circ$，说明 $l_1$ 和 $l_2$ 的方向相同。

方向相同，就是平行。

用生活来感受：

你在地上放一根棍子。

在棍子的不同位置，各竖一根旗杆（和棍子垂直）。

这两根旗杆，互相平行。

\textbf{如果一条直线垂直于两条平行线中的一条，那么它也垂直于另一条。}

\[
l_1 \parallel l_2 \text{ 且 } l_3 \perp l_1 \quad \Rightarrow \quad l_3 \perp l_2
\]

$l_1$ 和 $l_2$ 平行，方向一样。

$l_3$ 垂直于 $l_1$。

既然 $l_2$ 的方向和 $l_1$ 一样，$l_3$ 垂直于 $l_1$，自然也垂直于 $l_2$。

\section{在空间里——多了一种关系，但要非常小心}

\subsection{先把一件事说死}

\textbf{在同一个平面上，两条不重合的直线，要么相交，要么平行。}

\textbf{没有第三种可能。}

这一点，在平面几何里是铁律，不可动摇。

\subsection{那“异面直线”是怎么回事？}

当我们从平面走进三维空间——注意，是\textbf{三维空间}——情况才会发生变化。

在三维空间里，两条直线之间，确实存在一种新的关系：

\textbf{两条直线不在同一个平面上。}

这种情况，叫做\textbf{异面}。

\subsection{“不在同一个平面上”是什么意思？}

先说“在同一个平面上”是什么意思。

两条直线在同一个平面上，意思是：你能找到一张无限大的纸，把这两条线\textbf{同时}画在这张纸上。

如果你能找到这样一张纸，这两条线就“共面”。

共面的两条线，只能相交或者平行，没有其他可能。

如果你\textbf{找不到}任何一张纸能同时装下这两条线呢？

那这两条线就是“异面”的。

\subsection{用一个非常具体的例子来感受}

想象一座立交桥。

上面有一条路，从东往西。

下面有一条路，从南往北。

这两条路相交吗？

不相交。上面的路在上层，下面的路在下层，它们没有公共点——没有任何一个点，既在上面那条路上，又在下面那条路上。

这两条路平行吗？

也不平行。一条东西方向，一条南北方向，方向完全不同。

这两条路在同一个平面上吗？

不在。你没办法找到一张无限大的纸，同时把这两条路都画在上面。

\textbf{既不相交，又不平行，也不在同一个平面上。}

这就是\textbf{异面直线}。

\subsection{“看起来交叉”不等于“相交”}

这一点非常重要。

你站在远处看立交桥，上面的路和下面的路好像交叉了。

但它们真的相交吗？

\textbf{不相交。}

上面的路在上面那层，下面的路在下面那层。

它们之间隔着高度差，没有公共点。

\textbf{相交必须是两条线真的碰到同一个点。}

存在一个点 $P$，这个点同时在两条线上。

$P$ 在直线 $l_1$ 上：✓

$P$ 在直线 $l_2$ 上：✓

这样的点 $P$ 存在，两条线才叫相交。

如果找不到这样的点 $P$，两条线就不相交。

\subsection{那为什么方向不同，还可能碰不到？}

在平面上（二维），方向不同的两条线，一定会碰到。

因为在一张纸上，两条不平行的线，延伸下去，迟早会交叉。

这是平面几何的铁律。

但在三维空间里，两条线可以“错开”。

一条在上面，一条在下面。

一条朝东西，一条朝南北。

它们各走各的路，方向不同，但因为不在同一个“高度”（或者说不在同一个平面），永远碰不到。

\subsection{异面直线只存在于三维空间}

这一点必须反复强调——

\textbf{在一张纸上（二维平面），不存在异面直线。}

在一张纸上，所有的线都在同一个平面上。

两条不平行的线，一定会在纸上的某个点交叉。

\textbf{只有在三维空间里，两条线才可能“错开”，既不相交，又不平行。}

\subsection{异面直线所成的角}

虽然异面直线不相交，但我们还是可以定义它们之间的“夹角”。

方法是——

把其中一条直线\textbf{平移}（方向不变，位置移动），移到和另一条直线相交的位置。

然后量这两条线在交点处形成的角度。

这个角度，叫做\textbf{异面直线所成的角}。

\textbf{规定这个角度取锐角或直角，也就是在 $0^\circ$ 到 $90^\circ$ 之间。}

\subsection{总结：空间里两条直线的关系}

\begin{center}
\begin{tabular}{c|c|c|c}
情况 & 是否有公共点 & 方向是否相同 & 是否共面 \\
\hline
\textbf{相交} & 有一个公共点 & 方向不同 & 共面 \\
\textbf{平行} & 没有公共点 & 方向相同 & 共面 \\
\textbf{异面} & 没有公共点 & 方向不同 & 不共面
\end{tabular}
\end{center}

\textbf{如果你做题的时候，题目说“在平面内”，就不用考虑异面，只有相交和平行两种可能。}

\textbf{如果题目说“在空间中”，才需要考虑异面。}

这个区分非常重要。搞混了，后面所有的推理都会出问题。

\section{本章小结}

\textbf{线段、射线、直线：}

线段有两个端点，有长度。

射线有一个端点，朝一个方向无限延伸。

直线没有端点，朝两个方向无限延伸。

\textbf{两条直线的关系（平面上）：}

相交：有一个交点。

平行：永不相交，方向相同，处处等距。

重合：完全叠在一起。

\textbf{垂直：}

相交的特殊情况，交点处四个角都是 $90^\circ$。

符号：$\perp$。

\textbf{平行：}

永不相交，方向相同。

符号：$\parallel$。

\textbf{用斜率判断：}

平行：$k_1 = k_2$（斜率相等）。

垂直：$k_1 \times k_2 = -1$（斜率相乘等于 $-1$）。

\textbf{用向量判断：}

平行：方向向量的行列式等于零（$a_1 b_2 - a_2 b_1 = 0$）。

垂直：方向向量的点积等于零（$a_1 a_2 + b_1 b_2 = 0$）。

\textbf{重要性质：}

两条线都垂直于同一条线 $\rightarrow$ 这两条线平行。

两条线都平行于同一条线 $\rightarrow$ 这两条线也平行。

一条线垂直于两条平行线中的一条 $\rightarrow$ 它也垂直于另一条。

\textbf{三维空间：}

多了一种关系——异面直线：不在同一平面上，既不相交也不平行。

异面只存在于三维空间，平面上不存在异面直线。

\chapter{\texorpdfstring{双重积分——难道 $y$ 是空气吗？}{双重积分——难道 y 是空气吗？}}

上一章，我们认识了积分符号 $\int$。

记得吗？

\[
\int_a^b f(x) \, dx
\]

翻译成人话：

\begin{quote}
“从 $a$ 到 $b$，把每根竖条的面积【高度 $f(x)$ $\times$ 宽度 $dx$】全部加起来。”
\end{quote}

$x$ 是方向，从左走到右。

$f(x)$ 是高度，每个位置有多高。

$dx$ 是宽度，每根竖条有多窄。

全部加起来，就是曲线下面的面积。

但你有没有发现一件奇怪的事？

$y$ 呢？

平面直角坐标系里，明明有 $x$ 和 $y$ 两个方向。

$x$ 当了“路”——我们沿着 $x$ 方向切、走、加。

$y$ 从头到尾没露过面。

难道 $y$ 是空气？

\section{$y$ 不是空气，它一直都在}

其实，$y$ 在上一章里一直都在。

只是它没用“$y$”这个名字出场。

还记得吗？

函数 $y = f(x)$。

$y$ 就是 $f(x)$。

$y$ 就是\textbf{高度}。

上一章里，$y$ 的任务是什么？

\textbf{当高度。}

每根竖条有多高？$y$ 来决定。

$y$ 不是路，它是被量的那个东西。

它不是方向，它是大小。

所以上一章的世界是这样的：

\begin{itemize}
\item $x$：路（方向）
\item $y$：高度（大小）
\end{itemize}

只有\textbf{一条路}。

$y$ 没有路可走，只能当高度。

\section{但如果，$y$ 也想当路呢？}

上一章的图，你画过。

一根横着的 $x$ 轴，上面有一条曲线，曲线的高度就是 $y$。

整个图是\textbf{平的}。画在一张纸上，是二维的。

但现在，我们要多加一个方向。

\textbf{拿出你的纸和笔！}

认真的，拿出来。下面的东西不画出来，光看文字会很难想象。

\subsection{画一个立体的坐标系}

\textbf{第一步：画一个扁扁的叉。}

在纸的中间偏下一点，画一个 ×。

但是不要画成正正方方的 ×。

把它\textbf{压扁一点}，画得矮矮的、宽宽的。

像一个被人坐扁了的 ×。

\textbf{第二步：看叉的下半部分。}

这个扁叉有四个角。

我们只看\textbf{下面的两条线}。

一条往\textbf{右}伸出去——在它的末端写上 \textbf{$x$}。

一条往\textbf{左}伸出去——在它的末端写上 \textbf{$y$}。

上面那两条线不用管了，用橡皮擦掉，只留下面两条。

\textbf{第三步：从交叉点往上画一条竖线。}

从叉的正中间（两条线交叉的那个点），\textbf{笔直往上}画一条线。

在这条竖线的顶端写上 \textbf{$z$}。

在交叉点旁边标一个 \textbf{$O$}（这是原点）。

\textbf{画完了！}

你现在应该有三条线，从同一个点 $O$ 出发：

\begin{itemize}
\item \textbf{$x$}：往右
\item \textbf{$y$}：往左
\item \textbf{$z$}：往上
\end{itemize}

$x$ 和 $y$ 这两条线像是铺在\textbf{地面}上的两个方向。

$z$ 是从地面\textbf{竖直往上}的方向，代表高度。

就像你低头看地板——地板上可以往左走，也可以往右走。

然后抬头看天花板——那个方向就是 $z$。

\subsection{为什么要画成这样？}

以前的平面坐标系，$x$ 往右，$y$ 往上，两条线是横平竖直的。

但那是因为只有两个方向，一张平纸就够了。

现在有三个方向。纸还是平的，画不下三条互相垂直的线。

怎么办？

\textbf{把地面画成斜的。}

$x$ 和 $y$ 斜着画，\textbf{假装}它们是铺在地上的。

$z$ 竖着画，代表高度。

这样，一张平纸上就能\textbf{假装}画出一个立体空间了。

\section{三个字母，三个任务}

现在，三个字母各管各的：

\begin{center}
\begin{tabular}{c|c|c}
字母 & 任务 & 在图上的方向 \\
\hline
$x$ & 第一条路 & 往右 \\
$y$ & 第二条路 & 往左 \\
$z$ & 高度 & 往上
\end{tabular}
\end{center}

请你仔细看这张表。

\textbf{上一章}：$y$ 是高度。

\textbf{这一章}：$y$ 是路，$z$ 才是高度。

$y$ 这个字母没变。

但它的\textbf{任务变了}。

这是很多人学双重积分时最困惑的地方。

“上一章 $y$ 是高度，这一章 $y$ 怎么变成路了？”

原因很简单：

上一章只有\textbf{一条路}（$x$），$y$ 没路可走，只能当高度。

这一章有\textbf{两条路}，$y$ 升职了，从“高度”变成了“第二条路”。

高度的活儿，交给了新人 $z$。

\section{怎么在这个坐标系上标点？}

\subsection{先说清楚：坐标怎么写？}

在这个立体坐标系里，一个点需要\textbf{三个数字}来确定位置。

写成 $(x, y, z)$。

\textbf{这三个数字分别是什么意思？}

\textbf{第一个数字 $x$}：这个点在 $x$ 轴方向的位置。也就是从原点出发，沿着 $x$ 轴走了多远。

\textbf{第二个数字 $y$}：这个点在 $y$ 轴方向的位置。也就是沿着 $y$ 轴走了多远。

\textbf{第三个数字 $z$}：这个点在 $z$ 轴方向的位置。也就是往上走了多高。

\textbf{顺序很重要！}

永远是：先 $x$，再 $y$，最后 $z$。

$(3, 1, 5)$ 的意思是：$x$ 方向走 3，$y$ 方向走 1，$z$ 方向走 5。

$(1, 3, 5)$ 就不一样了：$x$ 方向走 1，$y$ 方向走 3，$z$ 方向走 5。

\textbf{数字一样，位置不同，顺序搞反了就是完全不同的点！}

\subsection{在轴上标刻度}

在标点之前，先把三条轴上的刻度画好。

\textbf{在 $x$ 轴上标刻度：}

从原点 $O$ 开始，沿着 $x$ 轴方向（往右），每隔一小段距离（比如 1 厘米），做一个小记号。

第一个记号标 $1$，第二个记号标 $2$，第三个记号标 $3$……

\textbf{在 $y$ 轴上标刻度：}

从原点 $O$ 开始，沿着 $y$ 轴方向（往左），也是每隔一小段距离做记号。

每段距离要和 $x$ 轴上的\textbf{一样长}。

标上 $1$、$2$、$3$……

\textbf{在 $z$ 轴上标刻度：}

从原点 $O$ 开始，沿着 $z$ 轴方向（往上），每隔一小段距离做记号。

每段距离也要和 $x$ 轴上的\textbf{一样长}。

标上 $1$、$2$、$3$……

\subsection{标点的步骤}

\textbf{问题：怎么标出点 $(3, 1, 5)$？}

这个点的意思是：$x$ 方向走 3，$y$ 方向走 1，$z$ 方向走 5。

\textbf{第一步：从原点出发，沿 $x$ 轴方向走 3 个单位。}

在 $x$ 轴上找到刻度 $3$ 的位置。

暂停，在这里做个记号，或者用手指按住。

\textbf{第二步：从刚才的位置，沿 $y$ 轴方向走 1 个单位。}

注意！不是回到原点重新走。

是从你刚才停下的位置（$x = 3$ 那个点），沿着\textbf{和 $y$ 轴平行}的方向走。

$y$ 轴是往左的，所以你也往左走。

走多远？走 1 个单位（和 $y$ 轴上一个刻度的距离一样长）。

暂停，在这里做个记号。

\textbf{第三步：从刚才的位置，沿 $z$ 轴方向走 5 个单位。}

还是不要回原点！

从你刚才停下的位置，沿着\textbf{和 $z$ 轴平行}的方向走。

$z$ 轴是往上的，所以你也往上走。

走多远？走 5 个单位。

\textbf{到了！}

你现在停下的位置，就是点 $(3, 1, 5)$。

在这里标上 $(3, 1, 5)$。

\textbf{标点的口诀：先 $x$，再 $y$，最后 $z$。每次从上一步停下的地方继续走，不要回原点！}

\subsection{地面上的点（$z = 0$）}

\textbf{问题：怎么标出点 $(3, 2, 0)$？}

$z = 0$ 的意思是：不往上走。这个点贴着地面。

\textbf{第一步：} 从原点出发，沿 $x$ 轴方向走 \textbf{3} 个单位。

\textbf{第二步：} 从这里，沿 $y$ 轴方向走 \textbf{2} 个单位。

\textbf{第三步：} 沿 $z$ 轴方向走 \textbf{0} 个单位。（就是不走，原地不动。）

\textbf{到了！} 这就是点 $(3, 2, 0)$。

因为 $z = 0$，这个点就在地面上。

\subsection{简化写法}

如果一个点在地面上（$z = 0$），有时候我们会省略 $z$，只写 $(x, y)$。

比如 $(3, 2)$ 就是 $(3, 2, 0)$ 的简化写法。

\textbf{但要记住：只有地面上的点（$z = 0$）才能这样简化！}

如果 $z$ 不是 0，就必须写三个数字。

\section{画底面长方形}

现在我们要在地面上画一块长方形区域。

比如：$x$ 从 $0$ 到 $3$，$y$ 从 $0$ 到 $2$。

\subsection{找到四个角}

这个长方形有四个角，全都在地面上（$z = 0$）：

\begin{enumerate}
\item \textbf{第一个角}：$(0, 0, 0)$——就是原点 $O$
\item \textbf{第二个角}：$(3, 0, 0)$——在 $x$ 轴上，$x=3$ 那个刻度
\item \textbf{第三个角}：$(0, 2, 0)$——在 $y$ 轴上，$y=2$ 那个刻度
\item \textbf{第四个角}：$(3, 2, 0)$——刚才标的那个点
\end{enumerate}

\subsection{连起来}

\textbf{第一条边：} 从 $(0, 0, 0)$ 到 $(3, 0, 0)$，沿着 $x$ 轴画一条线段。

\textbf{第二条边：} 从 $(0, 0, 0)$ 到 $(0, 2, 0)$，沿着 $y$ 轴画一条线段。

\textbf{第三条边：} 从 $(3, 0, 0)$ 到 $(3, 2, 0)$。这条边和 $y$ 轴平行，往左画。

\textbf{第四条边：} 从 $(0, 2, 0)$ 到 $(3, 2, 0)$。这条边和 $x$ 轴平行，往右画。

\textbf{画完了！}

四条边围成了一个平行四边形。

（为什么看起来是平行四边形？因为纸是平的，我们在假装画立体图。其实在真实的三维空间里，这是一个正正方方的长方形。）

\section{画长方体}

现在，我们要从这个底面，往上长出一个长方体。

假设高度是 $1$（也就是 $z$ 方向长 $1$）。

\subsection{从四个角往上画竖线}

\textbf{第一根竖线：} 从 $(0, 0, 0)$ 往正上方画，长度 $1$。顶端是点 $(0, 0, 1)$。

\textbf{第二根竖线：} 从 $(3, 0, 0)$ 往正上方画，长度 $1$。顶端是点 $(3, 0, 1)$。

\textbf{第三根竖线：} 从 $(0, 2, 0)$ 往正上方画，长度 $1$。顶端是点 $(0, 2, 1)$。

\textbf{第四根竖线：} 从 $(3, 2, 0)$ 往正上方画，长度 $1$。顶端是点 $(3, 2, 1)$。

四根竖线的\textbf{长度要一样}，方向也要一样（都是竖直往上）。

\subsection{连顶面}

把四根竖线的顶端连起来，顶面也是一个平行四边形，和底面形状一样。

\textbf{长方体画完了！}

\section{那问题来了}

这个长方体的\textbf{体积}是多少？

如果高度一直是 $1$（就是我们刚才画的），那很简单。

底面的 $x$ 方向长度 $= 3$

底面的 $y$ 方向长度 $= 2$

高度 $= 1$

\textbf{体积 = 长 $\times$ 宽 $\times$ 高 = $3 \times 2 \times 1 = 6$}

但如果高度不是固定的呢？

如果从底面的每一个点 $(x, y, 0)$，往上长的高度都不一样呢？

高度由一个函数 $z = f(x, y)$ 决定。

有些地方高，有些地方矮。

顶面不再是平的，而是\textbf{弯弯曲曲的曲面}。

底面和曲面之间围起来的空间，形状就不规则了。

\textbf{这个不规则立体的体积怎么算？}

\textbf{老办法：切！}

\section{怎么切？}

上一章，我们沿着\textbf{一个方向}切，把面积切成了一根根竖条。

这一章，我们有\textbf{两个方向}。

所以要切\textbf{两次}。

\subsection{第一刀：沿 $x$ 方向切}

想象你拿着一把很薄很薄的刀片。

刀片的方向：和 $y$ 轴平行，竖直立着。

从左到右，每隔一小段距离 $dx$，切一刀。

切完之后，整个立体被切成了一片一片的。

每一片在 $x$ 方向的厚度是 $dx$。

\subsection{第二刀：沿 $y$ 方向切}

现在换个方向。

刀片的方向：和 $x$ 轴平行，竖直立着。

从前到后，每隔一小段距离 $dy$，切一刀。

切完之后，每一片又被切成了一根根小柱子。

\subsection{切完之后}

两刀切完，整个立体变成了\textbf{无数个小柱子}。

每根小柱子的底面是一个小长方形：

\begin{itemize}
\item $x$ 方向的边长：$dx$
\item $y$ 方向的边长：$dy$
\item 底面积 = $dx \times dy$
\end{itemize}

每根小柱子的高度呢？

就是这个位置 $(x, y)$ 的高度，也就是 $f(x, y)$。

\subsection{为什么小柱子可以当成长方体？}

你可能会问：顶面不是弯曲的吗？切出来的小柱子顶部不应该也是弯的吗？

确实，如果切得不够细，小柱子的顶部还是弯的。

但是，记得吗？我们切到\textbf{无限细}。

$dx$ 无限小，$dy$ 也无限小。

每根小柱子的底面\textbf{无限小}。

想象一下。

一个弯弯曲曲的山坡。

你站得很远，能看出它是弯的。

但如果你只看脚下很小很小的一块地——比如一平方厘米——那块地几乎是\textbf{平的}。

曲面也一样。

当你看的范围\textbf{无限小}时，曲面在那一小块上几乎是\textbf{平的}。

所以，每根小柱子的顶部几乎是平的。

\textbf{切得无限细时，每根小柱子可以看成一个小长方体。}

\subsection{小长方体的体积}

每个小长方体的体积是什么？

\textbf{体积 = 底面积 $\times$ 高}

底面积 = $dx \times dy$

高 = $f(x, y)$

所以：

\textbf{每个小长方体的体积 = $f(x, y) \, dx \, dy$}

\textbf{把所有小长方体的体积加起来，就是整个立体的总体积！}

这就是\textbf{双重积分}。

\section{双重积分的符号}

\[
\iint f(x, y) \, dy \, dx
\]

你看，有\textbf{两个} $\int$。

因为有\textbf{两条路}，每条路加一次。

两条路，两次加法，两个 $\int$。

写完整一点，带上起点和终点：

\[
\int_a^b \int_c^d f(x, y) \, dy \, dx
\]

一个一个翻译。

\subsection{\texorpdfstring{外层：$\int_a^b \cdots \, dx$}{外层：int_a^b ... dx}}

这是\textbf{第一条路}。

沿 $x$ 方向，从 $x = a$ 走到 $x = b$。

（$a$ 和 $b$ 告诉你：$x$ 这条路从哪开始，到哪结束。）

\subsection{\texorpdfstring{内层：$\int_c^d \cdots \, dy$}{内层：int_c^d ... dy}}

这是\textbf{第二条路}。

沿 $y$ 方向，从 $y = c$ 走到 $y = d$。

（$c$ 和 $d$ 告诉你：$y$ 这条路从哪开始，到哪结束。）

\subsection{$f(x, y)$}

这是每个小长方体的\textbf{高度}。

上一章里，高度只看 $x$ 就够了，写成 $f(x)$。

这一章，你站的位置由 $x$ 和 $y$ \textbf{两个东西}决定。

所以高度要看 $x$ 和 $y$ 才能知道，写成 $f(x, y)$。

\subsection{$dy \, dx$}

$dy$ 是小长方体底面在 $y$ 方向的边长。

$dx$ 是小长方体底面在 $x$ 方向的边长。

$dy \times dx$ 是小长方体的\textbf{底面积}。

底面积 $\times$ 高度 = 小长方体的体积。

\subsection{完整翻译}

\[
\int_a^b \int_c^d f(x, y) \, dy \, dx
\]

\begin{quote}
“沿 $x$ 方向从 $a$ 走到 $b$，沿 $y$ 方向从 $c$ 走到 $d$，
每个位置 $(x, y)$ 有一个小长方体，
它的体积 = 高度 $f(x, y)$ $\times$ 底面积 $dy \, dx$，
把所有小长方体的体积，\textbf{全部加起来}。”
\end{quote}

\section{几个需要注意的事情}

\subsection{能不能反过来算？}

能！

$\int_a^b \int_c^d f(x, y) \, dy \, dx$ 和 $\int_c^d \int_a^b f(x, y) \, dx \, dy$

结果是\textbf{一样的}。

先走 $y$ 再走 $x$，和先走 $x$ 再走 $y$，最后加起来的总体积不变。

（就像你数苹果，先数第一排再数第二排，和先数第一列再数第二列，总数一样。）

\subsection{上下限反过来呢？}

如果写成 $\int_3^0$ 而不是 $\int_0^3$，意思是\textbf{倒着走}。

结果会\textbf{变成负数}。

通常情况下，让下限比上限小，这样结果是正数。

\subsection{$f(x, y)$ 是负数呢？}

如果高度是负数，意思是往\textbf{下}长，到地面下面去了。

这时候算出来的体积会是负数，表示在地面下方的体积。

通常情况下，我们考虑的是正的高度，地面上方的体积。

\section{怎么算？一次走一条路}

你可能会想：无限多个小长方体，怎么加？

关键技巧：\textbf{不要试图一次全加完。}

\textbf{先沿一条路走完，再沿另一条路走。}

分两步。

\subsection{为什么可以分两步？}

想象你在超市买了一堆苹果。

你要算总价。

\textbf{方法一：} 一个一个数，每个苹果 5 元。

第一个 5 元，第二个 5 元，第三个 5 元……一直加。

\textbf{方法二：} 先数有多少个苹果，比如 10 个。

然后算 $5 \times 10 = 50$ 元。

双重积分也是一样。

\textbf{方法一（一次全算）：} 把所有小长方体一个一个加。太累了。

\textbf{方法二（分两步）：} 先沿一个方向加一遍，得到一个中间结果。再沿另一个方向加一遍。

\subsection{第一步：先走里面那条路}

看符号：

\[
\int_a^b \int_c^d f(x, y) \, dy \, dx
\]

\textbf{里面}那层是 $\int_c^d f(x, y) \, dy$。

这一步，沿 $y$ 方向走。

但是，$x$ 怎么办？

\textbf{$x$ 不动。}

假装 $x$ 站在某个固定的位置，比如 $x = 1$。

你只沿着 $y$，从 $y = c$ 走到 $y = d$。

一边走，一边把这条路上所有小长方体加起来。

\textbf{为什么 $x$ 可以不动？}

因为我们在算“当 $x$ 固定在某个值时，沿 $y$ 方向的总和”。

就像切蛋糕，先切一刀，拿出一片。

在这一片里，$x$ 是固定的（这一片的位置不变）。

你只需要看这一片里，$y$ 方向有多少。

这就是一个\textbf{普通的单重积分}！

只不过式子里有个 $x$。

但 $x$ 现在是固定的，它就像一个\textbf{普通的数字}，不用管它。

走完之后，$y$ 就消失了。

\textbf{为什么 $y$ 会消失？}

因为 $y$ 从 $c$ 走到 $d$，这条路\textbf{走完了}。

就像上一章里，$x$ 从 $a$ 加到 $b$，加完之后 $x$ 的任务完成了，$x$ 消失了。

$y$ 也是一样。

路走完了，就用完了，没了。

剩下的结果里，只剩 $x$。

（因为 $x$ 还没走呢，它还站在那里等着。）

\subsection{第二步：再走外面那条路}

$y$ 消失了。

你手里有一个只跟 $x$ 有关的东西（第一步算出来的结果）。

现在沿 $x$ 方向，从 $x = a$ 走到 $x = b$，再加一次。

\[
\int_a^b \text{（第一步的结果）} \, dx
\]

这又是一个\textbf{普通的单重积分}。

走完之后，$x$ 也消失了。

\textbf{两条路都走完了，两个变量都用完了。}

剩下的，是一个\textbf{数字}。

这个数字就是总体积。

\section{三个例子}

\subsection{例子一：高度永远是 1}

\[
\int_0^3 \int_0^2 1 \, dy \, dx
\]

\textbf{翻译：}

\begin{itemize}
\item 外层 $\int_0^3 \cdots \, dx$：沿 $x$ 方向，从 $x = 0$ 走到 $x = 3$
\item 内层 $\int_0^2 \cdots \, dy$：沿 $y$ 方向，从 $y = 0$ 走到 $y = 2$
\item $1$：每个小长方体的高度永远是 1，不管你站在哪里
\end{itemize}

\textbf{第一步：先走 $y$ 这条路。}

$x$ 不动。只看 $y$。

\[
\int_0^2 1 \, dy
\]

这是什么意思？

从 $y = 0$ 到 $y = 2$，每一小块的高度是 $1$，每一小块的宽度是 $dy$。

全部加起来：

\[
\int_0^2 1 \, dy = 1 \times (2 - 0) = 1 \times 2 = 2
\]

（快速计算规律：$\int$ 上面的数减下面的数 = 总长度。常数乘以总长度就是结果。）

$y$ 走完了，消失了。剩下一个数字：$2$。

\textbf{第二步：再走 $x$ 这条路。}

\[
\int_0^3 2 \, dx
\]

从 $x=0$ 到 $x=3$。

\[
= 2 \times (3 - 0) = 2 \times 3 = 6
\]

$x$ 也走完了，也消失了。

\textbf{答案：$\int_0^3 \int_0^2 1 \, dy \, dx = 6$}

\textbf{验证：}

$x$ 从 0 到 3，长度是 3。

$y$ 从 0 到 2，长度是 2。

高度一直是 1。

这就是一个\textbf{长方体}。

长方体的体积 = 长 $\times$ 宽 $\times$ 高 = $3 \times 2 \times 1 = 6$ ✓

\subsection{例子二：高度永远是 5}

\[
\int_0^4 \int_0^3 5 \, dy \, dx
\]

\textbf{翻译：}

\begin{itemize}
\item 沿 $x$ 方向：从 $0$ 到 $4$
\item 沿 $y$ 方向：从 $0$ 到 $3$
\item 高度永远是 $5$
\end{itemize}

\textbf{第一步：走 $y$。}

\[
\int_0^3 5 \, dy = 5 \times (3 - 0) = 5 \times 3 = 15
\]

$y$ 消失了，剩下 $15$。

\textbf{第二步：走 $x$。}

\[
\int_0^4 15 \, dx = 15 \times (4 - 0) = 15 \times 4 = 60
\]

$x$ 也消失了。

\textbf{答案：$\int_0^4 \int_0^3 5 \, dy \, dx = 60$}

\textbf{验证：}

长方体，长 4，宽 3，高 5。

体积 = $4 \times 3 \times 5 = 60$ ✓

\subsection{例子三：起点不是 0}

\[
\int_1^4 \int_2^5 3 \, dy \, dx
\]

\textbf{翻译：}

\begin{itemize}
\item 沿 $x$ 方向：从 $1$ 到 $4$（注意！不是从 0 开始）
\item 沿 $y$ 方向：从 $2$ 到 $5$（也不是从 0 开始！）
\item 高度永远是 $3$
\end{itemize}

\textbf{第一步：走 $y$。}

\[
\int_2^5 3 \, dy = 3 \times (5 - 2) = 3 \times 3 = 9
\]

（不是 $3 \times 5$。是 $3 \times (5 - 2)$。就像从 2 楼走到 5 楼，不是走了 5 层，是走了 $5-2=3$ 层。）

$y$ 消失了，剩下 $9$。

\textbf{第二步：走 $x$。}

\[
\int_1^4 9 \, dx = 9 \times (4 - 1) = 9 \times 3 = 27
\]

\textbf{答案：$\int_1^4 \int_2^5 3 \, dy \, dx = 27$}

\textbf{验证：}

$x$ 从 1 到 4，长度是 $4 - 1 = 3$。

$y$ 从 2 到 5，长度是 $5 - 2 = 3$。

高度是 3。

长方体体积 = $3 \times 3 \times 3 = 27$ ✓

\section{一句话总结}

\textbf{单重积分：一条路，切一次，加一次。算面积。}

\textbf{双重积分：两条路，切两次，加两次。算体积。}

\[
\int_a^b \int_c^d f(x,y) \, dy \, dx
\]

\begin{quote}
“先走 $y$（里面），$y$ 走完消失。再走 $x$（外面），$x$ 也走完消失。剩下一个数字，就是体积。”
\end{quote}

$y$ 不是空气。

上一章，它藏在 $f(x)$ 里当高度。

这一章，它站出来当第二条路。

\textbf{名字没变，任务变了。}

这三个例子里，高度都是一个固定的数字（1、5、3）。

所以每次都切出来一个规规矩矩的长方体。

但如果高度会变呢？

\textbf{如果 $f(x, y)$ 不是一个数字，而是一个真正的函数呢？}

比如 $f(x, y) = x + y$，或者 $f(x, y) = x^2 + y^2$。

那就不能用长方体公式验证了。

必须\textbf{老老实实地积分}。

下一章，我们来试试。

\chapter{三重积分——全都是“方向”}

上一章，我们学了双重积分。

还记得发生了什么吗？

$y$ 升职了。

它不再当高度，变成了\textbf{第二条路}。

高度的工作，交给了新人 $z$。

\begin{center}
\begin{tabular}{c|c|c}
 & 单重积分 & 双重积分 \\
\hline
\textbf{路} & $x$ & $x$ 和 $y$ \\
\textbf{高度} & $y$（藏在 $f(x)$ 里） & $z$（藏在 $f(x,y)$ 里） \\
\textbf{算的是} & 面积 & 体积
\end{tabular}
\end{center}

现在我要问你一个问题。

$y$ 能升职，$z$ 为什么不能？

\textbf{如果 $z$ 也不想当高度了呢？}

\textbf{如果 $z$ 也想当路呢？}

\section{\texorpdfstring{$z$ 也升职了}{z 也升职了}}

$z$ 不当高度了。

$z$ 也变成了一条\textbf{路}。

现在三个字母全是路：

\begin{center}
\begin{tabular}{c|c|c}
字母 & 上一章的任务 & 这一章的任务 \\
\hline
$x$ & 第一条路 & 第一条路（没变） \\
$y$ & 第二条路 & 第二条路（没变） \\
$z$ & 高度 & \textbf{第三条路}（升职了！）
\end{tabular}
\end{center}

三条路，三个方向。

$x$ 往右。

$y$ 往左。

$z$ 往上。

\textbf{全都是方向，没有人当高度了。}

\section{等一下，那高度谁来管？}

好问题。

单重积分里，$y$ 当高度。

双重积分里，$z$ 当高度。

三重积分里，$x$、$y$、$z$ 全去当路了。

\textbf{高度呢？}

答案是：\textbf{不需要“高度”了。}

慢慢解释。

单重积分的时候，我们有一条路（$x$），路上每个位置有个高度（$y$）。

高度和路围成了一个\textbf{面}。面有\textbf{面积}。

双重积分的时候，我们有两条路（$x$ 和 $y$），每个位置有个高度（$z$）。

高度和底面围成了一个\textbf{立体}。立体有\textbf{体积}。

三重积分的时候，我们有三条路（$x$、$y$、$z$），三条路本身就围成了一个\textbf{立体空间}。

不需要额外的高度来“撑起”一个立体了。

三个方向本身就是立体的。

\textbf{那 $f(x, y, z)$ 是什么？}

它不是“高度”了。

它变成了：\textbf{空间里每个点的“浓度”。}

\subsection{“浓度”是什么意思？}

想象一个房间。

房间里每个位置都有空气。

但有些位置的空气\textbf{浓}一点（比如窗户旁边），有些位置的空气\textbf{淡}一点。

$f(x, y, z)$ 就是在告诉你：在位置 $(x, y, z)$ 这个点，浓度是多少。

\textbf{如果 $f(x, y, z) = 1$ 呢？}

每个点的浓度都是 1，到处一样。

那把所有“浓度 $\times$ 小块体积”加起来，就等于\textbf{整个空间的体积}。

\textbf{如果 $f(x, y, z) = 5$ 呢？}

每个点的浓度都是 5。

那结果就是\textbf{体积 $\times 5$}。

所以，三重积分算的不再是“体积”这一种东西了。

它算的是\textbf{“浓度在空间里的累积”}。

当浓度恰好是 1 的时候，累积的结果就是体积。

\textbf{通常情况下}，你见到的三重积分，$f(x, y, z)$ 要么是个常数（比如 1、3、5），要么是一个简单的函数（比如 $x + y + z$）。

\textbf{如果 $f(x, y, z) = 1$}，三重积分就是在算体积。这是\textbf{最常见}的情况。

\textbf{如果 $f(x, y, z)$ 是别的数}，那就是在算“加权体积”——每个小块不光算体积，还要乘上那个点的浓度。

\section{老办法：切！}

三个方向，就切\textbf{三刀}。

\subsection{第一刀：沿 $x$ 方向切}

从左到右，每隔 $dx$ 切一刀。

整个立体被切成了一片一片的薄片。

\subsection{第二刀：沿 $y$ 方向切}

换个方向，每隔 $dy$ 切一刀。

每片薄片又被切成一根根细条。

\subsection{第三刀：沿 $z$ 方向切}

再换个方向，从下到上，每隔 $dz$ 切一刀。

每根细条又被切成一块块\textbf{小方块}。

\subsection{切完之后}

三刀切完，整个立体变成了\textbf{无数个小方块}。

每个小方块的三个尺寸：

\begin{itemize}
\item $x$ 方向的边长：$dx$
\item $y$ 方向的边长：$dy$
\item $z$ 方向的边长：$dz$
\end{itemize}

\textbf{小方块的体积 = $dx \times dy \times dz$}

每个小方块的浓度呢？

就是 $f(x, y, z)$——由这个小方块所在的位置决定。

\textbf{每个小方块的贡献 = 浓度 $\times$ 体积 = $f(x, y, z) \times dx \times dy \times dz$}

\textbf{把所有小方块的贡献加起来，就是三重积分！}

\section{三重积分的符号}

\[
\iiint f(x, y, z) \, dz \, dy \, dx
\]

三个 $\int$。

三条路，三次加法，三个 $\int$。

写完整一点，带上起点和终点：

\[
\int_a^b \int_c^d \int_e^f f(x, y, z) \, dz \, dy \, dx
\]

字母有点多，别怕，一个一个翻译。

\subsection{\texorpdfstring{最外层：$\int_a^b \cdots \, dx$}{最外层：int_a^b ... dx}}

\textbf{第一条路。} 沿 $x$ 方向，从 $x = a$ 走到 $x = b$。

\subsection{\texorpdfstring{中间层：$\int_c^d \cdots \, dy$}{中间层：int_c^d ... dy}}

\textbf{第二条路。} 沿 $y$ 方向，从 $y = c$ 走到 $y = d$。

\subsection{\texorpdfstring{最里层：$\int_e^f \cdots \, dz$}{最里层：int_e^f ... dz}}

\textbf{第三条路。} 沿 $z$ 方向，从 $z = e$ 走到 $z = f$。

\subsection{\texorpdfstring{$f(x, y, z)$}{f(x, y, z)}}

每个小方块的\textbf{浓度}。

由位置 $(x, y, z)$ 决定。

如果浓度是常数（比如 1），就是在算纯体积。

\subsection{\texorpdfstring{$dz \, dy \, dx$}{dz dy dx}}

$dz$ 是小方块在 $z$ 方向的边长。

$dy$ 是小方块在 $y$ 方向的边长。

$dx$ 是小方块在 $x$ 方向的边长。

$dz \times dy \times dx$ = 小方块的\textbf{体积}。

\subsection{完整翻译}

\[
\int_a^b \int_c^d \int_e^f f(x, y, z) \, dz \, dy \, dx
\]

\begin{quote}
“沿 $x$ 方向从 $a$ 走到 $b$，
沿 $y$ 方向从 $c$ 走到 $d$，
沿 $z$ 方向从 $e$ 走到 $f$，
每个位置 $(x, y, z)$ 有一个小方块，
它的贡献 = 浓度 $f(x, y, z)$ $\times$ 体积 $dz \, dy \, dx$，
把所有小方块的贡献，\textbf{全部加起来}。”
\end{quote}

\section{怎么算？还是一次走一条路}

三条路，分\textbf{三步}走。

和双重积分的思路完全一样，只不过多了一步。

\subsection{第一步：先走最里面那条路（$z$）}

看符号：

\[
\int_a^b \int_c^d \int_e^f f(x, y, z) \, dz \, dy \, dx
\]

最里面那层是 $\int_e^f f(x, y, z) \, dz$。

这一步，沿 $z$ 方向走。

\textbf{$x$ 不动，$y$ 也不动。}

假装 $x$ 和 $y$ 都站在某个固定的位置。

你只沿着 $z$，从 $z = e$ 走到 $z = f$，一边走一边加。

走完之后，$z$ 消失了。

（路走完了，就用完了，没了。还记得吗？）

剩下的结果里，只有 $x$ 和 $y$。

\subsection{第二步：走中间那条路（$y$）}

$z$ 消失了。

你手里有一个只跟 $x$ 和 $y$ 有关的东西。

现在沿 $y$ 方向，从 $y = c$ 走到 $y = d$，再加一次。

走完之后，$y$ 也消失了。

剩下的结果里，只有 $x$。

\subsection{第三步：走最外面那条路（$x$）}

$y$ 也消失了。

你手里只剩一个跟 $x$ 有关的东西。

沿 $x$ 方向，从 $x = a$ 走到 $x = b$，最后一次加。

走完之后，$x$ 也消失了。

\textbf{三条路全走完了，三个变量全消失了。}

剩下的，是一个\textbf{数字}。

\section{几个需要注意的事情}

\subsection{顺序能换吗？}

能！

先走 $z$ 再走 $y$ 再走 $x$，和先走 $x$ 再走 $y$ 再走 $z$，结果\textbf{一样}。

（就像你数一箱苹果，先数行再数列再数层，和先数层再数行再数列，总数不变。）

\subsection{\texorpdfstring{$f(x, y, z) = 1$ 的时候}{f(x, y, z) = 1 的时候}}

这是最常见的特殊情况。

当浓度恒等于 1 时，三重积分算的就是\textbf{纯体积}。

每个小方块的贡献 = $1 \times dx \times dy \times dz$ = 就是它自身的体积。

全部加起来 = 总体积。

\section{三个例子}

\subsection{例子一：浓度永远是 1}

\[
\int_0^2 \int_0^3 \int_0^4 1 \, dz \, dy \, dx
\]

\textbf{翻译：}

\begin{itemize}
\item 最外层 $\int_0^2 \cdots \, dx$：沿 $x$ 方向，从 $0$ 到 $2$
\item 中间层 $\int_0^3 \cdots \, dy$：沿 $y$ 方向，从 $0$ 到 $3$
\item 最里层 $\int_0^4 \cdots \, dz$：沿 $z$ 方向，从 $0$ 到 $4$
\item $1$：浓度永远是 1
\end{itemize}

\textbf{第一步：走 $z$。}

$x$ 不动，$y$ 也不动。只看 $z$。

\[
\int_0^4 1 \, dz
\]

每一小块 = $1 \times dz$。

全部加起来：

\[
= 1 \times (dz + dz + dz + \cdots)
\]

用规律：$\int$ 上面的数减下面的数 = 总长度。

总长度 = $4 - 0 = 4$。

\[
= 1 \times 4 = 4
\]

$z$ 走完了，消失了。剩下 $4$。

\textbf{第二步：走 $y$。}

\[
\int_0^3 4 \, dy
\]

（现在的数字是 $4$，这是第一步的结果。）

\[
= 4 \times (dy + dy + dy + \cdots)
\]

总长度 = $3 - 0 = 3$。

\[
= 4 \times 3 = 12
\]

$y$ 消失了。剩下 $12$。

\textbf{第三步：走 $x$。}

\[
\int_0^2 12 \, dx
\]

\[
= 12 \times (dx + dx + dx + \cdots)
\]

总长度 = $2 - 0 = 2$。

\[
= 12 \times 2 = 24
\]

$x$ 也消失了。

\textbf{答案：$\int_0^2 \int_0^3 \int_0^4 1 \, dz \, dy \, dx = 24$}

\textbf{验证：}

$x$ 从 0 到 2，长度 2。

$y$ 从 0 到 3，长度 3。

$z$ 从 0 到 4，长度 4。

浓度是 1，所以三重积分就是纯体积。

这是一个长方体，体积 = $2 \times 3 \times 4 = 24$ ✓

\subsection{例子二：浓度永远是 5}

\[
\int_0^2 \int_0^3 \int_0^4 5 \, dz \, dy \, dx
\]

\textbf{翻译：}

和例子一的空间一样（$x$ 从 0 到 2，$y$ 从 0 到 3，$z$ 从 0 到 4）。

但浓度不是 1 了，是 \textbf{5}。

\textbf{第一步：走 $z$。}

\[
\int_0^4 5 \, dz = 5 \times (4 - 0) = 5 \times 4 = 20
\]

$z$ 消失了。剩下 $20$。

\textbf{第二步：走 $y$。}

\[
\int_0^3 20 \, dy = 20 \times (3 - 0) = 20 \times 3 = 60
\]

$y$ 消失了。剩下 $60$。

\textbf{第三步：走 $x$。}

\[
\int_0^2 60 \, dx = 60 \times (2 - 0) = 60 \times 2 = 120
\]

\textbf{答案：$\int_0^2 \int_0^3 \int_0^4 5 \, dz \, dy \, dx = 120$}

\textbf{验证：}

体积 = $2 \times 3 \times 4 = 24$。

浓度是 5。

$24 \times 5 = 120$ ✓

\subsection{例子三：起点不是 0}

\[
\int_1^3 \int_2^4 \int_1^6 2 \, dz \, dy \, dx
\]

\textbf{翻译：}

\begin{itemize}
\item 沿 $x$ 方向：从 $1$ 到 $3$
\item 沿 $y$ 方向：从 $2$ 到 $4$
\item 沿 $z$ 方向：从 $1$ 到 $6$
\item 浓度永远是 $2$
\end{itemize}

\textbf{第一步：走 $z$。}

\[
\int_1^6 2 \, dz = 2 \times (6 - 1) = 2 \times 5 = 10
\]

（不是 $2 \times 6$。是 $2 \times (6 - 1)$。就像从 1 楼走到 6 楼，走了 $6-1=5$ 层。）

$z$ 消失了。剩下 $10$。

\textbf{第二步：走 $y$。}

\[
\int_2^4 10 \, dy = 10 \times (4 - 2) = 10 \times 2 = 20
\]

$y$ 消失了。剩下 $20$。

\textbf{第三步：走 $x$。}

\[
\int_1^3 20 \, dx = 20 \times (3 - 1) = 20 \times 2 = 40
\]

\textbf{答案：$\int_1^3 \int_2^4 \int_1^6 2 \, dz \, dy \, dx = 40$}

\textbf{验证：}

$x$ 方向长度 = $3 - 1 = 2$。

$y$ 方向长度 = $4 - 2 = 2$。

$z$ 方向长度 = $6 - 1 = 5$。

体积 = $2 \times 2 \times 5 = 20$。

浓度是 2。

$20 \times 2 = 40$ ✓

\section{回头看：一个完整的规律}

\begin{center}
\begin{tabular}{c|c|c|c}
 & 单重积分 & 双重积分 & 三重积分 \\
\hline
\textbf{几条路} & 1 条（$x$） & 2 条（$x, y$） & 3 条（$x, y, z$） \\
\textbf{几个 $\int$} & 1 个 & 2 个 & 3 个 \\
\textbf{切几次} & 1 次 & 2 次 & 3 次 \\
\textbf{切成什么} & 竖条 & 小柱子 & 小方块 \\
\textbf{$f$ 是什么} & 高度 $f(x)$ & 高度 $f(x,y)$ & 浓度 $f(x,y,z)$ \\
\textbf{算什么} & 面积 & 体积 & 浓度的累积
\end{tabular}
\end{center}

\textbf{规律：}

\textbf{几条路，就几个 $\int$，就切几次，就分几步算。}

\textbf{每走完一条路，那个变量就消失。}

\textbf{全走完，剩一个数字。}

\section{一句话总结}

\textbf{三重积分：三条路，切三次，加三次。}

\[
\int_a^b \int_c^d \int_e^f f(x,y,z) \, dz \, dy \, dx
\]

\begin{quote}
“先走 $z$（最里面），$z$ 消失。
再走 $y$（中间），$y$ 消失。
最后走 $x$（最外面），$x$ 也消失。
剩下一个数字。”
\end{quote}

\textbf{快速计算规律：每一层，$\int$ 上面的数减下面的数 = 总长度。}

$x$、$y$、$z$，全都是方向了。

没有谁再当“高度”了。

\textbf{三个人全去走路了，没人留守。}

$f(x, y, z)$ 接管了“每个点有多少东西”的任务。

\section{下一章预告：假如 $f$ 也会走路}

这三章里，$f$ 一直是个常数。

1、3、5……不管你站在哪里，它都不变。

它就像一个钉在原地的路灯，永远那么亮。

但如果 $f$ 也开始动了呢？

你往右走一步，$f$ 变大了。

你往左走一步，$f$ 又变小了。

$f$ 不再是一个固定的数字，它变成了一个\textbf{函数}。

比如 $f(x) = x$，你走到哪，$f$ 就跟到哪，而且它的大小跟着你的位置变。

比如 $f(x, y) = x + y$，你走两条路，$f$ 就看着两条路一起变。

这时候，每根竖条的高度不一样了，每个小方块的浓度也不一样了。

\textbf{$f$ 也会走路了。}

下一章见！

\chapter{\texorpdfstring{假如 $f$ 也会走路}{假如 f 也会走路}}

\section{📌 一个小小的题外话}

在学数学的路上，你会遇到很多新名字。

有些名字听起来很吓人。

比如今天，你会遇到一个叫\textbf{“原函数”}的东西。

你可能一下子记不住它是什么意思。

\textbf{没关系。}

其实啊，很多时候，名字变了，但东西还是那个东西。

就像薯条，有人叫它“炸土豆条”，有人叫它“薯条”，还有人叫它“French fries”。

名字不一样，但吃起来都一样香，蘸番茄酱都一样好吃。

就像你的好朋友，在学校老师叫他“张小明”，回家妈妈叫他“明明”，你叫他“小明”。

三个名字，同一个人。他还是那个会跟你分享零食、一起打游戏的好朋友。

数学也一样。

\textbf{名字只是名字。}

你暂时不记得名字叫什么？没事。

你暂时不理解为什么要这样做？也没事。

\textbf{只要你会算，会做，就好。}

先学会怎么做，理解会慢慢跟上来的。

就像你学骑自行车的时候，不需要先学什么“力矩平衡”“重心原理”。

你先骑上去，摔几次，骑着骑着就会了。

等你会骑了，再回头看，哦，原来是这么回事。

数学也一样。

\textbf{先会做，再慢慢懂。}

所以今天，如果有什么地方你觉得“我好像没完全懂”，别慌。

\textbf{会做就好。懂，会跟上来的。}

好，题外话结束。我们开始。

\section{\texorpdfstring{$f$ 站着不动的时候}{f 站着不动的时候}}

前面三章里，$f$ 一直是个常数。

$f = 1$，$f = 3$，$f = 5$……

不管你站在哪个位置，$f$ 永远是同一个数字。

它\textbf{不动}。

那时候，算积分很简单。

\[
\int_0^4 3 \, dx
\]

翻译一下：从 $x=0$ 到 $x=4$，每根竖条的高度都是 $3$。

每根竖条的高度一样！

所以把 $3$ 提出来，乘以总长度就行了。

总长度是多少？$\int$ 上面的数减下面的数 = $4 - 0 = 4$。

\[
\int_0^4 3 \, dx = 3 \times 4 = 12
\]

这是老方法。\textbf{提出来，乘以长度。}

\begin{quote}
\textbf{📌 通俗理解：老方法}

就像买苹果。

每个苹果都是 3 块钱（价格不变），你买了 4 个。

总价 = 3 × 4 = 12 块。

价格一样，直接乘就行。
\end{quote}

\section{\texorpdfstring{$f$ 开始走了}{f 开始走了}}

但如果 $f$ 不是常数呢？

如果 $f$ 是 $x$ 呢？

\[
\int_0^4 x \, dx
\]

翻译一下：从 $x=0$ 到 $x=4$，每根竖条的高度是 $x$。

什么意思？

意思是，\textbf{高度跟着你的位置变}。

你站在 $x = 0$ 的地方，高度是多少？是 $0$。

你往右走一步，站在 $x = 1$ 的地方，高度是多少？是 $1$。

再往右走一步，站在 $x = 2$ 的地方，高度是多少？是 $2$。

再走，$x = 3$，高度是 $3$。

再走，$x = 4$，高度是 $4$。

\textbf{看到了吗？}

$f$ 不再站着不动了。

\textbf{它跟着 $x$ 一起走。}

你往右走一步，$f$ 就变大一点。

你走到哪，$f$ 就跟到哪。

\begin{quote}
\textbf{📌 通俗理解：$f$ 跟着走}

就像买苹果，但这次每个苹果的价格不一样。

第 1 个苹果 1 块钱，第 2 个苹果 2 块钱，第 3 个苹果 3 块钱……

价格一直在变，不能直接乘了。

得一个一个加：$1 + 2 + 3 + 4 + \cdots$
\end{quote}

这时候，你还能把 $f$ 提出来吗？

试试看：

\[
\int_0^4 x \, dx = x \times (4 - 0) = x \times 4 = 4x
\]

等等。

$4x$？

答案应该是一个\textbf{数字}啊。

怎么还有 $x$ 在里面？

\textbf{不对。$x$ 不是常数，不能提出来。}

以前 $f = 3$ 的时候，每根竖条的高度都是 $3$，都一样，所以能提出来。

现在每根竖条的高度都不一样——第一根高度是 $0$，第二根高度是 $1$，第三根高度是 $2$……

不一样的东西，怎么能提出来？

\textbf{老方法失效了。}

\section{我们需要一个新方法}

老方法只适合 $f$ 不动的时候。

$f$ 一旦开始走，老方法就不管用了。

我们需要一个新方法。

这个新方法，和\textbf{求导}有关。

求导是什么？

\textbf{切一刀，看看变成什么。}

\begin{quote}
\textbf{📌 通俗理解：求导}

求导就像切苹果。

你有一个完整的苹果，切一刀，变成了苹果片。

苹果 $\rightarrow$ 切一刀 $\rightarrow$ 苹果片

在数学里：$x^2$ $\rightarrow$ 切一刀 $\rightarrow$ $2x$
\end{quote}

让我们复习一下求导。

\section{复习：切一刀}

\textbf{例子1：把 $x^2$ 切一刀。}

\[
(x^2)' = 2x
\]

$x^2$ 切一刀，变成了 $2x$。

\textbf{例子2：把 $x^3$ 切一刀。}

\[
(x^3)' = 3x^2
\]

$x^3$ 切一刀，变成了 $3x^2$。

\textbf{例子3：把 $5x$ 切一刀。}

\[
(5x)' = 5
\]

$5x$ 切一刀，变成了 $5$。

好，复习完了。

现在，我要问你一个\textbf{反过来}的问题。

\section{反过来的问题}

以前我们问：把 $x^2$ 切一刀，变成什么？

答案是 $2x$。

现在我反过来问：

\begin{quote}
\textbf{谁，切一刀之后，会变成 $x$？}
\end{quote}

不是问“$x$ 切一刀变成什么”。

是问“\textbf{谁}切一刀会变成 $x$”。

\begin{quote}
\textbf{📌 通俗理解：反过来的问题}

以前问：苹果切一刀，变成什么？答：苹果片。

现在问：什么东西切一刀，会变成苹果片？答：苹果。

以前是：知道切之前的样子，问切完变成什么。

现在是：知道切完的样子，猜切之前是什么。

\textbf{反过来！}
\end{quote}

\section{找“反过来”}

\subsection{第一个：谁切一刀变成 $x$？}

我们来猜。

\textbf{猜测1：是不是 $x^2$？}

把 $x^2$ 切一刀，看看变成什么。

\[
(x^2)' = 2x
\]

变成了 $2x$，不是 $x$。

多了个 $2$。

\textbf{不对。$x^2$ 不是答案。}

\textbf{猜测2：那 $2$ 是多出来的，能不能想办法消掉？}

如果我切之前就把 $x^2$ \textbf{除以 2} 呢？

也就是 $\frac{x^2}{2}$。

把 $\frac{x^2}{2}$ 切一刀，看看变成什么。

$\frac{x^2}{2}$ 是什么意思？

就是 $x^2 \div 2$。

也可以写成 $\frac{1}{2} \times x^2$。

就是“半个 $x^2$”。

现在切一刀。

$x^2$ 切一刀变成 $2x$，对吧？

那“半个 $x^2$”切一刀呢？

就变成“半个 $2x$”。

半个 $2x$ 是多少？$\frac{1}{2} \times 2x = x$。

\textbf{变成了 $x$！}

\[
\left(\frac{x^2}{2}\right)' = x
\]

\textbf{找到了！}

$\frac{x^2}{2}$ 切一刀，就变成了 $x$。

所以，$x$ 的\textbf{“反过来”}是 $\frac{x^2}{2}$。

\begin{quote}
\textbf{📌 通俗理解：为什么要除以 2}

$x^2$ 切一刀变成 $2x$，多了个 2。

我们想要的是 $x$，不是 $2x$。

怎么消掉那个 2？

切之前先除以 2！

这样切完之后，多出来的 2 和除以的 2 就抵消了。

$\frac{x^2}{2}$ 切一刀 $\rightarrow$ $\frac{2x}{2}$ = $x$ ✓
\end{quote}

\subsection{第二个：谁切一刀变成 $x^2$？}

用同样的方法。

\textbf{猜测1：是不是 $x^3$？}

把 $x^3$ 切一刀。

\[
(x^3)' = 3x^2
\]

变成了 $3x^2$，不是 $x^2$。

多了个 $3$。

\textbf{不对。}

\textbf{猜测2：把 $x^3$ 除以 3 试试？}

也就是 $\frac{x^3}{3}$。

把 $\frac{x^3}{3}$ 切一刀。

$x^3$ 切一刀变成 $3x^2$。

“三分之一个 $x^3$”切一刀，就变成“三分之一个 $3x^2$”。

$\frac{1}{3} \times 3x^2 = x^2$。

\textbf{变成了 $x^2$！}

\[
\left(\frac{x^3}{3}\right)' = x^2
\]

\textbf{找到了！}

$\frac{x^3}{3}$ 切一刀，就变成了 $x^2$。

所以，$x^2$ 的\textbf{“反过来”}是 $\frac{x^3}{3}$。

\subsection{第三个：谁切一刀变成 $2x$？}

\textbf{猜测：是不是 $x^2$？}

把 $x^2$ 切一刀。

\[
(x^2)' = 2x
\]

变成了 $2x$！

\textbf{正好！不用除以什么。}

\textbf{找到了！}

$x^2$ 切一刀，就变成了 $2x$。

所以，$2x$ 的\textbf{“反过来”}是 $x^2$。

\subsection{第四个：谁切一刀变成常数 $3$？}

常数 $3$ 就是 $3$，不跟 $x$ 走，永远是 $3$。

谁切一刀会变成 $3$？

\textbf{猜测：是不是 $3x$？}

把 $3x$ 切一刀。

\[
(3x)' = 3
\]

\textbf{变成了 $3$！}

\textbf{找到了！}

$3x$ 切一刀，就变成了 $3$。

所以，$3$ 的\textbf{“反过来”}是 $3x$。

\begin{quote}
\textbf{📌 通俗理解：常数的反过来}

$3$ 的反过来是 $3x$。

为什么？

因为 $3x$ 切一刀就是 $3$。

$x$ 切一刀变成 $1$（常数）。

$3x$ 就是 $3$ 个 $x$，切一刀变成 $3$ 个 $1$ = $3$。
\end{quote}

\subsection{列个表}

\begin{center}
\begin{tabular}{c|c|c}
这个东西 & 切一刀变成 & 反过来说 \\
\hline
$\dfrac{x^2}{2}$ & $x$ & $x$ 的反过来是 $\dfrac{x^2}{2}$ \\
$\dfrac{x^3}{3}$ & $x^2$ & $x^2$ 的反过来是 $\dfrac{x^3}{3}$ \\
$x^2$ & $2x$ & $2x$ 的反过来是 $x^2$ \\
$3x$ & $3$ & $3$ 的反过来是 $3x$ \\
$5x$ & $5$ & $5$ 的反过来是 $5x$
\end{tabular}
\end{center}

\subsection{发现规律了吗？}

$x$ 的反过来是 $\frac{x^2}{2}$

$x^2$ 的反过来是 $\frac{x^3}{3}$

$x^3$ 的反过来是 $\frac{x^4}{4}$

\textbf{规律：指数加 1，然后除以新的指数。}

$x^1$ $\rightarrow$ 指数变成 $1+1=2$，除以 $2$ $\rightarrow \frac{x^2}{2}$

$x^2$ $\rightarrow$ 指数变成 $2+1=3$，除以 $3$ $\rightarrow \frac{x^3}{3}$

$x^3$ $\rightarrow$ 指数变成 $3+1=4$，除以 $4$ $\rightarrow \frac{x^4}{4}$

\begin{quote}
\textbf{📌 通俗理解：为什么是“指数加1，除以新指数”}

还记得求导的规律吗？

求导是：指数乘到前面，指数减 1。

$x^3$ 切一刀 $\rightarrow$ $3x^2$（指数 3 乘到前面，指数变成 2）

反过来呢？

反过来就把所有操作都反过来！

求导：指数减 1，乘到前面

反过来：指数加 1，除掉新指数

\textbf{正好相反！}
\end{quote}

\textbf{常数的反过来呢？}

$3$ 的反过来是 $3x$。

$5$ 的反过来是 $5x$。

$7$ 的反过来是 $7x$。

\textbf{常数 $c$ 的反过来是 $cx$。}

\subsection{这个“反过来”有个名字}

数学家给这个“反过来”起了个名字，叫\textbf{原函数}。

“原”就是“原来”的意思。

“原函数”就是“原来的函数”——切之前的样子。

但你不用记这个名字。

\textbf{记住“反过来”就行。}

“谁切一刀变成它”——那个“谁”，就是它的反过来。

\textbf{会找就好。}

\section{新方法来了}

现在，你已经会找“反过来”了。

接下来，告诉你新方法怎么用。

\subsection{新方法的步骤}

\textbf{计算 $\int_a^b f(x) \, dx$：}

\textbf{第一步：找 $f(x)$ 的“反过来”。}

问自己：谁，切一刀之后，会变成 $f(x)$？

把找到的那个东西叫做 $F(x)$。

\textbf{第二步：把 $\int$ 上面的数代入 $F(x)$。}

$\int$ 上面的数是 $b$。

把 $F(x)$ 里的 $x$ 换成 $b$，算出结果。

这个结果叫 $F(b)$。

\textbf{第三步：把 $\int$ 下面的数代入 $F(x)$。}

$\int$ 下面的数是 $a$。

把 $F(x)$ 里的 $x$ 换成 $a$，算出结果。

这个结果叫 $F(a)$。

\textbf{第四步：上面减下面。}

\[
\text{答案} = F(b) - F(a)
\]

就这四步。

\textbf{找反过来 $\rightarrow$ 代入上面的数 $\rightarrow$ 代入下面的数 $\rightarrow$ 相减。}

\begin{quote}
\textbf{📌 通俗理解：为什么是“代入上下，相减”？}

想象你开车。

车上有个里程表，显示你总共开了多少公里。

你从 A 地开到 B 地，想知道这段路有多长。

怎么算？

到 B 地时，看一眼里程表：320 公里。

想一想出发时（在 A 地），里程表是多少：280 公里。

这段路的长度 = 320 - 280 = 40 公里。

\textbf{上面减下面！}

积分也是一样。

$F(x)$ 就像里程表，记录着“从起点到现在的累积”。

$F(b)$ 是到 $b$ 时的累积。

$F(a)$ 是到 $a$ 时的累积。

从 $a$ 到 $b$ 这段的累积 = $F(b) - F(a)$。
\end{quote}

\section{例子一}

\[
\int_0^2 x \, dx
\]

\textbf{翻译}：从 $x=0$ 到 $x=2$，每根竖条的高度是 $x$，全部加起来。

\subsection{第一步：找 $x$ 的反过来}

问自己：谁切一刀变成 $x$？

答案：$\frac{x^2}{2}$。

（因为 $\left(\frac{x^2}{2}\right)' = x$，前面验证过了。）

所以 $F(x) = \frac{x^2}{2}$。

\subsection{第二步：代入上面的数 $2$}

\[
F(2) = \frac{2^2}{2}
\]

先算 $2^2$。

$2^2$ 是什么意思？就是 $2 \times 2 = 4$。

再算 $\frac{4}{2}$。

$\frac{4}{2}$ 是什么意思？就是 $4 \div 2 = 2$。

所以 $F(2) = 2$。

\subsection{第三步：代入下面的数 $0$}

\[
F(0) = \frac{0^2}{2}
\]

先算 $0^2$。

$0^2$ 是什么意思？就是 $0 \times 0 = 0$。

再算 $\frac{0}{2}$。

$\frac{0}{2}$ 是什么意思？就是 $0 \div 2 = 0$。

所以 $F(0) = 0$。

\subsection{第四步：上面减下面}

\[
\text{答案} = F(2) - F(0) = 2 - 0 = 2
\]

\textbf{答案：$\int_0^2 x \, dx = 2$}

\subsection{验证}

高度从 $0$ 升到 $2$，和地面围成一个三角形。

底 = $2$，高 = $2$。

三角形面积 = 底 $\times$ 高 $\div 2 = 2 \times 2 \div 2 = 2$ ✓

\section{例子二}

\[
\int_0^4 x \, dx
\]

\subsection{第一步：找 $x$ 的反过来}

和例子一一样：$F(x) = \frac{x^2}{2}$。

\subsection{第二步：代入上面的数 $4$}

\[
F(4) = \frac{4^2}{2}
\]

先算 $4^2$。$4 \times 4 = 16$。

再算 $\frac{16}{2}$。$16 \div 2 = 8$。

$F(4) = 8$。

\subsection{第三步：代入下面的数 $0$}

\[
F(0) = \frac{0^2}{2} = \frac{0}{2} = 0
\]

$F(0) = 0$。

\subsection{第四步：上面减下面}

\[
\text{答案} = F(4) - F(0) = 8 - 0 = 8
\]

\textbf{答案：$\int_0^4 x \, dx = 8$}

\subsection{验证}

三角形。底 = $4$，高 = $4$。

面积 = $4 \times 4 \div 2 = 8$ ✓

\section{例子三}

\[
\int_1^3 x \, dx
\]

注意！这次下面的数不是 $0$，是 $1$。

\subsection{第一步：找 $x$ 的反过来}

$F(x) = \frac{x^2}{2}$。

\subsection{第二步：代入上面的数 $3$}

\[
F(3) = \frac{3^2}{2}
\]

先算 $3^2$。$3 \times 3 = 9$。

再算 $\frac{9}{2}$。$9 \div 2 = 4.5$。

$F(3) = 4.5$。

\subsection{第三步：代入下面的数 $1$}

\[
F(1) = \frac{1^2}{2}
\]

先算 $1^2$。$1 \times 1 = 1$。

再算 $\frac{1}{2}$。$1 \div 2 = 0.5$。

$F(1) = 0.5$。

\subsection{第四步：上面减下面}

\[
\text{答案} = F(3) - F(1) = 4.5 - 0.5 = 4
\]

\textbf{答案：$\int_1^3 x \, dx = 4$}

\begin{quote}
\textbf{📌 注意：下面的数不是 0 的时候}

例子一和例子二，下面的数都是 $0$。

$F(0)$ 算出来也是 $0$，所以“上面减下面”就等于“上面减 0”= 上面。

但例子三不一样！下面的数是 $1$，$F(1) = 0.5$，不是 $0$。

所以一定要记得减！

\textbf{上面减下面，两个都要算，不能偷懒。}
\end{quote}

\section{例子四：$f$ 是 $x^2$}

\[
\int_0^2 x^2 \, dx
\]

这次 $f$ 不是 $x$ 了，是 $x^2$。

$f$ 走得更快了——不是一步一步匀速走，而是越走越快。

\subsection{第一步：找 $x^2$ 的反过来}

问自己：谁切一刀变成 $x^2$？

答案：$\frac{x^3}{3}$。

（因为 $\left(\frac{x^3}{3}\right)' = x^2$，前面验证过了。）

所以 $F(x) = \frac{x^3}{3}$。

\subsection{第二步：代入上面的数 $2$}

\[
F(2) = \frac{2^3}{3}
\]

先算 $2^3$。

$2^3$ 是什么意思？就是 $2 \times 2 \times 2$。

$2 \times 2 = 4$。

$4 \times 2 = 8$。

所以 $2^3 = 8$。

再算 $\frac{8}{3}$。

$8 \div 3$ 除不尽，就留着分数 $\frac{8}{3}$。

$F(2) = \frac{8}{3}$。

\subsection{第三步：代入下面的数 $0$}

\[
F(0) = \frac{0^3}{3}
\]

$0^3 = 0 \times 0 \times 0 = 0$。

$\frac{0}{3} = 0$。

$F(0) = 0$。

\subsection{第四步：上面减下面}

\[
\text{答案} = F(2) - F(0) = \frac{8}{3} - 0 = \frac{8}{3}
\]

\textbf{答案：$\int_0^2 x^2 \, dx = \frac{8}{3}$}

（$\frac{8}{3}$ 大约等于 $2.67$。）

\section{新方法和老方法是一家人}

让我证明给你看：新方法没有推翻老方法。

上一章里，$f = 3$（常数），我们用老方法算过：

\[
\int_0^4 3 \, dx = 3 \times (4 - 0) = 3 \times 4 = 12
\]

现在用新方法算同一道题：

\textbf{第一步：} $3$ 的反过来是什么？是 $3x$。

（因为 $(3x)' = 3$。）

所以 $F(x) = 3x$。

\textbf{第二步：} 代入上面的数 $4$：$F(4) = 3 \times 4 = 12$。

\textbf{第三步：} 代入下面的数 $0$：$F(0) = 3 \times 0 = 0$。

\textbf{第四步：} 上面减下面：$12 - 0 = 12$。

\textbf{一模一样！} ✓

老方法是新方法在“$f$ 不动”时的简化版。

新方法更强——$f$ 动不动都能算。

\begin{quote}
\textbf{📌 通俗理解：老方法和新方法的关系}

老方法：$f$ 不动（每个苹果都一个价），直接乘。

新方法：$f$ 动不动都能用。

老方法是新方法的\textbf{特殊情况}。

就像走路和跑步都是移动方式，走路是跑步的慢速版。

会新方法，老方法自然也会了。
\end{quote}

\section{\texorpdfstring{双重积分里，$f$ 也走路}{双重积分里，f 也走路}}


\[
\int_0^2 \int_0^3 x \, dy \, dx
\]

\textbf{翻译：}

\begin{itemize}
\item 外层 $\int_0^2 \cdots \, dx$：沿 $x$ 方向，从 $0$ 到 $2$
\item 内层 $\int_0^3 \cdots \, dy$：沿 $y$ 方向，从 $0$ 到 $3$
\item $f = x$：高度等于 $x$ 的位置
\end{itemize}

\subsection{第一步：先走里面的路（$y$）}

$x$ 不动。只看 $y$。

\[
\int_0^3 x \, dy
\]

\textbf{等一下。从 $y$ 的角度看，$x$ 是什么？}

$x$ 不跟 $y$ 走。$x$ 站在旁边看，不动。

所以从 $y$ 的角度看，$x$ 就是个\textbf{常数}。

常数可以用\textbf{老方法}！

提出来，乘以长度：

\[
\int_0^3 x \, dy = x \times (3 - 0) = x \times 3 = 3x
\]

$y$ 走完了，消失了。剩下 $3x$。

\begin{quote}
\textbf{📌 通俗理解：为什么 $x$ 变成了常数？}

现在走的是 $y$ 这条路。

走 $y$ 的时候，$x$ 站在旁边看，不动。

既然 $x$ 不动，它就像个普通的数字。

就像有个人站在路边，你从他身边走过。

对你来说，他的位置是固定的，不会变。
\end{quote}

\subsection{第二步：走外面的路（$x$）}

\[
\int_0^2 3x \, dx
\]

现在轮到 $x$ 走了。

$3x$ 是常数吗？

不是！$x$ 在变，$3x$ 也在变。

不能用老方法，要用\textbf{新方法}。

\textbf{找 $3x$ 的反过来。}

谁切一刀变成 $3x$？

我们知道 $x$ 的反过来是 $\frac{x^2}{2}$。

那 $3x$ 呢？$3x$ 就是 $3 \times x$。

\textbf{前面的常数可以直接乘进去。}

$3 \times x$ 的反过来 = $3 \times \frac{x^2}{2} = \frac{3x^2}{2}$。

\begin{quote}
\textbf{📌 通俗理解：为什么常数可以直接乘？}

$x$ 的反过来是 $\frac{x^2}{2}$。

$3x$ 就是 3 个 $x$。

3 个 $x$ 的反过来，就是 3 个“$x$ 的反过来”。

= $3 \times \frac{x^2}{2} = \frac{3x^2}{2}$。

就像：1 个苹果 5 块钱，3 个苹果就是 $3 \times 5 = 15$ 块钱。
\end{quote}

\textbf{验证一下：} $\left(\frac{3x^2}{2}\right)'$ 是多少？

$\frac{3x^2}{2}$ 就是 $\frac{3}{2} \times x^2$。

$x^2$ 切一刀是 $2x$。

$\frac{3}{2} \times x^2$ 切一刀，是 $\frac{3}{2} \times 2x = 3x$ ✓

确实变成 $3x$ 了。

所以 $F(x) = \frac{3x^2}{2}$。

\textbf{代入上面的数 $2$：}

\[
F(2) = \frac{3 \times 2^2}{2}
\]

先算 $2^2 = 4$。

再算 $3 \times 4 = 12$。

再算 $\frac{12}{2} = 6$。

$F(2) = 6$。

\textbf{代入下面的数 $0$：}

\[
F(0) = \frac{3 \times 0^2}{2} = \frac{0}{2} = 0
\]

$F(0) = 0$。

\textbf{上面减下面：}

\[
\text{答案} = 6 - 0 = 6
\]

\textbf{答案：$\int_0^2 \int_0^3 x \, dy \, dx = 6$}

\section{一句话总结}

以前：$f$ 站着不动 $\rightarrow$ \textbf{老方法}：提出来，乘以长度。

现在：$f$ 跟着走 $\rightarrow$ \textbf{新方法}：找反过来，代入上下，相减。

\begin{center}
\begin{tabular}{c|c|c}
情况 & $f$ 动不动 & 用什么方法 \\
\hline
$\int_0^4 3 \, dx$ & 不动 & 老方法 \\
$\int_0^4 x \, dx$ & 动了 & 新方法
\end{tabular}
\end{center}

\textbf{新方法四步走：}

\begin{enumerate}
\item \textbf{找反过来}：谁切一刀变成 $f(x)$？
\item \textbf{代入上面的数}
\item \textbf{代入下面的数}
\item \textbf{上面减下面} = 答案
\end{enumerate}

\textbf{找反过来的规律：}

$x^n$ 的反过来 = $\frac{x^{n+1}}{n+1}$

\textbf{指数加 1，除以新指数。}

常数 $c$ 的反过来 = $cx$

\section{下一章预告：$f$ 同时跟着两条路走}

这一章里，$f$ 学会走路了。

$f = x$，它跟着 $x$ 走。$x$ 往右，它也往右。$x$ 变大，它也变大。

但它只会跟着\textbf{一条路}走。

下一章，$f$ 要学会同时跟着\textbf{两条路}走了。

$f = x + y$

$x$ 往右走一步，$f$ 变大一点。

$y$ 往左走一步，$f$ 也变大一点。

两条路同时影响它。

听起来很复杂？

其实已经准备好了。

你会单重积分——一条路。

你会双重积分——两条路。

你会找“反过来”。

你会“代入上下，相减”。

\textbf{工具都在手里了，只是组合起来用而已。}

能学到这里的人，已经超过了绝大多数同龄人。

不，不只是同龄人。

很多大人学到这里都会晕。

但你没有晕。

下一章见！

\chapter{\texorpdfstring{当 $f$ 同时跟着两条路走：$f(x, y) = x + y$}{当 f 同时跟着两条路走：f(x, y) = x + y}}

上一章，$f$ 学会了跟着一条路走。

$f = x$，你往右走一步，$f$ 就变大一点。

这一章，$f$ 要同时跟着\textbf{两条路}走。

$f = x + y$

$x$ 变，$f$ 变。$y$ 变，$f$ 也变。两个方向同时影响它。

\section{先复习一下：求导是什么}

求导就是看\textbf{变化有多快}。

$x$ 变了一点点，$f$ 跟着变了多少？

$f$ 的变化量除以 $x$ 的变化量，就是导数。

\section{\texorpdfstring{$x + y$ 对 $x$ 求导}{x + y 对 x 求导}}

先看一个简单的情况：$y$ 不动。

$y$ 是个固定的数字，比如 $y = 5$。

那 $f(x) = x + 5$。

$x$ 从 10 变到 13，变了 3。

$f$ 呢？从 $10 + 5 = 15$ 变到 $13 + 5 = 18$，也变了 3。

$3 \div 3 = 1$。

$y$ 全程没动。所有变化都是 $x$ 造成的。

$x$ 变多少，$f$ 就变多少。变化速度是 1。

\[
\frac{d}{dx}(x + y) = 1
\]

不管 $y$ 是 7，是 100，还是 1000000，只要 $y$ 不动，结果都是 1。

因为 $y$ 站着不动，对变化毫无贡献。

\subsection{练一练}

\textbf{例子 1：} $y = 7$

\[
\frac{d}{dx}(x + 7) = 1
\]

验算：$x$ 从 0 变到 5，$f$ 从 7 变到 12，变了 5。$5 \div 5 = 1$。✓

\textbf{例子 2：} $y = 1000000$

\[
\frac{d}{dx}(x + 1000000) = 1
\]

不管 $y$ 多大，只要它不动，导数还是 1。

\section{\texorpdfstring{积分：把 $x + y$ 收集起来}{积分：把 x + y 收集起来}}

\subsection{不定积分}

\[
\int (x + y) \, dx
\]

$x + y$ 是两个东西加在一起。加法可以分开算——先收集 $x$，再收集 $y$。

\textbf{收集 $x$：}

\[
\int x \, dx = \frac{x^2}{2}
\]

（谁切一刀变成 $x$？答案是 $\frac{x^2}{2}$。上一章学过了。）

\textbf{收集 $y$：}

$y$ 不动，它就是个常数。

常数的反过来是常数乘以 $x$。

\[
\int y \, dx = yx
\]

\textbf{合在一起：}

\[
\int (x + y) \, dx = \frac{x^2}{2} + yx + C
\]

最后加一个 $C$。这个 $C$ 表示答案不唯一，可能还差一个常数。

（为什么？因为任何常数切一刀都变成 0，所以不定积分的答案里总会有一个“不确定的常数”。）

\subsection{验算}

把答案切一刀（求导），应该得到 $x + y$。

\[
\left(\frac{x^2}{2} + yx + C\right)' = x + y
\]

✓ 装回去的东西，拆开还是原样。

\subsection{练一练}

\textbf{例子 1：} $y = 3$

\[
\int (x + 3) \, dx = \frac{x^2}{2} + 3x + C
\]

\textbf{例子 2：} $y = 100$

\[
\int (x + 100) \, dx = \frac{x^2}{2} + 100x + C
\]

\subsection{定积分}

现在给一个范围：$x$ 从 $a$ 到 $b$。

\[
\int_a^b (x + y) \, dx = \left[\frac{x^2}{2} + yx\right]_a^b
\]

先算 $x = b$ 时的值，再算 $x = a$ 时的值，然后相减。

\textbf{例子 1：} $y = 3$，$x$ 从 0 到 2

\[
\int_0^2 (x + 3) \, dx
\]

第一步：反过来是 $\frac{x^2}{2} + 3x$

第二步：$x = 2$ 代进去

\[
\frac{2^2}{2} + 3 \times 2 = \frac{4}{2} + 6 = 2 + 6 = 8
\]

第三步：$x = 0$ 代进去

\[
\frac{0^2}{2} + 3 \times 0 = 0 + 0 = 0
\]

第四步：上面减下面

\[
8 - 0 = 8
\]

\textbf{答案：$\int_0^2 (x + 3) \, dx = 8$}

\textbf{例子 2：} $y = 5$，$x$ 从 1 到 4

\[
\int_1^4 (x + 5) \, dx
\]

第一步：反过来是 $\frac{x^2}{2} + 5x$

第二步：$x = 4$ 代进去

\[
\frac{16}{2} + 20 = 8 + 20 = 28
\]

第三步：$x = 1$ 代进去

\[
\frac{1}{2} + 5 = 5.5
\]

第四步：上面减下面

\[
28 - 5.5 = 22.5
\]

\textbf{答案：$\int_1^4 (x + 5) \, dx = 22.5$}

\section{\texorpdfstring{二元微积分：$x$ 和 $y$ 都能动了}{二元微积分：x 和 y 都能动了}}

前面的例子里，$y$ 一直站着不动，当一个固定的数字。

但现在，$x$ 和 $y$ 都可以变。

$f(x, y) = x + y$

$x$ 变，$f$ 变。$y$ 变，$f$ 也变。两个方向同时影响结果。

\subsection{偏导数：一次只看一个}

两个变量都在动，怎么办？

\textbf{一次只看一个，假装另一个不动。}

这就是\textbf{偏导数}。

\subsection{\texorpdfstring{对 $x$ 求偏导}{对 x 求偏导}}

假装 $y$ 不动（比如 $y = 5$），只看 $x$ 的变化。

$f = x + 5$

$x$ 从 10 变到 12，变了 2。

$f$ 从 15 变到 17，也变了 2。

$2 \div 2 = 1$。

\[
\frac{\partial}{\partial x}(x + y) = 1
\]

那个弯弯的 $\partial$ 符号就是偏导的标志，表示“只看这一个变量”。

\subsection{\texorpdfstring{对 $y$ 求偏导}{对 y 求偏导}}

现在反过来。假装 $x$ 不动（比如 $x = 7$），只看 $y$ 的变化。

$f = 7 + y$

$y$ 从 3 变到 8，变了 5。

$f$ 从 10 变到 15，也变了 5。

$5 \div 5 = 1$。

\[
\frac{\partial}{\partial y}(x + y) = 1
\]

\subsection{这说明什么？}

$x + y$ 对 $x$ 求偏导是 1，对 $y$ 求偏导也是 1。

$x$ 和 $y$ 在这个函数里地位平等。$x$ 变 1，答案变 1。$y$ 变 1，答案也变 1。

谁也不比谁重要。

\section{二重积分：收集两个方向}

以前积分只沿着一条路收集。

现在有两条路，就要收集两次。先沿一个方向扫一遍，再沿另一个方向扫一遍。

这就是\textbf{二重积分}。

\subsection{例子 1：正方形区域}

$x$ 从 0 到 1，$y$ 从 0 到 1。

\[
\int_0^1 \int_0^1 (x + y) \, dx \, dy
\]

\textbf{第一步：先对 $x$ 积分（$y$ 不动）}

\[
\int_0^1 (x + y) \, dx
\]

反过来是 $\frac{x^2}{2} + yx$

$x = 1$ 代进去：$\frac{1}{2} + y$

$x = 0$ 代进去：$0$

相减：$\frac{1}{2} + y$

第一遍扫完，剩下 $\frac{1}{2} + y$。

\textbf{第二步：再对 $y$ 积分}

\[
\int_0^1 \left(\frac{1}{2} + y\right) \, dy
\]

反过来是 $\frac{y}{2} + \frac{y^2}{2}$

$y = 1$ 代进去：$\frac{1}{2} + \frac{1}{2} = 1$

$y = 0$ 代进去：$0$

相减：$1 - 0 = 1$

\[
\int_0^1 \int_0^1 (x + y) \, dx \, dy = 1
\]

\subsection{例子 2：更大的矩形}

$x$ 从 0 到 2，$y$ 从 0 到 3。

\[
\int_0^3 \int_0^2 (x + y) \, dx \, dy
\]

\textbf{第一步：先对 $x$ 积分}

\[
\int_0^2 (x + y) \, dx
\]

反过来是 $\frac{x^2}{2} + yx$

$x = 2$ 代进去：$\frac{4}{2} + 2y = 2 + 2y$

$x = 0$ 代进去：$0$

相减：$2 + 2y$

\textbf{第二步：再对 $y$ 积分}

\[
\int_0^3 (2 + 2y) \, dy
\]

反过来是 $2y + y^2$

$y = 3$ 代进去：$6 + 9 = 15$

$y = 0$ 代进去：$0$

相减：$15$

\[
\int_0^3 \int_0^2 (x + y) \, dx \, dy = 15
\]

\subsection{例子 3：起点不是 0}

$x$ 从 1 到 3，$y$ 从 2 到 5。

\[
\int_2^5 \int_1^3 (x + y) \, dx \, dy
\]

\textbf{第一步：先对 $x$ 积分}

\[
\int_1^3 (x + y) \, dx
\]

反过来是 $\frac{x^2}{2} + yx$

$x = 3$ 代进去：$\frac{9}{2} + 3y$

$x = 1$ 代进去：$\frac{1}{2} + y$

相减：$\frac{9}{2} - \frac{1}{2} + 3y - y = 4 + 2y$

\textbf{第二步：再对 $y$ 积分}

\[
\int_2^5 (4 + 2y) \, dy
\]

反过来是 $4y + y^2$

$y = 5$ 代进去：$20 + 25 = 45$

$y = 2$ 代进去：$8 + 4 = 12$

相减：$45 - 12 = 33$

\[
\int_2^5 \int_1^3 (x + y) \, dx \, dy = 33
\]

\subsection{例子 4：三角形区域}

三角形的三个角是 $(0, 0)$、$(1, 0)$、$(0, 1)$。

这个三角形的斜边是 $x + y = 1$，也就是 $y = 1 - x$。

三角形不是规整的矩形，$y$ 的范围会随着 $x$ 变化。

对于每一个 $x$（从 0 到 1），$y$ 的范围是从 0 到 $1 - x$。

比如：

当 $x = 0$ 时，$y$ 可以从 0 到 1（完整的一段）。

当 $x = 0.5$ 时，$y$ 可以从 0 到 0.5（短了一半）。

当 $x = 1$ 时，$y$ 只能是 0（只剩一个点了）。

\[
\int_0^1 \int_0^{1-x} (x + y) \, dy \, dx
\]

注意：这次先对 $y$ 积分，因为 $y$ 的范围依赖于 $x$。

\textbf{第一步：先对 $y$ 积分}

\[
\int_0^{1-x} (x + y) \, dy
\]

把 $x$ 当成常数。

反过来是 $xy + \frac{y^2}{2}$

$y = 1 - x$ 代进去：

先算 $xy$：$x(1-x) = x - x^2$

再算 $\frac{y^2}{2}$：

$(1-x)^2 = 1 - 2x + x^2$

$\frac{(1-x)^2}{2} = \frac{1}{2} - x + \frac{x^2}{2}$

加起来：

$x - x^2 + \frac{1}{2} - x + \frac{x^2}{2}$

整理一下：

$x$ 项：$x - x = 0$

$x^2$ 项：$-x^2 + \frac{x^2}{2} = -\frac{x^2}{2}$

常数项：$\frac{1}{2}$

结果是：$\frac{1}{2} - \frac{x^2}{2}$

$y = 0$ 代进去：$0$

相减：$\frac{1}{2} - \frac{x^2}{2}$

\textbf{第二步：再对 $x$ 积分}

\[
\int_0^1 \left(\frac{1}{2} - \frac{x^2}{2}\right) \, dx
\]

反过来是 $\frac{x}{2} - \frac{x^3}{6}$

$x = 1$ 代进去：$\frac{1}{2} - \frac{1}{6}$

把 $\frac{1}{2}$ 变成 $\frac{3}{6}$：$\frac{3}{6} - \frac{1}{6} = \frac{2}{6} = \frac{1}{3}$

$x = 0$ 代进去：$0$

相减：$\frac{1}{3}$

\[
\int_0^1 \int_0^{1-x} (x + y) \, dy \, dx = \frac{1}{3}
\]

\subsection{交换积分顺序，验证一下}

先扫 $x$ 再扫 $y$，和先扫 $y$ 再扫 $x$，结果应该一样。

同样的三角形，换一种写法：

对于每一个 $y$（从 0 到 1），$x$ 的范围是从 0 到 $1 - y$。

\[
\int_0^1 \int_0^{1-y} (x + y) \, dx \, dy
\]

\textbf{第一步：先对 $x$ 积分}

\[
\int_0^{1-y} (x + y) \, dx
\]

反过来是 $\frac{x^2}{2} + yx$

$x = 1 - y$ 代进去：

\[
\frac{(1-y)^2}{2} + y(1-y)
\]

\[
= \frac{1 - 2y + y^2}{2} + y - y^2
\]

\[
= \frac{1}{2} - y + \frac{y^2}{2} + y - y^2
\]

\[
= \frac{1}{2} - \frac{y^2}{2}
\]

$x = 0$ 代进去：$0$

相减：$\frac{1}{2} - \frac{y^2}{2}$

\textbf{第二步：再对 $y$ 积分}

\[
\int_0^1 \left(\frac{1}{2} - \frac{y^2}{2}\right) \, dy = \frac{1}{3}
\]

和之前一模一样。

\textbf{不管从哪个方向开始扫，结果都是 $\frac{1}{3}$。} ✓

\section{本章小结}

\subsection{一元微积分（$y$ 是常数，不动）}

\begin{center}
\begin{tabular}{c|c|c}
操作 & 公式 & 答案 \\
\hline
求导 & $\dfrac{d}{dx}(x + y)$ & $1$ \\
不定积分 & $\displaystyle\int (x + y) \, dx$ & $\dfrac{x^2}{2} + yx + C$ \\
定积分 & $\displaystyle\int_a^b (x + y) \, dx$ & $\left(\dfrac{b^2}{2} + yb\right) - \left(\dfrac{a^2}{2} + ya\right)$
\end{tabular}
\end{center}

\subsection{二元微积分（$x$ 和 $y$ 都能动）}

\begin{center}
\begin{tabular}{c|c|c}
操作 & 公式 & 答案 \\
\hline
对 $x$ 偏导 & $\dfrac{\partial}{\partial x}(x + y)$ & $1$ \\
对 $y$ 偏导 & $\dfrac{\partial}{\partial y}(x + y)$ & $1$ \\
二重积分（正方形 $[0,1]\times[0,1]$） & $\displaystyle\int_0^1 \int_0^1 (x + y) \, dx \, dy$ & $1$ \\
二重积分（矩形 $[0,2]\times[0,3]$） & $\displaystyle\int_0^3 \int_0^2 (x + y) \, dx \, dy$ & $15$ \\
二重积分（三角形） & $\displaystyle\int_0^1 \int_0^{1-x} (x + y) \, dy \, dx$ & $\dfrac{1}{3}$
\end{tabular}
\end{center}

\subsection{常用积分规则速查表}

\[
\int c \, dx = cx + C \quad \text{（常数的积分）}
\]

\[
\int x \, dx = \frac{x^2}{2} + C
\]

\[
\int x^2 \, dx = \frac{x^3}{3} + C
\]

\[
\int x^3 \, dx = \frac{x^4}{4} + C
\]

\[
\int x^n \, dx = \frac{x^{n+1}}{n+1} + C \quad \text{（一般规律：指数加 1，除以新指数）}
\]

\[
\int (a + b) \, dx = \int a \, dx + \int b \, dx \quad \text{（加法可以分开算）}
\]

\subsection{为什么总用 $x + y$ 练手？}

因为 $x + y$ 是最简单的组合。

它不会给你添乱——求导是 1，积分也很规整。

当你在学新方法的时候，不希望函数本身太复杂让你分心。

所以先用 $x + y$ 把方法练熟，等方法会了，再去挑战更复杂的函数。

\chapter{\texorpdfstring{$x + y + z$：三元微积分篇}{x + y + z：三元微积分篇}}

\section{\texorpdfstring{开场：第三个变量 $z$}{开场：第三个变量 z}}

在上一篇里，我们已经和 $x$、$y$ 这两个变量打过交道了。

现在，第三个变量来了：\textbf{$z$}。

$z$ 也是一个会变的数。今天可能是 1，明天可能是 99，后天可能是 $-7$。它和 $x$、$y$ 一样善变。

那 \textbf{$x + y + z$} 是什么？就是三个变量凑在一起的总和。

比如说：$x = 2$，$y = 3$，$z = 5$。

那 $x + y + z = 2 + 3 + 5 = 10$。

三个善变的数加在一起，得到一个更更善变的结果。

\section{\texorpdfstring{一、一元微积分里的 $x + y + z$}{一、一元微积分里的 x + y + z}}

在一元微积分里，只有 $x$ 能动。$y$ 和 $z$ 都必须固定不变，当常数。

（没办法，一元微积分的规矩就是这样，只能有一个变量。）

假设 $y$ 今天是 3，$z$ 今天是 7。

那 $f(x) = x + y + z$ 就变成了 $f(x) = x + 3 + 7 = x + 10$。

\subsection{\texorpdfstring{对 $x + y + z$ 求导}{对 x + y + z 求导}}

我们来看看，$x$ 变化的时候，答案变多快。

\subsubsection{原式}

\[
f(x) = x + y + z
\]

\subsubsection{公式}

\[
\frac{d}{dx}(x + y + z) = \frac{d}{dx}(x) + \frac{d}{dx}(y) + \frac{d}{dx}(z)
\]

\subsubsection{计算过程}

$x$ 对 $x$ 求导等于 1（老朋友了）：

\[
\frac{d}{dx}(x) = 1
\]

$y$ 是常数，固定不动，导数等于 0：

\[
\frac{d}{dx}(y) = 0
\]

$z$ 也是常数，也固定不动，导数也等于 0：

\[
\frac{d}{dx}(z) = 0
\]

三个加起来：

\[
\frac{d}{dx}(x + y + z) = 1 + 0 + 0 = 1
\]

$y$ 和 $z$ 全程没动，对变化毫无贡献。所有变化都是 $x$ 带来的。变化速度还是 1。

\subsubsection{最终答案}

\[
\frac{d}{dx}(x + y + z) = 1
\]

\subsubsection{例子}

\textbf{例子 1：} $y = 3$，$z = 7$

\[
f(x) = x + 3 + 7 = x + 10
\]

\[
\frac{d}{dx}(x + 10) = 1
\]

验算：$x$ 从 5 变到 8，变了 3。$f(x)$ 从 15 变到 18，也变了 3。$3 \div 3 = 1$。完美！

\textbf{例子 2：} $y = 100$，$z = 200$

\[
f(x) = x + 100 + 200 = x + 300
\]

\[
\frac{d}{dx}(x + 300) = 1
\]

不管 $y$ 和 $z$ 加起来有多大，它们固定不动就是不动，导数还是 1。

\subsection{\texorpdfstring{对 $x + y + z$ 求不定积分}{对 x + y + z 求不定积分}}

\subsubsection{原式}

\[
\int (x + y + z) \, dx
\]

\subsubsection{公式}

\[
\int (x + y + z) \, dx = \int x \, dx + \int y \, dx + \int z \, dx
\]

\subsubsection{计算过程}

$x$ 的积分是 $\frac{x^2}{2}$（老规矩）：

\[
\int x \, dx = \frac{x^2}{2}
\]

$y$ 是常数，常数的积分是常数乘以 $x$：

\[
\int y \, dx = yx
\]

$z$ 也是常数，同样道理：

\[
\int z \, dx = zx
\]

三个加起来，别忘了 $C$：

\[
\int (x + y + z) \, dx = \frac{x^2}{2} + yx + zx + C
\]

可以写得更整齐：

\[
\int (x + y + z) \, dx = \frac{x^2}{2} + (y + z)x + C
\]

\subsubsection{最终答案}

\[
\int (x + y + z) \, dx = \frac{x^2}{2} + yx + zx + C
\]

\subsubsection{例子}

\textbf{例子 1：} $y = 3$，$z = 7$

\[
\int (x + 3 + 7) \, dx = \int (x + 10) \, dx = \frac{x^2}{2} + 10x + C
\]

验算：对答案求导

\[
\frac{d}{dx}\left(\frac{x^2}{2} + 10x + C\right) = x + 10
\]

对了！

\textbf{例子 2：} $y = 2$，$z = 5$

\[
\int (x + 2 + 5) \, dx = \int (x + 7) \, dx = \frac{x^2}{2} + 7x + C
\]

\subsection{\texorpdfstring{对 $x + y + z$ 求定积分}{对 x + y + z 求定积分}}

现在给一个范围，来算具体的总量。

\subsubsection{原式}

\[
\int_a^b (x + y + z) \, dx
\]

\subsubsection{例子 1}

$y = 3$，$z = 7$，$x$ 从 0 到 2。

\[
\int_0^2 (x + 10) \, dx
\]

第一步：不定积分是 $\frac{x^2}{2} + 10x$

第二步：$x = 2$ 代进去

\[
\frac{4}{2} + 20 = 2 + 20 = 22
\]

第三步：$x = 0$ 代进去

\[
0 + 0 = 0
\]

第四步：上限减下限

\[
22 - 0 = 22
\]

\subsubsection{完整计算写出来}

\[
\int_0^2 (x + 3 + 7) \, dx = 22
\]

\subsubsection{例子 2}

$y = 1$，$z = 4$，$x$ 从 1 到 3。

\[
\int_1^3 (x + 1 + 4) \, dx = \int_1^3 (x + 5) \, dx
\]

第一步：不定积分是 $\frac{x^2}{2} + 5x$

第二步：$x = 3$ 代进去

\[
\frac{9}{2} + 15 = 4.5 + 15 = 19.5
\]

第三步：$x = 1$ 代进去

\[
\frac{1}{2} + 5 = 0.5 + 5 = 5.5
\]

第四步：上限减下限

\[
19.5 - 5.5 = 14
\]

\subsubsection{完整计算写出来}

\[
\int_1^3 (x + 1 + 4) \, dx = 14
\]

\section{\texorpdfstring{二、二元微积分里的 $x + y + z$}{二、二元微积分里的 x + y + z}}

现在 $x$ 和 $y$ 都可以动了！但 $z$ 还是要固定不变，当常数。

假设 $z$ 今天是 5。

那 $f(x, y) = x + y + z$ 就变成了 $f(x, y) = x + y + 5$。

\subsection{\texorpdfstring{对 $x$ 求偏导}{对 x 求偏导}}

让 $y$ 和 $z$ 都固定不动，单独观察 $x$ 的变化。

\subsubsection{原式}

\[
f(x, y) = x + y + z
\]

\subsubsection{计算}

$x$ 对 $x$ 求导等于 1：

\[
\frac{\partial}{\partial x}(x) = 1
\]

$y$ 被当成常数，导数等于 0：

\[
\frac{\partial}{\partial x}(y) = 0
\]

$z$ 本来就是常数，导数等于 0：

\[
\frac{\partial}{\partial x}(z) = 0
\]

加起来：

\[
\frac{\partial}{\partial x}(x + y + z) = 1 + 0 + 0 = 1
\]

\subsubsection{最终答案}

\[
\frac{\partial f}{\partial x} = 1
\]

\subsection{\texorpdfstring{对 $y$ 求偏导}{对 y 求偏导}}

让 $x$ 和 $z$ 都固定不动，单独观察 $y$ 的变化。

\subsubsection{计算}

$x$ 被当成常数，导数等于 0：

\[
\frac{\partial}{\partial y}(x) = 0
\]

$y$ 对 $y$ 求导等于 1：

\[
\frac{\partial}{\partial y}(y) = 1
\]

$z$ 本来就是常数，导数等于 0：

\[
\frac{\partial}{\partial y}(z) = 0
\]

加起来：

\[
\frac{\partial}{\partial y}(x + y + z) = 0 + 1 + 0 = 1
\]

\subsubsection{最终答案}

\[
\frac{\partial f}{\partial y} = 1
\]

\subsection{\texorpdfstring{二重积分（$z$ 是常数）}{二重积分（z 是常数）}}

在一块平面区域上，把 $x + y + z$ 收集起来。$z$ 是固定的。

\subsubsection{例子：正方形区域}

$z = 5$，区域是 $x$ 从 0 到 1，$y$ 从 0 到 1。

\[
\int_0^1 \int_0^1 (x + y + 5) \, dx \, dy
\]

\textbf{第一步：先对 $x$ 积分}

\[
\int_0^1 (x + y + 5) \, dx
\]

不定积分是 $\frac{x^2}{2} + yx + 5x$

$x = 1$ 代进去：$\frac{1}{2} + y + 5$

$x = 0$ 代进去：$0$

相减：$\frac{1}{2} + y + 5 = y + 5.5$

\textbf{第二步：再对 $y$ 积分}

\[
\int_0^1 (y + 5.5) \, dy
\]

不定积分是 $\frac{y^2}{2} + 5.5y$

$y = 1$ 代进去：$\frac{1}{2} + 5.5 = 6$

$y = 0$ 代进去：$0$

相减：$6$

\subsubsection{完整计算写出来}

\[
\int_0^1 \int_0^1 (x + y + 5) \, dx \, dy = 6
\]

\subsubsection{和 $x + y$ 的对比}

还记得吗？$x + y$ 在同样的正方形上的二重积分是 1。

现在 $x + y + 5$ 的结果是 6。

为什么多了 5？

因为 $z = 5$ 是一个常数，它在整个 $1 \times 1$ 的正方形上都贡献了 5。

$1 \times 1$ 的面积是 1，5 乘以面积 1 就是 5。

原来 $x + y$ 的积分是 1，加上 5，就是 6。

这里有个小技巧：\textbf{常数的二重积分等于常数乘以面积！}

\section{\texorpdfstring{三、三元微积分（$x$、$y$、$z$ 都是变量）}{三、三元微积分（x、y、z 都是变量）}}

从现在开始，$x$、$y$、$z$ 三个都可以动了！

现在 $f(x, y, z) = x + y + z$ 真正地依赖于三个变量了。$x$ 变，答案变。$y$ 变，答案变。$z$ 变，答案也变。三个变量都在影响最终结果。

\subsection{三个偏导数}

三个变量都在动，怎么办？还是老办法：一次只盯一个，其他两个当常数。只不过现在要写三份，一个变量一份。

\subsubsection{\texorpdfstring{对 $x$ 求偏导}{对 x 求偏导}}

让 $y$ 和 $z$ 都固定不动，单独观察 $x$。

\[
\frac{\partial}{\partial x}(x + y + z) = 1 + 0 + 0 = 1
\]

\subsubsection{\texorpdfstring{对 $y$ 求偏导}{对 y 求偏导}}

让 $x$ 和 $z$ 都固定不动，单独观察 $y$。

\[
\frac{\partial}{\partial y}(x + y + z) = 0 + 1 + 0 = 1
\]

\subsubsection{\texorpdfstring{对 $z$ 求偏导}{对 z 求偏导}}

让 $x$ 和 $y$ 都固定不动，单独观察 $z$。

\[
\frac{\partial}{\partial z}(x + y + z) = 0 + 0 + 1 = 1
\]

\subsubsection{总结}

\[
\frac{\partial f}{\partial x} = 1
\]

\[
\frac{\partial f}{\partial y} = 1
\]

\[
\frac{\partial f}{\partial z} = 1
\]

$x + y + z$ 这个函数对谁求偏导都是 1。

这说明 $x$、$y$、$z$ 在这个函数里地位完全平等。谁变 1，答案就变 1。三个公平的合伙人，谁也不比谁重要。

\subsubsection{例子}

设 $f(x, y, z) = x + y + z$，求 $f$ 在点 $(2, 3, 5)$ 处的三个偏导数。

对 $x$ 的偏导数：$f_x(2, 3, 5) = 1$

对 $y$ 的偏导数：$f_y(2, 3, 5) = 1$

对 $z$ 的偏导数：$f_z(2, 3, 5) = 1$

不管在哪个点，三个偏导数都是 1。这个函数真的很老实。

\subsection{梯度向量（三个偏导数打包在一起）}

每次写三个偏导数太麻烦了，可以把它们打包成一个向量。

这个向量叫做\textbf{梯度}，符号是 $\nabla f$（读作“nabla f”或者“del f”）。

\[
\nabla f = \left(\frac{\partial f}{\partial x}, \frac{\partial f}{\partial y}, \frac{\partial f}{\partial z}\right)
\]

对于 $f(x, y, z) = x + y + z$：

\[
\nabla f = (1, 1, 1)
\]

这个向量 $(1, 1, 1)$ 告诉你：不管往 $x$、$y$、$z$ 哪个方向走，上升的速度都是一样的。

这个函数就像一个均匀倾斜的斜坡，三个方向都同样陡。

\section{四、三重积分}

以前沿着一条线收集（一重积分），后来在一块平面上收集（二重积分）。

现在，要在一个\textbf{立体空间}里收集！

所以要积分三次。先沿着一个方向扫，再沿着第二个方向扫，最后沿着第三个方向扫。

这就是\textbf{三重积分}。

\subsection{三重积分的符号}

\[
\iiint_V (x + y + z) \, dV
\]

或者写成三次积分的形式：

\[
\int \int \int (x + y + z) \, dx \, dy \, dz
\]

$V$ 代表一个三维的立体区域。

\subsection{例子 1：单位立方体}

$V$ 是一个立方体：$x$ 从 0 到 1，$y$ 从 0 到 1，$z$ 从 0 到 1。

这是一个边长为 1 的正方体。

\[
\int_0^1 \int_0^1 \int_0^1 (x + y + z) \, dx \, dy \, dz
\]

要扫三遍。

\textbf{第一步：先对 $x$ 积分（$y$ 和 $z$ 都固定不动）}

\[
\int_0^1 (x + y + z) \, dx
\]

不定积分是 $\frac{x^2}{2} + yx + zx$

$x = 1$ 代进去：$\frac{1}{2} + y + z$

$x = 0$ 代进去：$0$

相减：$\frac{1}{2} + y + z$

第一遍扫完，收集到了 $\frac{1}{2} + y + z$。

\textbf{第二步：再对 $y$ 积分（$z$ 还是固定不动）}

\[
\int_0^1 \left(\frac{1}{2} + y + z\right) \, dy
\]

不定积分是 $\frac{y}{2} + \frac{y^2}{2} + zy$

$y = 1$ 代进去：$\frac{1}{2} + \frac{1}{2} + z = 1 + z$

$y = 0$ 代进去：$0$

相减：$1 + z$

第二遍扫完，收集到了 $1 + z$。

\textbf{第三步：最后对 $z$ 积分（$z$ 终于可以动了）}

\[
\int_0^1 (1 + z) \, dz
\]

不定积分是 $z + \frac{z^2}{2}$

$z = 1$ 代进去：$1 + \frac{1}{2} = \frac{3}{2}$

$z = 0$ 代进去：$0$

相减：$\frac{3}{2}$

这个单位立方体上的三重积分结果是 \textbf{$\frac{3}{2}$}！

\subsubsection{完整计算写出来}

\[
\int_0^1 \int_0^1 \int_0^1 (x + y + z) \, dx \, dy \, dz = \frac{3}{2}
\]

\subsubsection{另一种写法}

$\frac{3}{2}$ 就是 1.5。

\subsection{例子 2：更大的立方体}

$V$ 是一个立方体：$x$ 从 0 到 2，$y$ 从 0 到 2，$z$ 从 0 到 2。

这是一个边长为 2 的正方体。

\[
\int_0^2 \int_0^2 \int_0^2 (x + y + z) \, dx \, dy \, dz
\]

\textbf{第一步：先对 $x$ 积分}

\[
\int_0^2 (x + y + z) \, dx
\]

不定积分是 $\frac{x^2}{2} + yx + zx$

$x = 2$ 代进去：$\frac{4}{2} + 2y + 2z = 2 + 2y + 2z$

$x = 0$ 代进去：$0$

相减：$2 + 2y + 2z$

\textbf{第二步：再对 $y$ 积分}

\[
\int_0^2 (2 + 2y + 2z) \, dy
\]

不定积分是 $2y + y^2 + 2zy$

$y = 2$ 代进去：$4 + 4 + 4z = 8 + 4z$

$y = 0$ 代进去：$0$

相减：$8 + 4z$

\textbf{第三步：最后对 $z$ 积分}

\[
\int_0^2 (8 + 4z) \, dz
\]

不定积分是 $8z + 2z^2$

$z = 2$ 代进去：$16 + 8 = 24$

$z = 0$ 代进去：$0$

相减：$24$

这个边长为 2 的立方体上的三重积分结果是 \textbf{24}！

\subsubsection{完整计算写出来}

\[
\int_0^2 \int_0^2 \int_0^2 (x + y + z) \, dx \, dy \, dz = 24
\]

\subsection{例子 3：长方体}

$V$ 是一个长方体：$x$ 从 0 到 1，$y$ 从 0 到 2，$z$ 从 0 到 3。

\[
\int_0^3 \int_0^2 \int_0^1 (x + y + z) \, dx \, dy \, dz
\]

\textbf{第一步：先对 $x$ 积分}

\[
\int_0^1 (x + y + z) \, dx
\]

不定积分是 $\frac{x^2}{2} + yx + zx$

$x = 1$ 代进去：$\frac{1}{2} + y + z$

$x = 0$ 代进去：$0$

相减：$\frac{1}{2} + y + z$

\textbf{第二步：再对 $y$ 积分}

\[
\int_0^2 \left(\frac{1}{2} + y + z\right) \, dy
\]

不定积分是 $\frac{y}{2} + \frac{y^2}{2} + zy$

$y = 2$ 代进去：$1 + 2 + 2z = 3 + 2z$

$y = 0$ 代进去：$0$

相减：$3 + 2z$

\textbf{第三步：最后对 $z$ 积分}

\[
\int_0^3 (3 + 2z) \, dz
\]

不定积分是 $3z + z^2$

$z = 3$ 代进去：$9 + 9 = 18$

$z = 0$ 代进去：$0$

相减：$18$

这个长方体上的三重积分结果是 \textbf{18}！

\subsubsection{完整计算写出来}

\[
\int_0^3 \int_0^2 \int_0^1 (x + y + z) \, dx \, dy \, dz = 18
\]

\subsection{例子 4：四面体（三维版的三角形）}

这是一个更有挑战的任务：在一个四面体上对 $x + y + z$ 积分。

先别急着算。我们得先搞清楚四面体是什么东西。

三角形是平面上最简单的图形，三条边围起来的。四面体呢，就是立体空间里最简单的形状，四个面围起来的。

所以四面体就是三角形的立体版。三角形有 3 个顶点，四面体有 4 个顶点。

光听描述有点抽象。在算积分之前，我们先一起动手画一个四面体出来！

拿好纸和笔，跟着一步一步画。

\subsubsection{认识三个方向}

画四面体之前，得先学会在纸上画立体的图。

纸是平的，怎么画立体的东西？我们先从认识三个方向开始。

想象一个房间的墙角。你就站在那个墙角里。那个角落就是原点，我们叫它 $(0, 0, 0)$。

从那个角落出发，有三个方向可以走。

\textbf{$x$ 方向就是往右走}，沿着地面，往右边的墙走。

\textbf{$y$ 方向就是往前走}，沿着地面，往前面的墙走。

\textbf{$z$ 方向就是往上走}，离开地面，往天花板走。

可以伸出右手来感受一下：

大拇指指向右边——这是 $x$ 方向。

食指指向前面——这是 $y$ 方向。

中指指向上面——这是 $z$ 方向。

三个手指互相垂直，就像房间的三条边一样。记住这个手势，后面画图的时候会用到。

\subsubsection{坐标是“走路说明书”}

现在来学最重要的一件事：坐标 $(x, y, z)$ 到底是什么意思。

坐标就是三个数，但更具体地说，它是一份\textbf{走路说明书}。告诉你从原点出发，怎么走到那个点。

\begin{quote}
\textbf{坐标 $(x, y, z)$ 的意思是：}

先往右走 $x$ 步，再往前走 $y$ 步，最后往上走 $z$ 步。

走完这三步，你就到了。
\end{quote}

来试试看：

坐标 \textbf{(1, 0, 0)} 怎么走？往右走 1 步，往前走 0 步——也就是不走，往上走 0 步——也不走。所以只需要往右走 1 步就到了。这个点在地面上，在右边。

那 \textbf{(0, 1, 0)} 呢？往右走 0 步——不走，往前走 1 步，往上走 0 步——不走。所以只需要往前走 1 步就到了。这个点在地面上，在前面。

\textbf{(0, 0, 1)} 呢？往右走 0 步——不走，往前走 0 步——不走，往上走 1 步。所以得飘到角落正上方 1 步的位置。这个点在空中。

那 \textbf{(2, 3, 1)} 呢？往右走 2 步，再往前走 3 步，最后往上走 1 步。就到了一个在空中的点。

坐标就是走路说明书。三个数告诉你三个方向各走几步。

\subsubsection{在纸上画坐标轴}

好了，现在拿出纸和笔，跟着画。

\textbf{第一步：画 $x$ 轴。}

在纸的中间偏左下的位置，画一条往右的箭头。在箭头旁边写上 $x$。这条线代表往右走。

\textbf{第二步：画 $y$ 轴。}

从同一个起点，画一条往左上方斜着的箭头。在箭头旁边写上 $y$。这条线代表往前走。

（为什么往前走要画成斜的？因为纸是平的，我们要假装它是立体的。在纸上画立体的图，往前的方向看起来就是斜着往左上方的。这是一种约定俗成的画法。）

\textbf{第三步：画 $z$ 轴。}

从同一个起点，画一条往正上方的箭头。在箭头旁边写上 $z$。这条线代表往上走。

\textbf{第四步：标出原点。}

三条线交汇的那个点，就是原点。在旁边写上 $(0, 0, 0)$。

现在纸上应该有三条线从一个点出发：$x$ 往右，$y$ 往左上方斜着，$z$ 往正上方。

\subsubsection{标出四面体的四个顶点}

坐标轴画好了！现在来标出四面体的四个角。

我们的四面体有四个顶点。

\textbf{第一个顶点：$(0, 0, 0)$。}

这就是原点，已经标好了！三条轴交汇的那个点就是。

\textbf{第二个顶点：$(1, 0, 0)$。}

从原点出发，沿着 $x$ 轴往右走 1 步。在 $x$ 轴上量 1 个单位的长度，画一个点，旁边写上 $(1, 0, 0)$。这个点在 $x$ 轴上，距离原点 1 个单位。

\textbf{第三个顶点：$(0, 1, 0)$。}

从原点出发，沿着 $y$ 轴往前走 1 步。在 $y$ 轴上量 1 个单位的长度，画一个点，旁边写上 $(0, 1, 0)$。

\textbf{第四个顶点：$(0, 0, 1)$。}

从原点出发，沿着 $z$ 轴往上走 1 步。在 $z$ 轴上量 1 个单位的长度，画一个点，旁边写上 $(0, 0, 1)$。

现在纸上应该有四个点了。一个在中间（原点），一个在右边（$x$ 轴上），一个在左上方（$y$ 轴上），一个在正上方（$z$ 轴上）。

\subsubsection{把四个顶点连起来}

四个顶点标好了，现在把它们连起来！四面体有 6 条边，就是把四个顶点两两相连。

\textbf{第一条边：连接 $(0,0,0)$ 和 $(1,0,0)$。}

一个在原点，一个在右边。画一条直线连起来。这两个点都在地面上。

\textbf{第二条边：连接 $(0,0,0)$ 和 $(0,1,0)$。}

一个在原点，一个在前面。画一条直线连起来。这两个点也都在地面上。

\textbf{第三条边：连接 $(0,0,0)$ 和 $(0,0,1)$。}

一个在地面的原点，一个在空中正上方。画一条竖直的直线连起来。

\textbf{第四条边：连接 $(1,0,0)$ 和 $(0,1,0)$。}

一个在右边，一个在前面，都在地面上。画一条直线连起来。这条线在地面上，从右边斜着连到前面。

\textbf{第五条边：连接 $(1,0,0)$ 和 $(0,0,1)$。}

一个在地面右边，一个在空中。画一条直线连起来。

\textbf{第六条边：连接 $(0,1,0)$ 和 $(0,0,1)$。}

一个在地面前面，一个在空中。画一条直线连起来。

画完之后，看起来像一个三角锥！没错，四面体就是三角锥。四个三角形的面围成一个立体。

\subsubsection{检查你的四面体}

画好了，来检查一下。

\textbf{数顶点：4 个。}

$(0, 0, 0)$，$(1, 0, 0)$，$(0, 1, 0)$，$(0, 0, 1)$。

\textbf{数边：6 条。}

四个顶点两两相连，一共 6 条边，看看图上是不是 6 条线都画了。

\textbf{数面：4 个。}

第一个面：\textbf{底面}。由 $(0,0,0)$、$(1,0,0)$、$(0,1,0)$ 三个点围成的三角形。这个面在地面上。

第二个面：\textbf{后面}。由 $(0,0,0)$、$(1,0,0)$、$(0,0,1)$ 三个点围成的三角形。

第三个面：\textbf{左面}。由 $(0,0,0)$、$(0,1,0)$、$(0,0,1)$ 三个点围成的三角形。

第四个面：\textbf{斜面}。由 $(1,0,0)$、$(0,1,0)$、$(0,0,1)$ 三个点围成的三角形。

特别注意第四个面——那个斜面。

\subsubsection{理解斜面}

现在看着自己的图，注意那个斜面。

四面体有四个面，其中三个面是“靠墙”的。

底面靠着地面，就是 $z = 0$ 的平面。

后面靠着后墙，就是 $y = 0$ 的平面。

左面靠着左墙，就是 $x = 0$ 的平面。

但是第四个面——由 $(1,0,0)$、$(0,1,0)$、$(0,0,1)$ 围成的那个三角形——它是斜着的，不靠着任何一面墙。

这个斜面有什么特别的？我们来看看这三个顶点有什么共同点。

$(1, 0, 0)$：把三个数加起来，$1 + 0 + 0 = 1$。

$(0, 1, 0)$：把三个数加起来，$0 + 1 + 0 = 1$。

$(0, 0, 1)$：把三个数加起来，$0 + 0 + 1 = 1$。

三个点的 $x + y + z$ 都等于 1！

这个斜面就是所有满足 \textbf{$x + y + z = 1$} 的点组成的平面。

注意到了吗？$x + y + z = 1$……这正好就是我们要积分的函数 $x + y + z$ 等于常数 1 的情况。这个斜面和我们要积分的函数有着天然的联系。

\subsubsection{四面体里面是什么？}

外壳画好了，但里面呢？四面体内部的点要满足什么条件？

看着图，想象四面体内部的空间。

在四面体内部的任何一点，都要满足四个条件：

\textbf{第一个：$x \geq 0$。} 不能穿过左墙。

\textbf{第二个：$y \geq 0$。} 不能穿过后墙。

\textbf{第三个：$z \geq 0$。} 不能穿到地下。

\textbf{第四个：$x + y + z \leq 1$。} 不能穿过斜面。

来试试：\textbf{(0.2, 0.2, 0.2)} 在不在里面？

检查一下：$x = 0.2 \geq 0$，通过！$y = 0.2 \geq 0$，通过！$z = 0.2 \geq 0$，通过！$x + y + z = 0.6 \leq 1$，通过！

四个条件都满足，这个点在里面！在图上，它大概在原点和斜面之间，比较靠近原点那边。

那 \textbf{(0.5, 0.5, 0.5)} 呢？

检查一下：$x = 0.5 \geq 0$，通过。$y = 0.5 \geq 0$，通过。$z = 0.5 \geq 0$，通过。$x + y + z = 1.5 \leq 1$……等等，$1.5$ 比 $1$ 大！不通过！

这个点在四面体外面了。它穿过了斜面。$0.5 + 0.5 + 0.5 = 1.5 > 1$，超出了斜面的范围。

所以斜面 $x + y + z = 1$ 就像一道门，超过了就出界。

\subsubsection{画图总结}

好，现在四面体已经画好了。

四个顶点在哪儿，六条边怎么连，四个面长什么样，都已经亲手画过了。

特别是那个斜面 $x + y + z = 1$，要记住。

把这张图放在旁边。接下来算积分的时候，随时看着它，会更容易理解。

\subsubsection{确定积分范围}

画好了四面体，现在要解决最关键的问题：积分的范围是什么？

看着你的图。我们要搞清楚 $x$、$y$、$z$ 各自能取什么值。

\textbf{先看 $x$ 的范围。}

$x$ 最小是 0，因为 $x = 0$ 是左墙，是边界。

$x$ 最大是多少呢？看图，$x$ 轴上最远的那个顶点是 $(1, 0, 0)$。所以 $x$ 最大是 1。

\textbf{$x$ 的范围：从 0 到 1。} 这个范围是固定的，不依赖于别的变量。

\textbf{再看 $y$ 的范围。}

这里出现了一个新东西：$y$ 的范围不是固定的，它取决于 $x$ 已经是多少。

为什么？看图。当 $x = 0$ 的时候，在原点这一侧。$y$ 可以从 0 一直走到 $(0, 1, 0)$，也就是 $y$ 可以到 1。

但是当 $x = 0.5$ 的时候，已经往右走了一半。如果 $y$ 还想到 1，那 $x + y = 1.5$，已经超过斜面了！所以 $y$ 最多只能到 0.5。

当 $x = 1$ 的时候，已经到了 $(1, 0, 0)$ 这个顶点。$y$ 只能是 0，否则就出界了。

规律就是：斜面方程 $x + y + z = 1$，当 $z = 0$ 的时候，变成 $x + y = 1$，也就是 $y = 1 - x$。

\textbf{$y$ 的范围：从 0 到 $1 - x$。}

\textbf{最后看 $z$ 的范围。}

同样的道理。$z$ 的范围取决于 $x$ 和 $y$ 已经是多少。

斜面方程 $x + y + z = 1$ 告诉我们 $z$ 最大是 $1 - x - y$。

举几个例子：当 $x = 0$，$y = 0$ 时，$z$ 可以从 0 到 1。当 $x = 0.2$，$y = 0.3$ 时，$z$ 可以从 0 到 0.5。当 $x = 0.5$，$y = 0.5$ 时，$z$ 只能是 0。

\textbf{$z$ 的范围：从 0 到 $1 - x - y$。}

注意看这三个范围的特点：$x$ 的范围是固定的，$y$ 的范围依赖于 $x$，$z$ 的范围依赖于 $x$ 和 $y$。这比立方体复杂多了——立方体里每个变量的范围都是固定的。

这就是四面体比立方体难算的根本原因。

\subsubsection{写出三重积分}

把范围写出来，完整的三重积分是：

\[
\int_0^1 \int_0^{1-x} \int_0^{1-x-y} (x + y + z) \, dz \, dy \, dx
\]

这个式子从右往左读。

最里面是对 $z$ 积分，$z$ 从 0 到 $1-x-y$。

中间是对 $y$ 积分，$y$ 从 0 到 $1-x$。

最外面是对 $x$ 积分，$x$ 从 0 到 1。

我们从里往外算，先算 $z$，再算 $y$，最后算 $x$。

好了，开工！

\subsubsection{第一步：对 $z$ 积分}

$x$ 和 $y$ 先固定不动。

我们要算最里面的积分：

\[
\int_0^{1-x-y} (x + y + z) \, dz
\]

在这一步，$x$ 和 $y$ 都被当成常数，只有 $z$ 在变。

\textbf{求不定积分：}

$x + y + z$ 对 $z$ 积分是什么？把 $x + y$ 当成一个常数。

常数对 $z$ 积分，就是常数乘以 $z$：

\[
\int (x + y) \, dz = (x+y)z
\]

$z$ 对 $z$ 积分：

\[
\int z \, dz = \frac{z^2}{2}
\]

加起来：

\[
\int (x + y + z) \, dz = (x+y)z + \frac{z^2}{2}
\]

\textbf{代入上下限：}

上限是 $z = 1 - x - y$，下限是 $z = 0$。

代入上限 $z = 1 - x - y$。这个有点长，我们慢慢算。

设 $s = x + y$，那么 $1 - x - y = 1 - s$。

先算 $(x+y)z$ 这部分：

\[
s \cdot (1 - s) = s - s^2
\]

再算 $\frac{z^2}{2}$ 这部分：

\[
\frac{(1 - s)^2}{2}
\]

\[
(1-s)^2 = 1 - 2s + s^2
\]

\[
\frac{(1-s)^2}{2} = \frac{1}{2} - s + \frac{s^2}{2}
\]

两部分加起来：

\[
(s - s^2) + \left(\frac{1}{2} - s + \frac{s^2}{2}\right)
\]

一项一项整理：

$s$ 项：$s - s = 0$，消掉了。

$s^2$ 项：$-s^2 + \frac{s^2}{2} = -\frac{s^2}{2}$。

常数项：$\frac{1}{2}$。

所以上限处的值是 $\frac{1}{2} - \frac{s^2}{2}$，也就是 $\frac{1}{2} - \frac{(x+y)^2}{2}$。

代入下限 $z = 0$：$(x+y) \cdot 0 + \frac{0}{2} = 0$。

上限减下限：$\frac{1}{2} - \frac{(x+y)^2}{2}$。

\textbf{第一步的结果：}

\[
\int_0^{1-x-y} (x + y + z) \, dz = \frac{1}{2} - \frac{(x+y)^2}{2}
\]

好，第一遍扫完了！

\subsubsection{第二步：对 $y$ 积分}

现在轮到 $y$ 了。$x$ 继续固定不动。

我们要算：

\[
\int_0^{1-x} \left(\frac{1}{2} - \frac{(x+y)^2}{2}\right) \, dy
\]

在这一步，$x$ 被当成常数，只有 $y$ 在变。

\textbf{先展开 $(x+y)^2$：}

$(x+y)^2 = x^2 + 2xy + y^2$。所以：

\[
\frac{(x+y)^2}{2} = \frac{x^2}{2} + xy + \frac{y^2}{2}
\]

那么：

\[
\frac{1}{2} - \frac{(x+y)^2}{2} = \frac{1}{2} - \frac{x^2}{2} - xy - \frac{y^2}{2}
\]

一共四项，我们一项一项积分。

\textbf{第一项：$\frac{1}{2}$ 对 $y$ 积分}

$\frac{1}{2}$ 是常数。

不定积分是 $\frac{y}{2}$

代入上限 $y = 1-x$：$\frac{1-x}{2}$

代入下限 $y = 0$：$0$

相减：$\frac{1-x}{2}$

\textbf{第二项：$\frac{x^2}{2}$ 对 $y$ 积分}

$\frac{x^2}{2}$ 也是常数（因为 $x$ 不动）。

不定积分是 $\frac{x^2 y}{2}$

代入上限 $y = 1-x$：$\frac{x^2(1-x)}{2} = \frac{x^2 - x^3}{2}$

代入下限 $y = 0$：$0$

相减：$\frac{x^2 - x^3}{2}$

\textbf{第三项：$xy$ 对 $y$ 积分}

$x$ 是常数，所以这相当于 $x \cdot \int y \, dy$

不定积分是 $\frac{xy^2}{2}$

代入上限 $y = 1-x$：$\frac{x(1-x)^2}{2}$

代入下限 $y = 0$：$0$

相减：$\frac{x(1-x)^2}{2}$

\textbf{第四项：$\frac{y^2}{2}$ 对 $y$ 积分}

不定积分是 $\frac{y^3}{6}$

代入上限 $y = 1-x$：$\frac{(1-x)^3}{6}$

代入下限 $y = 0$：$0$

相减：$\frac{(1-x)^3}{6}$

\textbf{合起来（注意正负号）：}

原式是 $+\frac{1}{2} - \frac{x^2}{2} - xy - \frac{y^2}{2}$

所以积分结果是：

\[
+\frac{1-x}{2} - \frac{x^2 - x^3}{2} - \frac{x(1-x)^2}{2} - \frac{(1-x)^3}{6}
\]

这也太长了，能化简吗？

当然能。虽然看起来复杂，但我们一步一步展开，再合并同类项就行了。

\textbf{化简过程：}

我们需要先算出 $(1-x)^2$ 和 $(1-x)^3$。

\[
(1-x)^2 = 1 - 2x + x^2
\]

\[
(1-x)^3 = (1-x)(1-2x+x^2) = 1 - 3x + 3x^2 - x^3
\]

来验算一下 $(1-x)^3$：$1 \cdot 1 = 1$，$1 \cdot (-2x) = -2x$，$1 \cdot x^2 = x^2$，$(-x) \cdot 1 = -x$，$(-x) \cdot (-2x) = 2x^2$，$(-x) \cdot x^2 = -x^3$。加起来：$1 - 3x + 3x^2 - x^3$。对了！

好，现在把四项都展开。

第一项：$\frac{1-x}{2} = \frac{1}{2} - \frac{x}{2}$

第二项：$-\frac{x^2 - x^3}{2} = -\frac{x^2}{2} + \frac{x^3}{2}$

第三项：$-\frac{x(1 - 2x + x^2)}{2} = -\frac{x - 2x^2 + x^3}{2} = -\frac{x}{2} + x^2 - \frac{x^3}{2}$

第四项：$-\frac{1 - 3x + 3x^2 - x^3}{6} = -\frac{1}{6} + \frac{x}{2} - \frac{x^2}{2} + \frac{x^3}{6}$

现在来合并同类项！

\textbf{常数项：} $\frac{1}{2} - \frac{1}{6} = \frac{3}{6} - \frac{1}{6} = \frac{1}{3}$

\textbf{$x$ 项：} $-\frac{1}{2} - \frac{1}{2} + \frac{1}{2} = -\frac{1}{2}$

\textbf{$x^2$ 项：} $-\frac{1}{2} + 1 - \frac{1}{2} = 0$ ——消掉了！

\textbf{$x^3$ 项：} $\frac{1}{2} - \frac{1}{2} + \frac{1}{6} = \frac{1}{6}$

漂亮！$x^2$ 项消掉了，结果比想象的简单。

\textbf{第二步的结果：}

\[
\int_0^{1-x} \left(\frac{1}{2} - \frac{(x+y)^2}{2}\right) \, dy = \frac{1}{3} - \frac{x}{2} + \frac{x^3}{6}
\]

第二遍扫完了！现在只剩最后一步。

\subsubsection{第三步：对 $x$ 积分}

压轴的来了。

我们要算：

\[
\int_0^1 \left(\frac{1}{3} - \frac{x}{2} + \frac{x^3}{6}\right) \, dx
\]

这是最后一步了，没有什么东西是常数了，就是普通的一元积分。而且上下限是固定的 0 到 1，清清爽爽。

\textbf{对每一项分别积分：}

第一项：$\int \frac{1}{3} \, dx = \frac{x}{3}$

第二项：$\int \frac{x}{2} \, dx = \frac{x^2}{4}$

第三项：$\int \frac{x^3}{6} \, dx = \frac{x^4}{24}$

不定积分合起来：$\frac{x}{3} - \frac{x^2}{4} + \frac{x^4}{24}$

\textbf{代入上下限：}

代入上限 $x = 1$：$\frac{1}{3} - \frac{1}{4} + \frac{1}{24}$

三个分数要通分。最小公分母是 24。

\[
\frac{1}{3} = \frac{8}{24}
\]

\[
\frac{1}{4} = \frac{6}{24}
\]

\[
\frac{1}{24} = \frac{1}{24}
\]

\[
\frac{8}{24} - \frac{6}{24} + \frac{1}{24} = \frac{8 - 6 + 1}{24} = \frac{3}{24} = \frac{1}{8}
\]

代入下限 $x = 0$：$0$

上限减下限：$\frac{1}{8} - 0 = \frac{1}{8}$

\subsubsection{最终答案}

\[
\int_0^1 \int_0^{1-x} \int_0^{1-x-y} (x + y + z) \, dz \, dy \, dx = \frac{1}{8}
\]

这个四面体上的三重积分结果是 \textbf{$\frac{1}{8}$}！

\subsubsection{回顾：为什么四面体比立方体难？}

来对比一下。

\textbf{立方体：}

立方体的边界是直直的。$x$ 从 0 到 1，$y$ 从 0 到 1，$z$ 从 0 到 1。每个变量的范围都是固定的数，互不影响。

\textbf{四面体：}

看你画的四面体，边界是斜的。$x$ 从 0 到 1 还算正常，但 $y$ 从 0 到 $1-x$——依赖于 $x$，$z$ 从 0 到 $1-x-y$——依赖于 $x$ 和 $y$。后面的变量范围依赖于前面的变量，所以积分完之后会出现更复杂的表达式。

这就是为什么四面体的计算量大得多。但方法还是一样的：一次只对一个变量积分，把其他的当成常数。耐心一点，一步一步来，就能算出来。

\subsubsection{完整计算过程总结}

\begin{center}
\begin{tabular}{c|c|c}
步骤 & 积分 & 结果 \\
\hline
第一步 & $\displaystyle\int_0^{1-x-y} (x + y + z) \, dz$ & $\dfrac{1}{2} - \dfrac{(x+y)^2}{2}$ \\
第二步 & $\displaystyle\int_0^{1-x} \left(\dfrac{1}{2} - \dfrac{(x+y)^2}{2}\right) \, dy$ & $\dfrac{1}{3} - \dfrac{x}{2} + \dfrac{x^3}{6}$ \\
第三步 & $\displaystyle\int_0^1 \left(\dfrac{1}{3} - \dfrac{x}{2} + \dfrac{x^3}{6}\right) \, dx$ & $\dfrac{1}{8}$
\end{tabular}
\end{center}

\section{公式汇总}

\subsection{一元微积分（$y$ 和 $z$ 都是常数）}

\begin{center}
\begin{tabular}{c|c|c}
操作 & 公式 & 答案 \\
\hline
求导 & $\dfrac{d}{dx}(x + y + z)$ & $1$ \\
不定积分 & $\displaystyle\int (x + y + z) \, dx$ & $\dfrac{x^2}{2} + yx + zx + C$ \\
定积分 & $\displaystyle\int_a^b (x + y + z) \, dx$ & $\left(\dfrac{b^2}{2} + yb + zb\right) - \left(\dfrac{a^2}{2} + ya + za\right)$
\end{tabular}
\end{center}

\subsection{二元微积分（$z$ 是常数）}

\begin{center}
\begin{tabular}{c|c|c}
操作 & 公式 & 答案 \\
\hline
对 $x$ 偏导 & $\dfrac{\partial}{\partial x}(x + y + z)$ & $1$ \\
对 $y$ 偏导 & $\dfrac{\partial}{\partial y}(x + y + z)$ & $1$ \\
二重积分（正方形 $[0,1]\times[0,1]$，$z=5$） & $\displaystyle\int_0^1 \int_0^1 (x + y + 5) \, dx \, dy$ & $6$
\end{tabular}
\end{center}

\subsection{三元微积分（$x$、$y$、$z$ 都是变量）}

\begin{center}
\begin{tabular}{c|c|c}
操作 & 公式 & 答案 \\
\hline
对 $x$ 偏导 & $\dfrac{\partial}{\partial x}(x + y + z)$ & $1$ \\
对 $y$ 偏导 & $\dfrac{\partial}{\partial y}(x + y + z)$ & $1$ \\
对 $z$ 偏导 & $\dfrac{\partial}{\partial z}(x + y + z)$ & $1$ \\
梯度 & $\nabla(x + y + z)$ & $(1, 1, 1)$ \\
三重积分（单位立方体 $[0,1]^3$） & $\displaystyle\int_0^1 \int_0^1 \int_0^1 (x + y + z) \, dx \, dy \, dz$ & $\dfrac{3}{2}$ \\
三重积分（立方体 $[0,2]^3$） & $\displaystyle\int_0^2 \int_0^2 \int_0^2 (x + y + z) \, dx \, dy \, dz$ & $24$ \\
三重积分（长方体 $[0,1]\times[0,2]\times[0,3]$） & $\displaystyle\int_0^3 \int_0^2 \int_0^1 (x + y + z) \, dx \, dy \, dz$ & $18$ \\
三重积分（四面体） & $\displaystyle\int_0^1 \int_0^{1-x} \int_0^{1-x-y} (x + y + z) \, dz \, dy \, dx$ & $\dfrac{1}{8}$
\end{tabular}
\end{center}

\section{$x + y$ 和 $x + y + z$ 的对比}

\begin{center}
\begin{tabular}{c|c|c}
项目 & $x + y$ & $x + y + z$ \\
\hline
变量个数 & 2 个 & 3 个 \\
所有偏导数 & 都是 1 & 都是 1 \\
梯度 & $(1, 1)$ & $(1, 1, 1)$ \\
单位区域上的积分 & 1（正方形） & 3/2（立方体） \\
“三角形/四面体”上的积分 & 1/3 & 1/8
\end{tabular}
\end{center}

看到规律了吗？

维度越高，积分要做的次数越多。

但基本方法都是一样的：一次只对一个变量积分，把其他的当成常数。

\section{结语}

$x + y + z$ 和 $x + y$ 一样，是最简单的三元函数。

当你学三重积分的时候，教科书会用 $x + y + z$ 让你练手。

因为它的偏导数都是 1，积分也很直接，不会让你在计算中迷失方向。

等你用 $x + y + z$ 把方法学会了，再去挑战更复杂的函数。

\section{你学到了什么}

读完这一篇，你知道了这些事情：

\textbf{第一}，$x + y + z$ 就是三个会变的数加在一起。

\textbf{第二}，一元微积分里，只有 $x$ 动，$y$ 和 $z$ 都是常数。导数是 1，积分是 $\frac{x^2}{2} + yx + zx + C$。

\textbf{第三}，二元微积分里，$x$ 和 $y$ 都动，$z$ 是常数。两个偏导数都是 1。

\textbf{第四}，三元微积分里，$x$、$y$、$z$ 都动。三个偏导数都是 1，梯度是 $(1, 1, 1)$。

\textbf{第五}，三重积分就是积分做三次：先扫一个方向，再扫第二个方向，最后扫第三个方向。

\textbf{第六}，四面体的积分比立方体难，因为边界是斜的，积分上限会变化。画一张图放在旁边，随时对照着看，会更容易理解。

\textbf{第七}，$x + y + z$ 是最简单的三元函数，用来帮你学会方法。

\chapter{$xy$：乘法搭档闯微积分}

\section{开场：从 AA 制到合伙做生意}

还记得 $x$ 和 $y$ 吗？

上一篇里，它们俩是 \textbf{AA 制}。$x$ 出 3 块钱，$y$ 出 7 块钱，$x + y = 10$ 块钱。各出各的，互不干扰。

但这一篇，它们决定\textbf{合伙做生意}了。

$x$ 负责卖柠檬水，今天卖了 $x$ 杯。$y$ 负责定价，每杯 $y$ 块钱。

那 \textbf{$xy$} 呢？就是一共赚了多少钱。

今天卖了 4 杯，每杯定价 5 块钱。

$xy = 4 \times 5 = 20$ 块钱。

就这么简单。两个数相乘，仅此而已。

\section{🎨 动手画一画：$xy$ 就是长方形的面积}

拿出一张纸和一支笔。

\textbf{第一步：} 在纸上画一个长方形。不用太大，占纸的四分之一就行。

\textbf{第二步：} 在长方形的底边（横着的那条边）下面写上 \textbf{$x$}。

\textbf{第三步：} 在长方形的左边（竖着的那条边）旁边写上 \textbf{$y$}。

\textbf{第四步：} 在长方形的肚子里写上大大的 \textbf{$xy$}。

你画好了吗？看看你的图——

底边是 $x$（卖了几杯），左边是 $y$（每杯多少钱），肚子里的面积 $xy$ 就是\textbf{总收入}。

这就是 $xy$ 最直观的样子：\textbf{一个长方形的面积}。

记住这张图，后面会反复用到它。

但接下来，神奇的事情发生了。

$x$ 多卖了 1 杯（从 4 变成 5），赚的钱多了多少？

多了 \textbf{5 块钱}——刚好就是 $y$ 的值，也就是每杯的价格！

$y$ 涨了 1 块（从 5 变成 6），赚的钱多了多少？

多了 \textbf{4 块钱}——刚好就是 $x$ 的值，也就是卖的杯数！

原来每多卖一杯，赚多少，取决于定了多高的价。每涨一块钱，赚多少，取决于卖了多少杯。

\textbf{这就是 $xy$ 和 $x + y$ 最大的不同。}

\textbf{$x + y$：两个数各干各的，互不相干。}

\textbf{$xy$：两个数绑在一起，互相成就。}

听起来有点复杂？别怕。我们慢慢来。

\section{一、$xy$ 对 $x$ 求导}

\subsection{老规矩：$y$ 先别动，当常数。}

\subsection{用数字看一看}

假设 $y$ 今天是 5（固定不动），那 $f(x) = 5x$。

\textbf{实验 1：}

$x$ 从 2 变到 3，变了 1。

$f(x)$ 呢？从 $5 \times 2 = 10$ 变到 $5 \times 3 = 15$，变了 5。

$5 \div 1 = 5$！

\textbf{实验 2：}

$x$ 从 10 变到 13，变了 3。

$f(x)$ 呢？从 $5 \times 10 = 50$ 变到 $5 \times 13 = 65$，变了 15。

$15 \div 3 = 5$！

\textbf{实验 3：}

$x$ 从 100 变到 100.1，变了 0.1。

$f(x)$ 呢？从 500 变到 500.5，变了 0.5。

$0.5 \div 0.1 = 5$！

不管 $x$ 怎么变，不管变了多少，答案永远是 5。

5 是什么？就是 $y$ 的值！

\section{🎨 动手画一画：为什么导数是 $y$}

翻到你刚才画的那张长方形的图。我们要在旁边画一张新的。

\textbf{第一步：} 再画一个长方形，和之前那个一样高，但是稍微宽一点点。

\textbf{第二步：} 在原来的长方形底边下面写 \textbf{$x = 3$}，在新的、更宽的长方形底边下面写 \textbf{$x + 1$}。

\textbf{第三步：} 两个长方形的左边对齐。它们的高度一样，都是 \textbf{$y = 5$}。

\textbf{第四步：} 现在看——新的长方形比旧的多出来一小条，在右边。用铅笔把那一小条涂上阴影。

\textbf{第五步：} 那一小条有多宽？宽 \textbf{1}（因为 $x$ 变到了 $x + 1$）。有多高？高 \textbf{$y = 5$}。在阴影条里写上 \textbf{$y \times 1 = y$}。

看看你的图——

$x$ 增加了 1，面积多出了一小条，这一小条的面积恰好就是 \textbf{$y$}。

所以 $xy$ 对 $x$ 的导数就是 $y$。

不管原来的长方形多宽，每次往右多长 1 格，多出来的那一条永远是 $y$ 那么大。这就是为什么导数是 $y$！

\subsection{换个 $y$ 再试试}

假设 $y$ 今天是 100，那 $f(x) = 100x$。

$x$ 从 1 变到 2，变了 1。

$f(x)$ 从 100 变到 200，变了 100。

$100 \div 1 = 100$。看，导数就是 $y$！

再来一个。$y = 0.5$，那 $f(x) = 0.5x$。

$x$ 从 6 变到 10，变了 4。

$f(x)$ 从 3 变到 5，变了 2。

$2 \div 4 = 0.5$。没错，导数还是 $y$！

\subsection{正式结论}

\[
\frac{d}{dx}(xy) = y
\]

\textbf{$x$ 变化 1 个单位，答案就变化 $y$ 个单位。$y$ 虽然不动，但它的大小决定了变化的速度。}

\textbf{就像柠檬水生意：每多卖 1 杯，收入增加多少？取决于每杯的价格。}

\subsection{和 $x + y$ 的对比}

\begin{center}
\begin{tabular}{c|c|c}
函数 & 对 $x$ 求导 & 含义 \\
\hline
$x + y$ & 1 & $x$ 变 1，答案就变 1，跟 $y$ 多大没关系 \\
$xy$ & $y$ & $x$ 变 1，答案变 $y$。$y$ 越大，变得越猛
\end{tabular}
\end{center}

$x + y$ 的导数永远是 1，太无聊了。$xy$ 就有意思了，答案取决于搭档的实力！

\subsection{例子}

\textbf{例子 1：} $y = 7$

\[
f(x) = 7x
\]

\[
\frac{d}{dx}(7x) = 7
\]

验算：$x$ 从 3 变到 5（变了 2），$f(x)$ 从 21 变到 35（变了 14）。$14 \div 2 = 7$。完美！

\textbf{例子 2：} $y = -3$（$y$ 变成了负数）

\[
f(x) = -3x
\]

\[
\frac{d}{dx}(-3x) = -3
\]

负数意味着 $x$ 增大时，答案反而减小了。$y$ 是负数的时候，把一切都搞反了！

验算：$x$ 从 1 变到 4（变了 3），$f(x)$ 从 $-3$ 变到 $-12$（变了 $-9$）。$-9 \div 3 = -3$。没错！

\section{特别插曲：积分符号上下的数字是什么意思？}

在继续讲积分之前，我们要先解决一个可能困惑你很久的问题。

你可能已经注意到，有些公式长这样：

\[
\int_0^2 3x \, dx
\]

这个弯弯的大符号 $\int$ 上面蹲着一个 \textbf{2}，下面蹲着一个 \textbf{0}。

它们是谁？在那里干什么？

\subsection{想象一条马路}

想象一条笔直的马路，马路上每隔一段就有一根路灯柱，上面标着数字：0 号、1 号、2 号、3 号……一直排下去。

积分的工作是\textbf{沿着这条马路收集东西}。

但不可能收集整条马路上的所有东西——那条路无限长，收不完。

所以必须指定：\textbf{从哪里开始收集，到哪里停下来}。

\textbf{下面的数字}就是“出发点”——从这里开始收集。

\textbf{上面的数字}就是“终点”——走到这里就停下。

\[
\int_{\text{从这里开始}}^{\text{到这里停下}} \text{收集的东西} \, dx
\]

\section{🎨 动手画一画：积分的起点和终点}

\textbf{第一步：} 在纸上画一条水平的线，这是一条“数轴”（就是一条标着数字的马路）。

\textbf{第二步：} 在这条线上均匀地标上几个数字：0、1、2、3、4、5。像尺子上的刻度一样。

\textbf{第三步：} 在数字 \textbf{0} 的位置画一面小旗子（或者一个向下的小箭头），旁边写上“\textbf{出发！}”

\textbf{第四步：} 在数字 \textbf{2} 的位置画一面小旗子，旁边写上“\textbf{停下！}”

\textbf{第五步：} 用彩色笔或者阴影，把 0 到 2 之间的那一段涂上颜色。

看看你的图——涂了颜色的那一段，就是积分要工作的范围。

所以 $\int_0^2$ 的意思是：\textbf{从 0 号路灯出发，走到 2 号路灯停下}。

\subsection{再举几个例子}

\[
\int_1^4
\]

从 1 号路灯出发，走到 4 号路灯停下。

\textbf{在你的数轴上，把 1 到 4 之间涂上不同的颜色，试试看。}

\[
\int_0^{10}
\]

从 0 出发，走到 10 停下。范围很大，要收集一大段。

\[
\int_3^3
\]

从 3 出发……走到 3 停下？那不就是原地不动吗！

没错！从 3 到 3，距离为零，什么也没收集到。所以 $\int_3^3$ 任何东西 $dx = 0$。

连门都没出，当然什么都没收集到。

\subsection{那具体“收集”了什么呢？}

光知道“从哪到哪”还不够。你还得知道沿途有什么东西可以收集。

$\int_0^2 3x \, dx$ 的意思是：

\begin{itemize}
\item \textbf{从 $x = 0$ 走到 $x = 2$}
\item \textbf{沿途收集 $3x$ 这个东西}
\end{itemize}

当 $x = 0$ 时，$3x = 0$（这里什么都没有）。

当 $x = 1$ 时，$3x = 3$（这里有 3 个东西）。

当 $x = 2$ 时，$3x = 6$（这里有 6 个东西）。

从 0 走到 2，沿路把 $3x$ 一点一点捡起来，全部加在一起。

最后能收集到多少？这就是定积分要回答的问题。

\subsection{回忆一下“上下没有数字”的情况}

\[
\int 3x \, dx = \frac{3x^2}{2} + C
\]

这个 $\int$ 上下\textbf{没有数字}。

这叫\textbf{不定积分}。意思是：\textbf{不指定范围，只问“收集规则是什么”}。

就像写了一本《收集手册》：“不管你从哪走到哪，$3x$ 的收集方式都是 $\frac{3x^2}{2}$。”

而上下\textbf{有数字}的叫\textbf{定积分}，是指定了具体范围，算出一个\textbf{确切的数}。

\begin{center}
\begin{tabular}{c|c|c|c}
写法 & 名字 & 意思 & 结果 \\
\hline
$\int 3x \, dx$ & 不定积分 & 收集手册（规则） & 一个公式 + C \\
$\int_0^2 3x \, dx$ & 定积分 & 从 0 到 2 具体收集 & 一个数字
\end{tabular}
\end{center}

\subsection{上限减下限：怎么从“手册”算出“具体数字”}

定积分的计算方法：

\textbf{第一步：} 先写出不定积分（查手册）。

\textbf{第二步：} 把\textbf{上面的数字}代入手册公式，算出一个值。

\textbf{第三步：} 把\textbf{下面的数字}代入手册公式，算出一个值。

\textbf{第四步：} \textbf{上面的值减去下面的值}，就是答案。

为什么是“上减下”？

这就像看水表。走到终点时读一次表，回头看看出发时的读数，相减就知道这一段用了多少水。

\section{🎨 动手画一画：上限减下限}

\textbf{第一步：} 画一条数轴，标上 0、1、2、3。

\textbf{第二步：} 在数轴上方画一条斜线。这条斜线从左下往右上走——$x$ 越大，线越高。（这代表 $f(x) = 3x$，$x$ 越大值越大。）

\textbf{第三步：} 在 $x = 0$ 的位置，斜线的高度是 0（因为 $3 \times 0 = 0$）。在 $x = 1$ 的位置，高度是 3。在 $x = 2$ 的位置，高度是 6。把这几个点标出来。

\textbf{第四步：} 从 $x = 0$ 到 $x = 2$，把斜线下方、数轴上方的区域涂上阴影。

看看你的图——涂了阴影的区域是一个\textbf{三角形}。

这个三角形的面积，就是 $\int_0^2 3x \, dx$ 的值！

底边长 2，高度 6，三角形面积 $= \frac{1}{2} \times 2 \times 6 = 6$。

现在用公式验算：

\[
\int_0^2 3x \, dx = \left[\frac{3x^2}{2}\right]_0^2 = \frac{3 \times 4}{2} - \frac{3 \times 0}{2} = 6 - 0 = 6
\]

一模一样！

\textbf{积分的几何意思就是：曲线下方的面积。}

上下有数字，就是告诉你从哪里量到哪里。

\subsection{那个中括号加上下标是什么意思？}

你可能注意到了这个写法：

\[
\left[\frac{3x^2}{2}\right]_0^2
\]

这个中括号 $[ \;]$ 的右边，下面写着 0，上面写着 2。

它的意思特别简单：

\textbf{“先把 $x = 2$ 代进去，再把 $x = 0$ 代进去，然后上面的减下面的。”}

\[
\left[\frac{3x^2}{2}\right]_0^2 = \frac{3 \times 2^2}{2} - \frac{3 \times 0^2}{2} = 6 - 0 = 6
\]

就是一个简写而已。省得每次都要写“把 $x =$ 多少代进去”这么一长串话。

这个中括号是一个速记符号。左边是公式，右上是终点，右下是起点。一目了然。

\section{二、$xy$ 的积分}

\subsection{什么叫“把常数提出来”}

$y$ 还是固定不动。那 $y$ 就是个常数，可以把它\textbf{提}出来！

打个比方：

假如 $y = 3$，那收集 $3x$，就等于收集 $x$ 之后再乘以 3。

就像去果园摘苹果。每棵树上有 3 个苹果，你先数有多少棵树，最后乘以 3 就行了。没必要一个苹果一个苹果地数。

\[
\int xy \, dx = y \int x \, dx
\]

$y$ 被提到了积分号外面，因为它是常数，在 $x$ 的世界里纹丝不动。

\subsection{完成计算}

$x$ 的积分是 $x^2/2$（还记得吗？）：

\[
\int xy \, dx = y \cdot \frac{x^2}{2} + C = \frac{x^2 y}{2} + C
\]

\subsection{验算}

把答案求导，应该得到 $xy$。

\[
\frac{d}{dx}\left(\frac{x^2 y}{2} + C\right) = xy \quad \checkmark
\]

看，求导把积分装好的东西又拆回去了，一模一样！

\subsection{例子（不定积分）}

\textbf{例子 1：} $y = 4$

\[
\int 4x \, dx = 4 \cdot \frac{x^2}{2} + C = 2x^2 + C
\]

验算：$(2x^2)' = 4x$。没错！

\textbf{例子 2：} $y = 10$

\[
\int 10x \, dx = 10 \cdot \frac{x^2}{2} + C = 5x^2 + C
\]

验算：$(5x^2)' = 10x$。完美！

\subsection{定积分来了}

现在给一个范围，只收集这一段。

\subsubsection{例子 1：$y = 3$，$x$ 从 0 到 2}

\[
\int_0^2 3x \, dx
\]

\textbf{第一步：} 不定积分（查手册）：$\frac{3x^2}{2}$

\textbf{第二步：} 把上面的数字 $x = 2$ 代进去：

\[
\frac{3 \times 4}{2} = 6
\]

\textbf{第三步：} 把下面的数字 $x = 0$ 代进去：

\[
\frac{3 \times 0}{2} = 0
\]

\textbf{第四步：} 上面的减下面的：

\[
6 - 0 = 6
\]

从 $x = 0$ 收集到 $x = 2$（就是沿着数轴从点 \textbf{0} 走到点 \textbf{2}，算这段线下面的面积），一共收集了 \textbf{6}！

\section{🎨 动手画一画：例子 1 的图}

\textbf{第一步：} 画一条数轴，标上 0、1、2。

\textbf{第二步：} 在 $x = 0$ 的位置标一个点，高度为 0（因为 $3 \times 0 = 0$）。

\textbf{第三步：} 在 $x = 1$ 的位置标一个点，高度为 3（因为 $3 \times 1 = 3$）。

\textbf{第四步：} 在 $x = 2$ 的位置标一个点，高度为 6（因为 $3 \times 2 = 6$）。

\textbf{第五步：} 把这三个点连成一条直线（它们恰好在一条线上）。

\textbf{第六步：} 把这条直线下面、数轴上面的区域涂上阴影。

看看你的图——这是一个三角形！底边长 2，高度 6。

三角形面积 $= \frac{1}{2} \times 2 \times 6 = 6$。

和我们算出来的一模一样！

\subsubsection{例子 2：$y = 5$，$x$ 从 1 到 3}

\[
\int_1^3 5x \, dx
\]

\textbf{第一步：} 不定积分：$\frac{5x^2}{2}$

\textbf{第二步：} $x = 3$ 代进去（上面的数字）：

\[
\frac{5 \times 9}{2} = \frac{45}{2} = 22.5
\]

\textbf{第三步：} $x = 1$ 代进去（下面的数字）：

\[
\frac{5 \times 1}{2} = \frac{5}{2} = 2.5
\]

\textbf{第四步：} 上面的减下面的：

\[
22.5 - 2.5 = 20
\]

从 $x = 1$ 收集到 $x = 3$（就是沿着数轴从点 \textbf{1} 走到点 \textbf{3}，算这段线下面的面积），一共收集了 \textbf{20}！

\section{🎨 动手画一画：例子 2 的图}

\textbf{第一步：} 画一条数轴，标上 0、1、2、3。

\textbf{第二步：} 画出 $f(x) = 5x$ 这条线。在 $x = 0$ 时高度 0，$x = 1$ 时高度 5，$x = 2$ 时高度 10，$x = 3$ 时高度 15。把这些点连成直线。

\textbf{第三步：} 注意！这次不是从 0 开始，而是从 1 开始。所以只把 \textbf{$x = 1$ 到 $x = 3$} 之间、直线下方的区域涂上阴影。

看看你的图——涂了阴影的部分是一个\textbf{梯形}（不是三角形，因为不是从 0 开始的）。

梯形的上底是 5（$x = 1$ 时的高度），下底是 15（$x = 3$ 时的高度），高是 2（从 1 到 3 的距离）。

梯形面积 $= \frac{(5 + 15) \times 2}{2} = \frac{40}{2} = 20$。

又对上了！

你看，积分就是在算面积。上下的数字告诉你量哪一段的面积，仅此而已。

\subsubsection{例子 3：$y = 2$，$x$ 从 0 到 4}

\[
\int_0^4 2x \, dx
\]

\textbf{第一步：} 不定积分：$\frac{2x^2}{2} = x^2$

哦？2 除以 2 直接消掉了！

\textbf{第二步：} $x = 4$ 代进去：$16$

\textbf{第三步：} $x = 0$ 代进去：$0$

\textbf{第四步：} $16 - 0 = 16$

这个特别好算，$y = 2$ 的时候分母刚好消掉了。

\section{🎨 动手画一画：例子 3 的图}

\textbf{第一步：} 画一条数轴，标上 0、1、2、3、4。

\textbf{第二步：} 画出 $f(x) = 2x$ 这条线。$x = 0$ 时高度 0，$x = 4$ 时高度 8。连成直线。

\textbf{第三步：} 把 $x = 0$ 到 $x = 4$ 之间、直线下方涂上阴影。

这是一个三角形，底边 4，高度 8。

面积 $= \frac{1}{2} \times 4 \times 8 = 16$。完美！

\subsection{一般公式（定积分）}

\[
\int_a^b xy \, dx = \left[\frac{x^2 y}{2}\right]_a^b = \frac{y(b^2 - a^2)}{2}
\]

把上面的数字 $b$ 和下面的数字 $a$ 代进去，上减下，搞定。

\section{三、二元微积分——$x$ 和 $y$ 都能动了}

从现在开始，$x$ 和 $y$ 都可以动了！

现在 $f(x, y) = xy$ 变得更精彩了。$x$ 动，答案变。$y$ 动，答案也变。而且——\textbf{一个变量动的时候，影响有多大，取决于另一个变量}。

\subsection{\texorpdfstring{对 $x$ 求偏导}{对 x 求偏导}}

老规矩，$y$ 先暂停，固定不动。

假设 $y = 5$，那 $f = 5x$。

$x$ 从 3 变到 4（变了 1），$f$ 从 15 变到 20（变了 5）。

$5 \div 1 = 5 = y$。

\[
\frac{\partial}{\partial x}(xy) = y
\]

\subsection{\texorpdfstring{对 $y$ 求偏导}{对 y 求偏导}}

现在 $x$ 别动，只看 $y$。

假设 $x = 4$，那 $f = 4y$。

$y$ 从 2 变到 5（变了 3），$f$ 从 8 变到 20（变了 12）。

$12 \div 3 = 4 = x$！

\[
\frac{\partial}{\partial y}(xy) = x
\]

\section{🎨 动手画一画：偏导数的意思}

我们要画两张图，体会“对 $x$ 偏导”和“对 $y$ 偏导”的不同。

\subsection{第一张图：对 $x$ 求偏导}

\textbf{第一步：} 画一个长方形，底边标 \textbf{$x = 3$}，左边标 \textbf{$y = 5$}。在肚子里写上面积 \textbf{15}。

\textbf{第二步：} 在它右边紧挨着画一个\textbf{窄窄的小条}，宽度为 \textbf{1}（$x$ 从 3 变成 4），高度还是 \textbf{$y = 5$}。

\textbf{第三步：} 在小条里写上面积 \textbf{5}。

\textbf{第四步：} 在小条旁边写上：\textbf{“$x$ 增加 1 $\to$ 面积增加 $y = 5$”}。

这就是对 $x$ 的偏导 $= y$ 的意思：\textbf{长方形往右边长一格，多出来的面积就是 $y$}。

\subsection{第二张图：对 $y$ 求偏导}

\textbf{第一步：} 画一个长方形，底边标 \textbf{$x = 4$}，左边标 \textbf{$y = 2$}。在肚子里写上面积 \textbf{8}。

\textbf{第二步：} 在它上面紧挨着画一个\textbf{矮矮的小条}，高度为 \textbf{1}（$y$ 从 2 变成 3），宽度还是 \textbf{$x = 4$}。

\textbf{第三步：} 在小条里写上面积 \textbf{4}。

\textbf{第四步：} 在小条旁边写上：\textbf{“$y$ 增加 1 $\to$ 面积增加 $x = 4$”}。

这就是对 $y$ 的偏导 $= x$ 的意思：\textbf{长方形往上长一格，多出来的面积就是 $x$}。

\subsection{重大发现：互相成就}

\begin{center}
\begin{tabular}{c|c|c}
操作 & 结果 & 含义 \\
\hline
对 $x$ 偏导 & $y$ & $x$ 每变化 1，答案变化 $y$ 个单位 \\
对 $y$ 偏导 & $x$ & $y$ 每变化 1，答案变化 $x$ 个单位
\end{tabular}
\end{center}

\textbf{你的偏导数是我，我的偏导数是你}——这就是互相成就。

就像跳双人舞：你跳得越好，我的每一步就越出彩。我跳得越好，你的每一步也越出彩。

\subsection{和 $x + y$ 的对比}

\begin{center}
\begin{tabular}{c|c|c}
 & $x + y$ & $xy$ \\
\hline
关系 & AA 制，各管各的 & 合伙人，互相影响 \\
对 $x$ 偏导 & 1（永远是 1） & $y$（取决于搭档的大小） \\
对 $y$ 偏导 & 1（永远是 1） & $x$（取决于搭档的大小） \\
一句话 & 谁变 1，答案就变 1 & 谁变 1，答案的变化量就是另一个
\end{tabular}
\end{center}

\section{四、二重积分——在一块面上收集}

以前沿着一条马路收集东西——一重积分，上下各一个数字，告诉你从哪走到哪。

现在要在\textbf{一整块地}上收集东西——二重积分。不光要知道 $x$ 从哪到哪，还要知道 $y$ 从哪到哪。

所以需要\textbf{两组上下数字}，一组管 $x$ 的范围，一组管 $y$ 的范围。

\[
\int_{\text{y 从哪}}^{\text{y 到哪}} \int_{\text{x 从哪}}^{\text{x 到哪}} \text{收集的东西} \, dx \, dy
\]

两个积分号，每个都有自己的上下数字。先积 $x$（里面那个），再积 $y$（外面那个）。

\subsection{例子 1：正方形 $[0, 1] \times [0, 1]$}

工作区域是一个正方形：$x$ 从 0 到 1，$y$ 从 0 到 1。

要收集的东西是 $xy$。

\[
\int_0^1 \int_0^1 xy \, dx \, dy
\]

\section{🎨 动手画一画：正方形区域}

\textbf{第一步：} 在纸上画一个正方形。

\textbf{第二步：} 底边的左端标上 \textbf{0}，右端标上 \textbf{1}。在底边下面写上 \textbf{$x$}。

\textbf{第三步：} 左边的下端标上 \textbf{0}，上端标上 \textbf{1}。在左边旁边写上 \textbf{$y$}。

\textbf{第四步：} 把正方形里面涂上淡淡的阴影。

\textbf{第五步：} 在正方形肚子里写上 \textbf{$xy$}。

看看你的图——这个正方形就是要收集的区域。要在这块地上把 $xy$ 的值一点一点全部加起来。

\subsection{第一步：先对 $x$ 积分（$y$ 固定不动）}

先沿着 $x$ 方向扫一遍。

\[
\int_0^1 xy \, dx
\]

$y$ 是常数，提到外面：

\[
y \int_0^1 x \, dx = y \cdot \left[\frac{x^2}{2}\right]_0^1 = y \cdot \left(\frac{1}{2} - 0\right) = \frac{y}{2}
\]

从 $x = 0$ 收集到 $x = 1$（就是固定某一行的 $y$ 值不动，从点 $(0, y)$ 水平走到点 $(1, y)$，算这一行上 $xy$ 的总量）。

第一遍扫完，收集到了 $\frac{y}{2}$。

\section{🎨 动手画一画：先扫 $x$ 方向}

回到你刚才画的正方形。

\textbf{第一步：} 在正方形内部，画几条\textbf{水平的虚线}，把正方形分成几条横条。

\textbf{第二步：} 选其中一条横条，用箭头标出方向：\textbf{从左到右}。旁边写“先沿 $x$ 方向扫”。

这就是“先对 $x$ 积分”的意思——对于每一个固定的 $y$ 值（每一条横条），沿着 $x$ 方向从 0 扫到 1。

\subsection{第二步：再对 $y$ 积分}

现在 $y$ 可以动了，沿着 $y$ 方向再扫一遍。

\[
\int_0^1 \frac{y}{2} \, dy = \frac{1}{2} \int_0^1 y \, dy = \frac{1}{2} \cdot \left[\frac{y^2}{2}\right]_0^1 = \frac{1}{2} \cdot \frac{1}{2} = \frac{1}{4}
\]

从 $y = 0$ 收集到 $y = 1$（就是从底边 $(x, 0)$ 那一行开始，一行一行往上扫，一直扫到顶边 $(x, 1)$ 那一行，把所有行的结果加起来）。

这块正方形上的二重积分结果是 $\dfrac{1}{4}$！

\section{🎨 动手画一画：再扫 $y$ 方向}

\textbf{在刚才的图旁边画一个竖箭头，从下到上，旁边写“再沿 $y$ 方向扫”。}

意思是：第一遍扫完每一条横条之后，把所有横条的结果从下到上全部加起来。

先扫 $x$，再扫 $y$——两遍扫完，整块地就收集干净了。

\subsection{例子 2：更大的矩形，$x$ 从 0 到 2，$y$ 从 0 到 3}

\[
\int_0^3 \int_0^2 xy \, dx \, dy
\]

\section{🎨 动手画一画：矩形区域}

\textbf{第一步：} 画一个长方形（不是正方形），横着比竖着短一点。

\textbf{第二步：} 底边左端标 \textbf{0}，右端标 \textbf{2}，下面写 \textbf{$x$}。

\textbf{第三步：} 左边下端标 \textbf{0}，上端标 \textbf{3}，旁边写 \textbf{$y$}。

\textbf{第四步：} 里面涂上阴影，写上 \textbf{$xy$}。

\subsection{第一步：先对 $x$ 积分}

\[
\int_0^2 xy \, dx = y \cdot \left[\frac{x^2}{2}\right]_0^2 = y \cdot \left(\frac{4}{2} - 0\right) = 2y
\]

里面的积分号上面是 2，下面是 0：从 $x = 0$ 收集到 $x = 2$（就是固定某一行的 $y$ 值不动，从点 $(0, y)$ 水平走到点 $(2, y)$，算这一行上 $xy$ 的总量）。

\subsection{第二步：再对 $y$ 积分}

\[
\int_0^3 2y \, dy = 2 \cdot \left[\frac{y^2}{2}\right]_0^3 = \left[y^2\right]_0^3 = 9 - 0 = 9
\]

外面的积分号上面是 3，下面是 0：从 $y = 0$ 收集到 $y = 3$（就是从底边 $(x, 0)$ 那一行开始，一行一行往上扫到顶边 $(x, 3)$ 那一行，把所有行的结果加起来）。

这块矩形上的二重积分结果是 \textbf{9}！

\subsection{例子 3：矩形，$x$ 从 1 到 3，$y$ 从 2 到 5}

\[
\int_2^5 \int_1^3 xy \, dx \, dy
\]

\section{🎨 动手画一画：不从零开始的矩形}

这次的矩形不从原点开始了！

\textbf{第一步：} 画一条水平数轴，标上 0、1、2、3。再画一条竖直数轴，标上 0、1、2、3、4、5。

\textbf{第二步：} 找到 $x = 1$、$x = 3$ 在水平轴上的位置，以及 $y = 2$、$y = 5$ 在竖直轴上的位置。

\textbf{第三步：} 以这四个边界画一个长方形：左边在 $x = 1$，右边在 $x = 3$，底边在 $y = 2$，顶边在 $y = 5$。

\textbf{第四步：} 把这个长方形涂上阴影，里面写上 \textbf{$xy$}。

看看你的图——这个矩形“飘”在坐标系的中间，不靠着原点。里面的积分号下面是 1 上面是 3（$x$ 的范围），外面的积分号下面是 2 上面是 5（$y$ 的范围）。

\subsection{第一步：先对 $x$ 积分}

\[
\int_1^3 xy \, dx = y \cdot \left[\frac{x^2}{2}\right]_1^3 = y \cdot \left(\frac{9}{2} - \frac{1}{2}\right) = y \cdot 4 = 4y
\]

从 $x = 1$ 收集到 $x = 3$（就是固定某一行的 $y$，从点 $(1, y)$ 水平走到点 $(3, y)$，算这一行的总量）。注意：上面代进去得 $\frac{9}{2}$，下面代进去得 $\frac{1}{2}$，\textbf{上减下}得到 4。

\subsection{第二步：再对 $y$ 积分}

\[
\int_2^5 4y \, dy = 4 \cdot \left[\frac{y^2}{2}\right]_2^5 = 4 \cdot \left(\frac{25}{2} - \frac{4}{2}\right) = 4 \cdot \frac{21}{2} = 42
\]

从 $y = 2$ 收集到 $y = 5$（就是从 $(x, 2)$ 那一行开始，一行行往上扫到 $(x, 5)$ 那一行，把所有行的结果加起来——整块矩形就扫完了）。

这块矩形上的二重积分结果是 \textbf{42}！

\subsection{例子 4：三角形区域（进阶挑战）}

三角形的三个角是 $(0, 0)$、$(1, 0)$、$(0, 1)$。斜边是 $y = 1 - x$。

\section{🎨 动手画一画：三角形区域}

\textbf{第一步：} 画两条互相垂直的线——水平的是 $x$ 轴，竖直的是 $y$ 轴。在交叉点写上 \textbf{(0, 0)}。

\textbf{第二步：} 在 $x$ 轴上往右量 1 个单位，画一个点，标上 \textbf{(1, 0)}。

\textbf{第三步：} 在 $y$ 轴上往上量 1 个单位，画一个点，标上 \textbf{(0, 1)}。

\textbf{第四步：} 用直线把 $(1, 0)$ 和 $(0, 1)$ 连起来。这就是斜边。

\textbf{第五步：} 现在有三条边围成的三角形：$x$ 轴上的一段、$y$ 轴上的一段、还有那条斜线。把三角形内部涂上阴影。

\textbf{第六步：} 在斜边旁边写上 \textbf{$x + y = 1$}（或者 $y = 1 - x$）。

看看你的图——这就是要收集的区域。

注意：这个三角形不是矩形！$y$ 的范围随着 $x$ 变化——$x$ 越大，$y$ 的范围越小。

\textbf{第七步（额外练习）：} 在三角形内部画三条竖直的虚线，分别在 $x = 0.2$、$x = 0.5$、$x = 0.8$ 的位置。看看每条虚线有多长？

\begin{itemize}
\item $x = 0.2$ 时，虚线从 $y = 0$ 到 $y = 0.8$，长度 0.8。
\item $x = 0.5$ 时，虚线从 $y = 0$ 到 $y = 0.5$，长度 0.5。
\item $x = 0.8$ 时，虚线从 $y = 0$ 到 $y = 0.2$，长度 0.2。
\end{itemize}

$x$ 越大，竖线越短。这就是为什么 $y$ 的上限是 $1 - x$，而不是一个固定的数！

\[
\int_0^1 \int_0^{1-x} xy \, dy \, dx
\]

看看上下的数字：

\begin{itemize}
\item \textbf{里面的积分号：} 下面是 0，上面是 \textbf{$1 - x$}。这不是一个固定的数字！它随着 $x$ 变化。$x = 0$ 的时候上限是 1，$x = 0.5$ 的时候上限是 0.5，$x = 1$ 的时候上限是 0。
\item \textbf{外面的积分号：} 下面是 0，上面是 1。这是固定的。
\end{itemize}

里面那个积分号的上面写着 $1 - x$，说明 $y$ 的终点线不是平的，而是斜的——就是你画的那条斜边！

\subsection{第一步：先对 $y$ 积分（$x$ 当常数）}

\[
\int_0^{1-x} xy \, dy
\]

$x$ 是常数，提出来：

\[
x \int_0^{1-x} y \, dy = x \cdot \left[\frac{y^2}{2}\right]_0^{1-x} = x \cdot \frac{(1-x)^2}{2}
\]

从 $y = 0$ 收集到 $y = 1 - x$（就是从底边上的点 $(x, 0)$ 竖直往上走，走到斜边上的点 $(x, 1-x)$ 停下——$x$ 越大，这条竖线越短，因为离斜边越近了）。

上面代进去的是 \textbf{$1 - x$}（不是一个普通数字，而是一个包含 $x$ 的表达式——因为三角形的边界就是斜的）。

展开 $(1-x)^2 = 1 - 2x + x^2$：

\[
= \frac{x(1 - 2x + x^2)}{2} = \frac{x - 2x^2 + x^3}{2}
\]

第一遍扫完了。结果有点长，但别慌，一项一项来。

\subsection{第二步：再对 $x$ 积分}

\[
\int_0^1 \frac{x - 2x^2 + x^3}{2} \, dx = \frac{1}{2} \int_0^1 \left(x - 2x^2 + x^3\right) dx
\]

分成三块，逐个击破：

\[
\int_0^1 x \, dx = \left[\frac{x^2}{2}\right]_0^1 = \frac{1}{2} - 0 = \frac{1}{2}
\]

\[
\int_0^1 2x^2 \, dx = 2 \cdot \left[\frac{x^3}{3}\right]_0^1 = 2 \cdot \left(\frac{1}{3} - 0\right) = \frac{2}{3}
\]

\[
\int_0^1 x^3 \, dx = \left[\frac{x^4}{4}\right]_0^1 = \frac{1}{4} - 0 = \frac{1}{4}
\]

合在一起：

\[
\frac{1}{2} - \frac{2}{3} + \frac{1}{4}
\]

通分（分母取 12）：

\[
\frac{6}{12} - \frac{8}{12} + \frac{3}{12} = \frac{1}{12}
\]

别忘了外面还有个 $\frac{1}{2}$：

\[
\frac{1}{2} \times \frac{1}{12} = \frac{1}{24}
\]

这块三角形上的二重积分结果是 $\dfrac{1}{24}$！

\subsection{彩蛋：交换积分顺序}

换个方向扫，结果应该一样。来验证一下。

\section{🎨 动手画一画：换个方向扫}

回到你刚才画的三角形。

\textbf{之前：} 你画了竖直的虚线（先固定 $x$，沿 $y$ 方向扫）。

\textbf{现在：} 擦掉竖直虚线，改画几条\textbf{水平的虚线}。

\begin{itemize}
\item $y = 0.2$ 时，水平线从 $x = 0$ 到 $x = 0.8$。
\item $y = 0.5$ 时，水平线从 $x = 0$ 到 $x = 0.5$。
\item $y = 0.8$ 时，水平线从 $x = 0$ 到 $x = 0.2$。
\end{itemize}

看到了吗？这次是 $x$ 的范围随着 $y$ 变化。每条水平线的右端点在斜边上，也就是 $x = 1 - y$。

所以换一种写法：

\[
\int_0^1 \int_0^{1-y} xy \, dx \, dy
\]

里面的积分号：$x$ 从 0 到 \textbf{$1 - y$}（上面的数字随 $y$ 变化）。
外面的积分号：$y$ 从 0 到 1。

\subsection{第一步：先对 $x$ 积分}

\[
\int_0^{1-y} xy \, dx = y \cdot \left[\frac{x^2}{2}\right]_0^{1-y} = y \cdot \frac{(1-y)^2}{2} = \frac{y - 2y^2 + y^3}{2}
\]

\subsection{第二步：再对 $y$ 积分}

\[
\frac{1}{2}\int_0^1 \left(y - 2y^2 + y^3\right) dy = \frac{1}{2}\left[\frac{y^2}{2} - \frac{2y^3}{3} + \frac{y^4}{4}\right]_0^1
\]

\[
= \frac{1}{2}\left(\frac{1}{2} - \frac{2}{3} + \frac{1}{4}\right) = \frac{1}{2} \cdot \frac{1}{12} = \frac{1}{24}
\]

果然！不管先扫 $x$ 再扫 $y$，还是先扫 $y$ 再扫 $x$，答案都是 $\frac{1}{24}$。

而且你发现没有？两次的计算\textbf{长得几乎一模一样}，只是 $x$ 和 $y$ 换了位置。这是因为 $xy$ 对 $x$ 和 $y$ 完全对称——换个名字，公式不变。

\subsection{超级彩蛋：$xy$ 的独家秘密——矩形上可以拆开算}

这里有一件漂亮的事情。

\textbf{正方形 $[0, 1] \times [0, 1]$：}

\[
\int_0^1 \int_0^1 xy \, dx \, dy = \frac{1}{4}
\]

如果分开算：

\[
\int_0^1 x \, dx = \frac{1}{2}, \qquad \int_0^1 y \, dy = \frac{1}{2}
\]

\[
\frac{1}{2} \times \frac{1}{2} = \frac{1}{4} \quad \checkmark
\]

\textbf{矩形 $[0, 2] \times [0, 3]$：}

\[
\int_0^3 \int_0^2 xy \, dx \, dy = 9
\]

分开算：

\[
\int_0^2 x \, dx = 2, \qquad \int_0^3 y \, dy = \frac{9}{2}
\]

\[
2 \times \frac{9}{2} = 9 \quad \checkmark
\]

大发现：

\[
\int_c^d \int_a^b xy \, dx \, dy = \left(\int_a^b x \, dx\right) \times \left(\int_c^d y \, dy\right)
\]

\textbf{在矩形区域上，$xy$ 的二重积分可以拆成两个一重积分的乘积！}

但要注意：三角形那种区域就不能这样拆了，因为一个变量的范围依赖另一个变量——上面的数字里包含了另一个变量。

\section{🎨 动手画一画：为什么矩形能拆，三角形不能}

翻回你画的两张图，仔细比较：

\textbf{矩形那张图：} $x$ 的范围是从 1 到 3（固定的），$y$ 的范围是从 2 到 5（也是固定的）。不管 $y$ 在哪里，$x$ 的范围都一样。两个变量各管各的。

\textbf{三角形那张图：} $y$ 的范围是从 0 到 $1 - x$（随 $x$ 变化的）。$x$ 在左边的时候 $y$ 可以很长，$x$ 在右边的时候 $y$ 很短。两个变量的范围\textbf{互相纠缠}。

在图上标注一下：矩形旁边写 \textbf{“可以拆”}，三角形旁边写 \textbf{“不能拆”}。

\section{公式汇总}

\subsection{一元微积分（$y$ 是常数，固定不动）}

\begin{center}
\begin{tabular}{c|c|c}
操作 & 公式 & 答案 \\
\hline
求导 & $\dfrac{d}{dx}(xy)$ & $y$ \\
不定积分 & $\displaystyle\int xy \, dx$ & $\dfrac{x^2 y}{2} + C$ \\
定积分 & $\displaystyle\int_a^b xy \, dx$ & $\dfrac{y(b^2 - a^2)}{2}$
\end{tabular}
\end{center}

\subsection{二元微积分（$x$ 和 $y$ 都是变量）}

\begin{center}
\begin{tabular}{c|c|c}
操作 & 公式 & 答案 \\
\hline
对 $x$ 偏导 & $\dfrac{\partial}{\partial x}(xy)$ & $y$ \\
对 $y$ 偏导 & $\dfrac{\partial}{\partial y}(xy)$ & $x$ \\
正方形 $[0,1]^2$ & $\displaystyle\int_0^1 \int_0^1 xy \, dx \, dy$ & $\dfrac{1}{4}$ \\
矩形 $[0,2]\times[0,3]$ & $\displaystyle\int_0^3 \int_0^2 xy \, dx \, dy$ & $9$ \\
矩形 $[1,3]\times[2,5]$ & $\displaystyle\int_2^5 \int_1^3 xy \, dx \, dy$ & $42$ \\
三角形 & $\displaystyle\int_0^1 \int_0^{1-x} xy \, dy \, dx$ & $\dfrac{1}{24}$
\end{tabular}
\end{center}

\subsection{上下数字速查}

\begin{center}
\begin{tabular}{c|c}
符号 & 意思 \\
\hline
$\displaystyle\int$ （光秃秃的） & 不定积分，只问收集规则，不指定范围 \\
$\displaystyle\int_0^2$ & 定积分，从 0 出发，到 2 停下 \\
$\displaystyle\int_1^3$ & 定积分，从 1 出发，到 3 停下 \\
$\left[f(x)\right]_a^b$ & 把 $x = b$ 代入减去 $x = a$ 代入（上减下） \\
$\displaystyle\int_0^{1-x}$ & 上限随 $x$ 变化！说明边界是斜的
\end{tabular}
\end{center}

\section{你学到了什么}

读完这一篇，你知道了这些事情：

\textbf{第一}，$xy$ 就是两个数相乘，就是长方形的面积。

\textbf{第二}，积分号上下的数字是“从哪到哪”——下面是起点，上面是终点，上减下得到答案。

\textbf{第三}，中括号右边的上下数字也是一样的意思：上面的代进去减去下面的代进去。

\textbf{第四}，$xy$ 对 $x$ 求导等于 $y$。画长方形就能看到：往右多长一格，多出来的面积就是 $y$。

\textbf{第五}，$xy$ 的积分是 $\frac{x^2 y}{2} + C$。把常数 $y$ 提出来，只对 $x$ 做积分。

\textbf{第六}，偏导数特别漂亮：对 $x$ 偏导是 $y$，对 $y$ 偏导是 $x$。你的导数是我，我的导数是你。

\textbf{第七}，二重积分就是积分做两次，先沿一个方向扫，再沿另一个方向扫。在矩形上可以拆成两个一重积分的乘积。

\textbf{第八}，当上面的数字里包含另一个变量（比如 $1 - x$），说明积分区域的边界是斜的，画出来看一看就明白了。

\chapter{$xyz$：只是多了个 $z$！}

\section{开场：第三个变量加入}

之前我们知道了，$xy$ 可以理解成\textbf{两个数相乘}。比如卖柠檬水，$x$ 是杯数，$y$ 是单价，$xy$ 就是一天的收入。

现在再加一个变量 $z$。$z$ 负责记\textbf{天数}。如果连续卖 $z$ 天，那总收入就是——

\[
xyz = \text{杯数} \times \text{单价} \times \text{天数} = \text{总收入}
\]

比如一天赚 20 块钱，连续卖 3 天，那就是 $20 \times 3 = 60$ 块钱！

\subsection{$xyz$ 还有另一个身份：长方体的体积}

上一篇我们说过，$xy$ 就是\textbf{长方形的面积}——底边是 $x$，高是 $y$。

现在 $z$ 加入了，长方形变成了\textbf{长方体}（就是一个盒子、一块砖头的形状）。

\[
xyz = \text{长} \times \text{宽} \times \text{高} = \text{体积}
\]

\section{🎨 动手画一画：$xyz$ 就是长方体的体积}

\textbf{第一步：} 在纸上画一个长方形，这是盒子的正面。

\textbf{第二步：} 从长方形的右上角，往右上方画一条斜线。从左上角也画一条同样方向同样长度的斜线。把两条斜线的顶端连起来。

\textbf{第三步：} 从右上角的斜线顶端，往下画一条竖线，连到从右下角出发的另一条斜线的顶端。

现在你画出了一个立体的盒子！

\textbf{第四步：} 在盒子的底边下面写 \textbf{$x$（长）}。

\textbf{第五步：} 在盒子的左边旁边写 \textbf{$y$（高）}。

\textbf{第六步：} 在盒子顶面的斜边旁边写 \textbf{$z$（宽/深度）}。

\textbf{第七步：} 在盒子的肚子里写上大大的 \textbf{$xyz$}。

看看你的图——$x$ 管长度，$y$ 管高度，$z$ 管深度。三个数乘在一起，就是盒子的体积。

记住这张图，后面会反复用到它。

\section{一、$xyz$ 的求导}

\subsection{一元微积分：两个变量当常数}

在一元微积分里，只有 $x$ 能动。$y$ 和 $z$ 都必须固定不变，当常数。

\subsection{用数字看一看}

假设 $y = 5$，$z = 3$（都固定不动），那 $f(x) = 15x$。

$x$ 从 2 变到 4，变了 2。

$f(x)$ 从 $15 \times 2 = 30$ 变到 $15 \times 4 = 60$，变了 30。

$30 \div 2 = 15$。15 是什么？就是 $yz = 5 \times 3$！

再试一个。$y = 10$，$z = 7$，$f(x) = 70x$。

$x$ 从 1 变到 2，变了 1。

$f(x)$ 从 70 变到 140，变了 70。

$70 \div 1 = 70 = yz$。没错，导数就是 $yz$！

\[
\frac{d}{dx}(xyz) = yz
\]

\section{🎨 动手画一画：为什么导数是 $yz$}

回到你刚才画的长方体。

\textbf{第一步：} 在长方体的右边，紧挨着它，画一片\textbf{薄薄的板}。这片板的厚度是 \textbf{1}（$x$ 增加了 1），高度是 \textbf{$y$}，深度是 \textbf{$z$}。

\textbf{第二步：} 在这片薄板上写上 \textbf{$y \times z \times 1 = yz$}。

看看你的图——长方体往右多长了一格，多出来的那一片薄板的体积恰好就是 \textbf{$yz$}。

\textbf{第三步：} 在薄板旁边写上：“$x$ 增加 1 $\to$ 体积增加 $yz$”。

上一篇里，$xy$ 对 $x$ 的导数是 $y$（多出来一个\textbf{窄条}，面积是 $y$）。

这一篇里，$xyz$ 对 $x$ 的导数是 $yz$（多出来一片\textbf{薄板}，体积是 $yz$）。

从窄条变成薄板，只是多了一个维度，道理完全一样。

\subsection{和前两篇的对比}

\begin{center}
\begin{tabular}{c|c|c}
函数 & 对 $x$ 求导 & 多出来的是什么 \\
\hline
$x + y$ & 1 & 一个点，大小固定是 1 \\
$xy$ & $y$ & 一个窄条，面积是 $y$ \\
$xyz$ & $yz$ & 一片薄板，体积是 $yz$
\end{tabular}
\end{center}

看到规律了吗？每次对 $x$ 求导，就是把 $x$ 拿掉，留下其余的。$x + y$ 拿掉 $x$ 剩 1（因为 $x$ 前面的系数是 1），$xy$ 拿掉 $x$ 剩 $y$，$xyz$ 拿掉 $x$ 剩 $yz$。

\section{二、偏导数——三个方向各看一次}

\subsection{三个变量全都能动}

现在进入三元微积分，$x$、$y$、$z$ 三个变量全都可以变化了。

还是老办法：一次只盯一个，其他两个当常数。

\subsection{\texorpdfstring{对 $x$ 求偏导}{对 x 求偏导}}

让 $y$ 和 $z$ 固定不动，只看 $x$。

\[
\frac{\partial}{\partial x}(xyz) = yz
\]

和刚才一样——把 $x$ 拿掉，留下 $yz$。

\subsection{\texorpdfstring{对 $y$ 求偏导}{对 y 求偏导}}

让 $x$ 和 $z$ 固定不动，只看 $y$。

假设 $x = 4$，$z = 3$（都不动），那 $f = 12y$。

$y$ 从 2 变到 5（变了 3），$f$ 从 24 变到 60（变了 36）。

$36 \div 3 = 12 = xz = 4 \times 3$。导数就是 $xz$！

\[
\frac{\partial}{\partial y}(xyz) = xz
\]

\section{🎨 动手画一画：对 $y$ 求偏导}

回到你画的长方体。

\textbf{第一步：} 在长方体的上面，紧挨着它，画一片\textbf{薄薄的板}。这片板的高度（厚度）是 \textbf{1}（$y$ 增加了 1），长度是 \textbf{$x$}，深度是 \textbf{$z$}。

\textbf{第二步：} 在这片薄板上写 \textbf{$x \times z \times 1 = xz$}。

\textbf{第三步：} 旁边写：“$y$ 增加 1 $\to$ 体积增加 $xz$”。

\subsection{对 $z$ 求偏导}

让 $x$ 和 $y$ 固定不动，只看 $z$。

假设 $x = 4$，$y = 5$（都不动），那 $f = 20z$。

$z$ 从 1 变到 3（变了 2），$f$ 从 20 变到 60（变了 40）。

$40 \div 2 = 20 = xy$。导数是 $xy$！

\[
\frac{\partial}{\partial z}(xyz) = xy
\]

$z$ 的偏导数是 $xy$——就是另外两个变量的乘积！

\section{🎨 动手画一画：对 $z$ 求偏导}

\textbf{第一步：} 在长方体的后面（斜着的方向），紧挨着它，画一片\textbf{薄薄的板}。这片板的深度（厚度）是 \textbf{1}（$z$ 增加了 1），长度是 \textbf{$x$}，高度是 \textbf{$y$}。

\textbf{第二步：} 在这片薄板上写 \textbf{$x \times y \times 1 = xy$}。

\textbf{第三步：} 旁边写：“$z$ 增加 1 $\to$ 体积增加 $xy$”。

\subsection{重大发现：三个变量互相成就}

\begin{center}
\begin{tabular}{c|c|c}
操作 & 结果 & 含义 \\
\hline
对 $x$ 求偏导 & $yz$ & $x$ 每变 1，体积变 $yz$ \\
对 $y$ 求偏导 & $xz$ & $y$ 每变 1，体积变 $xz$ \\
对 $z$ 求偏导 & $xy$ & $z$ 每变 1，体积变 $xy$
\end{tabular}
\end{center}

上一篇 $xy$ 的规律是——\textbf{你的偏导数是我，我的偏导数是你}。

这一篇 $xyz$ 的规律更漂亮——\textbf{你的偏导数是另外两个变量的乘积}！

每个变量退场的时候，留下的影响力就是另外两个搭档的合力。

$x$ 走了，留下 $yz$。$y$ 走了，留下 $xz$。$z$ 走了，留下 $xy$。

三个变量，谁也少不了谁。

\subsection{例子}

设 $f(x, y, z) = xyz$，求 $f$ 在点 $(2, 3, 5)$ 处的三个偏导数。

\[
f_x(2, 3, 5) = yz = 3 \times 5 = 15
\]

含义：在这个点上，$x$ 每增加 1，$f$ 大约增加 15。

\[
f_y(2, 3, 5) = xz = 2 \times 5 = 10
\]

含义：在这个点上，$y$ 每增加 1，$f$ 大约增加 10。

\[
f_z(2, 3, 5) = xy = 2 \times 3 = 6
\]

含义：在这个点上，$z$ 每增加 1，$f$ 大约增加 6。

在 $(2, 3, 5)$ 这个点上，$x$ 的影响力最大（15），$z$ 的影响力最小（6）。因为 $x$ 的搭档们（$y = 3, z = 5$）乘积最大，而 $z$ 的搭档们（$x = 2, y = 3$）乘积最小。

换个点就不一样了！比如 $(10, 1, 1)$，那 $z$ 的偏导是 $xy = 10$，而 $x$ 的偏导才 $yz = 1$。谁强谁弱，取决于搭档的实力。公平得很。

\section{三、$xyz$ 的积分}

\subsection{不定积分}

$y$ 和 $z$ 都不动？那把它们俩一起提出来！

\[
\int xyz \, dx = yz \int x \, dx = yz \cdot \frac{x^2}{2} + C = \frac{x^2 yz}{2} + C
\]

验算：对 $\frac{x^2 yz}{2}$ 求导，$x^2/2$ 的导数是 $x$，乘以常数 $yz$，得到 $xyz$。没错！

和上一篇一模一样。$xy$ 的积分是 $\frac{x^2 y}{2}$，$xyz$ 的积分就是 $\frac{x^2 yz}{2}$。只是多乘了一个 $z$，因为 $z$ 也是常数，一起提出来就行了。

\subsection{定积分}

\subsubsection{例子：$y = 3$，$z = 2$，$x$ 从 0 到 4}

\[
\int_0^4 6x \, dx
\]

（因为 $yz = 3 \times 2 = 6$，所以 $xyz = 6x$。）

\textbf{第一步：} 不定积分：$\frac{6x^2}{2} = 3x^2$

\textbf{第二步：} 把 $x = 4$ 代进去：$3 \times 16 = 48$

\textbf{第三步：} 把 $x = 0$ 代进去：$0$

\textbf{第四步：} 上限减下限：$48 - 0 = 48$

从 $x = 0$ 收集到 $x = 4$（就是沿着数轴从点 \textbf{0} 走到点 \textbf{4}，算这段线下面的面积），一共收集了 \textbf{48}！

\section{🎨 动手画一画}

\textbf{第一步：} 画一条数轴，标上 0、1、2、3、4。

\textbf{第二步：} 画出 $f(x) = 6x$ 这条线。$x = 0$ 时高度 0，$x = 4$ 时高度 24。连成直线。

\textbf{第三步：} 把 $x = 0$ 到 $x = 4$ 之间、直线下方涂上阴影。这是一个三角形，底边 4，高 24。

\textbf{第四步：} 验算面积：$\frac{1}{2} \times 4 \times 24 = 48$。和公式算出来的一模一样！

\section{四、三重积分}

\subsection{三个积分号？别慌！}

以前一重积分，沿着一条\textbf{线}收集——一个积分号，一组上下限。

后来二重积分，在一块\textbf{面}上收集——两个积分号，两组上下限。

现在三重积分，要在一个\textbf{立体空间}里收集——三个积分号，三组上下限。

每多一个积分号，就多一个方向要扫。一个方向一个方向地扫，扫完整个空间。

三个积分号是什么意思？

\[
\int_0^1 \int_0^1 \int_0^1 xyz \, dz \, dy \, dx
\]

三个 $\int$ 排在一起，每个都有自己的上下限。

\textbf{最里面那个}（配着 $dz$）：先沿 $z$ 方向扫——竖直的一根\textbf{柱子}。

\textbf{中间那个}（配着 $dy$）：再沿 $y$ 方向扫——把一排柱子合成一面\textbf{墙}。

\textbf{最外面那个}（配着 $dx$）：最后沿 $x$ 方向扫——把一面面墙合成一整栋\textbf{楼}。

\textbf{柱子 $\rightarrow$ 墙 $\rightarrow$ 楼。三步扫完整个空间。}

\section{🎨 动手画一画：三个方向怎么扫}

回到你画的长方体。

\textbf{第一步：} 在长方体内部，画一根\textbf{竖直的线}（从底面到顶面）。旁边写“\textbf{第一步：扫一根柱子（沿 $z$ 方向）}”。

\textbf{第二步：} 在长方体内部，画几根并排的竖线，组成一个\textbf{平面}（像一面墙）。旁边写“\textbf{第二步：把一排柱子合成一面墙（沿 $y$ 方向）}”。

\textbf{第三步：} 在长方体内部，画几面并排的“墙”，填满整个盒子。旁边写“\textbf{第三步：把所有墙合成整栋楼（沿 $x$ 方向）}”。

看看你的图——三次扫描，从一根柱子到一面墙到一整个盒子，整个空间就收集干净了。

\subsection{例子 1：单位正方体 $[0, 1] \times [0, 1] \times [0, 1]$}

这是一个边长为 1 的正方体：$x$ 从 0 到 1，$y$ 从 0 到 1，$z$ 从 0 到 1。

\[
\int_0^1 \int_0^1 \int_0^1 xyz \, dz \, dy \, dx
\]

\section{🎨 动手画一画：正方体}

\textbf{第一步：} 画一个立体的正方体（和你刚才画长方体的方法一样，只不过长宽高都一样长）。

\textbf{第二步：} 三条边分别标上：底边写 \textbf{$x$，从 0 到 1}；左边写 \textbf{$y$，从 0 到 1}；竖边写 \textbf{$z$，从 0 到 1}。

\textbf{第三步：} 在正方体肚子里写上 \textbf{$xyz$}。

\subsection{第一步：先对 $z$ 积分（$x$ 和 $y$ 都固定不动）}

沿 $z$ 方向扫一根柱子。

\[
\int_0^1 xyz \, dz
\]

$xy$ 是常数，提出来：

\[
xy \int_0^1 z \, dz = xy \cdot \left[\frac{z^2}{2}\right]_0^1 = xy \cdot \frac{1}{2} = \frac{xy}{2}
\]

从 $z = 0$ 收集到 $z = 1$（就是从地板上的点 $(x, y, 0)$ 竖直往上走到天花板上的点 $(x, y, 1)$，算这根柱子上 $xyz$ 的总量）。

第一遍扫完，每根柱子的价值是 $\frac{xy}{2}$。

\subsection{第二步：再对 $y$ 积分（$x$ 固定不动）}

把一排柱子合成一面墙。

\[
\int_0^1 \frac{xy}{2} \, dy = \frac{x}{2} \int_0^1 y \, dy = \frac{x}{2} \cdot \frac{1}{2} = \frac{x}{4}
\]

从 $y = 0$ 收集到 $y = 1$（就是从后墙上的那根柱子开始，一根根往前扫，把这一排柱子的总量加起来，合成一面墙的价值）。

第二遍扫完，每面墙的价值是 $\frac{x}{4}$。

\subsection{第三步：最后对 $x$ 积分}

把所有墙合在一起！

\[
\int_0^1 \frac{x}{4} \, dx = \frac{1}{4} \int_0^1 x \, dx = \frac{1}{4} \cdot \frac{1}{2} = \frac{1}{8}
\]

从 $x = 0$ 收集到 $x = 1$（就是从左墙那面墙开始，一面面墙往右推，整个正方体就扫完了）。

这个正方体的三重积分结果是 $\dfrac{1}{8}$！

\subsection{和之前的对比}

\begin{center}
\begin{tabular}{c|c|c}
函数 & 正方形 $[0,1]^2$ 上的二重积分 & 正方体 $[0,1]^3$ 上的三重积分 \\
\hline
$x + y$ & 1 & —（这篇不算这个） \\
$xy$ & 1/4 & —（$xy$ 不需要三重积分） \\
$xyz$ & — & \textbf{1/8}
\end{tabular}
\end{center}

$xy$ 在单位正方形上的积分是 $1/4 = (1/2)^2$。$xyz$ 在单位正方体上的积分是 $1/8 = (1/2)^3$。看到规律了吗？每多一个变量，就多乘一个 $1/2$！

\subsection{例子 2：大盒子 $[0, 2] \times [0, 3] \times [0, 4]$}

\[
\int_0^4 \int_0^3 \int_0^2 xyz \, dx \, dy \, dz
\]

\section{🎨 动手画一画：大盒子}

\textbf{第一步：} 画一个长方体，这次不是正方体了——长宽高不一样。

\textbf{第二步：} 底边标 \textbf{$x$，从 0 到 2}。左边标 \textbf{$y$，从 0 到 3}。竖边标 \textbf{$z$，从 0 到 4}（最高）。

\textbf{第三步：} 肚子里写 \textbf{$xyz$}。

\subsection{第一步：先对 $x$ 积分}

\[
\int_0^2 xyz \, dx = yz \cdot \left[\frac{x^2}{2}\right]_0^2 = yz \cdot 2 = 2yz
\]

从 $x = 0$ 收集到 $x = 2$（就是固定 $y$ 和 $z$，从点 $(0, y, z)$ 水平往右走到点 $(2, y, z)$，算这根横棍上 $xyz$ 的总量）。

\subsection{第二步：再对 $y$ 积分}

\[
\int_0^3 2yz \, dy = 2z \cdot \left[\frac{y^2}{2}\right]_0^3 = 2z \cdot \frac{9}{2} = 9z
\]

从 $y = 0$ 收集到 $y = 3$（就是固定 $z$，把一排横棍全部加起来，合成一片水平的楼层）。

\subsection{第三步：最后对 $z$ 积分}

\[
\int_0^4 9z \, dz = 9 \cdot \left[\frac{z^2}{2}\right]_0^4 = 9 \cdot 8 = 72
\]

从 $z = 0$ 收集到 $z = 4$（就是从地面一层层往上叠，整个盒子就扫完了）。

这个大盒子的三重积分结果是 \textbf{72}！

\subsection{例子 3：不从零开始的盒子 $[1, 3] \times [2, 5] \times [0, 2]$}

\[
\int_0^2 \int_2^5 \int_1^3 xyz \, dx \, dy \, dz
\]

\section{🎨 动手画一画：飘在空中的盒子}

这个盒子的 $x$ 不从 0 开始（从 1 到 3），$y$ 也不从 0 开始（从 2 到 5），只有 $z$ 从 0 开始（从 0 到 2）。

\textbf{第一步：} 画三条坐标轴：$x$ 往右，$y$ 往左上斜，$z$ 往上。

\textbf{第二步：} 在 $x$ 轴上标出 1 和 3。在 $y$ 轴上标出 2 和 5。在 $z$ 轴上标出 0 和 2。

\textbf{第三步：} 以这六个边界画一个盒子。这个盒子不靠着原点——它的左边在 $x = 1$，右边在 $x = 3$，前面在 $y = 2$，后面在 $y = 5$，底面在 $z = 0$，顶面在 $z = 2$。

\textbf{第四步：} 在盒子肚子里写 \textbf{$xyz$}。

\subsection{第一步：先对 $x$ 积分}

\[
\int_1^3 xyz \, dx = yz \cdot \left[\frac{x^2}{2}\right]_1^3 = yz \cdot \left(\frac{9}{2} - \frac{1}{2}\right) = yz \cdot 4 = 4yz
\]

从 $x = 1$ 收集到 $x = 3$（就是从点 $(1, y, z)$ 水平走到点 $(3, y, z)$，算这一段上 $xyz$ 的总量）。注意：上限代进去得 $9/2$，下限代进去得 $1/2$，\textbf{上减下}得到 4。

\subsection{第二步：再对 $y$ 积分}

\[
\int_2^5 4yz \, dy = 4z \cdot \left[\frac{y^2}{2}\right]_2^5 = 4z \cdot \left(\frac{25}{2} - \frac{4}{2}\right) = 4z \cdot \frac{21}{2} = 42z
\]

从 $y = 2$ 收集到 $y = 5$（就是一排排往前扫，把这一片楼层的结果加起来）。

\subsection{第三步：最后对 $z$ 积分}

\[
\int_0^2 42z \, dz = 42 \cdot \left[\frac{z^2}{2}\right]_0^2 = 42 \cdot 2 = 84
\]

从 $z = 0$ 收集到 $z = 2$（就是从地面一层层往上叠，整个盒子就扫完了）。

这个盒子的三重积分结果是 \textbf{84}！

\subsection{超级彩蛋：盒子上可以拆开算}

这里有一个美妙的规律。

\textbf{单位正方体 $[0, 1]^3$：}

\[
\int_0^1 \int_0^1 \int_0^1 xyz \, dz \, dy \, dx = \frac{1}{8}
\]

分开算：

\[
\int_0^1 x \, dx = \frac{1}{2}, \quad \int_0^1 y \, dy = \frac{1}{2}, \quad \int_0^1 z \, dz = \frac{1}{2}
\]

\[
\frac{1}{2} \times \frac{1}{2} \times \frac{1}{2} = \frac{1}{8} \quad \checkmark
\]

\textbf{大盒子 $[0, 2] \times [0, 3] \times [0, 4]$：}

\[
\int_0^4 \int_0^3 \int_0^2 xyz \, dx \, dy \, dz = 72
\]

分开算：

\[
\int_0^2 x \, dx = 2, \quad \int_0^3 y \, dy = \frac{9}{2}, \quad \int_0^4 z \, dz = 8
\]

\[
2 \times \frac{9}{2} \times 8 = 72 \quad \checkmark
\]

大发现：

\[
\int\!\!\!\int\!\!\!\int_{\text{盒子}} xyz \, dx \, dy \, dz = \left(\int x \, dx\right) \times \left(\int y \, dy\right) \times \left(\int z \, dz\right)
\]

\textbf{在盒子（长方体）区域上，$xyz$ 的三重积分可以拆成三个一重积分的乘积！}

上一篇的规律是：在矩形上，$xy$ 的二重积分可以拆成两个一重积分的乘积。

这一篇完全一样，只是多拆了一个。

三个变量各自独立干活，最后把成果乘在一起。

\subsection{例子 4：四面体——终极挑战}

\subsubsection{四面体是什么}

（还记得上一篇画过的四面体吗？）就是那个有四个三角形面的立体！

这个四面体有四个顶点：

\begin{itemize}
\item \textbf{(0, 0, 0)}：原点，房间的角落
\item \textbf{(1, 0, 0)}：往右走 1 步
\item \textbf{(0, 1, 0)}：往前走 1 步
\item \textbf{(0, 0, 1)}：往上走 1 步
\end{itemize}

三个面分别贴着三面“墙”（$x = 0$ 的墙，$y = 0$ 的墙，$z = 0$ 的地板），第四个面是一个斜面，满足 \textbf{$x + y + z = 1$}。

四面体内部的每一个点都满足：

\[
x \geq 0, \quad y \geq 0, \quad z \geq 0, \quad x + y + z \leq 1
\]

\section{🎨 动手画一画：四面体}

如果你之前已经画过四面体，翻出来看看。如果没有，现在画一个：

\textbf{第一步：} 在纸的中间偏左下画一个点，标上 \textbf{(0, 0, 0)}。

\textbf{第二步：} 从这个点往右画一条线段，终点标上 \textbf{(1, 0, 0)}。

\textbf{第三步：} 从原点往左上方斜着画一条线段，终点标上 \textbf{(0, 1, 0)}。

\textbf{第四步：} 从原点往正上方画一条线段，终点标上 \textbf{(0, 0, 1)}。

\textbf{第五步：} 把 $(1,0,0)$、$(0,1,0)$、$(0,0,1)$ 三个点两两连起来。这个三角形就是斜面。

\textbf{第六步：} 在斜面旁边写上 \textbf{$x + y + z = 1$}。

\textbf{第七步：} 把四面体内部涂上淡淡的阴影。

现在看看你的图——这就是要积分的区域。

\subsubsection{积分的范围怎么写？}

这不是一个规规矩矩的盒子。三个方向的范围互相纠缠——$z$ 的范围取决于 $x$ 和 $y$，$y$ 的范围取决于 $x$。

\textbf{$x$ 的范围：} 从 0 到 1（固定的，最简单）。

\textbf{$y$ 的范围：} 对于固定的 $x$，$y$ 从 0 到 $1 - x$。

为什么是 $1 - x$？因为 $x + y$ 不能超过 1（否则就超出四面体了）。所以 $y$ 最多走到 $1 - x$。$x$ 越大，$y$ 的活动空间越小。

\textbf{$z$ 的范围：} 对于固定的 $x$ 和 $y$，$z$ 从 0 到 $1 - x - y$。

为什么是 $1 - x - y$？因为 $x + y + z$ 不能超过 1。$x$ 和 $y$ 已经用掉了一些“名额”，$z$ 只能用剩下的。

\section{🎨 动手画一画：$z$ 的范围随 $x$ 和 $y$ 变化}

在你画的四面体上：

\textbf{第一步：} 选一个底面上的点，比如 $(0.2, 0.3, 0)$。从这个点竖直往上画一条线段，画到斜面上停下。斜面在这个点上方的位置是 $z = 1 - 0.2 - 0.3 = 0.5$。所以这根线段长 0.5。

\textbf{第二步：} 再选一个点，比如 $(0.5, 0.3, 0)$。竖直往上画到 $z = 1 - 0.5 - 0.3 = 0.2$。这根线段只有 0.2，短多了！

\textbf{第三步：} 再选 $(0.8, 0.1, 0)$。竖直往上画到 $z = 1 - 0.8 - 0.1 = 0.1$。更短了！

\textbf{旁边写上：} “$x$ 和 $y$ 越大，$z$ 的柱子越短——因为离斜面越近了。”

\subsubsection{正式写出三重积分}

\[
\int_0^1 \int_0^{1-x} \int_0^{1-x-y} xyz \, dz \, dy \, dx
\]

三个积分号，从里到外：

\begin{itemize}
\item \textbf{最里面：} $z$ 从 0 到 $1 - x - y$
\item \textbf{中间：} $y$ 从 0 到 $1 - x$
\item \textbf{最外面：} $x$ 从 0 到 1
\end{itemize}

深呼吸。一步一步来。每一步都是学过的简单积分，只是步数多了一点。

\subsubsection{第一步：先对 $z$ 积分（$x$ 和 $y$ 都不动）}

\[
\int_0^{1-x-y} xyz \, dz
\]

$xy$ 是常数，提出来：

\[
xy \int_0^{1-x-y} z \, dz = xy \cdot \left[\frac{z^2}{2}\right]_0^{1-x-y} = xy \cdot \frac{(1-x-y)^2}{2}
\]

从 $z = 0$ 收集到 $z = 1 - x - y$（就是从地板上的点 $(x, y, 0)$ 竖直往上走到斜面上的点 $(x, y, 1-x-y)$ 停下——$x$ 和 $y$ 越大，这根柱子越短，因为离斜面越近了）。

\[
= \frac{xy(1-x-y)^2}{2}
\]

第一遍扫完。每根柱子的价值是 $\frac{xy(1-x-y)^2}{2}$。看着有点长，但别怕，下一步会消化它。

\subsubsection{第二步：再对 $y$ 积分（$x$ 不动）}

\[
\int_0^{1-x} \frac{xy(1-x-y)^2}{2} \, dy
\]

$x$ 是常数，提出来：

\[
\frac{x}{2} \int_0^{1-x} y(1-x-y)^2 \, dy
\]

从 $y = 0$ 收集到 $y = 1 - x$（就是从后墙那根柱子开始，一根根往前扫到斜边上停下——和上一篇三角形里的道理一样，$x$ 越大，能往前扫的距离越短）。

这一步的计算有点长，但每一步都是简单的。来，跟着走。

\textbf{展开 $(1-x-y)^2$}

为了方便，先记 $a = 1 - x$（因为 $x$ 在这一步是常数，所以 $a$ 也是常数）。

那么 $1 - x - y = a - y$。

$(a - y)^2 = a^2 - 2ay + y^2$

所以：

$y(a-y)^2 = y(a^2 - 2ay + y^2) = a^2 y - 2ay^2 + y^3$

\textbf{逐项积分}

\[
\int_0^{a} (a^2 y - 2ay^2 + y^3) \, dy
\]

第一项：

\[
\int_0^a a^2 y \, dy = a^2 \cdot \left[\frac{y^2}{2}\right]_0^a = a^2 \cdot \frac{a^2}{2} = \frac{a^4}{2}
\]

第二项：

\[
\int_0^a 2ay^2 \, dy = 2a \cdot \left[\frac{y^3}{3}\right]_0^a = 2a \cdot \frac{a^3}{3} = \frac{2a^4}{3}
\]

第三项：

\[
\int_0^a y^3 \, dy = \left[\frac{y^4}{4}\right]_0^a = \frac{a^4}{4}
\]

\textbf{合在一起}

\[
\frac{a^4}{2} - \frac{2a^4}{3} + \frac{a^4}{4}
\]

提出公因子 $a^4$：

\[
a^4 \left(\frac{1}{2} - \frac{2}{3} + \frac{1}{4}\right)
\]

通分（分母取 12）：

\[
\frac{6}{12} - \frac{8}{12} + \frac{3}{12} = \frac{1}{12}
\]

所以：

\[
\int_0^{a} y(a-y)^2 \, dy = \frac{a^4}{12}
\]

\textbf{把 $a = 1 - x$ 代回去}

\[
= \frac{(1-x)^4}{12}
\]

\textbf{别忘了外面还有 $\frac{x}{2}$}

第二步的最终结果：

\[
\frac{x}{2} \cdot \frac{(1-x)^4}{12} = \frac{x(1-x)^4}{24}
\]

好！最难的一步搞定了。现在只剩最后一步！

\subsubsection{第三步：最后对 $x$ 积分}

\[
\int_0^1 \frac{x(1-x)^4}{24} \, dx = \frac{1}{24} \int_0^1 x(1-x)^4 \, dx
\]

需要算 $\int_0^1 x(1-x)^4 \, dx$。展开 $(1-x)^4$ 再逐项积分就行。

\textbf{展开 $(1-x)^4$}

$(1-x)^4 = 1 - 4x + 6x^2 - 4x^3 + x^4$

\textbf{乘以 $x$}

$x(1-x)^4 = x - 4x^2 + 6x^3 - 4x^4 + x^5$

\textbf{逐项积分}

\[
\int_0^1 x \, dx = \left[\frac{x^2}{2}\right]_0^1 = \frac{1}{2}
\]

\[
\int_0^1 4x^2 \, dx = 4 \cdot \left[\frac{x^3}{3}\right]_0^1 = \frac{4}{3}
\]

\[
\int_0^1 6x^3 \, dx = 6 \cdot \left[\frac{x^4}{4}\right]_0^1 = \frac{6}{4} = \frac{3}{2}
\]

\[
\int_0^1 4x^4 \, dx = 4 \cdot \left[\frac{x^5}{5}\right]_0^1 = \frac{4}{5}
\]

\[
\int_0^1 x^5 \, dx = \left[\frac{x^6}{6}\right]_0^1 = \frac{1}{6}
\]

\textbf{合在一起（注意正负号）}

\[
\frac{1}{2} - \frac{4}{3} + \frac{3}{2} - \frac{4}{5} + \frac{1}{6}
\]

通分（分母取 30）：

\[
\frac{15}{30} - \frac{40}{30} + \frac{45}{30} - \frac{24}{30} + \frac{5}{30}
\]

\[
= \frac{15 - 40 + 45 - 24 + 5}{30} = \frac{1}{30}
\]

\textbf{别忘了外面还有 $\frac{1}{24}$}

\[
\frac{1}{24} \times \frac{1}{30} = \frac{1}{720}
\]

这个四面体上的三重积分结果是 $\dfrac{1}{720}$！

\subsubsection{这个数字有多特别？}

$1/720$？$720$ 是什么来头？

$720 = 1 \times 2 \times 3 \times 4 \times 5 \times 6 = 6!$（6 的阶乘）。

来看看之前的记录：

\begin{center}
\begin{tabular}{c|c|c}
函数 & 区域 & 积分值 \\
\hline
$x$ & 线段 $[0, 1]$ & $\frac{1}{2} = \frac{1}{2!}$ \\
$xy$ & 三角形 & $\frac{1}{24} = \frac{1}{4!}$ \\
$xyz$ & 四面体 & $\frac{1}{720} = \frac{1}{6!}$
\end{tabular}
\end{center}

$1/2!$，$1/4!$，$1/6!$……分母是 2、4、6，每次多 2！

一个变量，分母是 $2! = 2$。两个变量，分母是 $4! = 24$。三个变量，分母是 $6! = 720$。

如果有四个变量 $xyzw$，在更高维的“四面体”上积分，答案应该是 $1/8! = 1/40320$。

数学的规律，真是美得让人起鸡皮疙瘩。

\subsection{例子 5：四面体的两种扫法}

\section{🎨 动手画一画：两种不同的扫描顺序}

翻回你画的四面体。

\textbf{第一种扫法（$dz\, dy\, dx$）：} 我们刚才用的就是这种。

\begin{itemize}
\item 在四面体内部画几根\textbf{竖直的短线}（柱子），从底面到斜面。
\item 旁边写“先扫 $z$（竖直柱子）$\to$ 再扫 $y$（把柱子合成墙）$\to$ 最后扫 $x$（把墙合成整个四面体）”
\end{itemize}

\textbf{第二种扫法（$dy\, dz\, dx$）：} 换个顺序试试。

\begin{itemize}
\item 在四面体内部画几根\textbf{水平往前的短线}，从后墙到斜面。
\item 旁边写“先扫 $y$（水平横棍）$\to$ 再扫 $z$（把横棍合成墙）$\to$ 最后扫 $x$（把墙合成整个四面体）”
\end{itemize}

不管用哪种扫法，最终答案都是 $1/720$。因为四面体还是那个四面体，里面的 $xyz$ 还是那个 $xyz$。扫法不同，只是换了个打扫顺序而已。

\section{终极对比：$xyz$ vs $xy$ vs $x + y$}

\begin{center}
\begin{tabular}{c|c|c|c}
 & $x + y$ & $xy$ & $xyz$ \\
\hline
\textbf{比喻} & AA 制 & 两人合伙 & 三人合伙 \\
\textbf{几何意义} & — & 长方形面积 & 长方体体积 \\
\textbf{对 $x$ 求导/偏导} & 1 & $y$ & $yz$ \\
\textbf{对 $y$ 求导/偏导} & 1 & $x$ & $xz$ \\
\textbf{对 $z$ 偏导} & — & — & $xy$ \\
\textbf{导数的含义} & 谁变 1，答案变 1 & 谁变 1，答案变另一个 & 谁变 1，答案变另外两个的乘积 \\
\textbf{正方形/正方体上的积分} & 1 & 1/4 & 1/8 \\
\textbf{三角形/四面体上的积分} & 1/3 & 1/24 & 1/720
\end{tabular}
\end{center}

\section{公式汇总}

\subsection{导数和偏导数}

\[
\frac{d}{dx}(xyz) = yz \quad \text{（$y$, $z$ 是常数时）}
\]

\[
\frac{\partial}{\partial x}(xyz) = yz, \quad \frac{\partial}{\partial y}(xyz) = xz, \quad \frac{\partial}{\partial z}(xyz) = xy
\]

\subsection{积分}

\[
\int xyz \, dx = \frac{x^2 yz}{2} + C
\]

\[
\int_a^b xyz \, dx = \frac{yz(b^2 - a^2)}{2}
\]

\subsection{三重积分}

\[
\int_0^1 \int_0^1 \int_0^1 xyz \, dz \, dy \, dx = \frac{1}{8}
\]

\[
\int_0^4 \int_0^3 \int_0^2 xyz \, dx \, dy \, dz = 72
\]

\[
\int_0^1 \int_0^{1-x} \int_0^{1-x-y} xyz \, dz \, dy \, dx = \frac{1}{720}
\]

\subsection{盒子上的分离规则}

\[
\int\!\!\!\int\!\!\!\int_{\text{盒子}} xyz \, dx \, dy \, dz = \left(\int x \, dx\right) \times \left(\int y \, dy\right) \times \left(\int z \, dz\right)
\]

\section{你学到了什么}

读完这一篇，你知道了这些事情：

\textbf{第一}，$xyz$ 就是三个数相乘，就是长方体的体积。

\textbf{第二}，$xyz$ 对 $x$ 求导等于 $yz$。画长方体就能看到：往右多长一格，多出来的薄板体积就是 $yz$。

\textbf{第三}，$xyz$ 的三个偏导数特别漂亮：\textbf{你的偏导数是另外两个变量的乘积}。对 $x$ 偏导是 $yz$，对 $y$ 偏导是 $xz$，对 $z$ 偏导是 $xy$。

\textbf{第四}，$xyz$ 的积分就是把不动的变量提出来，该怎么积就怎么积。和 $xy$ 完全一样，只是多提了一个常数。

\textbf{第五}，三重积分就是积分做三次——柱子 $\rightarrow$ 墙 $\rightarrow$ 楼，三步扫完整个空间。每多一个积分号，就多一组上下限，多一个方向要扫。

\textbf{第六}，在盒子（长方体）区域上，$xyz$ 的三重积分可以拆成三个一重积分的乘积。但四面体那种边界互相纠缠的区域不能拆。

\textbf{第七}，$xyz$ 在四面体上的积分是 $1/720 = 1/6!$，这和 $xy$ 在三角形上的 $1/24 = 1/4!$、$x$ 在线段上的 $1/2 = 1/2!$ 形成了一个美丽的规律。

\textbf{第八}，从 $x + y$ 到 $xy$ 到 $xyz$，核心方法从来没变过：\textbf{一次只看一个变量，其余的当常数。} 两个变量是这样，三个变量还是这样，一百个变量也是这样。

\chapter{$x^2y$：当一个变量开始抢戏}

\section{开场：$x$ 的贡献翻倍了}

之前 $xy$ 的生意模型里，$x$ 是杯数，$y$ 是单价，$xy$ 就是总收入。

但现在情况变了。假设不光卖柠檬水，每卖出一杯柠檬水还能搭售 $x$ 块饼干。这样 $x$ 的贡献变成了“杯数 $\times$ 杯数”。

总收入变成了：

\[
x^2 y = \text{杯数} \times \text{杯数} \times \text{单价} = x \times x \times y
\]

$x$ 出现了两次，$y$ 只出现一次。$x$ 的戏份明显比 $y$ 重了。

\section{🎨 动手画一画：$x^2y$ 的几何意义}

$x^2y$ 不再是简单的长方体了。它是一个\textbf{底面是正方形的长方体}。

\textbf{第一步：} 画一个立体的盒子。

\textbf{第二步：} 底面是一个\textbf{正方形}——长和宽都是 $x$。在底面的两条边上都写 \textbf{$x$}。

\textbf{第三步：} 盒子的高度是 \textbf{$y$}。在竖边上写 \textbf{$y$}。

\textbf{第四步：} 在盒子肚子里写上 \textbf{$x^2y$}。

看看你的图——底面积是 $x \times x = x^2$，高度是 $y$，体积就是 $x^2y$。

\textbf{$x$ 出现了两次，$y$ 只出现一次。$x$ 确实在抢戏。}

\section{一、$x^2y$ 的求导}

\subsection{一元微积分：$y$ 当常数}

$y$ 固定不动，只有 $x$ 能动。

\subsection{用数字看一看}

假设 $y = 3$（固定不动），那 $f(x) = 3x^2$。

\textbf{实验 1：}

$x$ 从 2 变到 3，变了 1。

$f(x)$ 从 $3 \times 4 = 12$ 变到 $3 \times 9 = 27$，变了 15。

$15 \div 1 = 15$。

\textbf{实验 2：}

$x$ 从 4 变到 5，变了 1。

$f(x)$ 从 $3 \times 16 = 48$ 变到 $3 \times 25 = 75$，变了 27。

$27 \div 1 = 27$。

等等……两次的导数不一样？第一次是 15，第二次是 27？

没错！这就是 $x^2$ 的特点。$x$ 越大，变化越猛。

\subsection{仔细研究一下}

当 $x = 2$ 时，导数应该是多少？

$f(x) = 3x^2$，根据求导公式，$f'(x) = 3 \times 2x = 6x$。

当 $x = 2$ 时，$f'(2) = 6 \times 2 = 12$。

当 $x = 4$ 时，$f'(4) = 6 \times 4 = 24$。

之前用 $x$ 从 2 变到 3 来估计 $x = 2$ 处的导数，得到 15，比真实值 12 大一点。这是因为变化量太大了，估计不够精确。

但规律是清楚的：\textbf{导数不再是一个固定的数，而是一个关于 $x$ 的函数。}

\subsection{正式结论}

\[
\frac{d}{dx}(x^2 y) = 2xy
\]

以前 $xy$ 对 $x$ 求导，得到 $y$——一个常数。

现在 $x^2y$ 对 $x$ 求导，得到 $2xy$——\textbf{$x$ 还在里面！}

这意味着：\textbf{$x$ 的位置不同，变化速度就不同。} $x$ 越大，变化越快。

\section{🎨 动手画一画：为什么导数是 $2xy$}

回到你画的那个底面是正方形的盒子。

\textbf{第一步：} 原来的盒子，底面边长是 $x$，高度是 $y$。体积是 $x^2y$。

\textbf{第二步：} 现在 $x$ 要增加一点点（从 $x$ 变成 $x + 1$）。底面的边长要变长。

\textbf{第三步：} 画一个新的盒子，底面边长是 $x + 1$。

\textbf{第四步：} 新盒子比旧盒子多出来的部分是什么？

这有点复杂。底面从 $x \times x$ 变成 $(x+1) \times (x+1)$。

$(x+1)^2 = x^2 + 2x + 1$

多出来的底面积是 $2x + 1$。

当变化量很小的时候（不是 1，而是趋近于 0），那个 $+1$ 可以忽略，多出来的底面积约等于 \textbf{$2x$}。

乘以高度 $y$，多出来的体积约等于 \textbf{$2xy$}。

\textbf{第五步：} 在图上标出多出来的部分。它不是一片薄板，而是\textbf{两片薄板加上一个小角落}：

\begin{itemize}
\item 右边一片：长 $x$，宽 1，高 $y$ $\to$ 体积 $xy$
\item 前面一片：长 $x$，宽 1，高 $y$ $\to$ 体积 $xy$
\item 角落一小条：长 1，宽 1，高 $y$ $\to$ 体积 $y$
\end{itemize}

两片加起来是 $2xy$，角落那一小条是 $y$。当变化量趋近于 0 时，角落那条可以忽略。

\textbf{所以导数是 $2xy$。}

\subsection{和之前的对比}

\begin{center}
\begin{tabular}{c|c|c}
函数 & 对 $x$ 求导 & 导数的特点 \\
\hline
$x + y$ & 1 & 常数，在哪都一样 \\
$xy$ & $y$ & 常数（因为 $y$ 不动），在哪都一样 \\
$x^2y$ & \textbf{$2xy$} & \textbf{随 $x$ 变化}，$x$ 越大，导数越大
\end{tabular}
\end{center}

$x^2$ 带来了一个本质的变化——\textbf{导数不再是常数}。

以前的函数，不管 $x$ 在哪里，变化速度都一样。

现在这个函数，$x$ 在 1 的时候慢悠悠，$x$ 在 100 的时候飞速狂奔。

\subsection{例子}

\textbf{例子 1：} $y = 5$，求 $f(x) = 5x^2$ 的导数

\[
\frac{d}{dx}(5x^2) = 5 \times 2x = 10x
\]

验算：

当 $x = 3$ 时，导数 $= 10 \times 3 = 30$。

这意味着：在 $x = 3$ 这个点，$x$ 每增加一丁点，$f(x)$ 大约增加 30 倍那么多。

当 $x = 10$ 时，导数 $= 10 \times 10 = 100$。

$x$ 越大，变化越猛。

\textbf{例子 2：} $y = 1$，求 $f(x) = x^2$ 的导数

\[
\frac{d}{dx}(x^2) = 2x
\]

这就是最基础的 $x^2$ 的导数——你可能在别的地方已经见过了。

\section{二、$x^2y$ 的积分}

\subsection{不定积分}

$y$ 是常数，提出来。然后对 $x^2$ 积分。

\[
\int x^2 y \, dx = y \int x^2 \, dx
\]

\subsection{$x^2$ 的积分是什么}

\[
\int x^2 \, dx = \frac{x^3}{3} + C
\]

为什么？因为对 $\frac{x^3}{3}$ 求导，得到 $\frac{3x^2}{3} = x^2$。

所以：

\[
\int x^2 y \, dx = y \cdot \frac{x^3}{3} + C = \frac{x^3 y}{3} + C
\]

\subsection{验算}

把答案求导：

\[
\frac{d}{dx}\left(\frac{x^3 y}{3}\right) = \frac{y}{3} \cdot 3x^2 = x^2 y \quad \checkmark
\]

完美。求导把积分装好的东西又拆回去了。

\subsection{和之前的对比}

\begin{center}
\begin{tabular}{c|c}
函数 & 不定积分 \\
\hline
$x$ & $\dfrac{x^2}{2} + C$ \\
$xy$ & $\dfrac{x^2 y}{2} + C$ \\
$x^2y$ & $\dfrac{x^3 y}{3} + C$
\end{tabular}
\end{center}

这里有个规律：

$x$ 的积分是 $x^2/2$——指数加 1，除以新指数。

$x^2$ 的积分是 $x^3/3$——还是指数加 1，除以新指数。

这个规律永远成立：$\displaystyle\int x^n \, dx = \frac{x^{n+1}}{n+1} + C$。

\subsection{定积分}

\subsubsection{例子 1：$y = 6$，$x$ 从 0 到 2}

\[
\int_0^2 6x^2 \, dx
\]

\textbf{第一步：} 不定积分：$6 \cdot \frac{x^3}{3} = 2x^3$

\textbf{第二步：} 把上面的数字 $x = 2$ 代进去：$2 \times 8 = 16$

\textbf{第三步：} 把下面的数字 $x = 0$ 代进去：$2 \times 0 = 0$

\textbf{第四步：} 上面的减下面的：$16 - 0 = 16$

从 $x = 0$ 收集到 $x = 2$（就是沿着数轴从点 \textbf{0} 走到点 \textbf{2}，算曲线 $6x^2$ 下面的面积），一共收集了 \textbf{16}。

\section{🎨 动手画一画：$x^2$ 的曲线不是直线！}

\textbf{第一步：} 画一条数轴，标上 0、1、2。

\textbf{第二步：} 画出 $f(x) = 6x^2$ 这条曲线。

注意：\textbf{这不是直线！} 它是一条\textbf{抛物线}，越往右越陡。

\begin{itemize}
\item $x = 0$ 时，$f(x) = 0$
\item $x = 1$ 时，$f(x) = 6$
\item $x = 2$ 时，$f(x) = 24$
\end{itemize}

\textbf{第三步：} 把这三个点标出来，然后用一条光滑的曲线连起来（向上弯曲的曲线）。

\textbf{第四步：} 把曲线下面、数轴上面的区域涂上阴影。

看看你的图——这不是三角形，也不是梯形，是一个\textbf{弯曲的区域}。

这就是为什么不能用简单的“底乘高除以 2”来算面积——曲线下的面积必须用积分来算。

\subsubsection{例子 2：$y = 3$，$x$ 从 1 到 4}

\[
\int_1^4 3x^2 \, dx
\]

\textbf{第一步：} 不定积分：$3 \cdot \frac{x^3}{3} = x^3$

\textbf{第二步：} $x = 4$ 代进去：$64$

\textbf{第三步：} $x = 1$ 代进去：$1$

\textbf{第四步：} 上减下：$64 - 1 = 63$

从 $x = 1$ 收集到 $x = 4$（就是沿着数轴从点 \textbf{1} 走到点 \textbf{4}，算曲线 $3x^2$ 下面的面积），一共 \textbf{63}。

\section{三、偏导数}

\subsection{\texorpdfstring{对 $x$ 求偏导}{对 x 求偏导}}

$y$ 固定不动，只看 $x$。

这和一元微积分的情况一样：

\[
\frac{\partial}{\partial x}(x^2 y) = 2xy
\]

\subsection{\texorpdfstring{对 $y$ 求偏导}{对 y 求偏导}}

$x$ 固定不动，只看 $y$。

$x$ 不动的时候，$x^2$ 就是一个常数。

假设 $x = 3$，那 $x^2 = 9$，$f = 9y$。

$y$ 从 2 变到 5（变了 3），$f$ 从 18 变到 45（变了 27）。

$27 \div 3 = 9 = x^2$。导数就是 $x^2$！

\[
\frac{\partial}{\partial y}(x^2 y) = x^2
\]

\section{🎨 动手画一画：对 $y$ 求偏导}

回到你画的那个底面是正方形的盒子。

\textbf{第一步：} 在盒子的上面，紧挨着它，画一片薄薄的板。这片板的高度（厚度）是 1（$y$ 增加了 1），底面是 \textbf{$x \times x = x^2$}。

\textbf{第二步：} 在这片薄板上写 \textbf{$x^2 \times 1 = x^2$}。

\textbf{第三步：} 旁边写：“$y$ 增加 1 $\to$ 体积增加 $x^2$”。

所以对 $y$ 的偏导数是 $x^2$。

\subsection{不对称了！}

\begin{center}
\begin{tabular}{c|c}
操作 & 结果 \\
\hline
对 $x$ 求偏导 & \textbf{$2xy$} \\
对 $y$ 求偏导 & \textbf{$x^2$}
\end{tabular}
\end{center}

上一篇 $xy$ 的偏导数是 $y$ 和 $x$——\textbf{对称的}，你的导数是我，我的导数是你。

这一篇 $x^2y$ 的偏导数是 $2xy$ 和 $x^2$——\textbf{不对称了}！

这很好理解。$x$ 出现了两次，$y$ 只出现了一次。$x$ 的影响力比 $y$ 大多了。这就是平方的力量。

\subsection{例子}

设 $f(x, y) = x^2y$，求 $f$ 在点 $(3, 4)$ 处的两个偏导数。

\[
f_x(3, 4) = 2xy = 2 \times 3 \times 4 = 24
\]

含义：在 $(3, 4)$ 这个点，$x$ 每增加 1，$f$ 大约增加 24。

\[
f_y(3, 4) = x^2 = 9
\]

含义：在 $(3, 4)$ 这个点，$y$ 每增加 1，$f$ 大约增加 9。

在这个点上，$x$ 的影响力（24）是 $y$ 的影响力（9）的将近 3 倍。$x$ 确实在抢戏。

\section{四、二重积分}

\subsection{例子 1：矩形 $[0, 2] \times [0, 3]$}

\[
\int_0^3 \int_0^2 x^2 y \, dx \, dy
\]

\section{🎨 动手画一画：矩形区域}

\textbf{第一步：} 画一个长方形，底边从 0 到 2，左边从 0 到 3。

\textbf{第二步：} 底边下面写 \textbf{$x$}，左边旁边写 \textbf{$y$}。

\textbf{第三步：} 长方形里面写 \textbf{$x^2y$}。

\subsection{第一步：先对 $x$ 积分}

\[
\int_0^2 x^2 y \, dx = y \cdot \left[\frac{x^3}{3}\right]_0^2 = y \cdot \frac{8}{3} = \frac{8y}{3}
\]

从 $x = 0$ 收集到 $x = 2$（就是固定某一行的 $y$ 值不动，从点 $(0, y)$ 水平走到点 $(2, y)$，算这一行上 $x^2y$ 的总量）。

\subsection{第二步：再对 $y$ 积分}

\[
\int_0^3 \frac{8y}{3} \, dy = \frac{8}{3} \cdot \left[\frac{y^2}{2}\right]_0^3 = \frac{8}{3} \cdot \frac{9}{2} = \frac{72}{6} = 12
\]

从 $y = 0$ 收集到 $y = 3$（就是从底边 $(x, 0)$ 那一行开始，一行一行往上扫，一直扫到顶边 $(x, 3)$ 那一行，把所有行的结果加起来）。

这个矩形上的二重积分结果是 \textbf{12}。

\subsection{$x^2y$ 也能在矩形上拆分吗？}

来检验一下。

\[
\int_0^3 \int_0^2 x^2 y \, dx \, dy \stackrel{?}{=} \left(\int_0^2 x^2 \, dx\right) \times \left(\int_0^3 y \, dy\right)
\]

左边我们刚算过，是 12。

右边：

\[
\int_0^2 x^2 \, dx = \left[\frac{x^3}{3}\right]_0^2 = \frac{8}{3}
\]

\[
\int_0^3 y \, dy = \left[\frac{y^2}{2}\right]_0^3 = \frac{9}{2}
\]

\[
\frac{8}{3} \times \frac{9}{2} = \frac{72}{6} = 12 \quad \checkmark
\]

没错！$x^2y$ 在矩形上也可以拆分。因为 $x^2y = (x^2) \times (y)$，是一个只关于 $x$ 的函数乘以一个只关于 $y$ 的函数。

\subsection{例子 2：三角形区域}

三角形顶点：$(0, 0)$、$(1, 0)$、$(0, 1)$。

\[
\int_0^1 \int_0^{1-x} x^2 y \, dy \, dx
\]

\section{🎨 动手画一画：三角形}

\textbf{第一步：} 画三角形，顶点在 $(0,0)$、$(1,0)$、$(0,1)$。

\textbf{第二步：} 在三角形内部画三条竖线，在 $x = 0.2$、$x = 0.5$、$x = 0.8$ 的位置。

\textbf{第三步：} 观察每条竖线的长度：
\begin{itemize}
\item 当 $x = 0.2$ 时，从 $y = 0$ 到 $y = 0.8$
\item 当 $x = 0.5$ 时，从 $y = 0$ 到 $y = 0.5$
\item 当 $x = 0.8$ 时，从 $y = 0$ 到 $y = 0.2$
\end{itemize}

\textbf{第四步：} 在图旁边写：“$y$ 的上限是 $1 - x$，随 $x$ 变化。”

\subsection{第一步：先对 $y$ 积分}

\[
\int_0^{1-x} x^2 y \, dy
\]

$x^2$ 是常数，提出来：

\[
x^2 \int_0^{1-x} y \, dy = x^2 \cdot \left[\frac{y^2}{2}\right]_0^{1-x} = x^2 \cdot \frac{(1-x)^2}{2} = \frac{x^2(1-x)^2}{2}
\]

从 $y = 0$ 收集到 $y = 1 - x$（就是从底边上的点 $(x, 0)$ 竖直往上走到斜边上的点 $(x, 1-x)$ 停下）。

\subsection{第二步：再对 $x$ 积分}

\[
\int_0^1 \frac{x^2(1-x)^2}{2} \, dx = \frac{1}{2} \int_0^1 x^2(1-x)^2 \, dx
\]

从 $x = 0$ 收集到 $x = 1$（就是从左边 $(0, y)$ 那条竖线开始，一条条往右扫到右下角 $(1, 0)$ 那个点，把整个三角形扫完）。

\textbf{展开 $(1-x)^2$}

$(1-x)^2 = 1 - 2x + x^2$

\textbf{乘以 $x^2$}

$x^2(1-x)^2 = x^2 - 2x^3 + x^4$

\textbf{逐项积分}

\[
\int_0^1 x^2 \, dx = \left[\frac{x^3}{3}\right]_0^1 = \frac{1}{3}
\]

\[
\int_0^1 2x^3 \, dx = 2 \cdot \left[\frac{x^4}{4}\right]_0^1 = \frac{2}{4} = \frac{1}{2}
\]

\[
\int_0^1 x^4 \, dx = \left[\frac{x^5}{5}\right]_0^1 = \frac{1}{5}
\]

\textbf{合在一起}

\[
\frac{1}{3} - \frac{1}{2} + \frac{1}{5}
\]

通分（分母取 30）：

\[
\frac{10}{30} - \frac{15}{30} + \frac{6}{30} = \frac{10 - 15 + 6}{30} = \frac{1}{30}
\]

\textbf{别忘了外面还有 $\frac{1}{2}$}

\[
\frac{1}{2} \times \frac{1}{30} = \frac{1}{60}
\]

这个三角形上的二重积分结果是 $\dfrac{1}{60}$！

\subsection{$x^2y$ vs $xy$ vs $x+y$ 在三角形上的对比}

\begin{center}
\begin{tabular}{c|c}
函数 & 三角形上的积分 \\
\hline
$x + y$ & $\dfrac{1}{3}$ \\
$xy$ & $\dfrac{1}{24}$ \\
$x^2y$ & $\dfrac{1}{60}$
\end{tabular}
\end{center}

$x + y$ 最大，因为加法让值分布得比较均匀。

$xy$ 比较小，因为在三角形里 $x$ 和 $y$ 都小于 1，相乘会变更小。

$x^2y$ 更小，因为 $x^2$ 让小的 $x$ 变得更小。当 $x = 0.5$ 时，$x^2 = 0.25$，缩水了一半。

\section{五、三重积分（附加内容）}

如果把 $x^2y$ 扩展到三维，变成 $x^2yz$，在单位正方体上的积分是多少？

\[
\int_0^1 \int_0^1 \int_0^1 x^2 yz \, dx \, dy \, dz
\]

因为是正方体（可分离区域）且函数是可分离的（$x^2 \times y \times z$），可以拆分：

\[
= \left(\int_0^1 x^2 \, dx\right) \times \left(\int_0^1 y \, dy\right) \times \left(\int_0^1 z \, dz\right)
\]

\[
= \frac{1}{3} \times \frac{1}{2} \times \frac{1}{2} = \frac{1}{12}
\]

对比：

\begin{center}
\begin{tabular}{c|c}
函数 & 单位正方体上的积分 \\
\hline
$xyz$ & $\dfrac{1}{8} = \dfrac{1}{2} \times \dfrac{1}{2} \times \dfrac{1}{2}$ \\
$x^2yz$ & $\dfrac{1}{12} = \dfrac{1}{3} \times \dfrac{1}{2} \times \dfrac{1}{2}$
\end{tabular}
\end{center}

$x^2$ 的贡献是 $1/3$（而不是 $1/2$），所以最终结果变小了。

\section{公式汇总}

\subsection{导数和偏导数}

\[
\frac{d}{dx}(x^2 y) = 2xy \quad \text{（$y$ 是常数时）}
\]

\[
\frac{\partial}{\partial x}(x^2 y) = 2xy
\]

\[
\frac{\partial}{\partial y}(x^2 y) = x^2
\]

\subsection{积分}

\[
\int x^2 y \, dx = \frac{x^3 y}{3} + C
\]

\[
\int_a^b x^2 y \, dx = \frac{y(b^3 - a^3)}{3}
\]

\subsection{二重积分}

\[
\int_0^3 \int_0^2 x^2 y \, dx \, dy = 12
\]

\[
\int_0^1 \int_0^{1-x} x^2 y \, dy \, dx = \frac{1}{60}
\]

\section{$x^2y$ vs $xy$ vs $x+y$ 终极对比}

\begin{center}
\begin{tabular}{c|c|c|c}
 & $x + y$ & $xy$ & $x^2y$ \\
\hline
\textbf{关系} & AA 制 & 平等合伙 & $x$ 主导 \\
\textbf{对 $x$ 偏导} & 1 & $y$ & \textbf{$2xy$} \\
\textbf{对 $y$ 偏导} & 1 & $x$ & \textbf{$x^2$} \\
\textbf{对称性} & 对称 & 对称 & \textbf{不对称} \\
\textbf{积分（对 $x$）} & $x + C$ & $\frac{x^2}{2} + C$ & $\frac{x^3}{3} + C$ \\
\textbf{矩形 $[0,2]\times[0,3]$} & 15 & 9 & \textbf{12} \\
\textbf{三角形} & $\frac{1}{3}$ & $\frac{1}{24}$ & $\frac{1}{60}$
\end{tabular}
\end{center}

\section{你学到了什么}

\textbf{第一}，$x^2y$ 就是 $x \times x \times y$，几何上是一个底面为正方形（边长 $x$）、高度为 $y$ 的长方体。

\textbf{第二}，$x^2y$ 对 $x$ 求导等于 \textbf{$2xy$}——导数不再是常数，而是随 $x$ 变化。这是因为 $x^2$ 的导数是 $2x$，不是 1。

\textbf{第三}，$x^2y$ 对 $y$ 求偏导等于 \textbf{$x^2$}——把 $x^2$ 当成常数，$y$ 的系数就是 $x^2$。

\textbf{第四}，$x^2y$ 的偏导数是\textbf{不对称}的：对 $x$ 偏导是 $2xy$，对 $y$ 偏导是 $x^2$。这和 $xy$ 的对称性（$y$ 和 $x$）不同。

\textbf{第五}，$x^2y$ 的积分是 $\frac{x^3y}{3} + C$。规律是：\textbf{指数加 1，除以新指数}。

\textbf{第六}，在矩形上，$x^2y$ 仍然可以拆分，因为 $x^2y = (x^2) \times (y)$ 是两个单变量函数的乘积。

\textbf{第七}，在三角形上，$x^2y$ 不能拆分，必须老老实实一步步积分。

\textbf{第八}，当函数中一个变量的指数增加时，这个变量的“戏份”就增加了——导数更复杂，对结果的影响更大。

\chapter{$xy^2$ 与规律的总结：反过来，然后看清全貌}

\section{开场：$y$ 也可以平方}

上一篇，$x$ 把自己变成了 $x^2$，在 $x^2y$ 里大出风头。

但凭什么只有 $x$ 能平方？$y$ 也可以！

\[
xy^2 = x \times y \times y
\]

$y$ 出现了两次，$x$ 只出现一次。这次轮到 $y$ 抢戏了。

\section{一、$xy^2$ 的求导}

\subsection{一元微积分：对 $x$ 求导（$y$ 是常数）}

$y$ 固定不动。那 $y^2$ 就是一个常数。

假设 $y = 3$，那 $y^2 = 9$，$f(x) = 9x$。

这和 $f(x) = 9x$ 没有任何区别——就是一个常数乘以 $x$。

\[
\frac{d}{dx}(xy^2) = y^2
\]

$y^2$ 是常数，常数乘以 $x$ 的导数就是那个常数。

\subsection{一元微积分：对 $y$ 求导（$x$ 是常数）}

现在换一个视角。把 $y$ 当作变量，$x$ 当作常数。

$f(y) = xy^2$

假设 $x = 4$，那 $f(y) = 4y^2$。

\textbf{实验 1：}

$y$ 从 2 变到 3，变了 1。

$f(y)$ 从 $4 \times 4 = 16$ 变到 $4 \times 9 = 36$，变了 20。

\textbf{实验 2：}

$y$ 从 5 变到 6，变了 1。

$f(y)$ 从 $4 \times 25 = 100$ 变到 $4 \times 36 = 144$，变了 44。

两次的变化量不一样——20 和 44。这说明导数随 $y$ 变化。

根据公式，$y^2$ 的导数是 $2y$，所以 $4y^2$ 的导数是 $8y$。

当 $y = 2$ 时，导数 $= 8 \times 2 = 16$。

当 $y = 5$ 时，导数 $= 8 \times 5 = 40$。

\[
\frac{d}{dy}(xy^2) = 2xy
\]

\subsection{偏导数}

\[
\frac{\partial}{\partial x}(xy^2) = y^2
\]

\[
\frac{\partial}{\partial y}(xy^2) = 2xy
\]

\section{🎨 动手画一画：$xy^2$ 的几何意义}

$xy^2$ 也是一个长方体，但这次\textbf{高度是正方形的边长}。

\textbf{第一步：} 画一个立体的盒子。

\textbf{第二步：} 底面是一个长方形——长是 $x$，宽是 $y$。在底面的两条边上分别写 \textbf{$x$} 和 \textbf{$y$}。

\textbf{第三步：} 盒子的高度也是 \textbf{$y$}。在竖边上写 \textbf{$y$}。

\textbf{第四步：} 在盒子肚子里写上 \textbf{$xy^2$}。

看看你的图——底面积是 $x \times y = xy$，高度是 $y$，体积就是 $xy \times y = xy^2$。

\textbf{$y$ 出现了两次，$x$ 只出现一次。这次轮到 $y$ 抢戏了。}

\subsection{对比 $x^2y$ 和 $xy^2$}

\begin{center}
\begin{tabular}{c|c|c}
 & $x^2y$ & $xy^2$ \\
\hline
谁在抢戏 & $x$ & $y$ \\
几何形状 & 底面正方形（$x \times x$），高 $y$ & 底面长方形（$x \times y$），高 $y$ \\
对 $x$ 偏导 & $2xy$ & $y^2$ \\
对 $y$ 偏导 & $x^2$ & $2xy$
\end{tabular}
\end{center}

这里有一个美妙的事实：

\textbf{把 $x^2y$ 里所有的 $x$ 换成 $y$、所有的 $y$ 换成 $x$，就得到 $xy^2$。}

\textbf{把 $x^2y$ 的偏导数里所有的 $x$ 换成 $y$、所有的 $y$ 换成 $x$，就得到 $xy^2$ 的偏导数。}

这叫做\textbf{对称性}——两个公式互为镜像。

\section{二、$xy^2$ 的积分}

\subsection{对 $x$ 积分}

\[
\int xy^2 \, dx = y^2 \int x \, dx = y^2 \cdot \frac{x^2}{2} + C = \frac{x^2 y^2}{2} + C
\]

\subsection{对 $y$ 积分}

\[
\int xy^2 \, dy = x \int y^2 \, dy = x \cdot \frac{y^3}{3} + C = \frac{xy^3}{3} + C
\]

\subsection{对比 $x^2y$ 和 $xy^2$}

\begin{center}
\begin{tabular}{c|c|c}
 & $x^2y$ & $xy^2$ \\
\hline
对 $x$ 积分 & $\dfrac{x^3 y}{3} + C$ & $\dfrac{x^2 y^2}{2} + C$ \\
对 $y$ 积分 & $\dfrac{x^2 y^2}{2} + C$ & $\dfrac{xy^3}{3} + C$
\end{tabular}
\end{center}

看到了吗？$x^2y$ 对 $x$ 积分和 $xy^2$ 对 $y$ 积分，形式完全一样（只是字母互换）。

$x^2y$ 对 $y$ 积分和 $xy^2$ 对 $x$ 积分也一样。

这就是对称性的力量——学会一个，另一个自动就会了。

\section{三、二重积分}

\subsection{矩形 $[0, 2] \times [0, 3]$ 上的 $xy^2$}

\[
\int_0^3 \int_0^2 xy^2 \, dx \, dy
\]

\textbf{第一步：先对 $x$ 积分}

\[
\int_0^2 xy^2 \, dx = y^2 \cdot \left[\frac{x^2}{2}\right]_0^2 = y^2 \cdot 2 = 2y^2
\]

从 $x = 0$ 收集到 $x = 2$（就是固定某一行的 $y$ 值不动，从点 $(0, y)$ 水平走到点 $(2, y)$，算这一行上 $xy^2$ 的总量）。

\textbf{第二步：再对 $y$ 积分}

\[
\int_0^3 2y^2 \, dy = 2 \cdot \left[\frac{y^3}{3}\right]_0^3 = 2 \cdot 9 = 18
\]

从 $y = 0$ 收集到 $y = 3$（就是从底边 $(x, 0)$ 那一行开始，一行一行往上扫，一直扫到顶边 $(x, 3)$ 那一行，把所有行的结果加起来）。

这个矩形上的二重积分结果是 \textbf{18}。

\subsection{$x^2y$ 和 $xy^2$ 在同一矩形上}

\begin{center}
\begin{tabular}{c|c}
函数 & 矩形 $[0, 2] \times [0, 3]$ 上的积分 \\
\hline
$x^2y$ & 12 \\
$xy^2$ & 18
\end{tabular}
\end{center}

为什么 $xy^2$ 更大？

因为 $y$ 的范围是 0 到 3，比 $x$ 的范围 0 到 2 更大。$y^2$ 在 $y = 3$ 时等于 9，而 $x^2$ 在 $x = 2$ 时只等于 4。所以 $y^2$ 的贡献更大。

\subsection{验证可分离性}

\[
\int_0^3 \int_0^2 xy^2 \, dx \, dy \stackrel{?}{=} \left(\int_0^2 x \, dx\right) \times \left(\int_0^3 y^2 \, dy\right)
\]

\[
= 2 \times 9 = 18 \quad \checkmark
\]

没错，$xy^2$ 在矩形上也可以拆分。

\subsection{三角形上的 $xy^2$}

三角形顶点：$(0, 0)$、$(1, 0)$、$(0, 1)$。

\[
\int_0^1 \int_0^{1-x} xy^2 \, dy \, dx
\]

\section{🎨 动手画一画：三角形}

\textbf{第一步：} 画三角形，顶点在 $(0,0)$、$(1,0)$、$(0,1)$。

\textbf{第二步：} 在三角形内部画三条竖线，在 $x = 0.2$、$x = 0.5$、$x = 0.8$ 的位置。

\textbf{第三步：} 观察每条竖线的长度：
当 $x = 0.2$ 时，从 $y = 0$ 到 $y = 0.8$；
当 $x = 0.5$ 时，从 $y = 0$ 到 $y = 0.5$；
当 $x = 0.8$ 时，从 $y = 0$ 到 $y = 0.2$。

\textbf{第四步：} 在图旁边写：“$y$ 的上限是 $1 - x$，随 $x$ 变化。”

\textbf{第一步：先对 $y$ 积分}

\[
\int_0^{1-x} xy^2 \, dy = x \cdot \left[\frac{y^3}{3}\right]_0^{1-x} = x \cdot \frac{(1-x)^3}{3} = \frac{x(1-x)^3}{3}
\]

从 $y = 0$ 收集到 $y = 1 - x$（就是从底边上的点 $(x, 0)$ 竖直往上走到斜边上的点 $(x, 1-x)$ 停下）。

\textbf{第二步：再对 $x$ 积分}

\[
\int_0^1 \frac{x(1-x)^3}{3} \, dx = \frac{1}{3} \int_0^1 x(1-x)^3 \, dx
\]

从 $x = 0$ 收集到 $x = 1$（就是从左边 $(0, y)$ 那条竖线开始，一条条往右扫到右下角 $(1, 0)$ 那个点，把整个三角形扫完）。

\textbf{展开 $(1-x)^3$}

\[
(1-x)^3 = 1 - 3x + 3x^2 - x^3
\]

\textbf{乘以 $x$}

\[
x(1-x)^3 = x - 3x^2 + 3x^3 - x^4
\]

\textbf{逐项积分}

\[
\int_0^1 x \, dx = \frac{1}{2}
\]

\[
\int_0^1 3x^2 \, dx = 3 \cdot \frac{1}{3} = 1
\]

\[
\int_0^1 3x^3 \, dx = 3 \cdot \frac{1}{4} = \frac{3}{4}
\]

\[
\int_0^1 x^4 \, dx = \frac{1}{5}
\]

\textbf{合在一起}

\[
\frac{1}{2} - 1 + \frac{3}{4} - \frac{1}{5}
\]

通分（分母取 20）：

\[
\frac{10}{20} - \frac{20}{20} + \frac{15}{20} - \frac{4}{20} = \frac{10 - 20 + 15 - 4}{20} = \frac{1}{20}
\]

\textbf{别忘了外面还有 $\frac{1}{3}$}

\[
\frac{1}{3} \times \frac{1}{20} = \frac{1}{60}
\]

这个三角形上的二重积分结果是 $\dfrac{1}{60}$！

\subsection{惊人的发现}

\begin{center}
\begin{tabular}{c|c}
函数 & 三角形上的积分 \\
\hline
$x^2y$ & $\dfrac{1}{60}$ \\
$xy^2$ & $\dfrac{1}{60}$
\end{tabular}
\end{center}

\textbf{一模一样！}

为什么？

这个三角形是关于直线 $y = x$ 对称的——把 $x$ 和 $y$ 互换，三角形还是那个三角形。

$x^2y$ 和 $xy^2$ 也是关于 $x$ 和 $y$ 对称的——互换 $x$ 和 $y$，一个变成另一个。

对称的函数在对称的区域上积分，结果当然一样！

\section{🎨 动手画一画：三角形的对称性}

\textbf{第一步：} 画出三角形，顶点在 $(0,0)$、$(1,0)$、$(0,1)$。

\textbf{第二步：} 画一条从 $(0,0)$ 到 $(0.5, 0.5)$ 的虚线。这是直线 $y = x$ 的一部分。

\textbf{第三步：} 沿着这条虚线折叠三角形（想象把纸对折）。

\textbf{第四步：} 你会发现：$(1, 0)$ 落到了 $(0, 1)$ 的位置，$(0, 1)$ 落到了 $(1, 0)$ 的位置。三角形和自己重合了。

\textbf{第五步：} 在图旁边写：“这个三角形关于 $y = x$ 对称。交换 $x$ 和 $y$，三角形不变。”

\section{四、规律的总结}

现在，退后一步，看看我们学过的所有函数，找出其中的规律。

\section{规律一：求导/偏导的规律}

\subsection{对 $x$ 求偏导：把 $x$ 的指数搬下来，指数减 1}

\begin{center}
\begin{tabular}{c|c|c}
函数 & 对 $x$ 偏导 & 规律 \\
\hline
$x$ & 1 & $x^1 \to 1 \cdot x^0 = 1$ \\
$x^2$ & $2x$ & $x^2 \to 2 \cdot x^1 = 2x$ \\
$x^3$ & $3x^2$ & $x^3 \to 3 \cdot x^2 = 3x^2$ \\
$xy$ & $y$ & $x^1 y \to 1 \cdot x^0 \cdot y = y$ \\
$x^2y$ & $2xy$ & $x^2 y \to 2 \cdot x^1 \cdot y = 2xy$ \\
$xy^2$ & $y^2$ & $x^1 y^2 \to 1 \cdot x^0 \cdot y^2 = y^2$
\end{tabular}
\end{center}

\textbf{一般规律：}

\[
\frac{\partial}{\partial x}(x^m y^n) = m \cdot x^{m-1} \cdot y^n
\]

口诀：\textbf{“指数搬下来当系数，指数减一，其他不动。”}

\subsection{对 $y$ 求偏导：同样的规律，只是换成 $y$}

\begin{center}
\begin{tabular}{c|c|c}
函数 & 对 $y$ 偏导 & 规律 \\
\hline
$y$ & 1 & $y^1 \to 1 \cdot y^0 = 1$ \\
$y^2$ & $2y$ & $y^2 \to 2 \cdot y^1 = 2y$ \\
$xy$ & $x$ & $x y^1 \to x \cdot 1 \cdot y^0 = x$ \\
$x^2y$ & $x^2$ & $x^2 y^1 \to x^2 \cdot 1 \cdot y^0 = x^2$ \\
$xy^2$ & $2xy$ & $x y^2 \to x \cdot 2 \cdot y^1 = 2xy$
\end{tabular}
\end{center}

\textbf{一般规律：}

\[
\frac{\partial}{\partial y}(x^m y^n) = n \cdot x^m \cdot y^{n-1}
\]

\section{🎨 动手画一画：做一张偏导数速查表}

拿一张纸，画一个表格：

\begin{center}
\begin{tabular}{c|c|c}
函数 & 对 $x$ 偏导 & 对 $y$ 偏导 \\
\hline
$x$ & 1 & 0 \\
$y$ & 0 & 1 \\
$xy$ & $y$ & $x$ \\
$x^2$ & $2x$ & 0 \\
$y^2$ & 0 & $2y$ \\
$x^2y$ & $2xy$ & $x^2$ \\
$xy^2$ & $y^2$ & $2xy$ \\
$x^2y^2$ & $2xy^2$ & $2x^2y$
\end{tabular}
\end{center}

把这张表贴在墙上，以后查起来很方便。

\section{规律二：积分的规律}

\subsection{对 $x$ 积分：指数加 1，除以新指数}

\begin{center}
\begin{tabular}{c|c|c}
函数 & 对 $x$ 积分 & 规律 \\
\hline
$x^0 = 1$ & $\frac{x^1}{1} = x$ & $x^0 \to \frac{x^1}{1} = x$ \\
$x^1 = x$ & $\frac{x^2}{2}$ & $x^1 \to \frac{x^2}{2}$ \\
$x^2$ & $\frac{x^3}{3}$ & $x^2 \to \frac{x^3}{3}$ \\
$x^3$ & $\frac{x^4}{4}$ & $x^3 \to \frac{x^4}{4}$ \\
$y$ & $xy$ & $y$ 是常数，$\int y \cdot 1 \, dx = yx$ \\
$xy$ & $\frac{x^2 y}{2}$ & $xy^1 \to \frac{x^2}{2} \cdot y$ \\
$x^2y$ & $\frac{x^3 y}{3}$ & $x^2 y \to \frac{x^3}{3} \cdot y$ \\
$xy^2$ & $\frac{x^2 y^2}{2}$ & $x y^2 \to \frac{x^2}{2} \cdot y^2$
\end{tabular}
\end{center}

\textbf{一般规律：}

\[
\int x^m y^n \, dx = \frac{x^{m+1}}{m+1} \cdot y^n + C
\]

口诀：\textbf{“指数加一放分母，其他不动跟着走。”}

\subsection{对 $y$ 积分：同样的规律}

\[
\int x^m y^n \, dy = x^m \cdot \frac{y^{n+1}}{n+1} + C
\]

\section{规律三：对称性}

\subsection{什么时候 $x$ 和 $y$ 地位平等？}

\begin{center}
\begin{tabular}{c|c|c|c}
函数 & 对 $x$ 偏导 & 对 $y$ 偏导 & 对称吗？ \\
\hline
$x + y$ & 1 & 1 & ✅ 完全对称 \\
$xy$ & $y$ & $x$ & ✅ 互为镜像 \\
$xyz$ & $yz$ & $xz$ & ✅ 互为镜像 \\
$x^2 + y^2$ & $2x$ & $2y$ & ✅ 互为镜像 \\
$x^2y^2$ & $2xy^2$ & $2x^2y$ & ✅ 互为镜像
\end{tabular}
\end{center}

\subsection{什么时候不对称？}

\begin{center}
\begin{tabular}{c|c|c|c}
函数 & 对 $x$ 偏导 & 对 $y$ 偏导 & 对称吗？ \\
\hline
$x^2y$ & $2xy$ & $x^2$ & ❌ 不对称 \\
$xy^2$ & $y^2$ & $2xy$ & ❌ 不对称 \\
$x^3y$ & $3x^2y$ & $x^3$ & ❌ 不对称
\end{tabular}
\end{center}

\textbf{规律：当 $x$ 和 $y$ 的指数相同时，函数关于 $x$ 和 $y$ 对称。当指数不同时，不对称。}

\subsection{对称函数的妙处}

如果函数关于 $x$ 和 $y$ 对称，而且积分区域也关于 $y = x$ 对称（比如正方形、等腰直角三角形），那么 $x$ 和 $y$ 的贡献相等。

这就是为什么 $x^2y$ 和 $xy^2$ 在三角形上的积分相等——它们互为镜像，三角形也关于 $y = x$ 对称。

\section{规律四：可分离性}

\subsection{什么时候可以拆分？}

\textbf{条件一：区域可分离}

积分区域必须是“盒子”形状——每个变量的范围是固定的常数，与其他变量无关。

✅ 矩形 $[0, 2] \times [0, 3]$

✅ 长方体 $[0, 2] \times [0, 3] \times [0, 4]$

❌ 三角形（$y$ 的范围取决于 $x$）

❌ 圆形（$y$ 的范围取决于 $x$）

❌ 四面体

\textbf{条件二：函数可分离}

被积函数必须能写成单变量函数的乘积。

✅ $xy = (x) \cdot (y)$

✅ $x^2 y = (x^2) \cdot (y)$

✅ $xy^2 = (x) \cdot (y^2)$

✅ $x^2 y^2 = (x^2) \cdot (y^2)$

✅ $xyz = (x) \cdot (y) \cdot (z)$

❌ $x + y$（加法不能拆成乘法）

❌ $x^2 + y^2$

❌ $xy + 1$

\textbf{两个条件都满足时：}

\[
\iint_{[a,b] \times [c,d]} f(x) g(y) \, dA = \left(\int_a^b f(x) \, dx\right) \times \left(\int_c^d g(y) \, dy\right)
\]

\section{规律五：积分结果的大小比较}

\subsection{在 $[0, 1]$ 区间内}

\begin{center}
\begin{tabular}{c|c}
函数 & 从 0 到 1 的积分 \\
\hline
1 & 1 \\
$x$ & $\frac{1}{2}$ \\
$x^2$ & $\frac{1}{3}$ \\
$x^3$ & $\frac{1}{4}$ \\
$x^4$ & $\frac{1}{5}$
\end{tabular}
\end{center}

\textbf{规律：指数越高，积分越小。}

为什么？因为在 0 到 1 之间，$x < 1$，所以 $x^2 < x$，$x^3 < x^2$，以此类推。

\subsection{在单位正方形 $[0, 1]^2$ 上}

\begin{center}
\begin{tabular}{c|c|c}
函数 & 积分值 & 等于 \\
\hline
1 & 1 & $1 \times 1$ \\
$xy$ & $\frac{1}{4}$ & $\frac{1}{2} \times \frac{1}{2}$ \\
$x^2y$ & $\frac{1}{6}$ & $\frac{1}{3} \times \frac{1}{2}$ \\
$xy^2$ & $\frac{1}{6}$ & $\frac{1}{2} \times \frac{1}{3}$ \\
$x^2y^2$ & $\frac{1}{9}$ & $\frac{1}{3} \times \frac{1}{3}$
\end{tabular}
\end{center}

\subsection{在单位正方体 $[0, 1]^3$ 上}

\begin{center}
\begin{tabular}{c|c|c}
函数 & 积分值 & 等于 \\
\hline
$xyz$ & $\frac{1}{8}$ & $\frac{1}{2} \times \frac{1}{2} \times \frac{1}{2}$ \\
$x^2yz$ & $\frac{1}{12}$ & $\frac{1}{3} \times \frac{1}{2} \times \frac{1}{2}$ \\
$x^2y^2z$ & $\frac{1}{18}$ & $\frac{1}{3} \times \frac{1}{3} \times \frac{1}{2}$ \\
$x^2y^2z^2$ & $\frac{1}{27}$ & $\frac{1}{3} \times \frac{1}{3} \times \frac{1}{3}$
\end{tabular}
\end{center}

\section{规律六：三角形/四面体上的特殊结果}

\begin{center}
\begin{tabular}{c|c|c}
函数 & 区域 & 积分值 \\
\hline
$x$ & 线段 $[0, 1]$ & $\frac{1}{2} = \frac{1}{2!}$ \\
$xy$ & 三角形 & $\frac{1}{24} = \frac{1}{4!}$ \\
$xyz$ & 四面体 & $\frac{1}{720} = \frac{1}{6!}$
\end{tabular}
\end{center}

\begin{center}
\begin{tabular}{c|c|c}
函数 & 区域 & 积分值 \\
\hline
$x^2$ & 线段 $[0, 1]$ & $\frac{1}{3}$ \\
$x^2y$ & 三角形 & $\frac{1}{60}$ \\
$xy^2$ & 三角形 & $\frac{1}{60}$ \\
$x^2y^2$ & 三角形 & ?（留作练习）
\end{tabular}
\end{center}

\section{五、一般公式 $x^m y^n$}

\subsection{定义}

\[
f(x, y) = x^m y^n
\]

其中 $m$ 和 $n$ 是非负整数（0, 1, 2, 3, \ldots）。

\subsection{偏导数}

\[
\frac{\partial}{\partial x}(x^m y^n) = m \cdot x^{m-1} \cdot y^n
\]

\[
\frac{\partial}{\partial y}(x^m y^n) = n \cdot x^m \cdot y^{n-1}
\]

\textbf{特殊情况：当 $m = 0$ 时，对 $x$ 的偏导数是 0（因为 $x$ 根本不在函数里）。}

\subsection{积分}

\[
\int x^m y^n \, dx = \frac{x^{m+1} y^n}{m+1} + C
\]

\[
\int x^m y^n \, dy = \frac{x^m y^{n+1}}{n+1} + C
\]

\subsection{在单位正方形上的积分}

\[
\int_0^1 \int_0^1 x^m y^n \, dx \, dy = \frac{1}{(m+1)(n+1)}
\]

\textbf{证明（用分离性）：}

\[
= \left(\int_0^1 x^m \, dx\right) \times \left(\int_0^1 y^n \, dy\right) = \frac{1}{m+1} \times \frac{1}{n+1}
\]

\subsection{例子验证}

\begin{center}
\begin{tabular}{c|c|c|c|c}
$m$ & $n$ & 函数 & 单位正方形上的积分 & 公式计算 \\
\hline
1 & 1 & $xy$ & $\frac{1}{4}$ & $\frac{1}{2 \times 2} = \frac{1}{4}$ ✓ \\
2 & 1 & $x^2y$ & $\frac{1}{6}$ & $\frac{1}{3 \times 2} = \frac{1}{6}$ ✓ \\
1 & 2 & $xy^2$ & $\frac{1}{6}$ & $\frac{1}{2 \times 3} = \frac{1}{6}$ ✓ \\
2 & 2 & $x^2y^2$ & $\frac{1}{9}$ & $\frac{1}{3 \times 3} = \frac{1}{9}$ ✓ \\
3 & 2 & $x^3y^2$ & $\frac{1}{12}$ & $\frac{1}{4 \times 3} = \frac{1}{12}$ ✓
\end{tabular}
\end{center}

\section{🎨 动手画一画：做一张大表格}

画一个二维表格，横轴是 $m$（$x$ 的指数），纵轴是 $n$（$y$ 的指数）。

\begin{center}
\begin{tabular}{c|c|c|c|c}
 & $m=0$ & $m=1$ & $m=2$ & $m=3$ \\
\hline
$n=0$ & 1 & $x$ & $x^2$ & $x^3$ \\
$n=1$ & $y$ & $xy$ & $x^2y$ & $x^3y$ \\
$n=2$ & $y^2$ & $xy^2$ & $x^2y^2$ & $x^3y^2$ \\
$n=3$ & $y^3$ & $xy^3$ & $x^2y^3$ & $x^3y^3$
\end{tabular}
\end{center}

在每个格子里，可以写上：
函数本身；
对 $x$ 的偏导数；
对 $y$ 的偏导数；
在单位正方形上的积分值。

这张表就是 $x^m y^n$ 家族的“族谱”。

\section{终极对照表}

\subsection{所有函数的偏导数}

\begin{center}
\begin{tabular}{c|c|c|c}
函数 & 对 $x$ 偏导 & 对 $y$ 偏导 & 对 $z$ 偏导 \\
\hline
$x + y$ & 1 & 1 & — \\
$xy$ & $y$ & $x$ & — \\
$xyz$ & $yz$ & $xz$ & $xy$ \\
$x^2y$ & $2xy$ & $x^2$ & — \\
$xy^2$ & $y^2$ & $2xy$ & — \\
$x^2y^2$ & $2xy^2$ & $2x^2y$ & — \\
$x^m y^n$ & $m x^{m-1} y^n$ & $n x^m y^{n-1}$ & —
\end{tabular}
\end{center}

\subsection{所有函数的积分}

\begin{center}
\begin{tabular}{c|c|c}
函数 & 对 $x$ 积分 & 对 $y$ 积分 \\
\hline
$x + y$ & $\frac{x^2}{2} + xy + C$ & $xy + \frac{y^2}{2} + C$ \\
$xy$ & $\frac{x^2 y}{2} + C$ & $\frac{xy^2}{2} + C$ \\
$x^2y$ & $\frac{x^3 y}{3} + C$ & $\frac{x^2 y^2}{2} + C$ \\
$xy^2$ & $\frac{x^2 y^2}{2} + C$ & $\frac{xy^3}{3} + C$ \\
$x^m y^n$ & $\frac{x^{m+1} y^n}{m+1} + C$ & $\frac{x^m y^{n+1}}{n+1} + C$
\end{tabular}
\end{center}

\subsection{在各种区域上的积分}

\begin{center}
\begin{tabular}{c|c|c|c}
函数 & 正方形 $[0,1]^2$ & 三角形 & 矩形 $[0,2]\times[0,3]$ \\
\hline
$xy$ & $\frac{1}{4}$ & $\frac{1}{24}$ & 9 \\
$x^2y$ & $\frac{1}{6}$ & $\frac{1}{60}$ & 12 \\
$xy^2$ & $\frac{1}{6}$ & $\frac{1}{60}$ & 18 \\
$x^2y^2$ & $\frac{1}{9}$ & ?（留作练习） & 24
\end{tabular}
\end{center}

\section{你学到了什么}

\textbf{第一}，$xy^2$ 和 $x^2y$ 是\textbf{镜像关系}——把所有 $x$ 换成 $y$、所有 $y$ 换成 $x$，一个变成另一个。

\textbf{第二}，对称的函数在对称的区域上积分，结果相等。$x^2y$ 和 $xy^2$ 在三角形上的积分都是 $1/60$。

\textbf{第三}，求偏导的规律：\textbf{指数搬下来当系数，指数减一，其他不动。}

\[
\frac{\partial}{\partial x}(x^m y^n) = m \cdot x^{m-1} \cdot y^n
\]

\textbf{第四}，积分的规律：\textbf{指数加一放分母，其他不动跟着走。}

\[
\int x^m y^n \, dx = \frac{x^{m+1} y^n}{m+1} + C
\]

\textbf{第五}，在单位正方形上，$x^m y^n$ 的积分等于 $\dfrac{1}{(m+1)(n+1)}$。

\textbf{第六}，可分离的函数在可分离的区域上，可以拆成多个一重积分的乘积。

\textbf{第七}，所有的规律都可以用一张表格来记忆。画出 $x^m y^n$ 的“族谱”，任何问题都能快速查到。

有时候 $x$ 抢戏（$x^2y$），有时候 $y$ 抢戏（$xy^2$），这才公平。而且它们互为镜像，学会一个就会另一个。

规律掌握了，以后不管遇到 $x^3y^5$ 还是 $x^{100}y^{200}$，都不怕了。

方法永远是一样的：\textbf{指数加一放分母，指数搬下来当系数。} 记住这两句话，走遍天下都不怕。

\section*{附录：快速参考卡片}

\begin{quote}
\ttfamily
偏导数速记\\
\\
\quad \(\partial/\partial x (x^m y^n) = m \cdot x^{m-1} \cdot y^n\)\\
\quad \(\partial/\partial y (x^m y^n) = n \cdot x^m \cdot y^{n-1}\)\\
\\
\quad 口诀：指数搬下来，指数减一，其他不动\\
\\
积分速记\\
\\
\quad \(\int x^m y^n\,dx = \dfrac{x^{m+1} y^n}{m+1} + C\)\\
\quad \(\int x^m y^n\,dy = \dfrac{x^m y^{n+1}}{n+1} + C\)\\
\\
\quad 口诀：指数加一放分母，其他不动跟着走\\
\\
单位正方形上的积分\\
\\
\quad \(\int\!\!\int x^m y^n\,dA = \dfrac{1}{(m+1)(n+1)}\)\\
\\
\quad 例：\(xy \to 1/4,\ x^2y \to 1/6,\ x^2y^2 \to 1/9\)
\end{quote}

\chapter{即便加了 $z$，又能奈何？}

\section{开场：$z$ 也想加入}

前面我们玩了很多花样——$xy$、$x^2y$、$xy^2$、$x^2y^2$……

那如果 $z$ 也加入进来，变成 $x^2y^2z^3$ 这种“怪物”，规律还管用吗？

答案是：\textbf{规律永远管用。} 来，我们试试。

\section{一般形式：$x^m y^n z^p$}

现在我们面对的是三个变量的一般形式：

\[
f(x, y, z) = x^m y^n z^p
\]

其中 $m$、$n$、$p$ 都是非负整数。

这个家族有多大？无穷大。随便举几个例子：

\begin{center}
\begin{tabular}{c|c|c|c}
$m$ & $n$ & $p$ & 函数 \\
\hline
1 & 1 & 1 & $xyz$ \\
2 & 1 & 1 & $x^2yz$ \\
1 & 2 & 1 & $xy^2z$ \\
1 & 1 & 2 & $xyz^2$ \\
2 & 2 & 2 & $x^2y^2z^2$ \\
3 & 2 & 1 & $x^3y^2z$ \\
1 & 1 & 3 & $xyz^3$ \\
0 & 2 & 3 & $y^2z^3$
\end{tabular}
\end{center}

看着挺多的，但别怕。规律只有两条，记住就全会了。

\section{规律一：偏导数}

\subsection{公式}

\[
\frac{\partial}{\partial x}(x^m y^n z^p) = m \cdot x^{m-1} \cdot y^n \cdot z^p
\]

\[
\frac{\partial}{\partial y}(x^m y^n z^p) = n \cdot x^m \cdot y^{n-1} \cdot z^p
\]

\[
\frac{\partial}{\partial z}(x^m y^n z^p) = p \cdot x^m \cdot y^n \cdot z^{p-1}
\]

\subsection{口诀}

\textbf{“对谁求导，谁的指数搬下来当系数，谁的指数减一，其他变量纹丝不动。”}

\subsection{例子验证}

\textbf{例子 1：} $f = x^2 y^3 z^4$

\[
\frac{\partial f}{\partial x} = 2 \cdot x^1 \cdot y^3 \cdot z^4 = 2xy^3z^4
\]

\[
\frac{\partial f}{\partial y} = 3 \cdot x^2 \cdot y^2 \cdot z^4 = 3x^2y^2z^4
\]

\[
\frac{\partial f}{\partial z} = 4 \cdot x^2 \cdot y^3 \cdot z^3 = 4x^2y^3z^3
\]

\textbf{例子 2：} $f = xyz$（最简单的情况，$m = n = p = 1$）

\[
\frac{\partial f}{\partial x} = 1 \cdot x^0 \cdot y \cdot z = yz
\]

\[
\frac{\partial f}{\partial y} = 1 \cdot x \cdot y^0 \cdot z = xz
\]

\[
\frac{\partial f}{\partial z} = 1 \cdot x \cdot y \cdot z^0 = xy
\]

这正是我们在 $xyz$ 那一篇学过的！

\textbf{例子 3：} $f = x^3 y z^2$

\[
\frac{\partial f}{\partial x} = 3x^2 y z^2
\]

\[
\frac{\partial f}{\partial y} = x^3 z^2
\]

\[
\frac{\partial f}{\partial z} = 2x^3 y z
\]

不管有几个变量，对谁求导就动谁，其他变量当常数。

\section{🎨 动手画一画：做一张三变量偏导数表}

画一个表格，每一行是一个函数，每一列是对不同变量的偏导数：

\begin{center}
\begin{tabular}{c|c|c|c}
函数 & 对 $x$ 偏导 & 对 $y$ 偏导 & 对 $z$ 偏导 \\
\hline
$xyz$ & $yz$ & $xz$ & $xy$ \\
$x^2yz$ & $2xyz$ & $x^2z$ & $x^2y$ \\
$xy^2z$ & $y^2z$ & $2xyz$ & $xy^2$ \\
$xyz^2$ & $yz^2$ & $xz^2$ & $2xyz$ \\
$x^2y^2z^2$ & $2xy^2z^2$ & $2x^2yz^2$ & $2x^2y^2z$ \\
$x^3y^2z$ & $3x^2y^2z$ & $2x^3yz$ & $x^3y^2$
\end{tabular}
\end{center}

\textbf{规律验证：} 每一行里，三个偏导数的“字母组合”加起来的总次数都比原函数少 1。

比如 $x^2y^2z^2$ 的总次数是 $2+2+2 = 6$，它的每个偏导数的总次数都是 5。

\section{规律二：积分}

\subsection{公式}

\[
\int x^m y^n z^p \, dx = \frac{x^{m+1} y^n z^p}{m+1} + C
\]

\[
\int x^m y^n z^p \, dy = \frac{x^m y^{n+1} z^p}{n+1} + C
\]

\[
\int x^m y^n z^p \, dz = \frac{x^m y^n z^{p+1}}{p+1} + C
\]

\subsection{口诀}

\textbf{“对谁积分，谁的指数加一放分母，其他变量纹丝不动。”}

\subsection{例子验证}

\textbf{例子 1：} $\int x^2 y^3 z^4 \, dx$

\[
= \frac{x^3 y^3 z^4}{3} + C
\]

验算：对结果求 $\partial/\partial x$，得到 $\frac{3x^2 y^3 z^4}{3} = x^2 y^3 z^4$。✓

\textbf{例子 2：} $\int xyz \, dz$

\[
= \frac{xyz^2}{2} + C
\]

验算：对结果求 $\partial/\partial z$，得到 $\frac{2xyz}{2} = xyz$。✓

\textbf{例子 3：} $\int x^3 y z^2 \, dy$

\[
= \frac{x^3 y^2 z^2}{2} + C
\]

\section{规律三：单位正方体上的积分}

\subsection{公式}

\[
\int_0^1 \int_0^1 \int_0^1 x^m y^n z^p \, dx \, dy \, dz = \frac{1}{(m+1)(n+1)(p+1)}
\]

\subsection{为什么？}

因为在单位正方体上，三个变量的范围都是独立的（0 到 1），函数 $x^m y^n z^p$ 可以分离，所以：

\[
= \left(\int_0^1 x^m \, dx\right) \times \left(\int_0^1 y^n \, dy\right) \times \left(\int_0^1 z^p \, dz\right)
\]

\[
= \frac{1}{m+1} \times \frac{1}{n+1} \times \frac{1}{p+1}
\]

\subsection{例子验证}

\begin{center}
\begin{tabular}{c|c|c|c|c|c}
函数 & $m$ & $n$ & $p$ & 积分值 & 验算 \\
\hline
$xyz$ & 1 & 1 & 1 & $\frac{1}{8}$ & $\frac{1}{2 \times 2 \times 2} = \frac{1}{8}$ ✓ \\
$x^2yz$ & 2 & 1 & 1 & $\frac{1}{12}$ & $\frac{1}{3 \times 2 \times 2} = \frac{1}{12}$ ✓ \\
$xy^2z$ & 1 & 2 & 1 & $\frac{1}{12}$ & $\frac{1}{2 \times 3 \times 2} = \frac{1}{12}$ ✓ \\
$xyz^2$ & 1 & 1 & 2 & $\frac{1}{12}$ & $\frac{1}{2 \times 2 \times 3} = \frac{1}{12}$ ✓ \\
$x^2y^2z^2$ & 2 & 2 & 2 & $\frac{1}{27}$ & $\frac{1}{3 \times 3 \times 3} = \frac{1}{27}$ ✓ \\
$x^3y^2z$ & 3 & 2 & 1 & $\frac{1}{24}$ & $\frac{1}{4 \times 3 \times 2} = \frac{1}{24}$ ✓
\end{tabular}
\end{center}

规律太整齐了！不管指数怎么变，公式永远是那个公式。

\section{规律四：对称性}

\subsection{什么时候三个变量地位平等？}

当 \textbf{$m = n = p$} 时，函数关于 $x$、$y$、$z$ 完全对称。

\begin{center}
\begin{tabular}{c|c|c|c|c}
函数 & 对 $x$ 偏导 & 对 $y$ 偏导 & 对 $z$ 偏导 & 对称吗？ \\
\hline
$xyz$ & $yz$ & $xz$ & $xy$ & ✅ 轮换对称 \\
$x^2y^2z^2$ & $2xy^2z^2$ & $2x^2yz^2$ & $2x^2y^2z$ & ✅ 轮换对称 \\
$x^3y^3z^3$ & $3x^2y^3z^3$ & $3x^3y^2z^3$ & $3x^3y^3z^2$ & ✅ 轮换对称
\end{tabular}
\end{center}

“轮换对称”的意思是：把 $x \to y \to z \to x$ 循环替换，偏导数也跟着循环替换。

\subsection{什么时候不对称？}

当指数不全相等时，某个变量的“戏份”就比别的重。

\begin{center}
\begin{tabular}{c|c|c}
函数 & 谁的戏份最重？ & 原因 \\
\hline
$x^2yz$ & $x$ & $x$ 的指数最高 \\
$xy^2z$ & $y$ & $y$ 的指数最高 \\
$xyz^2$ & $z$ & $z$ 的指数最高 \\
$x^3y^2z$ & $x$ & $x$ 的指数最高
\end{tabular}
\end{center}

\section{🎨 动手画一画：对称性的几何直觉}

\textbf{第一步：} 画一个立方体（正方体）。

\textbf{第二步：} 在三条从同一个顶点出发的边上分别标上 $x$、$y$、$z$。

\textbf{第三步：} 把立方体绕着“体对角线”（从一个角到对角的角）旋转 120°。

\textbf{第四步：} 观察——$x$ 轴转到了 $y$ 轴的位置，$y$ 轴转到了 $z$ 轴的位置，$z$ 轴转到了 $x$ 轴的位置。但立方体看起来还是一样的！

\textbf{第五步：} 在图旁边写上：\textbf{“当 $m = n = p$ 时，函数 $x^m y^n z^p$ 也有这种旋转对称性。”}

\section{规律五：推广到任意维度}

如果继续加变量呢？

\subsection{四个变量}

\[
f(x, y, z, w) = x^a y^b z^c w^d
\]

\[
\frac{\partial f}{\partial x} = a \cdot x^{a-1} y^b z^c w^d
\]

\[
\int x^a y^b z^c w^d \, dx = \frac{x^{a+1} y^b z^c w^d}{a+1} + C
\]

\[
\int_0^1 \int_0^1 \int_0^1 \int_0^1 x^a y^b z^c w^d \, dw \, dz \, dy \, dx = \frac{1}{(a+1)(b+1)(c+1)(d+1)}
\]

\subsection{$n$ 个变量}

\[
f(x_1, x_2, \ldots, x_n) = x_1^{m_1} x_2^{m_2} \cdots x_n^{m_n}
\]

\[
\frac{\partial f}{\partial x_k} = m_k \cdot x_k^{m_k - 1} \cdot \prod_{j \neq k} x_j^{m_j}
\]

\[
\int_0^1 \cdots \int_0^1 f \, dx_1 \cdots dx_n = \frac{1}{\prod_{k=1}^{n}(m_k + 1)}
\]

别被符号吓到。不管有多少个变量，规律永远是同一条：\textbf{对谁求导，谁的指数搬下来减一；对谁积分，谁的指数加一放分母。其他变量，纹丝不动。}

\section{终极速查表}

\subsection{偏导数}

\[
\frac{\partial}{\partial x_k}\left(\prod_{j=1}^{n} x_j^{m_j}\right) = m_k \cdot x_k^{m_k-1} \cdot \prod_{j \neq k} x_j^{m_j}
\]

\textbf{大白话：对谁求导，把谁的指数搬下来，指数减一，其他不动。}

\subsection{积分}

\[
\int \left(\prod_{j=1}^{n} x_j^{m_j}\right) dx_k = \frac{x_k^{m_k+1}}{m_k+1} \cdot \prod_{j \neq k} x_j^{m_j} + C
\]

\textbf{大白话：对谁积分，把谁的指数加一放分母，其他不动。}

\subsection{单位超立方体上的积分}

\[
\int_{[0,1]^n} \prod_{j=1}^{n} x_j^{m_j} \, d\mathbf{x} = \prod_{j=1}^{n} \frac{1}{m_j + 1}
\]

\textbf{大白话：每个变量贡献 $1/(\text{指数}+1)$，全部乘起来。}

\section{🎨 动手画一画：做一张终极速查卡}

\begin{quote}
\ttfamily
多变量幂函数速查\\
\\
函数：$x^m \cdot y^n \cdot z^p \cdot \ldots$\\
\\
偏导数（对 $x$）：\\
\\
\quad \(\partial/\partial x = m \cdot x^{m-1} \cdot y^n \cdot z^p \cdot \ldots\)\\
\\
\quad 口诀：指数搬下来，指数减一，其他不动\\
\\
积分（对 $x$）：\\
\\
\quad \(\int dx = x^{m+1}/(m+1) \cdot y^n \cdot z^p \cdot \ldots + C\)\\
\\
\quad 口诀：指数加一放分母，其他不动\\
\\
单位正方体/正方体上的积分：\\
\\
\quad \(\int\!\!\int\!\!\int x^m y^n z^p dV = 1/[(m+1)(n+1)(p+1)]\)\\
\\
\quad 例：$xyz \to 1/8,\ x^2yz \to 1/12,\ x^2y^2z^2 \to 1/27$
\end{quote}

\section{三变量实战练习}

\subsection{练习 1：求偏导数}

\[
f(x, y, z) = x^4 y^2 z^3
\]

答案：

\[
\frac{\partial f}{\partial x} = 4x^3 y^2 z^3
\]

\[
\frac{\partial f}{\partial y} = 2x^4 y z^3
\]

\[
\frac{\partial f}{\partial z} = 3x^4 y^2 z^2
\]

\subsection{练习 2：求不定积分}

\[
\int x^2 y^3 z \, dx = ?
\]

答案：

\[
= \frac{x^3 y^3 z}{3} + C
\]

\[
\int x^2 y^3 z \, dy = ?
\]

答案：

\[
= \frac{x^2 y^4 z}{4} + C
\]

\[
\int x^2 y^3 z \, dz = ?
\]

答案：

\[
= \frac{x^2 y^3 z^2}{2} + C
\]

\subsection{练习 3：单位正方体上的积分}

\[
\int_0^1 \int_0^1 \int_0^1 x^2 y^3 z \, dx \, dy \, dz = ?
\]

答案：

\[
= \frac{1}{(2+1)(3+1)(1+1)} = \frac{1}{3 \times 4 \times 2} = \frac{1}{24}
\]

验算（一步步积分）：

\[
\int_0^1 x^2 \, dx = \frac{1}{3}
\]

\[
\int_0^1 y^3 \, dy = \frac{1}{4}
\]

\[
\int_0^1 z \, dz = \frac{1}{2}
\]

\[
\frac{1}{3} \times \frac{1}{4} \times \frac{1}{2} = \frac{1}{24} \quad \checkmark
\]

\subsection{练习 4：长方体上的积分}

\[
\int_0^2 \int_0^3 \int_0^4 xyz^2 \, dx \, dy \, dz = ?
\]

因为是长方体，可以拆分：

\[
= \left(\int_0^2 x \, dx\right) \times \left(\int_0^3 y \, dy\right) \times \left(\int_0^4 z^2 \, dz\right)
\]

\[
= 2 \times \frac{9}{2} \times \frac{64}{3} = \frac{2 \times 9 \times 64}{2 \times 3} = \frac{1152}{6} = 192
\]

\section{你学到了什么}

\textbf{第一}，加了 $z$ 没什么可怕的。规律完全一样——对谁操作，就动谁的指数，其他变量当常数。

\textbf{第二}，偏导数：\textbf{指数搬下来当系数，指数减一。}

\[
\frac{\partial}{\partial x}(x^m y^n z^p) = m \cdot x^{m-1} \cdot y^n \cdot z^p
\]

\textbf{第三}，积分：\textbf{指数加一放分母。}

\[
\int x^m y^n z^p \, dx = \frac{x^{m+1} y^n z^p}{m+1} + C
\]

\textbf{第四}，单位正方体上的积分：\textbf{每个变量贡献 $1/(\text{指数}+1)$，乘起来。}

\[
\int_0^1 \int_0^1 \int_0^1 x^m y^n z^p \, dV = \frac{1}{(m+1)(n+1)(p+1)}
\]

\textbf{第五}，这些规律可以推广到任意多个变量。四个变量、五个变量、一百个变量——规律永远不变。

\textbf{第六}，当所有指数相等时，函数关于所有变量对称；当指数不等时，指数高的变量“戏份”更重。

原来加入新的变量，只是让乘积多了一项，并没有改变任何规律。一个变量是这个规律，两个变量是这个规律，三个变量还是这个规律。就算来一百个变量，规律也不变。

\textbf{数学的美，就在于规律的普遍性。} 掌握了核心规律，变量再多也不怕。

记住两句口诀就够了——

\begin{quote}
\textbf{求导：指数搬下来，指数减一，其他不动。}

\textbf{积分：指数加一放分母，其他不动。}
\end{quote}

带着这两句口诀，可以走遍整个多变量微积分的世界。

\chapter{$x-y$：不就是最后开减吗？}

\section{开场：$y$ 反着来了}

之前 $x + y$ 是两个变量相加，各出各的力。

但现在，$y$ 决定\textbf{倒着来}。你加多少，我就减多少。

于是，$x + y$ 变成了 \textbf{$x - y$}。

这有什么大不了的？\textbf{减法就是加上一个负数}。$x - y = x + (-y)$。把 $y$ 换成 $-y$，其他全都不变。

不信？我们来验证。

\section{一、$x - y$ 的求导}

\subsection{一元微积分：对 $x$ 求导（$y$ 是常数）}

$y$ 固定不动。

假设 $y = 5$（固定不动），那 $f(x) = x - 5$。

\textbf{实验 1：}

$x$ 从 10 变到 13，变了 3。

$f(x)$ 从 $10 - 5 = 5$ 变到 $13 - 5 = 8$，变了 3。

$3 \div 3 = 1$。

\textbf{实验 2：}

$x$ 从 100 变到 101，变了 1。

$f(x)$ 从 95 变到 96，变了 1。

$1 \div 1 = 1$。

\[
\frac{d}{dx}(x - y) = 1
\]

和 $x + y$ 完全一样！导数是 1。

为什么？因为 $y$ 是\textbf{常数}。不管 $y$ 是 $+5$ 还是 $-5$，它都不动。$x$ 变 1，答案就变 1，跟 $y$ 的正负没关系。

\subsection{偏导数}

\subsubsection{对 $x$ 求偏导}

\[
\frac{\partial}{\partial x}(x - y) = 1
\]

和 $x + y$ 一样。

\subsubsection{对 $y$ 求偏导}

现在 $x$ 固定不动，看 $y$ 怎么影响结果。

假设 $x = 10$，那 $f = 10 - y$。

$y$ 从 3 变到 7（变了 4），$f$ 从 $10 - 3 = 7$ 变到 $10 - 7 = 3$（变了 $-4$）。

$-4 \div 4 = -1$。

\[
\frac{\partial}{\partial y}(x - y) = -1
\]

这就是区别！

$x + y$ 对 $y$ 的偏导是 \textbf{$+1$}——$y$ 增加 1，答案增加 1。

$x - y$ 对 $y$ 的偏导是 \textbf{$-1$}——$y$ 增加 1，答案\textbf{减少} 1。

$y$ 的影响是\textbf{反着来}的。

\subsection{🎨 动手画一画：理解 $-1$ 的含义}

\textbf{第一步：} 画一条数轴，标上 0 到 10。

\textbf{第二步：} 在数轴上标出 $x = 10$ 的位置。

\textbf{第三步：} 假设 $y = 3$，那 $f = 10 - 3 = 7$。在数轴上标出 7 的位置，写上“$f = 7$”。

\textbf{第四步：} 假设 $y$ 增加到 5，那 $f = 10 - 5 = 5$。在数轴上标出 5 的位置，写上“$f = 5$”。

\textbf{第五步：} 画一个箭头从 7 指向 5，旁边写“$y$ 增加 2 $\rightarrow$ $f$ 减少 2”。

看看你的图——$y$ 往上走，$f$ 往下走。它们的方向是\textbf{相反的}。

这就是偏导数 $-1$ 的含义：\textbf{$y$ 每增加 1，$f$ 就减少 1。}

\subsection{和 $x + y$ 的对比}

\begin{center}
\begin{tabular}{c|c|c}
 & $x + y$ & $x - y$ \\
\hline
对 $x$ 偏导 & $+1$ & $+1$ \\
对 $y$ 偏导 & \textbf{$+1$} & \textbf{$-1$} \\
$y$ 增加时 & $f$ 增加 & $f$ 减少
\end{tabular}
\end{center}

\textbf{唯一的区别就是 $y$ 前面的符号。}

$x + y$ 里 $y$ 的系数是 $+1$，所以对 $y$ 的偏导是 $+1$。

$x - y$ 里 $y$ 的系数是 $-1$，所以对 $y$ 的偏导是 $-1$。

就这么简单。

\section{二、$x - y$ 的积分}

\subsection{不定积分（对 $x$ 积分）}

\[
\int (x - y) \, dx
\]

$y$ 是常数，当作 $-y$ 处理。

\[
= \int x \, dx - \int y \, dx = \frac{x^2}{2} - yx + C
\]

验算：

\[
\frac{d}{dx}\left(\frac{x^2}{2} - yx + C\right) = x - y \quad \checkmark
\]

\subsection{和 $x + y$ 的对比}

\begin{center}
\begin{tabular}{c|c}
函数 & 对 $x$ 的不定积分 \\
\hline
$x + y$ & $\dfrac{x^2}{2} + yx + C$ \\
$x - y$ & $\dfrac{x^2}{2} - yx + C$
\end{tabular}
\end{center}

看到了吗？唯一的区别就是 \textbf{$yx$ 前面的符号}。加号变减号，仅此而已。

\subsection{定积分}

\subsubsection{例子 1：$y = 3$，$x$ 从 0 到 4}

\[
\int_0^4 (x - 3) \, dx
\]

\textbf{第一步：} 不定积分是 $\frac{x^2}{2} - 3x$

\textbf{第二步：} $x = 4$ 代进去：$\frac{16}{2} - 12 = 8 - 12 = -4$

\textbf{第三步：} $x = 0$ 代进去：$0 - 0 = 0$

\textbf{第四步：} 上面的减下面的：$-4 - 0 = -4$

从 $x = 0$ 收集到 $x = 4$（就是沿着数轴从点 \textbf{0} 走到点 \textbf{4}，算曲线 $x - 3$ 下面的“面积”），结果是 \textbf{$-4$}。

等等，面积怎么能是负的？

\subsection{🎨 动手画一画：负面积的含义}

\textbf{第一步：} 画一条数轴（$x$ 轴），标上 0、1、2、3、4、5。

\textbf{第二步：} 画一条竖直的轴（$y$ 轴），标上 $-3$、$-2$、$-1$、0、1、2。

\textbf{第三步：} 画出 $f(x) = x - 3$ 这条直线。

\begin{itemize}
\item $x = 0$ 时，$f = -3$（在 $x$ 轴下方）
\item $x = 3$ 时，$f = 0$（正好穿过 $x$ 轴）
\item $x = 4$ 时，$f = 1$（在 $x$ 轴上方）
\end{itemize}

\textbf{第四步：} 把 $x = 0$ 到 $x = 3$ 之间、直线和 $x$ 轴之间的区域涂上\textbf{红色}。这部分在 $x$ 轴\textbf{下方}。

\textbf{第五步：} 把 $x = 3$ 到 $x = 4$ 之间、直线和 $x$ 轴之间的区域涂上\textbf{蓝色}。这部分在 $x$ 轴\textbf{上方}。

\textbf{第六步：} 计算两部分的面积：

红色三角形：底边 3，高度 3（向下），面积 $= \frac{1}{2} \times 3 \times 3 = 4.5$，但因为在 $x$ 轴下方，算作 \textbf{$-4.5$}。

蓝色三角形：底边 1，高度 1（向上），面积 $= \frac{1}{2} \times 1 \times 1 = 0.5$，算作 \textbf{$+0.5$}。

总面积 $= -4.5 + 0.5 = -4$。

和我们用公式算出来的一模一样！

当函数值为负（在 $x$ 轴下方）时，积分也是负的。这叫做\textbf{有符号面积}。

\subsubsection{例子 2：$y = 2$，$x$ 从 0 到 6}

\[
\int_0^6 (x - 2) \, dx
\]

\textbf{第一步：} 不定积分是 $\frac{x^2}{2} - 2x$

\textbf{第二步：} $x = 6$ 代进去：$18 - 12 = 6$

\textbf{第三步：} $x = 0$ 代进去：$0$

\textbf{第四步：} 上减下：$6 - 0 = 6$

从 $x = 0$ 收集到 $x = 6$，结果是 \textbf{6}。

这次是正的，因为 $x$ 轴上方的部分比下方的部分大。

\subsection{🎨 动手画一画：例子 2 的图}

\textbf{第一步：} 画出 $f(x) = x - 2$ 的直线。

\begin{itemize}
\item $x = 0$ 时，$f = -2$
\item $x = 2$ 时，$f = 0$（穿过 $x$ 轴）
\item $x = 6$ 时，$f = 4$
\end{itemize}

\textbf{第二步：} 红色区域（$x$ 轴下方）：底边 2，高 2，面积 $= \frac{1}{2} \times 2 \times 2 = 2$，算作 \textbf{$-2$}。

\textbf{第三步：} 蓝色区域（$x$ 轴上方）：底边 4，高 4，面积 $= \frac{1}{2} \times 4 \times 4 = 8$，算作 \textbf{$+8$}。

\textbf{第四步：} 总面积 $= -2 + 8 = 6$。✓

\subsection{一般公式}

\[
\int_a^b (x - y) \, dx = \left[\frac{x^2}{2} - yx\right]_a^b = \left(\frac{b^2}{2} - yb\right) - \left(\frac{a^2}{2} - ya\right)
\]

和 $x + y$ 的公式对比：

\begin{center}
\begin{tabular}{c|c}
函数 & 定积分公式 \\
\hline
$x + y$ & $\left(\frac{b^2}{2} + yb\right) - \left(\frac{a^2}{2} + ya\right)$ \\
$x - y$ & $\left(\frac{b^2}{2} - yb\right) - \left(\frac{a^2}{2} - ya\right)$
\end{tabular}
\end{center}

还是只有符号不同。

\section{三、二重积分}

\subsection{矩形 $[0, 2] \times [0, 3]$ 上的 $x - y$}

\[
\int_0^3 \int_0^2 (x - y) \, dx \, dy
\]

\subsection{🎨 动手画一画：矩形区域}

\textbf{第一步：} 画一个长方形，底边从 0 到 2（$x$），左边从 0 到 3（$y$）。

\textbf{第二步：} 在长方形里标出几个点的 $x - y$ 的值：

\begin{itemize}
\item 在左下角 $(0, 0)$：$x - y = 0$
\item 在右下角 $(2, 0)$：$x - y = 2$（正）
\item 在左上角 $(0, 3)$：$x - y = -3$（负）
\item 在右上角 $(2, 3)$：$x - y = -1$（负）
\end{itemize}

画一条线 $y = x$（从 $(0,0)$ 到 $(2,2)$），这条线的左下方 $x - y > 0$，右上方 $x - y < 0$。

\textbf{第一步：先对 $x$ 积分}

\[
\int_0^2 (x - y) \, dx = \left[\frac{x^2}{2} - yx\right]_0^2 = \left(2 - 2y\right) - 0 = 2 - 2y
\]

从 $x = 0$ 收集到 $x = 2$（就是固定某一行的 $y$ 值不动，从点 $(0, y)$ 水平走到点 $(2, y)$，算这一行上 $x - y$ 的总量）。

\textbf{第二步：再对 $y$ 积分}

\[
\int_0^3 (2 - 2y) \, dy = \left[2y - y^2\right]_0^3 = (6 - 9) - 0 = -3
\]

从 $y = 0$ 收集到 $y = 3$（就是从底边 $(x, 0)$ 那一行开始，一行一行往上扫，一直扫到顶边 $(x, 3)$ 那一行，把所有行的结果加起来）。

这个矩形上的二重积分结果是 \textbf{$-3$}。

\subsection{和 $x + y$ 的对比}

\[
\int_0^3 \int_0^2 (x + y) \, dx \, dy = ?
\]

先对 $x$ 积分：$\left[\frac{x^2}{2} + yx\right]_0^2 = 2 + 2y$

再对 $y$ 积分：$\int_0^3 (2 + 2y) \, dy = \left[2y + y^2\right]_0^3 = 6 + 9 = 15$

\begin{center}
\begin{tabular}{c|c}
函数 & 矩形 $[0,2]\times[0,3]$ 上的积分 \\
\hline
$x + y$ & \textbf{$+15$} \\
$x - y$ & \textbf{$-3$}
\end{tabular}
\end{center}

差这么多？

$x + y$ 的时候，$x$ 和 $y$ 都在加分，所有地方都是正的贡献。

$x - y$ 的时候，$x$ 在加分但 $y$ 在减分。在 $y$ 大的地方，减分超过了加分，拖累了总成绩。

特别是这个矩形里 $y$ 最大到 3，而 $x$ 最大只到 2。$y$ 的减分效应占了上风，所以总量是负的。

\subsection{三角形上的 $x - y$}

三角形顶点：$(0, 0)$、$(1, 0)$、$(0, 1)$。

\[
\int_0^1 \int_0^{1-x} (x - y) \, dy \, dx
\]

\textbf{第一步：先对 $y$ 积分}

\[
\int_0^{1-x} (x - y) \, dy = \left[xy - \frac{y^2}{2}\right]_0^{1-x}
\]

把 $y = 1 - x$ 代进去：

\[
= x(1-x) - \frac{(1-x)^2}{2}
\]

\[
= x - x^2 - \frac{1 - 2x + x^2}{2}
\]

\[
= x - x^2 - \frac{1}{2} + x - \frac{x^2}{2}
\]

\[
= 2x - \frac{3x^2}{2} - \frac{1}{2}
\]

从 $y = 0$ 收集到 $y = 1 - x$（就是从底边上的点 $(x, 0)$ 竖直往上走到斜边上的点 $(x, 1-x)$ 停下）。

\textbf{第二步：再对 $x$ 积分}

\[
\int_0^1 \left(2x - \frac{3x^2}{2} - \frac{1}{2}\right) dx
\]

\[
= \left[x^2 - \frac{x^3}{2} - \frac{x}{2}\right]_0^1
\]

\[
= 1 - \frac{1}{2} - \frac{1}{2} = 0
\]

从 $x = 0$ 收集到 $x = 1$（就是从左边 $(0, y)$ 那条竖线开始，一条条往右扫到右下角 $(1, 0)$ 那个点，把整个三角形扫完）。

结果是 \textbf{0}？

\subsection{为什么是 0？}

\subsection{🎨 动手画一画：三角形的对称性}

\textbf{第一步：} 画出三角形，顶点在 $(0,0)$、$(1,0)$、$(0,1)$。

\textbf{第二步：} 画出直线 $y = x$（从 $(0,0)$ 到 $(0.5, 0.5)$）。

\textbf{第三步：} 这条线把三角形分成两半：

\begin{itemize}
\item 下半部分（$y < x$ 的区域）：$x - y > 0$
\item 上半部分（$y > x$ 的区域）：$x - y < 0$
\end{itemize}

\textbf{第四步：} 观察这两半——它们关于直线 $y = x$ \textbf{完全对称}！

\textbf{第五步：} 在图旁边写：“正的部分和负的部分完全抵消，所以积分是 0。”

这个三角形关于 $y = x$ 对称。在下半部分 $x$ 比 $y$ 大的部分，到了上半部分正好反过来。最后打个平手！

\subsection{和 $x + y$ 在三角形上的对比}

\begin{center}
\begin{tabular}{c|c}
函数 & 三角形上的积分 \\
\hline
$x + y$ & $\dfrac{1}{3}$ \\
$x - y$ & \textbf{0}
\end{tabular}
\end{center}

$x + y$ 在这个三角形上积分是 $1/3$——全是正贡献，加起来。

$x - y$ 在这个三角形上积分是 0——正负抵消。

这是因为三角形关于 $y = x$ 对称，而 $x - y$ 在这条线两侧的值恰好相反。

\section{四、规律总结}

\subsection{加法和减法的统一处理}

\[
f(x, y) = ax + by
\]

其中 $a$ 和 $b$ 是常数。

\textbf{特殊情况：}

\begin{center}
\begin{tabular}{c|c|c}
$a$ & $b$ & 函数 \\
\hline
1 & 1 & $x + y$ \\
1 & $-1$ & $x - y$ \\
2 & 3 & $2x + 3y$ \\
1 & 0 & $x$ \\
0 & 1 & $y$
\end{tabular}
\end{center}

\textbf{偏导数：}

\[
\frac{\partial}{\partial x}(ax + by) = a
\]

\[
\frac{\partial}{\partial y}(ax + by) = b
\]

\textbf{积分：}

\[
\int (ax + by) \, dx = \frac{ax^2}{2} + bxy + C
\]

\[
\int (ax + by) \, dy = axy + \frac{by^2}{2} + C
\]

\subsection{例子验证}

\textbf{例子 1：} $f(x, y) = 3x - 2y$

\[
\frac{\partial f}{\partial x} = 3, \quad \frac{\partial f}{\partial y} = -2
\]

\[
\int (3x - 2y) \, dx = \frac{3x^2}{2} - 2xy + C
\]

\textbf{例子 2：} $f(x, y) = -x + 4y$

\[
\frac{\partial f}{\partial x} = -1, \quad \frac{\partial f}{\partial y} = 4
\]

\[
\int (-x + 4y) \, dy = -xy + 2y^2 + C
\]

\subsection{$x - y$ 不能拆分！}

（还记得吗？$xy$ 在矩形上可以拆成 $(\int x)(\int y)$。）

那 $x - y$ 呢？

\textbf{不能！}

因为 $x - y$ 是\textbf{加法}（或减法），不是乘法。

\[
\iint (x - y) \, dA \neq \left(\int x\right) - \left(\int y\right)
\]

正确的做法是把加减法\textbf{分配到积分里}：

\[
\iint (x - y) \, dA = \iint x \, dA - \iint y \, dA
\]

这是另一条规则——\textbf{积分的线性性质}。

\subsection{🎨 动手画一画：做一张加减法速查表}

\begin{quote}
\ttfamily
加减法函数 $ax + by$ 速查\\
\\
偏导数：\\
\\
\quad \(\partial/\partial x (ax + by) = a\)\\
\quad \(\partial/\partial y (ax + by) = b\)\\
\\
\quad 口诀：对谁求导，就是谁前面的系数\\
\\
积分：\\
\\
\quad \(\int(ax + by)\,dx = (a/2)x^2 + bxy + C\)\\
\quad \(\int(ax + by)\,dy = axy + (b/2)y^2 + C\)\\
\\
\quad 口诀：谁被积分，谁的指数加一放分母\\
\\
特殊情况：\\
\\
\quad $x + y:\ a = 1,\ b = 1$\\
\quad $x - y:\ a = 1,\ b = -1$\\
\quad $x:\ a = 1,\ b = 0$\\
\quad $y:\ a = 0,\ b = 1$
\end{quote}

\section{五、$x - y$ 与 $xy$ 的本质区别}

$x - y$ 和 $xy$ 到底有什么本质区别？

\begin{center}
\begin{tabular}{c|c|c}
 & $x + y$ 或 $x - y$ & $xy$ \\
\hline
运算 & 加减法 & 乘法 \\
偏导数 & 常数（系数） & 包含另一个变量 \\
积分 & 各项分开积 & 整体一起积 \\
可分离？ & ❌ 不能拆成乘积 & ✅ 能拆成乘积 \\
几何意义 & 平面 & 双曲抛物面
\end{tabular}
\end{center}

$x + y$ 和 $x - y$ 是\textbf{线性函数}——图像是一个平面。

$xy$ 是\textbf{双线性函数}——图像是一个马鞍形的曲面。

这是两类完全不同的函数，处理方法的细节不同，但\textbf{核心思路都是一样的}：对谁操作，就动谁，其他的按规则处理。

\section{公式汇总}

\subsection{$x - y$ 的所有公式}

\begin{center}
\begin{tabular}{c|c|c}
操作 & 公式 & 结果 \\
\hline
对 $x$ 求导 & $\dfrac{d}{dx}(x - y)$ & 1 \\
对 $x$ 偏导 & $\dfrac{\partial}{\partial x}(x - y)$ & 1 \\
对 $y$ 偏导 & $\dfrac{\partial}{\partial y}(x - y)$ & \textbf{$-1$} \\
对 $x$ 不定积分 & $\displaystyle\int (x - y) \, dx$ & $\dfrac{x^2}{2} - xy + C$ \\
对 $y$ 不定积分 & $\displaystyle\int (x - y) \, dy$ & $xy - \dfrac{y^2}{2} + C$
\end{tabular}
\end{center}

\subsection{$x + y$ 与 $x - y$ 的对比}

\begin{center}
\begin{tabular}{c|c|c}
 & $x + y$ & $x - y$ \\
\hline
对 $y$ 偏导 & $+1$ & \textbf{$-1$} \\
对 $x$ 积分 & $\frac{x^2}{2} + xy + C$ & $\frac{x^2}{2} - xy + C$ \\
矩形 $[0,2]\times[0,3]$ 上的积分 & $+15$ & \textbf{$-3$} \\
对称三角形上的积分 & $\frac{1}{3}$ & \textbf{0}
\end{tabular}
\end{center}

\subsection{核心规律}

\textbf{规律一：减法就是加上负数}

\[
x - y = x + (-1) \cdot y
\]

所有 $x + y$ 的规律，把 $y$ 的系数从 $+1$ 换成 $-1$，就变成 $x - y$ 的规律。

\textbf{规律二：偏导数就是系数}

\[
\frac{\partial}{\partial y}(x + by) = b
\]

$b = +1$ 时得 $+1$，$b = -1$ 时得 $-1$。

\textbf{规律三：积分可以分配}

\[
\iint (x - y) \, dA = \iint x \, dA - \iint y \, dA
\]

减号可以提出来，分别积分再相减。

\section{你学到了什么}

\textbf{第一}，$x - y$ 就是 $x + (-y)$，减法就是加上一个负数。

\textbf{第二}，$x - y$ 对 $x$ 的偏导是 $+1$，对 $y$ 的偏导是 \textbf{$-1$}。$y$ 增加 1，$f$ 减少 1。

\textbf{第三}，$x - y$ 的积分只是把 $x + y$ 的积分里 $y$ 前面的符号变成负号。

\textbf{第四}，在关于 $y = x$ 对称的区域（比如等腰直角三角形）上，$x - y$ 的积分是 \textbf{0}，因为正负完全抵消。

\textbf{第五}，加减法的函数不能像乘法那样“拆分”，但可以利用积分的\textbf{线性性质}分别积分再相加减。

\textbf{第六}，减法带来的唯一变化就是\textbf{符号}。规律不变，只是有些地方正变负、负变正。

不管是正着来还是反着来，关系还是\textbf{线性}的。学会了 $x + y$，$x - y$ 自动就会了。

以后遇到 $3x - 5y + 2z$ 这种更复杂的，也是一样的道理——每一项分开处理，最后按符号加起来。

\chapter{$x-y-z$：难？规律别忘！}

\section{开场：两个变量都反着来}

上一篇，$y$ 闹脾气，把 $x + y$ 变成了 $x - y$。

现在，$z$ 也学坏了，跟着一起倒着来。

于是，$x + y + z$ 变成了 \textbf{$x - y - z$}。

这有什么难的？两个负号而已。规律还是一样的。

$x - y - z = x + (-1)\cdot y + (-1)\cdot z$。把系数换一换，其他全都不变。

\section{一、偏导数}

\subsection{公式}

\[
f(x, y, z) = x - y - z
\]

可以写成：

\[
f(x, y, z) = (+1) \cdot x + (-1) \cdot y + (-1) \cdot z
\]

根据上一章的规律：\textbf{偏导数就是系数}。

\[
\frac{\partial f}{\partial x} = +1
\]

\[
\frac{\partial f}{\partial y} = -1
\]

\[
\frac{\partial f}{\partial z} = -1
\]

就这么简单。

\subsection{用数字验证}

假设我们在点 $(10, 3, 2)$ 处，此时 $f = 10 - 3 - 2 = 5$。

\textbf{验证对 $x$ 的偏导：}

$x$ 从 10 变到 12（变了 2），$y$ 和 $z$ 不动。

$f$ 从 5 变到 $12 - 3 - 2 = 7$（变了 2）。

偏导数 $= 2 \div 2 = +1$ ✓

\textbf{验证对 $y$ 的偏导：}

$y$ 从 3 变到 7（变了 4），$x$ 和 $z$ 不动。

$f$ 从 5 变到 $10 - 7 - 2 = 1$（变了 $-4$）。

偏导数 $= -4 \div 4 = -1$ ✓

\textbf{验证对 $z$ 的偏导：}

$z$ 从 2 变到 5（变了 3），$x$ 和 $y$ 不动。

$f$ 从 5 变到 $10 - 3 - 5 = 2$（变了 $-3$）。

偏导数 $= -3 \div 3 = -1$ ✓

\subsection{🎨 动手画一画：三个方向的影响}

\textbf{第一步：} 画一个三维坐标轴（和之前画四面体一样）。$x$ 往右，$y$ 往左上斜，$z$ 往上。

\textbf{第二步：} 在原点标上 $f = 0$（因为 $0 - 0 - 0 = 0$）。

\textbf{第三步：} 沿 $x$ 方向画一个箭头往右，旁边写 \textbf{$+1$}（$x$ 增加，$f$ 增加）。

\textbf{第四步：} 沿 $y$ 方向画一个箭头，旁边写 \textbf{$-1$}（$y$ 增加，$f$ 减少）。

\textbf{第五步：} 沿 $z$ 方向画一个箭头，旁边写 \textbf{$-1$}（$z$ 增加，$f$ 减少）。

看看你的图——三个方向的影响一目了然。

\subsection{和其他函数的对比}

\begin{center}
\begin{tabular}{c|c|c|c}
函数 & 对 $x$ 偏导 & 对 $y$ 偏导 & 对 $z$ 偏导 \\
\hline
$x + y + z$ & $+1$ & $+1$ & $+1$ \\
$x - y + z$ & $+1$ & \textbf{$-1$} & $+1$ \\
$x + y - z$ & $+1$ & $+1$ & \textbf{$-1$} \\
$x - y - z$ & $+1$ & \textbf{$-1$} & \textbf{$-1$} \\
$-x - y - z$ & \textbf{$-1$} & \textbf{$-1$} & \textbf{$-1$}
\end{tabular}
\end{center}

\textbf{每个变量前面的符号是什么，它的偏导数就是什么。}

加号就是 $+1$，减号就是 $-1$。就这么简单。

\section{二、积分}

\subsection{不定积分}

\subsubsection{对 $x$ 积分}

\[
\int (x - y - z) \, dx
\]

$y$ 和 $z$ 都是常数。

\[
= \int x \, dx - \int y \, dx - \int z \, dx
\]

\[
= \frac{x^2}{2} - yx - zx + C
\]

\[
= \frac{x^2}{2} - (y + z)x + C
\]

验算：对结果求 $\partial/\partial x$：

\[
\frac{\partial}{\partial x}\left(\frac{x^2}{2} - yx - zx\right) = x - y - z \quad \checkmark
\]

\subsubsection{对 $y$ 积分}

\[
\int (x - y - z) \, dy
\]

$x$ 和 $z$ 都是常数。

\[
= \int x \, dy - \int y \, dy - \int z \, dy
\]

\[
= xy - \frac{y^2}{2} - zy + C
\]

\subsubsection{对 $z$ 积分}

\[
\int (x - y - z) \, dz
\]

$x$ 和 $y$ 都是常数。

\[
= \int x \, dz - \int y \, dz - \int z \, dz
\]

\[
= xz - yz - \frac{z^2}{2} + C
\]

\subsection{和 $x + y + z$ 的对比}

\begin{center}
\begin{tabular}{c|c|c}
操作 & $x + y + z$ & $x - y - z$ \\
\hline
对 $x$ 积分 & $\frac{x^2}{2} + yx + zx + C$ & $\frac{x^2}{2} - yx - zx + C$ \\
对 $y$ 积分 & $xy + \frac{y^2}{2} + zy + C$ & $xy - \frac{y^2}{2} - zy + C$ \\
对 $z$ 积分 & $xz + yz + \frac{z^2}{2} + C$ & $xz - yz - \frac{z^2}{2} + C$
\end{tabular}
\end{center}

看到规律了吗？\textbf{哪个变量前面是减号，积分结果里它相关的项就是减号。}

\subsection{定积分}

\subsubsection{例子：$x$ 从 0 到 2，$y = 3$，$z = 1$}

\[
\int_0^2 (x - 3 - 1) \, dx = \int_0^2 (x - 4) \, dx
\]

\textbf{第一步：} 不定积分：$\frac{x^2}{2} - 4x$

\textbf{第二步：} $x = 2$ 代进去：$\frac{4}{2} - 8 = 2 - 8 = -6$

\textbf{第三步：} $x = 0$ 代进去：$0$

\textbf{第四步：} 上减下：$-6 - 0 = -6$

从 $x = 0$ 收集到 $x = 2$（就是沿着数轴从点 \textbf{0} 走到点 \textbf{2}，算曲线 $x - 4$ 下面的“有符号面积”），结果是 \textbf{$-6$}。

\subsection{🎨 动手画一画：负面积}

\textbf{第一步：} 画数轴，标上 0、1、2、3、4、5。

\textbf{第二步：} 画出 $f(x) = x - 4$ 这条直线。

\begin{itemize}
\item $x = 0$ 时，$f = -4$（在 $x$ 轴下方）
\item $x = 2$ 时，$f = -2$（还是在 $x$ 轴下方）
\item $x = 4$ 时，$f = 0$（穿过 $x$ 轴）
\end{itemize}

\textbf{第三步：} 把 $x = 0$ 到 $x = 2$ 之间的区域涂上颜色。这整块区域都在 $x$ 轴\textbf{下方}。

\textbf{第四步：} 这是一个梯形，上底 2，下底 4，高 2。面积 $= \frac{(2+4) \times 2}{2} = 6$。

但因为在 $x$ 轴下方，算作 \textbf{$-6$}。

和公式算出来的一致！

\section{三、二重积分}

\subsection{矩形 $[0, 2] \times [0, 3]$ 上，$z = 1$}

\[
\int_0^3 \int_0^2 (x - y - 1) \, dx \, dy
\]

\textbf{第一步：先对 $x$ 积分}

\[
\int_0^2 (x - y - 1) \, dx = \left[\frac{x^2}{2} - yx - x\right]_0^2 = \left(2 - 2y - 2\right) - 0 = -2y
\]

从 $x = 0$ 收集到 $x = 2$（就是固定某一行的 $y$ 值不动，从点 $(0, y)$ 水平走到点 $(2, y)$，算这一行上 $x - y - 1$ 的总量）。

\textbf{第二步：再对 $y$ 积分}

\[
\int_0^3 (-2y) \, dy = \left[-y^2\right]_0^3 = -9 - 0 = -9
\]

从 $y = 0$ 收集到 $y = 3$（就是从底边 $(x, 0)$ 那一行开始，一行一行往上扫，一直扫到顶边 $(x, 3)$ 那一行，把所有行的结果加起来）。

这个矩形上的二重积分总量是 \textbf{$-9$}。

\subsection{🎨 动手画一画：哪里是正，哪里是负}

\textbf{第一步：} 画出矩形 $[0, 2] \times [0, 3]$。

\textbf{第二步：} 在矩形里标出几个点的 $x - y - 1$ 的值：

\begin{itemize}
\item $(0, 0)$：$0 - 0 - 1 = -1$（负）
\item $(2, 0)$：$2 - 0 - 1 = 1$（正）
\item $(0, 3)$：$0 - 3 - 1 = -4$（负）
\item $(2, 3)$：$2 - 3 - 1 = -2$（负）
\end{itemize}

\textbf{第三步：} 画出 $x - y - 1 = 0$ 这条线，也就是 $y = x - 1$。这条线在 $(1, 0)$ 进入矩形，在 $(2, 1)$ 离开矩形。

\textbf{第四步：} 这条线的下方（右下角那一小块）是正的，其他地方都是负的。用不同颜色标出来。

\textbf{第五步：} 观察——负的区域远远大于正的区域。所以总积分是负的。

\section{四、三重积分}

\subsection{长方体 $[0, 2] \times [0, 3] \times [0, 4]$}

\[
\int_0^4 \int_0^3 \int_0^2 (x - y - z) \, dx \, dy \, dz
\]

三重积分！别慌，一步一步来。

\textbf{第一步：先对 $x$ 积分}

\[
\int_0^2 (x - y - z) \, dx = \left[\frac{x^2}{2} - yx - zx\right]_0^2 = (2 - 2y - 2z) - 0 = 2 - 2y - 2z
\]

从 $x = 0$ 收集到 $x = 2$（就是固定 $y$ 和 $z$，从点 $(0, y, z)$ 水平往右走到点 $(2, y, z)$，算这一条线上 $x - y - z$ 的总量）。

\textbf{第二步：再对 $y$ 积分}

\[
\int_0^3 (2 - 2y - 2z) \, dy = \left[2y - y^2 - 2zy\right]_0^3 = (6 - 9 - 6z) - 0 = -3 - 6z
\]

从 $y = 0$ 收集到 $y = 3$（就是固定 $z$，把这一排横棍全部加起来，合成一片水平的楼层）。

\textbf{第三步：最后对 $z$ 积分}

\[
\int_0^4 (-3 - 6z) \, dz = \left[-3z - 3z^2\right]_0^4 = (-12 - 48) - 0 = -60
\]

从 $z = 0$ 收集到 $z = 4$（就是从地面那层楼开始，一层一层往上叠，整个盒子就扫完了）。

这个长方体的三重积分总量是 \textbf{$-60$}！

\subsection{和 $x + y + z$ 的对比}

\[
\int_0^4 \int_0^3 \int_0^2 (x + y + z) \, dx \, dy \, dz = ?
\]

用同样的方法计算：

\textbf{对 $x$ 积分：} $2 + 2y + 2z$

\textbf{对 $y$ 积分：} $(6 + 9 + 6z) = 15 + 6z$

\textbf{对 $z$ 积分：} $[15z + 3z^2]_0^4 = 60 + 48 = 108$

\begin{center}
\begin{tabular}{c|c}
函数 & 长方体 $[0,2]\times[0,3]\times[0,4]$ 上的积分 \\
\hline
$x + y + z$ & \textbf{$+108$} \\
$x - y - z$ & \textbf{$-60$}
\end{tabular}
\end{center}

相差 168！这就是符号的力量。

\subsection{为什么差这么多？}

$x + y + z$ 的时候，三个变量都在加分。整个盒子里每个点都是正贡献。

$x - y - z$ 的时候，只有 $x$ 在加分，$y$ 和 $z$ 都在减分。

而且在这个盒子里，$y$ 最大到 3，$z$ 最大到 4，$x$ 最大只到 2。

两个减分的变量（$y$ 和 $z$）的范围都比加分的变量（$x$）大。

所以减分效应占了绝对优势，总量是负的。

\subsection{🎨 动手画一画：谁在加分，谁在减分}

\textbf{第一步：} 画一个长方体，标上三条边：$x$ 从 0 到 2，$y$ 从 0 到 3，$z$ 从 0 到 4。

\textbf{第二步：} 在长方体里标出几个点的 $x - y - z$ 的值：

\begin{itemize}
\item $(0, 0, 0)$：$0 - 0 - 0 = 0$
\item $(2, 0, 0)$：$2 - 0 - 0 = 2$（正）
\item $(0, 3, 0)$：$0 - 3 - 0 = -3$（负）
\item $(0, 0, 4)$：$0 - 0 - 4 = -4$（负）
\item $(2, 3, 4)$：$2 - 3 - 4 = -5$（负）
\end{itemize}

\textbf{第三步：} 找到 $x - y - z = 0$ 这个平面。这个平面把长方体分成两部分。

\textbf{第四步：} 观察——在这个长方体里，负的区域比正的区域大得多。

\textbf{第五步：} 在图旁边写：“$x$ 太小（只到 2），$y$ 和 $z$ 太大（到 3 和 4），所以负的占优。”

\section{五、四面体上的 $x - y - z$}

\subsection{四面体的定义}

顶点：$(0, 0, 0)$、$(1, 0, 0)$、$(0, 1, 0)$、$(0, 0, 1)$。

内部满足：$x \geq 0$，$y \geq 0$，$z \geq 0$，$x + y + z \leq 1$。

\subsection{积分}

\[
\int_0^1 \int_0^{1-x} \int_0^{1-x-y} (x - y - z) \, dz \, dy \, dx
\]

看起来很复杂？别怕，一步步来。

\textbf{第一步：对 $z$ 积分}

\[
\int_0^{1-x-y} (x - y - z) \, dz = \left[(x-y)z - \frac{z^2}{2}\right]_0^{1-x-y}
\]

把 $z = 1 - x - y$ 代进去：

\[
= (x-y)(1-x-y) - \frac{(1-x-y)^2}{2}
\]

令 $u = 1 - x - y$（为了简化书写）：

\[
= (x-y)u - \frac{u^2}{2}
\]

从 $z = 0$ 收集到 $z = 1 - x - y$（就是从底面上的点 $(x, y, 0)$ 竖直往上走到斜面上的点 $(x, y, 1-x-y)$ 停下）。

展开（把 $u = 1 - x - y$ 代回去）：

\[
(x-y)(1-x-y) = (x-y) - (x-y)(x+y) = (x-y) - (x^2 - y^2) = x - y - x^2 + y^2
\]

\[
\frac{(1-x-y)^2}{2} = \frac{1 - 2x - 2y + x^2 + 2xy + y^2}{2} = \frac{1}{2} - x - y + \frac{x^2}{2} + xy + \frac{y^2}{2}
\]

相减：

\[
(x - y - x^2 + y^2) - \left(\frac{1}{2} - x - y + \frac{x^2}{2} + xy + \frac{y^2}{2}\right)
\]

\[
= x - y - x^2 + y^2 - \frac{1}{2} + x + y - \frac{x^2}{2} - xy - \frac{y^2}{2}
\]

\[
= 2x - \frac{3x^2}{2} - xy + \frac{y^2}{2} - \frac{1}{2}
\]

第一步完成。有点乱，但我们继续。

\textbf{第二步：对 $y$ 积分}

\[
\int_0^{1-x} \left(2x - \frac{3x^2}{2} - xy + \frac{y^2}{2} - \frac{1}{2}\right) dy
\]

从 $y = 0$ 收集到 $y = 1 - x$（就是从后墙上的点 $(x, 0, z)$ 往前扫到斜边上的点 $(x, 1-x, z)$ 停下）。

逐项积分：

\[
\int_0^{1-x} 2x \, dy = 2x(1-x) = 2x - 2x^2
\]

\[
\int_0^{1-x} \frac{3x^2}{2} \, dy = \frac{3x^2}{2}(1-x) = \frac{3x^2}{2} - \frac{3x^3}{2}
\]

\[
\int_0^{1-x} xy \, dy = x \cdot \frac{(1-x)^2}{2} = \frac{x(1-2x+x^2)}{2} = \frac{x - 2x^2 + x^3}{2}
\]

\[
\int_0^{1-x} \frac{y^2}{2} \, dy = \frac{1}{2} \cdot \frac{(1-x)^3}{3} = \frac{(1-x)^3}{6}
\]

\[
\int_0^{1-x} \frac{1}{2} \, dy = \frac{1-x}{2}
\]

合在一起（注意符号）：

\[
(2x - 2x^2) - \left(\frac{3x^2}{2} - \frac{3x^3}{2}\right) - \left(\frac{x - 2x^2 + x^3}{2}\right) + \frac{(1-x)^3}{6} - \frac{1-x}{2}
\]

这很复杂。让我们展开 $(1-x)^3 = 1 - 3x + 3x^2 - x^3$：

\[
\frac{(1-x)^3}{6} = \frac{1}{6} - \frac{x}{2} + \frac{x^2}{2} - \frac{x^3}{6}
\]

现在把所有项按 $x$ 的次数整理：

\textbf{常数项：} $\frac{1}{6} - \frac{1}{2} = \frac{1}{6} - \frac{3}{6} = -\frac{2}{6} = -\frac{1}{3}$

\textbf{$x$ 的一次项：} $2x - \frac{x}{2} - \frac{x}{2} + \frac{x}{2} = 2x - \frac{x}{2} = \frac{3x}{2}$

\textbf{$x$ 的二次项：} $-2x^2 - \frac{3x^2}{2} + x^2 + \frac{x^2}{2} = -2x^2 - \frac{3x^2}{2} + \frac{3x^2}{2} = -2x^2$

\textbf{$x$ 的三次项：} $\frac{3x^3}{2} - \frac{x^3}{2} - \frac{x^3}{6} = x^3 - \frac{x^3}{6} = \frac{5x^3}{6}$

所以第二步的结果是：

\[
-\frac{1}{3} + \frac{3x}{2} - 2x^2 + \frac{5x^3}{6}
\]

\textbf{第三步：对 $x$ 积分}

\[
\int_0^1 \left(-\frac{1}{3} + \frac{3x}{2} - 2x^2 + \frac{5x^3}{6}\right) dx
\]

从 $x = 0$ 收集到 $x = 1$（就是从左墙 $(0, y, z)$ 那一片开始，一片一片往右推，整个四面体就扫完了）。

逐项积分：

\[
\int_0^1 -\frac{1}{3} \, dx = -\frac{1}{3}
\]

\[
\int_0^1 \frac{3x}{2} \, dx = \frac{3}{2} \cdot \frac{1}{2} = \frac{3}{4}
\]

\[
\int_0^1 2x^2 \, dx = 2 \cdot \frac{1}{3} = \frac{2}{3}
\]

\[
\int_0^1 \frac{5x^3}{6} \, dx = \frac{5}{6} \cdot \frac{1}{4} = \frac{5}{24}
\]

合在一起（注意符号）：

\[
-\frac{1}{3} + \frac{3}{4} - \frac{2}{3} + \frac{5}{24}
\]

通分（分母取 24）：

\[
-\frac{8}{24} + \frac{18}{24} - \frac{16}{24} + \frac{5}{24} = \frac{-8 + 18 - 16 + 5}{24} = \frac{-1}{24}
\]

四面体上 $x - y - z$ 的积分是 $\mathbf{-\dfrac{1}{24}}$！

\subsection{惊人的发现}

回忆一下之前的结果：

\begin{center}
\begin{tabular}{c|c}
函数 & 四面体上的积分 \\
\hline
$x + y + z$ & ? \\
$xyz$ & $\frac{1}{720}$ \\
$x - y - z$ & $-\frac{1}{24}$
\end{tabular}
\end{center}

让我们算一下 $x + y + z$ 在四面体上的积分。

根据对称性，$\iiint_T x \, dV = \iiint_T y \, dV = \iiint_T z \, dV$（因为四面体关于 $x$、$y$、$z$ 轮换对称）。

所以：

\[
\iiint_T (x + y + z) \, dV = 3 \iiint_T x \, dV
\]

而 $\iiint_T x \, dV$ 可以计算出来是 $\frac{1}{24}$（留作练习，方法和之前一样）。

所以：

\[
\iiint_T (x + y + z) \, dV = 3 \times \frac{1}{24} = \frac{1}{8}
\]

现在来验证 $x - y - z$：

\[
\iiint_T (x - y - z) \, dV = \iiint_T x \, dV - \iiint_T y \, dV - \iiint_T z \, dV
\]

\[
= \frac{1}{24} - \frac{1}{24} - \frac{1}{24} = -\frac{1}{24} \quad \checkmark
\]

和我们用硬算得到的结果完全一致！

原来可以这样算！不用那么复杂的展开！

\section{六、规律总结}

\subsection{线性函数的一般形式}

\[
f(x, y, z) = ax + by + cz
\]

其中 $a$、$b$、$c$ 是常数。

\textbf{特殊情况：}

\begin{center}
\begin{tabular}{c|c|c|c}
$a$ & $b$ & $c$ & 函数 \\
\hline
1 & 1 & 1 & $x + y + z$ \\
1 & $-1$ & $-1$ & $x - y - z$ \\
1 & $-1$ & 1 & $x - y + z$ \\
2 & 3 & $-5$ & $2x + 3y - 5z$
\end{tabular}
\end{center}

\subsection{偏导数}

\[
\frac{\partial}{\partial x}(ax + by + cz) = a
\]

\[
\frac{\partial}{\partial y}(ax + by + cz) = b
\]

\[
\frac{\partial}{\partial z}(ax + by + cz) = c
\]

\textbf{口诀：偏导数就是系数。}

\subsection{积分}

\[
\int (ax + by + cz) \, dx = \frac{a x^2}{2} + bxy + cxz + C
\]

\[
\int (ax + by + cz) \, dy = axy + \frac{b y^2}{2} + cyz + C
\]

\[
\int (ax + by + cz) \, dz = axz + byz + \frac{c z^2}{2} + C
\]

\textbf{口诀：谁被积分，谁的指数加一放分母；其他变量就乘上来。}

\subsection{多重积分的分配律}

\[
\iiint (ax + by + cz) \, dV = a \iiint x \, dV + b \iiint y \, dV + c \iiint z \, dV
\]

\textbf{口诀：常数可以提出来，加减法可以分开算。}

\subsection{对称区域上的简化}

如果积分区域关于 $x$、$y$、$z$ \textbf{轮换对称}（比如四面体），那么：

\[
\iiint x \, dV = \iiint y \, dV = \iiint z \, dV
\]

这意味着：

\[
\iiint (x - y - z) \, dV = \iiint x \, dV - \iiint y \, dV - \iiint z \, dV = I - I - I = -I
\]

其中 $I = \iiint x \, dV$。

\subsection{🎨 动手画一画：做一张终极速查表}

\begin{quote}
\ttfamily
线性函数 $ax + by + cz$ 速查\\
\\
偏导数：\\
\\
\quad \(\partial/\partial x = a,\ \partial/\partial y = b,\ \partial/\partial z = c\)\\
\\
\quad 口诀：偏导数就是系数\\
\\
积分（对 $x$）：\\
\\
\quad \(\int dx = (a/2)x^2 + bxy + cxz + C\)\\
\\
\quad 口诀：谁被积分，指数加一放分母\\
\\
多重积分：\\
\\
\quad \(\int\!\!\int\!\!\int (ax+by+cz)\,dV = a\int\!\!\int\!\!\int x\,dV + b\int\!\!\int\!\!\int y\,dV + c\int\!\!\int\!\!\int z\,dV\)\\
\\
\quad 口诀：常数提出来，加减分开算\\
\\
对称区域：\\
\\
\quad 四面体上：\(\int\!\!\int\!\!\int x\,dV = \int\!\!\int\!\!\int y\,dV = \int\!\!\int\!\!\int z\,dV\)\\
\\
\quad 所以：\(\int\!\!\int\!\!\int (x-y-z)\,dV = I - I - I = -I\)
\end{quote}

\section{公式汇总}

\subsection{$x - y - z$ 的所有公式}

\begin{center}
\begin{tabular}{c|c}
操作 & 结果 \\
\hline
对 $x$ 偏导 & $+1$ \\
对 $y$ 偏导 & $-1$ \\
对 $z$ 偏导 & $-1$ \\
对 $x$ 积分 & $\frac{x^2}{2} - xy - xz + C$ \\
对 $y$ 积分 & $xy - \frac{y^2}{2} - yz + C$ \\
对 $z$ 积分 & $xz - yz - \frac{z^2}{2} + C$ \\
四面体上的三重积分 & $-\frac{1}{24}$ \\
长方体 $[0,2]\times[0,3]\times[0,4]$ 上的积分 & $-60$
\end{tabular}
\end{center}

\subsection{不同函数的对比}

\begin{center}
\begin{tabular}{c|c|c}
函数 & 四面体上的积分 & 特点 \\
\hline
$x + y + z$ & $\frac{1}{8}$ & 全正，最大 \\
$x - y + z$ & $\frac{1}{24}$ & 一正一负 \\
$x + y - z$ & $\frac{1}{24}$ & 一正一负 \\
$x - y - z$ & $-\frac{1}{24}$ & 两负，最小 \\
$xyz$ & $\frac{1}{720}$ & 乘法，更小
\end{tabular}
\end{center}

\section{你学到了什么}

\textbf{第一}，$x - y - z$ 就是 $x + (-y) + (-z)$，把系数换成负号，其他不变。

\textbf{第二}，偏导数就是系数：对 $x$ 是 $+1$，对 $y$ 是 $-1$，对 $z$ 是 $-1$。

\textbf{第三}，积分的时候，减号跟着变量走。哪个变量是负的，它相关的项就带负号。

\textbf{第四}，多重积分可以用\textbf{分配律}分开算：

\[
\iiint (x - y - z) \, dV = \iiint x \, dV - \iiint y \, dV - \iiint z \, dV
\]

\textbf{第五}，在对称区域（如四面体）上，三个单独的积分相等，可以大大简化计算。

\textbf{第六}，规律永远是一样的：\textbf{加减法改变符号，不改变方法。}

不管是加还是减，不管有几个负号，方法永远不变：

\begin{quote}
\textbf{偏导数 = 系数}

\textbf{积分 = 指数加一放分母}

\textbf{多重积分 = 分开算，再合起来}
\end{quote}

记住这三条，任何线性函数都能搞定。

\section*{附录：快速检验表}

用这张表检验你的计算：

\begin{center}
\begin{tabular}{c|c|c|c|c}
函数 & $\partial/\partial x$ & $\partial/\partial y$ & $\partial/\partial z$ & 四面体积分 \\
\hline
$x$ & 1 & 0 & 0 & $1/24$ \\
$y$ & 0 & 1 & 0 & $1/24$ \\
$z$ & 0 & 0 & 1 & $1/24$ \\
$x+y+z$ & 1 & 1 & 1 & $1/8 = 3/24$ \\
$x-y-z$ & 1 & $-1$ & $-1$ & $-1/24$ \\
$x+y-z$ & 1 & 1 & $-1$ & $1/24$ \\
$2x-3y+z$ & 2 & $-3$ & 1 & $(2-3+1)/24 = 0$
\end{tabular}
\end{center}

最后一个例子：$2x - 3y + z$ 在四面体上的积分是 0！因为 $2 - 3 + 1 = 0$。

这就是规律的力量。

\chapter{出现次方：通用解决方法}

\section{开场：当指数不再是 1}

前几篇里，我们处理的函数都很“温和”：

\begin{itemize}
\item $x + y + z$（每个变量的指数都是 1）
\item $x - y - z$（还是指数 1，只是符号变了）
\item $xy$、$xyz$（指数还是 1，只是变成了乘法）
\end{itemize}

但现在，我们要升级了。

$x$ 变成了\textbf{平方}，$y$ 变成了\textbf{立方}，$z$ 变成了\textbf{四次方}！

我们会遇到这样的“怪物”：

\[
x^2 + y^3 - z^4
\]

\[
x^2 y^3 z
\]

\[
x^3 - 2y^2 + 5z^4
\]

别慌。\textbf{规律还是那两条}，只是数字变大了而已。

\textbf{“指数加一放分母，指数搬下来当系数。”} 这两句口诀，永远管用。

\section{一、回顾两条核心规律}

在继续之前，让我们把两条核心规律刻在脑子里。

\subsection{规律一：求导（偏导数）}

\[
\frac{\partial}{\partial x}(x^n) = n \cdot x^{n-1}
\]

\textbf{口诀：指数搬下来当系数，指数减一。}

\begin{center}
\begin{tabular}{c|c|c}
原函数 & 对 $x$ 求导 & 怎么来的 \\
\hline
$x^1 = x$ & $1 \cdot x^0 = 1$ & 指数 1 搬下来，1−1=0 \\
$x^2$ & $2x^1 = 2x$ & 指数 2 搬下来，2−1=1 \\
$x^3$ & $3x^2$ & 指数 3 搬下来，3−1=2 \\
$x^4$ & $4x^3$ & 指数 4 搬下来，4−1=3 \\
$x^{10}$ & $10x^9$ & 指数 10 搬下来，10−1=9 \\
$x^{100}$ & $100x^{99}$ & 指数 100 搬下来，100−1=99
\end{tabular}
\end{center}

\subsection{规律二：积分}

\[
\int x^n \, dx = \frac{x^{n+1}}{n+1} + C
\]

\textbf{口诀：指数加一放分母。}

\begin{center}
\begin{tabular}{c|c|c}
原函数 & 对 $x$ 积分 & 怎么来的 \\
\hline
$x^0 = 1$ & $\frac{x^1}{1} = x$ & 0+1=1 放分母 \\
$x^1 = x$ & $\frac{x^2}{2}$ & 1+1=2 放分母 \\
$x^2$ & $\frac{x^3}{3}$ & 2+1=3 放分母 \\
$x^3$ & $\frac{x^4}{4}$ & 3+1=4 放分母 \\
$x^4$ & $\frac{x^5}{5}$ & 4+1=5 放分母 \\
$x^{10}$ & $\frac{x^{11}}{11}$ & 10+1=11 放分母
\end{tabular}
\end{center}

\subsection{🎨 动手画一画：两条规律的记忆卡}

拿一张卡片，正面写：

\begin{quote}
\ttfamily
求导：指数搬下来，指数减一\\
\quad $(x^n)' = n \cdot x^{(n-1)}$
\end{quote}

背面写：

\begin{quote}
\ttfamily
积分：指数加一放分母\\
\quad $\int x^n dx = x^{(n+1)}/(n+1) + C$
\end{quote}

把这张卡片贴在书桌上。这两条规律是\textbf{一切}的基础。

\section{二、加法型——各项分开处理}

\subsection{形式}

\[
f(x, y, z) = ax^m + by^n + cz^p
\]

这是“加法型”函数。每一项只包含一个变量，各项之间用加减号连接。

\subsection{处理方法：分开处理}

\textbf{偏导数：只对包含该变量的项求导，其他项当常数（导数为 0）。}

\textbf{积分：每一项分别积分，然后加起来。}

\subsection{例子 1：$f(x, y) = x^2 + y^3$}

\subsubsection{偏导数}

\[
\frac{\partial f}{\partial x} = \frac{\partial}{\partial x}(x^2) + \frac{\partial}{\partial x}(y^3) = 2x + 0 = 2x
\]

（$y^3$ 里面没有 $x$，对 $x$ 求导就是 0）

\[
\frac{\partial f}{\partial y} = \frac{\partial}{\partial y}(x^2) + \frac{\partial}{\partial y}(y^3) = 0 + 3y^2 = 3y^2
\]

（$x^2$ 里面没有 $y$，对 $y$ 求导就是 0）

\subsubsection{积分}

\[
\int (x^2 + y^3) \, dx = \frac{x^3}{3} + y^3 x + C
\]

（$x^2$ 积分得 $x^3/3$；$y^3$ 对 $x$ 来说是常数，积分得 $y^3 \cdot x$）

\[
\int (x^2 + y^3) \, dy = x^2 y + \frac{y^4}{4} + C
\]

\subsection{例子 2：$f(x, y, z) = x^3 - 2y^2 + 5z^4$}

\subsubsection{偏导数}

\[
\frac{\partial f}{\partial x} = 3x^2 - 0 + 0 = 3x^2
\]

\[
\frac{\partial f}{\partial y} = 0 - 2 \cdot 2y + 0 = -4y
\]

\[
\frac{\partial f}{\partial z} = 0 - 0 + 5 \cdot 4z^3 = 20z^3
\]

\subsubsection{积分（对 $x$）}

\[
\int (x^3 - 2y^2 + 5z^4) \, dx = \frac{x^4}{4} - 2y^2 x + 5z^4 x + C
\]

\subsection{例子 3：$f(x, y, z) = x^5 + y^5 + z^5$}

\subsubsection{偏导数}

\[
\frac{\partial f}{\partial x} = 5x^4, \quad \frac{\partial f}{\partial y} = 5y^4, \quad \frac{\partial f}{\partial z} = 5z^4
\]

这个函数是\textbf{对称的}：三个偏导数的形式完全一样，只是字母不同。

\subsubsection{在单位正方体上的积分}

\[
\int_0^1 \int_0^1 \int_0^1 (x^5 + y^5 + z^5) \, dx \, dy \, dz
\]

分开算：

\[
= \int_0^1 \int_0^1 \int_0^1 x^5 \, dx \, dy \, dz + \int_0^1 \int_0^1 \int_0^1 y^5 \, dx \, dy \, dz + \int_0^1 \int_0^1 \int_0^1 z^5 \, dx \, dy \, dz
\]

每一项都类似。以第一项为例：

\[
\int_0^1 x^5 \, dx = \frac{1}{6}
\]

\[
\int_0^1 dy = 1
\]

\[
\int_0^1 dz = 1
\]

所以第一项 $= \frac{1}{6} \times 1 \times 1 = \frac{1}{6}$

三项都一样（对称性），所以总和 $= 3 \times \frac{1}{6} = \frac{1}{2}$

\subsection{🎨 动手画一画：加法型的处理流程}

画一个流程图：

\begin{quote}
\ttfamily
加法型函数 $f = A + B + C$\\
\quad ↓\\
分别处理 $A$，处理 $B$，处理 $C$\\
\quad ↓\\
结果加起来
\end{quote}

每一项独立处理，最后合并。简单！

\section{三、乘法型——整体处理}

\subsection{形式}

\[
f(x, y, z) = x^m y^n z^p
\]

这是“乘法型”函数。所有变量乘在一起。

\subsection{处理方法：对谁操作，动谁；其他当常数}

\textbf{偏导数：}

\[
\frac{\partial}{\partial x}(x^m y^n z^p) = m \cdot x^{m-1} \cdot y^n \cdot z^p
\]

\textbf{积分：}

\[
\int x^m y^n z^p \, dx = \frac{x^{m+1}}{m+1} \cdot y^n \cdot z^p + C
\]

\subsection{例子 1：$f(x, y) = x^2 y^3$}

\subsubsection{偏导数}

\[
\frac{\partial f}{\partial x} = 2x \cdot y^3 = 2xy^3
\]

\[
\frac{\partial f}{\partial y} = x^2 \cdot 3y^2 = 3x^2 y^2
\]

\subsubsection{积分}

\[
\int x^2 y^3 \, dx = \frac{x^3}{3} \cdot y^3 + C = \frac{x^3 y^3}{3} + C
\]

\[
\int x^2 y^3 \, dy = x^2 \cdot \frac{y^4}{4} + C = \frac{x^2 y^4}{4} + C
\]

\subsection{例子 2：$f(x, y, z) = x^3 y^2 z^4$}

\subsubsection{偏导数}

\[
\frac{\partial f}{\partial x} = 3x^2 y^2 z^4
\]

\[
\frac{\partial f}{\partial y} = 2x^3 y z^4
\]

\[
\frac{\partial f}{\partial z} = 4x^3 y^2 z^3
\]

\subsubsection{在单位正方体上的积分}

\[
\int_0^1 \int_0^1 \int_0^1 x^3 y^2 z^4 \, dx \, dy \, dz
\]

因为是乘法型，可以拆分：

\[
= \left(\int_0^1 x^3 \, dx\right) \times \left(\int_0^1 y^2 \, dy\right) \times \left(\int_0^1 z^4 \, dz\right)
\]

\[
= \frac{1}{4} \times \frac{1}{3} \times \frac{1}{5} = \frac{1}{60}
\]

\subsection{例子 3：$f(x, y, z) = x^{10} y^{20} z^{30}$}

看起来很吓人？其实一点都不难。

\subsubsection{偏导数}

\[
\frac{\partial f}{\partial x} = 10 x^9 y^{20} z^{30}
\]

\[
\frac{\partial f}{\partial y} = 20 x^{10} y^{19} z^{30}
\]

\[
\frac{\partial f}{\partial z} = 30 x^{10} y^{20} z^{29}
\]

\subsubsection{在单位正方体上的积分}

\[
= \frac{1}{11} \times \frac{1}{21} \times \frac{1}{31} = \frac{1}{11 \times 21 \times 31} = \frac{1}{7161}
\]

就这么简单。指数再大，规律不变。

\section{四、混合型——先拆分，再分别处理}

\subsection{形式}

\[
f(x, y, z) = x^m y^n + y^p z^q - x^r z^s
\]

这是“混合型”函数。既有加减法，又有乘法。

\subsection{处理方法：先按加减法拆分，每一项再按乘法处理}

\subsection{例子 1：$f(x, y) = x^2 y + xy^2$}

\subsubsection{偏导数}

先拆成两项：$A = x^2 y$，$B = xy^2$

\[
\frac{\partial A}{\partial x} = 2xy, \quad \frac{\partial B}{\partial x} = y^2
\]

\[
\frac{\partial f}{\partial x} = 2xy + y^2
\]

\[
\frac{\partial A}{\partial y} = x^2, \quad \frac{\partial B}{\partial y} = 2xy
\]

\[
\frac{\partial f}{\partial y} = x^2 + 2xy
\]

\subsubsection{积分（对 $x$）}

\[
\int (x^2 y + xy^2) \, dx = \frac{x^3 y}{3} + \frac{x^2 y^2}{2} + C
\]

\subsection{例子 2：$f(x, y, z) = x^2 y - yz^2 + x^3 z$}

\subsubsection{偏导数}

拆成三项：$A = x^2 y$，$B = -yz^2$，$C = x^3 z$

\[
\frac{\partial f}{\partial x} = 2xy + 0 + 3x^2 z = 2xy + 3x^2 z
\]

\[
\frac{\partial f}{\partial y} = x^2 - z^2 + 0 = x^2 - z^2
\]

\[
\frac{\partial f}{\partial z} = 0 - 2yz + x^3 = x^3 - 2yz
\]

\subsubsection{积分（对 $z$）}

\[
\int (x^2 y - yz^2 + x^3 z) \, dz = x^2 yz - \frac{yz^3}{3} + \frac{x^3 z^2}{2} + C
\]

\subsection{例子 3：$f(x, y, z) = x^2 + y^2 + z^2 + xy + yz + xz$}

这是一个典型的混合型函数，既有纯次方项，又有交叉乘积项。

\subsubsection{偏导数}

\[
\frac{\partial f}{\partial x} = 2x + y + z
\]

怎么来的？

\begin{itemize}
\item $x^2$ 对 $x$ 求导：$2x$
\item $y^2$ 对 $x$ 求导：0
\item $z^2$ 对 $x$ 求导：0
\item $xy$ 对 $x$ 求导：$y$（$y$ 是常数，所以就是 $y$）
\item $yz$ 对 $x$ 求导：0
\item $xz$ 对 $x$ 求导：$z$
\end{itemize}

同理：

\[
\frac{\partial f}{\partial y} = 2y + x + z
\]

\[
\frac{\partial f}{\partial z} = 2z + y + x
\]

\textbf{注意对称性：三个偏导数的形式是轮换对称的。}

\subsection{🎨 动手画一画：混合型的处理流程}

画一个流程图：

\begin{quote}
\ttfamily
混合型函数 $f = A + B - C + D$\\
\quad ↓\\
先拆成乘法项 $A$、$B$、$C$、$D$\\
\quad ↓\\
每一项用乘法规律处理\\
\quad ↓\\
按符号加减起来
\end{quote}

\section{五、定积分的计算}

\subsection{单位正方形上的公式}

对于任意多项式函数，在单位正方形 $[0,1]^2$ 上的积分：

\textbf{加法型：} 每项分别积分，结果加起来。

\[
\int_0^1 \int_0^1 (x^2 + y^3) \, dx \, dy = \int_0^1 \int_0^1 x^2 \, dx \, dy + \int_0^1 \int_0^1 y^3 \, dx \, dy
\]

\[
= \frac{1}{3} \times 1 + 1 \times \frac{1}{4} = \frac{1}{3} + \frac{1}{4} = \frac{7}{12}
\]

\textbf{乘法型：} 拆成各变量的乘积。

\[
\int_0^1 \int_0^1 x^2 y^3 \, dx \, dy = \frac{1}{3} \times \frac{1}{4} = \frac{1}{12}
\]

\textbf{混合型：} 先拆成各项，每项用相应的规律。

\[
\int_0^1 \int_0^1 (x^2 + x^2 y^3) \, dx \, dy = \frac{1}{3} + \frac{1}{3} \times \frac{1}{4} = \frac{1}{3} + \frac{1}{12} = \frac{5}{12}
\]

\subsection{例子：$f(x, y) = x^3 - 2xy + y^2$ 在 $[0, 2] \times [0, 3]$ 上}

\[
\int_0^3 \int_0^2 (x^3 - 2xy + y^2) \, dx \, dy
\]

\textbf{第一步：先对 $x$ 积分}

\[
\int_0^2 (x^3 - 2xy + y^2) \, dx = \left[\frac{x^4}{4} - x^2 y + y^2 x\right]_0^2
\]

$x = 2$ 代入：$\frac{16}{4} - 4y + 2y^2 = 4 - 4y + 2y^2$

$x = 0$ 代入：0

结果：$4 - 4y + 2y^2$

\textbf{第二步：再对 $y$ 积分}

\[
\int_0^3 (4 - 4y + 2y^2) \, dy = \left[4y - 2y^2 + \frac{2y^3}{3}\right]_0^3
\]

$y = 3$ 代入：$12 - 18 + 18 = 12$

$y = 0$ 代入：0

结果：12

\subsection{例子：三重积分}

\[
\int_0^1 \int_0^1 \int_0^1 (x^2 + y^2 + z^2) \, dx \, dy \, dz
\]

因为是加法型，分开算：

\[
= \int_0^1 \int_0^1 \int_0^1 x^2 \, dx \, dy \, dz + \int_0^1 \int_0^1 \int_0^1 y^2 \, dx \, dy \, dz + \int_0^1 \int_0^1 \int_0^1 z^2 \, dx \, dy \, dz
\]

每一项都一样（对称性）：

\[
= 3 \times \frac{1}{3} \times 1 \times 1 = 1
\]

\section{六、通用解题步骤}

不管遇到什么样的多项式函数，按这个步骤来：

\textbf{步骤一：识别类型}

看函数的结构：

\begin{itemize}
\item \textbf{纯加法型}：$x^2 + y^3 + z^4$（每项只有一个变量）
\item \textbf{纯乘法型}：$x^2 y^3 z^4$（所有变量乘在一起）
\item \textbf{混合型}：$x^2 y + y^3 z - xz^2$（既有加减，又有乘法）
\end{itemize}

\textbf{步骤二：拆分}

如果是混合型，先按加减号拆成若干项。

\textbf{步骤三：对每一项应用规律}

\textbf{求导：} 指数搬下来，指数减一，其他变量当常数。

\textbf{积分：} 指数加一放分母，其他变量当常数。

\textbf{步骤四：合并结果}

按原来的加减号把各项的结果合并。

\textbf{步骤五：验算（可选）}

\begin{itemize}
\item 对积分结果求导，应该得到原函数。
\item 对导数结果积分，应该得到原函数（加常数）。
\end{itemize}

\subsection{🎨 动手画一画：通用解题流程图}

\begin{quote}
\ttfamily
通用解题步骤\\
\\
1. 识别类型\\
\quad 纯加法型？→ 每项分开处理\\
\quad 纯乘法型？→ 对谁操作动谁，其他当常数\\
\quad 混合型？→ 先拆分，再分别处理\\
\\
2. 应用规律\\
\quad 求导：指数搬下来，指数减一\\
\quad 积分：指数加一放分母\\
\\
3. 合并结果\\
\quad 按原来的加减号合并\\
\\
4. 验算\\
\quad 对积分求导，应该得到原函数
\end{quote}

\section{速查表}

\subsection{常用导数}

\begin{center}
\begin{tabular}{c|c}
原函数 & 导数 \\
\hline
$x^0 = 1$ & 0 \\
$x^1 = x$ & 1 \\
$x^2$ & $2x$ \\
$x^3$ & $3x^2$ \\
$x^4$ & $4x^3$ \\
$x^n$ & $nx^{n-1}$ \\
$x^m y^n$（对 $x$） & $mx^{m-1}y^n$ \\
$x^m y^n$（对 $y$） & $nx^m y^{n-1}$
\end{tabular}
\end{center}

\subsection{常用积分}

\begin{center}
\begin{tabular}{c|c}
原函数 & 积分 \\
\hline
1 & $x + C$ \\
$x$ & $\frac{x^2}{2} + C$ \\
$x^2$ & $\frac{x^3}{3} + C$ \\
$x^3$ & $\frac{x^4}{4} + C$ \\
$x^n$ & $\frac{x^{n+1}}{n+1} + C$ \\
$x^m y^n$（对 $x$） & $\frac{x^{m+1}y^n}{m+1} + C$ \\
$x^m y^n$（对 $y$） & $\frac{x^m y^{n+1}}{n+1} + C$
\end{tabular}
\end{center}

\subsection{单位区间/正方形/正方体上的积分}

\begin{center}
\begin{tabular}{c|c|c|c}
函数 & $[0,1]$ 上 & $[0,1]^2$ 上 & $[0,1]^3$ 上 \\
\hline
$x^n$ & $\frac{1}{n+1}$ & — & — \\
$x^m y^n$ & — & $\frac{1}{(m+1)(n+1)}$ & — \\
$x^m y^n z^p$ & — & — & $\frac{1}{(m+1)(n+1)(p+1)}$ \\
$x^n + y^n$ & — & $\frac{2}{n+1}$ & — \\
$x^n + y^n + z^n$ & — & — & $\frac{3}{n+1}$
\end{tabular}
\end{center}

\section{终极口诀}

\textbf{求导：指数搬下来当系数，指数减一，其他不动。}

\textbf{积分：指数加一放分母，其他不动。}

\textbf{加法型：各项分开处理，结果加起来。}

\textbf{乘法型：对谁操作动谁，其他当常数。}

\textbf{混合型：先拆分，再分别处理，最后合并。}

不管指数是 1、2、3，还是 100，规律永远是这几条。

不管是加法、减法、乘法，处理方法永远是先拆分、再应用规律、最后合并。

出现次方不可怕，规律别忘就能解！

\section{你学到了什么}

\textbf{第一}，求导的核心规律：\textbf{指数搬下来，指数减一。}

\[
\frac{d}{dx}(x^n) = n x^{n-1}
\]

\textbf{第二}，积分的核心规律：\textbf{指数加一放分母。}

\[
\int x^n \, dx = \frac{x^{n+1}}{n+1} + C
\]

\textbf{第三}，加法型函数：各项分开处理，结果加起来。

\textbf{第四}，乘法型函数：对谁操作动谁，其他变量当常数乘在旁边。

\textbf{第五}，混合型函数：先按加减号拆分，每项用乘法规律，最后合并。

\textbf{第六}，在单位正方体上：

\begin{itemize}
\item 乘法型 $x^m y^n z^p$ 的积分 $= \frac{1}{(m+1)(n+1)(p+1)}$
\item 加法型 $x^m + y^n + z^p$ 的积分 $= \frac{1}{m+1} + \frac{1}{n+1} + \frac{1}{p+1}$
\end{itemize}

\textbf{第七}，不管指数多大、多复杂，\textbf{规律永远不变}。掌握核心口诀，任何多项式都能搞定。

\section*{附录：自测练习}

试试用今天学的方法解决这些问题（先自己算，再看答案）：

\textbf{练习 1：} 求 $f(x, y) = x^4 + 3y^2$ 的两个偏导数。

\paragraph{点击查看答案}
\[
\frac{\partial f}{\partial x} = 4x^3, \quad \frac{\partial f}{\partial y} = 6y
\]

\textbf{练习 2：} 求 $\displaystyle\int (x^3 - 2y^2) \, dx$。

\paragraph{点击查看答案}
\[
= \frac{x^4}{4} - 2y^2 x + C
\]

\textbf{练习 3：} 求 $\displaystyle\int_0^1 \int_0^1 x^2 y^4 \, dx \, dy$。

\paragraph{点击查看答案}
\[
= \frac{1}{3} \times \frac{1}{5} = \frac{1}{15}
\]

\textbf{练习 4：} 求 $f(x, y, z) = x^2 yz + xy^2 z + xyz^2$ 对 $x$ 的偏导数。

\paragraph{点击查看答案}
\[
\frac{\partial f}{\partial x} = 2xyz + y^2 z + yz^2 = yz(2x + y + z)
\]

\textbf{练习 5：} 在单位正方体上，$\displaystyle\iiint (x^2 + y^2 + z^2) \, dV = ?$

\paragraph{点击查看答案}
\[
= \frac{1}{3} + \frac{1}{3} + \frac{1}{3} = 1
\]

\chapter{闭合，实则不同：一重闭合微积分}

\section{开场白：一条首尾相连的路}

前面所有的章节里，我们走的都是“有头有尾”的路：

\begin{itemize}
\item 从 x = 0 走到 x = 2
\item 从点 $(0, 0)$ 走到点 $(1, 1)$
\item 从地面走到天花板
\end{itemize}

起点和终点是\textbf{不同的两个地方}。

但今天，我们要走一条\textbf{闭合的路}——从某个点出发，绕一圈，\textbf{回到原点}。

这时候，积分号上会套一个圆圈：$\oint$

这个圆圈就表示\textbf{闭合积分}——走一圈回到起点。

走一圈回到起点？那不就等于没走吗？积分结果会是 0 吗？

\textbf{有时候是 0，有时候不是。} 这就是今天要讲的内容。

\section{第一部分：什么是闭合路径}

\subsection{想象你在操场上跑步}

你站在操场的起跑线上（假设这是点 A）。

发令枪响，你开始跑。

你沿着跑道跑了一圈，最后\textbf{回到了起跑线}——还是点 A。

你的位置变了吗？\textbf{没有。} 你还是站在点 A。

但你跑了一圈吗？\textbf{跑了。} 你确实经过了整个跑道。

这就是\textbf{闭合路径}的核心：

\begin{quote}
\textbf{起点 = 终点}，但中间经过了一段旅程。
\end{quote}

\subsection{🎨 动手画一画：闭合路径}

\textbf{第一步：} 在纸上画一个圆（或者椭圆、正方形，任何首尾相连的封闭图形都行）。

\textbf{第二步：} 在圆上选一个点，标上 \textbf{A}。

\textbf{第三步：} 画一个箭头，从 A 出发，沿着圆走一圈，回到 A。

\textbf{第四步：} 在图旁边写上：\textbf{“闭合路径：从 A 出发，回到 A”}。

看看你的图——这就是闭合路径。

\subsection{闭合积分的符号}

普通积分：

\[
\int_a^b f(x) \, dx
\]

闭合积分：

\[
\oint f \, ds
\]

那个圈圈 $\oint$ 就是“闭合”的意思——\textbf{绕一圈回到起点}。

普通积分有起点 $a$ 和终点 $b$。闭合积分的起点和终点是同一个点，所以不用写 $a$ 和 $b$，只用一个圆圈表示。

\section{第二部分：最简单的闭合积分——位置的变化}

\subsection{问题：你跑一圈后，位置变了多少？}

答案很简单：\textbf{0}。

因为你回到了起点，位置没有变化。

用数学语言：

\[
\oint dx = 0
\]

这是什么意思？

\begin{itemize}
\item $dx$ 表示“位置的微小变化”
\item $\oint dx$ 表示“把一圈上所有的位置变化加起来”
\item 结果是 0，因为你回到了原点
\end{itemize}

\subsection{🎨 动手画一画：为什么位置变化是 0}

\textbf{第一步：} 画一个圆，在上面标出点 A（起点/终点）。

\textbf{第二步：} 从 A 出发，沿着圆走。在圆的右边画一小段，旁边写 \textbf{$+dx$}（往右走，位置增加）。

\textbf{第三步：} 继续走到圆的顶部、左边。在左边画一小段，旁边写 \textbf{$-dx$}（往左走，位置减少）。

\textbf{第四步：} 继续走回到 A。

\textbf{第五步：} 数一数：往右走的 $+dx$ 和往左走的 $-dx$，正好\textbf{抵消}了！

所以 $\oint dx = 0$。

\subsection{更严格的理解}

假设你走一个闭合的正方形路径：

\begin{itemize}
\item 从 $(0, 0)$ 走到 $(1, 0)$：x 增加了 1，dx 的积分 = $+1$
\item 从 $(1, 0)$ 走到 $(1, 1)$：x 没变，dx 的积分 = $0$
\item 从 $(1, 1)$ 走到 $(0, 1)$：x 减少了 1，dx 的积分 = $-1$
\item 从 $(0, 1)$ 走到 $(0, 0)$：x 没变，dx 的积分 = $0$
\end{itemize}

总和：$+1 + 0 + (-1) + 0 = \textbf{0}$

不管是圆形、正方形、还是任何奇形怪状的闭合路径，$\oint dx$ 永远等于 0。

这是因为 x 是一个\textbf{“状态量”}——它只取决于你在哪里，不取决于你怎么走过来的。走一圈回到原点，x 的净变化必然是 0。

\section{第三部分：不为 0 的闭合积分}

\subsection{位置变化 vs 路程}

$\oint dx = 0$，这没什么稀奇。但有些量，绕一圈之后\textbf{不是 0}。

比如说——\textbf{你走过的路程}。

\textbf{位置变化：} 你的终点和起点之间的差。

\begin{itemize}
\item 跑一圈：位置变化 = 0（回到原点）
\end{itemize}

\textbf{路程：} 你实际走过的距离。

\begin{itemize}
\item 跑一圈：路程 = 整个圈的周长（不是 0！）
\end{itemize}

\subsection{🎨 动手画一画：位置变化 vs 路程}

\textbf{第一步：} 画一个圆，周长是 10 米。标出起点 A。

\textbf{第二步：} 在圆外面写两个问题：

\begin{itemize}
\item 问题 1：从 A 走一圈回到 A，\textbf{位置变化}是多少？
\item 问题 2：从 A 走一圈回到 A，\textbf{路程}是多少？
\end{itemize}

\textbf{第三步：} 写出答案：

\begin{itemize}
\item 位置变化 = 0（起点终点相同）
\item 路程 = 10 米（走了一整圈）
\end{itemize}

\textbf{第四步：} 在图旁边写上：\textbf{“同样的路径，不同的量，不同的结果！”}

\subsection{数学表达}

\textbf{位置变化的闭合积分：}

\[
\oint dx = 0
\]

\textbf{路程的闭合积分：}

\[
\oint |dx| = \text{周长} \neq 0
\]

（$|dx|$ 表示取绝对值，所有的小段都算正的，不管方向）

\section{第四部分：一个关键概念——路径相关 vs 路径无关}

\subsection{两类不同的量}

\textbf{路径无关的量：} 结果只取决于起点和终点，和中间怎么走没关系。

\begin{itemize}
\item 例子：位置、海拔高度、温度
\item 特点：$\oint = 0$（走一圈回来，净变化是 0）
\end{itemize}

\textbf{路径相关的量：} 结果取决于你走过的具体路径。

\begin{itemize}
\item 例子：路程、做的功（某些情况下）、花的时间
\item 特点：$\oint \neq 0$（走一圈回来，结果不是 0）
\end{itemize}

\subsection{登山的例子}

想象你在爬山。

\textbf{海拔高度（路径无关）：}

\begin{itemize}
\item 你从山脚出发，爬到山顶，再下来回到山脚。
\item 你的海拔变化是多少？\textbf{0}。（回到了原来的高度）
\item 不管你走哪条路上山、哪条路下山，海拔变化都是 0。
\end{itemize}

\[
\oint d(\text{海拔}) = 0
\]

\textbf{走过的垂直距离（路径相关）：}

\begin{itemize}
\item 你上山时爬升了 500 米，下山时下降了 500 米。
\item 你走过的\textbf{总垂直距离}是 500 + 500 = 1000 米。
\item 如果你中间还翻越了几个小山包，这个数字会更大。
\end{itemize}

\[
\oint |d(\text{海拔})| = 1000 \text{ 米} \neq 0
\]

\subsection{🎨 动手画一画：路径无关 vs 路径相关}

画两幅图：

\textbf{图一：海拔变化（路径无关）}

画一座山，山脚标 A，山顶标 B。画一条路从 A 到 B 再回到 A。

在图旁边写：
\begin{itemize}
\item 上山：海拔 $+500$
\item 下山：海拔 $-500$
\item 总和：0
\end{itemize}

\textbf{图二：垂直距离（路径相关）}

同样的山，但这次标出：
\begin{itemize}
\item 上山：$+500$ 米（取绝对值）
\item 下山：$+500$ 米（取绝对值）
\item 总和：1000 米
\end{itemize}

\section{第五部分：一重闭合积分的计算}

\subsection{在一维上，什么叫“闭合”？}

在一维的数轴上，闭合路径就是：\textbf{从 a 走到 b，再从 b 走回 a。}

\[
\oint f(x) \, dx = \int_a^b f(x) \, dx + \int_b^a f(x) \, dx
\]

\subsection{关键公式}

\[
\int_b^a f(x) \, dx = -\int_a^b f(x) \, dx
\]

这是因为：交换上下限，积分变号。

所以：

\[
\oint f(x) \, dx = \int_a^b f(x) \, dx + \left(-\int_a^b f(x) \, dx\right) = 0
\]

\textbf{结论：对于普通的函数 $f(x)$，一维闭合积分永远等于 0。}

\subsection{验证一下}

\textbf{例子 1：} $f(x) = x$，从 0 到 2 再回到 0

\[
\int_0^2 x \, dx = \left[\frac{x^2}{2}\right]_0^2 = 2
\]

\[
\int_2^0 x \, dx = \left[\frac{x^2}{2}\right]_2^0 = 0 - 2 = -2
\]

\[
\oint x \, dx = 2 + (-2) = 0 \quad \checkmark
\]

\textbf{例子 2：} $f(x) = x^2$，从 1 到 3 再回到 1

\[
\int_1^3 x^2 \, dx = \left[\frac{x^3}{3}\right]_1^3 = 9 - \frac{1}{3} = \frac{26}{3}
\]

\[
\int_3^1 x^2 \, dx = -\frac{26}{3}
\]

\[
\oint x^2 \, dx = \frac{26}{3} + \left(-\frac{26}{3}\right) = 0 \quad \checkmark
\]

\textbf{例子 3：} $f(x) = \sin x$，从 0 到 $\pi$ 再回到 0

\[
\int_0^\pi \sin x \, dx = [-\cos x]_0^\pi = -(-1) - (-1) = 2
\]

\[
\int_\pi^0 \sin x \, dx = -2
\]

\[
\oint \sin x \, dx = 2 + (-2) = 0 \quad \checkmark
\]

\subsection{为什么永远是 0？}

在一维上，从 a 到 b 只有一条路（数轴上的直线）。

从 a 到 b 积分得到某个值，从 b 到 a 就是走回去，积分得到负的那个值。

一正一负，恰好抵消。

这就是为什么一维闭合积分永远是 0——\textbf{没有其他路可走}。

\section{第六部分：那什么时候闭合积分不是 0？}

\subsection{答案：当路径有多种选择时}

在一维数轴上，从 a 到 b 只有一条路，所以闭合积分必然是 0。

但在\textbf{二维平面}或\textbf{三维空间}里，从 a 到 b 可以有\textbf{无数条路}！

比如，从 $(0, 0)$ 到 $(1, 1)$：

\begin{itemize}
\item 可以先往右走，再往上走
\item 可以先往上走，再往右走
\item 可以走一条斜线
\item 可以走一个弯弯曲曲的 S 形
\end{itemize}

当有多条路可走时，\textbf{去和回可以走不同的路}，积分结果就可能不抵消。

这是\textbf{二维闭合积分}（曲线积分）的内容，下一章会详细讲。

\subsection{预告：二维闭合积分}

在二维平面上，绕一个闭合曲线走一圈：

\[
\oint_C (P \, dx + Q \, dy)
\]

这个积分\textbf{不一定是 0}！

什么时候是 0，什么时候不是 0，有一个著名的定理——\textbf{格林公式}——会告诉我们答案。

这是下一章的内容。

\section{第七部分：一维闭合积分的意义}

\subsection{既然永远是 0，为什么还要学？}

一维闭合积分永远是 0，那学它有什么用？

\textbf{第一，它帮助我们理解闭合积分的概念。} 一维是最简单的情况，先搞懂一维，才能理解二维、三维。

\textbf{第二，它告诉我们一个重要的事实：能写成 $dF$ 形式的东西，闭合积分必然是 0。}

什么意思？如果存在一个函数 $F(x)$，使得 $f(x) dx = dF$，那么：

\[
\oint f(x) \, dx = \oint dF = F(\text{终点}) - F(\text{起点}) = F(A) - F(A) = 0
\]

这种情况下，$f(x)$ 叫做\textbf{保守场}的一维版本。

\subsection{物理意义：保守力}

在物理学中：

\begin{itemize}
\item \textbf{保守力}（如重力、弹簧力）：做功只取决于起点和终点的位置，闭合路径上做功为 0。
\item \textbf{非保守力}（如摩擦力）：做功取决于路径，闭合路径上做功不为 0。
\end{itemize}

一维的情况下，所有力都表现得像“保守力”，因为只有一条路可走。

到了二维、三维，区别就出来了。

\section{第八部分：规律总结}

\subsection{一维闭合积分}

\[
\oint f(x) \, dx = 0
\]

\textbf{永远是 0}，因为：

\begin{itemize}
\item 一维上从 a 到 b 只有一条路
\item 从 a 到 b 的积分和从 b 到 a 的积分恰好相反
\item 两者相加等于 0
\end{itemize}

\subsection{闭合积分的符号}

\begin{center}
\begin{tabular}{c|c}
符号 & 含义 \\
\hline
$\displaystyle\int_a^b$ & 普通积分，从 a 到 b \\
$\displaystyle\oint$ & 闭合积分，绕一圈回到起点
\end{tabular}
\end{center}

\subsection{路径相关 vs 路径无关}

\begin{center}
\begin{tabular}{c|c|c}
类型 & 特点 & $\oint$ 的结果 \\
\hline
路径无关 & 只取决于起点和终点 & = 0 \\
路径相关 & 取决于具体走过的路径 & $\neq 0$
\end{tabular}
\end{center}

\subsection{各维度的情况}

\begin{center}
\begin{tabular}{c|c|c}
维度 & $\oint dx$ 是否一定为 0 & 原因 \\
\hline
一维 & \textbf{是} & 只有一条路，去和回必然抵消 \\
二维 & \textbf{不一定} & 有多条路，去和回可以走不同的路 \\
三维 & \textbf{不一定} & 同上
\end{tabular}
\end{center}

\section{本章的核心图}

\subsection{🎨 动手画一画}

画一张图，包含以下内容：

\textbf{左边：一维}

画一条数轴，标出 a 和 b。画一个箭头从 a 到 b（上面），再画一个箭头从 b 到 a（下面）。

旁边写：
\begin{quote}
\ttfamily
\(\oint f(x)\,dx = \int_{a\to b} + \int_{b\to a}\)\\
\(\qquad = \int_{a\to b} + \left(-\int_{a\to b}\right)\)\\
\(\qquad = 0\)
\end{quote}

\textbf{右边：二维（预告）}

画一个圆形闭合曲线，箭头沿着圆走一圈。

旁边写：
\begin{quote}
\ttfamily
\(\oint (P\,dx + Q\,dy) = ?\)\\
可能是 0，可能不是 0\\
取决于 $P$ 和 $Q$ 的性质\\
（下一章揭晓）
\end{quote}

\section{本章公式汇总}

\begin{center}
\begin{tabular}{c|c}
公式 & 含义 \\
\hline
$\displaystyle\int_b^a f(x) \, dx = -\int_a^b f(x) \, dx$ & 交换上下限，积分变号 \\
$\displaystyle\oint f(x) \, dx = 0$ & 一维闭合积分等于 0 \\
$\displaystyle\oint dx = 0$ & 位置变化的闭合积分 \\
$\displaystyle\oint |dx| = \text{周长}$ & 路程的闭合积分（不为 0）
\end{tabular}
\end{center}

\section{本章小结}

\textbf{第一}，闭合路径就是\textbf{起点 = 终点}的路径，绕一圈回到原点。

\textbf{第二}，闭合积分用 $\oint$ 表示，意思是沿着闭合路径积分。

\textbf{第三}，在一维上，闭合积分\textbf{永远是 0}。因为从 a 到 b 只有一条路，去和回必然抵消。

\textbf{第四}，\textbf{路径无关}的量（如位置、海拔）：闭合积分 = 0。

\textbf{第五}，\textbf{路径相关}的量（如路程、做的功）：闭合积分 $\neq 0$。

\textbf{第六}，一维是最简单的情况。到了二维、三维，闭合积分就不一定是 0 了——这是下一章的内容。

\textbf{第七}，理解一维闭合积分是理解更高维闭合积分的\textbf{基础}。先搞懂简单的，才能理解复杂的。

一维的世界很简单——只有一条路，去了就得原路返回。

但二维、三维的世界不一样——可以走不同的路，去和回可以绕道。

那时候，闭合积分就会展现出它真正的威力。

\textbf{闭合，实则不同。} 同样的终点，不同的旅程，可能有不同的结果。

\chapter{双重闭合积分：在气球表面收集}

\section{开场白：从一圈到一整个壳}

上一章我们学过\textbf{一重闭合积分}：

沿着一个\textbf{圈}走一圈，从起点出发，绑着圈走，最后回到起点。

今天升级了。

不是沿着一个\textbf{圈}，而是在一个\textbf{壳}的表面上收集。

什么是壳？就是一个完全封闭的表面——像气球的皮、鸡蛋的壳、盒子的六个面。

这就是\textbf{双重闭合积分}。

\section{第一节：什么是闭合的表面}

\subsection{闭合 vs 不闭合}

\textbf{闭合的表面}：完全封闭，没有洞，没有开口。

\textbf{不闭合的表面}：有边缘，有开口。

\subsection{🎨 动手画一画：闭合 vs 不闭合}

画四张图：

\textbf{图一：气球（闭合）}

画一个圆，涂上颜色。

旁边写：“完全封闭，没有洞。往里面吹气，气跑不出来。✓ 闭合”

\textbf{图二：球壳（闭合）}

画一个圆，代表乒乓球的壳。

旁边写：“一个完整的壳，把里面和外面完全隔开。✓ 闭合”

\textbf{图三：碗（不闭合）}

画半个圆，像一个碗。

旁边写：“上面有开口。往里面倒水，水会从上面溢出来。✗ 不闭合”

\textbf{图四：一张纸（不闭合）}

画一个长方形，代表一张纸。

旁边写：“有四条边。不能把东西包起来。✗ 不闭合”

\subsection{再举几个例子}

\begin{center}
\begin{tabular}{c|c|c}
物体 & 闭合吗？ & 为什么 \\
\hline
足球的表面 & ✓ 闭合 & 完全封闭，气跑不出来 \\
纸盒子（六个面都封好） & ✓ 闭合 & 六个面把里面完全包住 \\
易拉罐（没开口） & ✓ 闭合 & 完全密封 \\
易拉罐（开了口） & ✗ 不闭合 & 上面有个洞 \\
帐篷 & ✗ 不闭合 & 有门，有缝隙 \\
鸡蛋壳（完整的） & ✓ 闭合 & 完全包住蛋黄蛋白 \\
鸡蛋壳（打碎了一半） & ✗ 不闭合 & 有开口
\end{tabular}
\end{center}

\subsection{一句话总结}

\begin{quote}
\textbf{闭合的表面 = 能把里面和外面完全隔开的表面}
\end{quote}

\section{第二节：在表面上收集}

\subsection{回忆一下之前学的}

\textbf{一重积分}：沿着一条\textbf{线}收集。

\[
\int_0^2 x \, dx
\]

从 $x = 0$ 走到 $x = 2$，沿着这条线收集。

\textbf{二重积分}：在一块\textbf{平面区域}上收集。

\[
\iint_D xy \, dA
\]

在区域 $D$ 上，一小块一小块地收集，全部加起来。

\subsection{现在：在一个弯曲的表面上收集}

想象一个气球。

气球的表面不是平的——它是弯曲的，是一个球面。

如果气球表面上每个位置都有一些灰尘，而且不同位置的灰尘量不同，\textbf{整个气球表面一共有多少灰尘？}

这就需要\textbf{曲面积分}——在一个弯曲的表面上收集。

如果这个表面恰好是\textbf{闭合的}（像气球一样完全封闭），那就是\textbf{闭合曲面积分}，也叫\textbf{双重闭合积分}。

\subsection{符号}

普通的二重积分：

\[
\iint_D f(x, y) \, dA
\]

在弯曲表面上的积分：

\[
\iint_S f \, dS
\]

在\textbf{闭合}表面上的积分：

\[
\oiint_S f \, dS
\]

看到那个套在积分号上的小圆圈 ○ 了吗？那就是“闭合”的标志。

\[
\oiint
\]

两个积分号（因为表面是二维的），加上一个小圆圈（因为表面是闭合的）。

\subsection{🎨 动手画一画：三种积分的区别}

画三张图：

\textbf{图一：一重积分（沿着线）}

画一条弯弯曲曲的线，从点 A 到点 B。

旁边写：“沿着一条线收集。$\int_A^B$”

\textbf{图二：二重积分（在平面区域上）}

画一个平面上的区域（比如一个长方形或三角形），涂上颜色。

旁边写：“在一块平的区域上收集。$\iint_D$”

\textbf{图三：双重闭合积分（在闭合表面上）}

画一个气球（一个圆，表示球面）。

旁边写：“在一个封闭的壳上收集。$\oiint_S$”

\section{第三节：一个简单的例子}

\subsection{问题}

一个气球，表面积是 100 平方厘米。

气球表面每个地方都均匀地撒了灰尘，每平方厘米有 3 克灰尘。

问：整个气球表面一共有多少灰尘？

\subsection{答案}

不用积分也能算出来：

\[
\text{总灰尘} = 3 \times 100 = 300 \text{ 克}
\]

\subsection{用积分写出来}

\[
\oiint_S 3 \, dS = 3 \times \oiint_S 1 \, dS = 3 \times \text{表面积} = 3 \times 100 = 300
\]

其中 $\oiint_S 1 \, dS$ 就是表面积（每一小块的面积加起来）。

\subsection{🎨 动手画一画：气球上的灰尘}

\textbf{第一步：} 画一个圆（代表气球）。

\textbf{第二步：} 在圆的表面上画很多小点点，代表灰尘。让灰尘分布均匀。

\textbf{第三步：} 在旁边写：“每平方厘米 3 克，表面积 100 平方厘米，总共 300 克。”

\section{第四节：如果灰尘不均匀呢？}

\subsection{问题}

同样的气球，表面积 100 平方厘米。

但这次灰尘不是均匀的——气球的顶部灰尘多，底部灰尘少。

不同位置的灰尘量用一个函数描述：$f(x, y, z)$——在位置 $(x, y, z)$ 处，每平方厘米有 $f(x, y, z)$ 克灰尘。

问：整个气球表面一共有多少灰尘？

\subsection{方法}

把气球表面切成很多很多\textbf{小片}。

每一小片的面积是 $dS$（无限小）。

每一小片上的灰尘 $\approx f(x, y, z) \times dS$。

把所有小片的灰尘加起来：

\[
\text{总灰尘} = \oiint_S f(x, y, z) \, dS
\]

\subsection{🎨 动手画一画：不均匀的灰尘}

\textbf{第一步：} 画一个圆（代表气球）。

\textbf{第二步：} 在圆的上部画很多密集的点点（灰尘多）。

\textbf{第三步：} 在圆的下部画很少的点点（灰尘少）。

\textbf{第四步：} 在旁边写：“上面灰尘多，下面灰尘少。不能简单地用‘每平方厘米 × 面积’来算了，必须用积分。”

\section{第五节：闭合 vs 不闭合——有什么区别？}

\subsection{数学上的区别}

闭合表面用 $\oiint$（带圆圈），不闭合表面用 $\iint$（不带圆圈）。

\subsection{真正的区别}

\textbf{不闭合的表面}有边界——像一张纸有四条边，像一个碗有一圈边缘。

\textbf{闭合的表面}没有边界——气球的表面上，你找不到任何一条“边”。你可以在气球表面上一直走，永远不会走到“尽头”。

\subsection{🎨 动手画一画：有边 vs 无边}

\textbf{图一：一张纸（有边界）}

画一个长方形。用红色笔把四条边描粗。

旁边写：“四条边就是边界。走到边上就到头了。”

\textbf{图二：气球表面（无边界）}

画一个圆。不要描边——整个圆都是表面，没有“边”。

旁边写：“没有边界。在表面上可以一直走，永远不会走到头。”

\subsection{类比一维的情况}

\begin{center}
\begin{tabular}{c|c}
一维 & 二维 \\
\hline
一条线段（有两个端点） & 一张纸（有边界） \\
一个圆圈（没有端点，首尾相连） & 一个气球表面（没有边界，完全封闭） \\
沿线段积分：$\int_a^b$ & 在平面区域积分：$\iint_D$ \\
沿圆圈积分：$\oint$（一重闭合） & 在闭合表面积分：$\oiint$（双重闭合）
\end{tabular}
\end{center}

\section{第六节：贡献}

\subsection{每一小块的贡献}

和之前学的一样，每一小块对总结果的贡献是：

\[
\text{贡献} = f(\text{这个位置}) \times dS(\text{这一小块的面积})
\]

$f$ 大的地方，贡献大。

$f$ 小的地方，贡献小。

$f = 0$ 的地方，没有贡献。

$f < 0$ 的地方，贡献是负的。

\subsection{🎨 动手画一画：贡献的大小}

\textbf{第一步：} 画一个气球（圆）。

\textbf{第二步：} 把气球分成四块（像切西瓜一样，画两条垂直的线）。

\textbf{第三步：} 在四块上分别写上 $f$ 的值：

\begin{itemize}
\item 左上块：$f = 5$（贡献大）
\item 右上块：$f = 3$（贡献中）
\item 左下块：$f = 1$（贡献小）
\item 右下块：$f = 0$（没贡献）
\end{itemize}

\textbf{第四步：} 把贡献大的那块涂深色，贡献小的涂浅色。

\textbf{第五步：} 在旁边写：“$f$ 大的地方贡献多，$f$ 小的地方贡献少。”

\section{第七节：例子——盒子的六个面}

\subsection{问题}

一个正方体盒子，边长是 1。

盒子表面的六个面，每个面上的灰尘量不同：

\begin{itemize}
\item 顶面：每平方厘米 10 克
\item 底面：每平方厘米 2 克
\item 前面：每平方厘米 5 克
\item 后面：每平方厘米 5 克
\item 左面：每平方厘米 3 克
\item 右面：每平方厘米 7 克
\end{itemize}

问：整个盒子表面一共有多少灰尘？

\subsection{解答}

每个面的面积 = $1 \times 1 = 1$ 平方厘米。

每个面的灰尘量 = 灰尘密度 × 面积。

\begin{itemize}
\item 顶面：$10 \times 1 = 10$ 克
\item 底面：$2 \times 1 = 2$ 克
\item 前面：$5 \times 1 = 5$ 克
\item 后面：$5 \times 1 = 5$ 克
\item 左面：$3 \times 1 = 3$ 克
\item 右面：$7 \times 1 = 7$ 克
\end{itemize}

总共：$10 + 2 + 5 + 5 + 3 + 7 = $ \textbf{32 克}

\subsection{用积分写出来}

\[
\oiint_S f \, dS = 10 \times 1 + 2 \times 1 + 5 \times 1 + 5 \times 1 + 3 \times 1 + 7 \times 1 = 32
\]

盒子的表面是闭合的（六个面完全封住了里面），所以用 $\oiint$。

\subsection{🎨 动手画一画：盒子的六个面}

\textbf{第一步：} 画一个立体的正方体。

\textbf{第二步：} 在每个面上写上灰尘密度：顶面 10，底面 2，前面 5，后面 5，左面 3，右面 7。

\textbf{第三步：} 在旁边写：“六个面加起来，总灰尘 = 32 克。”

\section{第八节：动手画——正方体}

上面我们用盒子算了灰尘，现在用坐标把正方体画出来。画准了，以后算积分心里更有底。

\subsection{第一步：画坐标轴}

在纸的中间偏左下，画一个点，这是原点 \textbf{(0, 0, 0)}。

从原点往右画一条箭头，旁边写 \textbf{x}。

从原点往左上方斜着画一条箭头，旁边写 \textbf{y}。

从原点往正上方画一条箭头，旁边写 \textbf{z}。

\subsection{第二步：标出八个顶点}

正方体边长是 1，一个角放在原点。

按照“往右走 x 步、往左上斜走 y 步、往上走 z 步”的方法，把八个顶点标出来：

\begin{center}
\begin{tabular}{c|c|c}
顶点 & 坐标 & 怎么走到那里 \\
\hline
A & $(0, 0, 0)$ & 原点，不用走 \\
B & $(1, 0, 0)$ & 沿 x 轴往右走 1 步 \\
C & $(1, 1, 0)$ & 往右 1 步，再往前 1 步 \\
D & $(0, 1, 0)$ & 沿 y 轴往前走 1 步 \\
E & $(0, 0, 1)$ & 沿 z 轴往上走 1 步 \\
F & $(1, 0, 1)$ & 往右 1 步，再往上 1 步 \\
G & $(1, 1, 1)$ & 三个方向各走 1 步 \\
H & $(0, 1, 1)$ & 往前 1 步，再往上 1 步
\end{tabular}
\end{center}

在每个点旁边写上字母和坐标。

\subsection{第三步：连线，画出十二条边}

\textbf{看得见的边，用实线连：}

\begin{itemize}
\item A $\to$ B，B $\to$ C，A $\to$ D，A $\to$ E
\item B $\to$ F，E $\to$ F，E $\to$ H，F $\to$ G
\end{itemize}

\textbf{看不见的边（被挡住的），用虚线连：}

\begin{itemize}
\item D $\to$ C，C $\to$ G，D $\to$ H，H $\to$ G
\end{itemize}

\subsection{第四步：标出六个面的法向量}

\textbf{法向量是什么？}

就是从表面上每个点，垂直于表面、朝外指的一个箭头。

想象一下：你站在正方体的某个面上（假装你很小很小），头顶朝外。你头顶指的方向，就是那个面的法向量方向。

正方体有六个面，每个面都有一个朝外的法向量：

\textbf{底面}（ABCD，$z = 0$ 的那一面）：

法向量朝\textbf{正下方}。在底面中心画一个朝下的箭头，旁边写 \textbf{(0, 0, $-1$)}。

\textbf{顶面}（EFGH，$z = 1$ 的那一面）：

法向量朝\textbf{正上方}。在顶面中心画一个朝上的箭头，旁边写 \textbf{(0, 0, $+1$)}。

\textbf{前面}（ABFE，$y = 0$ 的那一面）：

法向量朝\textbf{正前方}（朝着你）。在前面中心画一个朝前的箭头，旁边写 \textbf{(0, $-1$, 0)}。

\textbf{后面}（DCGH，$y = 1$ 的那一面）：

法向量朝\textbf{正后方}。用虚线箭头画，旁边写 \textbf{(0, $+1$, 0)}。

\textbf{右面}（BCGF，$x = 1$ 的那一面）：

法向量朝\textbf{正右方}。在右面中心画一个朝右的箭头，旁边写 \textbf{($+1$, 0, 0)}。

\textbf{左面}（ADHE，$x = 0$ 的那一面）：

法向量朝\textbf{正左方}。用虚线箭头画，旁边写 \textbf{($-1$, 0, 0)}。

\subsection{第五步：标上 $f$ 的值，写出积分}

在每个面旁边写上灰尘密度 $f$：

\begin{center}
\begin{tabular}{c|c|c|c}
面 & $f$ 值 & 面积 & 贡献 \\
\hline
底面 & 2 & 1 & 2 \\
顶面 & 8 & 1 & 8 \\
前面 & 3 & 1 & 3 \\
后面 & 3 & 1 & 3 \\
右面 & 5 & 1 & 5 \\
左面 & 5 & 1 & 5
\end{tabular}
\end{center}

在图的下方写上：

\[
\oiint_S f \, dS = 2 + 8 + 3 + 3 + 5 + 5 = 26
\]

\section{第九节：动手画——球面}

\subsection{第一步：画坐标轴和原点}

和正方体一样，先画三个坐标轴。

在原点标上 \textbf{O(0, 0, 0)}。

\subsection{第二步：标出六个特殊点}

这六个点在坐标轴上，分别是球面和三条轴的交点：

\begin{center}
\begin{tabular}{c|c|c}
点 & 坐标 & 在哪里 \\
\hline
$P_1$ & $(1, 0, 0)$ & x 轴正方向，最右边 \\
$P_2$ & $(-1, 0, 0)$ & x 轴负方向，最左边 \\
$P_3$ & $(0, 1, 0)$ & y 轴正方向，最前面 \\
$P_4$ & $(0, -1, 0)$ & y 轴负方向，最后面 \\
$P_5$ & $(0, 0, 1)$ & z 轴正方向，最顶部 \\
$P_6$ & $(0, 0, -1)$ & z 轴负方向，最底部
\end{tabular}
\end{center}

在坐标轴上，把这六个点都标出来，每个点距离原点 1 个单位。

\subsection{第三步：画球面的轮廓}

\textbf{先画大圆：} 以原点为圆心，画一个大圆。这是球面从正面看到的轮廓。

\textbf{再画赤道：} 在圆的中间，画一条\textbf{横着的椭圆弧}，连接 $P_3$ 和 $P_4$（看得见的半段用实线，看不见的半段用虚线）。这是球面的“赤道”。

\textbf{再画经线：} 在圆的中间，画一条\textbf{竖着的半圆弧}（只画右边看得见的部分），连接 $P_5$ 和 $P_6$ 经过 $P_1$。

六个特殊点都在球面上，球面看起来就立体了。

\subsection{第四步：画法向量}

球面的法向量有一个漂亮的规律：

\begin{quote}
\textbf{球面上每一点的法向量，就是从原点 O 指向那个点的方向。}
\end{quote}

（还是那个比喻：你站在球面上，头顶朝外——在球面上，“朝外”刚好就是远离球心的方向。）

在六个特殊点上，各画一个法向量箭头：

\begin{itemize}
\item \textbf{$P_1(1, 0, 0)$：} 从 $P_1$ 往右画箭头。旁边写 \textbf{(1, 0, 0)}。
\item \textbf{$P_2(-1, 0, 0)$：} 从 $P_2$ 往左画箭头。旁边写 \textbf{(-1, 0, 0)}。
\item \textbf{$P_3(0, 1, 0)$：} 从 $P_3$ 往左上方斜着画箭头（y 方向）。旁边写 \textbf{(0, 1, 0)}。
\item \textbf{$P_4(0, -1, 0)$：} 从 $P_4$ 往右下方斜着画箭头。旁边写 \textbf{(0, -1, 0)}。
\item \textbf{$P_5(0, 0, 1)$：} 从 $P_5$ 往正上方画箭头。旁边写 \textbf{(0, 0, 1)}。
\item \textbf{$P_6(0, 0, -1)$：} 从 $P_6$ 往正下方画箭头。旁边写 \textbf{(0, 0, -1)}。
\end{itemize}

在图旁边写上：\textbf{“每一点的法向量 = 从 O 指向那个点，方向就是那个点的坐标。”}

\subsection{第五步：标上球面方程}

在图旁边写上球面的方程：

\[
x^2 + y^2 + z^2 = 1
\]

这个方程的意思是：球面上每一个点 $(x, y, z)$，到原点的距离都等于 1。

再写上积分符号：

\[
\oiint_S f \, dS = \text{把球面上所有小块的 } f \times dS \text{ 加起来}
\]

\section{第十节：两张图放在一起比较}

在一张纸上，左边画正方体，右边画球面。

画好之后，在两张图下面填写这张对比表：

\begin{center}
\begin{tabular}{c|c|c}
 & 正方体 & 球面 \\
\hline
关键点 & 8 个顶点（坐标是 0 或 1） & 6 个特殊点（坐标轴上） \\
边或轮廓 & 12 条直线边 & 大圆 + 椭圆弧 \\
法向量 & 六个方向，每面一个 & 从 O 往外辐射，每点不同 \\
法向量方向 & 六个固定方向（$\pm x$、$\pm y$、$\pm z$） & 就是那个点的坐标 $(x, y, z)$ \\
闭合吗 & ✓ & ✓ \\
积分符号 & $\oiint$ & $\oiint$
\end{tabular}
\end{center}

\section{总结}

\subsection{双重闭合积分是什么}

在一个\textbf{完全封闭的表面}上收集东西。

表面可以是：气球、球壳、盒子、鸡蛋壳……任何没有开口的壳。

\subsection{符号}

\[
\oiint_S f \, dS
\]

\begin{itemize}
\item $\oiint$：两个积分号 + 一个小圆圈
\item 两个积分号：因为表面是二维的
\item 小圆圈：因为表面是闭合的（没有边界）
\item $S$：表面
\item $f$：每个位置的值
\item $dS$：一小块表面的面积
\end{itemize}

\subsection{计算方法}

把表面切成很多小片，每一小片的贡献是 $f \times dS$，全部加起来。

\subsection{贡献}

\begin{itemize}
\item $f$ 大 $\rightarrow$ 贡献大
\item $f$ 小 $\rightarrow$ 贡献小
\item $f = 0$ $\rightarrow$ 没贡献
\item $f < 0$ $\rightarrow$ 负贡献
\end{itemize}

\subsection{和其他积分的对比}

\begin{center}
\begin{tabular}{c|c|c|c}
类型 & 在什么上收集 & 符号 & 闭合吗 \\
\hline
一重积分 & 一条线 & $\int$ & 不闭合 \\
一重闭合积分 & 一个圈 & $\oint$ & 闭合 \\
二重积分 & 一块平面区域 & $\iint$ & 不闭合 \\
双重闭合积分 & 一个封闭的壳 & $\oiint$ & 闭合
\end{tabular}
\end{center}

\section*{速查卡片}

\begin{quote}
\ttfamily
双重闭合积分速查\\
\\
是什么：在一个封闭表面上收集东西\\
\\
封闭表面：气球、盒子、球壳……没有开口的壳\\
\\
符号：∯ 或 ∮∮\\
\quad 两个积分号 + 小圆圈\\
\\
计算：把表面切成小片\\
\quad 每片贡献 = f × dS\\
\quad 全部加起来\\
\\
特点：闭合表面没有边界\\
\quad 可以在表面上一直走，永远走不到头
\end{quote}

\section*{画图检查清单}

每次画完，用这个清单检查一遍：

\textbf{正方体：}

\begin{itemize}
\item [ ] 坐标轴画好了（x、y、z 三个方向）
\item [ ] 八个顶点都标上了坐标
\item [ ] 看得见的边用实线，看不见的边用虚线
\item [ ] 六个面都有法向量（朝外的箭头）
\item [ ] 每个面旁边写了 $f$ 的值
\item [ ] 图下面写了积分公式和答案
\end{itemize}

\textbf{球面：}

\begin{itemize}
\item [ ] 坐标轴画好了
\item [ ] 原点 O 标好了
\item [ ] 六个特殊点都标上了坐标
\item [ ] 大圆画好了（实线）
\item [ ] 椭圆弧画好了（看得见实线，看不见虚线）
\item [ ] 六个特殊点上各画了一个法向量
\item [ ] 球面方程 $x^2 + y^2 + z^2 = 1$ 写好了
\item [ ] 图下面写了积分公式
\end{itemize}

一个小建议：每次遇到双重闭合积分，先画图，再计算。图画对了，题目就做对了一半。

\chapter{一步一步算：$\oint (x+y)\,dx$}

\section{想象一个场景}

地上有一个正方形。

四个角分别在：$(0,0)$，$(1,0)$，$(1,1)$，$(0,1)$。

正方形的每一个点上，都\textbf{用粉笔写了一个数}。

写的是什么？就是那个点的 $x + y$。

比如你站在 $(0.3,\; 0.7)$ 这个点，低头一看，地上写着 $0.3 + 0.7 = 1$。

站在 $(1,\; 1)$，地上写着 $1 + 1 = 2$。

站在 $(0,\; 0)$，地上写着 $0 + 0 = 0$。

\textbf{每个点脚底下都有一个数。位置不同，数就不同。}

\section{你要做什么？}

沿着正方形的边，\textbf{逆时针走一圈}。

每走一小步 $dx$，就低头看看脚下写的数是多少，乘以你这一步的长度 $dx$。

走完一圈，把所有这些乘积\textbf{加起来}。

但有一个规矩：

\textbf{只有你往左右方向走的时候才算数。}

往上下方向走的那些步，$dx = 0$，不算。

这就是 $\oint (x+y)\,dx$ 的意思。

\section{现在一段一段走}

\subsection{第一段：底边}

从 $(0,0)$ 走到 $(1,0)$。\textbf{往右。}

你走在最下面那条线上。

脚始终踩在 $y = 0$ 的地方，没上也没下。

只有 $x$ 在动，从 $0$ 慢慢变到 $1$。

脚下写的数是 $x + y = x + 0 = x$。

所以这一段，你每一步看到的数就是当时的 $x$。

\[
\int_0^1 x \, dx = \frac{x^2}{2}\bigg|_0^1 = \frac{1}{2}
\]

\begin{quote}
\textbf{底边攒了：$+\dfrac{1}{2}$}
\end{quote}

\subsection{第二段：右边}

从 $(1,0)$ 走到 $(1,1)$。\textbf{往上。}

你贴着最右边那条线往上走。

$x$ 一直是 $1$，没左也没右。

\textbf{$x$ 没动，$dx = 0$。}

脚下的数不管多大，乘以 $0$ 都是 $0$。

这一段白走了。

\begin{quote}
\textbf{右边攒了：$0$}
\end{quote}

\subsection{第三段：顶边}

从 $(1,1)$ 走到 $(0,1)$。\textbf{往左。}

你走在最上面那条线上。

脚始终踩在 $y = 1$ 的地方。

$x$ 在动，但方向反了——\textbf{从 $1$ 走回 $0$}。

脚下写的数是 $x + y = x + 1$。

因为你在顶边，$y = 1$，所以每个数都比底边的大了 $1$。

$x$ 从 $1$ 走到 $0$（往左走，方向反了）：

\[
\int_1^0 (x+1)\,dx
\]

先找原函数：$\dfrac{x^2}{2} + x$

代入：

\[
\left(\frac{0}{2} + 0\right) - \left(\frac{1}{2} + 1\right) = 0 - \frac{3}{2} = -\frac{3}{2}
\]

\textbf{负号哪来的？你是往左走的，$x$ 在减小，方向反了，符号就反了。}

\begin{quote}
\textbf{顶边攒了：$-\dfrac{3}{2}$}
\end{quote}

\subsection{第四段：左边}

从 $(0,1)$ 走到 $(0,0)$。\textbf{往下。}

你贴着最左边那条线往下走。

$x$ 一直是 $0$，没左也没右。

\textbf{$dx = 0$}。又白走了。

\begin{quote}
\textbf{左边攒了：$0$}
\end{quote}

\section{四段加起来}

\begin{center}
\begin{tabular}{c|c|c}
哪条边 & 怎么走 & 攒了多少 \\
\hline
底边 & 往右 & $+\frac{1}{2}$ \\
右边 & 往上 & $0$ \\
顶边 & 往左 & $-\frac{3}{2}$ \\
左边 & 往下 & $0$
\end{tabular}
\end{center}

\[
\frac{1}{2} + 0 - \frac{3}{2} + 0 = -1
\]

\[
\boxed{\oint (x+y)\,dx = -1}
\]

\section{等等，之前算过 $\oint y\,dx = -1$}

现在 $\oint (x+y)\,dx$ 也是 $-1$？

多了一个 $x$，答案居然没变？

拆开看看：

\[
\oint (x+y)\,dx = \oint x\,dx + \oint y\,dx
\]

$\oint y\,dx = -1$，这个你已经知道了。

那 $\oint x\,dx$ 呢？

\textbf{底边}：脚下的数是 $x$，从左往右走。$\int_0^1 x\,dx = +\dfrac{1}{2}$

\textbf{顶边}：脚下的数\textbf{也是} $x$，从右往左走。$\int_1^0 x\,dx = -\dfrac{1}{2}$

\textbf{右边、左边}：$dx = 0$，都是 $0$。

\[
\oint x\,dx = \frac{1}{2} - \frac{1}{2} = 0
\]

\textbf{底边加了 $\frac{1}{2}$，顶边减了 $\frac{1}{2}$。刚好抵消。}

\section{为什么 $x$ 的部分会抵消？}

底边上，脚下写的数是 $x$。

顶边上，脚下写的数\textbf{也是} $x$。

一模一样的数！

只是底边往右走，顶边往左走——\textbf{方向反了，一正一负，加起来刚好是零。}

\textbf{来回走了一趟，脚下踩的数一模一样，什么也没留下。}

\section{那为什么 $y$ 的部分没抵消？}

底边上，$y = 0$。脚下写的是 $0$。你几乎什么都没踩到。

顶边上，$y = 1$。脚下写的是 $1$。你每一步都踩着一个大数。

\textbf{底边和顶边，脚下的数根本不一样！}

大的减小的，减不干净，留下了东西。

\begin{quote}
\textbf{规律}

如果脚下的数\textbf{只跟 $x$ 有关}（跟 $y$ 无关），那底边和顶边踩的数一模一样，来回一走就抵消了。

如果脚下的数\textbf{跟 $y$ 有关}，那底边和顶边踩的数不一样——底边矮，数小；顶边高，数大。来回一走，抵消不干净，就留下了东西。

\textbf{脚下的数会随位置变，才可能留下痕迹。不变的，绕一圈就消失了。}
\end{quote}

\chapter{一步一步算：$\oiint (x+y)\,dS$}

\section{先搞清楚：这是在算什么？}

上一章，你沿着一个\textbf{圈}走了一圈。

圈是一条线，一维的。

现在升级了。

不是沿着一个\textbf{圈}，而是在一个\textbf{盒子}的表面上收集。

一个立方体，六个面。

你要把六个面\textbf{全部摸一遍}。

每摸到一小块面积 $dS$，就看看这个位置的 $(x+y)$ 是多少，乘以这一小块的面积。

把六个面上所有的小块全部加起来。

\[
\oiint (x+y)\,dS
\]

那个圈圈加两条杠的积分号，意思是：\textbf{绕着一个封闭的表面，把整个表面都算一遍}。

\section{我们的盒子}

一个正方体。

八个角在：

$(0,0,0)$，$(1,0,0)$，$(0,1,0)$，$(1,1,0)$

$(0,0,1)$，$(1,0,1)$，$(0,1,1)$，$(1,1,1)$

边长是 $1$。

六个面，我给它们起名字：

\[
\begin{array}{c|c|c}
\text{面} & \text{位置} & \text{特点} \\
\hline
\text{底面} & z = 0 & \text{最下面} \\
\text{顶面} & z = 1 & \text{最上面} \\
\text{前面} & y = 0 & \text{靠近你} \\
\text{后面} & y = 1 & \text{远离你} \\
\text{左面} & x = 0 & \text{左边} \\
\text{右面} & x = 1 & \text{右边}
\end{array}
\]

\section{一个面一个面地算}

\subsection*{底面（$z = 0$）}

底面是一个正方形，平躺在最下面。

$x$ 从 $0$ 到 $1$，$y$ 从 $0$ 到 $1$，$z$ 始终是 $0$。

这个面上每个点，脚下写的数是 $x + y$。

把整个底面上的 $(x+y)$ 加起来：

\[
\int_0^1 \int_0^1 (x+y)\,dx\,dy
\]

先对 $x$ 积分（$y$ 当常数）：

\[
\int_0^1 (x+y)\,dx = \frac{x^2}{2} + xy \Big|_0^1 = \frac{1}{2} + y
\]

再对 $y$ 积分：

\[
\int_0^1 \left(\frac{1}{2} + y\right)dy = \frac{1}{2}y + \frac{y^2}{2}\Big|_0^1 = \frac{1}{2} + \frac{1}{2} = 1
\]

\textbf{底面贡献：$1$}

\subsection*{顶面（$z = 1$）}

顶面也是一个正方形，在最上面。

$x$ 从 $0$ 到 $1$，$y$ 从 $0$ 到 $1$，$z$ 始终是 $1$。

但 $(x+y)$ 里面没有 $z$！

所以顶面和底面，\textbf{脚下的数是一模一样的}。

都是 $x + y$。

\[
\int_0^1 \int_0^1 (x+y)\,dx\,dy = 1
\]

跟底面一模一样。

\textbf{顶面贡献：$1$}

\subsection*{前面（$y = 0$）}

前面是靠近你的那一面。

$x$ 从 $0$ 到 $1$，$z$ 从 $0$ 到 $1$，$y$ 始终是 $0$。

这个面上，脚下写的数是：

\[
x + y = x + 0 = x
\]

\[
\int_0^1 \int_0^1 x\,dx\,dz
\]

先对 $x$：

\[
\int_0^1 x\,dx = \frac{1}{2}
\]

再对 $z$：

\[
\int_0^1 \frac{1}{2}\,dz = \frac{1}{2}
\]

\textbf{前面贡献：$\dfrac{1}{2}$}

\subsection*{后面（$y = 1$）}

后面是远离你的那一面。

$x$ 从 $0$ 到 $1$，$z$ 从 $0$ 到 $1$，$y$ 始终是 $1$。

这个面上，脚下写的数是：

\[
x + y = x + 1
\]

比前面\textbf{每个点都大了 $1$}。

\[
\int_0^1 \int_0^1 (x+1)\,dx\,dz
\]

先对 $x$：

\[
\int_0^1 (x+1)\,dx = \frac{1}{2} + 1 = \frac{3}{2}
\]

再对 $z$：

\[
\int_0^1 \frac{3}{2}\,dz = \frac{3}{2}
\]

\textbf{后面贡献：$\dfrac{3}{2}$}

\subsection*{左面（$x = 0$）}

左边那一面。

$y$ 从 $0$ 到 $1$，$z$ 从 $0$ 到 $1$，$x$ 始终是 $0$。

脚下的数是：

\[
x + y = 0 + y = y
\]

\[
\int_0^1 \int_0^1 y\,dy\,dz = \int_0^1 \frac{1}{2}\,dz = \frac{1}{2}
\]

\textbf{左面贡献：$\dfrac{1}{2}$}

\subsection*{右面（$x = 1$）}

右边那一面。

$y$ 从 $0$ 到 $1$，$z$ 从 $0$ 到 $1$，$x$ 始终是 $1$。

脚下的数是：

\[
x + y = 1 + y
\]

比左面\textbf{每个点都大了 $1$}。

\[
\int_0^1 \int_0^1 (1+y)\,dy\,dz = \int_0^1 \frac{3}{2}\,dz = \frac{3}{2}
\]

\textbf{右面贡献：$\dfrac{3}{2}$}

\section{六个面加起来}

\[
\begin{array}{c|c}
\text{面} & \text{贡献} \\
\hline
\text{底面} & 1 \\
\text{顶面} & 1 \\
\text{前面} & \frac{1}{2} \\
\text{后面} & \frac{3}{2} \\
\text{左面} & \frac{1}{2} \\
\text{右面} & \frac{3}{2}
\end{array}
\]

\[
1 + 1 + \frac{1}{2} + \frac{3}{2} + \frac{1}{2} + \frac{3}{2} = 2 + 2 + 2 = 6
\]

\[
\boxed{\oiint (x+y)\,dS = 6}
\]

\section{发现规律了吗？}

底面和顶面：$1 + 1 = 2$

前面和后面：$\frac{1}{2} + \frac{3}{2} = 2$

左面和右面：$\frac{1}{2} + \frac{3}{2} = 2$

\textbf{每一对相对的面，加起来都是 $2$！}

这不是巧合。

前面 $y=0$，后面 $y=1$。

前面脚下是 $x$，后面脚下是 $x+1$。

后面比前面，每个点都大了 $1$。

后面的面积是 $1$，所以后面比前面多贡献了 $1 \times 1 = 1$。

前面贡献 $\frac{1}{2}$，后面贡献 $\frac{3}{2}$。

差了 $1$。✓

左右面也一样。

左面 $x=0$，右面 $x=1$。

右面比左面，每个点都大了 $1$。

差了 $1$。✓

\section{现在升级：$\oiint (x+y+z)\,dS$}

同样的盒子，但脚下的数变成了 $x + y + z$。

\[
\boxed{\oiint (x+y+z)\,dS = 9}
\]

\section{为什么是 9？}

把 $(x+y+z)$ 拆成三块：

\[
\oiint (x+y+z)\,dS = \oiint x\,dS + \oiint y\,dS + \oiint z\,dS
\]

\textbf{$\oiint x\,dS$ 是多少？}

\begin{itemize}
\item 左面 ($x=0$)：脚下是 $0$，贡献 $0$
\item 右面 ($x=1$)：脚下是 $1$，贡献 $1$
\item 其他四个面：$x$ 从 $0$ 到 $1$，平均是 $\frac{1}{2}$，每个面面积是 $1$，所以每个面贡献 $\frac{1}{2}$
\end{itemize}

\[
\oiint x\,dS = 0 + 1 + \frac{1}{2} \times 4 = 3
\]

同样的道理：

\[
\oiint y\,dS = 3
\]

\[
\oiint z\,dS = 3
\]

\[
\oiint (x+y+z)\,dS = 3 + 3 + 3 = 9
\]

\section{一维 vs 二维：对比一下}

\[
\begin{array}{c|c|c}
 & \text{一重闭合积分} & \text{二重闭合积分} \\
\hline
\text{走的是什么} & \text{一条线（圈）} & \text{一个面（泡泡表面）} \\
\text{围住的是什么} & \text{一块平面} & \text{一团空间} \\
\text{符号} & \oint & \oiint \\
\text{例子} & \text{正方形的边} & \text{正方体的表面}
\end{array}
\]

一重闭合积分里，$dx$ 只管 $x$ 方向。往上下走的时候 $dx = 0$，不算数。

二重闭合积分里，$dS$ 是面积，六个面都要算，没有谁不算数。

\section{三维的一重闭合：$\oint (x+y+z)\,dx$}

现在来个新花样。

还是走一条线，但这条线\textbf{在三维空间里}。

路线是一个竖着立在空中的正方形：

\[
(0,0,0) \to (1,0,0) \to (1,0,1) \to (0,0,1) \to (0,0,0)
\]

四个点都在 $y = 0$ 的平面上。

逆时针走一圈（从前面看）。

\subsection*{第一段：$(0,0,0) \to (1,0,0)$，往右}

$y=0$，$z=0$，脚下的数是 $x + 0 + 0 = x$

\[
\int_0^1 x\,dx = \frac{1}{2}
\]

\subsection*{第二段：$(1,0,0) \to (1,0,1)$，往上}

$x$ 不变，$dx = 0$。

贡献：$0$

\subsection*{第三段：$(1,0,1) \to (0,0,1)$，往左}

$y=0$，$z=1$，脚下的数是 $x + 0 + 1 = x + 1$

$x$ 从 $1$ 到 $0$：

\[
\int_1^0 (x+1)\,dx = -\frac{3}{2}
\]

\subsection*{第四段：$(0,0,1) \to (0,0,0)$，往下}

$x$ 不变，$dx = 0$。

贡献：$0$

\subsection*{总和}

\[
\frac{1}{2} + 0 - \frac{3}{2} + 0 = -1
\]

\[
\boxed{\oint (x+y+z)\,dx = -1}
\]

\section{惊不惊喜？}

跟二维的 $\oint (x+y)\,dx = -1$ 一模一样！

为什么？

$(x+y+z)$ 里面多了个 $z$，但 $z$ 的部分会抵消。

第一段在底部，$z=0$。

第三段在顶部，$z=1$。

底部脚下是 $x+0$，顶部脚下是 $x+1$。

底部和顶部的变化刚好按方向算进去，最后留下的数跟前面的例子一致。

$y$ 呢？四条边上 $y$ 始终是 $0$，根本没变过，贡献 $0$。

\textbf{只有 $x$ 的变化和走的方向配合起来，才真正留下了东西。}

\section{总结}

\[
\begin{array}{c|c|c|c}
\text{积分} & \text{维度} & \text{例子} & \text{答案} \\
\hline
\oint (x+y)\,dx & \text{二维的线} & \text{正方形边界} & -1 \\
\oint (x+y+z)\,dx & \text{三维的线} & \text{竖立的正方形} & -1 \\
\oiint (x+y)\,dS & \text{三维的面} & \text{正方体表面} & 6 \\
\oiint (x+y+z)\,dS & \text{三维的面} & \text{正方体表面} & 9
\end{array}
\]

\textbf{关键区别}

\begin{itemize}
\item \textbf{一重闭合积分}：只管一个方向（比如 $dx$），其他方向不算数。来回走可能抵消。
\item \textbf{二重闭合积分}：整个表面都算数（$dS$），每一块面积都有贡献。
\end{itemize}

\textbf{一个是走线，一个是摸面。}

\textbf{走线的时候，方向很重要，来回会抵消。}

\textbf{摸面的时候，面积都是正的，不存在抵消。}

\chapter{算 $x-y$}

\section{一重闭合：$\oint (x-y)\,dx$}

一个正方形。

四个角在：$(0,0)$，$(1,0)$，$(1,1)$，$(0,1)$。

边长是 $1$。

你沿着边，\textbf{逆时针}走一圈。

每走一小步，低头看看脚下写的 $(x-y)$ 是多少，乘以你这一步在 $x$ 方向走的 $dx$。

走完一圈，全部加起来。

\subsection*{底边：$(0,0) \to (1,0)$，往右}

你走在最下面那条线上。

$y$ 始终是 $0$。

脚下写的是 $x - y = x - 0 = x$。

$x$ 从 $0$ 走到 $1$。

\[
\int_0^1 x\,dx = \frac{x^2}{2}\bigg|_0^1 = \frac{1^2}{2} - \frac{0^2}{2} = \frac{1}{2} - 0 = \frac{1}{2}
\]

\textbf{底边攒了：$+\dfrac{1}{2}$}

\subsection*{右边：$(1,0) \to (1,1)$，往上}

你贴着最右边那条线往上走。

$x$ 始终是 $1$，没有左右移动，$dx = 0$。

\textbf{右边攒了：$0$}

\subsection*{顶边：$(1,1) \to (0,1)$，往左}

你走在最上面那条线上。

$y$ 始终是 $1$。

脚下写的是 $x - y = x - 1$。

$x$ 从 $1$ 走到 $0$（方向反了）。

\[
\int_1^0 (x-1)\,dx
\]

先找原函数。

$(x-1)$ 里面有两块：$x$ 和 $-1$。

$x$ 的原函数是 $\dfrac{x^2}{2}$。

$-1$ 的原函数是 $-x$。

合起来：原函数是 $\dfrac{x^2}{2} - x$。

代入上限 $x = 0$：

\[
\frac{0^2}{2} - 0 = 0 - 0 = 0
\]

代入下限 $x = 1$：

\[
\frac{1^2}{2} - 1 = \frac{1}{2} - 1 = -\frac{1}{2}
\]

上限减下限：

\[
0 - \left(-\frac{1}{2}\right) = 0 + \frac{1}{2} = +\frac{1}{2}
\]

\textbf{顶边攒了：$+\dfrac{1}{2}$}

\subsection*{左边：$(0,1) \to (0,0)$，往下}

你贴着最左边那条线往下走。

$x$ 始终是 $0$，没有左右移动，$dx = 0$。

\textbf{左边攒了：$0$}

\subsection*{四段加起来}

\[
\begin{array}{c|c|c}
\text{边} & \text{方向} & \text{贡献} \\
\hline
\text{底边} & \text{往右} & +\frac{1}{2} \\
\text{右边} & \text{往上} & 0 \\
\text{顶边} & \text{往左} & +\frac{1}{2} \\
\text{左边} & \text{往下} & 0
\end{array}
\]

\[
\frac{1}{2} + 0 + \frac{1}{2} + 0 = 1
\]

\[
\boxed{\oint (x-y)\,dx = +1}
\]

\section{之前算过 $\oint (x+y)\,dx = -1$}

现在换了个减号，答案从 $-1$ 变成了 $+1$。

\textbf{符号反了！}

为什么？拆开看：

\[
\oint (x-y)\,dx = \oint x\,dx - \oint y\,dx
\]

你已经知道：

\begin{itemize}
\item $\oint x\,dx = 0$（底边和顶边脚下的数一样，来回抵消）
\item $\oint y\,dx = -1$
\end{itemize}

所以：

\[
0 - (-1) = +1
\]

\textbf{减去一个负数，就变成了正数。}

\section{用“对面”来看}

上次 $(x + y)$：

\begin{itemize}
\item 底边上，$y = 0$，脚下是 $x + 0 = x$
\item 顶边上，$y = 1$，脚下是 $x + 1$
\item 顶边比底边\textbf{每个点都大了 $1$}
\end{itemize}

这次 $(x - y)$：

\begin{itemize}
\item 底边上，$y = 0$，脚下是 $x - 0 = x$
\item 顶边上，$y = 1$，脚下是 $x - 1$
\item 顶边比底边\textbf{每个点都小了 $1$}
\end{itemize}

上次顶边脚下的数更大，往左走（负方向），大数乘以负方向，贡献很负。

这次顶边脚下的数更小，甚至变成了负数，负数乘以负方向，反而变正了。

\section{二重闭合：$\oiint (x-y)\,dS$}

一个正方体。

八个角在：

$(0,0,0)$，$(1,0,0)$，$(0,1,0)$，$(1,1,0)$

$(0,0,1)$，$(1,0,1)$，$(0,1,1)$，$(1,1,1)$

边长是 $1$，六个面。

每摸到一小块面积 $dS$，就看看这个位置的 $(x-y)$ 是多少，乘以这一小块面积。

六个面全部加起来。

\subsection*{底面（$z = 0$）}

$x$ 从 $0$ 到 $1$，$y$ 从 $0$ 到 $1$，$z$ 始终是 $0$。

脚下写的是 $x - y$。

\[
\int_0^1 \int_0^1 (x - y)\,dx\,dy
\]

先对 $x$ 积分（把 $y$ 当成不动的常数）：

\[
\int_0^1 (x - y)\,dx
\]

里面有两块：$x$ 和 $-y$。

$x$ 的原函数是 $\dfrac{x^2}{2}$。

$-y$ 的原函数是 $-yx$。

合起来：原函数是 $\dfrac{x^2}{2} - yx$。

代入上限 $x = 1$：

\[
\frac{1^2}{2} - y \times 1 = \frac{1}{2} - y
\]

代入下限 $x = 0$：

\[
\frac{0^2}{2} - y \times 0 = 0 - 0 = 0
\]

上限减下限：

\[
\frac{1}{2} - y - 0 = \frac{1}{2} - y
\]

再对 $y$ 积分：

\[
\int_0^1 \left(\frac{1}{2} - y\right)dy
\]

里面有两块：$\dfrac{1}{2}$ 和 $-y$。

$\dfrac{1}{2}$ 是常数，原函数是 $\dfrac{1}{2} \times y = \dfrac{y}{2}$。

$-y$ 的原函数是 $-\dfrac{y^2}{2}$。

合起来：原函数是 $\dfrac{y}{2} - \dfrac{y^2}{2}$。

代入上限 $y = 1$：

\[
\frac{1}{2} - \frac{1}{2} = 0
\]

代入下限 $y = 0$：

\[
0 - 0 = 0
\]

上限减下限：

\[
0 - 0 = 0
\]

\textbf{底面贡献：$0$}

\subsection*{顶面（$z = 1$）}

$x$ 从 $0$ 到 $1$，$y$ 从 $0$ 到 $1$，$z$ 始终是 $1$。

脚下写的是 $x - y$。

$(x - y)$ 里面没有 $z$。

$z$ 是 $0$ 还是 $1$，脚下的数都一样。

所以跟底面的计算\textbf{一模一样}，结果也是 $0$。

\textbf{顶面贡献：$0$}

\subsection*{前面（$y = 0$）}

$x$ 从 $0$ 到 $1$，$z$ 从 $0$ 到 $1$，$y$ 始终是 $0$。

脚下写的是 $x - 0 = x$。

\[
\int_0^1 \int_0^1 x\,dx\,dz
\]

先对 $x$ 积分：

\[
\int_0^1 x\,dx = \frac{x^2}{2}\bigg|_0^1 = \frac{1}{2}
\]

再对 $z$ 积分：

\[
\int_0^1 \frac{1}{2}\,dz = \frac{1}{2}
\]

\textbf{前面贡献：$+\dfrac{1}{2}$}

\subsection*{后面（$y = 1$）}

$x$ 从 $0$ 到 $1$，$z$ 从 $0$ 到 $1$，$y$ 始终是 $1$。

脚下写的是 $x - 1$。

\[
\int_0^1 \int_0^1 (x - 1)\,dx\,dz
\]

先对 $x$ 积分：

\[
\int_0^1 (x-1)\,dx
\]

里面有两块：$x$ 和 $-1$。

$x$ 的原函数是 $\dfrac{x^2}{2}$。

$-1$ 的原函数是 $-x$。

合起来：$\dfrac{x^2}{2} - x$。

代入上限 $x = 1$：

\[
\frac{1}{2} - 1 = -\frac{1}{2}
\]

代入下限 $x = 0$：

\[
0 - 0 = 0
\]

上限减下限：

\[
-\frac{1}{2} - 0 = -\frac{1}{2}
\]

再对 $z$ 积分：

\[
\int_0^1 \left(-\frac{1}{2}\right)dz = -\frac{1}{2}
\]

\textbf{后面贡献：$-\dfrac{1}{2}$}

\subsection*{左面（$x = 0$）}

$y$ 从 $0$ 到 $1$，$z$ 从 $0$ 到 $1$，$x$ 始终是 $0$。

脚下写的是 $0 - y = -y$。

\[
\int_0^1 \int_0^1 (-y)\,dy\,dz
\]

先对 $y$ 积分：

\[
\int_0^1 (-y)\,dy = -\frac{y^2}{2}\bigg|_0^1 = -\frac{1}{2}
\]

再对 $z$ 积分：

\[
\int_0^1 \left(-\frac{1}{2}\right)dz = -\frac{1}{2}
\]

\textbf{左面贡献：$-\dfrac{1}{2}$}

\subsection*{右面（$x = 1$）}

$y$ 从 $0$ 到 $1$，$z$ 从 $0$ 到 $1$，$x$ 始终是 $1$。

脚下写的是 $1 - y$。

\[
\int_0^1 \int_0^1 (1 - y)\,dy\,dz
\]

先对 $y$ 积分：

\[
\int_0^1 (1 - y)\,dy
\]

里面有两块：$1$ 和 $-y$。

$1$ 的原函数是 $y$。

$-y$ 的原函数是 $-\dfrac{y^2}{2}$。

合起来：$y - \dfrac{y^2}{2}$。

代入上限 $y = 1$：

\[
1 - \frac{1}{2} = \frac{1}{2}
\]

代入下限 $y = 0$：

\[
0 - 0 = 0
\]

上限减下限：

\[
\frac{1}{2} - 0 = \frac{1}{2}
\]

再对 $z$ 积分：

\[
\int_0^1 \frac{1}{2}\,dz = \frac{1}{2}
\]

\textbf{右面贡献：$+\dfrac{1}{2}$}

\subsection*{六个面加起来}

\[
\begin{array}{c|c}
\text{面} & \text{贡献} \\
\hline
\text{底面} & 0 \\
\text{顶面} & 0 \\
\text{前面} & \frac{1}{2} \\
\text{后面} & -\frac{1}{2} \\
\text{左面} & -\frac{1}{2} \\
\text{右面} & \frac{1}{2}
\end{array}
\]

\[
0 + 0 + \frac{1}{2} + \left(-\frac{1}{2}\right) + \left(-\frac{1}{2}\right) + \frac{1}{2} = 0
\]

\[
\boxed{\oiint (x-y)\,dS = 0}
\]

\section{零？！}

上次 $\oiint (x+y)\,dS = 6$。

这次 $\oiint (x-y)\,dS = 0$。

拆开看：

\[
\oiint (x-y)\,dS = \oiint x\,dS - \oiint y\,dS
\]

上一章你算过：

\begin{itemize}
\item $\oiint x\,dS = 3$
\item $\oiint y\,dS = 3$
\end{itemize}

所以：

\[
3 - 3 = 0
\]

\textbf{$x$ 贡献了 $3$，$y$ 也贡献了 $3$。}

\textbf{上次加号，$3 + 3 = 6$，它们叠在一起。}

\textbf{这次减号，$3 - 3 = 0$，它们互相吃掉了。}

\chapter{算 $x-y-z$}

\section{一重闭合：$\oint (x-y-z)\,dx$}

一个正方形，竖着立在空中。

四个角在：$(0,0,0)$，$(1,0,0)$，$(1,0,1)$，$(0,0,1)$。

四个点都在 $y = 0$ 的平面上。

边长是 $1$。

你沿着边，\textbf{逆时针}走一圈（从前面看）。

每走一小步，低头看看脚下写的 $(x - y - z)$ 是多少，乘以你这一步在 $x$ 方向走的 $dx$。

走完一圈，全部加起来。

\subsection*{第一段：$(0,0,0) \to (1,0,0)$，往右}

你走在最下面那条线上。

$y$ 始终是 $0$，$z$ 始终是 $0$。

脚下写的是 $x - 0 - 0 = x$。

$x$ 从 $0$ 走到 $1$。

\[
\int_0^1 x\,dx = \frac{x^2}{2}\bigg|_0^1 = \frac{1^2}{2} - \frac{0^2}{2} = \frac{1}{2} - 0 = \frac{1}{2}
\]

\textbf{第一段攒了：$+\dfrac{1}{2}$}

\subsection*{第二段：$(1,0,0) \to (1,0,1)$，往上}

你从右下角往上走。

$x$ 始终是 $1$，没有左右移动，$dx = 0$。

\textbf{第二段攒了：$0$}

\subsection*{第三段：$(1,0,1) \to (0,0,1)$，往左}

你走在最上面那条线上。

$y$ 始终是 $0$，$z$ 始终是 $1$。

脚下写的是 $x - 0 - 1 = x - 1$。

$x$ 从 $1$ 走到 $0$（方向反了）。

\[
\int_1^0 (x - 1)\,dx
\]

先找原函数。

$(x - 1)$ 里面有两块：$x$ 和 $-1$。

$x$ 的原函数是 $\dfrac{x^2}{2}$。

$-1$ 的原函数是 $-x$。

合起来：原函数是 $\dfrac{x^2}{2} - x$。

代入上限 $x = 0$：

\[
\frac{0^2}{2} - 0 = 0 - 0 = 0
\]

代入下限 $x = 1$：

\[
\frac{1^2}{2} - 1 = \frac{1}{2} - 1 = -\frac{1}{2}
\]

上限减下限：

\[
0 - \left(-\frac{1}{2}\right) = 0 + \frac{1}{2} = +\frac{1}{2}
\]

\textbf{第三段攒了：$+\dfrac{1}{2}$}

\subsection*{第四段：$(0,0,1) \to (0,0,0)$，往下}

你从左上角往下走。

$x$ 始终是 $0$，没有左右移动，$dx = 0$。

\textbf{第四段攒了：$0$}

\subsection*{四段加起来}

\[
\begin{array}{c|c|c}
\text{段} & \text{方向} & \text{贡献} \\
\hline
\text{第一段} & \text{往右} & +\frac{1}{2} \\
\text{第二段} & \text{往上} & 0 \\
\text{第三段} & \text{往左} & +\frac{1}{2} \\
\text{第四段} & \text{往下} & 0
\end{array}
\]

\[
\frac{1}{2} + 0 + \frac{1}{2} + 0 = 1
\]

\[
\boxed{\oint (x - y - z)\,dx = +1}
\]

\section{跟之前比一比}

上一章，$\oint (x - y)\,dx = +1$。

现在，$\oint (x - y - z)\,dx = +1$。

\textbf{多了个 $-z$，答案居然没变？}

为什么？

拆开看：

\[
\oint (x - y - z)\,dx = \oint x\,dx - \oint y\,dx - \oint z\,dx
\]

\textbf{$\oint x\,dx$}：

第一段（底部，往右）：$\int_0^1 x\,dx = +\dfrac{1}{2}$

第三段（顶部，往左）：$\int_1^0 x\,dx = -\dfrac{1}{2}$

其他两段：$dx = 0$。

加起来：$\dfrac{1}{2} - \dfrac{1}{2} = 0$

底部和顶部脚下都是 $x$，一模一样。来回一走，抵消了。

\textbf{$\oint y\,dx$}：

四条边上，$y$ 始终是 $0$（整个正方形都在 $y = 0$ 的平面上）。

脚下永远是 $0$。

加起来：$0$

\textbf{$\oint z\,dx$}：

第一段（底部）：$z = 0$，脚下是 $0$。$\int_0^1 0\,dx = 0$

第三段（顶部）：$z = 1$，脚下是 $1$。$\int_1^0 1\,dx = 1 \times (0 - 1) = -1$

其他两段：$dx = 0$。

加起来：$0 + (-1) = -1$

所以：

\[
\oint (x - y - z)\,dx = 0 - 0 - (-1) = 0 - 0 + 1 = +1
\]

\textbf{$x$ 的部分：底部和顶部一样，抵消了，贡献 $0$。}

\textbf{$y$ 的部分：到处都是 $0$，贡献 $0$。}

\textbf{$z$ 的部分：底部是 $0$，顶部是 $1$，没抵消，贡献 $-1$。但前面有负号，$-(-1) = +1$。}

\section{二重闭合：$\oiint (x-y-z)\,dS$}

一个正方体。

八个角在：

$(0,0,0)$，$(1,0,0)$，$(0,1,0)$，$(1,1,0)$

$(0,0,1)$，$(1,0,1)$，$(0,1,1)$，$(1,1,1)$

边长是 $1$，六个面。

每摸到一小块面积 $dS$，就看看这个位置的 $(x - y - z)$ 是多少，乘以这一小块面积。

六个面全部加起来。

\subsection*{底面（$z = 0$）}

$x$ 从 $0$ 到 $1$，$y$ 从 $0$ 到 $1$，$z$ 始终是 $0$。

脚下写的是 $x - y - 0 = x - y$。

\[
\int_0^1 \int_0^1 (x - y)\,dx\,dy
\]

前面已经算过，这个结果是 $0$。

\textbf{底面贡献：$0$}

\subsection*{顶面（$z = 1$）}

$x$ 从 $0$ 到 $1$，$y$ 从 $0$ 到 $1$，$z$ 始终是 $1$。

脚下写的是 $x - y - 1$。

\[
\int_0^1 \int_0^1 (x - y - 1)\,dx\,dy
\]

先对 $x$ 积分（$y$ 当常数）：

\[
\int_0^1 (x - y - 1)\,dx
\]

三块：$x$，$-y$，$-1$。

$x$ 的原函数是 $\dfrac{x^2}{2}$。

$-y$ 的原函数是 $-yx$。

$-1$ 的原函数是 $-x$。

合起来：$\dfrac{x^2}{2} - yx - x$。

代入上限 $x = 1$：

\[
\frac{1}{2} - y - 1 = -\frac{1}{2} - y
\]

代入下限 $x = 0$：

\[
0
\]

上限减下限：

\[
-\frac{1}{2} - y
\]

再对 $y$ 积分：

\[
\int_0^1 \left(-\frac{1}{2} - y\right)dy
\]

两块：$-\dfrac{1}{2}$ 和 $-y$。

$-\dfrac{1}{2}$ 的原函数是 $-\dfrac{y}{2}$。

$-y$ 的原函数是 $-\dfrac{y^2}{2}$。

合起来：$-\dfrac{y}{2} - \dfrac{y^2}{2}$。

代入上限 $y = 1$：

\[
-\frac{1}{2} - \frac{1}{2} = -1
\]

代入下限 $y = 0$：

\[
0
\]

上限减下限：

\[
-1
\]

\textbf{顶面贡献：$-1$}

\subsection*{前面（$y = 0$）}

$x$ 从 $0$ 到 $1$，$z$ 从 $0$ 到 $1$，$y$ 始终是 $0$。

脚下写的是 $x - 0 - z = x - z$。

\[
\int_0^1 \int_0^1 (x - z)\,dx\,dz
\]

先对 $x$ 积分（$z$ 当常数）：

\[
\int_0^1 (x - z)\,dx
\]

两块：$x$ 和 $-z$。

$x$ 的原函数是 $\dfrac{x^2}{2}$。

$-z$ 的原函数是 $-zx$。

合起来：$\dfrac{x^2}{2} - zx$。

代入上限 $x = 1$：

\[
\frac{1}{2} - z
\]

代入下限 $x = 0$：

\[
0
\]

上限减下限：

\[
\frac{1}{2} - z
\]

再对 $z$ 积分：

\[
\int_0^1 \left(\frac{1}{2} - z\right)dz
\]

两块：$\dfrac{1}{2}$ 和 $-z$。

$\dfrac{1}{2}$ 的原函数是 $\dfrac{z}{2}$。

$-z$ 的原函数是 $-\dfrac{z^2}{2}$。

合起来：$\dfrac{z}{2} - \dfrac{z^2}{2}$。

代入上限 $z = 1$：

\[
\frac{1}{2} - \frac{1}{2} = 0
\]

代入下限 $z = 0$：

\[
0
\]

上限减下限：

\[
0
\]

\textbf{前面贡献：$0$}

\subsection*{后面（$y = 1$）}

$x$ 从 $0$ 到 $1$，$z$ 从 $0$ 到 $1$，$y$ 始终是 $1$。

脚下写的是 $x - 1 - z$。

\[
\int_0^1 \int_0^1 (x - 1 - z)\,dx\,dz
\]

先对 $x$ 积分（$z$ 当常数）：

\[
\int_0^1 (x - 1 - z)\,dx
\]

三块：$x$，$-1$，$-z$。

$x$ 的原函数是 $\dfrac{x^2}{2}$。

$-1$ 的原函数是 $-x$。

$-z$ 的原函数是 $-zx$。

合起来：$\dfrac{x^2}{2} - x - zx$。

代入上限 $x = 1$：

\[
\frac{1}{2} - 1 - z = -\frac{1}{2} - z
\]

代入下限 $x = 0$：

\[
0
\]

上限减下限：

\[
-\frac{1}{2} - z
\]

再对 $z$ 积分：

\[
\int_0^1 \left(-\frac{1}{2} - z\right)dz
\]

两块：$-\dfrac{1}{2}$ 和 $-z$。

$-\dfrac{1}{2}$ 的原函数是 $-\dfrac{z}{2}$。

$-z$ 的原函数是 $-\dfrac{z^2}{2}$。

合起来：$-\dfrac{z}{2} - \dfrac{z^2}{2}$。

代入上限 $z = 1$：

\[
-\frac{1}{2} - \frac{1}{2} = -1
\]

代入下限 $z = 0$：

\[
0
\]

上限减下限：

\[
-1
\]

\textbf{后面贡献：$-1$}

\subsection*{左面（$x = 0$）}

$y$ 从 $0$ 到 $1$，$z$ 从 $0$ 到 $1$，$x$ 始终是 $0$。

脚下写的是 $0 - y - z = -y - z$。

\[
\int_0^1 \int_0^1 (-y - z)\,dy\,dz
\]

先对 $y$ 积分（$z$ 当常数）：

\[
\int_0^1 (-y - z)\,dy
\]

两块：$-y$ 和 $-z$。

$-y$ 的原函数是 $-\dfrac{y^2}{2}$。

$-z$ 的原函数是 $-zy$。

合起来：$-\dfrac{y^2}{2} - zy$。

代入上限 $y = 1$：

\[
-\frac{1}{2} - z
\]

代入下限 $y = 0$：

\[
0
\]

上限减下限：

\[
-\frac{1}{2} - z
\]

再对 $z$ 积分：

\[
\int_0^1 \left(-\frac{1}{2} - z\right)dz
\]

两块：$-\dfrac{1}{2}$ 和 $-z$。

$-\dfrac{1}{2}$ 的原函数是 $-\dfrac{z}{2}$。

$-z$ 的原函数是 $-\dfrac{z^2}{2}$。

合起来：$-\dfrac{z}{2} - \dfrac{z^2}{2}$。

代入上限 $z = 1$：

\[
-\frac{1}{2} - \frac{1}{2} = -1
\]

代入下限 $z = 0$：

\[
0
\]

上限减下限：

\[
-1
\]

\textbf{左面贡献：$-1$}

\subsection*{右面（$x = 1$）}

$y$ 从 $0$ 到 $1$，$z$ 从 $0$ 到 $1$，$x$ 始终是 $1$。

脚下写的是 $1 - y - z$。

\[
\int_0^1 \int_0^1 (1 - y - z)\,dy\,dz
\]

先对 $y$ 积分（$z$ 当常数）：

\[
\int_0^1 (1 - y - z)\,dy
\]

三块：$1$，$-y$，$-z$。

$1$ 的原函数是 $y$。

$-y$ 的原函数是 $-\dfrac{y^2}{2}$。

$-z$ 的原函数是 $-zy$。

合起来：$y - \dfrac{y^2}{2} - zy$。

代入上限 $y = 1$：

\[
1 - \frac{1}{2} - z = \frac{1}{2} - z
\]

代入下限 $y = 0$：

\[
0
\]

上限减下限：

\[
\frac{1}{2} - z
\]

再对 $z$ 积分：

\[
\int_0^1 \left(\frac{1}{2} - z\right)dz
\]

两块：$\dfrac{1}{2}$ 和 $-z$。

$\dfrac{1}{2}$ 的原函数是 $\dfrac{z}{2}$。

$-z$ 的原函数是 $-\dfrac{z^2}{2}$。

合起来：$\dfrac{z}{2} - \dfrac{z^2}{2}$。

代入上限 $z = 1$：

\[
\frac{1}{2} - \frac{1}{2} = 0
\]

代入下限 $z = 0$：

\[
0
\]

上限减下限：

\[
0
\]

\textbf{右面贡献：$0$}

\subsection*{六个面加起来}

\[
\begin{array}{c|c|c}
\text{面} & \text{脚下的数} & \text{贡献} \\
\hline
\text{底面} & x - y & 0 \\
\text{顶面} & x - y - 1 & -1 \\
\text{前面} & x - z & 0 \\
\text{后面} & x - 1 - z & -1 \\
\text{左面} & -y - z & -1 \\
\text{右面} & 1 - y - z & 0
\end{array}
\]

\[
0 + (-1) + 0 + (-1) + (-1) + 0 = -3
\]

\[
\boxed{\oiint (x - y - z)\,dS = -3}
\]

\section{为什么是 $-3$？}

拆开看：

\[
\oiint (x - y - z)\,dS = \oiint x\,dS - \oiint y\,dS - \oiint z\,dS
\]

之前算过：

\begin{itemize}
\item $\oiint x\,dS = 3$
\item $\oiint y\,dS = 3$
\item $\oiint z\,dS = 3$
\end{itemize}

所以：

\[
3 - 3 - 3 = -3
\]

\section{看看“对面”}

\[
\begin{array}{c|c|c}
\text{对面} & \text{贡献} & \text{相加} \\
\hline
\text{底面 }0，\ \text{顶面 }-1 & 0 + (-1) & -1 \\
\text{前面 }0，\ \text{后面 }-1 & 0 + (-1) & -1 \\
\text{左面 }-1，\ \text{右面 }0 & (-1) + 0 & -1
\end{array}
\]

\textbf{每一对都是 $-1$。}

三对加起来：$-1 + (-1) + (-1) = -3$。

\section{大对比}

\[
\begin{array}{c|c|c}
 & \oint \cdot\,dx\ (\text{走线}) & \oiint \cdot\,dS\ (\text{摸面}) \\
\hline
x + y & -1 & 6 \\
x - y & +1 & 0 \\
x + y + z & -1 & 9 \\
x - y - z & +1 & -3
\end{array}
\]

走线那一列：

不管加几个还是减几个，答案都是 $+1$ 或 $-1$。

因为只有一个方向的变化能留下痕迹（从底到顶，脚下的数变了多少），其他方向 $dx = 0$，根本不算。

摸面那一列：

$x + y$：$3 + 3 = 6$

$x - y$：$3 - 3 = 0$

$x + y + z$：$3 + 3 + 3 = 9$

$x - y - z$：$3 - 3 - 3 = -3$

\textbf{每多一个 $+$，就多加一个 $3$。}

\textbf{每多一个 $-$，就多减一个 $3$。}

\textbf{加号让贡献叠加，减号让贡献互相吃掉。}

\chapter{\texorpdfstring{算 $xy$}{算 xy}}

\section{\texorpdfstring{一重闭合：$\oint xy\,dx$}{一重闭合：oint xy dx}}

一个正方形。

四个角在：$(0,0)$，$(1,0)$，$(1,1)$，$(0,1)$。

边长是 $1$。

你沿着边，\textbf{逆时针}走一圈。

每走一小步，低头看看脚下写的 $xy$ 是多少，乘以你这一步在 $x$ 方向走的 $dx$。

走完一圈，全部加起来。

\subsection*{底边：$(0,0) \to (1,0)$，往右}

你走在最下面那条线上。

$y$ 始终是 $0$。

脚下写的是 $x \times 0 = 0$。

\[
\int_0^1 0\,dx = 0
\]

\textbf{底边攒了：$0$}

\subsection*{右边：$(1,0) \to (1,1)$，往上}

你贴着最右边那条线往上走。

$x$ 始终是 $1$，没有左右移动，$dx = 0$。

\textbf{右边攒了：$0$}

\subsection*{顶边：$(1,1) \to (0,1)$，往左}

你走在最上面那条线上。

$y$ 始终是 $1$。

脚下写的是 $x \times 1 = x$。

$x$ 从 $1$ 走到 $0$（方向反了）。

\[
\int_1^0 x\,dx
\]

$x$ 的原函数是 $\dfrac{x^2}{2}$。

代入上限 $x = 0$：

\[
\frac{0^2}{2} = 0
\]

代入下限 $x = 1$：

\[
\frac{1^2}{2} = \frac{1}{2}
\]

上限减下限：

\[
0 - \frac{1}{2} = -\frac{1}{2}
\]

\textbf{顶边攒了：$-\dfrac{1}{2}$}

\subsection*{左边：$(0,1) \to (0,0)$，往下}

你贴着最左边那条线往下走。

$x$ 始终是 $0$，没有左右移动，$dx = 0$。

\textbf{左边攒了：$0$}

\subsection*{四段加起来}

\[
\begin{array}{c|c|c|c}
\text{边} & \text{方向} & \text{脚下的数} & \text{贡献} \\
\hline
\text{底边} & \text{往右} & x \times 0 = 0 & 0 \\
\text{右边} & \text{往上} & \text{—（$dx=0$）} & 0 \\
\text{顶边} & \text{往左} & x \times 1 = x & -\frac{1}{2} \\
\text{左边} & \text{往下} & \text{—（$dx=0$）} & 0
\end{array}
\]

\[
0 + 0 + \left(-\frac{1}{2}\right) + 0 = -\frac{1}{2}
\]

\[
\boxed{\oint xy\,dx = -\frac{1}{2}}
\]

\section{跟之前的 $\oint y\,dx = -1$ 比一比}

上次脚下写的是 $y$，答案是 $-1$。

这次脚下写的是 $xy$，答案是 $-\dfrac{1}{2}$。

\textbf{乘了个 $x$，答案变小了？}

为什么？

上次 $\oint y\,dx$：

\begin{itemize}
\item 底边上 $y = 0$，脚下是 $0$，贡献 $0$
\item 顶边上 $y = 1$，脚下是 $1$（一个常数），贡献 $-1$
\end{itemize}

这次 $\oint xy\,dx$：

\begin{itemize}
\item 底边上 $y = 0$，脚下是 $x \times 0 = 0$，贡献 $0$
\item 顶边上 $y = 1$，脚下是 $x \times 1 = x$（不是常数了！），贡献 $-\dfrac{1}{2}$
\end{itemize}

上次顶边上每个点都写着 $1$，走完一整条边，攒了 $1 \times 1 = 1$。

这次顶边上写的是 $x$，从 $0$ 变到 $1$，\textbf{平均只有 $\dfrac{1}{2}$}，走完一整条边，攒了 $\dfrac{1}{2}$。

\textbf{乘了个 $x$，等于给顶边上的数打了个折扣。平均值从 $1$ 降到了 $\dfrac{1}{2}$。}

\section{\texorpdfstring{二重闭合：$\oiint xy\,dS$}{二重闭合：oiint xy dS}}

一个正方体。

八个角在：

$(0,0,0)$，$(1,0,0)$，$(0,1,0)$，$(1,1,0)$

$(0,0,1)$，$(1,0,1)$，$(0,1,1)$，$(1,1,1)$

边长是 $1$，六个面。

每摸到一小块面积 $dS$，就看看这个位置的 $xy$ 是多少，乘以这一小块面积。

六个面全部加起来。

\subsection*{底面（$z = 0$）}

$x$ 从 $0$ 到 $1$，$y$ 从 $0$ 到 $1$，$z$ 始终是 $0$。

脚下写的是 $xy$。

\[
\int_0^1 \int_0^1 xy\,dx\,dy
\]

先对 $x$ 积分（把 $y$ 当成不动的常数）：

\[
\int_0^1 xy\,dx
\]

$xy$ 对 $x$ 来说就是 $y$ 乘以 $x$。$y$ 是常数，可以提到外面。

\[
y \int_0^1 x\,dx = y \times \frac{x^2}{2}\bigg|_0^1 = y \times \left(\frac{1}{2} - 0\right) = \frac{y}{2}
\]

再对 $y$ 积分：

\[
\int_0^1 \frac{y}{2}\,dy
\]

\[
\frac{1}{2} \int_0^1 y\,dy = \frac{1}{2} \times \frac{y^2}{2}\bigg|_0^1 = \frac{1}{2} \times \left(\frac{1}{2} - 0\right) = \frac{1}{4}
\]

\textbf{底面贡献：$\dfrac{1}{4}$}

\subsection*{顶面（$z = 1$）}

$x$ 从 $0$ 到 $1$，$y$ 从 $0$ 到 $1$，$z$ 始终是 $1$。

脚下写的是 $xy$。

$xy$ 里面没有 $z$。$z$ 是 $0$ 还是 $1$，脚下的数都一样。

所以跟底面的计算\textbf{一模一样}。

\textbf{顶面贡献：$\dfrac{1}{4}$}

\subsection*{前面（$y = 0$）}

$x$ 从 $0$ 到 $1$，$z$ 从 $0$ 到 $1$，$y$ 始终是 $0$。

脚下写的是 $x \times 0 = 0$。

\[
\int_0^1 \int_0^1 0\,dx\,dz = 0
\]

\textbf{前面贡献：$0$}

\subsection*{后面（$y = 1$）}

$x$ 从 $0$ 到 $1$，$z$ 从 $0$ 到 $1$，$y$ 始终是 $1$。

脚下写的是 $x \times 1 = x$。

\[
\int_0^1 \int_0^1 x\,dx\,dz
\]

先对 $x$ 积分：

\[
\int_0^1 x\,dx = \frac{x^2}{2}\bigg|_0^1 = \frac{1}{2}
\]

再对 $z$ 积分：

\[
\int_0^1 \frac{1}{2}\,dz = \frac{1}{2}
\]

\textbf{后面贡献：$\dfrac{1}{2}$}

\subsection*{左面（$x = 0$）}

$y$ 从 $0$ 到 $1$，$z$ 从 $0$ 到 $1$，$x$ 始终是 $0$。

脚下写的是 $0 \times y = 0$。

\[
\int_0^1 \int_0^1 0\,dy\,dz = 0
\]

\textbf{左面贡献：$0$}

\subsection*{右面（$x = 1$）}

$y$ 从 $0$ 到 $1$，$z$ 从 $0$ 到 $1$，$x$ 始终是 $1$。

脚下写的是 $1 \times y = y$。

\[
\int_0^1 \int_0^1 y\,dy\,dz
\]

先对 $y$ 积分：

\[
\int_0^1 y\,dy = \frac{y^2}{2}\bigg|_0^1 = \frac{1}{2}
\]

再对 $z$ 积分：

\[
\int_0^1 \frac{1}{2}\,dz = \frac{1}{2}
\]

\textbf{右面贡献：$\dfrac{1}{2}$}

\subsection*{六个面加起来}

\[
\begin{array}{c|c|c}
\text{面} & \text{脚下的数} & \text{贡献} \\
\hline
\text{底面} & xy & \frac{1}{4} \\
\text{顶面} & xy & \frac{1}{4} \\
\text{前面} & 0 & 0 \\
\text{后面} & x & \frac{1}{2} \\
\text{左面} & 0 & 0 \\
\text{右面} & y & \frac{1}{2}
\end{array}
\]

\[
\frac{1}{4} + \frac{1}{4} + 0 + \frac{1}{2} + 0 + \frac{1}{2}
= \frac{3}{2}
\]

\[
\boxed{\oiint xy\,dS = \frac{3}{2}}
\]

\section{看看“对面”}

\textbf{底面 vs 顶面：}

底面 $\dfrac{1}{4}$，顶面 $\dfrac{1}{4}$。

加起来：$\dfrac{1}{2}$。

两个一模一样！因为 $xy$ 里没有 $z$，上下看到的数完全相同。

\textbf{前面 vs 后面：}

前面 $0$，后面 $\dfrac{1}{2}$。

加起来：$\dfrac{1}{2}$。

前面上 $y = 0$，脚下全是 $0$。

后面上 $y = 1$，脚下是 $x$，平均值 $\dfrac{1}{2}$。

\textbf{左面 vs 右面：}

左面 $0$，右面 $\dfrac{1}{2}$。

加起来：$\dfrac{1}{2}$。

左面上 $x = 0$，脚下全是 $0$。

右面上 $x = 1$，脚下是 $y$，平均值 $\dfrac{1}{2}$。

\textbf{每一对加起来都是 $\dfrac{1}{2}$！}

三对一共：$\dfrac{1}{2} + \dfrac{1}{2} + \dfrac{1}{2} = \dfrac{3}{2}$。✓

\section{跟加法的例子比一比}

\[
\begin{array}{c|c|c}
\text{脚下的数} & \oint \cdot\,dx\ (\text{走线}) & \oiint \cdot\,dS\ (\text{摸面}) \\
\hline
x + y & -1 & 6 \\
x - y & +1 & 0 \\
xy & -\frac{1}{2} & \frac{3}{2}
\end{array}
\]

\textbf{加法：两个变量各自独立贡献，互不干扰，最后加起来。}

\textbf{乘法：两个变量纠缠在一起。一个是 $0$，另一个不管多大都没用。}

一重闭合的差别：

$\oint y\,dx = -1$：顶边上每个点都是 $1$，攒满了一整条边。

$\oint xy\,dx = -\dfrac{1}{2}$：顶边上写的是 $x$，从 $0$ 到 $1$，平均只有 $\dfrac{1}{2}$。

\textbf{乘法让脚下的数变得不均匀了。}

二重闭合的差别：

$\oiint (x+y)\,dS = 6$：六个面都有贡献，没有哪个面是零。

$\oiint xy\,dS = \dfrac{3}{2}$：前面和左面的贡献是 $0$（因为 $x = 0$ 或 $y = 0$，相乘就是 $0$）。

\textbf{乘法有个特点：只要有一个是 $0$，整个就是 $0$。这让某些面直接“消失”了。}

\chapter{\texorpdfstring{算 $xyz$}{算 xyz}}

\section{\texorpdfstring{一重闭合：$\oint xyz\,dx$}{一重闭合：oint xyz dx}}

一个正方形，竖着立在空中。

四个角在：$(0,0,0)$，$(1,0,0)$，$(1,0,1)$，$(0,0,1)$。

四个点都在 $y = 0$ 的平面上。

边长是 $1$。

你沿着边，\textbf{逆时针}走一圈（从前面看）。

每走一小步，低头看看脚下写的 $xyz$ 是多少，乘以你这一步在 $x$ 方向走的 $dx$。

走完一圈，全部加起来。

\section{先停一秒}

四个角的坐标，看看 $y$：

$(0,\mathbf{0},0)$，$(1,\mathbf{0},0)$，$(1,\mathbf{0},1)$，$(0,\mathbf{0},1)$

\textbf{$y$ 全是 $0$。}

整个正方形都贴在 $y = 0$ 的平面上。

不管你走到哪个点，$y$ 都是 $0$。

\[
xyz = x \times 0 \times z = 0
\]

\textbf{脚下永远写着 $0$。}

\subsection*{第一段：$(0,0,0) \to (1,0,0)$，往右}

$y = 0$。

脚下写的是 $x \times 0 \times 0 = 0$。

\[
\int_0^1 0\,dx = 0
\]

\textbf{第一段攒了：$0$}

\subsection*{第二段：$(1,0,0) \to (1,0,1)$，往上}

$x$ 始终是 $1$，没有左右移动，$dx = 0$。

\textbf{第二段攒了：$0$}

\subsection*{第三段：$(1,0,1) \to (0,0,1)$，往左}

$y = 0$。

脚下写的是 $x \times 0 \times 1 = 0$。

\[
\int_1^0 0\,dx = 0
\]

\textbf{第三段攒了：$0$}

\subsection*{第四段：$(0,0,1) \to (0,0,0)$，往下}

$x$ 始终是 $0$，没有左右移动，$dx = 0$。

\textbf{第四段攒了：$0$}

\subsection*{四段加起来}

\[
0 + 0 + 0 + 0 = 0
\]

\[
\boxed{\oint xyz\,dx = 0}
\]

\section{这也太无聊了}

四段全是 $0$。

不是因为什么深刻的抵消。

不是因为方向相反互相吃掉。

只是因为 $y$ 始终是 $0$。

\textbf{乘法的规矩：三个数相乘，只要有一个是 $0$，整个就是 $0$。}

上一章 $xy$ 的时候，底边上 $y = 0$，底边就没了。

这一章 $xyz$，整条路上 $y = 0$，\textbf{四条边全没了}。

\section{跟加法比一比}

$\oint (x + y + z)\,dx = +1$。三个变量相加，$y = 0$ 只是让 $y$ 那一份消失，$x$ 和 $z$ 照样有贡献。

$\oint xyz\,dx = 0$。三个变量相乘，$y = 0$ 把\textbf{所有人}一起拖下水。

\textbf{加法：一个是 $0$，其他的还能活。}

\textbf{乘法：一个是 $0$，全部陪葬。}

\section{\texorpdfstring{二重闭合：$\oiint xyz\,dS$}{二重闭合：oiint xyz dS}}

一个正方体。

八个角在：

$(0,0,0)$，$(1,0,0)$，$(0,1,0)$，$(1,1,0)$

$(0,0,1)$，$(1,0,1)$，$(0,1,1)$，$(1,1,1)$

边长是 $1$，六个面。

每摸到一小块面积 $dS$，就看看这个位置的 $xyz$ 是多少，乘以这一小块面积。

六个面全部加起来。

\subsection*{底面（$z = 0$）}

$x$ 从 $0$ 到 $1$，$y$ 从 $0$ 到 $1$，$z$ 始终是 $0$。

脚下写的是 $x \times y \times 0 = 0$。

三个数相乘，$z = 0$，整个就是 $0$。

\[
\int_0^1 \int_0^1 0\,dx\,dy = 0
\]

\textbf{底面贡献：$0$}

\subsection*{顶面（$z = 1$）}

$x$ 从 $0$ 到 $1$，$y$ 从 $0$ 到 $1$，$z$ 始终是 $1$。

脚下写的是 $x \times y \times 1 = xy$。

$z = 1$，乘以 $1$ 等于没乘，所以就剩下 $xy$。

\[
\int_0^1 \int_0^1 xy\,dx\,dy
\]

前面已经算过，这个值是 $\dfrac{1}{4}$。

\textbf{顶面贡献：$\dfrac{1}{4}$}

\subsection*{前面（$y = 0$）}

$x$ 从 $0$ 到 $1$，$z$ 从 $0$ 到 $1$，$y$ 始终是 $0$。

脚下写的是 $x \times 0 \times z = 0$。

$y = 0$，整个就是 $0$。

\[
\int_0^1 \int_0^1 0\,dx\,dz = 0
\]

\textbf{前面贡献：$0$}

\subsection*{后面（$y = 1$）}

$x$ 从 $0$ 到 $1$，$z$ 从 $0$ 到 $1$，$y$ 始终是 $1$。

脚下写的是 $x \times 1 \times z = xz$。

$y = 1$，乘以 $1$ 等于没乘，剩下 $xz$。

\[
\int_0^1 \int_0^1 xz\,dx\,dz
\]

先对 $x$ 积分（$z$ 当常数）：

\[
\int_0^1 xz\,dx
\]

$xz$ 对 $x$ 来说，就是“常数 $z$”乘以“$x$”。

常数提到外面：

\[
z \int_0^1 x\,dx = z \times \frac{x^2}{2}\bigg|_0^1 = z \times \left(\frac{1}{2} - 0\right) = \frac{z}{2}
\]

再对 $z$ 积分：

\[
\int_0^1 \frac{z}{2}\,dz
\]

常数 $\dfrac{1}{2}$ 提到外面：

\[
\frac{1}{2}\int_0^1 z\,dz = \frac{1}{2} \times \frac{z^2}{2}\bigg|_0^1 = \frac{1}{2} \times \left(\frac{1}{2} - 0\right) = \frac{1}{4}
\]

\textbf{后面贡献：$\dfrac{1}{4}$}

\subsection*{左面（$x = 0$）}

$y$ 从 $0$ 到 $1$，$z$ 从 $0$ 到 $1$，$x$ 始终是 $0$。

脚下写的是 $0 \times y \times z = 0$。

$x = 0$，整个就是 $0$。

\[
\int_0^1 \int_0^1 0\,dy\,dz = 0
\]

\textbf{左面贡献：$0$}

\subsection*{右面（$x = 1$）}

$y$ 从 $0$ 到 $1$，$z$ 从 $0$ 到 $1$，$x$ 始终是 $1$。

脚下写的是 $1 \times y \times z = yz$。

$x = 1$，乘以 $1$ 等于没乘，剩下 $yz$。

\[
\int_0^1 \int_0^1 yz\,dy\,dz
\]

先对 $y$ 积分（$z$ 当常数）：

\[
z\int_0^1 y\,dy = z \times \frac{y^2}{2}\bigg|_0^1 = z \times \frac{1}{2} = \frac{z}{2}
\]

再对 $z$ 积分：

\[
\int_0^1 \frac{z}{2}\,dz
\]

\[
\frac{1}{2}\int_0^1 z\,dz = \frac{1}{2} \times \frac{z^2}{2}\bigg|_0^1 = \frac{1}{4}
\]

\textbf{右面贡献：$\dfrac{1}{4}$}

\subsection*{六个面加起来}

\[
\begin{array}{c|c|c|c}
\text{面} & \text{哪个变量被钉死} & \text{钉死的值} & \text{贡献} \\
\hline
\text{底面} & z & 0 & 0 \\
\text{顶面} & z & 1 & \frac{1}{4} \\
\text{前面} & y & 0 & 0 \\
\text{后面} & y & 1 & \frac{1}{4} \\
\text{左面} & x & 0 & 0 \\
\text{右面} & x & 1 & \frac{1}{4}
\end{array}
\]

\[
0 + \frac{1}{4} + 0 + \frac{1}{4} + 0 + \frac{1}{4} = \frac{3}{4}
\]

\[
\boxed{\oiint xyz\,dS = \frac{3}{4}}
\]

\section{发现规律了吗？}

六个面，\textbf{三个是 $0$，三个是 $\dfrac{1}{4}$}。

哪三个是 $0$？

\textbf{底面}（$z=0$），\textbf{前面}（$y=0$），\textbf{左面}（$x=0$）。

它们的共同点：\textbf{被钉死的那个变量等于 $0$。}

$xyz$ 三个数相乘，一个是 $0$，整个就是 $0$。

整个面上脚下全是 $0$，加起来当然还是 $0$。

哪三个是 $\dfrac{1}{4}$？

\textbf{顶面}（$z=1$），\textbf{后面}（$y=1$），\textbf{右面}（$x=1$）。

它们的共同点：\textbf{被钉死的那个变量等于 $1$。}

$1$ 乘以什么都是什么本身，等于没乘。

所以三个变量的乘积，变成了两个变量的乘积。

$xy$，$xz$，$yz$——每一个的积分都是 $\dfrac{1}{4}$。

\section{跟 $xy$ 比一比}

\[
\begin{array}{c|c|c}
\text{脚下的数} & \oint \cdot\,dx\ (\text{走线}) & \oiint \cdot\,dS\ (\text{摸面}) \\
\hline
xy & -\frac{1}{2} & \frac{3}{2} \\
xyz & 0 & \frac{3}{4}
\end{array}
\]

\textbf{走线}：$xy$ 还有半条命（顶边上 $y = 1$，活了下来），$xyz$ 连这半条命都没了（$y = 0$ 把所有人拖下水）。

\textbf{摸面}：$xy$ 有两个面是 $0$（前面和左面），$xyz$ 有三个面是 $0$（底面、前面、左面）。\textbf{多乘一个变量，多死一个面。}

\section{跟加法比一比}

\[
\begin{array}{c|c|c}
\text{脚下的数} & \oiint \cdot\,dS & \text{有几个面是 }0 \\
\hline
x + y & 6 & 0\ \text{个} \\
xy & \frac{3}{2} & 2\ \text{个} \\
x + y + z & 9 & 0\ \text{个} \\
xyz & \frac{3}{4} & 3\ \text{个}
\end{array}
\]

\textbf{加法：没有一个面是 $0$。因为 $x = 0$ 的时候，$y$ 和 $z$ 还在，还有贡献。各自独立，互不拖累。}

\textbf{乘法：变量等于 $0$ 的面直接消失。因为 $0$ 乘以任何数都是 $0$。一个倒下，全部陪葬。}

\chapter{\texorpdfstring{闭合遇次方}{闭合遇次方}}

\section{\texorpdfstring{一重闭合：$\oint y^2\,dx$}{一重闭合：oint y^2 dx}}

一个正方形。

四个角在：$(0,0)$，$(1,0)$，$(1,1)$，$(0,1)$。

边长是 $1$。

你沿着边，\textbf{逆时针}走一圈。

每走一小步，低头看看脚下写的 $y^2$ 是多少，乘以你这一步在 $x$ 方向走的 $dx$。

走完一圈，全部加起来。

\subsection*{底边：$(0,0) \to (1,0)$，往右}

$y$ 始终是 $0$。

脚下写的是 $0^2 = 0$。

\[
\int_0^1 0\,dx = 0
\]

\textbf{底边攒了：$0$}

\subsection*{右边：$(1,0) \to (1,1)$，往上}

$x$ 始终是 $1$，没有左右移动，$dx = 0$。

\textbf{右边攒了：$0$}

\subsection*{顶边：$(1,1) \to (0,1)$，往左}

$y$ 始终是 $1$。

脚下写的是 $1^2 = 1$。

$x$ 从 $1$ 走到 $0$（方向反了）。

\[
\int_1^0 1\,dx
\]

$1$ 是常数，原函数是 $x$。

代入上限 $x = 0$：

\[
0
\]

代入下限 $x = 1$：

\[
1
\]

上限减下限：

\[
0 - 1 = -1
\]

\textbf{顶边攒了：$-1$}

\subsection*{左边：$(0,1) \to (0,0)$，往下}

$x$ 始终是 $0$，$dx = 0$。

\textbf{左边攒了：$0$}

\subsection*{四段加起来}

\[
0 + 0 + (-1) + 0 = -1
\]

\[
\boxed{\oint y^2\,dx = -1}
\]

\section{等等……}

之前算过 $\oint y\,dx = -1$。

现在 $\oint y^2\,dx = -1$。

\textbf{$y$ 换成 $y^2$，答案居然没变？}

为什么？

底边上，$y = 0$。$\quad 0^1 = 0$，$\quad 0^2 = 0$。

顶边上，$y = 1$。$\quad 1^1 = 1$，$\quad 1^2 = 1$。

\textbf{$0$ 的几次方都是 $0$。$1$ 的几次方都是 $1$。}

底边和顶边上，$y$ 要么是 $0$，要么是 $1$。

不管你是 $y$、$y^2$、$y^3$、$y^{100}$——

$0$ 的任何次方还是 $0$，$1$ 的任何次方还是 $1$。

\textbf{所以答案永远是 $-1$。}

\section{\texorpdfstring{再来：$\oint x^2 y\,dx$}{再来：oint x^2 y dx}}

同样的正方形，同样逆时针。

脚下写的是 $x^2 y$。

\subsection*{底边：$(0,0) \to (1,0)$，往右}

$y = 0$。

脚下写的是 $x^2 \times 0 = 0$。

\[
\int_0^1 0\,dx = 0
\]

\textbf{底边攒了：$0$}

\subsection*{右边和左边}

$dx = 0$，都是 $0$。

\subsection*{顶边：$(1,1) \to (0,1)$，往左}

$y = 1$。

脚下写的是 $x^2 \times 1 = x^2$。

$x$ 从 $1$ 走到 $0$。

\[
\int_1^0 x^2\,dx
\]

$x^2$ 的原函数是 $\dfrac{x^3}{3}$。

代入上限 $x = 0$：

\[
\frac{0^3}{3} = 0
\]

代入下限 $x = 1$：

\[
\frac{1^3}{3} = \frac{1}{3}
\]

上限减下限：

\[
0 - \frac{1}{3} = -\frac{1}{3}
\]

\textbf{顶边攒了：$-\dfrac{1}{3}$}

\subsection*{四段加起来}

\[
0 + 0 + \left(-\frac{1}{3}\right) + 0 = -\frac{1}{3}
\]

\[
\boxed{\oint x^2 y\,dx = -\frac{1}{3}}
\]

\section{\texorpdfstring{再来一个：$\oint x^3 y^2\,dx$}{再来一个：oint x^3 y^2 dx}}

\subsection*{底边}

$y = 0$。$\quad x^3 \times 0^2 = 0$。

贡献：$0$

\subsection*{右边和左边}

$dx = 0$。

贡献：$0$

\subsection*{顶边}

$y = 1$。$\quad x^3 \times 1^2 = x^3$。

$x$ 从 $1$ 到 $0$。

\[
\int_1^0 x^3\,dx
\]

$x^3$ 的原函数是 $\dfrac{x^4}{4}$。

代入上限 $x = 0$：

\[
\frac{0^4}{4} = 0
\]

代入下限 $x = 1$：

\[
\frac{1^4}{4} = \frac{1}{4}
\]

上限减下限：

\[
0 - \frac{1}{4} = -\frac{1}{4}
\]

\textbf{顶边攒了：$-\dfrac{1}{4}$}

\subsection*{总和}

\[
\boxed{\oint x^3 y^2\,dx = -\frac{1}{4}}
\]

\section{\texorpdfstring{再来一个：$\oint x^4 y^{100}\,dx$}{再来一个：oint x^4 y^100 dx}}


\subsection*{底边}

$y = 0$。$\quad x^4 \times 0^{100} = 0$。

贡献：$0$

\subsection*{顶边}

$y = 1$。$\quad x^4 \times 1^{100} = x^4$。

\[
\int_1^0 x^4\,dx = \frac{x^5}{5}\bigg|_1^0 = 0 - \frac{1}{5} = -\frac{1}{5}
\]

\subsection*{总和}

\[
\boxed{\oint x^4 y^{100}\,dx = -\frac{1}{5}}
\]

\section{摆在一起看}

\[
\begin{array}{c|c|c|c}
\text{脚下的数} & x\ \text{的次方} & y\ \text{的次方} & \text{绕一圈的结果} \\
\hline
y & 0 & 1 & -\dfrac{1}{1} = -1 \\
xy & 1 & 1 & -\dfrac{1}{2} \\
x^2 y & 2 & 1 & -\dfrac{1}{3} \\
x^3 y^2 & 3 & 2 & -\dfrac{1}{4} \\
x^4 y^{100} & 4 & 100 & -\dfrac{1}{5}
\end{array}
\]

\section{看出来了吗？}

\textbf{答案只跟 $x$ 的次方有关！}

$x$ 的次方是 $m$，答案就是 $-\dfrac{1}{m+1}$。

\textbf{$y$ 的次方是几，完全不影响答案。}

$y^1$，$y^2$，$y^{100}$——无所谓。

因为在这个正方形上，$y$ 只有两个值。

底边 $y = 0$，顶边 $y = 1$。

$0$ 的任何次方都是 $0$，$1$ 的任何次方都是 $1$。

\textbf{次方改变不了 $0$ 和 $1$。}

\section{通用公式}

\[
\oint x^m y^n\,dx = -\frac{1}{m+1}
\]

（只要 $n \geq 1$，也就是说脚下的数里面有 $y$。）

如果 $n = 0$（没有 $y$，只有 $x^m$）呢？

底边脚下是 $x^m$，顶边脚下也是 $x^m$，一模一样。

来回一走，抵消了。

\[
\oint x^m\,dx = 0
\]

\section{为什么是 $m+1$？}

因为最后你要算的就是 $\int_0^1 x^m\,dx$。

$x^m$ 的原函数是 $\dfrac{x^{m+1}}{m+1}$。

代入 $0$ 到 $1$：$\dfrac{1}{m+1} - 0 = \dfrac{1}{m+1}$。

但顶边是往左走的，要加负号，所以是 $-\dfrac{1}{m+1}$。

\section{\texorpdfstring{二重闭合：$\oiint x^2 yz\,dS$}{二重闭合：oiint x^2 yz dS}}

一个正方体。

八个角在：

$(0,0,0)$，$(1,0,0)$，$(0,1,0)$，$(1,1,0)$

$(0,0,1)$，$(1,0,1)$，$(0,1,1)$，$(1,1,1)$

边长是 $1$，六个面。

脚下写的是 $x^2 yz$。

三个变量乘在一起，其中 $x$ 是平方。

\subsection*{底面（$z = 0$）}

$z = 0$。

$x^2 \times y \times 0 = 0$。

三个数相乘，一个是 $0$，整个就是 $0$。

\textbf{底面贡献：$0$}

\subsection*{顶面（$z = 1$）}

$z = 1$。

$x^2 \times y \times 1 = x^2 y$。

$1$ 乘以什么都是什么本身，$z$ 消失了，剩下 $x^2 y$。

\[
\int_0^1 \int_0^1 x^2 y\,dx\,dy
\]

先对 $x$ 积分（$y$ 当常数）：

\[
\int_0^1 x^2 y\,dx
\]

$y$ 是常数，提到外面：

\[
y\int_0^1 x^2\,dx
\]

$x^2$ 的原函数是 $\dfrac{x^3}{3}$。

\[
y \times \frac{x^3}{3}\bigg|_0^1 = y \times \left(\frac{1}{3} - 0\right) = \frac{y}{3}
\]

再对 $y$ 积分：

\[
\int_0^1 \frac{y}{3}\,dy
\]

$\dfrac{1}{3}$ 是常数，提到外面：

\[
\frac{1}{3}\int_0^1 y\,dy
\]

$y$ 的原函数是 $\dfrac{y^2}{2}$。

\[
\frac{1}{3} \times \frac{y^2}{2}\bigg|_0^1 = \frac{1}{3} \times \left(\frac{1}{2} - 0\right) = \frac{1}{6}
\]

\textbf{顶面贡献：$\dfrac{1}{6}$}

\subsection*{前面（$y = 0$）}

$y = 0$。

$x^2 \times 0 \times z = 0$。

\textbf{前面贡献：$0$}

\subsection*{后面（$y = 1$）}

$y = 1$。

$x^2 \times 1 \times z = x^2 z$。

$y$ 消失了，剩下 $x^2 z$。

\[
\int_0^1 \int_0^1 x^2 z\,dx\,dz
\]

先对 $x$ 积分（$z$ 当常数）：

\[
\int_0^1 x^2 z\,dx
\]

$z$ 是常数，提到外面：

\[
z\int_0^1 x^2\,dx = z \times \frac{x^3}{3}\bigg|_0^1 = z \times \frac{1}{3} = \frac{z}{3}
\]

再对 $z$ 积分：

\[
\int_0^1 \frac{z}{3}\,dz
\]

$\dfrac{1}{3}$ 提到外面：

\[
\frac{1}{3}\int_0^1 z\,dz = \frac{1}{3} \times \frac{z^2}{2}\bigg|_0^1 = \frac{1}{6}
\]

\textbf{后面贡献：$\dfrac{1}{6}$}

\subsection*{左面（$x = 0$）}

$x = 0$。

$0^2 \times y \times z = 0$。

\textbf{左面贡献：$0$}

\subsection*{右面（$x = 1$）}

$x = 1$。

$1^2 \times y \times z = yz$。

$x$ 消失了，剩下 $yz$。

\[
\int_0^1 \int_0^1 yz\,dy\,dz
\]

先对 $y$ 积分（$z$ 当常数）：

\[
z\int_0^1 y\,dy = z \times \frac{y^2}{2}\bigg|_0^1 = z \times \frac{1}{2} = \frac{z}{2}
\]

再对 $z$ 积分：

\[
\int_0^1 \frac{z}{2}\,dz = \frac{1}{2}\int_0^1 z\,dz = \frac{1}{2} \times \frac{z^2}{2}\bigg|_0^1 = \frac{1}{4}
\]

\textbf{右面贡献：$\dfrac{1}{4}$}

\subsection*{六个面加起来}

\[
\begin{array}{c|c|c|c}
\text{面} & \text{哪个变量是 }0 & \text{脚下的数} & \text{贡献} \\
\hline
\text{底面} & z = 0 & 0 & 0 \\
\text{顶面} & \text{—} & x^2 y & \frac{1}{6} \\
\text{前面} & y = 0 & 0 & 0 \\
\text{后面} & \text{—} & x^2 z & \frac{1}{6} \\
\text{左面} & x = 0 & 0 & 0 \\
\text{右面} & \text{—} & yz & \frac{1}{4}
\end{array}
\]

\[
0 + \frac{1}{6} + 0 + \frac{1}{6} + 0 + \frac{1}{4}
\]

通分（公分母 $12$）：

\[
\frac{2}{12} + \frac{2}{12} + \frac{3}{12} = \frac{7}{12}
\]

\[
\boxed{\oiint x^2 yz\,dS = \frac{7}{12}}
\]

\section{看看每个面的结果从哪来}

\textbf{顶面}：$z = 1$，剩下 $x^2 y$，积分变成 $\int_0^1 x^2\,dx$ 乘以 $\int_0^1 y\,dy$。

\[
\frac{1}{3} \times \frac{1}{2} = \frac{1}{6}
\]

\textbf{后面}：$y = 1$，剩下 $x^2 z$，积分变成 $\int_0^1 x^2\,dx$ 乘以 $\int_0^1 z\,dz$。

\[
\frac{1}{3} \times \frac{1}{2} = \frac{1}{6}
\]

\textbf{右面}：$x = 1$，剩下 $yz$，积分变成 $\int_0^1 y\,dy$ 乘以 $\int_0^1 z\,dz$。

\[
\frac{1}{2} \times \frac{1}{2} = \frac{1}{4}
\]

每一个不为零的面，都是\textbf{两个简单积分相乘}。

而每个简单积分都长这样：$\int_0^1 t^k\,dt = \dfrac{1}{k+1}$。

\section{通用公式}

脚下写的是 $x^a \cdot y^b \cdot z^c$（$a, b, c$ 都 $\geq 1$）。

\textbf{三个面是 $0$}（变量 $= 0$ 的那些面）：

底面（$z=0$）、前面（$y=0$）、左面（$x=0$）。

$0$ 的任何次方是 $0$，乘法里一个 $0$ 全完蛋。

\textbf{三个面有贡献}（变量 $= 1$ 的那些面）：

\[
\begin{array}{c|c|c|c}
\text{面} & \text{谁 }=1 & \text{剩下什么} & \text{贡献} \\
\hline
\text{顶面} & z = 1 & x^a y^b & \dfrac{1}{(a+1)(b+1)} \\
\text{后面} & y = 1 & x^a z^c & \dfrac{1}{(a+1)(c+1)} \\
\text{右面} & x = 1 & y^b z^c & \dfrac{1}{(b+1)(c+1)}
\end{array}
\]

\[
\oiint x^a y^b z^c\,dS = \frac{1}{(a+1)(b+1)} + \frac{1}{(a+1)(c+1)} + \frac{1}{(b+1)(c+1)}
\]

\section{验算一下之前的例子}

\textbf{$xyz$：$a=1, b=1, c=1$}

\[
\frac{1}{2 \times 2} + \frac{1}{2 \times 2} + \frac{1}{2 \times 2} = \frac{1}{4} + \frac{1}{4} + \frac{1}{4} = \frac{3}{4} \quad \checkmark
\]

\textbf{$x^2 yz$：$a=2, b=1, c=1$}

\[
\frac{1}{3 \times 2} + \frac{1}{3 \times 2} + \frac{1}{2 \times 2} = \frac{1}{6} + \frac{1}{6} + \frac{1}{4} = \frac{7}{12} \quad \checkmark
\]

\section{再试一个：$x^3 y^2 z^4$}

不用算六个面了。直接用公式。

$a = 3$，$b = 2$，$c = 4$。

\[
\frac{1}{4 \times 3} + \frac{1}{4 \times 5} + \frac{1}{3 \times 5}
\]

\[
= \frac{1}{12} + \frac{1}{20} + \frac{1}{15}
\]

通分（公分母 $60$）：

\[
= \frac{5}{60} + \frac{3}{60} + \frac{4}{60} = \frac{12}{60} = \frac{1}{5}
\]

\[
\boxed{\oiint x^3 y^2 z^4\,dS = \frac{1}{5}}
\]

\section{全景对比}

\subsection*{一重闭合的通用公式}

\[
\oint x^m y^n\,dx = -\frac{1}{m+1} \qquad (n \geq 1)
\]

\textbf{$y$ 的次方不影响答案。}

因为 $y$ 在边界上只有 $0$ 和 $1$，$0^n = 0$，$1^n = 1$。

\subsection*{二重闭合的通用公式}

\[
\oiint x^a y^b z^c\,dS = \frac{1}{(a+1)(b+1)} + \frac{1}{(a+1)(c+1)} + \frac{1}{(b+1)(c+1)}
\]

（$a, b, c$ 都 $\geq 1$）

\textbf{每一项，都是“把一个变量去掉，剩下两个变量各自积分再相乘”。}

去掉 $z$，剩 $x^a$ 和 $y^b$：$\dfrac{1}{a+1} \times \dfrac{1}{b+1}$

去掉 $y$，剩 $x^a$ 和 $z^c$：$\dfrac{1}{a+1} \times \dfrac{1}{c+1}$

去掉 $x$，剩 $y^b$ 和 $z^c$：$\dfrac{1}{b+1} \times \dfrac{1}{c+1}$

\subsection*{为什么这么简单？}

\textbf{因为正方形和正方体太特殊了。}

边界上的坐标，不是 $0$ 就是 $1$。

$0$ 的任何次方是 $0$——杀死整个面。

$1$ 的任何次方是 $1$——等于没有。

\textbf{变量 $= 0$ 的面：贡献 $0$。}

\textbf{变量 $= 1$ 的面：那个变量消失，剩下的各自积分。}

\section*{遇到除法怎么办？——别怕，大部分是纸老虎}

\bigskip\hrule\bigskip

\subsection*{📌 一个聪明的疑问}

你已经学过了加法：$f = x + y$，我们把它\textbf{拆开}。

你也学过了乘法：$f = xy$，我们把不动的那部分\textbf{提出来}。

如果你是个爱思考的孩子，你脑子里肯定会冒出一个问题：

\textbf{“那除法呢？”}

如果 $f$ 里面有分数，上面一个东西，下面一个东西，除在一起，该怎么办？

你能想到这个问题，说明你的数学直觉已经非常敏锐了！

今天，我们就来专门对付微积分里的\textbf{除法（分数）}。

\subsection*{第一种：纸老虎（下面是数字）}

最简单的除法，是下面只跟着一个普通的数字。

比如：
$$\int_0^2 \frac{x}{2} \, dx$$

翻译一下：每一小块的高度是 $\frac{x}{2}$（也就是 $x$ 除以 $2$）。

遇到这种题，千万别慌。

你想想，$\frac{x}{2}$ 是什么意思？

一个苹果除以 2，就是“半个苹果”。
$x$ 除以 2，就是“半个 $x$”，也就是 $\frac{1}{2} \times x$。

看到了吗？
\textbf{除以 2，其实就是乘以 $\frac{1}{2}$！}

$\frac{1}{2}$ 是什么？
是一个\textbf{常数}！它永远不动！

上一章我们学过：遇到不动的常数，怎么办？
\textbf{提出来！}

$$\int_0^2 \frac{x}{2} \, dx = \int_0^2 \left(\frac{1}{2} \times x\right) \, dx = \frac{1}{2} \times \int_0^2 x \, dx$$

剩下的部分（$\int_0^2 x \, dx$），你闭着眼睛都会算了：
$x$ 的反过来是 $\frac{x^2}{2}$。
代入 $2$ 得 $2$，代入 $0$ 得 $0$，相减等于 $2$。

最后乘以提出来的那个 $\frac{1}{2}$：
$$\frac{1}{2} \times 2 = 1$$

\textbf{答案：$1$}

\begin{quote}
\textbf{📌 总结：下面是数字}

除以一个数字 $\rightarrow$ 等于乘以它的倒数 $\rightarrow$ 变成常数提出来。

\textbf{完全不用怕！}
\end{quote}

\subsection*{第二种：假装很凶的纸老虎（可以约分）}

如果下面不是数字，而是字母 $x$ 呢？

比如这道题，看起来很吓人：
$$\int_1^3 \frac{x^2}{x} \, dx$$

你可能会想：“完了，上面在动，下面也在动，怎么提出来？”

\textbf{等一下。先别急着算积分。}

你看那个分数：$\frac{x^2}{x}$。
它能不能变简单点？

$x^2$ 是什么意思？
就是两个 $x$ 乘在一起：$x \times x$。

所以这个分数其实是：
$$\frac{x \times x}{x}$$

上面有两个 $x$ 相乘，下面有一个 $x$。
就像“拿两个苹果，再分给一个人”，你可以\textbf{抵消（约分）}掉一个！

上面的一个 $x$，和下面的那个 $x$ 抵消了。
只剩下一个 $x$！

$$\frac{x^2}{x} = x$$

所以，这道题根本不是什么复杂的除法。
它就是：
$$\int_1^3 x \, dx$$

这是我们最熟悉的题了！
$x$ 的反过来是 $\frac{x^2}{2}$。
代入 $3$：$\frac{3^2}{2} = 4.5$。
代入 $1$：$\frac{1^2}{2} = 0.5$。
相减：$4.5 - 0.5 = 4$。

\textbf{答案：$4$}

\begin{quote}
\textbf{📌 总结：上下都有字母}

遇到这种，第一反应是：\textbf{能不能抵消？}

把上面的 $x$ 和下面的 $x$ 划掉。

把它变回最简单的样子，再算。
\end{quote}

\subsection*{第三种：混合在一起的纸老虎}

再来一个稍微复杂一点的：
$$\int_1^2 \frac{x^2 + 2x}{x} \, dx$$

上面是一长串，下面是一个 $x$。
怎么办？

\textbf{还记得加法的规律吗？拆开！}

把这个大分数，拆成两个小分数：
$$\frac{x^2 + 2x}{x} = \frac{x^2}{x} + \frac{2x}{x}$$

就像你有两堆糖（一堆 $x^2$ 个，一堆 $2x$ 个），都要分给 $x$ 个人。
你可以先分第一堆，再分第二堆。

现在，分别对它们进行\textbf{抵消（约分）}：
第一堆：$\frac{x^2}{x} = x$
第二堆：$\frac{2x}{x} = 2$（上面的 $x$ 和下面的 $x$ 划掉，只剩 2）

所以，原来那个可怕的分数，变成了：
$$x + 2$$

这道题就变成了：
$$\int_1^2 (x + 2) \, dx$$

加法！拆开分别算：
\begin{itemize}
\item $\int_1^2 x \, dx$ $\rightarrow$ $x$ 的反过来是 $\frac{x^2}{2}$，代入得 $2 - 0.5 = 1.5$。
\item $\int_1^2 2 \, dx$ $\rightarrow$ 常数 $2$ 不动，总长度是 $1$，所以是 $2 \times 1 = 2$。
\item 加起来：$1.5 + 2 = 3.5$。
\end{itemize}

\textbf{答案：$3.5$}

\subsection*{终极魔王：无法抵消的 $\frac{1}{x}$（宇宙大爆炸）}

刚才的题，要么能提出来，要么能抵消。都是纸老虎。

但总有一些题，是\textbf{真老虎}。

比如这个：
$$\int_1^2 \frac{1}{x} \, dx$$

上面是 $1$，下面是 $x$。
不能提出来（因为 $x$ 在动）。
也不能抵消（上面没有 $x$ 给你划掉）。

只能老老实实用\textbf{新方法}：找反过来。

\subsubsection*{试着找反过来}

回忆一下我们的万能规律：
\textbf{找反过来 = 指数加 1，除以新指数。}

那 $\frac{1}{x}$ 的指数是多少？
在数学里，除以 $x$，就等于把 $x$ 放到负数的楼层去。
$$\frac{1}{x} = x^{-1}$$

好，现在对 $x^{-1}$ 用我们的规律：

第一步：指数加 1。
$-1 + 1 = 0$。
于是变成了 $x^0$。

第二步：除以新指数。
新指数是 $0$，所以我们要除以 $0$。
变成了 $\frac{x^0}{0}$。

\subsubsection*{等等，警报响了！🚨}

$\frac{x^0}{0}$ ？？？

在数学里，有一条绝对不能碰的铁律：
\textbf{“永远，永远，永远不要除以 0！”}

想象你有 5 个苹果，分给 0 个人。每人得几个？
这根本说不通！
一旦除以 0，数学宇宙就会直接大爆炸，计算器会直接报错（Error）。

\textbf{我们的“万能规律”在这里崩溃了。}
指数加 1 再除以新指数的方法，对 $x^{-1}$ 完全不管用！

怎么办？难道 $\frac{1}{x}$ 没有“反过来”吗？

\subsubsection*{英雄登场：VIP 贵宾 $\ln x$}

数学家们发现规律崩溃后，也急得团团转。
他们翻遍了所有的函数，终于找到了一个神奇的东西。

这个东西，切一刀之后，竟然正好变成了 $\frac{1}{x}$！

它是谁？
它叫 \textbf{自然对数}，符号是 \textbf{$\ln x$}（读作“lon x”）。

不要被这个奇怪的名字吓到。
就像我们开头说的，\textbf{名字只是名字}。

你不需要管 $\ln x$ 到底长什么样，你只需要记住一条死规矩：
\begin{quote}
\textbf{谁切一刀会变成 $\frac{1}{x}$？}

\textbf{答案是：$\ln x$。}
\end{quote}

因为 $\frac{1}{x}$ 的身份太特殊了（它打破了指数规律），所以数学家给它的反过来发了一张“VIP 贵宾卡”，专门起了个名字叫 $\ln x$。

\subsubsection*{用英雄来算题}

现在我们能算这道题了：
$$\int_1^2 \frac{1}{x} \, dx$$

\textbf{第一步：找 $\frac{1}{x}$ 的反过来。}
万能规律失效，请出 VIP 贵宾。
反过来是：\textbf{$\ln x$}。

\textbf{第二步：代入上面的数 $2$。}
变成 \textbf{$\ln 2$}。
（$\ln 2$ 其实是个无限不循环的丑小数，大约是 0.693...，但数学家很懒，他们觉得写小数太丑了，就直接保留 $\ln 2$ 这个名字，就像保留 $\pi$ 一样。）

\textbf{第三步：代入下面的数 $1$。}
变成 \textbf{$\ln 1$}。
（这里有个小常识：$\ln 1$ 永远等于 $0$。就像任何数的 0 次方都是 1 一样，你记住就好。）

\textbf{第四步：上面减下面。}
$$\ln 2 - \ln 1 = \ln 2 - 0 = \ln 2$$

\textbf{答案：$\ln 2$}

\subsection*{一句话总结}

遇到微积分里的除法（分数）怎么办？

\textbf{三板斧：}

\begin{enumerate}
\item \textbf{下面是数字？} $\rightarrow$ 当成常数提出来！（比如 $\frac{x}{2}$ 提个 $\frac{1}{2}$ 出来）
\item \textbf{上下都有字母？} $\rightarrow$ 能抵消就抵消！（比如 $\frac{x^2}{x} = x$）
\item \textbf{抵消不了的 $\frac{1}{x}$ 怎么办？} $\rightarrow$ 万能规律失效，直接呼叫 VIP 英雄 \textbf{$\ln x$}。
\end{enumerate}

你看，除法并不可怕。
绝大多数时候，它只要划几下就能变成我们熟悉的加法和乘法。
就算是遇到大魔王，我们也已经有了专属的魔法武器。

微积分里的基本运算，你现在已经\textbf{全部打通关了！}

\section*{二重积分中的除法：别怕，套路是一样的}

\bigskip\hrule\bigskip

\subsection*{📌 先来一个灵魂拷问}

在一元微积分里，你已经掌握了对付除法的三板斧：
\begin{itemize}
\item 下面是数字？提出来！
\item 上下都有字母？约分！
\item 约不掉的 $\frac{1}{x}$？请出 VIP 英雄 $\ln x$！
\end{itemize}

现在问题来了：

\textbf{“二重积分里有除法，怎么办？”}

别慌。你猜怎么着？

\textbf{套路完全一样，只不过要用两次。}

一元积分是爬一层楼，二重积分是爬两层楼。
但每一层楼的楼梯，长得一模一样。

\subsection*{第一种：纸老虎（下面是数字）}

来看这道题：
$$\int_0^2 \int_0^1 \frac{xy}{3} \, dy \, dx$$

翻译一下：我们要算一个立体的体积。
\begin{itemize}
\item $x$ 从 $0$ 走到 $2$
\item $y$ 从 $0$ 走到 $1$
\item 每个位置的高度是 $\frac{xy}{3}$
\end{itemize}

\subsubsection*{先别急着积分，看看能不能化简}

$\frac{xy}{3}$ 是什么意思？

就是 $xy$ 除以 $3$，也就是 $\frac{1}{3} \times xy$。

$\frac{1}{3}$ 是什么？

\textbf{常数！永远不动！}

不管你积 $x$ 还是积 $y$，这个 $\frac{1}{3}$ 都纹丝不动。

既然它不动，那就把它\textbf{提到最外面去}，别让它碍事：

$$\int_0^2 \int_0^1 \frac{xy}{3} \, dy \, dx = \frac{1}{3} \int_0^2 \int_0^1 xy \, dy \, dx$$

\subsubsection*{现在来算里面的部分}

剩下的 $\int_0^2 \int_0^1 xy \, dy \, dx$，就是我们熟悉的乘法问题了！

\textbf{第一步：先算内层（对 $y$ 积分）}

内层是：$\int_0^1 xy \, dy$

在这一步里，$x$ 是个“旁观者”，它不动。
$y$ 才是主角，从 $0$ 走到 $1$。

不动的 $x$ 可以提出来：
$$\int_0^1 xy \, dy = x \int_0^1 y \, dy$$

现在算 $\int_0^1 y \, dy$：
\begin{itemize}
\item $y$ 的反过来是 $\frac{y^2}{2}$
\item 代入上面的 $1$：$\frac{1^2}{2} = \frac{1}{2}$
\item 代入下面的 $0$：$\frac{0^2}{2} = 0$
\item 相减：$\frac{1}{2} - 0 = \frac{1}{2}$
\end{itemize}

所以内层的结果是：$x \times \frac{1}{2} = \frac{x}{2}$

\textbf{第二步：再算外层（对 $x$ 积分）}

现在题目变成了：$\int_0^2 \frac{x}{2} \, dx$

$\frac{x}{2}$ 是什么？就是 $\frac{1}{2} \times x$。

$\frac{1}{2}$ 不动，提出来：
$$\int_0^2 \frac{x}{2} \, dx = \frac{1}{2} \int_0^2 x \, dx$$

现在算 $\int_0^2 x \, dx$：
\begin{itemize}
\item $x$ 的反过来是 $\frac{x^2}{2}$
\item 代入上面的 $2$：$\frac{2^2}{2} = \frac{4}{2} = 2$
\item 代入下面的 $0$：$\frac{0^2}{2} = 0$
\item 相减：$2 - 0 = 2$
\end{itemize}

所以外层的结果是：$\frac{1}{2} \times 2 = 1$

\textbf{第三步：别忘了最开始提出来的 $\frac{1}{3}$}

最终答案：$\frac{1}{3} \times 1 = \frac{1}{3}$

\textbf{答案：$\boxed{\frac{1}{3}}$}

\begin{quote}
\textbf{📌 总结：下面是数字}

跟一元积分一模一样：
\begin{itemize}
\item 除以数字 $\rightarrow$ 等于乘以它的倒数 $\rightarrow$ 变成常数
\item 常数提到最外面 $\rightarrow$ 最后再乘回去
\end{itemize}

\textbf{完全不用怕！}
\end{quote}

\subsection*{第二种：假装很凶的纸老虎（可以约分）}

来个看起来很吓人的：
$$\int_1^2 \int_1^3 \frac{x^2 y}{x} \, dy \, dx$$

上面是 $x^2 y$，下面是 $x$，看着挺复杂。

\textbf{但是！先别急着积分！}

遇到分数，第一反应是什么？

\textbf{看看能不能约分！}

\subsubsection*{先化简}

$\frac{x^2 y}{x}$ 这个分数，仔细看：
\begin{itemize}
\item 上面有 $x^2$，也就是 $x \times x$
\item 下面有一个 $x$
\end{itemize}

上面有两个 $x$，下面有一个 $x$。
就像“两个苹果除以一个苹果”，可以抵消掉一个！

$$\frac{x^2 y}{x} = \frac{x \cdot x \cdot y}{x} = x \cdot y = xy$$

上下各划掉一个 $x$，剩下 $xy$。

\subsubsection*{题目瞬间变简单了}

$$\int_1^2 \int_1^3 \frac{x^2 y}{x} \, dy \, dx = \int_1^2 \int_1^3 xy \, dy \, dx$$

这不就是最普通的乘法问题吗！

\textbf{第一步：先算内层（对 $y$ 积分）}

内层是：$\int_1^3 xy \, dy$

$x$ 不动，提出来：$x \int_1^3 y \, dy$

算 $\int_1^3 y \, dy$：
\begin{itemize}
\item $y$ 的反过来是 $\frac{y^2}{2}$
\item 代入 $3$：$\frac{9}{2}$
\item 代入 $1$：$\frac{1}{2}$
\item 相减：$\frac{9}{2} - \frac{1}{2} = \frac{8}{2} = 4$
\end{itemize}

内层结果：$x \times 4 = 4x$

\textbf{第二步：再算外层（对 $x$ 积分）}

$$\int_1^2 4x \, dx$$

$4$ 不动，提出来：$4 \int_1^2 x \, dx$

算 $\int_1^2 x \, dx$：
\begin{itemize}
\item $x$ 的反过来是 $\frac{x^2}{2}$
\item 代入 $2$：$\frac{4}{2} = 2$
\item 代入 $1$：$\frac{1}{2}$
\item 相减：$2 - \frac{1}{2} = \frac{3}{2}$
\end{itemize}

外层结果：$4 \times \frac{3}{2} = 6$

\textbf{答案：$\boxed{6}$}

\begin{quote}
\textbf{📌 总结：上下都有字母}

遇到这种，第一反应是：\textbf{能不能约分？}

把上面的 $x$ 和下面的 $x$ 划掉。

约完了再算，轻松很多！
\end{quote}

\subsection*{第三种：混合在一起的纸老虎}

再来一个更复杂的：
$$\int_1^2 \int_0^1 \frac{x^2 + xy}{x} \, dy \, dx$$

上面是一长串 $x^2 + xy$，下面是 $x$。

怎么办？

\subsubsection*{先化简：拆开再约分}

\textbf{还记得加法的规律吗？拆开！}

把这个大分数，拆成两个小分数：
$$\frac{x^2 + xy}{x} = \frac{x^2}{x} + \frac{xy}{x}$$

就像你有两堆糖（一堆 $x^2$ 个，一堆 $xy$ 个），都要分给 $x$ 个人。
你可以先分第一堆，再分第二堆。

现在分别\textbf{约分}：

第一堆：$\frac{x^2}{x} = x$（上下各划掉一个 $x$）

第二堆：$\frac{xy}{x} = y$（上下各划掉一个 $x$，只剩 $y$）

所以：
$$\frac{x^2 + xy}{x} = x + y$$

\subsubsection*{题目变成了加法}

$$\int_1^2 \int_0^1 \frac{x^2 + xy}{x} \, dy \, dx = \int_1^2 \int_0^1 (x + y) \, dy \, dx$$

\textbf{加法怎么办？拆开分别算！}

$$\int_1^2 \int_0^1 x \, dy \, dx + \int_1^2 \int_0^1 y \, dy \, dx$$

\subsubsection*{第一块：$\int_1^2 \int_0^1 x \, dy \, dx$}

\textbf{内层：$\int_0^1 x \, dy$}

注意！这里是对 $y$ 积分，但被积的东西是 $x$。

$x$ 跟 $y$ 没关系，$x$ 就是个“旁观者”，完全不动。

这就像问你：$\int_0^1 3 \, dy$ 等于多少？
常数 $3$ 乘以区间长度 $1$，等于 $3$。

同样的道理：
$$\int_0^1 x \, dy = x \times (1 - 0) = x \times 1 = x$$

\textbf{外层：$\int_1^2 x \, dx$}

\begin{itemize}
\item $x$ 的反过来是 $\frac{x^2}{2}$
\item 代入 $2$：$\frac{4}{2} = 2$
\item 代入 $1$：$\frac{1}{2}$
\item 相减：$2 - \frac{1}{2} = \frac{3}{2} = 1.5$
\end{itemize}

\textbf{第一块的结果：$1.5$}

\subsubsection*{第二块：$\int_1^2 \int_0^1 y \, dy \, dx$}

\textbf{内层：$\int_0^1 y \, dy$}

\begin{itemize}
\item $y$ 的反过来是 $\frac{y^2}{2}$
\item 代入 $1$：$\frac{1}{2}$
\item 代入 $0$：$0$
\item 相减：$\frac{1}{2}$
\end{itemize}

内层结果：$\frac{1}{2}$（这是一个数字，不含任何字母！）

\textbf{外层：$\int_1^2 \frac{1}{2} \, dx$}

$\frac{1}{2}$ 是常数，乘以区间长度：
$$\frac{1}{2} \times (2 - 1) = \frac{1}{2} \times 1 = \frac{1}{2} = 0.5$$

\textbf{第二块的结果：$0.5$}

\subsubsection*{两块加起来}

$$1.5 + 0.5 = 2$$

\textbf{答案：$\boxed{2}$}

\begin{quote}
\textbf{📌 总结：混合型}

\begin{enumerate}
\item 先把分数拆开
\item 分别约分
\item 变成加法问题
\item 拆开分别算，最后加起来
\end{enumerate}
\end{quote}

\subsection*{终极魔王：无法约分的 $\frac{1}{x}$ 或 $\frac{1}{y}$}

前面的题，要么能提出来，要么能约分。都是纸老虎。

但现在，来个\textbf{真老虎}：

$$\int_1^2 \int_1^3 \frac{y}{x} \, dy \, dx$$

\subsubsection*{先看看能不能化简}

$\frac{y}{x}$ 能约分吗？

上面是 $y$，下面是 $x$。

不一样的字母，\textbf{没法约！}

能提出来吗？

$x$ 在动，$y$ 也在动，谁也提不出去。

怎么办？

\subsubsection*{别慌，换个角度看}

虽然 $\frac{y}{x}$ 是个分数，但你可以把它\textbf{拆成乘法}：

$$\frac{y}{x} = y \times \frac{1}{x}$$

你看，这其实是两个东西\textbf{乘在一起}！

一个是 $y$，一个是 $\frac{1}{x}$。

\subsubsection*{乘法怎么处理？谁不动提谁！}

现在我们有选择：先积 $y$，还是先积 $x$？

\textbf{如果先积 $y$：}
\begin{itemize}
\item 在对 $y$ 积分的时候，$x$ 是旁观者，$\frac{1}{x}$ 不动
\item 可以把 $\frac{1}{x}$ 提出来！
\end{itemize}

\textbf{如果先积 $x$：}
\begin{itemize}
\item 在对 $x$ 积分的时候，$y$ 是旁观者，$y$ 不动
\item 可以把 $y$ 提出来！
\end{itemize}

两条路都能走。题目里内层是 $dy$，所以我们先积 $y$。

\subsubsection*{第一步：先算内层（对 $y$ 积分）}

内层是：$\int_1^3 \frac{y}{x} \, dy = \int_1^3 y \cdot \frac{1}{x} \, dy$

在这一步里，$x$ 是旁观者。$\frac{1}{x}$ 就像一个常数，可以提出来：

$$\frac{1}{x} \int_1^3 y \, dy$$

现在算 $\int_1^3 y \, dy$：
\begin{itemize}
\item $y$ 的反过来是 $\frac{y^2}{2}$
\item 代入 $3$：$\frac{9}{2}$
\item 代入 $1$：$\frac{1}{2}$
\item 相减：$\frac{9}{2} - \frac{1}{2} = \frac{8}{2} = 4$
\end{itemize}

内层结果：$\frac{1}{x} \times 4 = \frac{4}{x}$

\subsubsection*{第二步：再算外层（对 $x$ 积分）}

现在题目变成了：

$$\int_1^2 \frac{4}{x} \, dx$$

$\frac{4}{x}$ 是什么？就是 $4 \times \frac{1}{x}$。

$4$ 是常数，提出来：

$$4 \int_1^2 \frac{1}{x} \, dx$$

\subsubsection*{第三步：对付大魔王 $\frac{1}{x}$}

现在我们要算 $\int_1^2 \frac{1}{x} \, dx$。

这就是一元积分里的那个大魔王！

还记得吗？

我们试着用万能规律：$\frac{1}{x} = x^{-1}$，指数加 1 得到 $0$，然后要除以 $0$……

\textbf{宇宙爆炸！万能规律失效！}

怎么办？

\textbf{请出 VIP 英雄：$\ln x$！}

谁切一刀会变成 $\frac{1}{x}$？答案是 $\ln x$。

所以：
$$\int_1^2 \frac{1}{x} \, dx = \ln x \Big|_1^2 = \ln 2 - \ln 1 = \ln 2 - 0 = \ln 2$$

\subsubsection*{最终答案}

外层结果是：$4 \times \ln 2 = 4\ln 2$

\textbf{答案：$\boxed{4\ln 2}$}

（$\ln 2$ 大约等于 $0.693$，所以 $4\ln 2 \approx 2.77$。但数学家喜欢保留 $\ln 2$ 这个写法，就像保留 $\pi$ 一样。）

\subsection*{再来一个魔王：$\frac{1}{xy}$}

趁热打铁，再来一个：

$$\int_1^e \int_1^e \frac{1}{xy} \, dy \, dx$$

（$e$ 是那个著名的数学常数，大约等于 $2.718$。）

\begin{quote}
\textbf{🎯 关于 $\ln$ 和 $e$ 的死规矩}

你不需要理解为什么，只需要记住这几条：

\begin{center}
\begin{tabular}{|c|c|}
\hline
规矩 & 意思 \\
\hline
$\ln 1 = 0$ & $1$ 的对数永远是 $0$ \\
\hline
$\ln e = 1$ & $e$ 的对数永远是 $1$ \\
\hline
$\int \frac{1}{x} dx = \ln x$ & $\frac{1}{x}$ 的反过来是 $\ln x$ \\
\hline
\end{tabular}
\end{center}

就像你记住 $\sin 0 = 0$、$\cos 0 = 1$ 一样，这些是基础常识。
\end{quote}

\begin{quote}
\textbf{⚠️ 超级重要的提醒：别忘了要减！}

很多同学算 $\ln$ 的积分时会犯一个错：

看到 $\ln x \Big|_1^e$，直接写 $\ln e = 1$，然后就结束了。

\textbf{错！大错特错！}

定积分永远是“\textbf{代入上面的 − 代入下面的}”。

正确做法：
$$\ln x \Big|_1^e = \ln e - \ln 1 = 1 - 0 = 1$$

这次碰巧 $\ln 1 = 0$，减了跟没减一样。

但如果是 $\ln x \Big|_2^e$，那就是 $\ln e - \ln 2 = 1 - \ln 2$，\textbf{不是} $1$！

\textbf{养成好习惯：每次都写完整的“上面 − 下面”！}
\end{quote}

\subsubsection*{先看看能不能化简}

$\frac{1}{xy}$ 能约分吗？

上面是 $1$，什么也约不掉。

但是！我们可以把它\textbf{拆成乘法}：

$$\frac{1}{xy} = \frac{1}{x} \times \frac{1}{y}$$

就像 $\frac{1}{6} = \frac{1}{2} \times \frac{1}{3}$ 一样。

\subsubsection*{第一步：先算内层（对 $y$ 积分）}

内层是：$\int_1^e \frac{1}{xy} \, dy = \int_1^e \frac{1}{x} \cdot \frac{1}{y} \, dy$

在这一步里，$x$ 是旁观者。$\frac{1}{x}$ 不动，提出来：

$$\frac{1}{x} \int_1^e \frac{1}{y} \, dy$$

现在要算 $\int_1^e \frac{1}{y} \, dy$。

又是大魔王！但我们已经有武器了——请出 VIP 英雄 $\ln y$：

$$\int_1^e \frac{1}{y} \, dy = \ln y \Big|_1^e = \ln e - \ln 1 = 1 - 0 = 1$$

内层结果：$\frac{1}{x} \times 1 = \frac{1}{x}$

\subsubsection*{第二步：再算外层（对 $x$ 积分）}

现在题目变成了：

$$\int_1^e \frac{1}{x} \, dx$$

又是大魔王！再请一次 VIP 英雄：

$$\int_1^e \frac{1}{x} \, dx = \ln x \Big|_1^e = \ln e - \ln 1 = 1 - 0 = 1$$

\textbf{答案：$\boxed{1}$}

\subsubsection*{回头看看发生了什么}

这道题，我们用了\textbf{两次} $\ln$：
\begin{itemize}
\item 对 $y$ 积分时，$\frac{1}{y}$ 的反过来是 $\ln y$
\item 对 $x$ 积分时，$\frac{1}{x}$ 的反过来是 $\ln x$
\end{itemize}

每一层楼的楼梯，长得一模一样。
你在一元积分里学会的武器，在二重积分里\textbf{照样管用}，只不过要用两次。

\subsection*{再来一个变形：$\frac{x}{y}$}

$$\int_1^2 \int_1^4 \frac{x}{y} \, dy \, dx$$

\subsubsection*{化简}

$\frac{x}{y}$ 能约分吗？不能，字母不一样。

但可以拆成乘法：
$$\frac{x}{y} = x \times \frac{1}{y}$$

\subsubsection*{第一步：先算内层（对 $y$ 积分）}

内层是：$\int_1^4 \frac{x}{y} \, dy = \int_1^4 x \cdot \frac{1}{y} \, dy$

$x$ 不动，提出来：
$$x \int_1^4 \frac{1}{y} \, dy$$

算 $\int_1^4 \frac{1}{y} \, dy$：
$$\ln y \Big|_1^4 = \ln 4 - \ln 1 = \ln 4 - 0 = \ln 4$$

内层结果：$x \times \ln 4 = x \ln 4$

\subsubsection*{第二步：再算外层（对 $x$ 积分）}

$$\int_1^2 x \ln 4 \, dx$$

$\ln 4$ 是常数（大约 $1.386$），提出来：
$$\ln 4 \int_1^2 x \, dx$$

算 $\int_1^2 x \, dx$：
\begin{itemize}
\item $x$ 的反过来是 $\frac{x^2}{2}$
\item 代入 $2$：$\frac{4}{2} = 2$
\item 代入 $1$：$\frac{1}{2}$
\item 相减：$2 - \frac{1}{2} = \frac{3}{2}$
\end{itemize}

外层结果：$\ln 4 \times \frac{3}{2} = \frac{3}{2} \ln 4$

\textbf{答案：$\boxed{\frac{3}{2} \ln 4}$}

\subsection*{📌 终极总结：二重积分除法的完整攻略}

\begin{center}
\begin{tabular}{|c|c|c|}
\hline
情况 & 怎么办 & 例子 \\
\hline
\textbf{下面是数字} & 当成常数提到最外面 & $\frac{xy}{3} \to \frac{1}{3} \times xy$ \\
\hline
\textbf{上下有同一个字母} & 先约分！能划的先划掉 & $\frac{x^2 y}{x} \to xy$ \\
\hline
\textbf{上面是加法} & 拆成几个分数，分别约分 & $\frac{x^2+xy}{x} \to x + y$ \\
\hline
\textbf{约不掉的 $\frac{1}{x}$ 或 $\frac{1}{y}$} & 请出 VIP 英雄 $\ln$ & $\int \frac{1}{x} dx = \ln x$ \\
\hline
\end{tabular}
\end{center}

\begin{quote}
\textbf{🎯 关于 $\ln$ 的死规矩（考试救命用）}

\begin{center}
\begin{tabular}{|c|c|}
\hline
规矩 & 意思 \\
\hline
$\ln 1 = 0$ & $1$ 的对数永远是 $0$ \\
\hline
$\ln e = 1$ & $e$ 的对数永远是 $1$ \\
\hline
$\int \frac{1}{x} dx = \ln x$ & $\frac{1}{x}$ 的反过来是 $\ln x$ \\
\hline
\end{tabular}
\end{center}

\textbf{⚠️ 最容易犯的错：忘记“上面 − 下面”！}

正确写法：$\ln x \Big|_1^e = \ln e - \ln 1 = 1 - 0 = 1$

错误写法：$\ln x \Big|_1^e = \ln e = 1$（漏掉了减 $\ln 1$）

这次碰巧对了，下次就错了！\textbf{永远写完整！}
\end{quote}

\subsection*{🎯 最核心的思想}

\textbf{先化简，再用“乘法”的规则处理。}

$\frac{y}{x}$ 看起来是除法，但你把它想成 $y \times \frac{1}{x}$，它就变成了\textbf{乘法}。

乘法怎么处理？\textbf{谁不动提谁！}

提出来之后，剩下的那个 $\frac{1}{x}$ 或 $\frac{1}{y}$，就是一元积分里的老熟人了——召唤 $\ln$ 来救场。

你看，二重积分的除法，\textbf{本质上就是一元积分除法的“两遍版”}。

你已经打败过一次这个 boss，现在不过是\textbf{再打一遍}而已。

同样的招式，用两次，搞定！

\section*{三重积分中的除法？放心，依然是那三板斧}

\bigskip\hrule\bigskip

\subsection*{📌 你可能在担心什么}

刚学完二重积分的除法，你心里可能犯嘀咕：

“二重积分要算两层，我已经有点晕了……”

“三重积分要算三层，除法会不会更复杂？”

“会不会有什么新的、可怕的东西冒出来？”

我现在就给你吃一颗定心丸：

\textbf{不会。}

\textbf{完全不会。}

\textbf{三重积分的除法，跟二重积分的除法，跟一元积分的除法，用的是同一套武功。}

唯一的区别是什么？

一元积分：爬 1 层楼。

二重积分：爬 2 层楼。

三重积分：爬 3 层楼。

\textbf{但每一层楼的楼梯，长得一模一样。}

你已经学会了爬楼梯的方法，现在不过是多爬一层而已。

\subsection*{第一种：纸老虎（下面是数字）}

来看这道题：
$$\int_0^1 \int_0^2 \int_0^3 \frac{xyz}{6} \, dz \, dy \, dx$$

翻译一下：
\begin{itemize}
\item $x$ 从 $0$ 走到 $1$
\item $y$ 从 $0$ 走到 $2$
\item $z$ 从 $0$ 走到 $3$
\item 要把 $\frac{xyz}{6}$ 这个东西，在整个区域里累加起来
\end{itemize}

\subsubsection*{老规矩：先化简}

$\frac{xyz}{6}$ 是什么？

就是 $xyz$ 除以 $6$，也就是 $\frac{1}{6} \times xyz$。

$\frac{1}{6}$ 是常数，永远不动。

不管你积 $x$、积 $y$、还是积 $z$，它都纹丝不动。

\textbf{提到最外面去！}

$$\int_0^1 \int_0^2 \int_0^3 \frac{xyz}{6} \, dz \, dy \, dx = \frac{1}{6} \int_0^1 \int_0^2 \int_0^3 xyz \, dz \, dy \, dx$$

\subsubsection*{现在来算里面的部分}

剩下的是 $\int_0^1 \int_0^2 \int_0^3 xyz \, dz \, dy \, dx$。

三个字母乘在一起，怎么办？

\textbf{还是那句话：谁不动提谁！}

\textbf{第一步：先算最内层（对 $z$ 积分）}

最内层是：$\int_0^3 xyz \, dz$

在这一步里，$z$ 从 $0$ 走到 $3$，它是主角。

$x$ 和 $y$ 呢？它们是旁观者，完全不动。

不动的 $xy$ 可以提出来：
$$\int_0^3 xyz \, dz = xy \int_0^3 z \, dz$$

现在算 $\int_0^3 z \, dz$：
\begin{itemize}
\item $z$ 的反过来是 $\frac{z^2}{2}$
\item 代入 $3$：$\frac{9}{2}$
\item 代入 $0$：$0$
\item 相减：$\frac{9}{2} - 0 = \frac{9}{2}$
\end{itemize}

最内层结果：$xy \times \frac{9}{2} = \frac{9xy}{2}$

\textbf{第二步：再算中间层（对 $y$ 积分）}

现在题目变成了：$\int_0^1 \int_0^2 \frac{9xy}{2} \, dy \, dx$

中间层是：$\int_0^2 \frac{9xy}{2} \, dy$

$\frac{9}{2}$ 是常数，$x$ 也是旁观者（因为这一步积的是 $y$）。

把 $\frac{9x}{2}$ 提出来：
$$\int_0^2 \frac{9xy}{2} \, dy = \frac{9x}{2} \int_0^2 y \, dy$$

现在算 $\int_0^2 y \, dy$：
\begin{itemize}
\item $y$ 的反过来是 $\frac{y^2}{2}$
\item 代入 $2$：$\frac{4}{2} = 2$
\item 代入 $0$：$0$
\item 相减：$2 - 0 = 2$
\end{itemize}

中间层结果：$\frac{9x}{2} \times 2 = 9x$

\textbf{第三步：最后算外层（对 $x$ 积分）}

现在题目变成了：$\int_0^1 9x \, dx$

$9$ 是常数，提出来：$9 \int_0^1 x \, dx$

现在算 $\int_0^1 x \, dx$：
\begin{itemize}
\item $x$ 的反过来是 $\frac{x^2}{2}$
\item 代入 $1$：$\frac{1}{2}$
\item 代入 $0$：$0$
\item 相减：$\frac{1}{2} - 0 = \frac{1}{2}$
\end{itemize}

外层结果：$9 \times \frac{1}{2} = \frac{9}{2}$

\textbf{第四步：别忘了最开始提出来的 $\frac{1}{6}$}

最终答案：$\frac{1}{6} \times \frac{9}{2} = \frac{9}{12} = \frac{3}{4}$

\textbf{答案：$\boxed{\frac{3}{4}}$}

你看，跟二重积分有什么区别？

\textbf{没区别。就是多算了一层。}

\subsection*{第二种：假装很凶的纸老虎（可以约分）}

来个吓人的：
$$\int_1^2 \int_1^2 \int_1^3 \frac{x^2 yz}{x} \, dz \, dy \, dx$$

上面是 $x^2 yz$，下面是 $x$，看着挺复杂。

\subsubsection*{先别急！看看能不能约分}

$\frac{x^2 yz}{x}$ 这个分数：
\begin{itemize}
\item 上面有 $x^2$，也就是 $x \times x$
\item 下面有一个 $x$
\end{itemize}

上下各划掉一个 $x$：
$$\frac{x^2 yz}{x} = xyz$$

\textbf{题目瞬间变简单了！}

$$\int_1^2 \int_1^2 \int_1^3 \frac{x^2 yz}{x} \, dz \, dy \, dx = \int_1^2 \int_1^2 \int_1^3 xyz \, dz \, dy \, dx$$

\subsubsection*{现在一层一层算}

\textbf{第一步：最内层（对 $z$ 积分）}

$\int_1^3 xyz \, dz = xy \int_1^3 z \, dz$

算 $\int_1^3 z \, dz$：
\begin{itemize}
\item $z$ 的反过来是 $\frac{z^2}{2}$
\item 代入 $3$：$\frac{9}{2}$
\item 代入 $1$：$\frac{1}{2}$
\item 相减：$\frac{9}{2} - \frac{1}{2} = \frac{8}{2} = 4$
\end{itemize}

最内层结果：$xy \times 4 = 4xy$

\textbf{第二步：中间层（对 $y$ 积分）}

$\int_1^2 4xy \, dy = 4x \int_1^2 y \, dy$

算 $\int_1^2 y \, dy$：
\begin{itemize}
\item 代入 $2$：$\frac{4}{2} = 2$
\item 代入 $1$：$\frac{1}{2}$
\item 相减：$2 - \frac{1}{2} = \frac{3}{2}$
\end{itemize}

中间层结果：$4x \times \frac{3}{2} = 6x$

\textbf{第三步：外层（对 $x$ 积分）}

$\int_1^2 6x \, dx = 6 \int_1^2 x \, dx$

算 $\int_1^2 x \, dx$：
\begin{itemize}
\item 代入 $2$：$2$
\item 代入 $1$：$\frac{1}{2}$
\item 相减：$2 - \frac{1}{2} = \frac{3}{2}$
\end{itemize}

外层结果：$6 \times \frac{3}{2} = 9$

\textbf{答案：$\boxed{9}$}

\subsection*{终极魔王：无法约分的 $\frac{1}{x}$、$\frac{1}{y}$、$\frac{1}{z}$}

来个真正的挑战：
$$\int_1^e \int_1^e \int_1^e \frac{1}{xyz} \, dz \, dy \, dx$$

\subsubsection*{能约分吗？}

上面是 $1$，什么也约不掉。

但可以拆成乘法：
$$\frac{1}{xyz} = \frac{1}{x} \times \frac{1}{y} \times \frac{1}{z}$$

三个东西乘在一起！

\begin{quote}
\textbf{⚠️ 再次提醒：别忘了“上面 − 下面”！}

每次用 $\ln$ 算定积分，都要写完整：
$$\ln x \Big|_1^e = \ln e - \ln 1 = 1 - 0 = 1$$

不要偷懒只写 $\ln e = 1$！

\textbf{养成好习惯，考试不丢分！}
\end{quote}

\subsubsection*{第一步：最内层（对 $z$ 积分）}

$\int_1^e \frac{1}{xyz} \, dz = \int_1^e \frac{1}{x} \cdot \frac{1}{y} \cdot \frac{1}{z} \, dz$

$\frac{1}{x}$ 和 $\frac{1}{y}$ 都不动，提出来：

$$\frac{1}{x} \cdot \frac{1}{y} \int_1^e \frac{1}{z} \, dz = \frac{1}{xy} \int_1^e \frac{1}{z} \, dz$$

现在算 $\int_1^e \frac{1}{z} \, dz$：

请出 VIP 英雄 $\ln z$：
$$\ln z \Big|_1^e = \ln e - \ln 1 = 1 - 0 = 1$$

最内层结果：$\frac{1}{xy} \times 1 = \frac{1}{xy}$

\subsubsection*{第二步：中间层（对 $y$ 积分）}

$\int_1^e \frac{1}{xy} \, dy = \frac{1}{x} \int_1^e \frac{1}{y} \, dy$

算 $\int_1^e \frac{1}{y} \, dy$：
$$\ln y \Big|_1^e = \ln e - \ln 1 = 1 - 0 = 1$$

中间层结果：$\frac{1}{x} \times 1 = \frac{1}{x}$

\subsubsection*{第三步：外层（对 $x$ 积分）}

$\int_1^e \frac{1}{x} \, dx$

算：
$$\ln x \Big|_1^e = \ln e - \ln 1 = 1 - 0 = 1$$

\textbf{答案：$\boxed{1}$}

\subsubsection*{回头看看发生了什么}

这道题，我们用了\textbf{三次} $\ln$：
\begin{itemize}
\item 对 $z$ 积分时，$\frac{1}{z}$ 的反过来是 $\ln z$
\item 对 $y$ 积分时，$\frac{1}{y}$ 的反过来是 $\ln y$
\item 对 $x$ 积分时，$\frac{1}{x}$ 的反过来是 $\ln x$
\end{itemize}

\textbf{每一层都用同样的武器，同样的招式。}

三重积分？不过是把一元积分的方法\textbf{用三遍}而已。

\subsection*{来个混合型：$\frac{yz}{x}$}

$$\int_1^2 \int_1^3 \int_0^4 \frac{yz}{x} \, dz \, dy \, dx$$

\subsubsection*{化简}

$\frac{yz}{x}$ 能约分吗？

上面是 $yz$，下面是 $x$。字母不一样，约不了。

但可以拆成乘法：
$$\frac{yz}{x} = y \times z \times \frac{1}{x}$$

\subsubsection*{第一步：最内层（对 $z$ 积分）}

$\int_0^4 \frac{yz}{x} \, dz = \int_0^4 y \cdot z \cdot \frac{1}{x} \, dz$

$y$ 和 $\frac{1}{x}$ 都不动，提出来：

$$\frac{y}{x} \int_0^4 z \, dz$$

算 $\int_0^4 z \, dz$：
\begin{itemize}
\item $z$ 的反过来是 $\frac{z^2}{2}$
\item 代入 $4$：$\frac{16}{2} = 8$
\item 代入 $0$：$0$
\item 相减：$8 - 0 = 8$
\end{itemize}

最内层结果：$\frac{y}{x} \times 8 = \frac{8y}{x}$

\subsubsection*{第二步：中间层（对 $y$ 积分）}

$\int_1^3 \frac{8y}{x} \, dy = \frac{8}{x} \int_1^3 y \, dy$

算 $\int_1^3 y \, dy$：
\begin{itemize}
\item 代入 $3$：$\frac{9}{2}$
\item 代入 $1$：$\frac{1}{2}$
\item 相减：$\frac{9}{2} - \frac{1}{2} = 4$
\end{itemize}

中间层结果：$\frac{8}{x} \times 4 = \frac{32}{x}$

\subsubsection*{第三步：外层（对 $x$ 积分）}

$\int_1^2 \frac{32}{x} \, dx = 32 \int_1^2 \frac{1}{x} \, dx$

算 $\int_1^2 \frac{1}{x} \, dx$：
$$\ln x \Big|_1^2 = \ln 2 - \ln 1 = \ln 2 - 0 = \ln 2$$

外层结果：$32 \times \ln 2 = 32 \ln 2$

\textbf{答案：$\boxed{32 \ln 2}$}

\subsection*{📌 终极总结：三重积分除法，还是那三板斧}

\begin{center}
\begin{tabular}{|c|c|c|}
\hline
情况 & 怎么办 & 例子 \\
\hline
\textbf{下面是数字} & 当成常数提到最外面 & $\frac{xyz}{6} \to \frac{1}{6} \times xyz$ \\
\hline
\textbf{上下有同一个字母} & 先约分！能划的先划掉 & $\frac{x^2 yz}{x} \to xyz$ \\
\hline
\textbf{约不掉的 $\frac{1}{x}$、$\frac{1}{y}$、$\frac{1}{z}$} & 请出 VIP 英雄 $\ln$ & $\int \frac{1}{z} dz = \ln z$ \\
\hline
\end{tabular}
\end{center}

\begin{quote}
\textbf{🎯 关于 $\ln$ 的死规矩（最后再强调一遍）}

\begin{center}
\begin{tabular}{|c|c|}
\hline
规矩 & 意思 \\
\hline
$\ln 1 = 0$ & $1$ 的对数永远是 $0$ \\
\hline
$\ln e = 1$ & $e$ 的对数永远是 $1$ \\
\hline
$\int \frac{1}{x} dx = \ln x$ & $\frac{1}{x}$ 的反过来是 $\ln x$ \\
\hline
\end{tabular}
\end{center}

\textbf{⚠️ 最容易犯的错：忘记“上面 − 下面”！}

✅ 正确：$\ln x \Big|_1^e = \ln e - \ln 1 = 1 - 0 = 1$

❌ 错误：$\ln x \Big|_1^e = \ln e = 1$（漏掉了减 $\ln 1$）

\textbf{永远写完整！}
\end{quote}

\subsection*{🎯 最核心的思想}

\textbf{不管是一重、二重、还是三重积分，除法的处理方式完全一样：}

\begin{enumerate}
\item \textbf{能提就提}（下面是数字）
\item \textbf{能约就约}（上下有相同字母）
\item \textbf{约不了就召唤 $\ln$}（$\frac{1}{x}$ 类型）
\end{enumerate}

然后，把除法变成乘法，用“谁不动提谁”的规则，一层一层往外算。

\textbf{一元积分：一层楼。}

\textbf{二重积分：两层楼。}

\textbf{三重积分：三层楼。}

\textbf{但楼梯都是同一种楼梯，你已经会爬了。}

现在不过是多爬一层、再多爬一层而已。

\textbf{同样的招式，用三遍，搞定！}

\section*{一重闭合积分：只不过要多算一些结果而已！}

\bigskip\hrule\bigskip

\subsection*{📌 什么是“闭合”？}

在讲闭合积分之前，先搞清楚一个词：\textbf{闭合}。

想象你在操场上跑步。

\textbf{不闭合}：你从 A 点跑到 B 点，然后停下来。起点和终点不一样。

\textbf{闭合}：你从 A 点出发，绕操场跑一圈，又回到 A 点。起点和终点是同一个地方。

就这么简单。

\textbf{闭合 = 绕一圈回到原点。}

\subsection*{📌 闭合积分长什么样？}

普通的定积分长这样：
$$\int_a^b f(x) \, dx$$

从 $a$ 走到 $b$，一条直线，有头有尾。

闭合积分长这样：
$$\oint f(x) \, dx$$

看到那个积分号上面多了个\textbf{小圈圈}吗？

那个圈就是在告诉你：\textbf{这是绕一圈的积分！}

但是等等，一重积分是在一条线上走的，怎么“绕一圈”？

别急，这里的“闭合”意思稍微不一样。

在一重积分里，\textbf{闭合通常指的是：把一条路径分成几段，每段分别算，最后加起来。}

就像你绕操场一圈，但操场不是圆的，而是由几条直线拼成的。你要一段一段地算。

\subsection*{📌 一个最简单的例子}

假设有一条“路径”，由三段组成：

\textbf{第一段}：从 $x = 0$ 走到 $x = 1$，沿着 $y = 0$（就是 x 轴）

\textbf{第二段}：从 $x = 1$ 走到 $x = 2$，沿着 $y = x - 1$（一条斜线）

\textbf{第三段}：从 $x = 2$ 走回 $x = 0$，沿着 $y = 0$（又回到 x 轴）

你看，这三段拼起来，是不是形成了一个\textbf{封闭的三角形}？

从起点出发，走了一圈，又回到了起点。

现在要算：
$$\oint x \, dx$$

意思是：沿着这个闭合路径，把 $x$ 累加起来。

怎么算？

\textbf{每一段分别算，最后加起来！}

\subsubsection*{第一段：从 $x = 0$ 到 $x = 1$}

$$\int_0^1 x \, dx$$

这就是普通的定积分！

\begin{itemize}
\item $x$ 的反过来是 $\frac{x^2}{2}$
\item 代入 $1$：$\frac{1}{2}$
\item 代入 $0$：$0$
\item 相减：$\frac{1}{2} - 0 = \frac{1}{2}$
\end{itemize}

\textbf{第一段结果：$\frac{1}{2}$}

\subsubsection*{第二段：从 $x = 1$ 到 $x = 2$}

$$\int_1^2 x \, dx$$

\begin{itemize}
\item 代入 $2$：$\frac{4}{2} = 2$
\item 代入 $1$：$\frac{1}{2}$
\item 相减：$2 - \frac{1}{2} = \frac{3}{2}$
\end{itemize}

\textbf{第二段结果：$\frac{3}{2}$}

\subsubsection*{第三段：从 $x = 2$ 走回 $x = 0$}

注意！这一段是\textbf{往回走}的！

从 $2$ 走到 $0$，方向反了。

$$\int_2^0 x \, dx$$

这里有个小技巧：\textbf{调换上下限，加个负号。}

$$\int_2^0 x \, dx = -\int_0^2 x \, dx$$

现在算 $\int_0^2 x \, dx$：
\begin{itemize}
\item 代入 $2$：$2$
\item 代入 $0$：$0$
\item 相减：$2 - 0 = 2$
\end{itemize}

所以：$\int_2^0 x \, dx = -2$

\textbf{第三段结果：$-2$}

\subsubsection*{三段加起来}

$$\frac{1}{2} + \frac{3}{2} + (-2) = \frac{1}{2} + \frac{3}{2} - 2 = 2 - 2 = 0$$

\textbf{答案：$\boxed{0}$}

\subsection*{📌 咦？怎么是 0？}

你可能会惊讶：绕了一圈，结果居然是 $0$？

其实这很正常！

你往前走，积累了一些东西；往回走，又把它们还回去了。

一去一回，正好抵消。

（当然，不是所有闭合积分都是 $0$，这取决于你积的是什么函数。）

\subsection*{📌 再来一个例子：带除法的}

$$\oint \frac{1}{x} \, dx$$

路径：先从 $x = 1$ 走到 $x = e$，再从 $x = e$ 走回 $x = 1$。

\subsubsection*{第一段：从 $x = 1$ 到 $x = e$}

$$\int_1^e \frac{1}{x} \, dx$$

请出 VIP 英雄 $\ln x$：
$$\ln x \Big|_1^e = \ln e - \ln 1 = 1 - 0 = 1$$

\textbf{第一段结果：$1$}

\subsubsection*{第二段：从 $x = e$ 走回 $x = 1$}

$$\int_e^1 \frac{1}{x} \, dx$$

调换上下限，加负号：
$$\int_e^1 \frac{1}{x} \, dx = -\int_1^e \frac{1}{x} \, dx = -1$$

\textbf{第二段结果：$-1$}

\subsubsection*{两段加起来}

$$1 + (-1) = 0$$

\textbf{答案：$\boxed{0}$}

又是 $0$！

因为走出去是 $+1$，走回来是 $-1$，正好抵消。

\subsection*{📌 不是 0 的例子}

别以为闭合积分都是 $0$，来看一个不是 $0$ 的：

路径分两段：
\begin{itemize}
\item 第一段：从 $x = 0$ 到 $x = 2$，被积函数是 $x^2$
\item 第二段：从 $x = 2$ 到 $x = 0$，被积函数变成 $x$（函数变了！）
\end{itemize}

$$\oint = \int_0^2 x^2 \, dx + \int_2^0 x \, dx$$

\subsubsection*{第一段：$\int_0^2 x^2 \, dx$}

\begin{itemize}
\item $x^2$ 的反过来是 $\frac{x^3}{3}$
\item 代入 $2$：$\frac{8}{3}$
\item 代入 $0$：$0$
\item 相减：$\frac{8}{3}$
\end{itemize}

\textbf{第一段结果：$\frac{8}{3}$}

\subsubsection*{第二段：$\int_2^0 x \, dx$}

调换上下限：$-\int_0^2 x \, dx = -2$

\textbf{第二段结果：$-2$}

\subsubsection*{两段加起来}

$$\frac{8}{3} + (-2) = \frac{8}{3} - 2 = \frac{8}{3} - \frac{6}{3} = \frac{2}{3}$$

\textbf{答案：$\boxed{\frac{2}{3}}$}

这次不是 $0$ 了！

因为去的时候是 $x^2$，回来的时候是 $x$，两个函数不一样，所以没法完全抵消。

\subsection*{📌 带除法的复杂例子}

来个综合的：

路径分三段：
\begin{itemize}
\item 第一段：从 $x = 1$ 到 $x = 2$，被积函数是 $\frac{x}{2}$
\item 第二段：从 $x = 2$ 到 $x = 4$，被积函数是 $\frac{1}{x}$
\item 第三段：从 $x = 4$ 到 $x = 1$，被积函数是 $\frac{x^2}{x} = x$
\end{itemize}

\subsubsection*{第一段：$\int_1^2 \frac{x}{2} \, dx$}

$\frac{x}{2} = \frac{1}{2} \times x$，常数提出来：

$$\frac{1}{2} \int_1^2 x \, dx = \frac{1}{2} \times \frac{3}{2} = \frac{3}{4}$$

\textbf{第一段结果：$\frac{3}{4}$}

\subsubsection*{第二段：$\int_2^4 \frac{1}{x} \, dx$}

请出 $\ln x$：
$$\ln x \Big|_2^4 = \ln 4 - \ln 2$$

这里用一个小技巧：$\ln 4 = \ln(2^2) = 2\ln 2$

所以：$\ln 4 - \ln 2 = 2\ln 2 - \ln 2 = \ln 2$

\textbf{第二段结果：$\ln 2$}

\subsubsection*{第三段：$\int_4^1 x \, dx$}

从 $4$ 走回 $1$，调换上下限：
$$\int_4^1 x \, dx = -\int_1^4 x \, dx$$

算 $\int_1^4 x \, dx$：
\begin{itemize}
\item 代入 $4$：$\frac{16}{2} = 8$
\item 代入 $1$：$\frac{1}{2}$
\item 相减：$8 - \frac{1}{2} = \frac{15}{2}$
\end{itemize}

所以：$\int_4^1 x \, dx = -\frac{15}{2}$

\textbf{第三段结果：$-\frac{15}{2}$}

\subsubsection*{三段加起来}

$$\frac{3}{4} + \ln 2 + \left(-\frac{15}{2}\right) = \frac{3}{4} - \frac{15}{2} + \ln 2$$

把分数统一：$\frac{3}{4} - \frac{30}{4} = -\frac{27}{4}$

\textbf{答案：$\boxed{-\frac{27}{4} + \ln 2}$}

（大约等于 $-6.75 + 0.693 = -6.057$）

\subsection*{📌 终极总结：闭合积分的套路}

\begin{center}
\begin{tabular}{|c|c|}
\hline
步骤 & 做什么 \\
\hline
\textbf{第一步} & 把闭合路径拆成几段 \\
\hline
\textbf{第二步} & 每一段单独算（就是普通的定积分！） \\
\hline
\textbf{第三步} & 注意方向！往回走的要调换上下限、加负号 \\
\hline
\textbf{第四步} & 所有段的结果加起来 \\
\hline
\end{tabular}
\end{center}

\begin{quote}
\textbf{⚠️ 最容易犯的错：忘记方向！}

往回走的那一段，上下限是反的！

比如从 $4$ 走到 $1$，写成 $\int_4^1$。

处理方法：$\int_4^1 = -\int_1^4$

\textbf{别忘了那个负号！}
\end{quote}

\subsection*{🎯 最核心的思想}

\textbf{闭合积分没有任何新东西。}

它就是\textbf{多个普通定积分的累加}。

你已经会算定积分了，对吧？
你已经会处理除法了，对吧？

那你就已经会算闭合积分了！

\textbf{唯一的区别是：要分段算，要注意方向，最后加起来。}

就像吃饭一样——
普通积分是吃一碗饭；
闭合积分是吃三碗饭。

饭还是那个饭，只不过多吃几碗而已。

\textbf{同样的招式，多用几遍，搞定！}

\section*{二重闭合积分中的除法：六个面，还是那三板斧}

\bigskip\hrule\bigskip

\subsection*{📌 什么是二重闭合积分？}

普通的二重积分，是在一个\textbf{平面区域}上积分。

二重闭合积分呢？是在一个\textbf{封闭的曲面}上积分。

什么叫“封闭的曲面”？

想象一个纸盒子。

一个长方体的纸盒子，有六个面：上、下、前、后、左、右。

这六个面拼在一起，形成了一个\textbf{完全封闭}的表面——没有缝隙，没有洞，密封的。

这就是“闭合曲面”。

二重闭合积分长这样：
$$\oiint f(x,y,z) \, dS$$

看到那个积分号上面的\textbf{两个圈圈}了吗？

两个圈 = 二重 + 闭合。

意思是：在一个\textbf{封闭的表面}上，把每个位置的 $f$ 值累加起来。

\subsection*{📌 怎么算？拆成几个面！}

想象你要给一个纸盒子刷油漆。

你会怎么做？

\textbf{一个面一个面地刷！}

先刷顶面，再刷底面，再刷前面，再刷后面，再刷左面，再刷右面。

最后把六个面用的油漆量\textbf{加起来}，就是总量。

二重闭合积分也一样：

$$\oiint f \, dS = \underbrace{\int\int_{\text{顶面}}}_{\text{第1块}} + \underbrace{\int\int_{\text{底面}}}_{\text{第2块}} + \underbrace{\int\int_{\text{前面}}}_{\text{第3块}} + \underbrace{\int\int_{\text{后面}}}_{\text{第4块}} + \underbrace{\int\int_{\text{左面}}}_{\text{第5块}} + \underbrace{\int\int_{\text{右面}}}_{\text{第6块}}$$

\textbf{每一块都是一个普通的二重积分！}

\begin{quote}
\textbf{🔔 提醒 ①：每一个面，都只是普通的二重积分！}

“二重闭合积分”听起来吓人，好像是什么全新的怪物。

\textbf{其实不是！}

它就是\textbf{六个普通二重积分的总和}（对于长方体盒子来说）。

每一个面上，你要做的事情跟之前学的二重积分\textbf{一模一样}：
\begin{itemize}
\item 先积一个变量
\item 再积另一个变量
\item 遇到除法？还是三板斧！
\end{itemize}

唯一的新东西是：\textbf{你要算六次，然后加起来}。

量变大了，但质没变。
\end{quote}

\subsection*{📌 每个面上发生了什么？}

一个长方体盒子有六个面。

在每个面上，\textbf{三个变量中有一个是固定的}（因为你站在那个面上，某个方向不动了）。

\begin{center}
\begin{tabular}{|c|c|c|}
\hline
面 & 哪个变量固定 & 剩下哪两个变量在动 \\
\hline
顶面 & $z$ 固定（等于上边界值） & $x$ 和 $y$ 在动 \\
\hline
底面 & $z$ 固定（等于下边界值） & $x$ 和 $y$ 在动 \\
\hline
前面 & $y$ 固定（等于上边界值） & $x$ 和 $z$ 在动 \\
\hline
后面 & $y$ 固定（等于下边界值） & $x$ 和 $z$ 在动 \\
\hline
右面 & $x$ 固定（等于上边界值） & $y$ 和 $z$ 在动 \\
\hline
左面 & $x$ 固定（等于下边界值） & $y$ 和 $z$ 在动 \\
\hline
\end{tabular}
\end{center}

\textbf{一个变量固定，就把它的值代入 $f$ 里面。}

剩下两个变量在动，就对它们做普通的二重积分。

\subsection*{📌 第一个例子：纸老虎（下面是数字）}

计算：
$$\oiint \frac{xyz}{2} \, dS$$

盒子的范围：$0 \leq x \leq 1$，$0 \leq y \leq 2$，$0 \leq z \leq 1$

\subsubsection*{六个面，逐一检查}

\textbf{底面（$z = 0$）：}

把 $z = 0$ 代入 $\frac{xyz}{2}$：
$$\frac{xy \cdot 0}{2} = 0$$

函数直接变成 $0$！

不管 $x$ 和 $y$ 怎么动，$0$ 累加起来还是 $0$。

\textbf{底面结果：$0$} ✅ 不用算！

\textbf{后面（$y = 0$）：}

把 $y = 0$ 代入：$\frac{x \cdot 0 \cdot z}{2} = 0$

\textbf{后面结果：$0$} ✅ 不用算！

\textbf{左面（$x = 0$）：}

把 $x = 0$ 代入：$\frac{0 \cdot yz}{2} = 0$

\textbf{左面结果：$0$} ✅ 不用算！

太好了！六个面里，三个直接是 $0$！

为什么？因为 $\frac{xyz}{2}$ 里面有 $x$、$y$、$z$ 三个字母相乘。只要其中任何一个等于 $0$，整个东西就是 $0$。

而底面（$z=0$）、后面（$y=0$）、左面（$x=0$）恰好都有一个变量等于 $0$。

\textbf{省了一半的工作量！}

\begin{quote}
\textbf{🔔 提醒 ②：动手之前，先花 30 秒检查每个面！}

很多同学一看到六个面，就闷头从第一个面开始算，一直算到第六个面。

\textbf{别急！先把六个面的值都代入看一看！}

你会发现：

\begin{center}
\begin{tabular}{|c|c|}
\hline
情况 & 结果 \\
\hline
代入后函数变成 $0$ & 这个面直接跳过！ \\
\hline
代入后除法消失了 & 变成普通的乘法，更简单！ \\
\hline
代入后除法还在 & 才需要认真用三板斧 \\
\hline
\end{tabular}
\end{center}

\textbf{先检查，再动手！很多面根本不用算，30 秒就能跳过！}
\end{quote}

\subsubsection*{现在来算剩下的三个面}

\textbf{顶面（$z = 1$）：}

把 $z = 1$ 代入 $\frac{xyz}{2}$：
$$\frac{xy \cdot 1}{2} = \frac{xy}{2}$$

需要计算：
$$\int_0^1 \int_0^2 \frac{xy}{2} \, dy \, dx$$

$\frac{1}{2}$ 是常数，提出来：
$$\frac{1}{2} \int_0^1 \int_0^2 xy \, dy \, dx$$

\textbf{内层（对 $y$ 积分）：}

$x$ 不动，提出来：$x \int_0^2 y \, dy$

\begin{itemize}
\item $y$ 的反过来是 $\frac{y^2}{2}$
\item 代入 $2$：$\frac{4}{2} = 2$
\item 代入 $0$：$0$
\item 相减：$2 - 0 = 2$
\end{itemize}

内层结果：$x \times 2 = 2x$

\textbf{外层（对 $x$ 积分）：}

$\frac{1}{2} \int_0^1 2x \, dx = \frac{1}{2} \times 2 \int_0^1 x \, dx = 1 \times \frac{1}{2} = \frac{1}{2}$

\textbf{顶面结果：$\frac{1}{2}$}

\textbf{前面（$y = 2$）：}

把 $y = 2$ 代入 $\frac{xyz}{2}$：
$$\frac{x \cdot 2 \cdot z}{2} = \frac{2xz}{2}$$

等等！$\frac{2xz}{2}$？上面有个 $2$，下面也有个 $2$？

\textbf{约分！} 上下各划掉一个 $2$：
$$\frac{2xz}{2} = xz$$

除法消失了！变成了普通的乘法！

需要计算：
$$\int_0^1 \int_0^1 xz \, dz \, dx$$

\textbf{内层（对 $z$ 积分）：}

$x$ 不动，提出来：$x \int_0^1 z \, dz = x \times \frac{1}{2} = \frac{x}{2}$

\textbf{外层（对 $x$ 积分）：}

$\int_0^1 \frac{x}{2} \, dx = \frac{1}{2} \int_0^1 x \, dx = \frac{1}{2} \times \frac{1}{2} = \frac{1}{4}$

\textbf{前面结果：$\frac{1}{4}$}

\textbf{右面（$x = 1$）：}

把 $x = 1$ 代入 $\frac{xyz}{2}$：
$$\frac{1 \cdot yz}{2} = \frac{yz}{2}$$

需要计算：
$$\int_0^2 \int_0^1 \frac{yz}{2} \, dz \, dy$$

$\frac{1}{2}$ 是常数，提出来：
$$\frac{1}{2} \int_0^2 \int_0^1 yz \, dz \, dy$$

\textbf{内层（对 $z$ 积分）：}

$y$ 不动，提出来：$y \int_0^1 z \, dz = y \times \frac{1}{2} = \frac{y}{2}$

\textbf{外层（对 $y$ 积分）：}

$\frac{1}{2} \int_0^2 \frac{y}{2} \, dy = \frac{1}{2} \times \frac{1}{2} \int_0^2 y \, dy = \frac{1}{4} \times 2 = \frac{1}{2}$

\textbf{右面结果：$\frac{1}{2}$}

\subsubsection*{六个面全部加起来}

\begin{center}
\begin{tabular}{|c|c|c|}
\hline
面 & 结果 & 原因 \\
\hline
底面（$z=0$） & $0$ & 函数变成 $0$ \\
\hline
后面（$y=0$） & $0$ & 函数变成 $0$ \\
\hline
左面（$x=0$） & $0$ & 函数变成 $0$ \\
\hline
顶面（$z=1$） & $\frac{1}{2}$ & 除以 $2$，提出来 \\
\hline
前面（$y=2$） & $\frac{1}{4}$ & $\frac{2}{2}$ 约分，除法消失 \\
\hline
右面（$x=1$） & $\frac{1}{2}$ & 除以 $2$，提出来 \\
\hline
\end{tabular}
\end{center}

$$0 + 0 + 0 + \frac{1}{2} + \frac{1}{4} + \frac{1}{2} = \frac{2}{4} + \frac{1}{4} + \frac{2}{4} = \frac{5}{4}$$

\textbf{答案：$\boxed{\frac{5}{4}}$}

你看到了吗？

六个面看起来很吓人，但：
\begin{itemize}
\item 三个面直接是 $0$（白捡！）
\item 一个面约分后除法消失了
\item 两个面用“提出常数”就搞定了
\end{itemize}

\textbf{全是老招式，没有任何新东西！}

\subsection*{📌 终极魔王：$\frac{1}{x}$ 出现了}

来个有真老虎的：
$$\oiint \frac{yz}{x} \, dS$$

盒子的范围：$1 \leq x \leq e$，$0 \leq y \leq 1$，$0 \leq z \leq 2$

（$e \approx 2.718$，就是那个满足 $\ln e = 1$ 的数。）

\subsubsection*{先检查六个面}

\textbf{底面（$z = 0$）：}

$\frac{y \cdot 0}{x} = 0$

\textbf{底面结果：$0$} ✅ 跳过！

\textbf{后面（$y = 0$）：}

$\frac{0 \cdot z}{x} = 0$

\textbf{后面结果：$0$} ✅ 跳过！

好，两个面直接是 $0$。剩下四个面要算。

\subsubsection*{顶面（$z = 2$）}

把 $z = 2$ 代入 $\frac{yz}{x}$：
$$\frac{y \cdot 2}{x} = \frac{2y}{x}$$

$x$ 还在下面，约不掉。

但可以拆成乘法：$\frac{2y}{x} = 2y \times \frac{1}{x}$

需要计算：
$$\int_1^e \int_0^1 \frac{2y}{x} \, dy \, dx$$

\textbf{内层（对 $y$ 积分）：}

$\frac{2}{x}$ 跟 $y$ 没关系，提出来：
$$\frac{2}{x} \int_0^1 y \, dy = \frac{2}{x} \times \frac{1}{2} = \frac{1}{x}$$

\textbf{外层（对 $x$ 积分）：}

$$\int_1^e \frac{1}{x} \, dx$$

大魔王来了！请出 VIP 英雄 $\ln x$：

$$\ln x \Big|_1^e = \ln e - \ln 1 = 1 - 0 = 1$$

\textbf{顶面结果：$1$}

\subsubsection*{前面（$y = 1$）}

把 $y = 1$ 代入 $\frac{yz}{x}$：
$$\frac{1 \cdot z}{x} = \frac{z}{x}$$

需要计算：
$$\int_1^e \int_0^2 \frac{z}{x} \, dz \, dx$$

\textbf{内层（对 $z$ 积分）：}

$\frac{1}{x}$ 跟 $z$ 没关系，提出来：
$$\frac{1}{x} \int_0^2 z \, dz = \frac{1}{x} \times \frac{z^2}{2}\Big|_0^2 = \frac{1}{x} \times 2 = \frac{2}{x}$$

\textbf{外层（对 $x$ 积分）：}

$$\int_1^e \frac{2}{x} \, dx = 2 \int_1^e \frac{1}{x} \, dx = 2 \times (\ln e - \ln 1) = 2 \times (1 - 0) = 2$$

\textbf{前面结果：$2$}

\subsubsection*{左面（$x = 1$）}

把 $x = 1$ 代入 $\frac{yz}{x}$：
$$\frac{yz}{1} = yz$$

除以 $1$！\textbf{除法直接消失了！}

需要计算：
$$\int_0^1 \int_0^2 yz \, dz \, dy$$

这就是最普通的乘法，你闭着眼都会算。

\textbf{内层（对 $z$ 积分）：}

$y$ 不动，提出来：$y \int_0^2 z \, dz = y \times 2 = 2y$

\textbf{外层（对 $y$ 积分）：}

$\int_0^1 2y \, dy = 2 \times \frac{y^2}{2}\Big|_0^1 = 2 \times \frac{1}{2} = 1$

\textbf{左面结果：$1$}

\subsubsection*{右面（$x = e$）}

把 $x = e$ 代入 $\frac{yz}{x}$：
$$\frac{yz}{e}$$

$e$ 是个数字（大约 $2.718$），是常数！

\textbf{除以数字 $\rightarrow$ 提出来！}

需要计算：
$$\int_0^1 \int_0^2 \frac{yz}{e} \, dz \, dy = \frac{1}{e} \int_0^1 \int_0^2 yz \, dz \, dy$$

等一下……$\int_0^1 \int_0^2 yz \, dz \, dy$ 这个积分，我们刚才在左面已经算过了！结果是 $1$！

所以：$\frac{1}{e} \times 1 = \frac{1}{e}$

\textbf{右面结果：$\frac{1}{e}$}

\subsubsection*{六个面全部加起来}

\begin{center}
\begin{tabular}{|c|c|c|c|}
\hline
面 & 代入后变成 & 用了什么招 & 结果 \\
\hline
底面（$z=0$） & $0$ & 直接跳过 & $0$ \\
\hline
后面（$y=0$） & $0$ & 直接跳过 & $0$ \\
\hline
顶面（$z=2$） & $\frac{2y}{x}$ & 召唤 $\ln x$ & $1$ \\
\hline
前面（$y=1$） & $\frac{z}{x}$ & 召唤 $\ln x$ & $2$ \\
\hline
左面（$x=1$） & $yz$ & 除以 $1$，除法消失 & $1$ \\
\hline
右面（$x=e$） & $\frac{yz}{e}$ & 除以数字，提出来 & $\frac{1}{e}$ \\
\hline
\end{tabular}
\end{center}

$$0 + 0 + 1 + 2 + 1 + \frac{1}{e} = 4 + \frac{1}{e}$$

\textbf{答案：$\boxed{4 + \frac{1}{e}}$}

\subsubsection*{回头看看发生了什么}

六个面里，\textbf{除法的四种情况全出现了}：

\begin{center}
\begin{tabular}{|c|c|c|}
\hline
面 & 除法类型 & 怎么处理的 \\
\hline
底面、后面 & 分子等于 $0$ & $\frac{0}{x} = 0$，直接跳过 \\
\hline
左面 & 除以 $1$ & $\frac{yz}{1} = yz$，除法消失 \\
\hline
右面 & 除以数字 $e$ & 常数提出来 \\
\hline
顶面、前面 & 除以字母 $x$ & 召唤 VIP 英雄 $\ln x$ \\
\hline
\end{tabular}
\end{center}

\textbf{每一种情况，都是你早就学过的老招式！}

\begin{quote}
\textbf{🔔 提醒 ③：做完所有面之后，回头检查四件事！}

\textbf{检查一：六个面都算了吗？}

长方体有六个面。数一数，是不是每个面都有结果（哪怕结果是 $0$）？

少算一个面，整道题全错！

\bigskip

\textbf{检查二：代入的值对不对？}

顶面代入的是 $z$ 的\textbf{上}边界，底面代入的是 $z$ 的\textbf{下}边界。

别搞反了！

\bigskip

\textbf{检查三：$\ln$ 有没有写完整？}

✅ 正确：$\ln x \Big|_1^e = \ln e - \ln 1 = 1 - 0 = 1$

❌ 错误：$\ln x \Big|_1^e = \ln e = 1$（偷懒省掉了减 $\ln 1$）

这次碰巧 $\ln 1 = 0$，减了跟没减一样。

但如果是 $\ln x \Big|_2^e = \ln e - \ln 2 = 1 - \ln 2$，\textbf{不是} $1$！

\textbf{每次都写“上面 − 下面”，一步都不省！}

\bigskip

\textbf{检查四：最后加的时候，正负号对不对？}

六个结果加起来时，仔细看看有没有负数。

别把 $+$ 写成 $-$，或者漏掉某一项。
\end{quote}

\subsection*{📌 终极总结：二重闭合积分的套路}

\begin{center}
\begin{tabular}{|c|c|}
\hline
步骤 & 做什么 \\
\hline
\textbf{第一步} & 搞清楚盒子的六个面，每个面上哪个变量固定 \\
\hline
\textbf{第二步} & 把固定的那个变量的值代入 $f$ \\
\hline
\textbf{第三步} & 先检查！代入后是 $0$ 吗？除法能消掉吗？ \\
\hline
\textbf{第四步} & 每个面单独算（就是普通的二重积分！） \\
\hline
\textbf{第五步} & 遇到除法？还是三板斧！ \\
\hline
\textbf{第六步} & 六个面的结果全部加起来 \\
\hline
\textbf{第七步} & 回头检查：面数、代入值、$\ln$ 的减法、正负号 \\
\hline
\end{tabular}
\end{center}

\subsection*{🎯 最核心的思想}

\textbf{二重闭合积分 = 普通二重积分 × 六个面}

每个面上，一个变量变成常数，剩下两个变量做普通的二重积分。

遇到除法？还是那三板斧：提常数、约分、召唤 $\ln$。

\textbf{闭合只是让你多算几个面，武功还是那套武功。}

你现在已经把所有类型的积分除法都打通关了：

\begin{center}
\begin{tabular}{|c|c|}
\hline
类型 & 本质 \\
\hline
一重积分 & 三板斧用一次 \\
\hline
二重积分 & 三板斧用两次（两层楼） \\
\hline
三重积分 & 三板斧用三次（三层楼） \\
\hline
一重闭合积分 & 拆成几段，每段用三板斧 \\
\hline
二重闭合积分 & 拆成几个面，每个面用三板斧 \\
\hline
\end{tabular}
\end{center}

\textbf{万变不离其宗。同样的积木，不同的搭法而已！}

\chapter*{写在最后}
\addcontentsline{toc}{chapter}{写在最后}

\bigskip
\hrule
\bigskip

嘿。

你真的看到这里了。

我得停下来，认认真真地跟你说几句话。

\bigskip
\hrule
\bigskip

你知道吗，这本书里讲的那些东西——积分、闭合曲面、法向量——很多大学生到了二十岁才第一次见到。有些人见到之后，直接就放弃了。

但你没有。

你从第一页翻到了最后一页。中间那些弯弯绕绕的符号，那些奇奇怪怪的概念，那些让人头疼的公式——你一个一个地啃下来了。

也许有些地方你看了两遍。也许有些地方你画了好几张图才搞明白。也许有些地方你现在还有点模模糊糊的。

\textbf{这些都没关系。}

因为你已经做到了一件大多数人一辈子都不会去做的事情：你主动走进了一扇很难的门，然后没有转身跑掉。

\bigskip
\hrule
\bigskip

说实话，写这本书的时候，我最怕的一件事就是：把你吓跑。

数学这个东西，长得就很吓人。满页的符号，像外星人的语言。我每写一段，都在心里问自己：\textbf{“这样说，能听懂吗？会不会太绕了？要不要再打一个比方？”}

如果你看到了这里，那说明——至少没有完全失败。

这对我来说，比什么都重要。

\bigskip
\hrule
\bigskip

我想谢谢你的一样东西，不是你的聪明，而是你的\textbf{耐心}。

聪明这个东西，说真的，没那么稀罕。这个世界上聪明人很多。

但愿意对着一本讲积分的书，一页一页、一步一步、老老实实地看下去的人——真的不多。

你的耐心，比聪明值钱一百倍。

\bigskip
\hrule
\bigskip

最后，我想跟你说一个秘密。

这本书讲完的内容，其实只是数学世界的一个小角落。外面还有大得多、美得多、也难得多的东西等着你。

微分方程、线性代数、拓扑、群论……每一个名字背后都藏着一整个宇宙。

但你不用急。

因为你已经证明了一件事：\textbf{你有能力推开那些门。}

至于什么时候推、推哪一扇——随你。

\bigskip
\hrule
\bigskip

谢谢你读完这本书。

不是客套，是真的谢谢。

能有你这样的读者，是写书的人最大的运气。

\bigskip
\hrule
\bigskip

\textit{后会有期。}

\end{document}